\documentclass[11pt,openany]{book}
\setcounter{secnumdepth}{3}
\usepackage[margin=1.1in]{geometry}
% Chinese edition: build with XeLaTeX (see the Makefile)
\usepackage[UTF8,heading=true]{ctex}
\usepackage{amsmath,amssymb,amsthm}
\usepackage{graphicx}
\usepackage{tikz}
\usetikzlibrary{arrows.meta,decorations.markings,decorations.pathmorphing,decorations.pathreplacing,calc}
\input{figures/icons}
\usepackage[unicode,colorlinks=true,linkcolor=blue!60!black,citecolor=blue!60!black]{hyperref}
\usepackage{microtype}
\setlength{\emergencystretch}{3em} % long names (licences, references) in narrow lines
\usepackage{empheq}
\usepackage{array}
\usepackage{booktabs}
\usepackage{multirow}
\usepackage{longtable}
\usepackage{circuitikz}
\definecolor{keyyellow}{rgb}{1.0,0.97,0.78}
% key equations: \begin{empheq}[box=\keybox]{equation} ... \end{empheq} -- light-yellow box, number stays outside
\newcommand{\keybox}[1]{{\setlength{\fboxsep}{7pt}\colorbox{keyyellow}{#1}}}
\usepackage[most]{tcolorbox}
% key statements (propositions etc.): the same light-yellow background, breakable across pages
\newtcolorbox{keyblock}{enhanced,breakable,colback=keyyellow,colframe=keyyellow,boxrule=0pt,arc=0pt,left=6pt,right=6pt,top=2pt,bottom=2pt,boxsep=0pt}
\renewcommand{\chaptermark}[1]{\markboth{\small 第\thechapter 章\quad #1}{}}
\renewcommand{\sectionmark}[1]{\markright{\small\thesection.\ #1}}
\renewcommand{\topfraction}{0.95}
\renewcommand{\textfraction}{0.05}
\renewcommand{\floatpagefraction}{0.75}

% part pages: title at the top third, and text may follow on the same page (used for the part introductions)
\makeatletter
\renewcommand\part{\if@openright\cleardoublepage\else\clearpage\fi\thispagestyle{plain}\if@twocolumn\onecolumn\@tempswatrue\else\@tempswafalse\fi\null\vspace{0.18\textheight}\secdef\@part\@spart}
\renewcommand\@endpart{\par\vspace{3em}\if@tempswa\twocolumn\fi}
\makeatother
\newtheorem{theorem}{定理}[chapter]
\newtheorem{proposition}{命题}[chapter]
\newtheorem{lemma}{引理}[chapter]
\theoremstyle{remark}
\newtheorem{remark}{注}[chapter]
% exercises: \begin{exercise}[\lvlB] ... \end{exercise}, numbered "Exercise 4.3"; difficulty marks
\theoremstyle{definition}
\newtheorem{exercise}{习题}[chapter]
\newcommand{\lvlA}{$\star$}
\newcommand{\lvlB}{$\star\star$}
\newcommand{\lvlC}{$\star\star\star$}
% answer to an exercise in the Answers appendix: \answer{ex:NN:k}{text}
\newcommand{\answer}[2]{\par\noindent\textbf{\ref{#1}.}~#2\par\smallskip}
\newcommand{\dd}{\mathrm{d}}
% neuron type code: first four letters of the primary mediator
\newcommand{\nt}[1]{\textsf{\textbf{#1}}}
\newcommand{\mV}{\,\text{mV}}
\newcommand{\ms}{\,\text{ms}}
\newcommand{\mM}{\,\text{mM}}
\newcommand{\um}{\,\text{\textmu m}}
\newcommand{\ion}[2]{\mathrm{#1}^{#2}}
\newcommand{\Na}{\ion{Na}{+}}
\newcommand{\K}{\ion{K}{+}}
\newcommand{\Cl}{\ion{Cl}{-}}
\newcommand{\Ca}{\ion{Ca}{2+}}
\newcommand{\conc}[1]{[\mathrm{#1}]}

\title{\Huge 20 瓦的计算机\\[16pt] \Large 大脑如何计算：面向工程师的定量方法}
\hypersetup{pdftitle={20 瓦的计算机. 大脑如何计算：面向工程师的定量方法},pdfauthor={Konstantin Kladko}}
\author{\Large 康斯坦丁·克拉季科（Konstantin Kladko）}
\date{}

\begin{document}
\frontmatter
\maketitle
\thispagestyle{empty}
\vspace*{\fill}
\noindent{\small \copyright~2026 Konstantin Kladko（康斯坦丁·克拉季科）。\\[4pt]
本书的正文、插图、习题和答案采用知识共享“署名—非商业性使用—相同方式共享 4.0 国际”许可协议（CC BY-NC-SA 4.0）发布：\url{https://creativecommons.org/licenses/by-nc-sa/4.0/deed.zh-hans}。在注明作者并以相同许可协议发布衍生作品的前提下，可以为非商业目的自由复制、用于教学、翻译和改编本书。\\[4pt]
模型程序和构建脚本采用 MIT 许可协议。\\[4pt]
源文件、模型和新版本：\url{https://github.com/kladkogex/20-watt-computer}。}
\clearpage

\tableofcontents
\chapter*{如何阅读本书}
\addcontentsline{toc}{chapter}{如何阅读本书}

本书不必按顺序通读。第1--4章是共同的基础：神经元有哪些类型，\nt{GLUT}和\nt{GABA}计算什么，它们用什么语言交流。之后全书分成几条支线。图~\ref{fig:roadmap} 的地图标出了各章之间的依赖关系：从第 $A$ 章指向第 $B$ 章的箭头表示第 $B$ 章用到了第 $A$ 章的概念和公式。地图只画出主要的依赖；零星的回溯引用不影响阅读。

\begin{figure}[h]
\centering
\begin{tikzpicture}[>=Stealth,scale=0.95,
  ch/.style={draw,thick,rounded corners=3pt,align=center,font=\scriptsize,text width=1.85cm,minimum height=0.62cm,inner sep=2pt},
  base/.style={ch,fill=blue!8}, bus/.style={ch,fill=orange!14}, sum/.style={ch,fill=gray!18},
  io/.style={ch,fill=green!10}, sort/.style={ch,fill=cyan!8}, lab/.style={ch,fill=violet!10},
  dep/.style={->,thick,gray!60!black}]
% foundation
\node[base] (c1) at (0,0) {1神经元类型};
\node[base] (c2) at (2.3,0) {2 \nt{GLUT}};
\node[base] (c3) at (4.6,0) {3 \nt{GLUT}与\nt{GABA}};
\node[base] (c4) at (6.9,0) {4摩尔斯电码};
% broadcast channels and the synthesis
\node[bus] (c5) at (4.6,-1.5) {5 \nt{SERO}};
\node[bus] (c6) at (7.3,-1.5) {6 \nt{DOPA}};
\node[bus] (c7) at (9.2,-0.9) {7多巴胺理论};
\node[bus] (c8) at (9.2,-2.1) {8 \nt{NORA}};
\node[bus] (c9) at (11.5,-2.1) {9 \nt{ACET}};
\node[sum] (c10) at (13.8,-0.9) {10整台机器};
% inputs and outputs
\node[io] (c12) at (2.3,-4.1) {12识别};
\node[io] (c11) at (4.6,-4.1) {11输入};
\node[io] (c13) at (6.9,-4.1) {13肌肉};
\node[io] (c14) at (9.2,-3.05) {14疼痛};
\node[io] (c15) at (9.2,-4.1) {15运动学习};
\node[io] (c16) at (9.2,-5.15) {16身体化学};
% lab, sorting, sleep
\node[lab] (c17) at (11.5,-4.1) {17调试};
\node[lab] (c21) at (13.8,-4.1) {21活体机器};
\node[lab] (c22) at (13.8,-5.15) {22 DishBrain};
\node[sort] (c18) at (11.5,-5.15) {18仪器式\\神经元};
\node[sort] (c19) at (13.8,-6.2) {19树突数目};
\node[bus] (c20) at (11.5,-6.2) {20睡眠};
% dependencies
\draw[dep] (c1) -- (c2); \draw[dep] (c2) -- (c3); \draw[dep] (c3) -- (c4);
\draw[dep] (c1) -- (c5); \draw[dep] (c4) -- (c6); \draw[dep] (c5) -- (c6);
\draw[dep] (c6) -- (c7); \draw[dep] (c6) -- (c8); \draw[dep] (c8) -- (c9);
\draw[dep] (c7) -- (c10); \draw[dep] (c9) -- (c10);
\draw[dep] (c4) -- (c11); \draw[dep] (c11) -- (c12); \draw[dep] (c11) -- (c13); \draw[dep] (c5) -- (c13);
\draw[dep] (c13) -- (c14); \draw[dep] (c13) -- (c15); \draw[dep] (c13) -- (c16);
\draw[dep] (c15) -- (c17); \draw[dep] (c9) -- (c17); \draw[dep] (c17) -- (c21); \draw[dep] (c21) -- (c22);
\draw[dep] (c16) -- (c18); \draw[dep] (c18) -- (c19); \draw[dep] (c16) -- (c20);
% legend
\begin{scope}[yshift=-7.25cm]
\foreach \s/\t [count=\i] in {base/基础, bus/通道与工作模式, sum/总结, io/输入与输出, sort/神经元的分类, lab/实验室}{
  \pgfmathtruncatemacro{\col}{mod(\i-1,3)}\pgfmathtruncatemacro{\row}{(\i-1)/3}
  \node[\s,text width=0.3cm,minimum height=0.32cm] at ({\col*4.8+0.2},{-\row*0.55}) {};
  \node[anchor=west,font=\scriptsize] at ({\col*4.8+0.55},{-\row*0.55}) {\t};}
\end{scope}
\end{tikzpicture}
\caption{各章地图。箭头从一章指向依赖它的那一章。章与章之间的其他引用只是指引，不构成依赖。}
\label{fig:roadmap}
\end{figure}

\section*{阅读路线}

\begin{center}
\footnotesize
\setlength{\tabcolsep}{4pt}
\renewcommand{\arraystretch}{1.2}
\begin{tabular}{>{\raggedright}p{3.0cm}>{\raggedright}p{3.9cm}>{\raggedright\arraybackslash}p{6.4cm}}
\toprule
路线 & 适合谁 & 按顺序阅读的章节 \\
\midrule
一学季课程 & 教师，10周 & 第1周：第1--2章；2：第3章；3：第4章；4：第5--6章；5：第7--8章；6：第9章；7：第10章；8：第11--12章；9：第13和15章；10：第17章 \\
人工智能与机器学习 & 熟悉神经网络、想弄懂大脑如何学习的读者 & 1--4，5（略读），6，7，8，9，11（略读），12，13（略读），15 \\
电子与硬件 & 电路设计师、神经形态芯片开发者 & 1--4，5，11，13，16，18，19，17（第9和15章略读） \\
神经接口与活体计算 & 搭建脑机接口和活体神经元芯片的人 & 1，2，4，11，13，17（第9和15章略读），21，22（第12章略读） \\
快速浏览 & 只有一个周末的工程师 & 1，2，3，6，10；第6章对第4--5章的引用可从上下文理解 \\
\bottomrule
\end{tabular}
\end{center}

“略读”的意思是：读完一章的开头、黄色方框中的命题和末尾的总结表就够了。

每章末尾都有习题：估算、推导、建模和设计；简要答案汇集在附录~\ref{sec:answers}。全部符号见附录~\ref{sec:notation}；每章主要命题所依据的实验见附录~\ref{sec:supplement}。书中的数字都由\texttt{models/}目录中的小程序算出，可以运行，也可以修改。

\mainmatter

\chapter{神经元类型}
\label{sec:neurontypes}

\section{把神经元当作黑箱}

假设有人给你一袋八百六十亿个神经元~\cite{Azevedo2009}，让你把它们分成几堆。你会怎么分？

工程师拿到一袋陌生的芯片，不会按封装外形去分。他会问：这个器件有几个输入、几个输出，输出的是什么。我们也这样画神经元——画成电路图上的一个方框（图~\ref{fig:blackbox}）。

\begin{figure}[t]
\centering
\begin{tikzpicture}[>=Stealth,x=1cm,y=1cm]
% the box
\node[draw,very thick,minimum width=2.6cm,minimum height=2.6cm,fill=blue!6] (N) at (0,0) {};
\node at (0,0.55) {\small 神经元};
\node at (0,-0.15) {$\displaystyle u=\sum_i w_i x_i$};
\node at (0,-0.8) {\small $y=\Theta(u-\theta)$};
% inputs
\foreach \k/\y/\lab in {1/1.05/{x_1},2/0.55/{x_2},3/0.05/{x_3}}{
  \draw[->,thick,blue!60!black] (-4.2,\y) -- (-1.3,\y);
  \node[left] at (-4.2,\y) {$\lab$};
  \node[above,blue!60!black] at (-2.1,\y-0.06) {\scriptsize $w_{\k}$};}
\node at (-4.45,-0.4) {$\vdots$};
\node at (-2.1,-0.4) {$\vdots$};
\draw[->,thick,blue!60!black] (-4.2,-1.05) -- (-1.3,-1.05);
\node[left] at (-4.2,-1.05) {$x_n$};
\node[above,blue!60!black] at (-2.1,-1.11) {\scriptsize $w_{n}$};
\draw[decorate,decoration={brace,amplitude=5pt},gray!60!black] (-4.9,-1.2) -- (-4.9,1.2);
\node[left,align=right,gray!40!black] at (-5.1,0) {\small $n\sim10^3$--$10^5$\\[-2pt]\small 个输入};
% single output wire, then fan-out
\draw[very thick,red!70!black] (1.3,0) -- (3.2,0);
\foreach \x in {1.6,2.2,2.34,2.9}{\draw[thick,red!70!black] (\x,0) -- (\x,0.4);}
\node[above right,font=\tiny\itshape,gray!30!black,inner sep=1pt] at (1.42,0.45) {咔嗒\dots\ 咔嗒咔嗒\dots};
\node[below,red!70!black,align=center] at (2.25,-0.05) {\scriptsize 一个输出 $y$};
\fill[red!70!black] (3.2,0) circle (0.06);
\foreach \y in {1.3,0.65,0,-0.65,-1.3}{
  \draw[->,thick,red!70!black] (3.2,0) -- (5.0,\y);}
\draw[decorate,decoration={brace,amplitude=5pt},gray!60!black] (5.3,1.45) -- (5.3,-1.45);
\node[right,align=left,gray!40!black] at (5.5,0) {\small $y$ 的副本\\[-2pt]\small 送往 $10^3$--$10^4$\\[-2pt]\small 个其他神经元};
\end{tikzpicture}
\caption{工程师眼中的神经元。许多输入 $x_i$，各有自己的权重 $w_i$；内部是加权和 $u$ 以及与阈值 $\theta$ 的比较（$\Theta$ 是阶跃函数：自变量为正时取 $1$，否则取 $0$）；一个输出 $y$，即一串脉冲，经分支复制给成千上万个接收者。}
\label{fig:blackbox}
\end{figure}

方框的输出不是一个数，而是一串短促的电压脉冲：“咔嗒”，停顿，“咔嗒咔嗒”，长时间停顿。每个脉冲的高度都一样，所以全部信息都在于脉冲\emph{何时}出现。粗略地说，方框里是一个加法器加一个比较器：输入按权重相加，和一旦越过阈值，方框就发出一个脉冲。下一章我们会换成一个老老实实在连续时间中工作的模型。

输出只有一个，但导线会分叉，每个分支传送的都是信号\emph{完全相同}的副本。电子学里把这叫作扇出（fan-out）。皮层神经元的扇出在几千量级；而芯片里的一个逻辑门通常只驱动几个其他门。

那么输出导线把\emph{什么}送到接收者那里？脉冲跑到每个分支的末端，在那里变成化学信号：分支向细胞间的狭窄缝隙释放少许某种特定物质——\emph{神经递质}，它改变接收者的输入电流。这就是分类的依据：\emph{神经元的输出释放哪种神经递质}。

\section{一个输出，一种信号}

只有当每个神经元的输出只有一种类型时，这种分类才有意义。事实如此吗？实验表明，几乎如此。

\begin{keyblock}
\begin{proposition}[一个神经元，一种输出]
\label{prop:dale}
几乎每个神经元都只有一种主要神经递质，并在其输出的所有分支上释放它。因此，神经元的类型由它的主要神经递质决定~\cite{Strata1999}。
\end{proposition}
\end{keyblock}

我们约定：用\emph{神经元主要递质英文名称的前四个字母}来标记神经元类型。例如，主要递质是谷氨酸（glutamate）的神经元就是\nt{GLUT}神经元。所有类型汇总在表~\ref{tab:types} 中。

\begin{table}[t]
\centering
\footnotesize
\setlength{\tabcolsep}{3pt}
\renewcommand{\arraystretch}{1.15}
\begin{tabular}{>{\raggedright}p{1.0cm}>{\raggedright}p{2.05cm}>{\raggedright}p{1.75cm}>{\raggedright}p{1.6cm}>{\centering}p{0.7cm}>{\raggedright}p{1.55cm}>{\raggedright}p{1.8cm}>{\raggedright\arraybackslash}p{3.2cm}}
\toprule
类型 & 主要递质 & 总线 & 速度 & 符号 & 脑中神经元数 & 每个神经元的输出数 & 在计算中的作用 \\
\midrule
\nt{GLUT} & 谷氨酸 & 寻址 & 快和慢 & $+$ & $\sim8\cdot10^{10}$ & $\sim10^4$ & 主要的正权重 \\
\nt{GABA} & $\gamma$-氨基丁酸 & 寻址 & 快和慢 & $-$ & $\sim3\cdot10^{9}$ & $10^3$--$10^4$ & 主要的负权重 \\
\nt{GLYC} & 甘氨酸 & 寻址 & 快 & $-$ & $10^6$--$10^7$ & $10^2$--$10^3$ & 脊髓和脑干中的负权重；谷氨酸某种受体的第二把钥匙 \\
\nt{ACET} & 乙酰胆碱 & 两者 & 快和慢 & $\pm$ & $\sim10^6$ & 肌肉：$10$--$10^3$；脑：$\sim10^5$ & 通往肌肉的输出；在脑中是学习速率（假说） \\
\nt{DOPA} & 多巴胺 & 广播 & 慢 & $\pm$ & $\sim5\cdot10^{5}$ & $10^5$--$10^6$ & 奖励预测误差 $\delta$ \\
\nt{SERO} & 血清素 & 广播 & 慢（有一种受体是快的） & $\pm$ & $\sim3\cdot10^{5}$ & $\sim10^5$ & 规划视野（假说） \\
\nt{HIST} & 组胺 & 广播 & 慢 & $+$ & $\sim6\cdot10^{4}$ & 无估计 & 全局信号“保持清醒” \\
\nt{NORA} & 去甲肾上腺素 & 广播 & 慢 & $\pm$ & $\sim5\cdot10^{4}$ & $\sim10^5$ & 增益（假说） \\
\nt{ADRE} & 肾上腺素 & 广播 & 慢 & $\pm$ & $\sim10^4$ & 无估计 & 神经元很少；在脑中的作用不明 \\
\nt{ASPA} & 天冬氨酸 & 寻址 & 快 & $+$ & 无估计 & 无估计 & 较弱的正权重；作用有争议 \\
\bottomrule
\end{tabular}
\caption{神经元类型，按脑中神经元数目从多到少排列。\emph{总线}：寻址——递质释放到通往特定接收者的狭窄缝隙中；广播——从一小群源神经元送往大片脑区（\ref{sec:buses}~节）。\emph{速度}：快——经离子型受体，毫秒级；慢——经代谢型受体，数百毫秒乃至更长。\emph{符号}是通常的符号；最终由接收者决定（\ref{sec:receiver}~节），$\pm$ 表示两种符号都有。\emph{脑中神经元数}——整个人脑中这一类型神经元的数目，只给量级；神经元总数约为 $8.6\cdot10^{10}$~\cite{Azevedo2009}。\nt{GLUT}神经元绝大多数（约 $7\cdot10^{10}$）是小脑中细小的输入神经元；\nt{GLUT}和\nt{GABA}的数目由这一计数和皮层中\nt{GABA}神经元的比例得出~\cite{Hendry1987}。在少数类型中，只有\nt{HIST}~\cite{Airaksinen1991} 和\nt{DOPA}神经元两群中的主要一群~\cite{Pakkenberg1991} 做过直接计数，所以\nt{DOPA}的数字是下限；\nt{SERO}的数字是两次局部计数之和~\cite{Baker1990,Baker1991}。\nt{NORA}、\nt{GLYC}、\nt{ACET}和\nt{ADRE}的数字是没有直接计数的粗略量级估计。\emph{每个神经元的输出数}——一个神经元有多少突触末梢，即输出导线分支上释放递质的位点（图~\ref{fig:blackbox}），只给量级。广播类型的末梢没有直接数过：\nt{DOPA}的数字是由一根输出导线的分支覆盖其靶区的比例推算的~\cite{Matsuda2009}，其余类型的估计更粗。\nt{ACET}分两种情况：控制肌肉的神经元，以及脑中的广播神经元。}
\label{tab:types}
\end{table}

% the pie is a table-type float with a figure caption, so it can never float ahead of Table~\ref{tab:types}
\begin{table}[tb]
\makeatletter\def\@captype{figure}\makeatother
\centering
\definecolor{vizglut}{HTML}{2A78D6}
\definecolor{vizgaba}{HTML}{E34948}
\definecolor{vizrest}{HTML}{8A8986}
\begin{tikzpicture}[>=Stealth]
% pie: GLUT 96.4 %, GABA 3.6 % (13.0 degrees), the rest 0.006 % (a hairline at 0 degrees)
\fill[vizglut] (0,0) -- (13:1.9) arc (13:360:1.9) -- cycle;
\fill[vizgaba] (0,0) -- (0:1.9) arc (0:13:1.9) -- cycle;
\draw[white,line width=1pt] (0,0) -- (13:1.9);
\draw[vizrest,line width=1.2pt] (0,0) -- (0:1.9);
\node[white,align=center,font=\small] at (190:0.95) {\nt{GLUT}\\$96.4\,\%$};
\draw[gray!60!black] (6.5:1.75) -- (2.05,2.1);
\node[anchor=west,font=\small] at (2.05,2.1) {\nt{GABA} $3.6\,\%$};
% zoom box for the remaining types
\draw[densely dashed,vizrest] (1.9,0) -- (3.1,1.65);
\draw[densely dashed,vizrest] (1.9,0) -- (3.1,-2.2);
\draw[vizrest,rounded corners=3pt] (3.1,-2.2) rectangle (10.1,1.65);
\node[anchor=west,font=\scriptsize] at (3.2,1.4) {其余所有类型合计：$\approx0.006\,\%$；对数坐标};
\foreach \code/\lg/\y/\val/\lx in {GLYC/6.477/1.0/{$10^6$--$10^7$}/4.05,ACET/6.0/0.6/{$\sim10^6$}/3.0,DOPA/5.699/0.2/{$\sim5\cdot10^5$}/2.699,SERO/5.477/-0.2/{$\sim3\cdot10^5$}/2.477,HIST/4.778/-0.6/{$\sim6\cdot10^4$}/1.778,NORA/4.699/-1.0/{$\sim5\cdot10^4$}/1.699,ADRE/4.0/-1.4/{$\sim10^4$}/1.0}{
  \node[anchor=east,font=\scriptsize] at (4.35,\y) {\nt{\code}};
  \fill[vizrest] (4.45,\y-0.13) rectangle ({4.45+\lg-3},\y+0.13);
  \node[anchor=west,font=\scriptsize] at ({4.5+\lx},\y) {\val};}
\draw[vizrest,thick] (7.45,1.0) -- (8.45,1.0);
\draw[vizrest,thick] (7.45,0.92) -- (7.45,1.08);
\draw[vizrest,thick] (8.45,0.92) -- (8.45,1.08);
\draw[gray!60!black] (4.45,-1.72) -- (8.45,-1.72);
\foreach \e in {3,4,5,6,7}{
  \draw[gray!60!black] ({4.45+\e-3},-1.72) -- ({4.45+\e-3},-1.66);
  \node[below,font=\scriptsize] at ({4.45+\e-3},-1.72) {$10^{\e}$};}
\end{tikzpicture}
\caption{按表~\ref{tab:types} 的估计，各类神经元在脑中所占的比例。圆几乎全被\nt{GLUT}占满；\nt{GABA}是一个窄扇区，其余所有类型合起来只是边界上的一根细线。右侧把这根细线放大，按神经元数的对数坐标画出；\nt{GLYC}的柱子画到 $10^6$--$10^7$ 区间的中点，横线段表示整个区间。\nt{ASPA}没有估计值，图中未画出。}
\label{fig:typepie}
\end{table}

这几堆的大小相差多悬殊，看图~\ref{fig:typepie} 就知道：每一千个神经元中，$964$ 个是\nt{GLUT}，$36$ 个是\nt{GABA}，其余所有类型加起来每十万个神经元中只有约六个。

GABA的缩写本来就是四个字母。那些几乎从不充当主要递质的物质——神经肽、气体、脂类、嘌呤——不单独编码；它们在附录~\ref{sec:othermediators} 中介绍。命题~\ref{prop:dale} 中的“几乎”二字很重要。许多神经元在释放主要递质的同时还释放辅助物质~\cite{Yao2023}。有些还释放第二种快速递质，甚至符号相反：一部分\nt{DOPA}神经元同时释放GABA~\cite{Tritsch2012}。还有些神经元在同一输出的不同分支上释放不同的递质~\cite{Zhang2015}。这些都有测量依据，所以“一个神经元，一种递质”是有例外的规则，而不是定律。但作为第一级划分，它经得起检验：在小鼠脑细胞图谱中，神经元按其表达的基因被分成五千多个群，而大类的划分仍然首先取决于主要递质~\cite{Yao2023}。

这条规则有一个推论，它立刻把大脑和人工神经网络区分开。在人工网络中，同一个神经元可以对一个接收者有正权重、对另一个有负权重：权重 $w_{ij}$ 只是矩阵里的数。在大脑中，作用的符号主要由递质决定，而一个神经元只有一种递质。所以，如果神经元 $j$ 兴奋某一个接收者，它也兴奋所有其他接收者。从神经元 $j$ 出发的整个权重矩阵\emph{列}都是同一符号：
\begin{empheq}[box=\keybox]{equation}
\operatorname{sign} w_{ij} \;=\; s_j \quad\text{对所有接收者 } i, \qquad s_j\in\{+1,-1\}.
\label{eq:dalesign}
\end{empheq}
神经元分为兴奋性（$s_j=+1$）和抑制性（$s_j=-1$），这是神经元本身的属性，而不是连接的属性。如果网络需要从兴奋性神经元出发的负连接，就只能通过中介来实现：兴奋性神经元兴奋一个抑制性神经元，后者再抑制目标。这就是延迟一级的负权重。

已知的神经递质有一百多种，但大脑的几乎全部工作都靠其中五六种，它们分成两个截然不同的类别。

\section{两条总线：寻址总线与广播总线}
\label{sec:buses}

计算机网络中投递消息有两种方式。第一种是\emph{寻址}方式（unicast，单播）：数据包从一个发送者送到特定的接收者，速度快，别人看不到。第二种是\emph{广播}（broadcast）：发送者把同一条消息一次发给网络中的所有节点。神经元两种方式都用，而神经递质恰好按这一标准分成两类（图~\ref{fig:phone_radio}）。

\begin{figure}[t]
\centering
\begin{tikzpicture}[>=Stealth,scale=0.95]
% ---- unicast: point-to-point wiring
\node[above] at (2.0,3.3) {\small \textbf{寻址总线}：谷氨酸和GABA};
\foreach \i in {0,...,4}{
  \fill[blue!60!black] (0.2+0.9*\i,2.5) circle (0.1);
  \fill[blue!60!black] (0.2+0.9*\i,0.6) circle (0.1);}
\foreach \a/\b in {0/0,0/1,1/1,1/3,2/2,3/2,3/4,4/4,2/0}{
  \draw[->,blue!60!black] (0.2+0.9*\a,2.38) -- (0.2+0.9*\b,0.72);}
\fill[red!70!black] (1.1,1.55) circle (0.09);
\pic[scale=1.2] at (1.75,1.9) {letter};
\draw[->,red!70!black,thick] (1.1,1.55) -- (0.28,0.7);
\draw[->,red!70!black,thick] (1.1,1.55) -- (1.95,0.72);
\node[align=center,below] at (2.0,0.2) {\small 数十亿个神经元，\\[-2pt]\small 各有各的收件人，\\[-2pt]\small 时间：毫秒};
% ---- broadcast: one small source, global targets
\begin{scope}[xshift=7.2cm]
\node[above] at (1.8,3.3) {\small \textbf{广播总线}：多巴胺，\dots};
\foreach \i in {0,...,8}{\fill[blue!60!black] (-0.2+0.5*\i,2.75) circle (0.08);}
\fill[orange!85!black] (1.8,0.35) ellipse (0.35 and 0.18);
\pic[scale=1.5] at (2.7,0.75) {mast};
\foreach \x in {-0.2,0.3,0.8,1.3,1.8,2.3,2.8,3.3,3.8}{
  \draw[->,orange!85!black] (1.8,0.5) .. controls (1.8,1.3) and (\x,1.6) .. (\x,2.62);}
\node[align=center,below] at (1.8,0.1) {\small 数千至数十万个神经元，\\[-2pt]\small 同一条消息发给所有人，\\[-2pt]\small 时间：数百毫秒乃至更长};
\end{scope}
\end{tikzpicture}
\caption{两种传输方式。左：寻址总线，像信封上写着地址的邮件；红色标出一个神经元及其两个接收者。右：广播总线，像电台：一小群源神经元，所有人收到同一条消息。}
\label{fig:phone_radio}
\end{figure}

\emph{寻址总线}就是谷氨酸和GABA。绝大多数神经元释放的都是它们。每个输出通向特定的细胞，递质只在接触处的狭窄缝隙内起作用，而且作用很快：一毫秒内打开接收者的离子通道，几毫秒内一切结束。这是\emph{数据}总线：大脑计算的内容就在这条总线上传送。

\emph{广播总线}是多巴胺、血清素、去甲肾上腺素，部分也包括乙酰胆碱。释放它们的是极小的神经元群，但每个神经元的输出分支都铺得极广。递质往往不释放到狭窄缝隙里，而是释放到细胞间隙中，在那里扩散开，作用于附近所有细胞。作用很慢：它不直接打开通道，而是在接收者内部启动一串化学反应。这条总线传送的不是数据，而是大片网络共用的\emph{控制参数}，它们改变网络的工作模式：增益、学习速率、“刚才做得好”的信号。因此这类递质被称为\emph{调质}（modulator）。长期以来人们认为，每个这样的通道就是全脑一个数。现代测量修正了这一图景：通道不是一根导线，而是一小束并行线路。同一种递质的源神经元分成若干子系统，各有自己的接收者和自己的消息，而在局部释放多少递质，还要由局部信号来调节~\cite{Ren2018,Poe2020}。这类线路的数目是几条或几十条，而不是几十亿条，所以通道仍然是广播式的。

如果本章你只想记住一幅图，那就是这一幅：大脑是一张以寻址方式传送数据的巨大网络，上方悬着几条传送控制信号的广播通道。程序员会认出熟悉的结构：计算，加上一小组控制变量，每个变量由一大片网络共用。

\section{\nt{GLUT}：正权重}

谷氨酸是大脑中主要的兴奋性递质。当一个神经元的输出释放谷氨酸时，接收者的通道打开，钠离子流入，接收者的膜电压升高，它发出自己的脉冲的机会增大。用图~\ref{fig:blackbox} 的语言说，这是一个\emph{正}权重的输入，$w_i>0$。

谁释放谷氨酸？几乎所有远距离传送数据的神经元：皮层中四分之三到五分之四的神经元，以及连接不同脑区的大多数神经元。大脑的网络首先是一张谷氨酸连接的网络。

一个工程细节。释放之后，缝隙中的谷氨酸浓度猛增数千倍，达到约一毫摩尔，然后以约一毫秒的时间常数下降~\cite{Clements1992}：谷氨酸从缝隙中扩散出去，专门的转运蛋白泵再把它回收进细胞。为什么要这样？道理和通信线路中每个符号之后把信号归零一样。如果不清除递质，下一个脉冲就会落在“占线”的线路上，间隔几毫秒的脉冲就会粘连在一起。

\section{\nt{GABA}：负权重}

设想一个所有权重都为正的网络。如果平均每个脉冲引发不止一个后续脉冲，活动就会指数增长，几步之后席卷整个网络。电子工程师会认出这是增益大于一的正反馈环：电路进入饱和。为避免这种情况，需要负权重。负权重的主要来源是 $\gamma$-氨基丁酸，缩写为GABA。

GABA在接收者上打开允许氯离子通过的通道。氯的平衡电位位于静息电位附近或略低，因此打开的氯通道把膜电压钳在静息附近，不让它升到阈值。这是一个\emph{负}权重的输入，$w_i<0$。皮层中五分之一到四分之一的神经元是GABA能神经元~\cite{Hendry1987}。它们的输出很短，抑制的是最近的邻居。

大脑中的抑制远不止做减法。它是负反馈，把整个网络维持在工作点上。一般认为，在正常工作的皮层中，到达神经元的兴奋性电流和抑制性电流几乎精确地相互抵消：这种工作状态由理论预言~\cite{vanVreeswijk1996}，也在活体皮层的电流测量中观察到~\cite{Haider2006}，神经元始终停在阈值附近。在不稳定点附近用强负反馈稳定下来的电路，对小信号反应既快又灵敏——这是一个古老的工程手法。

\section{\nt{DOPA}：误差信号}

第一条广播通道是多巴胺。人的多巴胺神经元约有五十万个，大约每十七万个神经元中有一个；这是在它们两群中的主要一群里数出来的~\cite{Pakkenberg1991}，所以是下限。它们的输出通向负责选择行动的脑区。

这条通道传送什么？答案来自这样的实验：在动物获得奖励时记录多巴胺神经元的脉冲~\cite{Schultz1997}。不时出其不意地给动物一滴果汁，多巴胺神经元就以一个短的簇放电作出反应。接着在给果汁之前先点亮一盏灯，总是提前一秒。动物学会这一关联后，神经元改为对\emph{灯}发出簇放电，而对准时到来的果汁本身完全没有反应。如果灯亮了却不给果汁，那么在果汁本该到来的时刻，神经元\emph{安静下来}，放电频率降到平常水平以下。

能解释这三种情况的规则是（图~\ref{fig:rpe}）：神经元报告的不是结果有多好，而是结果\emph{比预期好多少或差多少}。

\begin{figure}[t]
\centering
\begin{tikzpicture}[>=Stealth,x=3.4cm,y=1cm]
\foreach \y/\lab in {2.8/{意外的果汁},1.4/{灯后的果汁},0/{灯亮但没有果汁}}{
  \draw[gray!70!black] (0,\y) -- (3,\y);
  \draw[densely dotted,gray!60!black] (0,\y+0.3) -- (3,\y+0.3);
  \node[left,align=right,font=\small] at (-0.03,\y+0.35) {\lab};}
% row 1: juice without a cue -> burst at the juice
\draw[very thick,orange!85!black] plot[domain=0:3,samples=120] (\x,{2.8+0.3+0.75*exp(-((\x-2.08)/0.07)^2)});
\pic[scale=0.8] at (2,2.58) {juice};
% row 2: cue then juice -> burst at the cue, nothing at the juice
\draw[very thick,orange!85!black] plot[domain=0:3,samples=120] (\x,{1.4+0.3+0.75*exp(-((\x-1.08)/0.07)^2)});
\pic[scale=0.85] at (1,1.15) {lamp}; \pic[scale=0.85] at (2,1.15) {juice};
% row 3: cue, juice omitted -> burst at the cue, dip when the juice was due
\draw[very thick,orange!85!black] plot[domain=0:3,samples=120] (\x,{0.3+0.75*exp(-((\x-1.08)/0.07)^2)-0.28*exp(-((\x-2.1)/0.12)^2)});
\pic[scale=0.85] at (1,-0.25) {lamp}; \pic[scale=0.85] at (2,-0.25) {nojuice};
\node[right,font=\scriptsize] at (2.36,0.1) {下陷};
\node[right,font=\scriptsize] at (2.13,3.75) {惊喜！};
\node[left,font=\scriptsize] at (0.97,2.25) {对灯的爆发};
\node[above,font=\scriptsize] at (2,1.75) {没有反应};
\draw[->,gray!70!black] (0,-0.75) -- (3.05,-0.75) node[right,black,font=\small] {$t$};
\draw[gray!60!black] (1,-0.8) -- (1,-0.7) node[below=2pt,font=\scriptsize] {灯，$1\,\text{s}$};
\draw[gray!60!black] (2,-0.8) -- (2,-0.7) node[below=2pt,font=\scriptsize] {果汁，$2\,\text{s}$};
\end{tikzpicture}
\caption{三个实验中\nt{DOPA}神经元的放电频率（示意图，据~\cite{Schultz1997}）。虚线是平稳的背景频率。灯是提示信号，液滴是果汁，打叉的液滴是许诺了却没给的果汁。响应就是预测误差~\eqref{eq:rpe}。}
\label{fig:rpe}
\end{figure}

熟悉强化学习的程序员~\cite{SuttonBarto2018}会立刻认出这个量。学习选择行动的智能体保存一个估计 $V(s)$：处于状态 $s$ 时预期能得到多少奖励。每走一步，它就把预期和实际所得进行比较：
\begin{empheq}[box=\keybox]{equation}
\delta_t \;=\; r_t \;+\; \gamma\,V(s_{t+1}) \;-\; V(s_t) ,
\label{eq:rpe}
\end{empheq}
并修正估计：$V(s_t)\leftarrow V(s_t)+\alpha\,\delta_t$。这里 $r_t$ 是得到的奖励，$\gamma<1$ 是折扣因子（未来的奖励比眼前的奖励便宜多少），$\alpha$ 是学习速率。量 $\delta_t$ 称为\emph{时间差分误差}。

它的行为和多巴胺神经元完全一致。意外的果汁：$V$ 还是零，$r>0$，$\delta>0$。学会之后的果汁：奖励已由灯预告，果汁之前的 $V(s_t)$ 已等于 $r$，于是 $\delta=0$。果汁没来：$V(s_t)=r$，而 $r_t=0$，$\delta<0$。爆发之所以移到灯上，是因为恰恰在灯亮的那一刻，预期突然跃升：$\gamma V(s_{t+1})-V(s_t)>0$。这里的步 $t,t+1,\dots$ 只是算法上的约定：生物体里并没有这样的节拍，同一个量在连续时间中的形式将在第\ref{sec:dofa}章得到。

所以，多巴胺就是把标量误差信号 $\delta$ 广播给所有需要据此学习的对象。一级近似下，它是全脑一根导线，每一时刻传一个数；这根导线实际上在多大程度上是一束线，将在第\ref{sec:dofaalt}章讨论。

\section{其余的广播通道}

对其余调质，只有一些可用的假说，把它们的作用翻译成同一种全局参数的语言~\cite{Doya2002}。请把它们当作假说看待。

\emph{\nt{NORA}，去甲肾上腺素：增益。}去甲肾上腺素神经元比多巴胺神经元还少，粗略估计有几万个，而它们的输出几乎遍布全脑。每当发生意外或重要的事情，它们就骤然激活。去甲肾上腺素改变接收神经元传递特性的陡度：弱输入变得更弱，强输入变得更强。测量的结果是：接收者的背景脉冲比它对输入的响应被抑制得更厉害，信噪比上升~\cite{Foote1975NE}。在简单模型中用一个数来描述：如果神经元对输入电流 $u$ 的响应频率为 $f=F\bigl(g\cdot u\bigr)$，那么去甲肾上腺素增大系数 $g$~\cite{AstonJones2005}（详见第\ref{sec:nora}章）。$g$ 大时网络工作得“反差鲜明”，决策更快；$g$ 小时则“模糊”，更多地尝试各种可能，就像处于探索模式的强化学习智能体。

\emph{\nt{ACET}，乙酰胆碱：学习速率。}乙酰胆碱控制网络在多大程度上听取当前输入，又在多大程度上依赖自身存储的数据。乙酰胆碱水平高时，来自感官的输入被加强，内部“回忆”连接被削弱，同时权重也更容易改变~\cite{Hasselmo2006}。粗略地说，它是一个“写入/读取”开关，也连带决定学习速率 $\alpha$。

\emph{\nt{SERO}，血清素：规划视野。}血清素传送什么，人们了解得最少。一种假说把它与~\eqref{eq:rpe} 中的折扣因子 $\gamma$ 联系起来：血清素水平高，对应于愿意等待大的奖励，而不是马上抓住小的奖励。血清素神经元放电平稳而缓慢，清醒时每秒几个脉冲，在睡眠的某一阶段几乎完全沉默~\cite{JacobsAzmitia1992}。

六种递质汇总在表~\ref{tab:transmitters} 中。

\begin{table}[t]
\centering
\small
\begin{tabular}{>{\raggedright}p{2.4cm}>{\raggedright}p{3.0cm}>{\raggedright}p{2.0cm}>{\raggedright\arraybackslash}p{5.2cm}}
\toprule
神经元类型 & 神经元比例 & 总线 & 在计算中的作用 \\
\midrule
\nt{GLUT}（谷氨酸） & 大多数 & 寻址 & 正权重，$w>0$ \\
\nt{GABA}（GABA） & 皮层中20--25\% & 寻址 & 负权重，$w<0$；稳定 \\
\nt{DOPA}（多巴胺） & $\sim 6\cdot10^{-6}$ & 广播 & 预测误差 $\delta$ \\
\nt{NORA}（去甲肾上腺素） & $\sim 6\cdot10^{-7}$ & 广播 & 增益 $g$（假说） \\
\nt{ACET}（乙酰胆碱） & 很小 & 两者 & 学习速率 $\alpha$；“写入/读取”（假说） \\
\nt{SERO}（血清素） & $\sim 3\cdot10^{-6}$ & 广播 & 折扣 $\gamma$（假说） \\
\bottomrule
\end{tabular}
\caption{用计算语言描述的大脑主要神经递质。右列是大幅简化：对谷氨酸、GABA和多巴胺是可靠的，对其余递质只是假说。}
\label{tab:transmitters}
\end{table}

\section{符号由接收者决定}
\label{sec:receiver}

前面说“谷氨酸是正权重，GABA是负权重”，我稍微骗了你们一下。没有哪种分子本身带着符号。分子只是一把钥匙。它被捕获之后会发生什么，由\emph{锁}决定，也就是接收细胞上的受体。

最好的例子是多巴胺。它的受体有两个家族~\cite{Beaulieu2011}。在一类受体上，它使细胞更容易兴奋；在另一类上，则使细胞更难兴奋。在负责选择行动的脑区，一些神经元带第一类受体，另一些带第二类，于是同一个多巴胺信号同时加强一群神经元、削弱另一群。这就像同一个广播数据包，被网络中不同节点作出不同的解读。

\begin{keyblock}
\begin{proposition}[消息的含义由接收者决定]
\label{prop:receptor}
神经递质作用的符号和速度不是由递质本身决定的，而是由与之结合的受体决定的。受体分两类。\emph{离子型}受体本身就是离子通道，在不到一毫秒内打开：这是对电流的直接控制，就像晶体管的栅极。\emph{代谢型}受体在细胞内启动一串化学反应，作用需要数百毫秒乃至更长：这更像是写入一个设置寄存器，改变细胞此后的工作方式。寻址总线主要通过前者传话，广播总线主要通过后者。
\end{proposition}
\end{keyblock}

那为什么还能说“谷氨酸起兴奋作用”？因为实际遇到的谷氨酸受体几乎都是兴奋性的，而GABA受体几乎都是抑制性的。这不是自然定律，而是一条非常可靠的惯例，规则~\eqref{eq:dalesign} 正是建立在它的基础上。

\section{要点}

绝大多数神经元构成寻址数据总线，每个神经元的权重符号相同；极少数神经元构成广播通道，向大片脑区传送少数几个共用的控制信号。每十万个神经元中大约只有两个属于后者，而其余整个网络如何学习、在何种模式下工作，却取决于它们。对工程师来说这并不奇怪：任何计算机中，控制寄存器都比存储单元少得不可比拟，可是没有它们，机器就无法运转。

下一章我们来检验：一个只由\nt{GLUT}神经元（数量最多的类型）组成、不依赖任何全局节拍的网络，能计算什么。

\section{习题}

\begin{exercise}[\lvlA]
\label{ex:01:1}
\emph{几堆与几根线。}根据表~\ref{tab:types}（取区间在对数坐标上的中点；\nt{ACET}取 $10^5$ 个输出）：(a)~如果一个神经元是一个人，\nt{GLUT}神经元相当于几个地球人口？全部广播神经元相当于多大的城市？(b)~\nt{DOPA}、\nt{SERO}、\nt{NORA}、\nt{ACET}在末梢中所占的比例是它们在神经元中所占比例的多少倍？
\end{exercise}

\begin{exercise}[\lvlA]
\label{ex:01:2}
\emph{归零。}缝隙中的谷氨酸按 $c_0e^{-t/\tau_c}$ 下降，$c_0=1\mM$，$\tau_c=1\ms$；接收者能察觉高于 $10\,\mu$M的浓度。(a)~一次释放后线路“占线”多长时间？(b)~转运泵被部分阻断，$\tau_c=10\ms$。释放频率最高到多少，线路仍来得及空出来？
\end{exercise}

\begin{exercise}[\lvlB]
\label{ex:01:3}
\emph{爆发移到哪里。}在图~\ref{fig:rpe} 的实验中，灯亮之后依次经过状态 $s_0\to\dots\to s_{K-1}$，步长为 $\Delta$，$T=K\Delta$；奖励 $r$ 在最后一次转移时到来，灯亮的时刻不可预测。求学会后的 $V(s_k)$ 以及对灯的响应 $\delta$。令 $\gamma=e^{-\Delta/\tau_\gamma}$，求 $T=1$~s、$\Delta=0.1$~s、$\gamma=0.95$ 时的 $\tau_\gamma$。
\end{exercise}

\begin{exercise}[\lvlB]
\label{ex:01:4}
\emph{建模：由误差学习。}对习题~\ref{ex:01:3} 的状态链进行模拟，$K=10$，$\gamma=0.95$，$r=1$，规则为 $V(s_t)\leftarrow V(s_t)+\alpha\delta_t$。当 $\alpha=0.05$、$0.1$、$0.2$、$0.5$ 时，对灯的响应在第几次试验 $n^*$ 首次超过对果汁的响应？$n^*$ 如何依赖于 $\alpha$？
\end{exercise}

\begin{exercise}[\lvlB]
\label{ex:01:5}
\emph{设计：分发一个数。}要把数 $\delta$ 送到全部 $10^{10}$ 个神经元；每个接收者（包括中继神经元）都必须从上一层得到 $10$ 个末梢，每一级增加 $1\ms$。当每个神经元有 $10^4$ 个输出（\nt{GLUT}）和 $10^{5.5}$ 个输出（\nt{DOPA}）时，最经济的中继树需要多少神经元、多少级？
\end{exercise}

\begin{exercise}[\lvlB]
\label{ex:01:6}
\emph{平衡为何灵敏。}网络中 $80\,\%$ 是\nt{GLUT}，$20\,\%$ 是\nt{GABA}，全连接，权重为 $J_E/N$ 和 $-J_I/N$；$\tau\dot r=-r+Wr+h$。当 $J_E=10$ 时，$W$ 唯一的非零特征值为 $0.5$。求抑制电流与兴奋电流之比，以及抑制只要减弱百分之几网络就会进入雪崩。
\end{exercise}

\begin{exercise}[\lvlC]
\label{ex:01:7}
\emph{研究：\nt{SERO}就是 $\gamma$ 吗？}由习题~\ref{ex:01:3}，\nt{DOPA}神经元对灯的响应等于 $re^{-T/\tau_\gamma}$。延迟 $T=1,\dots,5$~s，每个延迟呈现 $n$ 次，$\ln\delta$ 的测量散布为 $0.5$。要区分 $\tau_\gamma=2$~s和 $2.4$~s（显著性水平 $0.05$，检验功效 $0.8$），需要多大的 $n$？除了 $\gamma$，还有什么机制会给出同样的随 $T$ 衰减？
\end{exercise}

\chapter{\nt{GLUT}神经元计算什么}
\label{sec:glutcomputer}

\section{游戏规则}

在第\ref{sec:neurontypes}章中我们看到，大脑中绝大多数神经元是\nt{GLUT}：它们只兴奋，权重只取正值。取任意多个这样的神经元，然后问：如果只允许使用活体中真实存在的东西，这样的网络能计算什么？

活体中没有让所有元件同时翻转的时钟发生器。每个神经元各自在连续时间中工作：脉冲随时到来、随时离开，延迟各不相同。有\emph{输入}——对光、声或压力起反应的感觉神经元，也有\emph{输出}——控制肌肉的神经元。两者之间是网络，没有谁告诉它什么时候开始计算、什么时候结束。对电子工程师来说，这是一个\emph{异步}电路，元件的延迟事先并不精确知道。

神经元模型取在连续时间中老老实实工作的模型里最简单的一种。静息时神经元膜电压 $V$ 为零。来自神经元 $j$ 的每个脉冲使它跳升 $w_j\ge0$，之后电压以时间常数 $\tau$ 自行回落到零：
\begin{empheq}[box=\keybox]{equation}
\tau\,\frac{\dd V}{\dd t} \;=\; -V \;+\; \tau\sum_j w_j \sum_k \delta\bigl(t-t_j^{(k)}-d_j\bigr), \qquad w_j\ge 0 .
\label{eq:glutneuron}
\end{empheq}
其中 $t_j^{(k)}$ 是神经元 $j$ 发放脉冲的时刻，$d_j$ 是从它出发的路径延迟，$\delta$ 是狄拉克 $\delta$ 函数，表示瞬时跳变。当 $V$ 达到阈值 $\theta$ 时，神经元发出一个脉冲，$V$ 复位为零，此后约一毫秒内神经元不能再次发放。典型数值：不同神经元的 $\tau$ 从零点几毫秒到二十毫秒，延迟从零点几毫秒到几十毫秒。这就是\emph{泄漏积分发放神经元}：带阈值的RC电路。这是有意的简化：皮层神经元的输入分支本身会产生局部脉冲~\cite{Smith2013}，一个神经元的行为更像一个小型多层网络（第\ref{sec:synthesis}章）。对本章的问题而言，一个RC电路就够了。

它与逻辑门的主要区别在于泄漏：神经元对一个脉冲的记忆不是永久的，只持续大约 $\tau$，只有在这个窗口内到达的输入才会相加。

\section{或与与}

从逻辑门开始。如果一个脉冲就能把电压抬到阈值以上，$w\ge\theta$，那么神经元对任一输入的每个脉冲都会发放。这就是\emph{或}门。

与门更有意思。设 $w<\theta<2w$：一个脉冲不够，需要两个。但这时“两个都来了吗”这个问题，在没说明\emph{在多长时间内}之前是没有意义的。设第一个脉冲在时刻 $0$ 到达，第二个在 $\Delta$ 之后到达。第二个到达时，第一个还剩 $we^{-\Delta/\tau}$，若 $we^{-\Delta/\tau}+w\ge\theta$，神经元就发放。由此得到符合窗口：
\begin{empheq}[box=\keybox]{equation}
\Delta \;\le\; \Delta_{\max} \;=\; \tau\,\ln\frac{w}{\theta-w} .
\label{eq:window}
\end{empheq}
这不是逻辑与，而是\emph{符合检测器}：时间上的与（图~\ref{fig:coincidence}）。当 $\theta=1.5w$ 时，窗口等于 $\tau\ln2\approx0.7\tau$。对 $\tau=10\ms$ 的皮层神经元，这约为七毫秒。在听觉系统的神经元中，$\tau$ 小于一毫秒，窗口缩小到零点几毫秒。

\begin{figure}[t]
\centering
\begin{tikzpicture}[>=Stealth,x=0.9cm,y=1.6cm]
\foreach \dx/\lab/\c [count=\i] in {0.5/{符合}/red!70!black, 2.6/{不符合}/blue!60!black}{
 \begin{scope}[xshift={(\i-1)*7.2cm}]
  \draw[->,gray!70!black] (0,0) -- (6.3,0) node[below left,black] {\small $t$};
  \draw[->,gray!70!black] (0,0) -- (0,1.75) node[left,black] {\small $V$};
  \draw[densely dashed,gray!60!black] (0,1.5) -- (6.0,1.5) node[right,black] {\small $\theta$};
  \draw[very thick,\c] (0,0) -- (1,0) -- (1,1)
      plot[domain=1:1+\dx,samples=30] (\x,{exp(-(\x-1)/1.5)});
  \pgfmathsetmacro{\v}{exp(-\dx/1.5)+1}
  \draw[very thick,\c] (1+\dx,{exp(-\dx/1.5)}) -- (1+\dx,{min(\v,1.5)});
  \ifdim\v pt<1.5pt
    \draw[very thick,\c] plot[domain=1+\dx:6,samples=40] (\x,{\v*exp(-(\x-1-\dx)/1.5)});
  \else
    \draw[very thick,\c] (1+\dx,1.5) -- (1+\dx,0) -- (6,0);
    \draw[->,thick,\c] (1+\dx,1.55) -- (1+\dx,1.72) node[above,font=\scriptsize] {脉冲};
  \fi
  \draw[->,thick] (1,-0.35) -- (1,-0.08); \draw[->,thick] (1+\dx,-0.35) -- (1+\dx,-0.08);
  \draw[decorate,decoration={brace,amplitude=3pt,mirror}] (1,-0.42) -- (1+\dx,-0.42) node[midway,below=3pt] {\scriptsize $\Delta$};
  \node[\c] at (3.5,-0.95) {\small \lab};
 \end{scope}}
\end{tikzpicture}
\caption{符合检测器，阈值 $\theta=1.5w$。左：第二个脉冲在 $\Delta<\Delta_{\max}$ 后到达，和达到阈值，神经元发放。右：$\Delta>\Delta_{\max}$，第一个脉冲几乎什么也没剩下，神经元保持沉默。}
\label{fig:coincidence}
\end{figure}

得到与门的另一种办法是利用持续时间，而不是精确时刻。灯亮着的时候，感觉神经元频繁而均匀地发放脉冲；如果一个这样的输入不足以越过阈值，两个就足够，那么只要两个输入都活跃，神经元就一直工作。这样网络按\emph{电平}计算，脉冲到达的精确时间无关紧要。下面的大部分内容正是这样。

\section{做不到的事}

现在试着做一个非门：输入沉默时工作的神经元。做不到：没有输入脉冲，\eqref{eq:glutneuron} 中就没有任何跳变，电压停在零，低于阈值。由许多神经元组成的网络呢？也不行，而且可以对所有网络一次性证明。

用频率来讨论最方便。设 $r_i$ 是神经元 $i$ 的脉冲频率，$I_i$ 是来自输入的驱动，$f_i$ 是它的传递特性：频率如何依赖于总驱动。积分神经元的这一特性是不减的：驱动越大，脉冲越密。达到平衡的网络中，频率满足方程
\begin{equation}
r_i \;=\; f_i\Bigl(\sum_j w_{ij}\,r_j + I_i\Bigr), \qquad w_{ij}\ge 0 .
\label{eq:ratenet}
\end{equation}

\begin{keyblock}
\begin{proposition}[单调性]
\label{prop:monotone}
设网络~\eqref{eq:ratenet} 只由\nt{GLUT}神经元组成，并从完全沉默开始。若加强输入，则没有一个神经元的稳态频率会减小。这样的网络所计算的一切，都是输入的\emph{单调}函数。
\end{proposition}
\end{keyblock}

证明。从沉默 $r^{(0)}=0$ 开始，不断把频率代入~\eqref{eq:ratenet} 的右端：$r^{(k+1)}=F\bigl(r^{(k)}\bigr)$。右端 $F$ 对每个自变量都不减，因为权重非负，$f_i$ 不减。由于 $r^{(1)}\ge0=r^{(0)}$，由归纳法 $r^{(k+1)}\ge r^{(k)}$：频率只增不减，收敛到~\eqref{eq:ratenet} 的最小解。若把输入 $I$ 换成更大的 $I'$，则每一步都有 $r^{(k)}(I)\le r^{(k)}(I')$，极限也是如此。真实网络不是一步一步而是连续地趋于平衡，但结论同样成立：在正连接下，两次运行之间开始时成立的不等式始终保持。

对单个脉冲而言，这一结论并不严格成立：多出的一个输入脉冲可能使神经元提前一点发放，而由于一毫秒的不应期，它的下一个脉冲会丢失。但这个效应弱而不可靠，不能拿来计算。对频率而言，单调性是严格的。

后果与任何没有反相器的电路相同。没有非门。没有异或：加上第二个输入不能熄灭输出。最麻烦的是：\emph{活动不能被活动终止}，只能撤掉维持它的东西。

\section{两根导线：接通与断开}

出路由生物体自己给出。许多感觉器官用两组神经元传递刺激：在视觉中，一些细胞对光变强起反应，另一些对光变弱起反应~\cite{Schiller1992}。我们称之为\emph{接通}神经元和\emph{断开}神经元。网络得到的不是一根导线 $x$，而是一对 $(x,\bar x)$；网络自己算不出的“不存在”，由传感器直接告诉它。

有了一对导线，非门就是免费的：把两根导线对调即可。与门和或门各用两个神经元实现：
\begin{empheq}[box=\keybox]{equation}
\begin{aligned}
x\wedge y \;&\to\; \bigl(\;x\wedge y,\;\; \bar x\vee\bar y\;\bigr),\\
x\vee y \;&\to\; \bigl(\;x\vee y,\;\; \bar x\wedge\bar y\;\bigr).
\end{aligned}
\label{eq:dualrail}
\end{empheq}
右边所有运算都是单调的，所以没有违反命题~\ref{prop:monotone}。

对异步电路来说，双线编码还有第二个优点，而且更重要。一对导线有三种状态：“是”“否”，以及两根都沉默时的\emph{“尚无答案”}。接收者判断计算已经完成，靠的不是时钟，而是信号本身：每一对中恰好有一根导线开始工作。两次刺激之间传感器沉默，网络回到静默，为下一次做好准备。异步电子学正是用这一手法构造在元件延迟任意时都能正确工作的电路~\cite{Sparso2001}。

\begin{keyblock}
\begin{proposition}[异步计算]
\label{prop:dualrail}
如果每个输入都以“接通—断开”一对导线的形式到来，那么由\nt{GLUT}神经元组成的网络可以计算输入的任意逻辑函数。答案同样以成对导线的形式给出，答案是否就绪，看每一对中是否有一根导线开始工作即可。这不需要统一的节拍。代价是神经元数目大约翻倍。
\end{proposition}
\end{keyblock}

例子是异或：“两个刺激中恰好有一个发生了变化”。答案的正导线是 $(x\wedge\bar y)\vee(\bar x\wedge y)$，反导线是 $(x\wedge y)\vee(\bar x\wedge\bar y)$。一共六个神经元，没有一个需要抑制（图~\ref{fig:dualxor}）。

\begin{figure}[t]
\centering
\begin{tikzpicture}[>=Stealth,x=1cm,y=1cm,
  nrn/.style={circle,draw,very thick,fill=blue!6,minimum size=0.75cm,inner sep=0pt,font=\small},
  inp/.style={font=\small}]
\node[inp] (X)  at (0,3.3) {$x$};
\node[inp] (Xb) at (0,2.2) {$\bar x$};
\node[inp] (Y)  at (0,1.1) {$y$};
\node[inp] (Yb) at (0,0)   {$\bar y$};
\node[nrn] (A1) at (3.2,3.3) {$2$};
\node[nrn] (A2) at (3.2,2.2) {$2$};
\node[nrn] (A3) at (3.2,1.1) {$2$};
\node[nrn] (A4) at (3.2,0)   {$2$};
\node[nrn] (O)  at (7.2,2.75) {$1$};
\node[nrn] (Ob) at (7.2,0.55) {$1$};
\foreach \s/\t in {X/A1,Yb/A1,Xb/A2,Y/A2,X/A3,Y/A3,Xb/A4,Yb/A4}{\draw[->,thick,blue!60!black] (\s) -- (\t);}
\draw[->,thick,blue!60!black] (A1) -- node[pos=0.45,above,sloped,font=\scriptsize,black] {$x\wedge\bar y$} (O);
\draw[->,thick,blue!60!black] (A2) -- node[pos=0.45,below,sloped,font=\scriptsize,black] {$\bar x\wedge y$} (O);
\draw[->,thick,blue!60!black] (A3) -- node[pos=0.45,above,sloped,font=\scriptsize,black] {$x\wedge y$} (Ob);
\draw[->,thick,blue!60!black] (A4) -- node[pos=0.45,below,sloped,font=\scriptsize,black] {$\bar x\wedge\bar y$} (Ob);
\draw[->,very thick,red!70!black] (O) -- (8.6,2.75) node[right,black] {$x\oplus y$：“是”};
\draw[->,very thick,red!70!black] (Ob) -- (8.6,0.55) node[right,black] {$\overline{x\oplus y}$：“否”};
\node[below,font=\small,gray!40!black] at (3.2,-0.55) {与：阈值 $2$};
\node[below,font=\small,gray!40!black] at (7.2,-0.55) {或：阈值 $1$};
\node[below,font=\small,gray!40!black] at (0,-0.55) {成对导线};
\end{tikzpicture}
\caption{用双线编码、仅由\nt{GLUT}神经元实现的异或。圆圈中的数是以单个输入权重为单位的阈值。四个与神经元识别输入的四种组合，两个或神经元把它们汇集成一对答案导线。}
\label{fig:dualxor}
\end{figure}

\section{用延迟代替时钟}

涉及时间的计算，只需要导线中的延迟和符合检测器。最好的例子是大脑如何判断声音从哪里来。

来自侧面声源的声音先到达近侧的耳朵，后到达远侧的耳朵。时间差取决于方向：声源在正前方时为零，在正侧方时，对人来说为 $0.22\,\text{m}/343\,\text{m/s}\approx0.6\ms$。大脑能分辨约十\emph{微秒}的差别~\cite{KlumppEady1956,Grothe2010}，尽管神经元本身发放的精度不超过零点几毫秒。这是怎么做到的？

让来自左耳的导线沿一排符合检测器从左向右走，来自右耳的导线从右向左走（图~\ref{fig:itd}）。信号在导线上以速度 $v$ 传播。对于位于长度为 $L$ 的一排中、距左端 $x$ 处的检测器，左侧信号需要走 $x/v$，右侧信号需要走 $(L-x)/v$。如果声音到达右耳晚了 $\delta t$，两个信号就会在满足下式的位置同时相遇：
\begin{empheq}[box=\keybox]{equation}
\frac{x}{v} \;=\; \frac{L-x}{v} + \delta t
\quad\Longrightarrow\quad
x \;=\; \frac{L}{2} + \frac{v\,\delta t}{2} .
\label{eq:itd}
\end{empheq}
只有两个信号在窗口~\eqref{eq:window} 内到达的那个检测器才会发放。发放检测器的编号就是声源的方向。时间差变成了\emph{位置}。

\begin{figure}[t]
\centering
\begin{tikzpicture}[>=Stealth,
  nrn/.style={circle,draw,very thick,fill=blue!6,minimum size=0.7cm,inner sep=0pt}]
\foreach \k in {0,...,6}{\node[nrn] (D\k) at (1.4*\k,0) {\scriptsize $2$};}
\node[nrn,fill=red!15] at (D4) {\scriptsize $2$};
\draw[->,thick,blue!60!black] (-1.4,0.9) node[left,black] {\small 左耳} -- (9.2,0.9);
\draw[->,thick,orange!85!black] (9.8,-0.9) node[right,black] {\small 右耳} -- (-0.8,-0.9);
\foreach \k in {0,...,6}{
  \draw[->,blue!60!black] (1.4*\k,0.9) -- (D\k.north);
  \draw[->,orange!85!black] (1.4*\k,-0.9) -- (D\k.south);}
\draw[->,thick,red!70!black] (D4.north east) ++(0.1,0.05) -- ++(0.5,0.45) node[above right,font=\small,black] {发放};
\node[font=\small,gray!40!black] at (4.2,-1.5) {延迟增大 $\longrightarrow$ \hspace{2.2cm} $\longleftarrow$ 延迟增大};
\pic[scale=1.4] at (-2.3,1.75) {speaker};
\node[font=\scriptsize,align=center] at (-2.3,2.35) {声源在左};
\end{tikzpicture}
\caption{判断声音的方向。来自两耳的信号相向而行；每个圆圈是阈值为 $2$ 的符合检测器。如果声音到达右耳较晚，信号就按公式~\eqref{eq:itd} 在中点右侧相遇。网络的输出是发放神经元的编号。这里声源在左：声音先到达左耳。}
\label{fig:itd}
\end{figure}

这里既没有时钟，也没有抑制，甚至没有双线编码：网络不回答“是”或“否”，而是在许多神经元中指出一个。这种由延迟线和符合检测器组成的电路，看来确实在鸟类听觉脑干中工作~\cite{Carr1990}。微秒级的精度不是来自每个神经元的精度，而是来自检测器数量多、每个都盯着自己的延迟。在哺乳动物中没有找到这样一排延迟：那里同样的任务看来由符合检测器完成，它们由\nt{GLYC}神经元发出的、时间上精确安排的抑制来调谐~\cite{Brand2002,Grothe2010}。抑制如何收窄符合窗口，我们将在第\ref{sec:gaba}章以\nt{GABA}为例讨论；甘氨酸的抑制方式与之相同，也是打开接收者的氯通道。

\section{存储}

计算机需要存储。在这样的网络中，存储就是自我维持的活动。取一群彼此强连接的\nt{GLUT}神经元，用一个频率 $r$ 来描述它（图~\ref{fig:recurrentgroup}a）。平衡时 $r=f(wr+I)$，其中 $w$ 是这群神经元对自身的反馈强度，$I$ 是外来驱动。用图解法求解这个方程：让曲线 $f(wr+I)$ 与直线 $r$ 相交（图~\ref{fig:bistable}）。

\begin{figure}[t]
\centering
% (b): tau dr/dt = -r + f(w r + I), f as in fig:bistable, w=1, I=0.6; Euler dt=0.001 tau.
\begin{tikzpicture}[>=Stealth,
  nrn/.style={circle,draw,thick,fill=blue!6,minimum size=0.5cm,inner sep=0pt}]
\begin{scope}
\node[font=\small] at (-0.5,1.55) {(a)};
\foreach \k in {1,...,6}{\node[nrn] (n\k) at ({1.4+1.0*cos(60*\k)},{1.0*sin(60*\k)}) {};}
\foreach \a/\b in {1/2,2/3,3/4,4/5,5/6,6/1,1/4,2/5,3/6}{\draw[<->,blue!60!black] (n\a) -- (n\b);}
\foreach \k in {2,3,4}{\draw[->,thick,green!45!black] ($(n\k)+(-0.95,0)$) -- (n\k);}
\node[font=\small,green!45!black] at (-0.35,-0.45) {$I$};
\node[font=\Large] at (3.1,0) {$\approx$};
\node[nrn,minimum size=0.85cm,font=\small] (R) at (4.35,-0.1) {$r$};
\draw[->,thick,blue!60!black] (R) edge[loop above,looseness=6] node[above,font=\small] {$w$} (R);
\draw[->,thick,green!45!black] (3.55,-0.95) node[left,font=\small] {$I$} -- (R);
\end{scope}
\begin{scope}[xshift=6.3cm,y=2.6cm,x=1.05cm]
\node[font=\small] at (-0.6,1.25) {(b)};
\draw[->,gray!70!black] (0,0) -- (7.3,0) node[below left,font=\small,black] {$t/\tau$};
\draw[->,gray!70!black] (0,0) -- (0,1.12) node[left,font=\small,black] {$r$};
\foreach \t in {0,2,4,6}{\draw[gray!60!black] (\t,-0.02) -- (\t,0.02); \node[below,font=\scriptsize] at (\t,0) {\t};}
\draw[densely dotted,gray!60!black] (0,0.5) -- (7,0.5) node[right,font=\scriptsize,black,align=left] {写入\\阈值};
\fill[green!45!black,opacity=0.25] (1,-0.2) rectangle (2,-0.13);
\fill[gray!60!black,opacity=0.35] (1,-0.29) rectangle (1.5,-0.22);
\draw[very thick,blue!60!black] plot coordinates {(0.0,0.002) (0.1,0.002) (0.2,0.002) (0.3,0.002) (0.4,0.002) (0.5,0.002) (0.6,0.002) (0.7,0.002) (0.8,0.002) (0.9,0.002) (1.0,0.002) (1.1,0.083) (1.2,0.165) (1.3,0.242) (1.4,0.313) (1.5,0.378) (1.6,0.437) (1.7,0.491) (1.8,0.539) (1.9,0.583) (2.0,0.623) (2.1,0.643) (2.2,0.665) (2.3,0.688) (2.4,0.71) (2.5,0.732) (2.6,0.753) (2.7,0.773) (2.8,0.792) (2.9,0.81) (3.0,0.826) (3.2,0.855) (3.4,0.88) (3.6,0.9) (3.8,0.917) (4.0,0.931) (4.4,0.953) (4.8,0.967) (5.2,0.977) (5.6,0.984) (6.0,0.988) (6.5,0.992) (7.0,0.994)};
\draw[very thick,gray!60!black,densely dashed] plot coordinates {(0.0,0.002) (0.1,0.002) (0.2,0.002) (0.3,0.002) (0.4,0.002) (0.5,0.002) (0.6,0.002) (0.7,0.002) (0.8,0.002) (0.9,0.002) (1.0,0.002) (1.1,0.083) (1.2,0.165) (1.3,0.242) (1.4,0.313) (1.5,0.378) (1.6,0.358) (1.7,0.336) (1.8,0.313) (1.9,0.291) (2.0,0.269) (2.2,0.228) (2.4,0.191) (2.6,0.159) (2.8,0.133) (3.0,0.11) (3.2,0.091) (3.4,0.076) (3.6,0.063) (3.8,0.052) (4.0,0.043) (4.4,0.03) (4.8,0.021) (5.2,0.015) (5.6,0.011) (6.0,0.008) (6.5,0.006) (7.0,0.004)};
\node[font=\scriptsize,blue!60!black,anchor=west] at (3.3,0.8) {驱动 $1\tau$：已写入};
\node[font=\scriptsize,gray!50!black,anchor=west] at (2.3,0.3) {驱动 $0.5\tau$：熄灭};
\end{scope}
\end{tikzpicture}
\caption{与自身强连接的神经元群。(a)~许多相互兴奋的\nt{GLUT}神经元，表现得就像一个频率为 $r$、以强度 $w$ 自我兴奋的神经元；外部送来驱动 $I$。(b)~$w=1$、曲线 $f$ 与图~\ref{fig:bistable} 相同时，神经元群的频率随时间的变化。驱动 $I=0.6$ 施加的时段用坐标轴下方的色条标出。如果它持续 $1\tau$，神经元群来得及越过不稳定点 $r=0.5$，燃起来，驱动结束后仍保持活跃。如果只持续 $0.5\tau$，它够不到这一点，又熄灭下去：要写入一个比特，驱动必须持续约 $0.7\tau$ 以上。}
\label{fig:recurrentgroup}
\end{figure}

\begin{figure}[t]
\centering
\begin{tikzpicture}[>=Stealth,x=5cm,y=3.2cm]
\draw[->,gray!70!black] (0,0) -- (1.1,0) node[below left,black] {\small $r$};
\draw[->,gray!70!black] (0,0) -- (0,1.1);
\draw[thick,gray!60!black] (0,0) -- (1.05,1.05) node[right,black] {\small $r$};
\draw[very thick,blue!60!black] plot[domain=0:1.05,samples=80] (\x,{1/(1+exp(-(\x-0.5)/0.08))});
\draw[very thick,red!70!black,densely dashed] plot[domain=0:1.05,samples=80] (\x,{1/(1+exp(-(\x+0.3-0.5)/0.08))});
\fill[blue!60!black] (0.002,0.002) circle (0.018);
\draw[blue!60!black,fill=white,thick] (0.5,0.5) circle (0.018);
\fill[blue!60!black] (0.998,0.998) circle (0.018);
\node[below right,font=\small] at (0.02,0.0) {关};
\node[right,font=\small] at (0.52,0.46) {不稳定};
\node[below,font=\small] at (0.99,0.96) {开};
\node[anchor=east,font=\small,red!70!black] at (0.19,0.62) {$f(wr+I)$};
\node[anchor=west,font=\small,blue!60!black] at (0.45,0.1) {$f(wr)$};
\end{tikzpicture}
\caption{由强反馈\nt{GLUT}神经元群构成的存储单元。实线是没有外部驱动的情形：与直线 $r$ 有两个稳定交点（“关”和“开”），中间还有一个不稳定交点。短暂的驱动 $I$（虚线）只留下上面那个交点。}
\label{fig:bistable}
\end{figure}

如果反馈足够强，就有三个交点：下面的是静默，上面的是神经元群全力发放，中间的不稳定。这就得到了一个有两个稳定状态的元件——触发器。往里写入1很容易：短暂的驱动抬高曲线，下面的交点消失，神经元群燃起来，并在驱动结束后继续保持活跃（图~\ref{fig:recurrentgroup}b）。驱动太短，来不及把神经元群推到不稳定点，它就又熄灭下去。这种在刺激消失后还能持续数秒的自我维持活动，确实在大脑中观察到，被认为是短时记忆的基础~\cite{Wang2001}。但这并不是记住几秒钟的唯一方式。记录也可以存放在突触本身的慢变量中：一串脉冲之后的易化（就像第\ref{sec:code}章中的“划”突触）能维持约一秒，网络在短暂的爆发之间几乎可以保持沉默——这不是触发器，而是一个充了电的电容器~\cite{Mongillo2008}。在记忆保持期间确实观察到这种“安静”的间歇~\cite{Stokes2015}，而且不靠持续脉冲的存储更省能量。皮层在什么情况下用哪种方式，尚有争论。

那么怎么擦除？按命题~\ref{prop:monotone}，用活动是做不到的；只能\emph{撤掉}某样东西。第一种办法是撤掉支撑：如果这群神经元只有在另一群神经元的驱动下才保持活跃，记忆就随驱动一起熄灭。第二种是疲劳。真实突触在长时间工作时递质储备会耗竭~\cite{Tsodyks1997}，神经元中还会累积降低兴奋性的慢电流。反馈变弱，上面的交点消失，神经元群在几秒后自行熄灭。这样得到的存储由信号开启，却只能由时间关闭。

\section{序列}

可以不要时钟，而用\emph{由事件启动的定时器}。把若干\nt{GLUT}神经元群连成一条链：第一群兴奋第二群，第二群兴奋第三群，依此类推。外部信号点燃第一群，活动波沿链传播，每一环比前一环晚一两个延迟，而且只发放一次：脉冲刚过后神经元处于不应期，而波已经向前传去。

这样的链度量从事件开始经过的时间：第 $k$ 环发放了，就说明从启动算起已经过去了 $k$ 个延迟。把它的输出接到肌肉上，就得到一串自行展开的动作。鸣禽唱歌时，某个脑区中的单个\nt{GLUT}神经元严格依次发放，各自对应歌曲中的一个时刻——非常像这样的链~\cite{Hahnloser2002}。

与时钟不同，这条链在未被启动时保持沉默，只运行一次，而且只为连在它上面的对象计时。网络仍然没有共同的节拍。

\section{还缺什么}

仅由\nt{GLUT}神经元、不用抑制，就能做出很多东西。但还缺三样，而且三样都源于命题~\ref{prop:monotone}。

\emph{选择。}网络必须选定一个方向、一个动作。如果两个刺激几乎相等，图~\ref{fig:itd} 的网络中两个输出都会发放，却没有办法压下弱的、保留强的。

\emph{复位。}不能用信号擦除存储。

\emph{稳定。}在连接并非由工程师规划的网络中，正反馈环不可避免，其中的活动可能雪崩式增长（第\ref{sec:neurontypes}章）。唯一的刹车就是同样的疲劳，既慢又粗糙。

这三种毛病用同一剂药就能治好：用脉冲熄灭脉冲的神经元。这就是\nt{GABA}，需要它不是为了计算，而是为了选择、复位，以及让计算处于控制之下。

\section{习题}

\begin{exercise}[\lvlA]
\label{ex:02:1}
\emph{听觉检测器阵列。}来自两耳的信号沿图~\ref{fig:itd} 的导线以 $5$~m/s的速度传播，最大到达时间差为 $0.64\ms$。检测器是~\eqref{eq:glutneuron} 型神经元，$\tau=0.5\ms$，$\theta=1.5w$，排列得使相邻检测器能区分 $10\,\mu$s。一排有多少个检测器？一个声音会使其中多少个发放？
\end{exercise}

\begin{exercise}[\lvlB]
\label{ex:02:2}
\emph{输入不等使窗口偏移。}符合检测器~\eqref{eq:glutneuron}，$\tau=0.5\ms$，$\theta=1.5$，从权重分别为 $1$ 和 $0.8$ 的两个输入各收到一个脉冲。求符合窗口及其中心。如果来自一只耳朵的输入比另一只弱 $20\,\%$，图~\ref{fig:itd} 的阵列会产生多少微秒的误差？
\end{exercise}

\begin{exercise}[\lvlB]
\label{ex:02:3}
\emph{哪些网络不振荡。}网络 $\tau\dot r_i=-r_i+f_i\bigl(\textstyle\sum_jw_{ij}r_j+I_i\bigr)$，$f_i$ 不减，$i\ne j$ 时 $w_{ij}\ge0$，是单调的，稳定且不振荡。证明：每个环中抑制连接数为偶数的网络，可通过替换 $r_i\to\pm r_i$ 化为这种网络。两群神经元的相互抑制会振荡吗？\nt{GLUT}--\nt{GABA}环呢？
\end{exercise}

\begin{exercise}[\lvlB]
\label{ex:02:4}
\emph{设计：双线加法器。}每个比特用一对导线 $(x,\bar x)$ 传送。用\nt{GLUT}神经元搭建三个比特的全加器，输出和 $(S,\bar S)$ 与进位 $(C,\bar C)$ 两对导线，给出权重和阈值。限定八个神经元以内。若像图~\ref{fig:dualxor} 那样用与、或元件搭建，需要多少个？
\end{exercise}

\begin{exercise}[\lvlB]
\label{ex:02:5}
\emph{建模：写入要多久。}图~\ref{fig:recurrentgroup} 的神经元群：$\tau\dot r=-r+f(r+I)$，$f(u)=1/\bigl(1+e^{-(u-0.5)/0.08}\bigr)$，写入前处于下方状态；驱动 $I$ 施加时长 $T$。数值求出 $I=0.6$ 时能写入比特的最小 $T$，以及阈值 $I_c$：低于它时无论 $T$ 多大都写不进比特。
\end{exercise}

\begin{exercise}[\lvlA]
\label{ex:02:6}
\emph{链作为定时器。}链的每一环是 $100$ 个\nt{GLUT}神经元，延迟 $2\ms$，独立散布 $0.2\ms$。(a)~量出 $1$~s需要多少神经元？相对误差多大？它如何依赖于时间间隔？(b)~真实定时器的误差在任何间隔下都是 $5\,\%$。什么样的散布来源会造成这种情况？
\end{exercise}

\begin{exercise}[\lvlC]
\label{ex:02:7}
\emph{研究：触发器还是电容器。}$100$ 个神经元把一个比特保持 $10$~s。触发器：每个神经元每秒发放 $20$ 个脉冲。电容器：易化量 $u$ 以 $\tau_F=1$~s衰减，$u\ge u_{\max}/2$ 时比特可读，刷新是每个神经元发一串三个脉冲。电容器省多少倍？神经元群被压制 $200\ms$ 后，比特会怎样？
\end{exercise}

\chapter{\nt{GLUT}与\nt{GABA}协同工作}
\label{sec:gaba}

\section{游戏规则}

仅由\nt{GLUT}神经元组成的网络不会选择、不能复位存储、也无法阻止活动发生雪崩，原因都在单调性（命题~\ref{prop:monotone}）。现在加入数量第二多的神经元类型——\nt{GABA}。规则不变：只用活体中真实存在的东西，不设时钟发生器。

神经元模型与~\eqref{eq:glutneuron} 相同，只是来自\nt{GABA}神经元的脉冲不是抬高、而是压低接收者的电压。权重现在有了两种符号，符号由发送者决定，即规则~\eqref{eq:dalesign}。

电路的几个参数。皮层中五分之一到四分之一的神经元是\nt{GABA}神经元~\cite{Hendry1987}。它们的输出很短：抑制的是邻居，而不是远处的脑区。抑制在一毫秒后到达，持续几毫秒。没有输入时\nt{GABA}神经元保持沉默：要抑制别人，它自己必须先受到兴奋，通常来自同一网络中的\nt{GLUT}神经元。

因此，从\nt{GLUT}神经元 $A$ 到神经元 $B$ 的负连接必须通过中介实现：$A$ 兴奋一个\nt{GABA}神经元，后者再抑制 $B$。改变符号的代价是一个神经元和一个延迟，约一毫秒。

对抑制而言，重要的不只是连接\emph{来自谁}，还有它\emph{落在哪里}：落点在很大程度上决定了连接的作用~\cite{Markram2004}。\nt{GLUT}输出主要落在接收者的输入分支即\emph{树突}上，携带数据：每个这样的输入是总和中的一项。落在神经元胞体上的输出作用于整个总和：许多快速\nt{GABA}神经元正是以这种方式对接收者的响应做除法，并决定它何时发放。落在接收者产生脉冲之处的输出是最后的裁决，一次否决整个输出。而落在别的导线末梢上的输出，只改变一条输入线路的释放概率 $p$，不影响其他线路~\cite{Rudomin1999}；第\ref{sec:pain}章中的疼痛闸门就是这样工作的。在脑外，输出落在肌纤维上（第\ref{sec:muscles}章）和器官上（第\ref{sec:chemoutput}章），还有一些输出哪里也不落，而是把递质释放到组织间隙或血液中（第\ref{sec:serocomputer}章和第\ref{sec:chemoutput}章）。

\section{非与否决}

现在能做非门了吗？差不多。输入沉默时还能工作的神经元，必须从某处获得兴奋。活体大脑中有这种兴奋：每个神经元不断收到来自成千上万其他神经元的背景活动。让背景使神经元 $Y$ 保持工作，输入 $x$ 经\nt{GABA}中介抑制它。这就是非门。

不过更有用的是\emph{否决}：“$a$，但 $b$ 出现时除外”：
\begin{empheq}[box=\keybox]{equation}
y \;=\; a \wedge \bar b .
\label{eq:veto}
\end{empheq}
神经元 $Y$ 从 $a$ 得到兴奋，从被 $b$ 兴奋的\nt{GABA}神经元得到抑制。这里不需要背景：兴奋由 $a$ 自己提供。用否决和或可以搭出一切。例如异或：$x\oplus y=(x\wedge\bar y)\vee(\bar x\wedge y)$，即两个否决神经元、两个\nt{GABA}中介和一个或神经元。

\begin{keyblock}
\begin{proposition}[\nt{GLUT}+\nt{GABA}的完备性]
\label{prop:gabacomplete}
由\nt{GLUT}和\nt{GABA}神经元组成的网络可以计算输入的任意逻辑函数，而且每个比特只用一根导线传送。第\ref{sec:glutcomputer}章中的“接通—断开”传感器不再需要。如果函数在所有输入都沉默时应输出1，就需要一个背景活动源。
\end{proposition}
\end{keyblock}

证明：只要有常数1，否决~\eqref{eq:veto} 加上或就是函数完备的，因为非 $x$ 就是否决“1，但 $x$ 出现时除外”。而在所有输入沉默时输出为0的函数，不需要常数1，只用否决和或就能搭出。

\section{运动方向}

时间上的否决，就像第\ref{sec:glutcomputer}章中时间上的与一样，给出运动方向检测器。

设两个相邻的传感器 $A$ 和 $B$ 相距 $\Delta x$。输出神经元 $Y$ 从 $B$ 得到兴奋，从 $A$ 经\nt{GABA}中介得到抑制，延迟为 $d$（图~\ref{fig:direction}）。以速度 $v$ 从 $A$ 移向 $B$ 的光斑在时刻 $0$ 经过 $A$，抑制在时刻 $d$ 到达 $Y$；光斑在时刻 $\Delta x/v$ 经过 $B$，兴奋也在这时到达。如果这两个时刻在窗口~\eqref{eq:window} 的精度内重合，抑制就抵消兴奋，$Y$ 保持沉默。这发生在速度约为
\begin{empheq}[box=\keybox]{equation}
v^* \;=\; \frac{\Delta x}{d}
\label{eq:vstar}
\end{empheq}
时。从 $B$ 移向 $A$ 时，来自 $B$ 的兴奋在时刻 $0$ 到达，而来自 $A$ 的抑制要到时刻 $\Delta x/v+d$ 才到，这时 $Y$ 已经发放了。

\begin{figure}[t]
\centering
\begin{tikzpicture}[>=Stealth,
  nrn/.style={circle,draw,very thick,fill=blue!6,minimum size=0.9cm,inner sep=0pt},
  gab/.style={nrn,fill=red!10},
  inh/.style={-{Bar[width=7pt]},thick,red!70!black}]
\node[nrn] (A) at (0,2) {$A$};
\node[nrn] (B) at (3,2) {$B$};
\node[gab] (G) at (0.9,0.6) {\small \nt{GABA}};
\node[nrn] (Y) at (3,-0.6) {$Y$};
\draw[->,thick] (A) -- (G);
\draw[inh] (G) -- node[below left=-2pt] {\scriptsize 延迟 $d$} (Y);
\draw[->,thick] (B) -- (Y);
\draw[->,very thick,red!70!black] (Y) -- (4.6,-0.6) node[right,black] {\small 输出};
\draw[->,thick,orange!85!black] (-0.6,2.9) -- (3.6,2.9) node[right,black] {\small 沉默};
\draw[<-,thick,orange!85!black] (-0.6,3.4) -- (3.6,3.4) node[right,black] {\small 响应};
\foreach \yy/\dd in {2.9/-1,3.4/1}{
  \foreach \k in {1,2,3}{\draw[orange!70!yellow,line width=0.8pt] (1.5+\dd*0.12+\dd*0.13*\k,\yy+0.1-0.1*\k+0.1) -- ++(\dd*0.25,0);}
  \fill[yellow!70!orange,draw=orange!70!black] (1.5,\yy) circle (0.14);}
\node[font=\scriptsize,align=center] at (1.5,3.95) {光斑};
\end{tikzpicture}
\caption{运动方向检测器。$B$ 兴奋 $Y$，$A$ 经\nt{GABA}中介以延迟 $d$ 抑制 $Y$。以约 $\Delta x/d$ 的速度从 $A$ 移向 $B$ 时，抑制抵消兴奋；反向运动时，抑制来迟了。带短线的黄色光斑表示在传感器上方移动的光。}
\label{fig:direction}
\end{figure}

在视网膜中，对运动方向的选择性看来正是这样产生的：来自不应引起响应的运动方向一侧的不对称抑制~\cite{Barlow1965}。

\section{减法还是除法}

到目前为止，我们的抑制做的是减法。实际上它的机制更精巧。

突触打开一个通道，也就是接入一个电导。兴奋性突触的平衡电位远高于阈值，比静息电位高约 $70\mV$：打开的通道把电压往上拉。抑制性\nt{GABA}突触的通道让氯离子通过，而氯的平衡电位在静息电位附近，所以打开的通道不是把电压往下拉，而是拉\emph{向静息}。以静息电位为零点，具有泄漏电导 $g_L$、兴奋性电导 $g_E$ 和抑制性电导 $g_I$ 的膜的稳态电压为
\begin{empheq}[box=\keybox]{equation}
V_\infty \;=\; \frac{g_E\,E_E}{g_L+g_E+g_I} ,
\label{eq:shunt}
\end{empheq}
其中 $E_E\approx70\mV$ 是兴奋性突触的平衡电位。这就是分压器的欧姆定律：$g_E$ 把电压拉向 $E_E$，$g_L$ 和 $g_I$ 把它拉向零。

$g_I$ 不是带着负号出现在分子里，而是出现在分母里：抑制不是从兴奋中\emph{减去}，而是把兴奋\emph{除以}某个数（图~\ref{fig:shunt}）。这种抑制称为分流抑制：打开的氯通道把膜短路了。对网络来说，这是增益调节。如果 $g_I$ 与邻居的总活动成正比，每个神经元响应的就不是自己输入的绝对强度，而是它在总活动中所占的份额。这种除以总活动的运算见于大脑的许多部位，被认为是大脑的标准运算之一~\cite{Carandini2012}：同一个网络在输入弱时和输入强时都能工作。

需要说明：公式~\eqref{eq:shunt} 讲的是不发放脉冲的神经元。对于正在发放的神经元，电压始终停在阈值附近，流过抑制性电导的电流几乎恒定，分流对脉冲\emph{频率}的作用主要是减法~\cite{HoltKoch1997}。只有同时给神经元加上相互平衡的兴奋性和抑制性背景电导，就像活体皮层那种抖动而平衡的状态，频率才会被除~\cite{Chance2002}。所以大脑的除法不是来自单独的分流，而是来自抑制加上背景噪声~\cite{Carandini2012}。

\begin{figure}[t]
\centering
\begin{tikzpicture}[>=Stealth,x=1.25cm,y=0.05cm]
\draw[->,gray!70!black] (0,0) -- (5.4,0) node[right,black] {\small $g_E/g_L$};
\draw[->,gray!70!black] (0,0) -- (0,68) node[above,black] {\small $V_\infty$，mV};
\foreach \v in {20,40,60}{\draw[gray!60!black] (-0.05,\v) -- (0.05,\v) node[left=3pt,font=\scriptsize] {\v};}
\foreach \x in {1,2,3,4,5}{\draw[gray!60!black] (\x,-1.5) -- (\x,1.5) node[below=5pt,font=\scriptsize] {\x};}
\draw[very thick,blue!60!black] plot[domain=0:5,samples=60] (\x,{70*\x/(1+\x)});
\draw[very thick,red!70!black] plot[domain=0:5,samples=60] (\x,{70*\x/(3+\x)});
\draw[thick,densely dashed,gray!50!black] plot[domain=0.56:5,samples=60] (\x,{70*\x/(1+\x)-25});
\node[right,font=\small,blue!60!black] at (5.02,58.3) {无抑制};
\node[right,font=\small,red!70!black] at (5.02,45.5) {除法：$g_I=2g_L$};
\node[right,font=\small,gray!40!black] at (5.02,32.5) {减去 $25\mV$};
\end{tikzpicture}
\caption{分流抑制对电压做除法，而不是减法。蓝色曲线是无抑制时的稳态电压~\eqref{eq:shunt}，红色曲线是加上抑制性电导 $g_I=2g_L$ 后的情形。红色曲线就是蓝色曲线沿水平方向拉伸三倍：仿佛输入被除以 $1+g_I/g_L=3$。虚线是减去恒定 $25\mV$ 的抑制：整条曲线下移，斜率不变。}
\label{fig:shunt}
\end{figure}

还有第二个推论。膜时间常数 $\tau_{\text{eff}}=C/(g_L+g_E+g_I)$，抑制会使它变短。而符合窗口~\eqref{eq:window} 与 $\tau$ 成正比，所以与兴奋同时到达的抑制会\emph{收窄窗口}。输入直接兴奋神经元，同时经\nt{GABA}中介以一毫秒的延迟抑制它，这是大脑中最常见的电路之一：神经元只来得及响应最初一两毫秒内到达的输入~\cite{Pouille2001}。

\section{选择}

现在来看\nt{GLUT}网络最缺的东西：从几个选项中选一个。取两群\nt{GLUT}神经元，输入分别为 $I_1$ 和 $I_2$。每群经自己的\nt{GABA}中介以强度 $\beta$ 抑制另一群（图~\ref{fig:wta}）。对频率 $r_1,r_2$ 有
\begin{equation}
\tau\,\frac{\dd r_1}{\dd t} = -r_1 + \bigl[I_1-\beta r_2\bigr]_+ ,\qquad
\tau\,\frac{\dd r_2}{\dd t} = -r_2 + \bigl[I_2-\beta r_1\bigr]_+ ,
\label{eq:wta}
\end{equation}
其中 $[u]_+=\max(u,0)$：频率不能为负。

\begin{figure}[t]
\centering
\begin{tikzpicture}[>=Stealth,
  nrn/.style={circle,draw,very thick,fill=blue!6,minimum size=0.9cm,inner sep=0pt},
  gab/.style={nrn,fill=red!10,minimum size=0.75cm},
  inh/.style={-{Bar[width=7pt]},thick,red!70!black}]
\node[nrn] (E1) at (0,0) {$r_1$};
\node[nrn] (E2) at (4,0) {$r_2$};
\node[gab] (G1) at (1.4,1.1) {};
\node[gab] (G2) at (2.6,-1.1) {};
\draw[->,thick] (E1) -- (G1); \draw[inh] (G1) -- (E2);
\draw[->,thick] (E2) -- (G2); \draw[inh] (G2) -- (E1);
\draw[->,thick,blue!60!black] (-1.6,0) node[left,black] {$I_1$} -- (E1);
\draw[->,thick,blue!60!black] (5.6,0) node[right,black] {$I_2$} -- (E2);
\end{tikzpicture}
\caption{二选一。两群\nt{GLUT}神经元（蓝）经\nt{GABA}中介（红）相互抑制。}
\label{fig:wta}
\end{figure}

只要两个神经元都在工作，系统~\eqref{eq:wta} 就是线性的，其行为由两个特征值决定：两个频率同向偏离时为 $-(1+\beta)/\tau$，反向偏离时为 $-(1-\beta)/\tau$。抑制弱时，$\beta<1$，两个特征值都是负的：两个神经元都工作，抑制只是压低弱者。抑制强时，$\beta>1$，第二个特征值变为正：$r_1$ 与 $r_2$ 之间的任何差别都会增长，直到一个神经元完全沉默。

\begin{keyblock}
\begin{proposition}[赢者通吃]
\label{prop:wta}
如果相互抑制强于1，即 $\beta>1$，两群都工作的状态就是不稳定的。网络进入一群工作、另一群沉默的状态。频率为 $r_1=I_1$ 的赢家只要 $I_2<\beta I_1$ 就能保住胜利。
\end{proposition}
\end{keyblock}

我们在模型上验证了这一点。$\beta=0.5$、输入为 $1.0$ 和 $0.9$ 时，两群都工作，频率分别为 $0.73$ 和 $0.53$。$\beta=2$、输入相同时，第二群完全沉默。这种选择带有记忆：如果之后把输入对调，原来的赢家仍然是赢家，因为 $1.0<2\cdot0.9$。

一百群神经元也是一样：每群抑制其余各群，只有一群存活。

\section{存储复位}

第\ref{sec:glutcomputer}章中的存储只能靠疲劳关闭。现在给它加一个\nt{GABA}神经元，由“复位”信号兴奋，并抑制这群神经元。抑制使图~\ref{fig:bistable} 中的曲线 $f(wr+I)$ 下移，上面的交点消失，神经元群熄灭，复位信号结束后仍停留在下方状态。现在存储由一个信号写入、由另一个信号擦除，就像带置位和复位输入的触发器。

更漂亮的做法是：把两个这样的神经元群像图~\ref{fig:wta} 那样用相互抑制连接起来，并给每群加上自反馈。送到其中一群输入端的信号使存储单元进入相应状态，单元记住最后一次决定。这正是由两个交叉耦合的或非门组成的锁存器的直接对应物。

\section{稳定性与节律}

剩下\nt{GLUT}网络的第三种毛病：雪崩。取一群自反馈为 $w_{EE}$ 的\nt{GLUT}神经元，以及一群\nt{GABA}神经元：前者以强度 $w_{IE}$ 兴奋后者，后者以强度 $w_{EI}$ 抑制前者。对频率 $E$ 和 $I$ 有
\begin{equation}
\tau_E\frac{\dd E}{\dd t} = -E + w_{EE}E - w_{EI}I + h, \qquad
\tau_I\frac{\dd I}{\dd t} = -I + w_{IE}E ,
\label{eq:ei}
\end{equation}
其中 $h$ 是外部输入~\cite{WilsonCowan1972}。没有抑制时，若 $w_{EE}>1$，活动会无限增长：每个脉冲引发不止一个后续脉冲。有了抑制，稳态频率
\begin{equation}
E^* \;=\; \frac{h}{1-w_{EE}+w_{EI}w_{IE}}
\label{eq:eistar}
\end{equation}
只要分母为正就是有限的。但这还不够：平衡点还必须是吸引的。对~\eqref{eq:ei} 做线性分析，得到两个条件。

\begin{keyblock}
\begin{proposition}[稳定性条件]
\label{prop:ei}
平衡点~\eqref{eq:eistar} 稳定的条件是抑制
\begin{enumerate}
\item[(i)] \emph{足够强}：$w_{EI}\,w_{IE} > w_{EE}-1$；
\item[(ii)] \emph{足够快}：$\tau_I < \tau_E/(w_{EE}-1)$。
\end{enumerate}
若违反 (i)，活动雪崩式增长。若满足 (i) 而违反 (ii)，抑制来迟，活动开始振荡。
\end{proposition}
\end{keyblock}

条件~(i) 来自系统~\eqref{eq:ei} 矩阵的行列式，条件~(ii) 来自它的迹。第二个条件的含义，任何调过控制器的人都明白：迟到的反馈不能消除偏差，反而会把它激起来。

\begin{figure}[t]
\centering
\begin{tikzpicture}[>=Stealth]
\draw[->,gray!70!black] (0,0) -- (10.9,0) node[below left,black] {\small $t$，ms};
\draw[->,gray!70!black] (0,0) -- (0,3.3) node[left,black] {\small $E$};
\foreach \t in {0,50,100,150}{\draw[gray!70!black] (\t*0.07,0.06) -- (\t*0.07,-0.06) node[below,font=\scriptsize] {\t};}
\input{figures/ei_traces}
\node[font=\small,gray!40!black,anchor=west] at (0.9,3.15) {无抑制：雪崩};
\node[font=\small,blue!60!black,anchor=west] at (7.4,1.25) {快速抑制};
\node[font=\small,red!70!black,anchor=west] at (7.4,2.55) {慢速抑制};
\end{tikzpicture}
\caption{接通输入后，反馈为 $w_{EE}=1.5$、$\tau_E=10\ms$ 的\nt{GLUT}神经元群的频率。无抑制时（虚线）活动无限增长。有强度 $w_{EI}w_{IE}=2$、$\tau_I=2\ms$ 的抑制时（蓝色曲线），活动达到 $E^*=h/(1-1.5+2)=0.67h$ 的水平。抑制强度相同但较慢，$\tau_I=30\ms$ 时（红色曲线），活动发生振荡。}
\label{fig:ei}
\end{figure}

一般认为，皮层正是工作在这种状态：它自身的反馈强到足以在没有抑制时进入雪崩，只靠快速抑制把它维持在工作点上~\cite{Tsodyks1997isn}。

这种状态有一个反直觉的特征，可以从外部辨认出来~\cite{Tsodyks1997isn}。在~\eqref{eq:ei} 中给\nt{GABA}神经元群直接加上外部输入 $h_I$。稳态频率变为
\begin{equation}
E^* \;=\; \frac{h-w_{EI}\,h_I}{D},\qquad
I^* \;=\; \frac{w_{IE}\,h+(1-w_{EE})\,h_I}{D},\qquad
D \;=\; 1-w_{EE}+w_{EI}w_{IE} .
\label{eq:isn}
\end{equation}
若 $w_{EE}>1$，第二个公式中 $h_I$ 的系数为负：推一下\nt{GABA}神经元，它们的频率反而\emph{下降}。原因在于环路。额外的抑制压低了\nt{GLUT}神经元群，后者对\nt{GABA}神经元的兴奋减少，于是\nt{GABA}神经元失去的比得到的更多。按图~\ref{fig:ei} 的数值（$w_{EE}=1.5$，$w_{EI}w_{IE}=2$，$D=1.5$），每增加一个单位的输入，$I^*$ 就降低三分之一个单位。皮层中确实发现了这个特征：当视皮层神经元的响应被周围刺激抑制时，它们收到的抑制性电流不增反降，与兴奋性电流一同下降~\cite{Ozeki2009}。后来在皮层的不同区域和不同层中都发现了同样的特征~\cite{Sanzeni2020}。

违反条件~(ii) 时的振荡在大脑中同样存在。“\nt{GLUT}兴奋\nt{GABA}，\nt{GABA}抑制\nt{GLUT}”这个带几毫秒延迟的环路，产生周期约 $12$--$50\ms$（频率 $20$--$80$\,Hz）的节律，皮层中这样的快速节律众所周知~\cite{Cardin2009,Buzsaki2012}。这不是计算机的时钟：节律是局部的，只在网络活跃时才出现。但在几个周期之内，它把时间划分成神经元可以发放的窗口。

\section{小结}

\begin{table}[t]
\centering
\small
\begin{tabular}{>{\raggedright}p{3.3cm}>{\raggedright}p{4.4cm}>{\raggedright\arraybackslash}p{4.4cm}}
\toprule
& 只有\nt{GLUT} & \nt{GLUT} + \nt{GABA} \\
\midrule
非 & 只能借助“接通—断开”传感器 & 否决 $a\wedge\bar b$；有背景活动时可实现非 \\
每比特导线数 & 两根 & 一根 \\
运动方向 & 无 & 带延迟的否决 \\
增益调节 & 无 & 除法：分流加背景噪声 \\
符合窗口 & $\sim\tau$ & 被抑制收窄 \\
多选一 & 无 & 相互抑制，$\beta>1$ \\
存储复位 & 只能靠疲劳 & 经\nt{GABA}用信号复位 \\
稳定性 & 只能靠疲劳 & 快速负反馈 \\
节律 & 不需要 & 抑制慢时出现 \\
\bottomrule
\end{tabular}
\caption{\nt{GABA}给\nt{GLUT}神经元网络带来了什么。}
\label{tab:gaba}
\end{table}

\nt{GABA}几乎没有给网络增加新的\emph{计算}——这些用双线传感器也能算出来——它增加的是\emph{控制}（表~\ref{tab:gaba}）：选择、复位、增益调节和稳定性。大脑中的兴奋承载数据，而抑制决定哪些数据通过、何时通过、以多大音量通过。对皮层中每四五个神经元中的一个而言，这是非常合理的分工。

\section{习题}

\begin{exercise}[\lvlA]
\label{ex:03:1}
\emph{分流的数值。}$E_E=70\mV$。由~\eqref{eq:shunt} 求 $g_E=g_L$ 和 $4g_L$ 时的 $V_\infty$，分别在无抑制和有分流 $g_I=2g_L$ 的情况下。一个在 $g_E=g_L$ 时与分流扣除等量电压的减法，在 $g_E=4g_L$ 时会给出多大的 $V_\infty$？分流使符合窗口收窄几倍？
\end{exercise}

\begin{exercise}[\lvlB]
\label{ex:03:2}
\emph{推导稳定性条件。}求线性系统~\eqref{eq:ei} 矩阵的迹和行列式，并由此推出命题~\ref{prop:ei} 的两个条件。在条件~(ii) 的边界上，平衡点以振荡方式失稳。按图~\ref{fig:ei} 的数值求这一振荡的周期。
\end{exercise}

\begin{exercise}[\lvlB]
\label{ex:03:3}
\emph{由延迟产生节律。}把模型~\eqref{eq:ei} 中的慢抑制换成带延迟 $d$ 的快抑制：$\tau_E\dot E=-E(t)-GE(t-d)+h$。$\tau_E=10\ms$ 时，对 $G=2$ 和 $G=5$ 求产生振荡的最小延迟 $d_c$ 及振荡周期。与习题~\ref{ex:03:2} 不同，它们是否落在皮层快速节律的 $12$--$50\ms$ 范围内？
\end{exercise}

\begin{exercise}[\lvlB]
\label{ex:03:4}
\emph{建模：带记忆的选择。}在 $\beta=2$、$\tau=1$ 时对~\eqref{eq:wta} 积分。(a)~$I_1=1$，把 $I_2$ 从 $0$ 缓慢升到 $3$ 再降回来：网络在哪些 $I_2$ 处切换？(b)~当输入差 $\varepsilon$ 缩小到十分之一时，从静默出发的决策时间增加多少？
\end{exercise}

\begin{exercise}[\lvlB]
\label{ex:03:5}
\emph{设计：速度带检测器。}图~\ref{fig:direction} 的检测器在从 $A$ 到 $B$ 的运动、渡越时间为 $5$--$20\ms$ 时必须沉默。延迟为 $d_k$ 的\nt{GABA}中介在 $A$ 发放后的区间 $[d_k,\,d_k+5\ms]$ 内否决 $Y$；来自 $B$ 的兴奋立即到达。需要几个中介、延迟各为多少？
\end{exercise}

\begin{exercise}[\lvlB]
\label{ex:03:6}
\emph{否决与背景。}在“每个输入一个\nt{GABA}中介、每个 $f=1$ 的输入组合一个否决、再加一个总的或”的通用方案中，奇偶校验 $x\oplus y\oplus z$ 需要多少个神经元？用两个直接接收输入的神经元（一个\nt{GABA}、一个\nt{GLUT}）实现它，给出权重和阈值。
\end{exercise}

\begin{exercise}[\lvlC]
\label{ex:03:7}
\emph{设计：百里挑一。}$100$ 群\nt{GLUT}中的每一群都必须同等地抑制所有其他群，但不抑制自己。线性\nt{GABA}中介由神经元群兴奋并抑制它们；神经元群可以相互兴奋，但不能自我兴奋。中介最少要几个？如果完全没有直接的兴奋性连接呢？
\end{exercise}

\begin{exercise}[\lvlC]
\label{ex:03:8}
\emph{研究：抑制何时反常。}在图~\ref{fig:ei} 的模型中（$w_{EI}=2$，$w_{IE}=1$），\nt{GLUT}群的传递函数改为 $f(u)=[u]_+^2$，\nt{GABA}群仍是线性的。输入超过哪个值 $h_c$ 时，给\nt{GABA}群增加的输入反而降低它自己的频率？用模拟加以验证。
\end{exercise}

\chapter{摩尔斯电码：关于\nt{GLUT}和\nt{GABA}神经元所说的语言，我们知道什么}
\chaptermark{摩尔斯电码}
\label{sec:code}

\section{只有一个符号的字母表}

摩尔斯电码有两个符号——点和划，符号之间还有三种长度的停顿。而神经元的字母表更贫乏，正如我们在第\ref{sec:neurontypes}章中看到的：只有一个符号。脉冲持续约一毫秒，同一个神经元的所有脉冲高度和形状都相同。因此全部消息都在停顿里，也就是在时刻序列 $t_1,t_2,t_3,\dots$ 之中。这是一种没有划的电码，意义只由点的节奏承载（图~\ref{fig:morse}）。

\begin{figure}[t]
\centering
\begin{tikzpicture}[>=Stealth,x=0.08cm,y=1cm]
% Morse letter: dot, dash, dot
\node[left,font=\small] at (-6,2.55) {摩尔斯：};
\fill[black] (0,2.45) rectangle (6,2.65);
\fill[black] (14,2.45) rectangle (32,2.65);
\fill[black] (40,2.45) rectangle (46,2.65);
\node[right,font=\small,gray!40!black] at (52,2.55) {点、划、点——意义在于时长和停顿};
% stimulus onset
\node[left,font=\small] at (-6,0.4) {神经元：};
\draw[->,thick,green!45!black] (0,-0.55) -- (0,-0.05);
\node[below,font=\scriptsize,green!45!black] at (0,-0.55) {刺激};
\draw[gray!70!black] (0,0) -- (152,0) node[right,black] {\small $t$};
% spikes: single ones blue, the burst red
\foreach \s in {14,26,41,88,112,131}{\draw[very thick,blue!60!black] (\s,0) -- (\s,0.8);}
\foreach \s in {60,63,66,69}{\draw[very thick,red!70!black] (\s,0) -- (\s,0.8);}
% latency
\draw[<->,thick] (0,1.1) -- node[above,font=\scriptsize] {$t_1$：首个脉冲的时间} (14,1.1);
% burst
\draw[decorate,decoration={brace,amplitude=4pt},red!70!black] (58,0.95) -- (71,0.95) node[midway,above=4pt,font=\scriptsize] {簇放电——“划”};
% counting windows
\foreach \a/\b/\n in {0/50/3,50/100/5,100/150/2}{
  \draw[decorate,decoration={brace,amplitude=4pt,mirror},gray!50!black] (\a+1,-0.2) -- (\b-1,-0.2) node[midway,below=4pt,font=\small] {$n=\n$};}
\node[font=\scriptsize,gray!40!black] at (75,-1.05) {频率编码：每个窗口内的脉冲数 $n$};
\foreach \x in {0,50,100,150}{\draw[gray!60!black] (\x,-0.06) -- (\x,0.06);}
\end{tikzpicture}
\caption{摩尔斯电码与神经元的电码。同一串脉冲可以有不同的读法：数出每 $50\ms$ 窗口内的脉冲数（\ref{sec:rate}~节），测量延迟 $t_1$~\eqref{eq:latency}，或者注意到簇放电——“划”~\eqref{eq:burst}。}
\label{fig:morse}
\end{figure}

停顿可以有多长？下限由不应期决定：一个脉冲之后，神经元约一毫秒内不能再次发放，所以每秒不会超过几百个脉冲。上限则没有：神经元可以沉默好几分钟。皮层\nt{GLUT}神经元的频率从每秒零点几个到几十个脉冲不等，大多数神经元大部分时间工作在下限附近。

在第\ref{sec:glutcomputer}章和第\ref{sec:gaba}章中，我们顺带采用过这样或那样用脉冲记录数的方式。现在直接发问：\nt{GLUT}和\nt{GABA}神经元彼此用什么语言交谈？完整的答案还没有，这种语言尚未破译。但关于它，有一些事情是确知的，而且几乎所有已知的内容，都可以像通信工程师那样推理出来：信道容量、噪声、接收机、传输代价。

\section{能传多少}

先看上限。设接收者分辨脉冲时刻的精度为 $\Delta t$：小于 $\Delta t$ 的偏移它分辨不出。把时间切成长度为 $\Delta t$ 的格子，格子比不应期短，所以每个格子里要么有一个脉冲，要么没有（图~\ref{fig:cells}）。平均频率为 $r$ 时，脉冲落在某个格子里的概率为 $p=r\Delta t$。当各格子相互独立时，脉冲流携带的信息最多，这时每个格子贡献的熵为 $-p\log_2p-(1-p)\log_2(1-p)$，当 $p\ll1$ 时约等于 $p\log_2(e/p)$ 比特。除以 $\Delta t$，就得到极限传输速率及其摊到每个脉冲上的份额~\cite{MacKay1952}：
\begin{empheq}[box=\keybox]{equation}
R_{\max} \;\approx\; r\,\log_2\frac{e}{r\,\Delta t}\ \ \text{bit/s},
\qquad
\frac{R_{\max}}{r} \;\approx\; \log_2\frac{e}{r\,\Delta t}\ \ \text{bit/脉冲}.
\label{eq:spikeentropy}
\end{empheq}
当 $r=10$ 个脉冲每秒、$\Delta t=1\ms$ 时，约为每个脉冲八比特、每秒八十比特。

\begin{figure}[t]
\centering
\begin{tikzpicture}[>=Stealth,x=0.5cm,y=1cm]
% time cut into cells of width Delta t; a cell holds one impulse or none
\foreach \k in {0,...,23}{\draw[gray!55] (\k,0) rectangle (\k+1,0.7);}
\foreach \k in {2,7,8,15,21}{
  \fill[red!10] (\k,0) rectangle (\k+1,0.7);
  \draw[very thick,red!70!black] (\k+0.5,0.05) -- (\k+0.5,0.65);}
\foreach \k in {0,...,23}{
  \pgfmathtruncatemacro{\bit}{(\k==2)||(\k==7)||(\k==8)||(\k==15)||(\k==21)}
  \ifnum\bit=1 \def\bitcol{red!70!black}\else \def\bitcol{gray!60!black}\fi
  \node[font=\scriptsize,text=\bitcol] at (\k+0.5,-0.25) {\bit};}
\draw[<->,gray!60!black] (0,0.95) -- (1,0.95) node[midway,above,font=\scriptsize,black] {$\Delta t$};
\node[font=\small,anchor=east] at (-0.3,0.35) {脉冲};
\node[font=\small,anchor=east] at (-0.3,-0.25) {比特};
\draw[decorate,decoration={brace,amplitude=5pt,mirror},gray!60!black] (0,-0.55) -- (24,-0.55)
  node[midway,below=6pt,font=\scriptsize,black,align=center] {只会计数的接收者：“$n=5$”；\\ 能看见格子的接收者：$24$ 格中具体是哪 $5$ 格——这是 $\log_2\binom{24}{5}\approx15$ 比特};
\end{tikzpicture}
\caption{容量~\eqref{eq:spikeentropy} 从何而来。时间被切成长度为 $\Delta t$ 的格子，每个格子里要么有脉冲（$1$），要么没有（$0$）：脉冲流就是一个二进制串。脉冲落入一个格子的概率为 $p=r\Delta t$。格子越细，串越长，比特越多；脉冲计数器几乎丢掉了“1 \emph{在哪里}”所记录的全部信息。}
\label{fig:cells}
\end{figure}

每个脉冲的比特数只以对数方式依赖于时间精度：精度为 $10\ms$ 时，每个脉冲仍有约五比特。但如果接收者只数每秒的脉冲数，那么在 $r=10$ 时，它在整整一秒内得到的还不到两比特（\ref{sec:rate}~节）。

\begin{keyblock}
\begin{proposition}[脉冲信道的容量]
\label{prop:capacity}
平均频率为 $r$、以时间精度 $\Delta t$ 读取的脉冲流，携带的信息不超过~\eqref{eq:spikeentropy}。这个容量几乎全部在于脉冲的\emph{时间}：只会计数的接收者从同一脉冲流中提取的信息，比能以毫秒精度分辨脉冲时刻的接收者少几十倍。
\end{proposition}
\end{keyblock}

这是上限，不是测量值。在已知刺激的感觉神经元上测得的是每个脉冲一到几比特；最好的情况下可达按其精度计算的上限的一半~\cite{Rieke1997}。可见，至少在神经系统的入口处，脉冲的时间是被利用了的，而不是白白浪费。

\section{有噪声的信道}

容量~\eqref{eq:spikeentropy} 是在无噪声的条件下算出的。关于噪声产生在哪里，有三个确知的事实。

\emph{脉冲发生器本身是精确的。}把皮层神经元连同一小块组织从脑中取出，反复给它注入同一个随时间快速变化的电流，脉冲以约一毫秒的精度重复出现~\cite{Mainen1995}。如果电流恒定，脉冲时刻每次都会漂移开来：产生精确时间的是输入的快速跳变，而不是神经元本身。

\emph{突触不可靠。}到达\nt{GLUT}突触的脉冲只以一定概率 $p$ 释放一份递质。在海马突触上，一半以上的脉冲根本到不了接收者，原因正是释放的随机性，而不是传导故障~\cite{AllenStevens1994}。对通信工程师来说，这是删除信道；两个神经元之间有多个突触，能部分挽回局面。

\emph{在活体皮层中，脉冲流看起来是随机的。}脉冲间隔的散布与随机泊松流大致相同；重复同一刺激时，窗口内脉冲数的变化使其方差接近均值~\cite{Softky1993}。

精确的元件怎么会产生随机的脉冲流？神经元接收数千个输入，每个都不可靠，而兴奋和抑制像第\ref{sec:gaba}章那样相互平衡。膜电压在阈值附近抖动，脉冲时刻由这些抖动决定。这是噪声，还是我们不会读的消息——在很大程度上仍是悬而未决的问题。一部分“噪声”已经被证明是消息：在视皮层中，相当大一部分散布跟随动物自身的运动~\cite{Stringer2019}。这里的困难是原则性的：压缩得好的消息，在不知道编码的人看来与随机流无法区分。

\section{频率：慢，但可靠}
\label{sec:rate}

最简单的编码是\emph{频率编码}：数记录在神经元在窗口 $T$ 内发出多少个脉冲之中。删除和时刻抖动都伤不到这种编码。但它有代价。如果脉冲流是泊松流，窗口内的脉冲数 $n$ 的均值为 $rT$，散布为 $\sqrt{rT}$：
\begin{empheq}[box=\keybox]{equation}
\frac{\sigma_n}{\bar n} \;=\; \frac{1}{\sqrt{r\,T}} .
\label{eq:rateerror}
\end{empheq}
要以百分之十的精度测出频率，需要一百个脉冲，每秒二十个脉冲时就是五秒。相邻电平相隔两个散布左右才可分辨，所以在时间 $T$ 内、频率不超过 $r$ 时，大约只能分辨 $\sqrt{rT}$ 个电平：$r=10$、$T=1\,\text{s}$ 时是三到四个电平，不到两比特。
大脑不会工作得这么慢。人在一瞥之间就能认出图片上的动物，大脑中的响应大约在 $150\ms$ 后出现~\cite{Thorpe1996}。粗略估计，在这段时间里信号要经过十来级相继的处理，每级 $10$--$20\ms$。按皮层通常的频率，每个神经元在一级处理中只来得及发出零个或一个脉冲，它的频率在这样的窗口里是无定义的。出路有两条：用很多神经元，或者看时间。

\emph{很多神经元。}设 $N$ 个神经元携带同一个数，接收者把它们的脉冲加起来。如果它们的噪声相互独立，\eqref{eq:rateerror} 中的 $rT$ 就换成 $NrT$：一百个神经元在 $r=20$ 时，$15\ms$ 内给出三十个脉冲，精度约百分之二十。空间代替了时间。然而皮层中相邻神经元的噪声并不独立：一对相邻神经元脉冲数的相关系数通常约为 $0.1$~\cite{Zohary1994}。若相关系数为 $c$，$N$ 个神经元平均值的方差为
\begin{empheq}[box=\keybox]{equation}
\sigma^2_N \;=\; \sigma^2_1\,\frac{1+(N-1)\,c}{N}
\;\xrightarrow[N\to\infty]{}\; c\,\sigma^2_1 .
\label{eq:correlated}
\end{empheq}

\begin{keyblock}
\begin{proposition}[平均的极限]
\label{prop:pooling}
在噪声相关时，对一群神经元求平均，误差最多只能减小到 $1/\sqrt{c}$ 分之一。$c=0.1$ 时是三分之一，而这个收益在几十个神经元时就已达到；再多加神经元也没有用。
\end{proposition}
\end{keyblock}

并非所有相关都同样有害。限制精度的只是共同噪声中\emph{与信号本身相似}的那一部分，即像被测量的变化那样推动整个神经元群的那部分噪声~\cite{MorenoBote2014}。大脑中这样的噪声实际上有多少，仍在研究之中。

利用很多神经元还有另一种方式：把数记录在\emph{哪一个}神经元发放上，就像声源方向检测器那样（图~\ref{fig:itd}）。这是\emph{位置}编码，是已知最可靠的编码：在感觉神经元中，导线的编号告诉我们皮肤的哪一点被触碰、声音的音高是多少，这已被反复证实。它在神经元数目上代价高昂，但速度快：一条消息只需一个延迟。

\section{时间：快，但从何时算起}

第二条出路是把数写进脉冲的时间。取神经元~\eqref{eq:glutneuron}，从时刻 $t=0$ 起给它一个恒定输入，这个输入本可把电压抬到 $U$。电压按 $V=U\bigl(1-e^{-t/\tau}\bigr)$ 上升，在时刻
\begin{empheq}[box=\keybox]{equation}
t_1 \;=\; \tau\,\ln\frac{U}{U-\theta}, \qquad U>\theta
\label{eq:latency}
\end{empheq}
达到阈值 $\theta$。强输入，脉冲早；弱输入，脉冲晚；阈下输入，没有脉冲。$\tau=10\ms$ 时，输入 $U=4\theta$ 在 $2.9\ms$ 后给出第一个脉冲，$U=2\theta$ 在 $6.9\ms$ 后，$U=1.2\theta$ 在 $17.9\ms$ 后。仅仅一个脉冲就携带了一个数。

但 $t_1$ 从什么时刻算起？没有公共的计时基准，接收者不知道刺激何时开始。有三种答案。

\emph{相对延迟。}两个神经元看到同一个刺激，它们的延迟差 $t_1^A-t_1^B$ 只取决于各自的输入，与开始时刻无关。比较时刻的是带延迟的符合检测器（图~\ref{fig:itd}）。如果 $N$ 个神经元各发出一个脉冲，它们到达的\emph{顺序}本身就携带多达 $\log_2N!$ 比特：十个神经元约二十二比特，而“发放与否”的答案最多给出十比特。在视网膜中已经证明，根据其输出神经元首个脉冲的相对延迟，可以快速而可靠地重建图像~\cite{Gollisch2008}。

\emph{齐射作为词首。}刺激的突变使大量神经元几乎同时发放。这样的齐射本身就是起点标记，就像摩尔斯电码中词与词之间的长停顿，接收者从它开始计算延迟。

\emph{局部节律。}在第\ref{sec:gaba}章中，\nt{GLUT}--\nt{GABA}环路产生了周期为 $12$--$50\ms$ 的局部节律。这不是时钟，但在它的一个周期内，输入强的神经元比输入弱的神经元先发放，于是~\eqref{eq:latency} 变成\emph{相位}编码：数记录在脉冲落在周期的哪一段。这种编码已被观察到：负责空间中特定位置的神经元，随着动物穿过“自己的”位置，在慢速局部节律的周期中发放得越来越早~\cite{OKeefe1993}。大脑在多大范围内使用相位，目前还不清楚。

\section{编码由接收者选择}

一串脉冲中既有频率，也有时间，还有顺序。读出其中哪一种，由接收者决定，而且只由一个参数决定——它的时间常数 $\tau$。

对方程~\eqref{eq:glutneuron} 做时间平均。如果神经元收到频率为 $r_j$ 的独立脉冲流，则平均电压及其相对抖动为
\begin{empheq}[box=\keybox]{equation}
\bar V \;=\; \tau\sum_j w_j\,r_j ,
\qquad
\frac{\sigma_V}{\bar V} \;\approx\; \frac{1}{\sqrt{2\,r_{\text{in}}\,\tau}} ,
\label{eq:reader}
\end{empheq}
其中第二个等式是对相等权重写的，$r_{\text{in}}$ 是全部输入的总频率。$\tau$ 大时，电压几乎不抖动，复现频率的加权和：接收者是\emph{积分器}，它读出频率而忘掉时间。$\tau$ 小时，电压在脉冲之间来得及回落，只有当几个脉冲在窗口~\eqref{eq:window} 内到达时才能达到阈值：接收者是\emph{符合检测器}，它读出时间（图~\ref{fig:reader}）。一千个输入、每个每秒五个脉冲，$\tau=10\ms$ 时抖动约百分之十，几乎是纯粹的积分。如果像第\ref{sec:gaba}章那样，抑制把 $\tau$ 缩短到 $2\ms$，抖动就增大到百分之二十以上，神经元开始注意到符合。

\begin{figure}[t]
\centering
\begin{tikzpicture}[>=Stealth,x=0.115cm,y=1cm]
% common input: the same spike train for both receivers
\foreach \s in {6,21,22,38,52,55.5,59,62.5,66,69.5,73,76.5,80,83.5,87,90.5,94,97.5}{\draw[thick,gray!40!black] (\s,5.25) -- (\s,5.65);}
\draw[gray!70!black] (0,5.25) -- (105,5.25);
\node[left,font=\small] at (-1,5.45) {输入};
\node[above,font=\scriptsize,gray!40!black] at (13.5,5.65) {稀疏};
\node[above,font=\scriptsize,gray!40!black] at (21.5,6.0) {一对};
\node[above,font=\scriptsize,gray!40!black] at (75,5.65) {密集};
% axes and thresholds
\foreach \y/\th in {2.7/2.5,0/1.5}{
  \draw[gray!70!black] (0,\y) -- (106,\y) node[right,black] {\small $t$};
  \draw[densely dashed,gray!60!black] (0,\y+0.6*\th) -- (104,\y+0.6*\th) node[right,black,font=\small] {$\theta$};
}
\node[anchor=south west,font=\small] at (0,4.3) {$\tau=10\ms$：积分器};
\node[anchor=south east,font=\small] at (104,1.0) {$\tau=2\ms$：符合检测器};
% integrator: tau=10 ms, theta=2.5w
\begin{scope}[yshift=2.7cm]
\foreach \a/\b/\v in {0/6/0.000,6/21/1.000,21/22/1.223,22/38/2.107,38/52/1.425,52/55.5/1.351,55.5/59/1.952,59/62.5/2.376,62.5/66/0.000,66/69.5/1.000,69.5/73/1.705,73/76.5/2.201,76.5/80/0.000,80/83.5/1.000,83.5/87/1.705,87/90.5/2.201,90.5/94/0.000,94/97.5/1.000,97.5/104/1.705}{\draw[very thick,blue!60!black] plot[domain=\a:\b,samples=14] (\x,{0.6*\v*exp(-(\x-\a)/10)});}
\foreach \s/\p/\q in {6/0.000/1.000,21/0.223/1.223,22/1.107/2.107,38/0.425/1.425,52/0.351/1.351,55.5/0.952/1.952,59/1.376/2.376,62.5/1.674/2.500,62.5/2.500/0.000,66/0.000/1.000,69.5/0.705/1.705,73/1.201/2.201,76.5/1.551/2.500,76.5/2.500/0.000,80/0.000/1.000,83.5/0.705/1.705,87/1.201/2.201,90.5/1.551/2.500,90.5/2.500/0.000,94/0.000/1.000,97.5/0.705/1.705}{\draw[very thick,blue!60!black] (\s,0.6*\p) -- (\s,0.6*\q);}
\foreach \s in {62.5,76.5,90.5}{\draw[->,thick,red!70!black] (\s,1.58) -- (\s,1.95);}
\end{scope}
% coincidence: tau=2 ms, theta=1.5w
\begin{scope}[yshift=0cm]
\foreach \a/\b/\v in {0/6/0.000,6/21/1.000,21/22/1.001,22/38/0.000,38/52/1.000,52/55.5/1.001,55.5/59/1.174,59/62.5/1.204,62.5/66/1.209,66/69.5/1.210,69.5/73/1.210,73/76.5/1.210,76.5/80/1.210,80/83.5/1.210,83.5/87/1.210,87/90.5/1.210,90.5/94/1.210,94/97.5/1.210,97.5/104/1.210}{\draw[very thick,blue!60!black] plot[domain=\a:\b,samples=14] (\x,{0.6*\v*exp(-(\x-\a)/2)});}
\foreach \s/\p/\q in {6/0.000/1.000,21/0.001/1.001,22/0.607/1.500,22/1.500/0.000,38/0.000/1.000,52/0.001/1.001,55.5/0.174/1.174,59/0.204/1.204,62.5/0.209/1.209,66/0.210/1.210,69.5/0.210/1.210,73/0.210/1.210,76.5/0.210/1.210,80/0.210/1.210,83.5/0.210/1.210,87/0.210/1.210,90.5/0.210/1.210,94/0.210/1.210,97.5/0.210/1.210}{\draw[very thick,blue!60!black] (\s,0.6*\p) -- (\s,0.6*\q);}
\foreach \s in {22}{\draw[->,thick,red!70!black] (\s,0.98) -- (\s,1.35);}
\end{scope}
\foreach \x in {0,50,100}{\draw[gray!60!black] (\x,-0.06) -- (\x,0.06) node[below,font=\scriptsize] {\x\,ms};}
\end{tikzpicture}
\caption{同一个输入，两种不同的接收者。每个脉冲把电压抬高 $w$，之后电压像神经元模型~\eqref{eq:glutneuron} 那样回落；红色箭头是接收者的脉冲。慢接收者（$\tau=10\ms$，$\theta=2.5w$）在脉冲\emph{多}的地方发放，注意不到那一对。快接收者（$\tau=2\ms$，$\theta=1.5w$）只对窗口~\eqref{eq:window} 内的一对脉冲发放，而放过密集却不符合的脉冲流。}
\label{fig:reader}
\end{figure}

测量表明，在活跃的皮层中，膜电导比静息时增大数倍~\cite{Destexhe2003}。因此那里的 $\tau$ 较短，工作状态下的皮层神经元更接近符合检测器，而不是积分器。

\begin{keyblock}
\begin{proposition}[编码由接收者决定]
\label{prop:reader}
同一串脉冲，对慢接收者携带的是频率，对快接收者携带的是时间，对被释放了递质的组织体积携带的是在数秒内平滑过的浓度水平（第\ref{sec:serocomputer}章）。使用哪种编码，不由发送者决定，而由接收者的时间常数决定，而这个时间常数又由抑制来调节。
\end{proposition}
\end{keyblock}
\section{点与划：作为滤波器的突触}

到目前为止，突触在我们这里只是带权重的导线。实际上，它的响应取决于先前的脉冲，这给神经元的语言带来了第二个符号。

\emph{压抑：突触传送变化。}每次释放都会消耗随时可释放的递质储备的一个份额 $U$~\cite{Tsodyks1997}，储备以数百毫秒量级的时间常数 $\tau_{\text{rec}}$ 恢复。频率为 $r$ 时，每个脉冲的响应幅度大致稳定在 $A(r)=A_0/(1+U r\tau_{\text{rec}})$，其中 $A_0$ 是充分休息的突触的响应。接收者在单位时间内得到的电流为
\begin{empheq}[box=\keybox]{equation}
r\,A(r) \;=\; \frac{A_0\,r}{1+U r\,\tau_{\text{rec}}}
\;\xrightarrow[r\to\infty]{}\; \frac{A_0}{U\tau_{\text{rec}}} .
\label{eq:depression}
\end{empheq}
当 $U=0.5$、$\tau_{\text{rec}}=0.8\,\text{s}$ 时，从每秒 $1/(U\tau_{\text{rec}})=2.5$ 个脉冲起就开始饱和。如果频率从五升到二十，稳态电流只增加三分之一。但新频率下的最初几个脉冲到来时，储备仍是原来的水平，电流瞬间跃升到四倍，然后在 $\tau_{\text{rec}}/(1+Ur\tau_{\text{rec}})\approx90\ms$ 内回落（图~\ref{fig:depression}）。

这样的突触传送的不是频率，而是频率的\emph{变化}，本质上是相对变化 $\Delta r/r$~\cite{Abbott1997depression}。对电子工程师来说，这是直接装在导线里的高通滤波器。

\begin{figure}[t]
\centering
\begin{tikzpicture}[>=Stealth,x=12.5cm,y=1cm]
% input rate
\draw[->,gray!70!black] (0,2.6) -- (0.9,2.6) node[right,black] {\small $t$};
\draw[->,gray!70!black] (0,2.6) -- (0,4.3) node[above,black] {\small $r$};
\draw[very thick,blue!60!black] (0,2.9) -- (0.2,2.9) -- (0.2,3.8) -- (0.8,3.8);
\node[left,font=\scriptsize] at (0,2.9) {$5$};
\node[left,font=\scriptsize] at (0,3.8) {$20$};
\node[right,font=\small,blue!60!black] at (0.4,4.1) {突触输入端的频率};
% output current
\draw[->,gray!70!black] (0,0) -- (0.9,0) node[right,black] {\small $t$};
\draw[->,gray!70!black] (0,0) -- (0,2.45) node[left,black] {\small $rA$};
\draw[densely dashed,gray!60!black] (0,0.667) -- (0.8,0.667);
\draw[very thick,red!70!black] (0,0.5) -- (0.2,0.5) -- (0.2,2.0)
   plot[domain=0.2:0.8,samples=60] (\x,{0.667+(2.0-0.667)*exp(-(\x-0.2)/0.089)});
\node[left,font=\scriptsize] at (0,0.5) {$1.7$};
\node[left,font=\scriptsize] at (0,2.0) {$6.7$};
\node[right,font=\scriptsize] at (0.8,0.667) {$2.2$};
\node[right,font=\small,red!70!black] at (0.28,1.6) {送往接收者的电流：跃升后在 $\approx90\ms$ 内回落};
\draw[gray!60!black] (0.2,-0.08) -- (0.2,0.08); \node[below,font=\scriptsize] at (0.2,0) {$0.2\,\text{s}$};
\draw[gray!60!black] (0.8,-0.08) -- (0.8,0.08); \node[below,font=\scriptsize] at (0.8,0) {$0.8\,\text{s}$};
\end{tikzpicture}
\caption{按公式~\eqref{eq:depression} 工作的压抑型突触，$U=0.5$，$\tau_{\text{rec}}=0.8\,\text{s}$；电流以每秒 $A_0$ 为单位，虚线是稳态电流：它只从 $1.7$ 升到 $2.2$，而在跃变瞬间达到 $6.7$。突触告诉接收者的是频率变了，而不是频率是多少。}
\label{fig:depression}
\end{figure}

\emph{易化：簇放电就是划。}另一些突触的释放概率很小，但每个脉冲之后会在几十毫秒内升高。单个脉冲经过这样的突触多半会丢失。而\emph{簇放电}——间隔几毫秒的一串连续脉冲——几乎一定能通过。即使没有易化，簇放电中 $k$ 个脉冲全部丢失的概率也只有
\begin{empheq}[box=\keybox]{equation}
P_{\text{loss}} \;=\; (1-p)^k ,
\label{eq:burst}
\end{empheq}
$p=0.2$ 时，单个脉冲为 $0.8$，四个脉冲的簇放电为 $0.41$；易化使它更小。于是神经元有了第二个符号：单个脉冲是点，簇放电是划。压抑型突触最善于传送停顿后的第一个脉冲；易化型突触放行簇放电而压制单个的点。簇放电传送得更可靠，这是测量过的；大脑确实把它当作一个独立的符号来用，则是一个可信的假说~\cite{Lisman1997}。

\section{\nt{GABA}神经元怎样说话}

皮层中数量最多的\nt{GABA}神经元是快速型的~\cite{Tremblay2016}。它们的脉冲比\nt{GLUT}神经元的短，能以每秒几百次的频率长时间均匀发放，几乎不疲劳。它们的突触更可靠：与每个接收者的连接由许多释放位点组成，脉冲几乎从不丢失。它们的输出落在接收者产生脉冲的位置附近，所以它们控制的不是接收者如何对输入求和，而是接收者是否发放、何时发放。

它们自己从许多邻近的\nt{GLUT}神经元得到兴奋，报告周围的总活动——这正是第\ref{sec:gaba}章的除法所需要的。最后，快速\nt{GABA}神经元彼此相连，容易一起发放；上面谈到的局部节律正是由它们设定的~\cite{Cardin2009}。

用摩尔斯电码的语言说，\nt{GABA}神经元负责的不是字母，而是\emph{停顿}。它们把连续的\nt{GLUT}脉冲流切成接收者可以发放的窗口，收窄符合窗口，并在脉冲流太响时把整体音量压低。

\begin{keyblock}
\begin{proposition}[角色分工]
\label{prop:punctuation}
\nt{GLUT}神经元携带消息的内容：\emph{什么}和\emph{多少}。快速\nt{GABA}神经元携带消息的标记：\emph{何时}可以说话、说得\emph{多响}；另有一类神经元通过抑制抑制性神经元，决定闸门\emph{为谁}打开。这种分工的各个部分已有测量；作为一般规则，它是一个工作假说。
\end{proposition}
\end{keyblock}

还有其他较慢的\nt{GABA}神经元，它们的输出落在接收者的个别输入上。它们像第\ref{sec:gaba}章的否决~\eqref{eq:veto} 那样工作：关闭一路\nt{GLUT}脉冲流，而不影响其余各路。

把\nt{GABA}神经元分为快速和慢速，是一幅旧图景，而且并不完整。现在皮层中至少区分出三大类~\cite{Tremblay2016}，而第三类的构造出人意料：它抑制的不是\nt{GLUT}神经元，而是其他\nt{GABA}神经元，首先是那些关闭输入的神经元~\cite{Pfeffer2013}。抑制的抑制，就是双重否定。这样的神经元不向接收者发送任何兴奋性脉冲，就能\emph{打开}闸门。开启它的是来自皮层其他区域的信号和广播通道，所以它像整片网络的使能输入：它活跃时，否决被解除，脉冲流得以通过~\cite{Pi2013}。在附录~\ref{sec:othermediators} 中，这一手法称为去抑制。
\section{节俭的语言}

关于神经元的语言，最后一件确知的事是：它很昂贵。一个脉冲连同它引起的突触活动，大约要消耗 $10^9$ 个ATP分子（对啮齿类皮层的估计~\cite{Attwell2001}），而一个分子提供约 $10^{-19}$ 焦耳。因此一个脉冲的代价约为 $E_{\text{spike}}\sim10^{-10}$ 焦耳。整个大脑消耗约 $20$ 瓦，如果一半用于传送信号，$P\sim10$ 瓦，那么单个神经元的平均频率为
\begin{empheq}[box=\keybox]{equation}
\bar r \;\approx\; \frac{P}{N\,E_{\text{spike}}}
\;\sim\; \frac{10\ \text{W}}{8.6\cdot10^{10}\cdot10^{-10}\ \text{J}}
\;\approx\; 1\ \text{个脉冲每秒}.
\label{eq:energy}
\end{empheq}
对皮层更细致的估计给出的数值还要小~\cite{Lennie2003}。大脑的编码必须是\emph{稀疏}的：只有一小部分神经元能快速工作。

从~\eqref{eq:spikeentropy} 可以看出，这并不那么糟。精度为 $1\ms$ 时，以每秒 $100$ 次工作的神经元，每个脉冲最多携带五比特，而每秒发放一次的神经元，每个脉冲可携带多达十一比特。如果要为脉冲付费，那么少说话更划算。皮层确实用精度换能量：食物不足时，神经元保持脉冲频率不变，但在突触电流上花得更少，它们的响应随之变得不那么精确~\cite{Padamsey2022}。

在摩尔斯电码中，常用字母很短：英文文本中最常见的字母E只是一个点。据目前所知，神经元的编码以另一种形式遵循同样的规则：预料之中的事只花很少的脉冲~\cite{Barlow1961}。神经元在恒定刺激下逐渐降低频率，第\ref{sec:glutcomputer}章的传感器响应的是变化而不是电平，第\ref{sec:neurontypes}章的\nt{DOPA}神经元只报告奖励预测误差。

\begin{keyblock}
\begin{proposition}[节约]
\label{prop:economy}
能量把神经元的平均频率限制在每秒约一个脉冲的量级。因此\nt{GLUT}神经元的语言是稀疏的，把脉冲花在新闻上：花在变化了的或未被预测的事情上。能量上限是测量过的；大脑的编码被调整到每焦耳比特数最多，则是一个论据充分的假说。
\end{proposition}
\end{keyblock}

\section{小结}

\begin{table}[t]
\centering
\footnotesize
\setlength{\tabcolsep}{4pt}
\renewcommand{\arraystretch}{1.15}
\begin{tabular}{>{\raggedright}p{2.6cm}>{\raggedright}p{2.3cm}>{\raggedright}p{3.0cm}>{\raggedright}p{2.0cm}>{\raggedright\arraybackslash}p{3.6cm}}
\toprule
记录方式 & 携带什么 & 谁来读 & 每条消息的时间 & 已知情况 \\
\midrule
发放神经元的编号（位置编码） & 哪个值 & 任何接收者；符合检测器 & 一个延迟 & 在感觉神经元和声源定向中已证实 \\
单个神经元的频率 & 量值 & $\tau$ 大的积分器；组织体积（第\ref{sec:serocomputer}章） & 数百毫秒到数秒 & 已证实；对 $10$--$20\ms$ 的一级处理而言太慢 \\
神经元群的频率 & 量值 & 有数千个输入的神经元 & $10$--$50\ms$ & 已证实；收益受相关性限制~\eqref{eq:correlated} \\
首个脉冲的时间、顺序 & 量值、顺序 & 带延迟的符合检测器 & 一个延迟 & 在神经系统入口处已证实；在皮层中尚有争论 \\
局部节律中的相位 & 量值、位置 & 有共同节律的神经元 & 一个周期 $12$--$50\ms$ 或更慢 & 在个别脑区观察到；普遍作用是假说 \\
簇放电（“划”） & “重要”：可靠投递 & 易化型突触 & $\sim10\ms$ & 可靠性已测量；作为独立符号是假说 \\
频率的变化 & 新闻 & 压抑型突触 & $\sim0.1\,\text{s}$ & 已证实 \\
快速\nt{GABA}神经元的脉冲 & 何时、多响 & 所有邻近的\nt{GLUT}神经元 & $1$--$30\ms$ & 节律和除法已测量；作为规则的“标点”是假说 \\
\bottomrule
\end{tabular}
\caption{\nt{GLUT}和\nt{GABA}神经元用脉冲记录消息的方式。右列把测量到的与推测的区分开来。}
\label{tab:code}
\end{table}

表~\ref{tab:code} 汇集了我们已知的东西；从中也能看出我们不知道的东西。我们不知道皮层在多大程度上利用脉冲的精确时间，而不仅仅是神经元群的频率。我们不知道神经元是否有“词”——许多神经元脉冲的稳定组合，作为一个整体被读出。我们也不知道大脑不同部位的语言是否相同。摩尔斯电码一个晚上就能学会，因为它的码表是已知的。神经元语言的码表还有待编写，而编写它的，多半会是通信工程师。

\section{习题}

\begin{exercise}[\lvlA]
\label{ex:04:1}
\emph{容量的数值。}$\Delta t=1\ms$。(a)~按~\eqref{eq:energy} 中的脉冲代价，频率为每秒 $1$ 个脉冲的神经元每焦耳给出多少比特？(b)~$r=10$ 时，只数每秒脉冲数的接收者提取的信息比上限~\eqref{eq:spikeentropy} 少多少倍？
\end{exercise}

\begin{exercise}[\lvlB]
\label{ex:04:2}
\emph{最优频率。}神经元消耗功率 $P_0+rE_{\text{spike}}$，其中 $P_0$ 是 $20$~W中的另一半除以全部神经元数，$E_{\text{spike}}=10^{-10}$~J。按~\eqref{eq:spikeentropy}、取 $\Delta t=1\ms$，频率 $r$ 为多少时每焦耳比特数 $R_{\max}(r)/(P_0+rE_{\text{spike}})$ 最大？
\end{exercise}

\begin{exercise}[\lvlB]
\label{ex:04:3}
\emph{什么样的相关有害。}$N=1000$ 个神经元，方差 $\sigma^2=1$，两两相关系数 $c=0.1$。分别对相同斜率 $f'=\mathbf 1$ 和满足 $\mathbf 1^{\mathsf T}f'=0$ 的斜率 $\pm1$，计算费希尔信息 $\mathcal I=f'^{\mathsf T}\Sigma^{-1}f'$（$\Sigma$ 为协方差矩阵）。两者相差多少倍？
\end{exercise}

\begin{exercise}[\lvlB]
\label{ex:04:4}
\emph{建模：符合检测器。}神经元~\eqref{eq:glutneuron} 接收 $1000$ 个泊松输入，每个每秒 $5$ 个脉冲，$w=1$；阈值 $\theta$ 取为使自由电压每秒越过它一次。用模型求出：$\tau=10$ 和 $2\ms$ 时，多少个同时到达的输入 $m$ 能以 $1/2$ 的概率点燃神经元？
\end{exercise}

\begin{exercise}[\lvlB]
\label{ex:04:5}
\emph{设计：相对变化检测器。}选择突触~\eqref{eq:depression} 的参数：每秒 $40$ 个脉冲时的稳态电流比每秒 $10$ 个时高出不超过 $10\,\%$，而频率跃升到 $40$ 后，电流以不超过 $50\ms$ 的时间常数达到新水平。最小的 $\tau_{\text{rec}}$ 是多少？对应的 $U$ 呢？
\end{exercise}

\begin{exercise}[\lvlA]
\label{ex:04:6}
\emph{首个脉冲时间与簇放电。}(a)~由~\eqref{eq:latency} 求出使首个脉冲恰好在一个时间常数后到来的输入。它取决于什么？(b)~释放概率 $p=0.2$ 时，簇放电中要有多少个脉冲，才能使全部丢失的概率低于 $5\,\%$？
\end{exercise}

\begin{exercise}[\lvlC]
\label{ex:04:7}
\emph{研究：排列中有多少比特。}神经元是相互独立的格子~\eqref{eq:spikeentropy}，$r=10$，$\Delta t=1\ms$。(a)~把 $20$ 个格子组成的“词”的熵分解为脉冲数的熵与剩余部分。(b)~模拟用 $N=30$ 和 $10^4$ 条记录的词频来估计熵。估计值偏低多少？
\end{exercise}

\chapter{\nt{SERO}神经元计算什么}
\label{sec:serocomputer}

\section{第一个广播通道}

到目前为止，书中所有的神经元都通过寻址总线通信：\nt{GLUT}和\nt{GABA}把脉冲发给确定的接收者，我们也弄清了这样的网络计算什么、用什么语言说话。现在转向第~\ref{sec:buses}~节中的第二条总线——广播总线。它的第一个通道是\nt{SERO}。和以前一样，我们只允许自己使用活体中真实存在的东西。

本章会出现三种器件，以后所有广播通道都会用到它们：一种神经元，在平稳的背景输入下自己工作，直到被叫停；一根导线，实际上是一块组织体积；以及用缓慢的浓度代替单个脉冲。第~\ref{sec:dofa}--\ref{sec:acet}~章中的\nt{DOPA}、\nt{NORA}和\nt{ACET}也由同样的零件搭成。

从工程师的角度看，\nt{SERO}神经元有四个重要的不同之处。

\emph{第一：符号由接收者选择。}血清素的受体两种符号都有，规则~\eqref{eq:dalesign}在这里不适用：符号不由发送者决定，而由接收者的受体决定（命题~\ref{prop:receptor}）~\cite{BarnesSharp1999}。

\emph{第二：只要有平稳的背景输入，\nt{SERO}神经元就自己工作。}清醒时它平稳地每秒发放几个脉冲，并不需要每个脉冲都有一次推动：这是一个起搏器。但它需要持续的背景输入。在脑片中没有这种输入，大多数这类细胞保持沉默；如果通过$\alpha_1$受体给予去甲肾上腺素，平稳的节律就会恢复，而阻断这些受体则使节律熄灭~\cite{VandermaelenAghajanian1983}。在体内，这个背景来自\nt{NORA}。对电路而言，这是一个需要供电的振荡器：只要有电，它就自己工作，直到被抑制叫停。

\emph{第三：它很慢。}血清素受体几乎都是代谢型的：它们自己不打开通道，而是在细胞内启动一连串反应。响应很慢：经主要受体的抑制性电流在大约一秒内升起又衰减~\cite{CourtneyFord2016}——比\nt{GLUT}快突触的响应长几十倍。

\emph{第四：它把神经递质释放到体积中，而不是导线里。}在\nt{SERO}神经元输出到达的区域，血清素常常不是释放到狭窄的突触间隙，而是释放到细胞间隙中并向四周扩散~\cite{Bunin1998}。接收者感受到的是浓度，无法知道分子来自哪个发送者。

在这样的时间尺度上，单个脉冲已经不重要，因此我们用频率$r_j$和浓度来描述网络。神经元$j$释放递质的那个体积中，血清素浓度$c$因释放而上升，又因递质被泵回而下降：
\begin{empheq}[box=\keybox]{equation}
\tau_c\,\frac{\dd c}{\dd t} \;=\; -c \;+\; \kappa\sum_j r_j ,
\qquad
r_i \;=\; f_i\bigl(b_i + s_i\,g\,c\bigr), \quad s_i=\pm1 .
\label{eq:seroneuron}
\end{empheq}
这是读取脉冲流的又一种方式，补充了命题~\ref{prop:reader}：体积既看不到单个脉冲，也看不到它们的顺序，只看到在时间$\tau_c$内平均的频率。第一个方程就是与膜相同的RC电路：$\tau_c$是递质在体积中的寿命，它没有被直接测量，我们取其量级为一秒；$\kappa$是一个脉冲贡献的递质量。第二个方程描述接收者：$b_i$是它自身的输入，起搏器的$b_i$为正（其中也包括来自\nt{NORA}的平稳背景输入）；$g$是受体的灵敏度；符号$s_i$由接收者选择；$f_i$是不减的传递特性。

\section{一个神经元实现“非”}

取一个带抑制性受体的起搏器，$s=-1$。周围没有血清素时，它靠自身输入$b$工作。血清素出现后，输入减少$gc$，若$gc>b$，神经元就沉默。这就是“非”：输入没有时，输出才有。它只用了一个神经元（图~\ref{fig:serogates}）。命题~\ref{prop:monotone}在这里不适用：它要求权重为正。

如果同一个起搏器受到来自两个不同体积的血清素作用，那么只要其中至少一个有血清素，它就沉默。这是“或非”，而用“或非”可以搭出任何逻辑函数。这里不需要\nt{GLUT}网络的双导线编码：信号的缺失由那些一直工作、直到被叫停的神经元来计算。

\begin{figure}[t]
\centering
\begin{tikzpicture}[>=Stealth,
  nrn/.style={circle,draw,very thick,fill=blue!6,minimum size=0.95cm,inner sep=0pt},
  pace/.style={nrn,double,double distance=1.2pt},
  inh/.style={-{Bar[width=7pt]},thick,red!70!black}]
\node[pace] (N1) at (0,0) {};
\draw[inh] (-1.6,0) node[left,black] {$c$} -- (N1);
\draw[->,very thick,red!70!black] (N1) -- (1.3,0) node[right,black] {$\bar c$};
\node[below=0.55cm] at (0,0) {\small 非};
\node[pace] (N2) at (5,0) {};
\draw[inh] (3.4,0.45) node[left,black] {$c_1$} -- (N2);
\draw[inh] (3.4,-0.45) node[left,black] {$c_2$} -- (N2);
\draw[->,very thick,red!70!black] (N2) -- (6.3,0) node[right,black] {$\overline{c_1\vee c_2}$};
\node[below=0.55cm] at (5,0) {\small 或非};
\end{tikzpicture}
\caption{由\nt{SERO}神经元构成的逻辑。双圆圈是起搏器：在平稳的背景输入下它自己工作，直到被抑制。带横杠的线是抑制性受体，$s=-1$；$c$、$c_1$、$c_2$是接收者周围各体积中的血清素浓度。左为“非”，右为“或非”。}
\label{fig:serogates}
\end{figure}

\begin{keyblock}
\begin{proposition}[\nt{SERO}逻辑的完备性]
\label{prop:serocomplete}
如果网络中有带抑制性受体的起搏器，且释放到不同体积中的递质互不混合，那么网络~\eqref{eq:seroneuron}可以计算任何逻辑函数，而且每一位只用一根导线传送。
\end{proposition}
\end{keyblock}

在逻辑上，\nt{SERO}网络比\nt{GLUT}更强，但“释放到不同体积中的递质互不混合”这一条件在脑中并不成立（第~\ref{sec:volume}~节）。

\section{负反馈}

血清素神经元自身带有抑制性受体~\cite{Sprouse1987}。它们释放的血清素回到它们自己身上，使它们放慢。对工程师来说这是熟悉的电路：带负反馈的放大器（图~\ref{fig:serofeedback}）。

哪些是测量结果，哪些是模型。在\nt{SERO}神经元所在的中缝核内部，血清素抑制邻居并不经过公共体积，而是通过突触点对点进行：转运体迅速把它清除，不让它在间隙外积累。单次释放只造成放电中的短暂停顿，平均频率并不因此改变~\cite{CourtneyFord2016}。因此，下面这个经公共体积闭合的回路是一个模型：我们计算的是，如果这种反馈在平均频率的层面上工作，它会做什么。

\begin{figure}[t]
\centering
\begin{tikzpicture}[>=Stealth,
  nrn/.style={circle,draw,very thick,fill=blue!6,minimum size=0.95cm,inner sep=0pt},
  pace/.style={nrn,double,double distance=1.2pt},
  blk/.style={draw,very thick,rounded corners=2pt,fill=orange!10,minimum height=0.9cm,minimum width=2.2cm,align=center,font=\small},
  inh/.style={-{Bar[width=7pt]},thick,red!70!black}]
\node[pace] (N) at (0,0) {};
\node[blk] (V) at (3.6,0) {体积\\$\tau_c$};
\draw[->,very thick] (N) -- node[above] {\small $r$} (V);
\draw[->,very thick] (V) -- (6.0,0) node[right] {\small $c$——送往所有接收者};
\draw[inh] (V.south) -- ++(0,-0.9) -| node[pos=0.25,below] {\small $-g\,c$} (N.south);
\draw[->,thick,blue!60!black] (-1.7,0) node[left,black] {\small $b$} -- (N);
\end{tikzpicture}
\caption{经自身递质形成反馈的\nt{SERO}神经元：浓度$c$送往所有接收者，同时抑制神经元自己。在模型中这是一个电平调节器；在脑中回路是否起这种作用，还是假说。}
\label{fig:serofeedback}
\end{figure}

我们来算一算这个回路做什么。设传递特性是线性的，$u>0$时$f(u)=au$。平衡时$c=\kappa r$且$r=a(b-gc)$。把一个代入另一个，得到
\begin{empheq}[box=\keybox]{equation}
r \;=\; \frac{a\,b}{1+G}, \qquad G \;=\; a\,g\,\kappa .
\label{eq:feedback}
\end{empheq}
量$G$是环路增益。这和任何负反馈放大器的公式相同，并带来同样的两个好处。

\emph{对参数离散的稳健性。}设神经元的兴奋性$a$变化了比例$\varepsilon$。没有反馈时，频率也会变化同样的比例。有反馈且$G\gg1$时，频率$r\approx b/(g\kappa)$几乎与$a$无关：它的变化是无反馈时的$1+G$分之一。$G=9$时，离散被压低到十分之一。

\emph{速度。}把$r=a(b-gc)$代入第一个方程~\eqref{eq:seroneuron}，得到$\tau_c\,\dd c/\dd t=-(1+G)\,c+\kappa ab$。浓度以时间常数$\tau_c/(1+G)$趋近自己的水平——比没有反馈时快$1+G$倍。

这就是\nt{SERO}系统可能执行的第一种计算：把脑中血清素的总体水平保持在设定值上，就像恒温器保持温度一样。这是假说：在实验中，自抑制目前只表现为短暂的停顿，并不改变平均频率。

\section{从脉冲到电平}

第二种计算藏在第一个方程~\eqref{eq:seroneuron}中。神经元发放的是单个脉冲，而接收者感受到的是浓度，它在时间$\tau_c$内把脉冲平滑。如果源的频率变化，浓度会滞后地跟随它：
\begin{equation}
c(t) \;=\; \frac{\kappa}{\tau_c}\int_{-\infty}^{t} e^{-(t-t')/\tau_c}\,\sum_j r_j(t')\,\dd t' .
\label{eq:lowpass}
\end{equation}
这是源活动在最近几个$\tau_c$内的滑动平均——一个低通滤波器（图~\ref{fig:lowpass}）。最简单的数模转换器正是这样把脉冲送过RC电路，把它们的频率变成电平。\nt{SERO}系统同时对几十万个神经元做同样的事。

这个量同时也是记忆。源沉默之后，浓度还会维持量级为$\tau_c$的时间。这不是一个比特，而是一条平缓消退的痕迹。

\begin{figure}[t]
\centering
\begin{tikzpicture}[>=Stealth,x=1.45cm,y=1cm]
\draw[gray!70!black] (0,3.2) -- (8.1,3.2);
\node[left,font=\small] at (-0.05,3.4) {脉冲};
\node[above,font=\scriptsize,gray!40!black] at (1,3.65) {每秒$2$个};
\node[above,font=\scriptsize,gray!40!black] at (3.5,3.65) {每秒$8$个};
\node[above,font=\scriptsize,gray!40!black] at (6.5,3.65) {每秒$2$个};
\draw[->,gray!70!black] (0,0) -- (8.3,0) node[right,black] {\small $t$};
\draw[->,gray!70!black] (0,0) -- (0,2.75) node[left,black] {\small $c$};
\foreach \s in {0.136,0.529,0.760,1.223,1.714,1.748,1.755,2.663,2.701,2.734,3.414,3.493,3.719,3.800,3.928,3.948,4.074,4.327,4.420,4.589,4.728,4.736,4.914,5.025,5.205,5.220,6.224,6.544,7.178}{\draw[thick,blue!60!black] (\s,3.2) -- (\s,3.6);}
\foreach \a/\b/\v in {0.000/0.136/0.000,0.136/0.529/1.000,0.529/0.760/1.675,0.760/1.223/2.330,1.223/1.714/2.466,1.714/1.748/2.509,1.748/1.755/3.425,1.755/2.663/4.402,2.663/2.701/2.775,2.701/2.734/3.672,2.734/3.414/4.553,3.414/3.493/3.306,3.493/3.719/4.055,3.719/3.800/4.235,3.800/3.928/4.905,3.928/3.948/5.316,3.948/4.074/6.211,4.074/4.327/6.476,4.327/4.420/6.028,4.420/4.589/6.493,4.589/4.728/6.483,4.728/4.736/6.642,4.736/4.914/7.589,4.914/5.025/7.351,5.025/5.205/7.579,5.205/5.220/7.331,5.220/6.224/8.221,6.224/6.544/4.012,6.544/7.178/3.914,7.178/8.000/3.076}{\draw[very thick,orange!85!black] plot[domain=\a:\b,samples=10] (\x,{0.25*\v*exp(-(\x-\a))});}
\foreach \s/\p/\q in {0.136/0.000/1.000,0.529/0.675/1.675,0.760/1.330/2.330,1.223/1.466/2.466,1.714/1.509/2.509,1.748/2.425/3.425,1.755/3.402/4.402,2.663/1.775/2.775,2.701/2.672/3.672,2.734/3.553/4.553,3.414/2.306/3.306,3.493/3.055/4.055,3.719/3.235/4.235,3.800/3.905/4.905,3.928/4.316/5.316,3.948/5.211/6.211,4.074/5.476/6.476,4.327/5.028/6.028,4.420/5.493/6.493,4.589/5.483/6.483,4.728/5.642/6.642,4.736/6.589/7.589,4.914/6.351/7.351,5.025/6.579/7.579,5.205/6.331/7.331,5.220/7.221/8.221,6.224/3.012/4.012,6.544/2.914/3.914,7.178/2.076/3.076}{\draw[very thick,orange!85!black] (\s,0.25*\p) -- (\s,0.25*\q);}
\draw[thick,densely dashed,black] plot[domain=0:2,samples=30] (\x,{0.25*2*(1-exp(-\x))})
  plot[domain=2:5,samples=40] (\x,{0.25*(8-6.271*exp(-(\x-2)))})
  plot[domain=5:8,samples=40] (\x,{0.25*(2+5.688*exp(-(\x-5)))});
\foreach \x in {0,2,5,8}{\draw[gray!60!black] (\x,-0.05) -- (\x,0.05) node[below=3pt,font=\scriptsize] {\x\,s};}
\end{tikzpicture}
\caption{从脉冲到电平。上方是\nt{SERO}神经元的脉冲：两秒稀疏，三秒密集，然后又变稀疏。下方是$\tau_c=1\,\text{s}$时的血清素浓度~\eqref{eq:lowpass}。虚线是脉冲完全均匀时的同一曲线。}
\label{fig:lowpass}
\end{figure}

\section{导线就是体积}
\label{sec:volume}

回到命题~\ref{prop:serocomplete}的条件。不同发送者在附近释放的血清素会叠加成同一个浓度。

\emph{这里的导线不是纤维，而是一块组织体积。}两个信号只有在释放区域之间的距离大于递质来得及扩散的距离时，才能保持区分。在时间$\tau$内，扩散系数为$D$的分子会走出
\begin{empheq}[box=\keybox]{equation}
\ell \;\sim\; \sqrt{D\,\tau}\, .
\label{eq:diffusionlength}
\end{empheq}
细胞间隙的曲折使小分子的扩散减慢约$2.6$倍~\cite{Sykova2008}，而血清素本身在组织中的扩散系数没有测量过；按分子大小估计，它为几个$10^{-6}$~cm$^2$/s。如果信号存在$\tau\sim1$~s（也是估计），那么$\ell\sim\sqrt{3\cdot10^{-6}}$~cm $\approx20\um$。这种网络中的地址就是\emph{位置}：接收者只能凭自己所在的位置判断信号来自谁。

真实的\nt{SERO}神经元的构造使这些独立地址几乎不复存在。每个神经元的输出散布在脑的巨大区域上：一个细胞的分支能到达许多彼此相距很远的区域~\cite{GagnonParent2014}。每个细胞的释放位点估计约有$10^5$个（表~\ref{tab:types}）；直接计数尚未做过。不同\nt{SERO}神经元的“导线”相互重叠、彼此混合。不过它们并没有完全合并成一根导线：\nt{SERO}神经元分成几个子系统，把输出送往不同区域，对同一事件的响应也不同~\cite{Ren2018}。但这样的子系统只有几个，而不是几十万个。

\begin{keyblock}
\begin{proposition}[独立导线的数目]
\label{prop:volumewires}
在采用容积传递的网络中，独立信号的数目不受神经元数目限制，而受它们释放递质的、互不重叠的、尺寸为$\ell\sim\sqrt{D\tau}$的区域数目限制。如果每个神经元的输出散布在很大的体积中，整个网络只能传送少数几个公共变量。
\end{proposition}
\end{keyblock}

因此，在脑中没有材料来搭建命题~\ref{prop:serocomplete}中的逻辑电路。而调节器~\eqref{eq:feedback}和滤波器~\eqref{eq:lowpass}不需要单独的导线：一个公共量就够了。滤波器直接来自已测量到的、在目标区域向体积中的释放~\cite{Bunin1998}；电平调节器则是假说。

\section{规划视野}
\label{sec:horizon}

一个公共的慢变量能派什么用场？一个好的候选者是这样的参数：它对整个网络必须相同，变化不快于秒或分钟的尺度。在第~\ref{sec:neurontypes}~章中已经出现过这样的参数：误差公式~\eqref{eq:rpe}中的系数$\gamma$，它表示未来的奖励比眼前的奖励便宜多少。在连续时间中，它的位置由\emph{视野}$\tau_\gamma$占据：需要等待时间$D$才能得到的奖励$R$，现在值$R\,e^{-D/\tau_\gamma}$。

视野决定是否值得等待。假设可以立即拿到小奖励$r_0$，或者等待大奖励$R$。只要满足下式，等待就是划算的：
\begin{empheq}[box=\keybox]{equation}
D \;<\; \tau_\gamma\,\ln\frac{R}{r_0} .
\label{eq:patience}
\end{empheq}
当$\tau_\gamma=2\,\text{s}$、延迟奖励大三倍时，值得等到$2.2\,\text{s}$；若视野长一倍，则可等到$4.4\,\text{s}$。耐心与视野成正比。

公式~\eqref{eq:patience}依赖于指数折扣。在秒级延迟上测得的折扣更接近双曲线$R/(1+D/\tau_\gamma)$：猴子在选择中对延迟奖励的估值，以及\nt{DOPA}神经元对其预示信号的响应，都是这样下降的~\cite{KobayashiSchultz2008}。对双曲线，条件~\eqref{eq:patience}换成$D<\tau_\gamma\,(R/r_0-1)$：对奖励之比的依赖不再是对数的，而是线性的，在同样的数字下等待的界限几乎移动一倍。如果延迟是随机的，同一个指数也会给出双曲线（习题~\ref{ex:05:7}），而耐心与时间尺度成正比的关系依然保持。

按照一种假说，视野正是由\nt{SERO}设定的~\cite{Doya2002}。它具备担当这一角色的一切条件：整个网络共享一个公共水平，在几秒内变化。也有直接的实验：人为激活血清素神经元，动物会等待延迟奖励更久才放弃~\cite{Miyazaki2014}，在奖励变得稀少的地方也会更久地尝试获取奖励~\cite{Lottem2018}。血清素浓度在网络内部究竟怎样变成$\tau_\gamma$，目前还不知道。下一章中，视野将进入\nt{DOPA}广播的误差。

\section{小结}

\begin{table}[t]
\centering
\small
\begin{tabular}{>{\raggedright}p{3.6cm}>{\raggedright}p{4.3cm}>{\raggedright\arraybackslash}p{4.3cm}}
\toprule
& \nt{GLUT} & \nt{SERO} \\
\midrule
反应时间 & 毫秒 & 零点几秒到一秒 \\
作用符号 & 只有$+$ & $+$和$-$，由接收者选择 \\
没有输入时 & 沉默 & 在平稳背景输入下自己工作 \\
寻址 & 点对点 & 体积，地址即位置 \\
“非” & 只能通过“断开”传感器 & 一个神经元 \\
记忆 & 自持活动，因疲劳而熄灭 & 浓度，在时间$\tau_c$内消退 \\
主要计算 & 符合、延迟、逻辑、序列 & 平滑；电平调节与视野$\tau_\gamma$（假说） \\
\bottomrule
\end{tabular}
\caption{仅由\nt{GLUT}神经元和仅由\nt{SERO}神经元构成的网络各计算什么。}
\label{tab:glutvssero}
\end{table}

作为逻辑元件（表~\ref{tab:glutvssero}），\nt{SERO}神经元比\nt{GLUT}更丰富，但作为通信系统，它是一种广播：缓慢，几乎没有地址。因此它不去充当处理器，而是平滑并广播第~\ref{sec:neurontypes}~章谈到的几个公共量——按照假说，其中也包括规划视野。起搏器、体积和缓慢的浓度，我们还会再遇到三次：在\nt{DOPA}、\nt{NORA}和\nt{ACET}那里。

\section{习题}

\begin{exercise}[\lvlA]
\label{ex:05:1}
\emph{导线就是体积。}取血清素的$D=3\cdot10^{-6}$~cm$^2$/s（第~\ref{sec:volume}~节）。按~\eqref{eq:diffusionlength}求寿命为$1\,\text{s}$的信号的$\ell$，以及递质扩散到$5$~cm所需的时间。那么\nt{SERO}的广播性是靠扩散实现的，还是靠具有$\sim10^5$个释放位点的导线分支实现的？
\end{exercise}

\begin{exercise}[\lvlA]
\label{ex:05:2}
\emph{电平调节器。}证明在~\eqref{eq:seroneuron}中，当$f(u)=au$时，相对灵敏度$S=(a/r)\,\partial r/\partial a$恰好等于$1/(1+G)$，而不仅仅在$G\gg1$时成立。设$G=9$，兴奋性$a$增大了$20\,\%$。不作线性化，$r$会变化百分之几？
\end{exercise}

\begin{exercise}[\lvlB]
\label{ex:05:3}
\emph{回路中的慢受体。}\nt{SERO}受体的响应有延迟：$r(t)=a\bigl(b-gc(t-T)\bigr)$，体积按~\eqref{eq:seroneuron}。证明当$G>1$时，回路仅在$T<T^*=\tau_c\arccos(-1/G)/\sqrt{G^2-1}$时稳定。求$\tau_c=1\,\text{s}$时$G=2$和$G=9$的$T^*$。延迟为几百毫秒时，能实现多大的离散抑制？
\end{exercise}

\begin{exercise}[\lvlB]
\label{ex:05:4}
\emph{电平的散粒噪声。}按~\eqref{eq:lowpass}，每个脉冲给$c$加上一个$\tau_c=1\,\text{s}$的指数。若全部$3\cdot10^5$个\nt{SERO}神经元都以每秒$2$个的频率向同一体积释放泊松脉冲，求$c$的变异系数。对单个频率为每秒$4$个脉冲的神经元，多大的$\tau_c$能给出$\mathrm{CV}=10\,\%$？
\end{exercise}

\begin{exercise}[\lvlB]
\label{ex:05:5}
\emph{仿真：反馈对抗噪声。}仿真$N=50$个泊松起搏器，频率$r_i=a_i[b-gc]_+$，$a_i$在$[2.8,\,5.2]$上均匀分布，体积按~\eqref{eq:seroneuron}，$\kappa=1/N$，$\tau_c=1\,\text{s}$。与$b=0.1$、$g=0$相比，反馈（$b=1$，$g=2.25$）使电平$c$的CV降低多少倍？它能使各神经元的频率趋于一致吗？
\end{exercise}

\begin{exercise}[\lvlB]
\label{ex:05:6}
\emph{设计：用\nt{SERO}逻辑实现XOR。}门是一个起搏器，若$b-g\sum_kc_k>0$则输出$r_0$，否则输出$0$，送入自己的体积$\tau_c\,\dd c/\dd t=-c+\kappa r$；它计算“或非”。用五个这样的门搭出XOR，并求$\tau_c=1\,\text{s}$、$G'=g\kappa r_0/b=2$时它的建立时间。
\end{exercise}

\begin{exercise}[\lvlB]
\label{ex:05:7}
\emph{延迟未知时的耐心。}(a)~动物愿意等待三倍大的奖励，只要等待不超过$6\,\text{s}$。这意味着视野$\tau_\gamma$是多少？(b)~设延迟$D$服从均值为$\bar D$的指数分布。证明当$\bar D<\tau_\gamma(R/r_0-1)$时等待是划算的，并在$\tau_\gamma=2\,\text{s}$、$R=3r_0$时与固定延迟作比较。
\end{exercise}

\begin{exercise}[\lvlC]
\label{ex:05:8}
\emph{浓度如何设定视野。}\nt{DOPA}神经元直接以权重$k_1$接收$V$，又经过一个以$\tau_d=0.1\,\text{s}$平滑的\nt{GABA}中介以权重$-k_2$接收$V$；对慢变的$V$，总和为$(k_1-k_2)V+k_2\tau_d\dot V$。选取$k_1$、$k_2$，使之符合~\eqref{eq:ctd}且$\tau_\gamma=2\,\text{s}$。\nt{SERO}把$k_1$乘以$m$：要使视野加倍，$m$需要改变多少？
\end{exercise}

\chapter{会学习的神经网络：加入\nt{DOPA}}
\chaptermark{会学习的网络}
\label{sec:dofa}

\section{游戏规则}

在前几章的所有电路中，权重都是由工程师设定的。体内没有工程师。权重必须在工作过程中自己确定下来，并使网络做出有用的事。这就叫学习。

在原有规则之外再加一条。突触自己改变自己的权重，而它能知道的只有\emph{它附近}发生的事：发送者的活动$x$、接收者的活动$y$，以及能到达它的那些物质。谁也不能给某个突触单独寄来一份修正。

这就排除了人工神经网络的主要方法——误差反向传播。在那里，输出误差沿着同样的权重反向穿过网络，在前向传递之后的一个单独阶段进行，每个权重都得到自己的修正。这需要与前向导线精确对应的反向导线，以及统一的时钟节拍。二者在脑中都没有找到，至少没有以人工网络中的那种形式存在~\cite{Lillicrap2020,WhittingtonBogacz2019}。

我们把权重的变化写成一个缓慢的连续过程：
\begin{equation}
\frac{\dd w}{\dd t} \;=\; \eta\,\Phi(\text{突触所知道的信息}),
\label{eq:plasticity}
\end{equation}
其中$\eta$是学习速率，它小到在一个脉冲的时间内权重几乎不变。整章讨论的就是函数$\Phi$可以是什么样子。

\section{权重存储在哪里}
\label{sec:weightstore}

一个“改变自己权重”的突触是什么？连接有两侧：发送者导线的末梢（\emph{突触前}侧）和接收者带受体的一片膜（\emph{突触后}侧），二者之间是一条窄缝。在\nt{GLUT}连接中，接收者的膜通常位于\emph{树突棘}上——接收者输入分支上的一个小突起。连接的权重可以方便地分解为三个因子~\cite{delCastillo1954}：
\begin{empheq}[box=\keybox]{equation}
w \;=\; N_s\,p\,A_1 .
\label{eq:quantal}
\end{empheq}
这里$N_s$是两个神经元之间释放位点的数目，$p$是一个脉冲在一个位点释放一份递质的概率（就是第~\ref{sec:code}~章中的那个$p$），$A_1$是接收者对一份递质的响应，即有多少受体捕获它。因子$p$位于发送者一侧，$A_1$位于接收者一侧，而$N_s$需要两侧共同参与：新的释放位点长出来，旧的消失，这一过程持续终生。

\emph{由接收者决定。}在典型的\nt{GLUT}突触中，有一种谷氨酸受体只有在两个条件同时满足时才导通电流：谷氨酸从发送者到来，并且接收者的膜已经被兴奋~\cite{Nowak1984}。这是内置于一个分子中的赫布规则——第~\ref{sec:glutcomputer}~章的符合检测器，安装在突触本身的输入端。受体一旦触发，就让钙进入接收者，钙启动权重的变化。接收者是唯一能同时看到输入和输出的地方，所以决定在这里做出。

\emph{执行决定的可以是任何一侧。}研究最多的长时程增强发生在接收者一侧：新的受体嵌入膜中~\cite{Malinow2002}，树突棘长大~\cite{Matsuzaki2004}——$A_1$增大。但接收者也能改变$p$：它释放一种物质，穿过间隙逆向回到发送者——即附录~\ref{sec:othermediators}中的反向通道——于是发送者开始释放更少的递质~\cite{WilsonNicoll2001}。而权重的短时变化，即第~\ref{sec:code}~章中的耗竭和易化，是发送者根据自己近期的活动自行调整的$p$的变化。

对工程师来说，权重寄存器分布在两块芯片上。学习规则由接收者执行，而结果它既可以写在自己这边，也可以通过反向通道写在发送者那边。因此，本章规则中“局部”一词指的不是连接的某一侧，而是整个连接。

\section{无教师学习}

突触能做的最简单的事，是在发送者和接收者同时活跃时增强：
\begin{equation}
\frac{\dd w}{\dd t} \;=\; \eta\,x\,y .
\label{eq:hebb}
\end{equation}
这就是赫布规则~\cite{Hebb1949}。它是局部的，也不需要时钟。但它和第~\ref{sec:glutcomputer}~章中的\nt{GLUT}网络有同样的毛病：正反馈。强的权重使$y$增大，大的$y$又进一步增强权重，权重无限增长。解药是对权重施加负反馈：让权重同时按$y^2$成比例地衰减。对输入为$\mathbf x$、输出为线性$y=\mathbf w\cdot\mathbf x$的神经元，得到
\begin{empheq}[box=\keybox]{equation}
\frac{\dd\mathbf w}{\dd t} \;=\; \eta\,y\,\bigl(\mathbf x - y\,\mathbf w\bigr) .
\label{eq:oja}
\end{empheq}
这条规则仍然是局部的：突触知道自己的输入$x_i$、输出$y$和权重$w_i$。

\begin{keyblock}
\begin{proposition}[赫布规则找到什么]
\label{prop:oja}
规则~\eqref{eq:oja}使权重向量收敛到单位长度，并指向输入变化最强的方向——即输入相关矩阵$C=\langle\mathbf x\mathbf x^{\mathsf T}\rangle$的主特征向量~\cite{Oja1982}。神经元学会输出最有表现力的那一个量，把它的输入压缩到其中。
\end{proposition}
\end{keyblock}

证明。对输入求~\eqref{eq:oja}的平均：$\dd\mathbf w/\dd t=\eta\bigl(C\mathbf w-(\mathbf w^{\mathsf T}C\mathbf w)\,\mathbf w\bigr)$。不动点是$C$的单位长度特征向量。若$\mathbf w$接近特征值为$\lambda_k$的特征向量$\mathbf e_k$，沿另一个向量$\mathbf e_j$的小扰动按$e^{\eta(\lambda_j-\lambda_k)t}$增长。只有所有$\lambda_j<\lambda_k$的那个点是稳定的，也就是主向量。

赫布规则还有一种看脉冲时间的变体。如果发送者的脉冲比接收者的脉冲早$s$毫秒到达，权重增加$A_+e^{-s/\tau_+}$；如果晚$s$毫秒，则减少$A_-e^{-s/\tau_-}$，$\tau_\pm$约为二十毫秒。这样的规则已在\nt{GLUT}神经元的突触上测得~\cite{BiPoo1998}。它增强那些\emph{可能是}响应原因的连接，削弱迟到的连接（图~\ref{fig:stdp}）。让同一个兴奋序列多次通过网络，网络中就会自己长出从每一环指向下一环的连接——这正是第~\ref{sec:glutcomputer}~章的链，只不过无需工程师就搭成了。

\begin{figure}[t]
\centering
\begin{tikzpicture}[>=Stealth,x=0.07cm,y=1.3cm]
\draw[->,gray!70!black] (-66,0) -- (68,0) node[right,black] {\small $s$};
\draw[->,gray!70!black] (0,-1.2) -- (0,1.35) node[above,black] {\small 权重变化};
\draw[very thick,blue!60!black] plot[domain=0.5:60,samples=50] (\x,{exp(-\x/20)});
\draw[very thick,red!70!black] plot[domain=-60:-0.5,samples=50] (\x,{-0.6*exp(\x/20)});
\draw[blue!60!black,thick] (0.5,0) -- (0.5,1);
\draw[red!70!black,thick] (-0.5,0) -- (-0.5,-0.6);
\foreach \x in {-40,-20,20,40}{\draw[gray!60!black] (\x,-0.04) -- (\x,0.04);}
\node[below,font=\scriptsize] at (20,-0.04) {$20\ms$};
\node[above,font=\scriptsize] at (-20,0.04) {$-20\ms$};
\node[align=left,font=\small,blue!60!black,anchor=west] at (22,0.95) {发送者在先：\\ 连接增强};
\node[align=right,font=\small,red!70!black,anchor=east] at (-12,-0.95) {发送者在后：\\ 连接减弱};
\end{tikzpicture}
\caption{按时间的赫布规则。横轴为$s=t_{\text{接收}}-t_{\text{发送}}$，即接收者的脉冲比发送者的脉冲晚多少；$\tau_\pm=20\ms$，$A_-=0.6A_+$。宽几十毫秒的窗口与符合窗口~\eqref{eq:window}同一量级。}
\label{fig:stdp}
\end{figure}

但本节所有的规则教会的都是\emph{常见}的东西，而不是\emph{有用}的东西。一个只知道$x$和$y$的突触，不知道这个动作带来的是果汁还是疼痛。要从结果中学习，突触需要第三个信号。

\section{第三个因子}

第三个信号由\nt{DOPA}带来：多巴胺神经元广播一个数——奖励预测误差$\delta$~\eqref{eq:rpe}（第~\ref{sec:neurontypes}~章）。设权重的变化与三个因子的乘积成正比：发送者的活动、接收者的活动和$\delta$。

难点在于，奖励并不是在网络选择动作的时刻到来，而是在几百毫秒或几秒之后。等$\delta$到来时，乘积$xy$早已为零，突触需要记住自己曾经参与过。最简单的记忆就是我们熟悉的RC电路：
\begin{empheq}[box=\keybox]{equation}
\tau_e\,\frac{\dd e}{\dd t} \;=\; -e \;+\; x\,y ,
\qquad
\frac{\dd w}{\dd t} \;=\; \eta\,\delta(t)\,e(t) .
\label{eq:threefactor}
\end{empheq}
量$e$称为\emph{资格迹}~\cite{Fremaux2016}。共同活动在突触中留下一个标记，它在时间$\tau_e$内消退。如果在此期间多巴胺到来，权重就按$\delta$的符号变化；如果没有，标记便消失得无影无踪（图~\ref{fig:trace}）。在\nt{GLUT}突触上测得：在共同活动之后一两秒内到来的多巴胺会把它变成连接的增强，更晚到来的则不会~\cite{Yagishita2014}。在第三个信号不是多巴胺、而是接收者自身强烈兴奋的地方，也发现了长达数秒的可塑性窗口~\cite{Bittner2017}。

\begin{figure}[t]
\centering
\begin{tikzpicture}[>=Stealth,x=7cm,y=1cm]
\fill[orange!12] (0.7,-0.2) rectangle (0.8,5.2);
% rows: x*y, delta, traces, weights
\foreach \y/\lab in {4.4/{$x\,y$},3.1/{$\delta$},1.4/{资格迹$e$},0/{权重$w$}}{
  \draw[gray!70!black] (0,\y) -- (1.42,\y);
  \node[left,font=\small] at (-0.02,\y+0.3) {\lab};}
\draw[very thick,blue!60!black] (0,4.4) -- (0,4.9) -- (0.2,4.9) -- (0.2,4.4);
\node[right,font=\scriptsize,blue!60!black] at (0.21,4.8) {发送者和接收者同时活跃};
\draw[very thick,orange!85!black] (0,3.1) -- (0.7,3.1) -- (0.7,3.7) -- (0.8,3.7) -- (0.8,3.1) -- (1.4,3.1);
\node[right,font=\scriptsize,orange!85!black] at (0.81,3.6) {多巴胺到达};
% traces, each in units of its own maximum
\draw[very thick,blue!60!black] plot[domain=0:0.2,samples=20] (\x,{1.4+1.1*(1-exp(-\x))/(1-exp(-0.2))})
  plot[domain=0.2:1.4,samples=60] (\x,{1.4+1.1*exp(-(\x-0.2))});
\draw[thick,densely dashed,gray!50!black] plot[domain=0:0.2,samples=30] (\x,{1.4+1.1*(1-exp(-\x/0.05))/(1-exp(-4))})
  plot[domain=0.2:1.4,samples=80] (\x,{1.4+1.1*exp(-(\x-0.2)/0.05)});
\node[right,font=\scriptsize,blue!60!black] at (1.0,2.15) {$\tau_e=1\,\text{s}$};
\node[right,font=\scriptsize,gray!40!black] at (0.3,1.65) {$\tau_e=50\ms$};
% weights
\draw[very thick,blue!60!black] (0,0.25) -- (0.7,0.25) -- (0.8,0.8) -- (1.4,0.8) node[right,font=\scriptsize] {$\tau_e=1\,\text{s}$：权重增大};
\draw[thick,densely dashed,gray!50!black] (0,0.2) -- (1.4,0.2) node[right,font=\scriptsize,gray!40!black] {$\tau_e=50\ms$：不变};
% time axis marks
\foreach \x/\l in {0/0,0.2/0.2,0.7/0.7,1.4/1.4}{\draw[gray!60!black] (\x,-0.05) -- (\x,0.05) node[below=3pt,font=\scriptsize] {\l\,s};}
\draw[<->,thick,gray!50!black] (0.2,5.35) -- node[above,font=\scriptsize,black] {奖励延迟$0.5\,\text{s}$} (0.7,5.35);
\end{tikzpicture}
\caption{按规则~\eqref{eq:threefactor}的资格迹。共同活动持续$0.2\,\text{s}$，多巴胺在其结束半秒后到达。资格迹以各自最大值的比例表示。慢的资格迹（$\tau_e=1\,\text{s}$）此刻仍然存在，快的（$\tau_e=50\ms$）早已消退。}
\label{fig:trace}
\end{figure}

\begin{keyblock}
\begin{proposition}[用广播信号学习]
\label{prop:threefactor}
规则~\eqref{eq:threefactor}只改变在最近$\tau_e$时间内参与了网络工作的那些突触，改变方向由公共信号$\delta$给定。广播信号加上局部资格迹，就得到有地址的修正。只要动作与结果之间的延迟不超过$\tau_e$，这种延迟就是允许的。
\end{proposition}
\end{keyblock}

\section{为什么行得通：噪声即搜索}
\label{sec:dofagradient}

为什么规则~\eqref{eq:threefactor}会使结果变好？答案来自第~\ref{sec:code}~章的噪声。设神经元的输出$y=f(u)+\xi$，其中$u=\sum_jw_jx_j$，$\xi$是均值为零、方差为$\sigma^2$的随机抖动。奖励$R$取决于$y$。我们不取$x_jy$本身作为资格迹，而取它的随机部分$x_j\xi$，并取误差$\delta=R-\bar R$作为第三个因子，其中$\bar R$是期望奖励。噪声很小时$R\approx R\bigl(f(u)\bigr)+R'\,\xi$，因此
\begin{empheq}[box=\keybox]{equation}
\bigl\langle \delta\,\xi\,x_j \bigr\rangle \;=\; \sigma^2\,R'\,x_j \;=\; \frac{\sigma^2}{f'(u)}\,\frac{\partial\langle R\rangle}{\partial w_j} .
\label{eq:perturbation}
\end{empheq}
平均步长沿着期望奖励的梯度~\cite{Williams1992}。谁也没有计算导数：网络随机地偏离了一下，多巴胺报告结果是否变好，参与的突触就向有利的方向移动（图~\ref{fig:perturbation}）。第~\ref{sec:code}~章中妨碍传送信息的噪声，在这里变成了搜索。

\begin{figure}[t]
\centering
\begin{tikzpicture}[>=Stealth,x=2cm,y=3cm]
\draw[->,gray!70!black] (0,0) -- (4.2,0) node[below left,font=\small,black] {输出$y$};
\draw[->,gray!70!black] (0,0) -- (0,1.2) node[left,font=\small,black] {奖励$R$};
% reward hill
\draw[very thick,blue!60!black] plot[domain=0:4,samples=60] (\x,{max(0,1-((\x-3)/2.2)^2)});
\node[font=\scriptsize,blue!60!black,anchor=south] at (3,1.02) {最佳输出};
% noise around f(u)
\draw[thick,gray!60!black] plot[domain=0.6:2.4,samples=50] (\x,{0.28*exp(-((\x-1.5)/0.3)^2/2)});
\draw[densely dashed,gray!60!black] (1.5,0) -- (1.5,0.535);
\node[below,font=\small] at (1.5,0) {$f(u)$};
\node[font=\scriptsize,gray!50!black] at (1.5,-0.2) {抖动$\xi$};
% expected reward
\draw[densely dotted,gray!60!black] (0.5,0.535) node[left,font=\scriptsize,black] {$\bar R$} -- (2.1,0.535);
% two random deviations
\fill[green!45!black] (2.1,0.833) circle (0.035);
\draw[very thick,green!45!black] (2.1,0.535) -- (2.1,0.833);
\node[font=\scriptsize,green!45!black,anchor=west,align=left] at (2.18,0.6) {$\xi>0$：更好，\\ $\delta>0$};
\fill[red!70!black] (0.9,0.089) circle (0.035);
\draw[very thick,red!70!black] (0.9,0.535) -- (0.9,0.089);
\node[font=\scriptsize,red!70!black,anchor=east,align=right] at (0.83,0.27) {$\xi<0$：更差，\\ $\delta<0$};
% resulting step
\draw[->,very thick,orange!85!black] (1.55,0.4) -- (1.95,0.4) node[right,font=\scriptsize,black] {步长};
\end{tikzpicture}
\caption{噪声即搜索~\eqref{eq:perturbation}。神经元的输出在$f(u)$附近抖动。向右的随机偏离（绿色）得到的比预期多，$\delta>0$；向左的（红色）得到的少，$\delta<0$。两种情况下乘积$\delta\,\xi$都是正的，权重向右移动，即向奖励增大的方向。平均而言，步长与斜率$R'$成正比：在山顶，向两侧的偏离同样糟糕，步长相互抵消。}
\label{fig:perturbation}
\end{figure}

\begin{keyblock}
\begin{proposition}[无需反向导线的梯度下降]
\label{prop:gradient}
如果网络中存在活动的随机偏离，而广播信号等于奖励预测误差、并且在每种情境下平均为零，那么规则~\eqref{eq:threefactor}平均而言沿期望奖励的梯度改变权重。这不需要反向导线，也不需要单独的学习阶段。
\end{proposition}
\end{keyblock}

现在就明白了，为什么多巴胺报告的不是奖励本身，而是奖励预测的\emph{误差}。如果第三个因子是奖励本身$R\ge0$，式~\eqref{eq:perturbation}中的平均不会变，但会出现两个麻烦。第一，有用部分之外会多出一项$\bar R\,\xi x_j$：平均为零，但噪声很强。第二，更糟的是，真实的突触无法把资格迹$x_jy$中的随机部分和非随机部分分开。在给定输入下，$x_jf(u)$在每次尝试中都是同一个数；如果$\delta$在这一情境下平均为零，这部分资格迹不产生任何改变，而若$\delta\ge0$，它就会一次又一次地增强网络所做的一切——对的和错的。因此$\bar R$必须是\emph{在给定情境下}的期望，而不是整个一生的平均。计算它是评论器的事，下面会讲到。

我们在模型上检验了这一点。两组\nt{GLUT}神经元通过相互抑制在两个动作中选一个，如图~\ref{fig:wta}；从“开始”信号到两组的权重起初相等，由噪声决定谁获胜。动作$A$以概率$0.8$带来果汁，动作$B$以概率$0.2$带来果汁，果汁在选择后半秒到达。当$\tau_e=1\,\text{s}$、$\delta=R-\bar R$时，五百次尝试中选择$A$的比例从前一百次的$0.65$--$0.8$升到最后一百次的$0.96$--$1.0$，而权重保持在$0.85$--$1.35$范围内。当$\tau_e=50\ms$、小于奖励延迟时，网络什么也没学会：选择$A$的比例始终在一半左右。当$\delta=R$、不减去期望时，网络也开始选择$A$，但获胜者的权重在同样五百次尝试中增长了一千倍，并且还在继续增长；而在同样$\delta=R$、学习速率大十倍时，三次运行中有一次网络永久固定了自己最初的随机选择——更差的那个。

\section{一根导线的代价}

信号$\delta$对所有人只有一个，而它所评估的噪声是影响结果的全部$N$个神经元抖动之和。每个突触在$\delta$中看到的是自己的有用部分和$N-1$个邻居的外来噪声。为了把自己的部分分离出来，必须在许多次尝试上求平均，学习时间与$N$成正比地增长~\cite{Werfel2005}。反向传播无需付出这种代价。

\begin{keyblock}
\begin{proposition}[广播的代价]
\label{prop:broadcastcost}
依靠一个公共信号学习，所需时间与其噪声影响结果的神经元数目成正比地增长。这种学习适合从少数几个选项中做选择的小电路，但不适合调节数十亿个权重。
\end{proposition}
\end{keyblock}

脑看来正是这样分工的：\nt{DOPA}神经元的输出主要通向选择动作的区域（第~\ref{sec:neurontypes}~章），而皮层庞大的\nt{GLUT}网络——绝大多数权重都在这里——主要依靠局部规则学习，不需要外部教师。过去人们把这些规则归结为赫布规则~\eqref{eq:oja}。现在越来越多的证据表明，皮层有自己的目标——\emph{预测}自己的下一个输入，于是误差就在原地计算，即预测值与实际到来值之差~\cite{Keller2018}：教师内置于网络本身。在第三个信号是接收者自身强烈兴奋的地方，长达数秒的可塑性窗口~\cite{Bittner2017}表明，突触处确实存在局部误差信号。这在多大程度上是皮层的普遍规则，还是假说。

\section{误差从何而来：连续时间中的评论器}

剩下的问题是，\nt{DOPA}神经元自己如何计算$\delta$。在强化学习程序中，预测误差按步$t,t+1,\dots$计算。体内没有步，所以我们在连续时间中写出它~\cite{Doya2000}。设$r(t)$为奖励流，$V(t)$为对全部未来奖励的期望，其中越远的奖励价值越低：
\begin{equation}
V(t) \;=\; \int_t^{\infty} e^{-(s-t)/\tau_\gamma}\,r(s)\,\dd s .
\end{equation}
求导得$\dd V/\dd t=V/\tau_\gamma-r$。这个等式的残差就是预测误差：
\begin{empheq}[box=\keybox]{equation}
\delta(t) \;=\; r(t) \;+\; \frac{\dd V}{\dd t} \;-\; \frac{V}{\tau_\gamma} .
\label{eq:ctd}
\end{empheq}
常数$\tau_\gamma$是规划视野，即系数$\gamma$的连续时间对应物；按照假说，它由\nt{SERO}设定（第~\ref{sec:horizon}~节），这样~\eqref{eq:patience}就是值得等多久的规则。我们检查三种情况（图~\ref{fig:ctdcases}）。预示果汁的灯亮了：$V$跃升，$\dd V/\dd t$很大，出现一个爆发。承诺的果汁到了：$r>0$，但期望在同一时刻耗尽，$\dd V/\dd t\approx-r$，没有爆发。果汁没来：期望在没有奖励的情况下消失，$\dd V/\dd t<0$，出现低谷。

\begin{figure}[t]
\centering
\begin{tikzpicture}[>=Stealth,x=1.15cm,y=1cm]
\foreach \col/\ttl/\lampon/\juice in {0/{意外的果汁}/0/1, 1/{承诺的果汁}/1/1, 2/{果汁没来}/1/0}{
 \begin{scope}[xshift=\col*4.6cm]
  \node[font=\small] at (1.5,3.55) {\ttl};
  \foreach \yy in {2.4,1.2,0}{\draw[gray!60!black] (0,\yy) -- (3.1,\yy);}
  % reward r
  \ifnum\juice=1
    \fill[green!45!black] (1.95,2.4) rectangle (2.05,2.85);
    \pic[scale=0.55] at (2.0,3.1) {juice};
  \else
    \pic[scale=0.55] at (2.0,3.1) {nojuice};
  \fi
  % expectation V and error delta
  \ifnum\lampon=1
    \pic[scale=0.55] at (1.0,3.1) {lamp};
    \draw[very thick,blue!60!black] (0,1.2) -- (1,1.2) -- (1,1.45)
      plot[domain=1:2,samples=20] (\x,{1.2+0.25*exp((\x-1)/1.5)}) -- (2,{1.2+0.25*exp(1/1.5)}) -- (2,1.2) -- (3.1,1.2);
    \draw[very thick,red!70!black] (0,0) -- (0.95,0) -- (0.95,0.5) -- (1.05,0.5) -- (1.05,0) -- (3.1,0);
    \ifnum\juice=0
      \draw[very thick,red!70!black] (1.95,0) -- (1.95,-0.45) -- (2.05,-0.45) -- (2.05,0);
      \node[font=\scriptsize,red!70!black,anchor=west] at (2.1,-0.35) {低谷};
    \else
      \node[font=\scriptsize,gray!50!black,anchor=north,align=center] at (2.0,-0.08) {$r$与$V$的下降\\ 相互抵消};
    \fi
  \else
    \draw[very thick,blue!60!black] (0,1.2) -- (3.1,1.2);
    \draw[very thick,red!70!black] (0,0) -- (1.95,0) -- (1.95,0.5) -- (2.05,0.5) -- (2.05,0) -- (3.1,0);
  \fi
  \ifnum\col=0
    \node[font=\small,anchor=east] at (-0.1,2.55) {$r$};
    \node[font=\small,anchor=east] at (-0.1,1.35) {$V$};
    \node[font=\small,anchor=east] at (-0.1,0.1) {$\delta$};
  \fi
 \end{scope}}
\end{tikzpicture}
\caption{误差~\eqref{eq:ctd}的三种情况；自上而下依次为奖励$r$、期望$V$和误差$\delta$。左：没有期望，果汁完全出乎意料。中：灯使$V$跃升，这是$\dd V/\dd t$的跃变，$\delta$的爆发落在灯亮时。之后$V$恰好按视野的要求增长，$\delta=0$；果汁到来时，奖励$r$与$V$的下降相互抵消。右：$V$在没有奖励时下降——低谷。这就是图~\ref{fig:rpe}中的爆发、转移和低谷，只是现在能看出它们从何而来。}
\label{fig:ctdcases}
\end{figure}

用我们已有的电路实现~\eqref{eq:ctd}需要什么？三样东西。

\emph{估计值$V$本身。}它由一组\nt{GLUT}神经元——评论器——输出，其权重用同样的规则~\eqref{eq:threefactor}、同样的$\delta$学习：若$\delta>0$，说明期望偏低，此前活跃的连接就增强。

\emph{导数$\dd V/\dd t$。}任何对变化而非对电平作出响应的电路都能计算它：直接兴奋加上经\nt{GABA}中介的延迟抑制，如第~\ref{sec:gaba}~章的方向检测器；或者第~\ref{sec:code}~章的耗竭型突触~\eqref{eq:depression}。

\emph{一根导线中的符号。}误差可正可负，而脉冲频率只能为正。解决办法与第~\ref{sec:serocomputer}~章中\nt{SERO}神经元的相同：\nt{DOPA}神经元是起搏器。它平稳的每秒几个脉冲表示$\delta=0$，簇放电表示$\delta>0$，停顿表示$\delta<0$。负误差的余地更小：频率不能降到零以下。真实的\nt{DOPA}神经元的低谷确实比爆发浅~\cite{Bayer2005}。

多巴胺的行为像$\delta$，这是测量结果。脑究竟如何计算$\dd V/\dd t$，还是假说。

\section{行动者与评论器}

把所有部分组装起来（图~\ref{fig:actorcritic}）。情境神经元兴奋各动作组，动作组通过\nt{GABA}选出获胜者（命题~\ref{prop:wta}）——这是\emph{行动者}；情境神经元同时兴奋输出$V$的评论器~\cite{SuttonBarto2018}。\nt{DOPA}神经元接收奖励$r$和$V$，计算~\eqref{eq:ctd}，并把$\delta$广播给行动者和评论器的所有突触。

\begin{figure}[t]
\centering
\begin{tikzpicture}[>=Stealth,
  nrn/.style={circle,draw,very thick,fill=blue!6,minimum size=0.9cm,inner sep=0pt,font=\small},
  gab/.style={circle,draw,very thick,fill=red!10,minimum size=0.6cm,inner sep=0pt},
  dofa/.style={circle,draw,very thick,double,double distance=1.2pt,fill=orange!15,minimum size=1.35cm,inner sep=0pt,font=\scriptsize},
  inh/.style={-{Bar[width=7pt]},thick,red!70!black},
  bc/.style={->,thick,densely dashed,orange!85!black}]
\node[nrn] (S1) at (0,2.4) {$x_1$};
\node[nrn] (S2) at (0,1.2) {$x_2$};
\node[nrn] (S3) at (0,0) {$x_3$};
\node[left=0.15cm,align=right,font=\small] at (-0.5,1.2) {情境};
\node[nrn] (A) at (4,2.6) {$A$};
\node[nrn] (B) at (4,1.0) {$B$};
\node[gab] (G) at (5.2,1.8) {};
\draw[->,thick] (A) -- (G); \draw[->,thick] (B) -- (G);
\draw[inh] (G) to[bend left=35] (A);
\draw[inh] (G) to[bend right=35] (B);
\node[above,font=\small] at (4.6,3.05) {行动者：选择动作};
\node[nrn] (V) at (4,-1.1) {$V$};
\node[below,font=\small] at (4,-1.6) {评论器};
\foreach \s in {S1,S2,S3}{\foreach \t in {A,B,V}{\draw[->,gray!60!black] (\s) -- (\t);}}
\node[dofa] (D) at (8.0,-1.1) {\nt{DOPA}};
\draw[->,thick] (V) -- node[above,font=\scriptsize] {$V,\ \dd V/\dd t$} (D);
\draw[->,thick,green!45!black] (10.0,-1.1) node[right,black,font=\small] {奖励$r$} -- (D);
\pic[scale=1.1] at (10.6,-0.5) {juice};
\draw[bc] (D.north) -- (8.0,4.0) -- node[above,font=\small,black] {$\delta$——送往行动者和评论器的所有突触} (2.2,4.0) -- (2.2,3.0);
\draw[bc] (D.south) -- (8.0,-2.4) -- (2.2,-2.4) -- (2.2,-0.6);
\end{tikzpicture}
\caption{学习选择动作的网络。灰色箭头是权重可变的\nt{GLUT}突触；红色小圆是\nt{GABA}神经元；双圆圈是按~\eqref{eq:ctd}计算$\delta$的\nt{DOPA}起搏器；虚线箭头是$\delta$的广播。}
\label{fig:actorcritic}
\end{figure}

递质作用的符号由接收者选择（命题~\ref{prop:receptor}），在动作选择区域，一部分神经元带有一个家族的多巴胺受体，另一部分带有另一个家族的。在脑片实验中，多巴胺与共同活动一起使前者的突触增强，使后者的突触减弱~\cite{Shen2008}；在模型中，这意味着对后者而言，\eqref{eq:threefactor}中$\delta$的符号是反的。前者学习\emph{值得做}什么，后者学习\emph{不值得做}什么。

\section{小结}

\begin{table}[t]
\centering
\footnotesize
\setlength{\tabcolsep}{4pt}
\renewcommand{\arraystretch}{1.15}
\begin{tabular}{>{\raggedright}p{2.6cm}>{\raggedright}p{2.7cm}>{\raggedright}p{3.1cm}>{\raggedright\arraybackslash}p{5.0cm}}
\toprule
规则 & 突触知道什么 & 网络学到什么 & 已知情况 \\
\midrule
按频率的赫布规则~\eqref{eq:oja} & $x$、$y$、自身权重 & 压缩输入：主方向 & 共同活动时的增强已测得；限制权重的机制仍在研究中 \\
按时间的赫布规则 & $\sim20\ms$内$x$和$y$脉冲的先后次序 & 因果联系、序列 & 已在\nt{GLUT}突触上测得 \\
三因子规则~\eqref{eq:threefactor} & $x$、$y$、资格迹$e$、公共信号$\delta$ & 带来奖励的动作 & 活动后数秒内多巴胺的作用已测得；作为学习选择的主要机制，是有充分依据的假说 \\
评论器~\eqref{eq:ctd} & 同上 & 奖励期望$V$ & 多巴胺与$\delta$的相似已测得；计算$\dd V/\dd t$的电路是假说 \\
误差反向传播 & 每个权重各自的修正 & 一切，但需要反向导线和时钟节拍 & 在脑中未发现这种形式；人们在寻找它的近似 \\
\bottomrule
\end{tabular}
\caption{由\nt{GLUT}、\nt{GABA}和\nt{DOPA}神经元构成的网络中的学习规则。}
\label{tab:learning}
\end{table}

学习规则汇总于表~\ref{tab:learning}。\nt{GLUT}携带数据，\nt{GABA}控制数据流，而\nt{DOPA}决定网络中保留哪些改变。

\section{习题}

\begin{exercise}[\lvlA]
\label{ex:06:1}
\emph{资格迹能存活多久。}共同活动$xy=1$持续$T_a=0.2\,\text{s}$，从$e=0$开始，多巴胺在其结束后$d=0.5\,\text{s}$到达。(a)~此刻$\tau_e=1\,\text{s}$时的资格迹~\eqref{eq:threefactor}是$\tau_e=50\ms$时的多少倍？(b)~$\tau_e$为多少时它最大？
\end{exercise}

\begin{exercise}[\lvlA]
\label{ex:06:2}
\emph{赫布规则的漂移。}发送者和接收者输出独立的泊松脉冲流，各为每秒$10$个脉冲，每一对脉冲都按图~\ref{fig:stdp}的窗口改变权重，$\tau_\pm=20\ms$，$A_-=0.6A_+$。完全随机的活动持续一小时，权重增长多少个$A_+$？$\tau_-$为多少时漂移消失？
\end{exercise}

\begin{exercise}[\lvlB]
\label{ex:06:3}
\emph{是相关，而不是协方差。}命题~\ref{prop:oja}的规则找到$C=\langle\mathbf x\mathbf x^{\mathsf T}\rangle$的主特征向量。频率为正：$\mathbf x=\mathbf m+\mathbf n$，$\mathbf m=(\mu,0)$，噪声协方差为$\operatorname{diag}(s_1^2,s_2^2)$。在什么条件下神经元会学到$\mathbf e_1$？当$\mu=1$、$s_1=0.1$、$s_2=0.5$时它学到什么？
\end{exercise}

\begin{exercise}[\lvlA]
\label{ex:06:4}
\emph{爆发中的视野。}灯在不可预测的时刻亮起，大小为$R$的果汁在$L=1\,\text{s}$后到达。证明对于按~\eqref{eq:ctd}学到的期望，灯与果汁之间$\delta\equiv0$，并求$\tau_\gamma=2\,\text{s}$和$\tau_\gamma=4\,\text{s}$时灯亮爆发的面积。
\end{exercise}

\begin{exercise}[\lvlB]
\label{ex:06:5}
\emph{广播的代价从何而来。}$N$个神经元独立抖动，$\xi_i\sim\mathcal N(0,\sigma^2)$，误差$\delta=a\sum_i\xi_i$，突触$j$用$\delta\,\xi_jx_j$估计梯度。求单次尝试的信噪比。单个神经元用$100$次、每次$5\,\text{s}$的尝试学会：$N=10^4$和$N=10^{10}$时学习需要多久？
\end{exercise}

\begin{exercise}[\lvlB]
\label{ex:06:6}
\emph{设计：计算$\delta$的电路。}\nt{DOPA}神经元接收$r$、权重为$k_1$的$V$，以及以$\tau_d$平滑、权重为$-k_2$的$V$的副本。(a)~选取$k_1$、$k_2$，使对慢变的$V$输出为~\eqref{eq:ctd}。(b)~当$V=V_0e^{t/\tau_\gamma}$时，真实的$\delta=0$。若$\tau_\gamma=2\,\text{s}$，$\tau_d$为多少时电路的虚假误差不超过$V/\tau_\gamma$的$5\,\%$？
\end{exercise}

\begin{exercise}[\lvlB]
\label{ex:06:7}
\emph{仿真：奖励延迟的极限。}在\texttt{models/\allowbreak dofa\_learning.py}副本的\texttt{run()}中加入奖励延迟$d$。求$d=0.5\,\text{s}$、$\tau_e=0.05$、$0.2$、$1$、$4\,\text{s}$时，以及$\tau_e=1\,\text{s}$、$d=3\,\text{s}$时，$500$次尝试中最后一百次选择$A$的比例。资格迹存活到奖励时刻的比例（习题~\ref{ex:06:1}）为多少时，学习消失？
\end{exercise}

\begin{exercise}[\lvlC]
\label{ex:06:8}
\emph{能否绕过广播的代价。}$N$个神经元：$y_i=w_i+\xi_i$，$\xi_i\sim\mathcal N(0,\sigma^2)$，奖励$R=-\sum_i(y_i-w_i^*)^2$，规则$\Delta w_i=\eta\,(R-\langle R\rangle)\,\xi_i$。在最优$\eta$下，误差$\sum_i(w_i-w_i^*)^2$缩小为$1/e$需要多少次尝试？如果$N$个神经元中只有$m$个随机选出的神经元抖动呢？稀疏搜索有帮助吗？
\end{exercise}

\chapter{基于多巴胺的其他现代学习理论}
\chaptermark{\nt{DOPA}的其他理论}
\label{sec:dofaalt}

\section{导线上实际传送的是什么}

在第~\ref{sec:dofa}~章中，\nt{DOPA}系统对我们来说是一根导线，上面传送一个数——奖励预测误差$\delta$~\eqref{eq:ctd}。这幅图景解释了很多事情，但多巴胺测得越精确，发现的与之不符的东西就越多。有些神经元对厌恶事件的响应和对奖励一样；动物奔向目标时多巴胺浓度不断上升，尽管并没有发生任何意外；不同的\nt{DOPA}神经元表现各不相同。

对工程师来说，所有这些发现都归结为一个问题：\emph{广播总线上信号的格式是什么？}是一个数还是一个向量？导线上是一个信号，还是按频率分开的几个？它像$\delta$那样谈论未来，还是谈论过去？下面给出六种现代答案；对每一种，我们都把已测量的与推测的分开。

\section{不是一个数，而是一个分布}

误差~\eqref{eq:ctd}教给评论器的是奖励的\emph{平均}期望。但必定到手的半滴果汁和概率为二分之一的一滴果汁平均值相同。要区分它们，需要奖励的整个分布。

取最简单的任务：预示信号之后到来一个大小随机的奖励$r$。设有许多\nt{DOPA}神经元，第$i$个负责自己的估计$V_i$，因此在奖励时刻它的误差为$\delta_i=r-V_i$。再设每个神经元以权重$\alpha_i^+$计入正误差，以权重$\alpha_i^-$计入负误差。每次奖励之后，估计移动
\begin{equation}
\Delta V_i \;=\; \eta\,\bigl(\alpha_i^+\,[\delta_i]_+ \;-\; \alpha_i^-\,[-\delta_i]_+\bigr).
\label{eq:asymmetric}
\end{equation}
当正修正与负修正在平均意义上相互平衡时，估计不再变化：
\begin{empheq}[box=\keybox]{equation}
\alpha_i^+\,\bigl\langle[r-V_i]_+\bigr\rangle \;=\; \alpha_i^-\,\bigl\langle[V_i-r]_+\bigr\rangle ,
\qquad
\kappa_i \;=\; \frac{\alpha_i^+}{\alpha_i^++\alpha_i^-} .
\label{eq:expectile}
\end{empheq}
$\kappa_i=1/2$时这就是普通的平均值。$\kappa_i>1/2$时神经元是“乐观者”：它的估计偏向分布的上端；$\kappa_i<1/2$时是“悲观者”。

例：果汁以概率$1/2$到来，每滴大小为$1$。于是~\eqref{eq:expectile}给出$\kappa_i\cdot\tfrac12(1-V_i)=(1-\kappa_i)\cdot\tfrac12V_i$，即$V_i=\kappa_i$（图~\ref{fig:expectiles}）。一组$\kappa_i$各不相同的神经元记录奖励分布，就像统计学中一组分位数记录一个分布一样。

\begin{figure}[t]
\centering
\begin{tikzpicture}[>=Stealth,x=7cm,y=3cm]
\draw[->,gray!70!black] (-0.08,0) -- (1.15,0) node[right,black] {\small 奖励};
\draw[->,gray!70!black] (-0.08,0) -- (-0.08,0.75) node[left,black] {\small 概率};
\fill[blue!25] (-0.03,0) rectangle (0.03,0.5);
\fill[blue!25] (0.97,0) rectangle (1.03,0.5);
\node[below,font=\small] at (0,0) {$0$};
\node[below,font=\small] at (1,0) {$1$};
\node[left,font=\scriptsize] at (-0.08,0.5) {$1/2$};
\foreach \k/\c/\lab/\h in {0.2/blue!60!black/{悲观者，$\kappa=0.2$}/0.62, 0.5/gray!50!black/{平均，$\kappa=0.5$}/0.42, 0.8/red!70!black/{乐观者，$\kappa=0.8$}/0.22}{
  \draw[very thick,\c] (\k,0) -- (\k,\h);
  \node[above,font=\scriptsize,\c] at (\k,\h) {\lab};}
\end{tikzpicture}
\caption{奖励分布由一组\nt{DOPA}神经元记录。奖励以概率$1/2$为$0$或$1$（蓝色柱）；比例为$\kappa$的神经元收敛到估计$V=\kappa$~\eqref{eq:expectile}。}
\label{fig:expectiles}
\end{figure}

已测量的是：不同\nt{DOPA}神经元的“翻转点”——响应从爆发变为低谷的奖励大小——各不相同，而对正负误差响应不对称性更大的神经元在更大的奖励处翻转，正如~\eqref{eq:expectile}所要求的~\cite{Dabney2020}。从一组神经元的响应中可以重建分布的形状。这是离散性的主要功能而不是副作用，则是假说。

\begin{keyblock}
\begin{proposition}[由不对称得到分布]
\label{prop:expectile}
如果\nt{DOPA}神经元计入正误差的比例$\kappa_i$各不相同，它们的估计就收敛到奖励分布的不同点~\eqref{eq:expectile}。一组神经元传送的不是平均值，而是分布，从而也传送了风险。
\end{proposition}
\end{keyblock}

\section{不只是误差，还有重要性}

按误差公式~\eqref{eq:ctd}，厌恶事件——电击、苦味——应当引起低谷：结果比预期差。一部分\nt{DOPA}神经元确实如此，另一部分却对厌恶事件以\emph{爆发}响应，就像对奖励一样~\cite{Matsumoto2009}。这些神经元报告的不是误差的符号，而是发生了某件\emph{重要}的事：意外的、强烈的、需要注意的~\cite{BrombergMartin2010}。

工程师可以把它看作两个信号。第一个是有符号的误差$\delta$：权重往哪个方向改。第二个是它的模$|\delta|$或事件的新异性：改变权重的\emph{幅度}，也就是学习速率；意外事件之后应当学得更快。就含义而言，它接近第~\ref{sec:neurontypes}~章中去甲肾上腺素的增益。

还有第三种可能：为单独的一个维度设一根单独的导线。投射到某个动作选择区域的\nt{DOPA}神经元不对奖励响应，而对刺激的新异性和强度响应，动物要学会躲避危险就需要它们~\cite{Menegas2018}：导线上传送的不是“更好还是更差”，而是“危险还是不危险”。

已测量的是：不同组的\nt{DOPA}神经元对厌恶事件和新异事件的响应不同，不同的输出教不同的东西。脑究竟如何使用重要性信号——作为学习速率、作为注意信号，还是作为危险维度上的单独误差——仍在讨论中。

\section{不只是教，还要催}

多巴胺还有即时作用：多巴胺越多，动物越愿意、越快地行动。这怎样与学习相协调？

工程上的回答是\emph{按频率分离}。持续几百毫秒的快速爆发和低谷携带$\delta$，负责教学。在几分钟内变化的慢电平携带另一个量——单位时间的平均奖励$\bar R$~\cite{Niv2007}。设在时间$\tau_a$内完成一个动作需要付出努力$C/\tau_a$：越快越贵。但每拖延一秒要损失$\bar R$——这一秒本可以得到的奖励。总代价$C/\tau_a+\bar R\,\tau_a$在下式时最小：
\begin{empheq}[box=\keybox]{equation}
\tau_a^{*} \;=\; \sqrt{\frac{C}{\bar R}} .
\label{eq:vigor}
\end{empheq}
环境越富足，时间越贵，快速行动就越划算：若$\bar R$加倍，动作时间就除以$\sqrt2\approx1.4$。注意，$\bar R$正是第~\ref{sec:dofa}~章中被减去的那个期望奖励。

多巴胺在目标处的释放可以在脉冲频率不变的情况下改变：输出端本身的局部信号会调节它~\cite{Mohebi2019}。但脉冲本身也有慢成分：在下一节的实验中，接近奖励时的平缓上升在\nt{DOPA}神经元的频率中也能看到~\cite{Kim2020}。可见，慢电平由两个来源组成——脉冲频率和释放的局部调节——哪一个为主，看来取决于输出和任务。

\begin{keyblock}
\begin{proposition}[一根导线上的两个信号]
\label{prop:multiplex}
\nt{DOPA}系统至少传送两个在时间上分开的量：负责教学的快速误差$\delta$，以及设定时间代价~\eqref{eq:vigor}、从而设定动作速度的慢电平。已测得快成分和慢成分表现不同；慢成分恰好等于$\bar R$，是假说。
\end{proposition}
\end{keyblock}

\section{接近时的上升}

如果动物沿走廊奔向远处的奖励，动作选择区域中的多巴胺浓度会在目标接近的几秒内平缓上升~\cite{Hamid2016}。按第~\ref{sec:dofa}~章，这不应该发生：一切都按计划进行，$\delta$应为零。

第一种解释：这种上升不是$\delta$，而是期望$V$本身，即上一节中“目标很近，值得努力”的信号~\cite{Hamid2016}。第二种解释保留了$\delta$~\cite{Kim2020}。如果期望是准确的，那么在视野为$\tau_\gamma$时，通往奖励$R$的途中它按$V=Re^{-(T-t)/\tau_\gamma}$增长，按误差公式~\eqref{eq:ctd}有$\dd V/\dd t-V/\tau_\gamma=0$。但设在远处期望偏低，随着接近估计不断修正。那么期望增长得更陡，比如$V=Re^{-(T-t)/\tau_v}$，$\tau_v<\tau_\gamma$，于是
\begin{empheq}[box=\keybox]{equation}
\delta(t) \;=\; \frac{\dd V}{\dd t}-\frac{V}{\tau_\gamma} \;=\; V\Bigl(\frac{1}{\tau_v}-\frac{1}{\tau_\gamma}\Bigr) \;>\; 0 .
\label{eq:ramp}
\end{empheq}
误差为正，并随$V$一起增长——正是那种平缓上升（图~\ref{fig:ramp}）。当$\tau_\gamma=4\,\text{s}$、$\tau_v=1\,\text{s}$时，得到每秒$\delta=0.75\,V$。这一版本在虚拟走廊中得到检验：把动物瞬间移到离目标更近处，多巴胺以爆发响应，这正是导数应有的表现，而不是平缓的电平~\cite{Kim2020}。

\begin{figure}[t]
\centering
\begin{tikzpicture}[>=Stealth,x=1.6cm,y=2.6cm]
\draw[->,gray!70!black] (-4.2,0) -- (0.5,0) node[below left,black] {\small $t$};
\draw[->,gray!70!black] (-4.2,0) -- (-4.2,1.2);
\foreach \x/\l in {-4/{$T-4$},-2/{$T-2$},0/{$T$}}{\draw[gray!60!black] (\x,-0.03) -- (\x,0.03); \node[below,font=\scriptsize] at (\x,0) {\l\,s};}
\draw[thick,densely dashed,gray!60!black] plot[domain=-4:0,samples=40] (\x,{exp(\x/4)});
\draw[very thick,blue!60!black] plot[domain=-4:0,samples=60] (\x,{exp(\x)});
\draw[very thick,red!70!black] plot[domain=-4:0,samples=60] (\x,{0.75*exp(\x)});
\node[anchor=south east,font=\small,gray!50!black] at (-1.3,0.77) {$\tau_v=\tau_\gamma$时的$V$：$\delta=0$};
\node[right,font=\small,blue!60!black] at (0.05,1.0) {$\tau_v=1\,\text{s}$时的$V$};
\node[right,font=\small,red!70!black] at (0.05,0.72) {按~\eqref{eq:ramp}的$\delta$};
\end{tikzpicture}
\caption{接近奖励时多巴胺的上升被解释为预测误差，$\tau_\gamma=4\,\text{s}$。虚线：恰好按视野要求增长的期望（$\delta=0$）；蓝线：增长更陡的期望；红线：误差~\eqref{eq:ramp}。}
\label{fig:ramp}
\end{figure}

已测量的是：平缓上升确实存在——在多巴胺浓度中，也在\nt{DOPA}神经元的脉冲频率中~\cite{Kim2020}——而环境的瞬间跳变会引起多巴胺的跳变。两种解释哪一种正确，或者两者在不同输出上都正确，仍在讨论中。

\section{向后看}

上面所有的理论都是向\emph{前}看的。一种现代理论把方向反了过来~\cite{Jeong2022}。设脑学习的不是根据预示信号预测奖励，而是在得到奖励之后\emph{回顾}：哪个事件最常出现在它之前？这种联系的度量是两个概率之差：
\begin{empheq}[box=\keybox]{equation}
M \;=\; P(\text{奖励前出现预示信号}) \;-\; P(\text{随机窗口中出现预示信号}) .
\label{eq:retro}
\end{empheq}
如果预示信号在没有奖励时也常出现，第二项就大，联系就弱。在这一理论中，多巴胺报告的不是预测误差，而是该事件在多大程度上可能是奖励的\emph{原因}。

回顾的机制要求的东西与第~\ref{sec:dofa}~章的三因子规则相同：每个事件都必须留下一条逐渐消退的迹，以便在奖励时刻能读出之前发生了什么。区别不在突触，而在\nt{DOPA}传送的内容。

在理论~\eqref{eq:retro}与误差~\eqref{eq:ctd}预测不同的实验中，多巴胺释放在若干情况下遵循前者~\cite{Jeong2022}。但在另一些实验中，学习过程中多巴胺的响应逐渐从奖励时刻移到预示信号时刻——正如按误差~\eqref{eq:ctd}学习所要求的~\cite{Amo2022}。哪一种理论描述了脑，目前还不知道。

\section{不只是关于奖励的误差}

最后一种扩展涉及误差是\emph{关于什么}的。如果把预期的果汁换成另一种同样有价值但味道不同的果汁，价值没有改变，误差~\eqref{eq:ctd}为零。然而\nt{DOPA}神经元对这种替换有响应~\cite{Takahashi2017}。人为诱发的多巴胺爆发还能在完全没有奖励的情况下，教会动物两个中性刺激之间的联系~\cite{Sharpe2017}。看来，\nt{DOPA}报告的不仅是奖励的预测误差，还有对将\emph{发生什么}的预测误差。

形式上这是~\eqref{eq:ctd}的简单推广：不再用单一的奖励流$r$，而用一组描述当前事件的特征$\varphi_k$——味道、位置、声音——每个特征都有自己的期望$M_k$和自己的误差~\cite{Gardner2018}：
\begin{empheq}[box=\keybox]{equation}
\delta_k(t) \;=\; \varphi_k(t) \;+\; \frac{\dd M_k}{\dd t} \;-\; \frac{M_k}{\tau_\gamma} ,
\qquad k=1,\dots,K .
\label{eq:vectortd}
\end{empheq}
奖励只是特征之一，误差向量教会的不仅是什么好，还有什么之后跟着什么——即世界模型。

\nt{DOPA}神经元的异质性与此一致：在同一个任务中，不同神经元携带不同的量——有的与奖励相关，有的与运动相关，有的与注意指向相关~\cite{Engelhard2019}。

\begin{keyblock}
\begin{proposition}[线束代替导线]
\label{prop:bundle}
\nt{DOPA}系统的信号不是一个数。它按神经元展开（分布~\eqref{eq:expectile}、不同特征~\eqref{eq:vectortd}、单独的危险维度），也按时间展开（快速误差和慢电平~\eqref{eq:vigor}）。标量误差~\eqref{eq:ctd}是这个信号的一阶近似，正如带阈值的加法器是神经元的一阶近似。
\end{proposition}
\end{keyblock}

\section{小结}

\begin{table}[t]
\centering
\footnotesize
\setlength{\tabcolsep}{4pt}
\renewcommand{\arraystretch}{1.15}
\begin{tabular}{>{\raggedright}p{2.5cm}>{\raggedright}p{3.2cm}>{\raggedright}p{3.6cm}>{\raggedright\arraybackslash}p{4.2cm}}
\toprule
理论 & 导线上传送什么 & 主要测量 & 已知情况 \\
\midrule
预测误差（第~\ref{sec:dofa}~章） & 标量$\delta$~\eqref{eq:ctd} & 爆发、向预示信号转移、低谷 & 已测量；一阶近似 \\
分布 & 一组不对称性各异的$\delta_i$ & 不同神经元翻转点不同 & 与实验相符；离散的作用是假说 \\
重要性 & $|\delta|$、新异性；危险维度 & 部分神经元对厌恶事件出现爆发 & 已测量；如何使用仍在讨论 \\
时间代价 & 慢电平$\bar R$ & 多巴胺加快动作；慢成分既见于输出端的释放，也见于脉冲频率 & 快慢分离已测得；$\bar R$是假说 \\
接近时的上升 & $V$，或估计修正时的$\delta$~\eqref{eq:ramp} & 平缓上升，走廊中瞬移时的跳变 & 上升已测得；解释仍在讨论 \\
向后看 & 因果联系的度量~\eqref{eq:retro} & 两种理论预测相悖的实验 & 结果相互矛盾 \\
误差向量 & 按特征的$\delta_k$~\eqref{eq:vectortd} & 对味道替换的响应；无奖励学习 & 已测量；向量形式是假说 \\
\bottomrule
\end{tabular}
\caption{\nt{DOPA}系统传送什么：第~\ref{sec:dofa}~章的标准图景及其现代扩展。}
\label{tab:dofaalt}
\end{table}

表~\ref{tab:dofaalt}中的理论并不相互排斥：几乎所有理论都把预测误差保留为基础，只是扩展了它的格式。真正与之竞争的只有向后看，这场争论要由实验来裁决。对工程师来说结论很简单：如果导线实际上是许多根通往不同目的地的导线，命题~\ref{prop:broadcastcost}中的广播代价就会下降。

\section{习题}

\begin{exercise}[\lvlA]
\label{ex:07:1}
\emph{单滴奖励的分布点。}大小为$1$的果汁以概率$p=0.2$到来，否则奖励为$0$。按~\eqref{eq:expectile}求$\kappa_i=0.2$、$0.5$、$0.8$时的$V_i$。这一奖励的中位数是多少？点~\eqref{eq:expectile}与分位数有何不同？
\end{exercise}

\begin{exercise}[\lvlB]
\label{ex:07:2}
\emph{由两个神经元得到分布。}奖励为$0$或$x$，取$x$的概率为$p$。两个按规则~\eqref{eq:asymmetric}工作的神经元报告$V(1/2)=0.25$和$V(0.8)=0.5$。求出$p$、$x$和奖励的标准差。$\kappa=0.2$的神经元会报告什么？
\end{exercise}

\begin{exercise}[\lvlB]
\label{ex:07:3}
\emph{仿真：由不对称得到分布。}仿真九个$\kappa_i=0.1,\dots,0.9$、按规则~\eqref{eq:asymmetric}工作的\nt{DOPA}神经元，$\alpha_i^+=2\kappa_i$，$\alpha_i^-=2(1-\kappa_i)$，奖励在$[0,1]$上均匀分布。$V_i$的离散如何依赖于$\eta$？要把这个分布与两个等概率的$0.25$和$0.75$区分开，需要多大的$\eta$？
\end{exercise}

\begin{exercise}[\lvlA]
\label{ex:07:4}
\emph{时间代价。}(a)~推导~\eqref{eq:vigor}，并求最优点的总代价。(b)~脑对$\bar R$的估计偏差为$k$倍。求$k=2$和$k=4$时的相对多付代价。多巴胺慢电平需要多精确地报告$\bar R$？
\end{exercise}

\begin{exercise}[\lvlB]
\label{ex:07:5}
\emph{瞬移时的爆发。}期望按~\eqref{eq:ramp}增长，$\tau_v=1\,\text{s}$，$\tau_\gamma=4\,\text{s}$；把动物从$T-2\,\text{s}$处瞬间移到$T-1\,\text{s}$处。以$R$为单位求$\delta$爆发的面积，以及它替代了瞬移前多少秒的上升。如果多巴胺报告的是$V$本身，瞬移会显示什么？
\end{exercise}

\begin{exercise}[\lvlB]
\label{ex:07:6}
\emph{设计：一根导线上的两个信号。}\nt{DOPA}导线上传送一个电平，其增长速度为每秒增加$s=1/60$个脉冲每秒，还有每次$A=3$个脉冲的事件。资格迹$\tau_e=1\,\text{s}$的突触从频率中减去其以$\tau_f$平滑的副本。$\tau_f$取何值时，事件损失不超过$10\,\%$，且在资格迹时间内电平的贡献小于事件贡献的$10\,\%$？
\end{exercise}

\begin{exercise}[\lvlB]
\label{ex:07:7}
\emph{实验设计：向后看对误差。}在一个窗口内预示信号出现的概率为$0.1$，其后奖励出现的概率为$0.8$。比较$\Delta P=P(\text{奖励}\mid\text{预示})-P(\text{奖励}\mid\text{无})$与~\eqref{eq:retro}中的$M$，如果(a)~加入无预示信号的奖励，每个窗口$0.08$；(b)~把无奖励的预示信号频率加倍。
\end{exercise}

\begin{exercise}[\lvlC]
\label{ex:07:8}
\emph{线束与广播的代价。}在习题~\ref{ex:06:8}的模型中，$N$个神经元分成$K$组，各有自己的导线，每根导线中渗入比例为$\rho$的其他组的误差。证明学习常数为$N/K+2+\rho^2N(1-1/K)$。当$K\to\infty$、$N=10^{10}$、$\rho=0.01$时，它对应多少次尝试？
\end{exercise}

\chapter{何时探索，何时决断：加入\nt{NORA}}
\chaptermark{加入\nt{NORA}}
\label{sec:nora}

\section{还缺什么}

在第~\ref{sec:dofa}~章中，网络借助噪声学习：活动的随机偏离就是搜索，而\nt{DOPA}报告这些偏离是否使结果变好。但那里还留下一个未设定的参数——搜索\emph{多少}。噪声太强，网络就会乱撞，不能利用已经知道的东西。噪声太弱，它就一次又一次地选择同一个选项，注意不到附近某处已经变好了。在强化学习中，这称为利用与探索之间的权衡，每个算法都为它设有一个旋钮。在第~\ref{sec:neurontypes}~章中我们推测，在脑中转动这个旋钮的是\nt{NORA}。

人类的去甲肾上腺素神经元有几万个（表~\ref{tab:types}），几乎全部集中在一个小核团中，而它们的输出几乎散布到全脑，不过已经发现，这个核团的不同部分在一定程度上服务于不同的区域~\cite{Poe2020}。这是一个比\nt{DOPA}还要窄的广播通道。它以两种模式工作~\cite{AstonJones2005}。在\emph{紧张性}模式下，神经元平稳地发放脉冲，频率从每秒零点几个到几个不等，这一频率随觉醒程度缓慢变化。在\emph{时相性}模式下，它们对当前任务中的重要事件以短促的簇放电响应，大约出现在做出决定的时刻。瞳孔直径是一个方便但不精确的测试点：它随这些神经元的活动变化，但也随其他区域的活动~\cite{Joshi2016}以及皮层中胆碱能输出的活动~\cite{Reimer2016}变化。瞳孔反映的是总体的唤醒水平，而不单单是\nt{NORA}。

\section{增益}

已测量的是：去甲肾上腺素对神经元背景活动的抑制强于对其刺激响应的抑制，信号与背景之比升高~\cite{Foote1975NE}。至于单个神经元特性曲线的斜率是否真的被乘上一个系数，实验并未给出结论。和第~\ref{sec:neurontypes}~章一样，我们采用一个能再现这一效应的简单模型：去甲肾上腺素改变接收者传递特性的斜率~\cite{ServanSchreiber1990}。于是弱的阈下输入（背景）产生的输出更小，强的输入产生的输出更大。设神经元以频率
\begin{empheq}[box=\keybox]{equation}
r \;=\; F\bigl(g\,(u-\theta)\bigr), \qquad F(z)=\frac{r_{\max}}{1+e^{-z}}
\label{eq:gain}
\end{empheq}
响应，其中$u$是输入电流，$\theta$是阈值，$g$是增益，在模型中由\nt{NORA}控制。$g$小时，频率在很宽的范围内平滑地跟随输入：神经元是\emph{模拟放大器}。$g$大时，几乎任何高于阈值的输入都给出满频率，低于阈值则为零：神经元是\emph{比较器}，几乎是数字元件（图~\ref{fig:gaincurves}）。一个旋钮就能让同一个神经元在两种模式间切换，而在电子学中这两种模式由不同的电路实现。

\begin{figure}[t]
\centering
\begin{tikzpicture}[>=Stealth,x=1.1cm,y=2.4cm]
\draw[->,gray!70!black] (-3.3,0) -- (3.4,0) node[below left,black] {\small $u$};
\draw[->,gray!70!black] (-3.3,0) -- (-3.3,1.2) node[left,black] {\small $r$};
\draw[densely dashed,gray!60!black] (0,0) -- (0,1.1);
\node[below,font=\small] at (0,0) {$\theta$};
\draw[densely dotted,gray!60!black] (-3.3,1) -- (3.2,1) node[right,black,font=\small] {$r_{\max}$};
\draw[very thick,blue!60!black] plot[domain=-3.2:3.2,samples=60] (\x,{1/(1+exp(-0.8*\x))});
\draw[very thick,orange!85!black] plot[domain=-3.2:3.2,samples=60] (\x,{1/(1+exp(-2*\x))});
\draw[very thick,red!70!black] plot[domain=-3.2:3.2,samples=120] (\x,{1/(1+exp(min(-8*\x,9)))});
\node[blue!60!black,font=\small,anchor=west] at (0.5,0.12) {小$g$：放大器};
\node[red!70!black,font=\small,anchor=east] at (-0.3,0.85) {大$g$：比较器};
\end{tikzpicture}
\caption{增益$g$取三个值时的传递特性~\eqref{eq:gain}。}
\label{fig:gaincurves}
\end{figure}

大增益能对抗噪声吗？不一定。设两个输入相差$\Delta u$，需要判断哪个更大。如果噪声$\sigma_{\text{前}}$是在放大器\emph{之前}加到输入上的，那么信号和噪声一起被放大，二者之比不变。如果噪声$\sigma_{\text{后}}$是在放大器\emph{之后}加入的——来自不可靠的突触和网络自身的随机脉冲，如第~\ref{sec:code}~章——那么信噪比
\begin{empheq}[box=\keybox]{equation}
d' \;=\; \frac{g\,\Delta u}{\sqrt{g^2\sigma_{\text{前}}^2+\sigma_{\text{后}}^2}}
\label{eq:gainsnr}
\end{empheq}
随$g$增大，直到$g\sigma_{\text{前}}$与$\sigma_{\text{后}}$相当为止~\cite{ServanSchreiber1990}。

\begin{keyblock}
\begin{proposition}[增益何时有用]
\label{prop:gainnoise}
\nt{NORA}提高增益，只能改善对放大器\emph{之后}、在网络内部产生的噪声的输入分辨。随输入一起到来的噪声，增益无法去除。
\end{proposition}
\end{keyblock}

\section{二选一中的增益}

回到第~\ref{sec:gaba}~章的选择：两组\nt{GLUT}神经元，输入为$I_1$、$I_2$，以强度$\beta$相互抑制。现在每组都有增益$g$：在选择方程~\eqref{eq:wta}中，括号内的表达式乘以$g$。只要两组都在工作，稳态频率由$r_1=g(I_1-\beta r_2)$、$r_2=g(I_2-\beta r_1)$求出。把两式相加和相减，得到和$r_1+r_2=g(I_1+I_2)/(1+g\beta)$与差$r_1-r_2=g(I_1-I_2)/(1-g\beta)$。二者之比就是网络偏好一个输入胜过另一个的鲜明程度：
\begin{empheq}[box=\keybox]{equation}
\frac{r_1-r_2}{r_1+r_2} \;=\; \varepsilon\,\frac{1+g\beta}{1-g\beta}, \qquad \varepsilon=\frac{I_1-I_2}{I_1+I_2}, \quad g\beta<1 .
\label{eq:wtagain}
\end{empheq}
输入的微小差别$\varepsilon$被放大$(1+g\beta)/(1-g\beta)$倍，当$g\beta$接近一时，这个倍数无限增大。当$g\beta>1$时，两组都工作的状态是不稳定的，正如命题~\ref{prop:wta}：一组获胜，另一组沉默（图~\ref{fig:pitchfork}）。

\begin{figure}[t]
\centering
\begin{tikzpicture}[>=Stealth,x=3.2cm,y=1.5cm]
\draw[->,gray!70!black] (0,0) -- (2.15,0) node[below left,black] {\small $g\beta$};
\draw[->,gray!70!black] (0,-1.25) -- (0,1.35) node[above,black] {\small $(r_1-r_2)/(r_1+r_2)$};
\draw[gray!60!black] (1,-0.04) -- (1,0.04); \node[below right,font=\small] at (1,0) {$1$};
\node[left,font=\scriptsize] at (0,1) {$1$}; \node[left,font=\scriptsize] at (0,-1) {$-1$};
\draw[very thick,blue!60!black] plot[domain=0:0.905,samples=60] (\x,{0.05*(1+\x)/(1-\x)}) -- (1,1) -- (2.05,1);
\draw[very thick,blue!60!black] (1.105,-1) -- (2.05,-1);
\draw[thick,densely dashed,gray!60!black] (1,0) -- (2.05,0);
\node[font=\small,align=center] at (0.45,-0.62) {在斟酌：\\两组\\都在工作};
\node[font=\small,align=center] at (1.55,0.45) {已决断：\\只有一组工作};
\node[font=\small,blue!60!black,anchor=south west] at (1.05,-0.97) {弱输入先胜出——就保持下去};
\end{tikzpicture}
\caption{不同增益下的二选一，输入相差$\varepsilon=0.05$。当$g\beta>1$时两个决定都是稳定的（实线；弱输入获胜在$g\beta>(1+\varepsilon)/(1-\varepsilon)\approx1.1$时成为可能），网络保持已经做出的那个；对称状态（虚线）不稳定。}
\label{fig:pitchfork}
\end{figure}

\begin{keyblock}
\begin{proposition}[增益切换选择模式]
\label{prop:gainwta}
在相互抑制为$\beta$的选择网络中，增益$g$使网络越过边界$g\beta=1$。低于它时网络在“斟酌”：两组都在工作，输入之差按~\eqref{eq:wtagain}被放大。高于它时网络“已决断”：只有一组工作，决定保持不变。\nt{NORA}改变$g$，就能在这两种模式之间切换网络，而不触动任何一个权重。
\end{proposition}
\end{keyblock}

\section{增益即温度}

现在给选择加上在网络内部、放大器之后产生的噪声。哪一组获胜，取决于量$g(u_A-u_B)+\eta$的符号，其中$u_A-u_B$是输入之差，即第~\ref{sec:dofa}~章学到的权重之差，$\eta$是尺度为$\sigma$的噪声。如果噪声服从逻辑斯谛分布（对高斯分布，曲线几乎相同），选择$A$的概率为
\begin{empheq}[box=\keybox]{equation}
P(A) \;=\; \frac{1}{1+e^{-g\,(u_A-u_B)/\sigma}} .
\label{eq:softmax}
\end{empheq}
程序员会认出这是温度为$T=\sigma/g$的 softmax 选择。增益就是逆温度。$g$小时，即使明显更好的选项也只比差的选项稍常被选中：网络在大量探索。$g$大时，已知的最佳选项几乎总被选中：网络在利用已知。

需要说明~\eqref{eq:wtagain}和~\eqref{eq:softmax}中的$g$是什么：它是选择时刻对当前任务输入之差的\emph{有效}增益。\nt{NORA}的两种模式对它的作用相反。决策时刻的时相性簇放电提高它，网络巩固选择。反之，高紧张性活动伴随着探索：人在随机选择之前基线瞳孔更大~\cite{JepmaNieuwenhuis2011}，而在大鼠中，\nt{NORA}向前扣带皮层的输入使选择转入随机模式~\cite{Tervo2014stochastic}。可见，高紧张性\nt{NORA}对应\emph{小}的有效$g$，簇放电则对应大的有效$g$。

\section{探索多少}

哪种增益最好？我们在模型上做了检验。网络按~\eqref{eq:softmax}从两个动作中选一个；每个动作以某个概率带来奖励，网络像第~\ref{sec:dofa}~章那样，根据所选动作的奖励来估计这个概率。世界不时变化：两个概率都重新抽取。变化的可能是所选的动作，也可能是网络很久没有尝试的那个；后者只有通过探索才能发现。图~\ref{fig:invertedU}显示了在不同的恒定$g$下网络获得的最佳可能奖励的比例。

\begin{figure}[t]
\centering
\begin{tikzpicture}[>=Stealth,x=1.2cm,y=1cm]
\draw[->,gray!70!black] (0,0) -- (7.6,0) node[below left,black] {\small $g$};
\draw[->,gray!70!black] (0,0) -- (0,4.4) node[above,black,font=\small] {最佳奖励的比例};
\foreach \x/\l in {0/0.5,1/1,2/2,3/4,4/8,5/16,6/32,7/64}{\draw[gray!60!black] (\x,-0.06) -- (\x,0.06); \node[below,font=\scriptsize] at (\x,0) {\l};}
\foreach \v/\y in {0.8/0.8,0.9/2.4,1.0/4.0}{\draw[gray!60!black] (-0.06,\y) -- (0.06,\y); \node[left,font=\scriptsize] at (0,\y) {\v};}
\draw[very thick,blue!60!black] plot[mark=*,mark size=1.3pt] coordinates {(0,0.368) (1,0.880) (2,1.728) (3,2.784) (4,3.456) (5,3.616) (6,3.488) (7,3.264)};
\draw[very thick,red!70!black] plot[mark=*,mark size=1.3pt] coordinates {(0,0.368) (1,0.672) (2,1.184) (3,1.840) (4,2.176) (5,1.920) (6,1.456) (7,1.104)};
\node[blue!60!black,font=\small,anchor=south] at (5,3.7) {平稳的世界};
\node[red!70!black,font=\small,anchor=north] at (5.1,1.25) {快速变化的世界};
\draw[->,thick,blue!60!black] (5,3.25) -- (5,3.55);
\draw[->,thick,red!70!black] (4,1.8) -- (4,2.1);
\node[font=\small\itshape,gray!35!black,anchor=west] at (0.15,2.3) {盲目探索};
\node[font=\small\itshape,gray!35!black] at (6.45,2.55) {察觉不到变化};
\end{tikzpicture}
\caption{恒定增益$g$下获得的最佳可能奖励的比例（$g$轴为对数刻度）。蓝线：世界平均每一千次尝试变化一次；红线：每二十次变化一次。箭头标出最佳$g$：在快速变化的世界中它小一半。为 60 次运行、每次 4000 次尝试的平均。}
\label{fig:invertedU}
\end{figure}

$g=0.5$时网络几乎不利用知识，只得到可能奖励的约四分之三；$g$很大时则错过变化。最佳值：世界每一千次尝试变化一次时为$g=16$（比例$0.976$），每二十次变化一次时为$g=8$（$0.886$）。

\begin{keyblock}
\begin{proposition}[倒 U 形]
\label{prop:explore}
网络收集到的奖励与增益呈倒 U 形关系：$g$太小，尝试浪费在探索上；$g$太大，察觉不到变化。世界变化越快，最佳$g$越小。
\end{proposition}
\end{keyblock}

倒 U 形在实验中也能观察到：在\nt{NORA}神经元紧张性活动适中、时相性响应清晰时，动物完成任务最好，而在紧张性活动很高时，它们会分心，在任务之间切换~\cite{AstonJones2005}。这与上一节中模式的区别相符：决策时刻的时相性簇放电恰在需要选择时给出大的有效$g$，而没有簇放电的高紧张性活动会不加区分地放大一切，包括无关输入，因此对当前任务的有效$g$下降——这就是探索。

\section{谁来转动旋钮}

既然最佳增益取决于世界变化得多快，最好能对它进行调节。但\nt{NORA}神经元怎么知道世界变了？一个自然的标志是\emph{意外}：预测误差比平常大了。重要的是区分两种不确定性。如果奖励只是随机的，误差总是很大，这是预料之中的。如果误差突然比平常的大小增大了，那就说明有什么东西变了。按照一种假说，\nt{NORA}报告的正是这种非预期的不确定性，而预期的不确定性由\nt{ACET}报告~\cite{YuDayan2005}。

拿这个信号做什么？可以降低增益，多做探索。可以提高学习速率，以便更快地忘掉过时的估计：实验表明，在突然的变化之后人学得更快，同时瞳孔扩大~\cite{Nassar2012}；提醒一下，瞳孔不仅是\nt{NORA}的指标，也是\nt{ACET}的指标~\cite{Reimer2016}。也可以二者同时进行：按照另一种假说，\nt{NORA}的时相性簇放电起着\emph{中断}的作用——网络据此复位当前状态，重新适应新环境~\cite{BouretSara2005}，就像处理器的公共中断线。

坦白地说说我们没找到什么。在模型中，我们试了两种简单的调节方式：当平滑后的误差超过其长期平均值时，降低$g$或提高学习速率。在一个时而长期稳定、时而频繁变化的世界中，两种方式的最佳设置都得到最佳奖励的$0.926$--$0.927$——几乎与最佳恒定$g$（$g=8$时为$0.925$）或恒定但足够大的学习速率（$0.928$）相同，而设置不当时则更低。图~\ref{fig:invertedU}中的曲线在顶点附近很平坦，在这样的世界中几乎没有什么可调的。调节应当在不同状态要求截然不同设置的地方才能带来回报。\nt{NORA}究竟实现了怎样的控制律，目前尚未确定。

\section{小结}

\begin{table}[t]
\centering
\footnotesize
\setlength{\tabcolsep}{4pt}
\renewcommand{\arraystretch}{1.15}
\begin{tabular}{>{\raggedright}p{3.0cm}>{\raggedright}p{4.4cm}>{\raggedright\arraybackslash}p{6.0cm}}
\toprule
\nt{NORA}做什么 & 怎样做 & 已知情况 \\
\midrule
改变神经元响应的斜率 & \eqref{eq:gain}中的增益$g$ & 信噪比提高已测得；乘法模型是一种简化 \\
切换选择 & 使网络越过$g\beta=1$ & 由模型~\eqref{eq:wtagain}推出；没有对阈值的直接测量 \\
决定探索多少 & $g$是逆温度~\eqref{eq:softmax}；紧张性活动对应小的有效$g$（探索），簇放电对应大的（决断） & 倒 U 形和两种工作模式已测得；与探索的联系是有充分依据的假说 \\
报告变化 & 非预期的不确定性 & 变化后瞳孔和学习速率增大——已测得；这是\nt{NORA}的信号——假说 \\
中断并复位 & 对意外事件的时相性簇放电 & 对重要事件的簇放电已测得；“中断”是假说 \\
\bottomrule
\end{tabular}
\caption{\nt{NORA}为\nt{GLUT}、\nt{GABA}和\nt{DOPA}网络增加了什么。}
\label{tab:nora}
\end{table}

\nt{DOPA}告诉网络权重\emph{往哪里}改。\nt{NORA}不改变任何一个权重，却决定网络\emph{如何}利用已知的东西（表~\ref{tab:nora}）：一个旋钮——增益——把神经元从放大器变成比较器，把选择网络从“斟酌”变成“已决断”，把选择从探索变成利用。\nt{NORA}如何得知世界变化得多快，仍是一个悬而未决的问题。

\section{习题}

\begin{exercise}[\lvlA]
\label{ex:08:1}
\emph{全脑一个旋钮。}（a）~根据表~\ref{tab:types}估计每个\nt{NORA}神经元对应多少个脑神经元。（b）~放大器前的噪声为$\sigma_{\text{前}}=0.5$，之后为$\sigma_{\text{后}}=4$。$g$为多少时，\eqref{eq:gainsnr}中的$d'$达到其极限的$90\,\%$？
\end{exercise}

\begin{exercise}[\lvlB]
\label{ex:08:2}
\emph{带增益的选择。}$\tau\,\dd r_1/\dd t=-r_1+g[I_1-\beta r_2]_+$，$\tau\,\dd r_2/\dd t=-r_2+g[I_2-\beta r_1]_+$，$\tau=20\ms$。（a）~$g\beta=0.5$和$0.9$时，$r_1-r_2$衰减需要多长时间？（b）~当$\varepsilon=0.05$时，求出一个$g\beta$值：低于它两组都工作，高于它弱输入的获胜是稳定的。
\end{exercise}

\begin{exercise}[\lvlA]
\label{ex:08:3}
\emph{由噪声得到 softmax。}$K$组中每组各得到自己的 Gumbel 噪声，$P(\eta_k<z)=\exp(-e^{-z/\sigma})$，$gu_k+\eta_k$最大者获胜。求$P(k)$。当$u_A-u_B=0.1$、$\sigma=1$时，$g$为多少才能使两个选项中较好的那个在$90\,\%$的情况下被选中？
\end{exercise}

\begin{exercise}[\lvlB]
\label{ex:08:4}
\emph{估计最佳增益。}在图~\ref{fig:invertedU}的模型中，变化之后的奖励概率在$[0,1]$上均匀分布。网络每$e^{g\Delta}$次尝试试一次较差的动作；令其等于$1/h$，得$g^*\approx\ln(1/h)/\bar\Delta$。求$\bar\Delta=\langle|p_A-p_B|\rangle$以及$h=0.001$、$0.01$、$0.05$时的$g^*$。
\end{exercise}

\begin{exercise}[\lvlB]
\label{ex:08:5}
\emph{仿真：增益与变化速度。}用\texttt{models/\allowbreak nora\_explore.py}的副本（它把文件写到当前目录）求$h$从$0.001$到$0.2$时的$g^*(h)$。习题~\ref{ex:08:4}的估计在哪里成立？为什么在快速变化时$g^*$不再减小？
\end{exercise}

\begin{exercise}[\lvlB]
\label{ex:08:6}
\emph{设计：选择网络的中断。}习题~\ref{ex:08:2}的网络，$\beta=2$，$\tau=20\ms$，$g=1$，保持着$r_1=0.9$、$r_2=0$，尽管现在输入为$I_1=0.9$、$I_2=1.0$。\nt{NORA}在时间$T_p$内把$g$降到$g_{\text{lo}}$。对$g_{\text{lo}}=0$解析地求出最小的$T_p$，对$g_{\text{lo}}=0.25$用模型求出。
\end{exercise}

\begin{exercise}[\lvlC]
\label{ex:08:7}
\emph{调节何时有回报。}一个先知知道当前是平稳时期还是多变时期，并在每个时期设定各自的$g$。（a）~用模型求它在本章的世界中（每个时期$1000$次尝试，$h=0.001$和$0.05$）相对于最佳恒定$g$的收益。（b）~构造一个收益更大的世界，并求按意外进行的调节能拿到其中多少。
\end{exercise}

\chapter{何时写入，何时读出：加入\nt{ACET}}
\chaptermark{加入\nt{ACET}}
\label{sec:acet}

\section{还缺什么}

存储芯片除了地址线和数据线，还有一条线——写使能。它关闭时，存储器只能读；打开时，数据线上的内容被写入。没有这条线，存储器就无法工作：如果每次读都同时写入点什么，读出的内容，包括读错的内容，会立即被写回去。

在脑中，这个问题比在芯片中更尖锐。在第~\ref{sec:glutcomputer}~章中，存储器是一组\nt{GLUT}神经元，由其自身的连接维持在开启状态（图~\ref{fig:bistable}）。在第~\ref{sec:dofa}~章中我们看到，这些连接在工作过程中直接按赫布规则学习。也就是说，同一批突触既保存旧内容又接收新内容，而且没有单独的写入阶段，也没有时钟。网络怎样区分需要记住眼前所见的时刻和需要回忆已知内容的时刻？在第~\ref{sec:neurontypes}~章中我们推测，这条线由\nt{ACET}掌控。本章将检验这样一条线应当如何工作，以及为什么缺了它不行。

\section{一根导线的三种作用}

在脑内部工作的乙酰胆碱神经元不多，集中在几个小核团中，而它们的输出遍布整个皮层。这是一个广播通道，和\nt{DOPA}、\nt{NORA}一样。皮层中的乙酰胆碱水平在专注清醒时高，在慢波睡眠时低~\cite{Hasselmo1999}。与多巴胺一样，信号的含义由接收者的受体决定（命题~\ref{prop:receptor}）：乙酰胆碱既有打开通道的快速受体，也有在细胞内启动反应的慢速受体。对我们重要的是三种作用。

\emph{第一：削弱内部连接。}在嗅皮层脑片上的实验表明，乙酰胆碱强烈抑制同一区域内神经元之间连接的传递，而几乎不影响从外部进入该区域的输入~\cite{HasselmoBower1992}。

\emph{第二：增强输入。}乙酰胆碱提高神经元的兴奋性，并易化某些输入通路的传递；来自感觉器官的信号变得比网络自身的活动更响亮~\cite{Hasselmo2006}。

\emph{第三：使权重更容易改变。}乙酰胆碱水平高时，连接更容易改变~\cite{Hasselmo2006}。

对工程师来说，这三种作用是同一个操作：从“读”模式切换到“写”模式。网络不再听自己，开始听输入，并记住它听到的东西。

\section{补全的网络}

取$N$个相互连接的\nt{GLUT}神经元。用赫布规则~\cite{Hebb1949}在它们的连接中存入几个模式——每个模式是$pN$个同时活跃的神经元的集合。如果向这样的网络输入模式的一半，连接会激发缺失的另一半：网络根据部分\emph{补全}模式，正如图~\ref{fig:bistable}中的神经元组自我维持一样。这就是联想记忆：内容的一部分充当地址。\nt{GABA}像第~\ref{sec:gaba}~章那样把活跃神经元的总数维持在$pN$左右，使两个模式不能同时点亮。

记乙酰胆碱水平为$a$，取值从$0$到$1$，并把它的三种作用直接写进模型：
\begin{empheq}[box=\keybox]{equation}
\tau\frac{\dd r_i}{\dd t} = -r_i + F\Bigl(\tfrac{1+a}{2}\,x_i + (1-a)\sum_j W_{ij}\,r_j - g\Bigr),
\qquad
\frac{\dd W_{ij}}{\dd t} = \eta\,a\,(r_i-p)(r_j-p) .
\label{eq:acetnet}
\end{empheq}
这里$r_i$是神经元$i$的频率，$x_i$是它来自感觉器官的输入，$F$是带阈值的传递特性，$g$是来自\nt{GABA}的总体抑制，当活跃神经元多于$pN$时增大。因子$\frac{1+a}{2}$是输入的增强，$1-a$是内部连接的削弱，$\eta a$是学习速率。第二个方程是赫布规则~\eqref{eq:hebb}的一种便于处理稀疏模式的形式~\cite{Tsodyks1988}：两个神经元都比平时活跃时权重增大，只有一个活跃时权重减小。权重的负值部分一如既往由\nt{GABA}中介实现。没有时钟：神经元和权重都连续变化。

\begin{figure}[t]
\centering
\begin{tikzpicture}[>=Stealth,
  nrn/.style={circle,draw,very thick,fill=blue!6,minimum size=0.8cm,inner sep=0pt,font=\small},
  acet/.style={circle,draw,very thick,double,double distance=1.2pt,fill=orange!15,minimum size=1.35cm,inner sep=0pt,font=\scriptsize},
  bc/.style={->,thick,densely dashed,orange!85!black}]
\foreach \i/\y in {1/2.4,2/1.2,3/0}{
  \node[nrn] (N\i) at (3,\y) {$r_\i$};
  \draw[->,thick,blue!60!black] (0,\y) node[left,black] {\small $x_\i$} -- (N\i);}
\node[left,align=right,font=\small] at (-0.5,1.2) {来自\\感觉器官};
\draw[->,thick,gray!60!black] (N1) to[bend left=30] (N2);
\draw[->,thick,gray!60!black] (N2) to[bend left=30] (N1);
\draw[->,thick,gray!60!black] (N2) to[bend left=30] (N3);
\draw[->,thick,gray!60!black] (N3) to[bend left=30] (N2);
\draw[->,thick,gray!60!black] (N1.east) .. controls (4.6,2.0) and (4.6,0.4) .. (N3.east);
\node[right,font=\small,gray!40!black,align=left] at (4.7,1.2) {内部\\连接$W$};
\node[acet] (A) at (1.2,-1.9) {\nt{ACET}};
\draw[bc] (A) -- (1.6,-0.15);
\node[font=\small,orange!60!black,anchor=east] at (0.95,-0.9) {$+$输入};
\draw[bc] (A) -- (4.3,0.6);
\node[font=\small,orange!60!black,anchor=west] at (4.3,0.1) {$-$内部连接，$+$学习速率};
\node[right,font=\small,align=left] at (2.2,-1.9) {高水平：“写”\\低水平：“读”};
\end{tikzpicture}
\caption{\nt{ACET}作为写使能线。神经元$r_i$从感觉器官接收输入$x_i$（蓝色箭头），并通过存有模式的内部连接$W$（灰色）相互兴奋。乙酰胆碱（虚线箭头）增强输入、削弱内部连接并使权重更容易改变，如~\eqref{eq:acetnet}所示。}
\label{fig:acetnet}
\end{figure}

\section{写入与读出相互妨碍}

我们在一个写入与读出确实相冲突的任务上检验模型~\eqref{eq:acetnet}。一个$200$个神经元的网络已经存有五个模式，每个$20$个活跃神经元。现在需要记住第六个，它与一个旧模式有一半重合：有十个神经元是共同的。这种情况经常发生：新房间像熟悉的房间，新面孔像旧面孔。

\emph{低$a$时的写入。}首先把乙酰胆碱的前两种作用与第三种分开：学习速率保持满值，只改变输入的增强和内部连接的削弱。$a=0$时，输入点亮共同的十个神经元，它们经内部连接点亮\emph{旧}模式的其余一半。\nt{GABA}不允许两者同时点亮，于是有自身连接支撑的旧模式挤掉只靠输入维持的新模式。网络看到的不是眼前的东西，而是记得的东西——而赫布规则勤勤恳恳地把看到的写了下来。

结果，新模式根本没有记住：凭它的特有一半，网络什么也想不起来。而旧模式被破坏了：用它自己的特有一半唤起时，它平均带出十个外来神经元中的六个。$a=0.1$时新模式已能记住，但与旧模式融合，破坏甚至更严重：两个模式各自用自己的一半唤起时，都会带出约八个外来神经元；从$a=0.2$开始写入是干净的：回忆时多余的神经元不到一个（十次运行的平均）。

\emph{采用真实学习速率的写入。}现在让乙酰胆碱像~\eqref{eq:acetnet}中那样也设定学习速率。一次呈现就完整记住新模式，只在$a\ge0.9$时发生；$a=0.75$时记住十分之八，$a\le0.5$时完全记不住。

\emph{读出。}向存有五个干净写入模式的网络输入某个模式的一半，然后再撤去。$a\le0.2$时网络把模式完整补全并自行维持；$a=0.3$时几乎完整；$a\ge0.4$时内部连接被削弱得太厉害，提示一消失，活动就熄灭（图~\ref{fig:acetsim}）。

\begin{figure}[t]
\centering
\begin{tikzpicture}[>=Stealth,x=9cm,y=3.2cm]
\draw[->,gray!70!black] (0,0) -- (1.1,0) node[right,black] {\small $a$};
\draw[->,gray!70!black] (0,0) -- (0,1.2);
\foreach \x in {0,0.2,0.4,0.6,0.8,1.0}{\draw[gray!60!black] (\x,-0.02) -- (\x,0.02); \node[below,font=\scriptsize] at (\x,0) {$\x$};}
\foreach \y in {0,0.5,1}{\draw[gray!60!black] (-0.01,\y) -- (0.01,\y); \node[left,font=\scriptsize] at (0,\y) {$\y$};}
\node[rotate=90,font=\small] at (-0.1,0.6) {模式恢复的比例};
\draw[very thick,blue!60!black] plot[mark=*,mark size=1.6pt] coordinates {(0,1.00) (0.1,1.00) (0.2,1.00) (0.3,0.96) (0.4,0.04) (0.5,0.00) (0.75,0.00) (1.0,0.00)};
\draw[very thick,red!70!black] plot[mark=*,mark size=1.6pt] coordinates {(0.1,0.00) (0.2,0.00) (0.3,0.00) (0.5,0.00) (0.75,0.80) (0.9,1.00) (1.0,1.00)};
\node[font=\small,blue!60!black,anchor=west,align=left] at (0.03,0.5) {读出：\\凭一半\\恢复模式};
\node[font=\small,red!70!black,anchor=east,align=right] at (1.0,0.35) {写入：\\一次记住\\新模式};
\end{tikzpicture}
\caption{模型~\eqref{eq:acetnet}：$200$个神经元，每个模式$20$个活跃神经元，十次运行的平均。蓝点为读出：撤去提示后，凭一半恢复了模式的多大比例。红点为写入：若乙酰胆碱也设定学习速率，一个与旧模式一半相似的新模式在一次呈现后能回忆起多大比例。读出在$a\le0.3$时有效，写入在$a\ge0.9$时有效。}
\label{fig:acetsim}
\end{figure}

\begin{keyblock}
\begin{proposition}[写入和读出需要不同的模式]
\label{prop:writeread}
在模型~\eqref{eq:acetnet}中，同一批连接既保存旧内容又学习新内容，而乙酰胆碱同时设定学习速率，此时不存在一个$a$值能使写入和读出同样顺利。写入时必须削弱内部连接，否则网络写下的不是输入，而是回忆出的东西；读出时必须加强它们，否则网络无法根据部分补全模式。需要一个开关，而一根广播导线就能充当它。
\end{proposition}
\end{keyblock}

如果乙酰胆碱不改变学习速率，那么窄带$0.2\le a\le0.3$对两件事都适用。但那样的话，网络在每次读出时也会学习，也就是把读出的内容连同其错误重新写入。在学习期间抑制内部连接这一想法，来自根据脑片测量建立的模型~\cite{HasselmoAndersonBower1992}。上面的数字属于我们简化的模型；脑中各模式的边界究竟在哪里，还不知道。

\section{该在多大程度上相信眼睛}

“写/读”开关是一个更一般问题的极端情况：在多大程度上相信输入，在多大程度上相信记忆？对这个问题工程师有确切的答案，它也说明了为什么一根导线能同时控制模式和学习速率。

设网络跟踪一个缓慢游走的量$s(t)$：它的方差每秒增加$q$。感觉器官报告$y(t)=s(t)+$噪声，噪声强度为$R$。连续时间中的最佳估计$\hat s$由卡尔曼—布西滤波器给出~\cite{KalmanBucy1961}：
\begin{empheq}[box=\keybox]{equation}
\frac{\dd\hat s}{\dd t} \;=\; \frac{P}{R}\,\bigl(y-\hat s\bigr),
\qquad
\frac{\dd P}{\dd t} \;=\; q-\frac{P^2}{R} ,
\label{eq:kalman}
\end{empheq}
其中$P$是估计误差的方差，即网络自己对所知内容有多不确定。第一个方程就是神经元模型~\eqref{eq:glutneuron}中膜的那种 RC 电路，只是时间常数$R/P$是变化的。网络不确定时，$P$大，时间常数小，估计迅速跟随输入：这是“写”。网络确定时，$P$小，估计几乎不听输入，依赖记忆：这是“读”。

在稳态下$P=\sqrt{qR}$，记忆的时间常数为$\sqrt{R/q}$：如果世界每一百秒漂移一个单位，$q=0.01$，而感觉噪声$R=1$，那么记忆应保持约十秒。

这里的关键是，一个数$P/R$同时决定了两件事：输入相对于记忆的权重，以及所存估计变化的速度。这恰恰就是乙酰胆碱在~\eqref{eq:acetnet}中所做的。按照一种假说~\cite{YuDayan2005}，乙酰胆碱向网络报告的正是一个类似$P$的量——\emph{预期的}不确定性，即网络事先就知道的那种。这个通道并不总是缓慢的：一部分\nt{ACET}神经元在二十毫秒左右内对奖励和惩罚作出响应，强化越出乎意料，响应越强~\cite{Hangya2015}。

\begin{keyblock}
\begin{proposition}[输入权重和学习速率共用一个旋钮]
\label{prop:kalman}
在最优滤波器~\eqref{eq:kalman}中，输入与记忆混合时的权重和学习速率是同一个量$P/R$。因此用一个信号控制它们是合理的。当世界的性质不变时，这个信号趋于恒定水平$\sqrt{q/R}$，只有当世界的性质变化时才需要调节它。
\end{proposition}
\end{keyblock}

命题的第二句解释了第~\ref{sec:nora}~章的结果：在那里，按意外调节学习速率并没有胜过最佳的恒定速率——在变化性质不变的世界中，最佳学习速率就是恒定的。当世界突然变得不同时，才会有收益。这时估计突然远离输入，而$P$对此一无所知。正确的做法是急剧提高$P$，从而进入写入模式。按照我们在第~\ref{sec:nora}~章讨论过的假说，报告这种\emph{非预期}变化的是\nt{NORA}~\cite{YuDayan2005}。分工便自然形成：\nt{NORA}发现世界变了，\nt{ACET}让网络保持在写入模式，直到新环境被学会。

\section{睡眠作为读出模式}

如果高乙酰胆碱水平是写入模式，那么在低水平时网络做什么？在慢波睡眠中，乙酰胆碱水平下降，内部连接摆脱抑制，来自感觉器官的输入几乎没有。处于无提示读出模式的网络回放其中记录的内容。活动记录显示，白天一起工作的神经元在慢波睡眠中再次一起发放~\cite{WilsonMcNaughton1994}，甚至按相同的顺序~\cite{Skaggs1996}。按照一种假说~\cite{Hasselmo1999}，这些回放把白天快速学到的内容转写到长期记忆中（人为延长短暂的回放爆发能改善记忆~\cite{FernandezRuiz2019}）：白天从输入写入，夜间从内部重写。对工程师来说，这类似于白天写入快速缓冲区、夜间把它转存到慢速的永久存储器。

\section{小结}

\begin{table}[t]
\centering
\footnotesize
\setlength{\tabcolsep}{4pt}
\renewcommand{\arraystretch}{1.15}
\begin{tabular}{>{\raggedright}p{3.2cm}>{\raggedright}p{4.2cm}>{\raggedright\arraybackslash}p{6.0cm}}
\toprule
\nt{ACET}做什么 & 怎样做 & 已知情况 \\
\midrule
削弱内部连接 & \eqref{eq:acetnet}中的因子$1-a$ & 已在脑片上测得 \\
增强输入 & 因子$\frac{1+a}{2}$ & 兴奋性和输入通路传递的提高已测得 \\
使学习更容易 & 速率$\eta a$ & 已测量 \\
切换“写/读” & 命题~\ref{prop:writeread} & 由模型推出；没有对模式的直接测量 \\
报告预期的不确定性 & \eqref{eq:kalman}中的$P$ & 假说 \\
让网络在睡眠中空出来进行回放 & 慢波睡眠中的低$a$ & 睡眠中的低水平和回放已测得；转写到长期记忆是假说 \\
\bottomrule
\end{tabular}
\caption{\nt{ACET}为\nt{GLUT}、\nt{GABA}、\nt{DOPA}和\nt{NORA}网络增加了什么。}
\label{tab:acet}
\end{table}

\nt{DOPA}告诉网络权重往哪里改，\nt{NORA}告诉它以多大的决断利用已知的东西，而\nt{ACET}告诉它该听世界还是听自己（表~\ref{tab:acet}）。这是表~\ref{tab:types}中最后一条广播控制线：写使能。没有它，用同一批连接既保存模式又读出模式的存储器，每读一次就会损坏自己一次。

\section{习题}

\begin{exercise}[\lvlA]
\label{ex:09:1}
\emph{记住多久。}$q=0.01$、$R=1$的滤波器~\eqref{eq:kalman}记忆约$10$~s。求$P_\infty$和记忆$\sqrt{R/q}$，如果（a）~传感器噪声的幅度增大一倍；（b）~世界游走的速度快一百倍。（c）~噪声需增大多少倍，网络才能记住一分钟？
\end{exercise}

\begin{exercise}[\lvlB]
\label{ex:09:2}
\emph{变化之后的写入模式。}世界变了，$q=0.01$、$R=1$的滤波器~\eqref{eq:kalman}的不确定性跃升到$P_0\to\infty$。（a）~$P$以什么时间常数回到$P_\infty$？（b）~多少秒后学习速率超出稳态值不超过$10\,\%$？
\end{exercise}

\begin{exercise}[\lvlB]
\label{ex:09:3}
\emph{恒定学习速率。}网络按$\dd\hat s/\dd t=(y-\hat s)/T$学习；世界游走，$\dd s/\dd t=w$，$y=s+n$，其中$w$和$n$是强度为$q$和$R$的白噪声。求最佳$T$及此时误差的方差。如果$T$错了$2$倍和$10$倍，方差增大多少倍？
\end{exercise}

\begin{exercise}[\lvlB]
\label{ex:09:4}
\emph{写入的边界在哪里。}在\texttt{models/\allowbreak acet\_memory.py}中，阈值$\theta=0.35$，模式内权重$w=0.041$（理想值$0.045$）。新模式与旧模式共有$m=10$个神经元。在~\eqref{eq:acetnet}中$a$为多少时，旧模式的其余神经元不会被这$m$个点亮（阈值是陡的）？
\end{exercise}

\begin{exercise}[\lvlB]
\label{ex:09:5}
\emph{仿真：记忆容量。}修改\texttt{models/\allowbreak acet\_memory.py}中的函数\texttt{read\_sweep}：在$a=1$时写入$M$个模式（$M$从$5$到$100$），然后在$a=0$时凭一半读出前十个。$M$为多少时开始出错？极限$M_{\max}$如何依赖于$N$和$p$？
\end{exercise}

\begin{exercise}[\lvlB]
\label{ex:09:6}
\emph{设计：无泄漏的写入。}学习速率为$\eta a$时，网络在读出时也在学习。提出一个单调的$\eta(a)\le\eta$，在$a\le0.3$时为零，在$a\ge0.9$时不差于$\eta a$。检验：在$a=0.2$下用含三个外来神经元的提示读出$20$次之后，有多少外来神经元进入了模式？
\end{exercise}

\begin{exercise}[\lvlC]
\label{ex:09:7}
\emph{夜间如何转存缓冲区。}（a）~睡眠中$a=0.05$，一次回放持续$0.1\,\text{s}$，而白天$a=1$时为$0.2\,\text{s}$。在速率$\eta a$和习题~\ref{ex:09:6}的$\eta(a)$下，多少次回放才能累积一次白天呈现所产生的权重变化？（b）~永久存储的学习速度比缓冲区慢$100$倍，每晚听到模式$20$次，每次$0.1\,\text{s}$。需要多少个夜晚？
\end{exercise}

\chapter{整台机器}
\chaptermark{整台机器}
\label{sec:synthesis}

\section{我们搭起了什么}

我们每次从表~\ref{tab:types}中取出一种神经元类型，追问它给这台计算机器增加了什么。规则每次都一样：只用活体中真实存在的东西，也不借助任何全局时钟。现在把各个部件装成一台机器，看看它们如何协同工作。

第~\ref{sec:neurontypes}~章的图景得到了证实，也变得更精确。庞大的\nt{GLUT}神经元寻址网络承载数据。\nt{GABA}神经元控制这些数据的流动。网络之上悬着四条广播通道，每条各有一个旋钮：\nt{DOPA}决定保留哪些权重变化，\nt{SERO}维持总体水平，并且按照假说还设定规划的时间范围，\nt{NORA}在探索与决断之间选择，\nt{ACET}在写入与读取之间选择（图~\ref{fig:machine}）。

\begin{figure}[t]
\centering
\begin{tikzpicture}[>=Stealth,
  blk/.style={draw,very thick,rounded corners=3pt,align=center,font=\small},
  reg/.style={draw,very thick,double,double distance=1.2pt,circle,minimum size=1.25cm,inner sep=0pt,font=\scriptsize,fill=orange!15},
  bc/.style={->,thick,densely dashed,orange!85!black}]
% sensors, bus, actions
\node[blk,fill=green!8,text width=1.7cm] (S) at (0.9,0) {传感器\\\scriptsize 开/关};
\draw[very thick,rounded corners=4pt,fill=blue!5] (2.6,-1.35) rectangle (9.0,1.35);
\node[font=\small] at (5.8,0.95) {\nt{GLUT}总线：数据};
\node[font=\scriptsize,align=center] at (4.3,0.25) {符合、延迟、\\逻辑（第\ref{sec:glutcomputer}章）};
\node[font=\scriptsize,align=center] at (7.4,0.25) {记忆、编码\\（第\ref{sec:code}、\ref{sec:acet}章）};
\draw[thick,red!70!black,fill=red!8,rounded corners=2pt] (2.8,-1.15) rectangle (8.8,-0.55);
\node[font=\scriptsize,red!60!black] at (5.8,-0.85) {\nt{GABA}（第\ref{sec:gaba}章）：禁止、选择、除法、节律};
\node[blk,fill=blue!5,text width=2.0cm] (A) at (10.9,0) {选择\\动作};
\draw[->,very thick] (S) -- (2.6,0);
\draw[->,very thick] (9.0,0) -- (A);
\draw[->,very thick] (A) -- (12.6,0) node[right,font=\small] {肌肉};
\pic[scale=1.1] at (13.2,-0.5) {spring};
\pic[scale=1.15] at (0.4,0.85) {eyepic};
\pic[scale=1.05] at (1.35,0.85) {speaker};
% registers
\node[reg] (D) at (3.4,3.2) {\nt{DOPA}};
\node[reg] (C) at (5.6,3.2) {\nt{SERO}};
\node[reg] (N) at (7.8,3.2) {\nt{NORA}};
\node[reg] (H) at (10.0,3.2) {\nt{ACET}};
\node[font=\scriptsize,align=center,above=1pt] at (D.north) {误差 $\delta$};
\node[font=\scriptsize,align=center,above=1pt] at (C.north) {水平，$\tau_\gamma$};
\node[font=\scriptsize,align=center,above=1pt] at (N.north) {增益 $g$};
\node[font=\scriptsize,align=center,above=1pt] at (H.north) {写入 $a$};
\draw[bc] (D.south) -- (3.4,1.35);
\draw[bc] (C.south) -- (5.6,1.35);
\draw[bc] (N.south) -- (7.8,1.35);
\draw[bc] (H.south) -- (10.0,2.1) -- (8.8,1.35);
\draw[bc] (H.south east) to[bend left=20] (A.north east);
\draw[->,thick] (2.85,1.35) -- (2.85,2.75) -- (D.west) node[pos=0.25,left,font=\scriptsize] {$V$};
\draw[->,thick,green!45!black] (0.4,3.2) node[left,black,font=\small] {奖励 $r$} -- (2.2,3.2) -- (D.west);
\pic[scale=0.9] at (-0.45,3.7) {juice};
\draw[->,thick,densely dotted,gray!50!black] ([yshift=0.45cm]D.north) to[bend left=18] node[above,font=\scriptsize,black] {意外} ([yshift=0.45cm]N.north);
\end{tikzpicture}
\caption{整台机器。\nt{GLUT}总线把数据从传感器送到动作选择；\nt{GABA}控制总线内部的数据流。双线圆圈是广播通道；虚线箭头表示把它们的信号发给整个网络，每条通道有自己的旋钮。\nt{DOPA}接收奖励$r$和来自总线内部评价者的预期$V$（第~\ref{sec:dofa}~章）；点线箭头是一个假说：\nt{NORA}通过异常大的误差得知意外（第~\ref{sec:nora}~章）。}
\label{fig:machine}
\end{figure}

\section{一个公式容纳所有旋钮}

所有旋钮可以汇总成一个示意性的公式。这不是新模型，而是第~\ref{sec:glutcomputer}--\ref{sec:acet}~章各个模型的汇总：每个符号都取自相应的一章。
\begin{empheq}[box=\keybox]{equation}
\begin{aligned}
\tau\,\frac{\dd r_i}{\dd t} \;&=\; -r_i + F\Bigl(g\,\Bigl[\tfrac{1+a}{2}\,x_i + (1-a)\sum_j W_{ij}\,r_j - \sum_k M_{ik}\,q_k - \theta\Bigr]\Bigr),\\
\frac{\dd W_{ij}}{\dd t} \;&=\; \eta\,a\,\delta(t)\,e_{ij}(t),
\qquad \tau_e\,\frac{\dd e_{ij}}{\dd t} = -e_{ij} + r_i\,r_j,
\qquad \delta = r + \frac{\dd V}{\dd t} - \frac{V}{\tau_\gamma} .
\end{aligned}
\label{eq:machine}
\end{empheq}
其中$r_i$是\nt{GLUT}神经元的发放频率，$x_i$是来自传感器的输入，$W_{ij}\ge0$是它们之间的权重（第~\ref{sec:glutcomputer}~章）；$q_k$是\nt{GABA}神经元的发放频率，$M_{ik}\ge0$是它们的抑制强度（第~\ref{sec:gaba}~章）；$e_{ij}$是资格迹，$\delta$是\nt{DOPA}误差（第~\ref{sec:dofa}~章）；$\tau_\gamma$是时间范围，按照假说由\nt{SERO}设定；$g$是\nt{NORA}的增益（第~\ref{sec:nora}~章）；$a$是\nt{ACET}的水平（第~\ref{sec:acet}~章）。

请注意式~\eqref{eq:machine}中缺了什么。没有节拍：一切都用微分方程写出，每个神经元自行变化。没有独立的存储器：负责记忆的正是参与计算的那些$W_{ij}$，以及同一个活动$r_i$。绝大多数突触没有各自的专属修正：它们的学习是局部资格迹与一个公共数的乘积；只有在学习动作的回路里，才有通往每个输出的单独误差导线（第~\ref{sec:motorlearning}~章）。控制信号总共只有四类：$\delta$、$\tau_\gamma$、$g$、$a$，却要管理八百亿个神经元。每一类其实不是一个数，而是一小束并行线路（命题~\ref{prop:bundle}），但每束只有几条线。

\begin{keyblock}
\begin{proposition}[架构]
\label{prop:architecture}
就我们目前的理解，大脑这台机器可以用三层来描述。数据沿\nt{GLUT}寻址总线传送。寻址式\nt{GABA}神经元引导、限制和归一化数据流。工作模式由几条广播通道设定，每条都是一小束并行线路，由网络的大片区域共用。前两层已经确立；各广播通道之间的分工则主要是假说。
\end{proposition}
\end{keyblock}

\section{所有类型汇于一表}

表~\ref{tab:synthesis}汇总了本书中所有的神经元类型。没有单独成章的四种类型在表中各占一行；其中两种在讲机器输出的几章里找到了角色：\nt{GLYC}在脊髓回路中（第~\ref{sec:muscles}~章），\nt{ADRE}在化学输出中（第~\ref{sec:chemoutput}~章）。其余两种在计算中的作用要么很简单，要么尚不清楚。不决定神经元类型的物质收在附录~\ref{sec:othermediators}中。

\begin{table}[t]
\centering
\footnotesize
\setlength{\tabcolsep}{3pt}
\renewcommand{\arraystretch}{1.15}
\begin{tabular}{>{\raggedright}p{1.3cm}>{\raggedright}p{1.0cm}>{\raggedright}p{3.9cm}>{\raggedright}p{1.5cm}>{\raggedright}p{1.8cm}>{\raggedright\arraybackslash}p{3.8cm}}
\toprule
类型 & 章 & 计算什么 & 时间 & 总线 & 已知程度 \\
\midrule
\nt{GLUT} & \ref{sec:glutcomputer}, \ref{sec:code}, \ref{sec:sensors}, \ref{sec:motorlearning}, \ref{sec:dendrites} & 数据：符合、延迟、逻辑、记忆、序列 & ms & 寻址 & 已确立 \\
\nt{GABA} & \ref{sec:gaba}, \ref{sec:code}, \ref{sec:pain}, \ref{sec:motorlearning} & 数据流控制：禁止、选择、除法、稳定、节律；疼痛闸门 & ms & 寻址 & 已确立；频率除法需与背景噪声共同作用 \\
\nt{SERO} & \ref{sec:serocomputer}, \ref{sec:pain}, \ref{sec:sleep} & 水平调节器、平滑；时间范围$\tau_\gamma$；疼痛音量；睡眠模式 & 秒 & 容积 & 自抑制已测量；时间范围是假说 \\
\nt{DOPA} & \ref{sec:dofa}, \ref{sec:dofaalt} & 预测误差$\delta$；慢水平是时间的代价 & $0.1$\,s和分钟 & 广播 & $\delta$已测量；扩展见第\ref{sec:dofaalt}章 \\
\nt{NORA} & \ref{sec:nora}, \ref{sec:pain}, \ref{sec:chemoutput}, \ref{sec:sleep} & 增益$g$：探索还是决断；意外；器官的“油门” & 秒 & 广播 & 信噪比提高已测量；在探索中的作用是假说 \\
\nt{ACET} & \ref{sec:acet}, \ref{sec:muscles}, \ref{sec:chemoutput}, \ref{sec:sleep} & 写入还是读取；学习速度；通往肌肉的输出；器官的“刹车” & 秒 & 两者皆有 & 三种效应已测量；模式切换是假说 \\
\nt{GLYC} & \ref{sec:muscles}, \ref{sec:pain} & 脊髓和脑干中的负权重：禁止拮抗肌 & ms & 寻址 & 已确立 \\
\nt{HIST} & --- & 全局“保持清醒”信号 & 缓慢 & 广播 & 在计算中的作用未经研究 \\
\nt{ADRE} & \ref{sec:chemoutput} & 经血液给全身踩“油门”；在脑内未知 & 缓慢 & 广播 & 脑外已确立；在脑内的作用不清楚 \\
\nt{ASPA} & --- & 弱正权重 & ms & 寻址 & 作用有争议 \\
\bottomrule
\end{tabular}
\caption{本书所有神经元类型：各自计算什么，以及对此了解多少。}
\label{tab:synthesis}
\end{table}

\section{旋钮如何协同工作}

到目前为止，每一章只转动一个旋钮。现在让我们跟踪一个事件穿过整台机器（图~\ref{fig:timeline}）。动物早已学会：动作$A$会带来果汁。某一时刻世界变了：$A$不再带来果汁，带来果汁的是动物几乎从不尝试的动作$B$。

\emph{误差。}$A$之后果汁没有来，\nt{DOPA}神经元以一次低谷作答，按误差公式~\eqref{eq:ctd}，$\delta<0$。刚刚选择了$A$的突触带着资格迹，于是被削弱。

\emph{意外。}一次低谷只是寻常的偶然。但低谷反复出现，误差变得比平常更大。按照第~\ref{sec:nora}~章的假说，这正是\nt{NORA}的信号：世界变了。增益$g$下降，选择变得柔和，噪声更常促成对$B$的尝试。

\emph{写入。}按照第~\ref{sec:acet}~章的假说，变化之后\nt{ACET}升高，把网络切换到写入模式：保存旧习惯的内部连接被削弱，来自传感器的输入被加强，权重容易改变。

\emph{新习惯。}尝试$B$带来了谁也没料到的果汁：一次爆发，$\delta>0$，选择了$B$的突触被加强。这样试上几次，预期$V$就追上了新世界，误差重新变小。

\emph{复原。}意外过去，\nt{NORA}恢复高增益，选择重新变得果断；\nt{ACET}回落，网络在读取模式下使用学到的东西。在整个过程中，\nt{SERO}按照假说设定时间范围$\tau_\gamma$：多远之后的果汁仍然值得等待。

\begin{figure}[t]
\centering
\begin{tikzpicture}[>=Stealth,x=1.1cm,y=0.95cm]
\foreach \y/\lab in {4/{$\delta$（\nt{DOPA}）},3/{$g$（\nt{NORA}）},2/{$a$（\nt{ACET}）},1/{选择 $B$}}{
  \draw[gray!60!black] (0,\y-0.35) -- (10,\y-0.35);
  \node[left,font=\scriptsize] at (0,\y) {\lab};}
\draw[densely dashed,gray!60!black] (2,0.4) -- (2,4.55) node[above,font=\scriptsize,black] {世界变了};
\draw[densely dashed,gray!60!black] (5,0.4) -- (5,4.55) node[above,font=\scriptsize,black] {$B$第一次换来果汁};
\pic[scale=0.85] at (2,5.25) {nojuice};
\pic[scale=0.85] at (5,5.25) {juice};
% delta: dips after the change, burst at the first juice for B, then back to zero
\draw[thick,red!70!black] (0,3.65) -- (2.3,3.65) -- (2.3,3.3) -- (2.5,3.3) -- (2.5,3.65) -- (3.2,3.65) -- (3.2,3.3) -- (3.4,3.3) -- (3.4,3.65) -- (4.1,3.65) -- (4.1,3.35) -- (4.3,3.35) -- (4.3,3.65) -- (5.0,3.65) -- (5.0,4.35) -- (5.2,4.35) -- (5.2,3.65) -- (5.8,3.65) -- (5.8,4.1) -- (6.0,4.1) -- (6.0,3.65) -- (6.6,3.65) -- (6.6,3.85) -- (6.8,3.85) -- (6.8,3.65) -- (10,3.65);
% gain: high, drops after surprise, recovers
\draw[thick,blue!60!black] (0,3.2) -- (3.0,3.2) -- (3.4,2.75) -- (5.8,2.75) -- (7.2,3.2) -- (10,3.2);
% ACh: low, rises after the change, falls after learning
\draw[thick,green!45!black] (0,1.75) -- (2.8,1.75) -- (3.3,2.25) -- (6.8,2.25) -- (7.8,1.75) -- (10,1.75);
% choice of B
\draw[thick,black] (0,0.7) -- (3.0,0.7) plot[domain=3:10,samples=40] (\x,{0.7+0.6/(1+exp(-(\x-5.3)*1.8))});
\draw[->,gray!70!black] (0,0.2) -- (10.2,0.2) node[below left,font=\scriptsize,black] {时间};
\end{tikzpicture}
\caption{一个事件穿过整台机器的全过程。这是依据第~\ref{sec:dofa}--\ref{sec:acet}~章的模型和假说画出的定性示意图，而不是计算结果。变化之后\nt{DOPA}以低谷作答；按照假说，反复的低谷抬高意外信号，\nt{NORA}降低增益，\nt{ACET}开启写入；$B$第一次换来果汁时出现爆发，选择转向$B$，各旋钮随后复原。}
\label{fig:timeline}
\end{figure}

请注意，没有哪个旋钮知道其他旋钮在做什么。每个旋钮只对自己的特征作出反应——误差、误差的异常大小、不确定性——动作的次序自行形成，就像异步电路中每个模块在输入就绪时触发一样。据我们所知，这种整体上的协同工作还没有人检验过：各个旋钮分别都测量过，而它们的联合工作是假说。

\section{这台机器哪里不像计算机}

\begin{table}[t]
\centering
\footnotesize
\setlength{\tabcolsep}{4pt}
\renewcommand{\arraystretch}{1.15}
\begin{tabular}{>{\raggedright}p{2.2cm}>{\raggedright}p{4.2cm}>{\raggedright\arraybackslash}p{6.6cm}}
\toprule
& 普通计算机 & 大脑 \\
\midrule
时间 & 统一节拍，约一吉赫兹 & 没有全局时钟；每个神经元自行变化，时间窗由事件和局部节律决定（第\ref{sec:glutcomputer}、\ref{sec:gaba}、\ref{sec:code}章） \\
元件 & 可靠的二进制逻辑门 & 模拟阈值元件；突触丢失大部分脉冲（第\ref{sec:code}章） \\
连接的符号 & 任意 & 由发送端神经元决定（第\ref{sec:neurontypes}章） \\
存储器 & 与处理器分离 & 存在于执行计算的同一批连接中，也存在于自我维持的活动中（第\ref{sec:glutcomputer}、\ref{sec:acet}章） \\
学习 & 误差反向传播 & 局部规则、一个广播数$\delta$（第\ref{sec:dofa}章），以及通往动作回路输出的误差导线（第\ref{sec:motorlearning}章） \\
控制 & 寄存器和指令总线 & 四条广播通道，每条是由几条线路组成的线束（第\ref{sec:serocomputer}、\ref{sec:dofa}--\ref{sec:acet}章） \\
能量 & 单个处理器几十到几百瓦 & 整个大脑约$20$瓦，平均每个神经元每秒约一个脉冲（第\ref{sec:code}章） \\
\bottomrule
\end{tabular}
\caption{大脑与普通计算机。}
\label{tab:computer}
\end{table}

表~\ref{tab:computer}汇总了主要区别。每一条都不是大自然来不及修正的缺陷，而是同一套方案的组成部分。没有统一节拍，就需要“开—关”传感器和符合检测器。元件不可靠，就需要大量神经元提供冗余。存储与计算合一，就需要\nt{ACET}，让写入不破坏读取。学习没有反向导线，就需要\nt{DOPA}和噪声。而每秒一个脉冲的能量上限，迫使整套语言精打细算，只把脉冲花在新消息上。

\section{我们还不知道什么}

结束这份汇总最诚实的方式，是列出我们还不知道的东西。每一条都是留给工程师的题目。

\emph{编码。}我们知道脉冲信道的容量，也知道接收方会选择读取消息的哪一部分（第~\ref{sec:code}~章）。但皮层在多大程度上利用脉冲的精确时间，神经元是否有“词”，目前还不知道。

\emph{跨越多层的学习。}广播信号学得很慢，而且每多一个影响结果的神经元就更慢（命题~\ref{prop:broadcastcost}）。那么皮层如何在多层回路中调整数十亿个权重，目前还不知道；有些模型让簇放电在层间传递误差~\cite{Payeur2021}。答案的一部分也许在于神经元自身分支内部的计算，在那里单个神经元表现得像一个小型多层网络：要复现一个皮层神经元模型的响应，人工网络需要五到八层~\cite{Beniaguev2021}，而人类神经元的分支甚至能计算异或~\cite{Gidon2020}。另一个候选是自监督学习：如果皮层学习预测自己的输入，误差就在每一层就地产生，不需要公共信号（第~\ref{sec:dofa}~章）~\cite{Keller2018}。

\emph{广播通道实际传送什么。}对\nt{DOPA}，有一幅标准图景及其六种扩展（第~\ref{sec:dofaalt}~章）；对\nt{NORA}和\nt{ACET}，有看似合理的假说（第~\ref{sec:nora}、\ref{sec:acet}~章）；对\nt{SERO}，则大多还是猜测。

\emph{时间上的协调。}我们已经看到如何计算一个函数、如何从事件起计时、如何用局部节律把数据流切成时间窗。但一个没有统一节拍的庞大网络如何协调相距遥远的脑区，我们只理解了一部分。

\emph{能量。}总共二十瓦。我们明白编码为什么必须稀疏（命题~\ref{prop:economy}），但目前还没有人能造出一台以同样功率完成同样工作的机器~\cite{MehonicKenyon2022}。

大脑是一台计算机器，它不是工程师组装的，却遵循任何工程师都认得出的规律：RC电路、阈值、反馈、滤波器、总线和广播寄存器。它的电路图我们只读懂了一部分。读完它的，将是会读电路图的人。

\section{本书接下来的内容}

本章组装了机器的计算核心。其余各章走出这个核心。第~\ref{sec:sensors}~章和第~\ref{sec:recognition}~章讲输入：传感器如何把光、声音和分子变成脉冲，机器又如何从中认出物体。第~\ref{sec:muscles}--\ref{sec:chemoutput}~章讲输出以及为输出服务的部分：肌肉、作为警报信号的疼痛、借助单独误差导线的动作学习，以及通往身体的慢速化学输出。第~\ref{sec:debugging}~章讲如何从外部调试这台机器，包括用脑机接口。第~\ref{sec:eetypes}~章和第~\ref{sec:dendrites}~章再次给神经元分类，这次依据的是电学功能和输入数目。第~\ref{sec:sleep}~章讲机器与外界断开时做什么，第~\ref{sec:livingmachine}~章和第~\ref{sec:dishbrain}~章讲用芯片上的活神经元搭成的机器。

\section{习题}

本章习题把不同章节的结果联系起来：每道题至少用到其中两章。

\begin{exercise}[\lvlA]
\label{ex:10:1}
\emph{总线与寄存器。}四条广播通道各是约五条线路的线束（命题~\ref{prop:bundle}）；每条线路每秒变化一次，可分辨四个电平。\nt{GLUT}总线（$8.6\cdot10^{10}$个神经元，频率按式~\eqref{eq:energy}，精度按式~\eqref{eq:spikeentropy}取$1\ms$）一秒传送的信息，控制系统要传多少年？
\end{exercise}

\begin{exercise}[\lvlB]
\label{ex:10:2}
\emph{同一回路上的两个旋钮。}在回路~\eqref{eq:ei}中，式~\eqref{eq:machine}的旋钮作用于兴奋：$\tau_E\dot E=-E+g[(1-a)w_{EE}E-w_{EI}I+h]$，\nt{GABA}不受其控制。取图~\ref{fig:ei}的数值，\nt{NORA}的簇放电把$g$升到$6$。$a$最小为多少时回路稳定？\nt{NORA}和\nt{ACET}的旋钮能否独立调节？
\end{exercise}

\begin{exercise}[\lvlB]
\label{ex:10:3}
\emph{广播的代价从何而来。}$N$个神经元的输出$y_k=f(u_k)+\xi_k$带有方差为$\sigma^2$的独立高斯噪声；$R\approx R_0+\sum_kR'_k\xi_k$，各$R'_k$相同，\nt{DOPA}广播$\delta=R-R_0$。求单次尝试中修正量$\delta\,\xi_jx_j$的信噪比。所需尝试次数如何随$N$增长？
\end{exercise}

\begin{exercise}[\lvlB]
\label{ex:10:4}
\emph{把声音和闪光绑定。}声音经$5\ms$到达总线，闪光经$30\ms$再加上首个脉冲的延迟~\eqref{eq:latency}（$\tau=10\ms$，$U$从$1.2\theta$到$4\theta$）；每个输入的背景是每秒$5$个脉冲。选定声音线路的延迟和适用于所有对比度的最小符合窗。背景每秒造成多少次错误绑定？
\end{exercise}

\begin{exercise}[\lvlB]
\label{ex:10:5}
\emph{建模：谁察觉到变化。}评价者看到$y=s+\text{噪声}$（$R=1$）；$s$以$1/200$的概率跳变到方差为$J^2=4$的新值。模拟$q=0$的卡尔曼滤波器，当$E\leftarrow0.9E+\delta^2/(P+R)$超过$20$时把$P$升到$J^2$。它的误差比最佳常数$\alpha$时小多少？
\end{exercise}

\begin{exercise}[\lvlA]
\label{ex:10:6}
\emph{改掉一个习惯要尝试几次。}网络学到$Q_A=0.8$、$Q_B=0.2$，按式~\eqref{eq:softmax}选择，$\sigma=1$，$g=8$。现在$A$以$0.2$的概率给出果汁；$Q_A\leftarrow Q_A+\alpha(r-Q_A)$，$Q_B$不变。$\alpha=0.1$和$0.5$时，选择$A$多少次后，选$A$的概率降到$0.6$？若\nt{NORA}把$g$降到$2$呢？
\end{exercise}

\begin{exercise}[\lvlC]
\label{ex:10:7}
\emph{用线束代替导线。}输出$y_k=w_k+\xi_k$，奖励$R=-\sum_ky_k^2$；线束中的一条线路向一组$G$个神经元广播该组输出的误差，$w_k\leftarrow w_k+\eta\,\delta_l\xi_k$。证明在最佳$\eta$下，要达到$\sum w_k^2=\varepsilon\sum w_k^2(0)$需要$\approx(G+2)\ln(1/\varepsilon)$次尝试。$\varepsilon=0.01$时，由$10^6$个神经元组成、只有五条线路、每三秒尝试一次的回路要学多久？
\end{exercise}

\chapter{机器的输入：眼睛、耳朵和其他传感器如何接入}
\chaptermark{机器的输入}
\label{sec:sensors}
\begin{chaptergoals}
\item 每种感觉器官如何作为换能器工作：受体闸门、增益调节、脉冲编码器；
\item 为什么光子散粒噪声限制暗处的视觉，为什么黄昏时我们看得更慢；
\item 自动增益控制如何把巨大的输入范围装进脉冲，使传感器报告变化；
\item 眼睛如何把图像压缩约一百倍，耳朵如何作为滤波器组工作。
\end{chaptergoals}

\section{传感器即换能器}

在图~\ref{fig:machine}中，数据从左侧的“传感器”模块进入机器。现在打开这个模块。对工程师来说，每个感觉器官都是一个\emph{换能器}：前端是模拟输入级，末端是模数转换器，只不过输出的数字不是数值，而是脉冲。

所有感觉器官的方案都相同（图~\ref{fig:transducer}）。物理量——光、压力、振动、分子——作用于感觉细胞膜上的受体蛋白。这个蛋白像一道闸门：它直接或经过一小串反应打开离子通道。于是产生\emph{感受器电流}，随之在膜上形成模拟电压。接着电压经过增益调节，进入脉冲编码器——也就是我们熟悉的积分神经元~\eqref{eq:glutneuron}，它把电平变成脉冲的频率和时刻。输出几乎总是一样的：眼睛、耳朵、嗅觉和皮肤的感觉神经元都释放谷氨酸，因此各路输入直接接到\nt{GLUT}总线上。

\begin{figure}[t]
\centering
\begin{tikzpicture}[>=Stealth,
  blk/.style={draw,very thick,rounded corners=2pt,align=center,font=\footnotesize,minimum height=1.0cm,text width=1.9cm,fill=blue!5}]
\node[align=center,font=\small] (P) at (0.1,0) {光、声、\\压力、\\分子};
\node[blk] (R) at (3.0,0) {受体\\闸门};
\node[blk] (G) at (6.5,0) {增益\\调节};
\node[blk] (E) at (10.0,0) {脉冲\\编码器};
\node[align=center,font=\small] (B) at (13.4,0) {\nt{GLUT}\\总线};
\draw[->,very thick] (P) -- (R);
\foreach \ic/\xx in {sun/-1.1,speaker/-0.3,press/0.5,molecule/1.3}{\pic[scale=0.85] at (\xx,1.05) {\ic};}
\draw[->,very thick] (R) -- node[above,font=\scriptsize] {电流} (G);
\draw[->,very thick] (G) -- node[above,font=\scriptsize] {电平} (E);
\draw[->,very thick,red!70!black] (E) -- node[above,font=\scriptsize,black] {脉冲} (B);
\draw[decorate,decoration={brace,amplitude=5pt,mirror},gray!60!black] (1.8,-0.8) -- (7.7,-0.8) node[midway,below=5pt,font=\scriptsize,black] {模拟部分};
\draw[decorate,decoration={brace,amplitude=5pt,mirror},gray!60!black] (8.8,-0.8) -- (11.2,-0.8) node[midway,below=5pt,font=\scriptsize,black] {脉冲};
\end{tikzpicture}
\caption{任何感觉器官都是一条变换链。受体蛋白打开通道并产生电流；增益调节使电流适应环境；编码器~\eqref{eq:glutneuron}把电平变成脉冲，送上\nt{GLUT}总线。在眼睛和耳朵里，最前面的几级根本不发放脉冲：信号一直保持模拟形式，直到需要远距离传送时为止。}
\label{fig:transducer}
\end{figure}

有一个细节电子工程师很熟悉。在眼睛和耳朵里，最前端的细胞根本不发放脉冲：它们向邻近细胞传递平滑变化的电压，就像电路板上的模拟级。只有在需要把信号送往远处的地方才出现脉冲。模拟信号传过几毫米就会衰减并混入噪声，而脉冲，正如我们在第~\ref{sec:neurontypes}~章中所见，到达导线末端时高度不变。数字化发生在长线开始的地方。

\section{在物理极限上}

大脑的传感器工作在物理灵敏度的极限附近。暗处的光传感器能响应\emph{单个光子}：吸收一个光粒子会产生约一皮安的电流，在自身噪声的背景上清晰可辨~\cite{Baylor1979}。声音传感器能察觉其敏感纤毛束小于一纳米的位移~\cite{Hudspeth1989}。

光线很弱时，限制精度的不是传感器，而是光本身：光子随机到达，构成泊松流。若在积累时间$T$内平均有$N=\Phi T$个光子落入传感器，其数目的涨落为$\sqrt N$。增量$\Delta N$要能被察觉，就必须比这个涨落大几倍，因此可分辨的亮度相对变化为
\begin{empheq}[box=\keybox]{equation}
\frac{\Delta I}{I} \;\approx\; \frac{k}{\sqrt{\Phi\,T}} .
\label{eq:photonnoise}
\end{empheq}
这与脉冲计数的公式~\eqref{eq:rateerror}相同：光子的散粒噪声与脉冲的散粒噪声服从同一规律。在暗处，眼睛提高精度只有一个办法——更长时间地积累光子，而积累时间确实会增加到零点几秒。所以在黄昏时我们看得更慢。

\section{动态范围：增益调节}

传感器最难的工程问题是范围。照度从星夜到阳光下的雪地大约变化十个数量级，声强从听阈到痛阈变化十二个数量级。而脉冲频率只在每秒零到几百次之间变化，不到三个数量级。这样的输入无法线性地装进这样的输出。

解决办法与任何接收机相同：\emph{自动增益控制}。眼睛感觉细胞的响应可以很好地用下式描述：
\begin{empheq}[box=\keybox]{equation}
r \;=\; r_{\max}\,\frac{I}{I+\sigma} ,
\label{eq:nakarushton}
\end{empheq}
其中$\sigma$是响应达到一半时的亮度~\cite{NakaRushton1966}。式~\eqref{eq:nakarushton}本身只覆盖两三个数量级。但$\sigma$不是常数：在明亮背景下长时间工作时，它随平均亮度一起增大。响应曲线沿亮度的对数轴移动，每次都把自己的陡峭段放到当前输入所在的位置（图~\ref{fig:adaptation}）。这与第~\ref{sec:gaba}~章的分流抑制一样，是除以总活动~\cite{Carandini2012}，只不过这里的分母是最近几秒的平均亮度。

\begin{figure}[t]
\centering
\begin{tikzpicture}[>=Stealth,x=1.25cm,y=2.8cm]
\draw[->,gray!70!black] (-2,0) -- (6.4,0) node[below left,font=\small,black] {$\log_{10} I$};
\draw[->,gray!70!black] (-2,0) -- (-2,1.15) node[left,font=\small,black] {$r/r_{\max}$};
\foreach \u in {-2,0,2,4,6}{\draw[gray!60!black] (\u,-0.02) -- (\u,0.02); \node[below,font=\scriptsize] at (\u,0) {$\u$};}
\draw[densely dashed,gray!60!black] (-2,0.5) -- (6.2,0.5);
\foreach \s/\c/\lab in {0/blue!60!black/{暗背景},2/green!45!black/{中等},4/red!70!black/{明亮}}{
  \pgfmathsetmacro{\lo}{max(-2,\s-3)}
  \draw[very thick,\c] (-2,0) -- (\lo,0) plot[domain=\lo:6.2,samples=80] (\x,{1/(1+10^(\s-\x))});
  \node[font=\scriptsize,\c,anchor=west] at (\s+0.15,0.42) {\lab};}
\foreach \s/\ic in {0/moon,2/cloud,4/sun}{\pic at (\s+0.6,0.64) {\ic};}
\end{tikzpicture}
\caption{按式~\eqref{eq:nakarushton}进行的光传感器增益调节。每条曲线覆盖两三个数量级的亮度；换到更亮的背景时$\sigma$增大，曲线沿对数轴移动。这些移动的曲线合起来覆盖整个范围。}
\label{fig:adaptation}
\end{figure}

增益调节有代价，这在第~\ref{sec:code}~章已经见过：传感器报告的不是亮度，而是\emph{相对于背景}的亮度。恒定的输入逐渐不再引起响应，而每一次变化都会引起响应。机器的输入从一开始就报告变化而不是电平，这与命题~\ref{prop:economy}一样是在节省脉冲~\cite{Barlow1961}。第~\ref{sec:glutcomputer}~章的接通传感器和断开传感器也由此而来~\cite{Schiller1992}：一类报告变亮了，另一类报告变暗了。

工程师在\emph{事件相机}中复现了这一手法。这种相机不按时钟拍摄帧：每个像素各自独立，不需要统一节拍，当亮度的对数变化了给定比例时，就输出一个“变亮”或“变暗”事件~\cite{Lichtsteiner2008}。这正是第~\ref{sec:glutcomputer}~章的开—关双线编码在硅片上的实现，它让相机获得了普通传感器阵列无法企及的动态范围和速度~\cite{Gallego2022}。

\section{眼睛：压缩到一根线缆}

眼睛里约有$10^8$个光传感器~\cite{Curcio1990}，而把图像送往大脑的输出神经元只有约一百万个。信号离开眼睛之前被压缩了大约一百倍。主要的压缩手法我们都熟悉。每个输出神经元响应的是一个小光斑内的亮度与其周围亮度之差——这是通过抑制减去邻居，与第~\ref{sec:gaba}~章相同。图像中均匀的区域几乎不占带宽，边缘和变化则被传送出去。不同的输出神经元各有专长：一些报告变亮，另一些报告变暗，还有一些报告朝某个方向的运动（图~\ref{fig:direction}）；在小鼠身上已经数出三十多种这样的输出通道~\cite{Baden2016}。

这样得到的线缆带宽有多大？在豚鼠身上，根据输出神经元的记录，测得$\sim10^5$个这类神经元的信息率约为每秒$875$千比特；换算到人的一百万个输出神经元，约为每秒$10^7$比特~\cite{Koch2006}——相当于一根老式的十兆网线。每个神经元的平均频率为每秒几个脉冲，于是每个脉冲携带几比特，与第~\ref{sec:code}~章得到的结果一致。

\section{耳朵：滤波器组}

耳朵要解决另一个问题：把声音按频率分解。工程师会装一组带通滤波器，而耳朵正是用机械方式做到了这一点。声振动沿一块弹性隔膜传播，隔膜的性质沿长度平滑变化，每一段只对自己的频率共振。位于该段的传感器只响应自己的频带（图~\ref{fig:filterbank}）。频率沿位置大致按对数排列：每个倍频程分到的传感器数量相近。被触发的传感器的序号就是频率，这是第~\ref{sec:code}~章的位置编码。

\begin{figure}[t]
\centering
\begin{tikzpicture}[>=Stealth,x=1.35cm,y=2.2cm]
\draw[->,gray!70!black] (-0.3,0) -- (7.6,0) node[above left,font=\small,black] {频率，Hz};
\draw[->,gray!70!black] (-0.3,0) -- (-0.3,1.15) node[left,font=\small,black] {响应};
\foreach \k/\f in {0/125,1/250,2/500,3/1000,4/2000,5/4000,6/8000}{
  \draw[gray!60!black] (\k,-0.02) -- (\k,0.02); \node[below,font=\scriptsize] at (\k,0) {\f};
  \draw[thick,blue!60!black] plot[domain={max(\k-1.2,-0.3)}:{\k+0.7},samples=60] (\x,{1/(1+((\x-\k)/(\x<\k?0.45:0.2))^4)});}
% a piano keyboard: seven white keys per octave, like the octaves on the axis
\foreach \i in {0,...,48}{\draw[gray!50!black,fill=white,line width=0.3pt] ({\i/7},-0.44) rectangle ({(\i+1)/7},-0.26);}
\foreach \o in {0,...,6}{\foreach \j in {0,1,3,4,5}{\fill[black] ({\o+(\j+1)/7-0.035},-0.35) rectangle ({\o+(\j+1)/7+0.035},-0.26);}}
\end{tikzpicture}
\caption{耳朵作为带通滤波器组的示意图。每条曲线是某一段上的传感器对不同频率声音的响应；中心频率相隔一个倍频程，就像对数刻度上的琴键。真实滤波器的高频侧陡峭，低频侧平缓。下方是键盘：与耳朵一样，每个倍频程在上面占据相同的位置。}
\label{fig:filterbank}
\end{figure}

耳朵的谐振器不是被动的。一部分传感器会随着声音的节拍主动推动隔膜振动，把微弱振动的振幅放大数千倍（$66$--$76$\,dB），而随着响度增大，增益下降~\cite{Ruggero1997,Robles2001}。在工程师看来，这是一个工作在自激边缘的正反馈放大器——与第~\ref{sec:gaba}~章的网络一样，生活在临界边缘：系统在边界附近对小信号特别敏感。响度压缩也由同一个放大器完成：轻声放大得多，响声放大得少。

耳朵还有第二种编码。在低频下，神经元的脉冲与声音同相：神经元在每个波的特定位置发放，尽管并非每个波都发放。在豚鼠身上，相位锁定在约$600$\,Hz开始减弱，直到约$3.5$\,kHz仍可察觉~\cite{PalmerRussell1986}。这样，脉冲时刻保存了声音的精确时间，进而保存了声音到达两耳的时间差，第~\ref{sec:glutcomputer}~章的网络正是据此找出方向~\cite{Grothe2010}。在高频下，脉冲跟不上声波，只剩下位置编码。一只耳朵同时使用第~\ref{sec:code}~章的两种编码：神经元跟得上的地方用时间，跟不上的地方用位置。

\section{其他传感器}

\begin{table}[t]
\centering
\footnotesize
\setlength{\tabcolsep}{3pt}
\renewcommand{\arraystretch}{1.15}
\begin{tabular}{>{\raggedright}p{1.9cm}>{\raggedright}p{2.6cm}>{\raggedright}p{4.0cm}>{\raggedright\arraybackslash}p{4.6cm}}
\toprule
感觉 & 测量什么 & 输入如何构成 & 本书中的手法 \\
\midrule
视觉 & 光 & $\sim10^8$个传感器，压缩$\sim100$倍，$\sim10^7$\,bit/s & 增益调节、减去邻居、双线、运动方向 \\
听觉 & 气压 & 机械滤波器组、主动放大器 & 位置编码和时间编码、延迟 \\
触觉 & 压力、振动 & 快适应和慢适应传感器 & “高通滤波器—低通滤波器”对 \\
温度觉 & 热 & 冷、热各有传感器 & 双线编码 \\
嗅觉 & 分子 & $\sim400$种受体，每个传感器只有一种 & 组合编码：气味是被触发类型的集合 \\
味觉 & 食物中的物质 & 五种味质：甜、苦、咸、酸、鲜 & 标记线路；部分细胞释放ATP或血清素，而不是谷氨酸 \\
本体感觉 & 肌肉牵拉 & 长度和牵拉速度传感器 & 为运动调节器提供反馈 \\
\bottomrule
\end{tabular}
\caption{传感器如何接入。第三列中的结构都已测量；右列把它们与前几章的手法联系起来。}
\label{tab:sensors}
\end{table}

其余感觉汇总在表~\ref{tab:sensors}中，几乎每一行都能认出前几章的某个手法。触觉使用一对滤波器：一些压力传感器在接触时响应并立即沉默，另一些只要压力存在就保持响应。温度由两根导线传送，冷和热各一根。最不寻常的输入是嗅觉。人类的嗅觉受体约有四百种，每个嗅觉神经元只带一种~\cite{Mombaerts2004}。一种气味激活的不是一种受体，而是一组，消息就是\emph{哪些}类型被触发了。对工程师来说，这是一个长约四百的二进制向量，可能取值的数目是天文数字。

\begin{keyblock}
\begin{proposition}[机器的输入]
\label{prop:sensors}
每种感觉都是一个换能器：它工作在物理灵敏度的极限附近，用自动增益控制压缩巨大的动态范围，并通过\nt{GLUT}总线首先传送\emph{变化}。为此，传感器使用与机器本身相同的手法：双线编码、位置编码、时间编码和除以背景。传感器的结构已经测量；它们的编码被调节到使每个脉冲携带最多信息，这是一个有充分依据的假说。
\end{proposition}
\end{keyblock}

输入和其他一切一样，是异步工作的。每个传感器在有东西要报告时才发出脉冲，而不同感觉到达的延迟各不相同：声音在几毫秒后到达，光在几十毫秒后到达。没有统一节拍，机器如何把它们合成一幅世界图景——这是第~\ref{sec:synthesis}~章中时间协调问题的一部分。

\section{习题}

\begin{exercise}[\lvlA]
\label{ex:11:1}
\emph{光子要积累多久。}黄昏时每秒有$100$个光子落入传感器；增量比涨落~\eqref{eq:photonnoise}大$3$倍时才可察觉。(a)~单个传感器察觉$10\,\%$的亮度变化需要多长时间？(b)~要在$0.2\,\text{s}$内完成，需要把多少个传感器相加？沿一个方向的分辨率下降多少倍？
\end{exercise}

\begin{exercise}[\lvlA]
\label{ex:11:2}
\emph{一个脉冲有多少比特。}(a)~眼睛通过$10^6$个输出神经元以每秒$3$--$5$个脉冲的频率传送$\sim10^7$\,bit/s。在$\Delta t=1\ms$时，它用了容量~\eqref{eq:spikeentropy}的多大比例？(b)~一种气味激活$\approx400$种受体中的$20$种。这样的集合携带多少比特？两种随机气味平均有几种共同类型？
\end{exercise}

\begin{exercise}[\lvlB]
\label{ex:11:3}
\emph{一条曲线~\eqref{eq:nakarushton}能做什么。}(a)~在什么亮度范围内响应从$r_{\max}$的$10\,\%$增加到$90\,\%$？这是几个数量级？(b)~要覆盖十个数量级的亮度，移动的曲线需要多少个位置？十二个数量级的响度呢？
\end{exercise}

\begin{exercise}[\lvlB]
\label{ex:11:4}
\emph{耳朵中时间编码的极限。}声音到达两耳的时间差最大为$0.22\,\text{m}/343\,\text{m/s}$。(a)~脉冲与纯音同相：在多高的频率以下，据此确定的方向是唯一的？(b)~脉冲抖动$0.5\ms$，需要分辨$20$\,\textmu s。在$50\ms$内需要平均多少个每秒$100$个脉冲的神经元？
\end{exercise}

\begin{exercise}[\lvlB]
\label{ex:11:5}
\emph{接到眼睛线缆上的事件相机。}一台$10^6$像素、输出$10^7$\,bit/s的相机，在像素的$\ln I$变化$\ln1.2$时发送一个事件（地址和符号）。长$500$像素、亮度跳变为$4$倍的边缘最快能以多大速度移动？每像素$8$\,bit的普通相机通过同样的输出每秒能传几帧？
\end{exercise}

\begin{exercise}[\lvlB]
\label{ex:11:6}
\emph{建模：增益调节。}传感器~\eqref{eq:nakarushton}以$\tau=1\,\text{s}$调节$\sigma$，或按对数，$\tau\,\dd\ln\sigma/\dd t=\ln I-\ln\sigma$，或按线性，$\tau\,\dd\sigma/\dd t=I-\sigma$；背景按$1\to100\to1$跳变。在每种规律下，向上和向下跳变后，对$+10\,\%$探测的响应要过几秒才恢复到适应后水平的一半？
\end{exercise}

\begin{exercise}[\lvlC]
\label{ex:11:7}
\emph{曲线适配于怎样的世界。}当输出噪声小而恒定时，若$r(I)=r_{\max}\Phi_I(I)$，关于亮度的信息最多，其中$\Phi_I$是亮度的分布函数。(a)~曲线~\eqref{eq:nakarushton}对哪种分布最优？其$\log_{10}I$的离散程度多大？(b)~输出为泊松噪声时，哪条曲线最优？
\end{exercise}

\chapter{图像和视频识别}
\chaptermark{识别}
\label{sec:recognition}
\begin{chaptergoals}
\item 为什么识别必须忽略平移、缩放和光照，为什么模板匹配行不通；
\item 对部件取“与”、对位置取“或”交替进行的流水线如何在$150\ms$内识别；
\item 带资格迹的赫布规则如何从无标注的视频中学会不变性；
\item 大脑与卷积网络有何不同，错觉揭示了它的什么。
\end{chaptergoals}

\section{问题：像素不同，东西相同}

第~\ref{sec:sensors}~章把图像送到了机器的输入端：约一百万个眼睛输出神经元，一路压缩过的数据流，其中主要传送边缘和变化。但在传感器和动作之间还有一步，我们至今没有谈到。图中的猫不是一组像素。把它平移半个画面、旋转、推远、关掉一半的灯——几乎所有像素都会变，而“猫”这个答案应当不变。工程师称之为\emph{不变性}：分类器的答案不应随平移、缩放、旋转和光照而改变。

一个简单的实验就能看出困难所在。取一维的“视网膜”，共$24$个点，再取两个各由五个点组成的图形，它们可以出现在任何位置。把输入与存储在某一位置的样板相比较的分类器，正确率为$49$--$52\%$——和抛硬币一样。平移彻底破坏了逐像素比较。需要另一种方案。

\section{多级流水线}

大脑做这件事有多快，我们已经从第~\ref{sec:code}~章知道：闪现的图片上的动物约在$150\ms$内被认出，信号经过大约十级，每级$10$--$20\ms$~\cite{Thorpe1996}。在每一级，神经元只来得及发出零个或一个脉冲。可见，快速识别是沿流水线走一遍，没有长时间的思考和反复。问题是每一级做什么。

视觉最初几级的测量提示了答案~\cite{HubelWiesel1962}：它由两种交替进行的操作组成（图~\ref{fig:conveyor}）。

\emph{选择性。}$S$级的神经元观察输入的一小块区域，只有当那里出现特定的部件组合时才发放：这里有这条边，旁边有那条边。这是对部件的“与”——第~\ref{sec:glutcomputer}~章的阈值神经元，其阈值高于任何单个部件所能提供的输入。

\emph{不变性。}$C$级的神经元汇集许多$S$神经元的输出，这些$S$神经元在不同位置寻找同一组合；只要在任何地方找到它，$C$神经元就发放。这是对位置的“或”——更确切地说，是取最大值。

对一维模型，这两种操作写作：
\begin{empheq}[box=\keybox]{equation}
s_{f,p} \;=\; \Bigl[\sum_i t_{f,i}\,x_{p+i} - \theta\Bigr]_+ ,
\qquad
c_f \;=\; \max_p\, s_{f,p} ,
\label{eq:scstage}
\end{empheq}
其中$x$是输入，$t_{f,i}$是特征$f$的样板，$p$是位置，$\theta$是阈值。第一个式子使响应具有选择性，第二个式子使响应与位置无关。网络用第~\ref{sec:gaba}~章的手段求出许多输入中的最大值：相互抑制只留下最强的输入（命题~\ref{prop:wta}），而除以总活动使不同亮度下的响应保持一致~\cite{Carandini2012}。

\begin{figure}[t]
\centering
\begin{tikzpicture}[>=Stealth,
  px/.style={draw,minimum size=0.36cm,inner sep=0pt},
  on/.style={px,fill=blue!55!black},
  snode/.style={circle,draw,very thick,minimum size=0.62cm,inner sep=0pt,font=\scriptsize,fill=blue!6},
  cnode/.style={circle,draw,very thick,minimum size=0.8cm,inner sep=0pt,font=\scriptsize,fill=red!10}]
% two inputs: the same shape at two positions
\foreach \row/\sh/\lab in {0/1/{第1帧},1/7/{第2帧}}{
  \begin{scope}[yshift=-\row*1.0cm]
  \node[left,font=\scriptsize] at (-0.25,0) {\lab};
  \foreach \i in {0,...,13}{\node[px] at (\i*0.4,0) {};}
  \foreach \d in {0,1,3,4}{\node[on] at ({(\sh+\d)*0.4},0) {};}
  \end{scope}}
% S stage: detectors of the shape at several positions
\foreach \k in {0,...,5}{
  \node[snode] (S\k) at ({0.8+0.8*\k},1.7) {$S$};}
\node[font=\scriptsize,left] at (-0.25,1.7) {$S$级：对部件取“与”};
\draw[->,thick,blue!60!black] (1.2,0.25) -- (S1.south);
\draw[->,thick,orange!85!black] (3.6,0.25) -- (S4.south);
\node[font=\scriptsize,blue!60!black,anchor=east] at (1.25,0.85) {第1帧};
\node[font=\scriptsize,orange!85!black,anchor=west] at (3.9,0.85) {第2帧};
% C stage
\node[cnode] (C) at (2.8,3.25) {$C$};
\node[font=\scriptsize,left] at (-0.25,3.25) {$C$级：对位置取“或”};
\foreach \k in {0,...,5}{\draw[->,gray!60!black] (S\k.north) -- (C);}
\draw[->,very thick,red!70!black] (C.north) -- ++(0,0.6) node[above,font=\scriptsize,black] {“有这个图形”——两帧都是};
% next stages
\draw[decorate,decoration={brace,amplitude=5pt},gray!60!black] (6.3,1.4) -- (6.3,-1.3);
\node[right,font=\scriptsize,align=left] at (6.5,0.05) {同一个物体，\\不同的像素};
\draw[->,thick,densely dashed,gray!60!black] (3.6,3.25) -- (5.9,3.25) node[right,font=\scriptsize,black,align=left] {后续的$S$、$C$对——\\更大、更复杂、\\更不变};
\end{tikzpicture}
\caption{流水线中的一对级。$S$神经元在不同位置寻找同一个部件组合；只要在任何地方找到它，$C$神经元就作出响应~\eqref{eq:scstage}。第1帧和第2帧中的图形激活不同的$S$神经元，却激活同一个$C$神经元。后续的级对在越来越大、越来越复杂的组合上重复同样的操作。}
\label{fig:conveyor}
\end{figure}

在同一个一维模型中，一对级~\eqref{eq:scstage}能在任何位置无误地认出图形；如果输入的每个点以$0.1$的概率随机翻转，正确率为$86\%$，概率为$0.2$时为$70\%$。硬币仍然只有一半。把几对这样的级叠起来：第一对寻找边缘，下一对寻找边缘的组合，再往上是物体的部件，顶部是完整的物体。每上一对，组合更大，而响应越来越不依赖物体在哪里、以什么方式摆放~\cite{Riesenhuber1999}。

\begin{keyblock}
\begin{proposition}[选择性与不变性流水线]
\label{prop:conveyor}
快速识别是沿大约十级走一遍。在每一对级中，对部件的“与”~\eqref{eq:scstage}使响应具有选择性，对位置的“或”使响应与平移无关。这两种操作交替进行，得到的响应既能区分物体，又几乎不依赖于物体在哪里、以什么方式被看到。最初几级的结构和通过速度已经测量；整条链可归结为这两种操作，则是一种简化。
\end{proposition}
\end{keyblock}

如果这种结构看起来眼熟，那是理所当然的。正是这种交替——在所有位置寻找特征，再在相邻位置间合并——被工程师移植到了人工\emph{卷积}网络中~\cite{Fukushima1980}。最近又反过来：为识别物体而训练的卷积网络，成了已知最好的视觉高层神经元响应模型~\cite{Yamins2014}。顶层的结果读起来很简单：根据最后一级约一百个神经元的总频率，线性判决器在图片呈现十分之一秒后就能认出物体~\cite{Hung2005}。这就是第~\ref{sec:code}~章的群体频率编码。

\section{不存在的老师：从视频中学习}

卷积网络要用数以百万计带标注的图片来训练。几乎没有人给大脑标注：孩子一连几小时看着各种物体，却只偶尔听到“杯子”这个词。那么$C$级怎么知道该合并哪些$S$神经元？要把“这里的这个图形”和“那里的这个图形”合并，就得先知道它们是同一个图形。

提示来自视频本身。世界是连续变化的：视野中的物体在平移、旋转、光照改变时仍是同一个物体，只是偶尔视线才转向另一个物体。凡是变化\emph{缓慢}的，多半属于物体本身；变化快的，则属于它的位置和外观~\cite{WiskottSejnowski2002}。因此，$C$神经元只需第~\ref{sec:dofa}~章带资格迹的赫布规则：把先后出现的东西联系起来，而不只是同时出现的东西~\cite{Foldiak1991}：
\begin{empheq}[box=\keybox]{equation}
\tau_{\text{tr}}\,\frac{\dd\bar y_k}{\dd t} \;=\; -\bar y_k + y_k ,
\qquad
\frac{\dd\mathbf w_k}{\dd t} \;=\; \eta\,\bar y_k\,\bigl(\mathbf x - \mathbf w_k\bigr) .
\label{eq:traceinvariance}
\end{empheq}
其中$y_k$是通过抑制竞争获胜的$C$神经元$k$的活动，$\bar y_k$是它的资格迹，与规则~\eqref{eq:threefactor}中的RC电路相同，$\mathbf x$是$S$神经元的输出。在物体的第一个视图上获胜的神经元，在第二个视图到来时仍然记得这一点，于是也把第二个视图拉向自己。

\begin{figure}[t]
\centering
\begin{tikzpicture}[>=Stealth,x=1.0cm,y=1.0cm]
\foreach \i/\lab in {0/1,1/2,2/3,3/4}{
  \node[draw,thick,fill=blue!8,minimum width=0.9cm,minimum height=0.55cm,font=\scriptsize] at (\i+0.5,2.2) {$A$在\lab};}
\foreach \i/\lab in {0/4,1/3,2/2,3/1}{
  \node[draw,thick,fill=orange!12,minimum width=0.9cm,minimum height=0.55cm,font=\scriptsize] at (\i+6.5,2.2) {$B$在\lab};}
\node[font=\scriptsize,gray!40!black] at (5.0,2.2) {停顿};
\draw[->,gray!70!black] (0,0) -- (10.4,0) node[below left,font=\scriptsize,black] {时间};
% traces
\draw[thick,blue!60!black] (0,0.05) .. controls (0.4,1.2) and (0.8,1.3) .. (1.2,1.35) -- (4,1.4) .. controls (4.6,0.5) and (5.2,0.15) .. (6,0.08) -- (10,0.05);
\draw[thick,orange!85!black] (0,0.05) -- (6,0.05) .. controls (6.4,1.2) and (6.8,1.3) .. (7.2,1.35) -- (10,1.4);
\node[font=\scriptsize,blue!60!black,anchor=west] at (1.4,1.65) {资格迹$\bar y_1$——在$A$的整个经过期间保持};
\node[font=\scriptsize,orange!85!black,anchor=west] at (7.2,1.65) {资格迹$\bar y_2$};
\draw[decorate,decoration={brace,amplitude=4pt},gray!60!black] (4.0,2.6) -- (0,2.6);
\draw[decorate,decoration={brace,amplitude=4pt,mirror},gray!60!black] (0,-0.35) -- (4,-0.35) node[midway,below=4pt,font=\scriptsize,black] {$A$的所有视图都与神经元1相连};
\end{tikzpicture}
\caption{视频如何教会不变性（示意图）。物体$A$经过四个位置，然后停顿，然后是物体$B$。在$A$的第一个视图上获胜的神经元，其资格迹~\eqref{eq:traceinvariance}在整个经过期间保持，于是$A$的所有视图都与同一个神经元相连。停顿抹去资格迹，$B$归另一个神经元。}
\label{fig:tracevideo}
\end{figure}

我们在模型上检验了这一点（图~\ref{fig:tracevideo}）。两个$C$神经元争夺来自$16$个$S$神经元的输入：两个图形，各有八个位置。图形依次经过所有位置，每个位置$50\ms$，然后停顿$0.3\,\text{s}$，再换下一个随机图形。没有任何标注。如果资格迹很短，$\tau_{\text{tr}}=5\ms$，神经元就按\emph{位置}瓜分输入，响应与图形的吻合程度不比硬币好：十次运行平均$52\%$。如果资格迹与一个位置的呈现时间相当，从$\tau_{\text{tr}}\approx20\ms$起，每个神经元都收集到自己图形的\emph{所有}位置（在$20$、$50$和$200\ms$时，十次运行全部如此）：仅凭视频就学会了不变性。物体之间的停顿很重要：没有停顿，资格迹就会延续到下一个物体上，把它们混在一起。

实验中也能看到同样的情况。如果在眼睛从一处跳到另一处时悄悄替换物体，那么高层神经元在一两个小时内就开始把被替换的物体当成同一个来响应：它们把先后出现的东西联系起来了~\cite{LiDiCarlo2008}。大脑中的不变性确实是从时间的连续性中学到的。这条规则能在多大程度上解释全部视觉学习，还是假说。

\section{视频就是运动}

视频不仅是用于学习的帧流，它本身也是任务：什么在动，往哪里动，动得多快。第一步我们已经熟悉——第~\ref{sec:gaba}~章的方向检测器（图~\ref{fig:direction}），其中延迟的抑制抵消朝某一方向的运动。这样的检测器遍布整个视野，接下来对它们的处理与对形状特征的处理相同：汇集不同位置上许多同方向检测器的级，对大物体的运动作出响应，而与物体所在位置无关。这还是那一对“与”和“或”~\eqref{eq:scstage}，只不过特征不是“边缘”，而是“在时间$d$内移动了的边缘”。

视频还是可预测的：下一帧几乎和上一帧相同。按照\emph{预测编码}假说，高层向低层发送预测，而向上传送的主要是误差——预测没有考虑到的部分~\cite{RaoBallard1999}。这与命题~\ref{prop:economy}中节省脉冲的道理相同：为新消息付费，而不为已知的东西付费。个别观察与这一假说相符，但作为流水线的总体结构，它尚未得到证明。

\section{大脑与卷积网络}

相似性能走多远？对单个物体的快速识别，相似性确实很深~\cite{Yamins2014}。但大脑也有简单卷积网络所没有的东西。

\emph{反馈。}大脑的流水线不只是前向的：每一级都有向后和向侧面的连接。对于困难的图片——有遮挡、光线差——高层的响应形成得较晚，而这些较晚的响应用带反馈的网络描述，比用前馈网络更好~\cite{Kar2019}。走一遍能解决容易的情况，困难的情况需要绕几圈。

\emph{稳健性。}人工网络可能被肉眼察觉不到的像素篡改所欺骗~\cite{Szegedy2014}：人看不出来的改动，会把“熊猫”变成“长臂猿”~\cite{Goodfellow2015}。人的视觉对这种篡改要稳健得多，但并非无懈可击：同时欺骗许多网络的图片在短暂呈现时也会改变人的选择~\cite{Elsayed2018}。为什么稳健性相差如此之大，仍是开放问题；候选因素有噪声、除以总活动和反馈。

\emph{老师。}网络需要数以百万计的标注样本，大脑主要依靠无标注的视频~\eqref{eq:traceinvariance}，外加\nt{DOPA}关于什么重要的少量提示。

\emph{能量。}整个大脑约$20$瓦（命题~\ref{prop:economy}），识别只占其中一部分；在图形加速器上运行的训练好的卷积网络要消耗几百瓦。

\section{错觉：机器在哪里出错，为什么}
\label{sec:illusions}

想弄懂别人电路的工程师，会给它输入不寻常的信号，看它在哪里出错。对视觉来说，这样的信号早已被发明出来——那就是视错觉。错觉不是故障。每种错觉都指向机器的某个特定机制，这个机制在寻常环境中工作正确，却在专门挑选的图片上暴露了自己。

\emph{合理的预期。}输入含有噪声，机器用世界上通常出现的情况来补全它。慢速运动比快速运动常见；当输入不可靠时——例如图片对比度很低——速度估计就偏向慢速，于是暗淡的物体看起来比明亮的物体移动得慢~\cite{Weiss2002}。许多亮度错觉也这样解释：我们看到的不是像素的亮度，而是过去最常与这种亮度对应的表面~\cite{Purves2011}。一张连衣裙的照片，有人看成白金色，有人看成蓝黑色，也是同一机制：图片有歧义，人们在下意识地假定何种光照上意见不一~\cite{LaferSousa2015,Wallisch2017}。人与人之间的差异已经测量；大脑在这里按照概率论的规则把输入与预期结合起来，是一个有充分依据的假说。

\emph{增益调节。}盯着瀑布看一分钟，再看静止的岩石：岩石会向上漂。“向下”和“向上”通道是一对导线，如式~\eqref{eq:dualrail}；增益调节~\eqref{eq:nakarushton}压低了疲劳的那个通道，这对导线的平衡随之偏移。周围环境改变中心外观的倾斜错觉和对比度错觉，用除以邻居的活动来解释~\cite{Schwartz2009}，与第~\ref{sec:gaba}~章相同。也有一个值得警惕的例子：白色栅格交叉处的暗斑长期被解释为简单的邻域抑制，但改变栅格形状的实验表明，这种解释并不成立~\cite{SchillerCarvey2005}。

\emph{延迟补偿。}从眼睛到高层的路径需要几十毫秒，运动的物体在这段时间里已经移动了。如果在它旁边闪现一个静止的点，物体看起来会在闪光的前方，尽管二者实际上是对齐的~\cite{Nijhawan1994}。一种解释是，机器把运动向前外推，以补偿延迟——这与第~\ref{sec:muscles}~章相同：反馈有延迟时，就不得不预测。这种错觉的解释有好几种，争论尚无定论。

\emph{时间编码中的错误。}一幅由色彩过渡鲜明的圆环构成的静止图片，看起来在旋转。高对比度区域比暗淡区域处理得快，延迟之差~\eqref{eq:latency}被方向检测器读成了运动~\cite{Conway2005}。触发这种错觉的是眼睛微小的不自主跳动和眨眼：每次跳动都刷新输入，检测器再次看到“运动”~\cite{OteroMillan2012}。这是第~\ref{sec:code}~章的时间编码被读错了。

\emph{一图两像。}立方体线框时而以一个面朝向我们，时而以另一个面；或者两只眼睛看到不同的图片：知觉每隔几秒就会跳变一次。最好的模型是两种解释之间的相互抑制、缓慢疲劳和噪声~\cite{MorenoBote2007}，也就是在第~\ref{sec:muscles}~章产生步行节律的那个方案~\eqref{eq:halfcentre}。切换的时间由噪声和疲劳决定；按照模型~\cite{MorenoBote2007}，起主要作用的是噪声，疲劳只起辅助作用。

那么人工网络呢？训练来预测视频下一帧的网络，在同样的静止圆环上看到了旋转，而在对照图片上看不到~\cite{Watanabe2018}。训练来去噪和锐化图像的网络，会受到与人相同的颜色和亮度错觉的影响~\cite{GomezVilla2020}。可见，一部分错觉是机器所学任务本身的结果，而不是神经组织的特性。哪些错觉是大脑独有的，这是对任何视觉模型的很好检验。

\begin{keyblock}
\begin{proposition}[错觉是机制的印记]
\label{prop:illusions}
当一个在寻常世界中有用的机制得到不寻常的输入时，就会产生错觉：预期压倒了不可靠的输入，增益调节使一对通道失衡，延迟补偿把物体推向前方，延迟之差被读成运动，带有噪声和疲劳的相互抑制发生跳变。每种错觉都是指向相应机制的测试信号。错觉本身已经测量；对它们的解释从有充分依据到存在争议不等。
\end{proposition}
\end{keyblock}

\section{小结}

\begin{table}[t]
\centering
\footnotesize
\setlength{\tabcolsep}{4pt}
\renewcommand{\arraystretch}{1.15}
\begin{tabular}{>{\raggedright}p{3.0cm}>{\raggedright}p{4.6cm}>{\raggedright\arraybackslash}p{5.2cm}}
\toprule
机器做什么 & 怎么做 & 已知程度 \\
\midrule
在$150\ms$内认出 & 沿大约十级走一遍 & 速度已测量 \\
使响应具有选择性 & 对部件取“与”，$S$级~\eqref{eq:scstage} & 在最初几级已测量 \\
使响应具有不变性 & 对位置取“或”，$C$级；抑制和除法 & 在最初几级已测量；对高层是简化 \\
给出答案 & 最后一级约一百个神经元的频率，线性读出 & 已测量 \\
无标注学习 & 在连续视频上使用带资格迹的赫布规则~\eqref{eq:traceinvariance} & 替换物体会改变响应——已测量；作为主要规则——假说 \\
看到运动 & 方向检测器，然后是同样的“与/或”对 & 已测量 \\
在可预测的内容上节省 & 向上传送预测误差 & 假说 \\
解决困难情况 & 反馈，绕几圈 & 较晚的响应和反馈的作用已测量 \\
以可预测的方式出错 & 错觉：预期、增益调节、延迟、带噪声和疲劳的相互抑制 & 错觉已测量；解释从有依据到有争议 \\
\bottomrule
\end{tabular}
\caption{机器如何识别图像和视频。}
\label{tab:recognition}
\end{table}

识别是由我们已有的部件组装起来的（表~\ref{tab:recognition}）：阈值给出对部件的“与”，相互抑制给出对位置的“或”，除法给出与亮度无关的特性，带资格迹的赫布规则代替了不存在的老师。工程师从视觉那里借来了选择性与不变性的交替，并据此建造了卷积网络；而大脑最重要的本领还没有交出来——几乎不耗能量、从无标注的视频中学习。

\section{习题}

\begin{exercise}[\lvlA]
\label{ex:12:1}
\emph{一级之内的频率。}流水线的一级持续$15\ms$，神经元的频率为每秒$20$个脉冲，噪声相关系数$c=0.1$~\eqref{eq:correlated}。一百个神经元的群体频率的相对误差是多少？群体任意大时的极限是多少？一级之内能分辨多少个频率等级？
\end{exercise}

\begin{exercise}[\lvlA]
\label{ex:12:2}
\emph{一次识别的能量。}在$150\ms$的一次通过中，十级、每级$10^7$个神经元，每个神经元至多发出一个脉冲，每个脉冲耗能$10^{-10}$\,J。(a)~识别占这$150\ms$内全脑能量的多大比例？(b)~一台$300$\,W、$10\ms$认出一张图片的加速器，耗能是它的多少倍？
\end{exercise}

\begin{exercise}[\lvlB]
\label{ex:12:3}
\emph{$S$级的阈值。}$24$点的“视网膜”，图形$A=11011$和$B=10101$，样板~\eqref{eq:scstage}在图形的1上取$+1$，在0上取$-1$。(a)~阈值取多少时，$S$神经元能容忍一个翻转点，却对另一个图形保持沉默？(b)~在$100\times100$的图片上检测十种$5\times5$特征，需要多少个$S$神经元？
\end{exercise}

\begin{exercise}[\lvlB]
\label{ex:12:4}
\emph{一图两像中的一像持续多久。}在跳变方案~\eqref{eq:halfcentre}中，设$\tau\ll\tau_a$，当第1组获胜时，$r_2=0$，$r_1=I-a_1$。(a)~求$I=1$、$\beta=2$、$b=2$、$\tau_a=0.4\,\text{s}$时获胜的持续时间。(b)~$\tau_a$取多少时，图像每三秒切换一次？
\end{exercise}

\begin{exercise}[\lvlB]
\label{ex:12:5}
\emph{建模：不变性的代价。}用\texttt{models/\allowbreak recognition\_models.py}中的函数\texttt{two\_stage}（$4000$次试验）求：点的翻转概率$p$为多少时，$N=24$的$S$、$C$级对每四次错一次？$p=0.1$时，从$N=8$到$N=192$，正确率如何变化？
\end{exercise}

\begin{exercise}[\lvlB]
\label{ex:12:6}
\emph{建模：资格迹的窗口。}用\texttt{models/\allowbreak recognition\_models.py}中的\texttt{trace\_learning}以$0.3\,\text{s}$的停顿运行十次：在$5\ms$到$2\,\text{s}$之间，$\tau_{\text{tr}}$取哪些值时，所有运行都无误地学会不变性？这个窗口的上限从何而来？
\end{exercise}

\begin{exercise}[\lvlB]
\label{ex:12:7}
\emph{任何位置的运动。}在一排间距为$0.1^\circ$的相邻点上放置第~\ref{sec:gaba}~章的方向检测器：来自$A$、延迟为$d$的抑制，若与来自$B$的兴奋相差不超过$5\ms$，就将其抵消；其上是$C$级。要使$C$对每秒$5$到$20^\circ$的反向运动保持沉默，需要哪些延迟？
\end{exercise}

\begin{exercise}[\lvlC]
\label{ex:12:8}
\emph{像素篡改。}要改变模型\texttt{two\_stage}（$N=24$）对干净图形的答案，需要翻转几个点？对同样具有平移不变性的分类器“输入中1多于$3.5$个就是$A$”呢？$p=0.1$时它的正确率是多少？级对为$0.86$。
\end{exercise}

\chapter{机器的输出：大脑如何控制肌肉}
\chaptermark{机器的输出}
\label{sec:muscles}
\begin{chaptergoals}
\item 肌肉如何把脉冲变成力，充当一个低通的数模转换器；
\item 为什么按大小启动单位就以浮点方式设定力，步长与力成正比；
\item 一对只会拉的肌肉如何以力之差设定力矩、以力之和设定刚度；
\item 为什么反馈延迟限制纠正速度，带疲劳的抑制如何产生节律。
\end{chaptergoals}

\section{快速输出}

机器在世界中快速完成的一切，都是用肌肉完成的。说话、注视、迈步、微笑——都是肌肉的收缩。第二种输出，即缓慢的身体化学输出，在第~\ref{sec:chemoutput}~章讨论。在图~\ref{fig:machine}中，输出只是一个标着“肌肉”的箭头；现在来看看箭头后面是什么。

链条的最后一环是\nt{ACET}神经元，也就是表~\ref{tab:types}中注明“通往肌肉的输出”的那一类。每根肌纤维只接受恰好一个这样的神经元的输入，而一个神经元控制一组纤维——从几百根到两千来根~\cite{Feinstein1955mu}。需要精细调节的肌肉中，单位较小；腿部大肌肉中，单位较大。神经元连同它的纤维称为\emph{运动单位}：这是大脑能够单独启动的最小一块肌肉。

神经元在肌肉上的突触与第~\ref{sec:code}~章的皮层突触完全不同。在那里，大部分脉冲都丢失了。在这里，每个脉冲释放的乙酰胆碱比使纤维兴奋所需的量多好几倍~\cite{WoodSlater2001}。这是一个带有\emph{安全系数}的突触：神经元的一个脉冲恰好引起纤维的一次收缩。在机器内部，可以容许不可靠的连接，靠冗余纠正错误；在输出端，一个错误的代价就是一次动作，于是可靠性直接在硬件上买下来。

\section{力：频率与单位数}

一个脉冲引起一次单收缩——抽搐：力在几十毫秒内上升，然后回落。它的一个良好近似是~\cite{Fuglevand1993}
\begin{equation}
h(t) \;=\; F_1\,\frac{t}{T_c}\,e^{\,1-t/T_c},
\label{eq:twitch}
\end{equation}
其中$T_c$是上升时间，量级为$50\ms$，$F_1$是一次抽搐的力。如果脉冲来得更频繁，抽搐就叠加起来：
\begin{empheq}[box=\keybox]{equation}
F(t) \;=\; \sum_k h\bigl(t-t_k\bigr) .
\label{eq:forcesum}
\end{empheq}
这与第~\ref{sec:serocomputer}~章把脉冲变成浓度的低通滤波器相同，只不过现在输出的不是浓度，而是力。肌肉是一个把脉冲转换为力的数模转换器（图~\ref{fig:tetanus}）。在我们的模型中，每秒五个脉冲时，抽搐几乎不重叠，力在$0.2F_1$到$1.1F_1$之间跳动；每秒十五个时，力已保持在$1.8F_1$到$2.2F_1$之间；每秒四十个时，力几乎平稳，约为$5.4F_1$。真实的肌肉在高频脉冲下会饱和，因此线性叠加~\eqref{eq:forcesum}只在每秒几十个脉冲以内成立。

\begin{figure}[t]
\centering
\begin{tikzpicture}[>=Stealth,x=20cm,y=0.55cm]
\draw[->,gray!70!black] (0,0) -- (0.66,0) node[right,font=\small,black] {$t$，s};
\draw[->,gray!70!black] (0,0) -- (0,6.4) node[left,font=\small,black] {$F/F_1$};
\foreach \t in {0,0.2,0.4,0.6}{\draw[gray!60!black] (\t,-0.1) -- (\t,0.1); \node[below,font=\scriptsize] at (\t,0) {\t};}
\foreach \f in {1,3,5}{\draw[gray!60!black] (-0.004,\f) -- (0.004,\f); \node[left,font=\scriptsize] at (0,\f) {\f};}
\draw[thick,blue!60!black] plot file {figures/muscle_twitch_5.dat};
\draw[thick,green!45!black] plot file {figures/muscle_twitch_15.dat};
\draw[thick,red!70!black] plot file {figures/muscle_twitch_40.dat};
\node[font=\scriptsize,blue!60!black,anchor=west] at (0.605,0.9) {5次/秒};
\node[font=\scriptsize,green!45!black,anchor=west] at (0.605,2.1) {15次/秒};
\node[font=\scriptsize,red!70!black,anchor=west] at (0.605,5.4) {40次/秒};
\end{tikzpicture}
\caption{肌肉作为数模转换器：一个运动单位规则发放时的力~\eqref{eq:forcesum}，$T_c=50\ms$。稀疏的脉冲产生分离的抽搐，密集的脉冲融合成平稳的力——正如图~\ref{fig:lowpass}中密集的脉冲融合成平稳的浓度。}
\label{fig:tetanus}
\end{figure}

大脑有两个控制力的旋钮：每个单位的脉冲频率，以及被启动的单位数。第二个旋钮的设计非常巧妙。肌肉中的单位各不相同：最弱的比最强的弱一百倍。当需要的力越来越大时，单位\emph{按大小顺序}启动，从小到大~\cite{Henneman1957}。没有谁去计算这个顺序：小神经元在同样的公共输入下就是更容易兴奋，所以启动顺序由硬件本身决定。

这有什么好处？设每个后续单位比前一个强$q$倍，$q$略大于1。所有已启动单位的力是一个等比数列之和，几乎全部来自最后几个单位。于是，每再启动一个单位，就增加
\begin{empheq}[box=\keybox]{equation}
\frac{\Delta F}{F} \;\approx\; q-1 \;=\; \text{const} ,
\label{eq:sizeprinciple}
\end{empheq}
也就是说，力的步长与力本身成正比。用力小时，大脑以细小的份额来调配；用力大时，份额也大。这是一个\emph{浮点数}：阶码由启动了多少单位决定，尾数由它们的脉冲频率决定。这与第~\ref{sec:sensors}~章传感器的对数刻度相同：机器的输入和输出都以相对比例来度量世界。

浮点也有反面。力的离散程度同样与力本身成正比：用力大的动作必然不如用力小的动作精确。按照一个颇有影响的假说，正是这种依赖于信号的噪声解释了我们的动作为何平滑：平滑的指令使终点的离散最小~\cite{HarrisWolpert1998}。

\section{只能拉，不能推：每个关节两根导线}

肌肉只会拉。它不能推，正如\nt{GLUT}神经元不能输出负信号。解决办法在第~\ref{sec:glutcomputer}~章已经见过：\emph{两根导线}。每个关节由一对向相反方向拉的肌肉负责（图~\ref{fig:antagonists}）。若它们的力为$F_1$和$F_2$，力臂为$\ell$，则关节的力矩和刚度为
\begin{empheq}[box=\keybox]{equation}
M \;=\; \ell\,(F_1-F_2), \qquad K \;\approx\; \kappa_m\,(F_1+F_2),
\label{eq:antagonists}
\end{empheq}
其中$\kappa_m$是把肌肉的力与其刚度联系起来的系数：绷紧的肌肉比放松的肌肉弹性更强。力之差决定力矩，力之和决定刚度~\cite{Hogan1984}。大脑用两根导线控制两个独立的量：关节转向哪里，以及把它握得多稳。要端住一杯水不洒出来时，我们同时绷紧两块肌肉：力矩不变，刚度更高，意外的碰撞就不那么容易使手偏移。

\begin{figure}[t]
\centering
\begin{tikzpicture}[>=Stealth,
  nrn/.style={circle,draw,very thick,minimum size=0.95cm,inner sep=0pt,font=\scriptsize},
  inh/.style={-{Bar[width=7pt]},thick,red!70!black}]
% the joint: a bar on a pivot, two springs pulling down on either side
\fill[gray!30] (4.6,-0.25) -- (5.4,-0.25) -- (5,0.15) -- cycle;
\draw[line width=3pt,gray!50!black] (2.4,0.2) -- (7.6,0.2);
\fill[black] (5,0.2) circle (0.08);
\draw[decorate,decoration={coil,segment length=5pt,amplitude=4pt},very thick,blue!60!black] (3.0,0.2) -- (3.0,-1.6);
\draw[decorate,decoration={coil,segment length=5pt,amplitude=4pt},very thick,orange!85!black] (7.0,0.2) -- (7.0,-1.6);
\draw[very thick] (2.5,-1.6) -- (3.5,-1.6); \draw[very thick] (6.5,-1.6) -- (7.5,-1.6);
\node[font=\small,blue!60!black] at (2.35,-0.75) {$F_1$};
\node[font=\small,orange!85!black] at (7.65,-0.75) {$F_2$};
\draw[->,thick] (5.9,0.75) arc (60:120:1.8) node[above,font=\scriptsize,pos=0.5] {$M=\ell(F_1-F_2)$};
% motor neurons and reflex circuit
\node[nrn,fill=green!12] (M1) at (1.6,-3.0) {\nt{ACET}};
\node[nrn,fill=green!12] (M2) at (8.4,-3.0) {\nt{ACET}};
\node[nrn,fill=red!10] (G) at (5.0,-3.0) {\nt{GLYC}};
\node[nrn,fill=blue!6] (S) at (1.6,-5.0) {\nt{GLUT}};
\draw[->,very thick] (M1) -- (3.0,-1.7);
\draw[->,very thick] (M2) -- (7.0,-1.7);
\draw[->,thick] (S) -- (M1);
\draw[->,thick] (S) -- (G);
\draw[inh] (G) -- (M2);
\draw[->,thick,densely dashed,gray!60!black] (3.15,-1.2) to[bend left=25] node[pos=0.35,right,font=\scriptsize,black] {牵拉} (S.north east);
\draw[->,thick,blue!60!black] (-0.3,-3.0) node[left,font=\scriptsize,black,align=right] {大脑\\指令} -- (M1);
\draw[->,thick,blue!60!black] (10.3,-3.0) node[right,font=\scriptsize,black,align=left] {大脑\\指令} -- (M2);
\end{tikzpicture}
\caption{一个关节上的两块肌肉和最快的反馈环路。\nt{ACET}神经元接收大脑的指令，把关节拉向不同方向；力之差使关节转动，力之和使它更僵硬。左侧肌肉的牵拉传感器（\nt{GLUT}）兴奋该肌肉自己的\nt{ACET}神经元，并经\nt{GLYC}中间神经元抑制拮抗肌的神经元——即第~\ref{sec:gaba}~章的否决。}
\label{fig:antagonists}
\end{figure}

\section{有延迟的反馈}

肌肉不是盲目工作的。每块肌肉里都有表~\ref{tab:sensors}中提到的牵拉传感器，最短的环路直接在脊髓中闭合（图~\ref{fig:antagonists}）。肌肉被意外拉长——牵拉传感器就兴奋该肌肉自己的\nt{ACET}神经元，肌肉收缩，把关节拉回原位。与此同时，经由抑制性中间神经元关闭拮抗肌，使它不来妨碍。担任这个中间神经元的是\nt{GLYC}——表~\ref{tab:types}中那个没有单独成章的类型：它的位置正在这里，在脊髓的快速环路中。

每个环路都有延迟$d$：对手臂而言，脊髓环路约为$25$--$30\ms$，经皮层的环路长一倍半到三倍，经眼睛的环路为$100$--$200\ms$。而我们在命题~\ref{prop:ei}中已经看到，有延迟的反馈会使系统振荡。对最简单的调节器，它根据$d$秒之前的信息以速率$G$纠正偏差$x$，$\dd x/\dd t=-G\,x(t-d)$，稳定条件为~\cite{Hayes1950}
\begin{empheq}[box=\keybox]{equation}
G\,d \;<\; \frac{\pi}{2} ,
\label{eq:delaystable}
\end{empheq}
在边界上系统以频率$1/(4d)$振荡。我们在模型上检验了这一点：$d=30\ms$时，$G=40$每秒时偏差衰减，而$G=55$每秒时，略高于边界值$52$，系统开始振荡。

由此得出两个推论。第一：环路越长，纠正就必须越慢。经过延迟约一百五十毫秒的眼睛，手的纠正不能快于大约十分之一秒，否则手就会颤抖。第二：脊髓环路的稳定边界$1/(4\cdot30\ms)\approx8$次/秒，接近每个健康人都有的细微颤动的频率——每秒八到十二次。牵拉环路是这种颤动的可能来源之一~\cite{McAuleyMarsden2000}。

既然反馈有延迟，大脑又如何做到动作既快又准？一个流行的假说是：\emph{预测}。大脑保存身体的内部模型，不等传感器回报就计算出指令的结果~\cite{WolpertGhahramani2000}。传感器的作用是修正模型，而不是修正每一个动作。工程师会称之为带状态观测器的前馈控制。

\section{运动节律发生器}

步行、呼吸、咀嚼都是节律性动作。它们需要一个时钟来打拍子吗？不需要。脊髓和脑干中有一些小网络，即使输入中没有任何节律性的东西，它们自己也会产生节律~\cite{MarderBucher2001}。最简单的这种网络由我们已有的部件组成：两组\nt{GLUT}神经元，如图~\ref{fig:wta}那样经\nt{GABA}或\nt{GLYC}中间神经元相互抑制，并像第~\ref{sec:glutcomputer}~章的记忆那样缓慢疲劳：
\begin{empheq}[box=\keybox]{equation}
\tau\,\frac{\dd r_i}{\dd t} = -r_i + \bigl[I - \beta\,r_j - a_i\bigr]_+ ,
\qquad
\tau_a\,\frac{\dd a_i}{\dd t} = -a_i + b\,r_i ,
\label{eq:halfcentre}
\end{empheq}
其中$a_i$是第$i$组的疲劳，$j$是另一组。$\beta>1$的相互抑制选出获胜者（命题~\ref{prop:wta}）。获胜者工作并逐渐疲劳；当疲劳吃掉了它的输入，已经休息过来的失败者便冲到前面，自己又开始疲劳。两组轮流工作，各自控制一对肌肉中的一块（图~\ref{fig:halfcentre}）。

\begin{figure}[t]
\centering
\begin{tikzpicture}[>=Stealth,x=3.2cm,y=2.4cm]
\draw[->,gray!70!black] (0,0) -- (4.2,0) node[right,font=\small,black] {$t$，s};
\draw[->,gray!70!black] (0,0) -- (0,1.2) node[left,font=\small,black] {$r$};
\foreach \t in {0,1,2,3,4}{\draw[gray!60!black] (\t,-0.03) -- (\t,0.03); \node[below,font=\scriptsize] at (\t,0) {\t};}
\draw[thick,blue!60!black] plot file {figures/halfcentre1.dat};
\draw[thick,orange!85!black] plot file {figures/halfcentre2.dat};
\node[font=\scriptsize,blue!60!black,anchor=west] at (1.6,1.28) {第1组：“向前”};
\node[font=\scriptsize,orange!85!black,anchor=west] at (2.7,1.28) {第2组：“向后”};
\end{tikzpicture}
\caption{按模型~\eqref{eq:halfcentre}的节律发生器：$I=1$，$\beta=2$，$b=2$，$\tau=20\ms$，$\tau_a=0.4\,\text{s}$。两组神经元相互抑制并会疲劳，以$0.84\,\text{s}$的周期轮流工作，尽管输入是恒定信号。}
\label{fig:halfcentre}
\end{figure}

在我们的模型中，输入恒定时，两组以$0.84\,\text{s}$的周期交替，约为疲劳时间$\tau_a$的两倍。来自上方的恒定指令——“走”——就地变成了节律。这是又一个局部节律，就像第~\ref{sec:gaba}~章\nt{GLUT}--\nt{GABA}环路的节律：它只在网络工作时出现，也只为它所控制的肌肉打拍子。

\begin{keyblock}
\begin{proposition}[机器的输出]
\label{prop:muscles}
大脑像一个分层的调节器体系那样控制肌肉。顶层选择动作（第~\ref{sec:dofa}~章）；下层的局部发生器把恒定指令变成节律，脊髓的快速环路稳住关节。力以浮点方式设定：阶码由按大小启动的单位数给出，尾数由脉冲频率给出。每个关节由一对只会拉的肌肉控制，其力之差决定力矩，力之和决定刚度。这些结构已经测量；大脑依赖身体的内部模型，是一个有充分依据的假说。
\end{proposition}
\end{keyblock}

\section{小结}

\begin{table}[t]
\centering
\footnotesize
\setlength{\tabcolsep}{4pt}
\renewcommand{\arraystretch}{1.15}
\begin{tabular}{>{\raggedright}p{3.0cm}>{\raggedright}p{4.6cm}>{\raggedright\arraybackslash}p{5.2cm}}
\toprule
输出做什么 & 怎么做 & 已知程度 \\
\midrule
把脉冲变成力 & 抽搐叠加~\eqref{eq:forcesum}，低通滤波器 & 已测量；线性叠加是简化 \\
调配力 & 按大小启动单位，浮点~\eqref{eq:sizeprinciple} & 启动顺序已测量 \\
只会拉却能推 & 一对肌肉：力矩为差，刚度为和~\eqref{eq:antagonists} & 已测量 \\
稳住关节 & 牵拉环路，经\nt{GLYC}否决拮抗肌 & 已测量 \\
不颤抖 & $Gd<\pi/2$~\eqref{eq:delaystable} & 来自控制理论；对颤动的贡献是可能的原因 \\
动作迅速 & 依据内部模型预测 & 有充分依据的假说 \\
有节律地行走和呼吸 & 节律发生器~\eqref{eq:halfcentre} & 发生器已测量；最简方案是模型 \\
\bottomrule
\end{tabular}
\caption{机器如何控制肌肉。}
\label{tab:muscles}
\end{table}

机器的输出采用与其他一切相同的手法（表~\ref{tab:muscles}）：用两根导线代替符号，用低通滤波器代替数模转换器，用对数刻度代替线性刻度，用相互抑制和疲劳代替时钟发生器。输入（第~\ref{sec:sensors}~章）和输出都以相对比例度量世界，而在二者之间工作的，就是本书所讲的这台机器。

\section{习题}

\begin{exercise}[\lvlA]
\label{ex:13:1}
\emph{浮点的具体数字。}一块肌肉有$N=100$个运动单位，力为$f_n=f_1q^{\,n-1}$；最强的比最弱的强$100$倍。(a)~求$q$、相对步长~\eqref{eq:sizeprinciple}以及以分贝计的动态范围。(b)~由$100$个相同单位组成、总力相同的“线性数模转换器”，在力为$10f_1$时步长是多少？
\end{exercise}

\begin{exercise}[\lvlA]
\label{ex:13:2}
\emph{安全系数的代价。}肌纤维上的突触每个脉冲释放泊松分布、均值为$m$的量子数，纤维在量子数达到$n_{\text{th}}=20$及以上时兴奋；安全系数$S=m/n_{\text{th}}$。以每秒$20$个脉冲工作的单位，在$S=2$和$S=4$时每天各失败几次？
\end{exercise}

\begin{exercise}[\lvlB]
\label{ex:13:3}
\emph{肌肉作为滤波器。}$T_c=50\ms$的抽搐~\eqref{eq:twitch}是一个线性系统的冲激响应。(a)~求其传递函数$H(s)$和截止频率。(b)~规则脉冲的频率为多少时，力的脉动幅度为平均力的$10\%$？
\end{exercise}

\begin{exercise}[\lvlB]
\label{ex:13:4}
\emph{纠正快不过延迟。}调节器$\dd x/\dd t=-G\,x(t-d)$。(a)~证明当$Gd=1/e$时，偏差以最快的速度无振荡衰减，速率为$1/d$。(b)~对脊髓环路$d=30\ms$和经眼睛的环路$d=150\ms$，求这个$G$和衰减时间常数。
\end{exercise}

\begin{exercise}[\lvlB]
\label{ex:13:5}
\emph{节律发生器的周期。}当$\tau\ll\tau_a$时，模型~\eqref{eq:halfcentre}的周期$T$由方程$(\beta-1)(1-E^{2+b})/(\beta-E)=b\,(1-E^{1+b})/(1+b)$给出，其中$E=e^{-T/(2\tau_a)}$。(a)~求$\beta=b=2$、$\tau_a=0.4$\,s时的$T$。(b)~证明$T$与输入$I$无关，且只有当$b>\beta-1$时才可能产生节律。
\end{exercise}

\begin{exercise}[\lvlB]
\label{ex:13:6}
\emph{发生器建模。}从\texttt{models/\allowbreak muscle\_models.py}导入函数\texttt{halfcentre}（不要直接运行该文件），在$b=0.8$、$1.05$、$1.2$、$1.5$、$2$、$4$以及$I=0.5$、$1$、$2$时测量周期，舍去前两个周期。“走快点”的指令能通过$I$改变步频吗？
\end{exercise}

\begin{exercise}[\lvlB]
\label{ex:13:7}
\emph{刚度与反射。}关节处于牵拉环路中：$\dd\theta/\dd t=-a\theta(t)-G\theta(t-d)$，$a=K/B$，$d=30\ms$。把它化为习题~\ref{ex:13:4}，求使偏差以$15\ms$的时间常数无振荡衰减的最小$a$及所需的$G$。当肌肉共同收缩增强时，应如何改变$G$？
\end{exercise}

\begin{exercise}[\lvlC]
\label{ex:13:8}
\emph{步频旋钮。}按习题~\ref{ex:13:5}和~\ref{ex:13:6}，发生器~\eqref{eq:halfcentre}的输入$I$只改变幅度，不改变步频。用本书的部件提出对模型的最小改动，使得低于某个$I_{\min}$时没有节律，而高于它时周期随$I$增大而缩短，$I$增至三倍时至少缩短一半。求$\tau\ll\tau_a$时的$I_{\min}$，并用模拟验证。
\end{exercise}

\chapter{疼痛：警报信号}
\chaptermark{疼痛}
\label{sec:pain}

\section{是警报，不是测量}

第~\ref{sec:sensors}~章的所有传感器都在测量世界：光有多少，声音多高，压力多大。疼痛则不同。它报告的不是“那里有什么”，而是“危险，放下一切”。对工程师来说，这不是测量通道，而是\emph{警报信号}：一条优先级最高的线路，必须能穿透机器当前正在做的任何工作。

警报与测量通道有三点不同，而疼痛三点都具备。第一，它很难被适应压低：痛觉传感器的适应远弱于触觉传感器，轻度损伤后甚至变得更敏感~\cite{LaMotte1982hyper}；强烈加热后响应减弱，只在其中一类中测到过~\cite{Treede1995heat}。光传感器在恒定背景下会沉默（图~\ref{fig:adaptation}），而一分钟后就沉默的警报毫无用处。第二，它的阈值很高：痛觉传感器只在刺激接近危险时才发放，例如对加热，不同传感器的阈值从$41^\circ$到$46$--$48^\circ$不等，取决于类型~\cite{Treede1995heat, Weidner1999cfib}。第三——这也是本章的重点——警报的音量不仅由传感器决定，也由机器本身决定。同一处伤口在不同情境下疼得不一样。让我们看看它是由哪些已经熟悉的部件组成的。

痛觉传感器是\nt{GLUT}神经元，和第~\ref{sec:sensors}~章的其他感觉神经元一样。其中一部分在释放谷氨酸的同时还释放附录~\ref{sec:othermediators}中的P物质，它缓慢地增强传递。

\section{两根导线：快痛与慢痛}

撞一下手指，你会感到两次疼：先是短促尖锐的，一秒后是钝而持久的。这是两根不同的导线。快导线有绝缘，以$5$到$30$米每秒的速度$v$传导脉冲；慢导线细而裸露，速度为$0.5$--$2$米每秒。来自脚部的信号要走约一米，到达时间
\begin{empheq}[box=\keybox]{equation}
t \;=\; \frac{L}{v}
\label{eq:paintime}
\end{empheq}
对快导线是几十毫秒，对慢导线约一秒。两根导线传送两条消息。快导线报告\emph{哪里}和\emph{何时}：它的脉冲到得更早、更准，如同第~\ref{sec:code}~章的时间编码。慢导线报告\emph{有多糟}，它持久的钝痛使机器在损伤未消退时一直处于警报模式。

\section{脊髓中的闸门}

疼痛信号首先到达脊髓，并在那里通过一道\emph{闸门}~\cite{MelzackWall1965}；这道闸门的具体神经元是在相对晚近才找到的~\cite{Duan2014}。把疼痛继续传往大脑的输出\nt{GLUT}神经元上，汇聚着痛觉导线和触觉导线。旁边还有一个抑制性中间神经元，\nt{GABA}或\nt{GLYC}，它被触觉兴奋（图~\ref{fig:gate}）。触摸关闭闸门。在最初的闸门理论~\cite{MelzackWall1965}中，疼痛反过来关闭这个中间神经元，于是强烈的疼痛打开闸门；这是理论的假设，在已找到的闸门神经元上并未直接证实。

\begin{figure}[t]
\centering
\begin{tikzpicture}[>=Stealth,
  nrn/.style={circle,draw,very thick,minimum size=1.0cm,inner sep=0pt,font=\scriptsize},
  inh/.style={-{Bar[width=7pt]},thick,red!70!black},
  bc/.style={->,thick,densely dashed,orange!85!black}]
\node[nrn,fill=blue!6] (Z) at (6,0) {\nt{GLUT}};
\node[nrn,fill=red!10] (Q) at (3.6,-1.6) {\nt{GABA}};
\draw[->,very thick,red!70!black] (0,0.6) node[left,font=\small,black,align=right] {疼痛 $\nu$} -- (5.45,0.15);
\draw[->,very thick,blue!60!black] (0,-2.6) node[left,font=\small,black,align=right] {触摸 $\mu$} -- (2.9,-2.0);
\draw[->,thick,blue!60!black] (1.8,-2.35) -- (5.5,-0.35) node[pos=0.8,below right=-2pt,font=\scriptsize] {$w_{z\mu}$};
\draw[inh] (1.6,0.5) -- (3.35,-1.1) node[pos=0.5,left=1pt,font=\scriptsize] {$w_{q\nu}$};
\draw[inh] (Q) -- (Z) node[pos=0.5,above left=-1pt,font=\scriptsize] {$\beta$};
\draw[->,very thick] (Z) -- (8.2,0) node[right,font=\small,align=left] {送往大脑：$z$};
\draw[bc] (6,2.1) node[above,font=\scriptsize,black,align=center] {来自上方：\nt{SERO}、\nt{NORA}、阿片类} -- (Z.north) node[pos=0.45,right,font=\scriptsize,black] {增益 $g$};
\end{tikzpicture}
\caption{按模型~\eqref{eq:gate}的脊髓疼痛闸门。输出\nt{GLUT}神经元接收疼痛$\nu$和来自触摸$\mu$的微弱兴奋。抑制性中间神经元（\nt{GABA}或\nt{GLYC}）被触摸兴奋，并按闸门理论的假设被疼痛关闭。增益$g$经广播通道从上方到来。}
\label{fig:gate}
\end{figure}

像第~\ref{sec:gaba}~章那样，用频率把它写出来。中间神经元的频率$q$和输出频率$z$为
\begin{empheq}[box=\keybox]{equation}
q \;=\; \bigl[w_{q\mu}\,\mu + q_0 - w_{q\nu}\,\nu\bigr]_+ ,
\qquad
z \;=\; \bigl[\,g\,(\nu + w_{z\mu}\,\mu) - \beta\,q\,\bigr]_+ ,
\label{eq:gate}
\end{empheq}
其中$\nu$是疼痛输入，$\mu$是触摸输入，$w$是连接权重（第一个下标是接收方，第二个是发送方），$\beta$是抑制强度，$q_0$是中间神经元自身的背景活动，$g$是由上方设定的增益（见下文）。这就是第~\ref{sec:gaba}~章的否决~\eqref{eq:veto}：“有疼痛，但没有触摸”——只有一处修正，作为假设取自闸门理论：强烈的疼痛自己关闭否决神经元。

我们在$w_{q\mu}=1$、$q_0=0.2$、$w_{q\nu}=0.5$、$w_{z\mu}=0.3$、$\beta=1.5$时计算了模型（图~\ref{fig:gatecurves}）。没有触摸时，疼痛从$\nu\approx0.18$开始通过。有触摸时，阈值升到$0.86$，在$\nu=1$时输出降为四分之一，从$1$降到$0.25$。而在剧痛$\nu=2$时，触摸已经无济于事：闸门大开。生活中也是这样：揉一揉撞伤处有帮助，严重创伤时则不然。单纯的触摸，没有疼痛，不会通过闸门：触摸不会引发警报。

\begin{figure}[t]
\centering
\begin{tikzpicture}[>=Stealth,x=5.2cm,y=1.35cm]
\draw[->,gray!70!black] (0,0) -- (2.12,0) node[right,font=\small,black] {疼痛 $\nu$};
\draw[->,gray!70!black] (0,0) -- (0,4.3) node[left,font=\small,black] {$z$};
\foreach \t in {0,0.5,1,1.5,2}{\draw[gray!60!black] (\t,-0.05) -- (\t,0.05); \node[below,font=\scriptsize] at (\t,0) {\t};}
\foreach \f in {1,2,3,4}{\draw[gray!60!black] (-0.01,\f) -- (0.01,\f); \node[left,font=\scriptsize] at (0,\f) {\f};}
\draw[very thick,red!70!black] plot file {figures/pain_gate_notouch.dat};
\draw[very thick,blue!60!black] plot file {figures/pain_gate_touch.dat};
\draw[very thick,orange!85!black,densely dashed] plot file {figures/pain_gate_sens.dat};
\draw[very thick,red!70!black] (0.08,3.9) -- (0.2,3.9) node[right,font=\scriptsize,black] {无触摸};
\draw[very thick,blue!60!black] (0.08,3.45) -- (0.2,3.45) node[right,font=\scriptsize,black] {有触摸};
\draw[very thick,orange!85!black,densely dashed] (0.08,3.0) -- (0.2,3.0) node[right,font=\scriptsize,black] {敏化后，$g=2$};
\end{tikzpicture}
\caption{闸门~\eqref{eq:gate}的输出与疼痛强度的关系。触摸提高阈值，减弱轻度和中度疼痛，但对剧痛无效。敏化（虚线）使增益加倍：同样的疼痛以两倍的音量传出，阈值则降低。}
\label{fig:gatecurves}
\end{figure}

\section{来自上方的音量旋钮}

闸门不仅受下方控制。从脑干下行到闸门的通路会改变式~\eqref{eq:gate}中的增益$g$：\nt{SERO}、\nt{NORA}以及附录~\ref{sec:othermediators}中的阿片肽~\cite{Fields2004}。这正是第~\ref{sec:serocomputer}--\ref{sec:acet}~章的那些广播通道，只不过在这里它们的旋钮控制警报的音量。它们能把旋钮往两个方向拧：同样的下行回路既能减轻疼痛，也能加重疼痛。

机器为什么要调低警报？因为警报是一种中断，而中断并不总是合时宜。逃避捕食者的动物不应因为爪子受伤而停下：先跑，再舔伤口。对缓解的预期也会减轻疼痛，而阻断阿片受体就能消除这一效应的一部分~\cite{Levine1978}：在大脑中算出的预期，沿阿片通道下传到闸门。这幅图景本身已经测量；大脑究竟如何决定音量应该多大，还是假说。

\section{敏化：会学习的警报}

损伤之后，闸门的输出神经元变得更敏感：同样的输入产生更大的输出，阈值降低~\cite{Woolf1983}。在模型~\eqref{eq:gate}中，这是增益$g$的增大，最简单的描述是一个由疼痛本身充电的慢RC电路：
\begin{empheq}[box=\keybox]{equation}
\tau_s\,\frac{\dd g}{\dd t} \;=\; -(g-1) + k_s\,z ,
\label{eq:sensitization}
\end{empheq}
其中$\tau_s$是敏感性恢复正常所需的时间；它比疼痛本身持续得更久，因为在损伤性刺激停止之后敏化仍然持续~\cite{Woolf1983}；$k_s$是充电强度。只要组织受损，闸门的输出就使$g$保持在1以上，伤口附近发生的一切都被放大传出（图~\ref{fig:gatecurves}中的虚线：$g=2$时阈值降到$0.11$，输出加倍）。这是一种合理的警报设置：伤未愈合时，宁可多加小心。

在工程师看来，这是带有慢泄漏的正反馈：疼痛提高增益，增益又加重疼痛。只要环路较弱，伤口愈合后$g$就会恢复正常。但如果环路很强，警报可能在损伤已经不存在之后仍卡在开启状态——就像图~\ref{fig:bistable}中带有强反馈的神经元群。模型~\eqref{eq:sensitization}是一种简化；中枢敏化的存在已经测量。

\section{疼痛既是老师，也是中断}

除了报警，疼痛在机器里还有两项工作。

\emph{老师。}疼痛是负奖励：在误差公式~\eqref{eq:ctd}中，它以$r<0$的形式进入。第~\ref{sec:dofa}~章所讲的一切在这里同样适用：疼痛的先兆会提前引起误差，网络学会回避导致疼痛的事物。对人的实验表明，从疼痛中学习时，大脑的活动正是跟随时间差分误差~\cite{Seymour2004}，而且一部分\nt{DOPA}神经元也对不愉快的事件作出响应（第~\ref{sec:dofaalt}~章）。

\emph{中断。}疼痛把注意力从当前任务上拉开——这是它的本职，而不是副作用~\cite{EcclestonCrombez1999}。在第~\ref{sec:nora}~章中，我们曾就\nt{NORA}的短促簇放电讨论过同样的“中断”作用。而最快的响应甚至不到达大脑：从热物上缩手的动作直接在脊髓中闭合，所用的正是图~\ref{fig:antagonists}中的那对肌肉和对拮抗肌的同样否决。

\begin{keyblock}
\begin{proposition}[警报信号]
\label{prop:pain}
疼痛不是测量，而是高优先级的警报信号。它的适应远弱于测量通道，沿两根导线传送（快的报告“哪里”，慢的报告“有多糟”），通过一道由触摸关闭的闸门（而按闸门理论的假设，剧痛会打开它），其音量由广播通道从上方设定。除了疼痛打开闸门这一点，这些结构都已测量；简单模型~\eqref{eq:gate}和~\eqref{eq:sensitization}是简化。
\end{proposition}
\end{keyblock}

\section{小结}

\begin{table}[t]
\centering
\footnotesize
\setlength{\tabcolsep}{4pt}
\renewcommand{\arraystretch}{1.15}
\begin{tabular}{>{\raggedright}p{3.0cm}>{\raggedright}p{4.9cm}>{\raggedright\arraybackslash}p{4.9cm}}
\toprule
疼痛做什么 & 怎么做 & 已知程度 \\
\midrule
发出警报 & 高阈值，适应弱于触觉 & 阈值已测量；适应的比较是估计 \\
报告“哪里”和“有多糟” & 两根速度不同的导线~\eqref{eq:paintime} & 已测量 \\
通过闸门 & 经\nt{GABA}/\nt{GLYC}的否决~\eqref{eq:gate} & 闸门已测量；疼痛开门是理论假设；模型是简化 \\
改变音量 & 下行的\nt{SERO}、\nt{NORA}、阿片类 & 已测量；选择音量的规则是假说 \\
损伤后学习 & 增益增大~\eqref{eq:sensitization} & 敏化已测量；模型是简化 \\
教会回避 & 式~\eqref{eq:ctd}中的负奖励 & 与时间差分误差的相似性已测量 \\
中断工作 & 夺取注意力；缩回反射 & 已测量 \\
\bottomrule
\end{tabular}
\caption{疼痛作为警报信号。}
\label{tab:pain}
\end{table}

疼痛由我们已经见过的部件组成（表~\ref{tab:pain}）：\nt{GLUT}传感器、两根导线、经抑制性中间神经元的否决、广播式音量旋钮、慢RC电路、预测误差和中断。它只有一点是新的：它是机器唯一一个由机器自己设定音量的输入——一套主人既能调低、也能重新配置的警报系统。

\section{习题}

\begin{exercise}[\lvlA]
\label{ex:14:1}
\emph{两根导线。}信号来自脚部，$L=1$\,m。(a)~一次齐射沿一束同类导线传送，其速度布满范围~\eqref{eq:paintime}。对快、慢两类导线，齐射到达脊髓时被拉长多少？(b)~沿$15$和$1$\,m/s导线传来的两波疼痛，时间差达到$0.1$\,s才可分辨。距离多近时它们融为一体？
\end{exercise}

\begin{exercise}[\lvlB]
\label{ex:14:2}
\emph{触摸何时会痛。}闸门~\eqref{eq:gate}取正文的参数，增益$g$随敏化增大。$g$在多大以内，无论强度$\mu$多大，无痛的触摸都不会通过闸门？$g$为多少时，触摸$\mu=1$开始通过？
\end{exercise}

\begin{exercise}[\lvlB]
\label{ex:14:3}
\emph{敏化的稳定性。}输入恒定，闸门输出对$g$是线性的：$z=gA-C$，$A=\nu+w_{z\mu}\mu$，$C=\beta q\ge0$。(a)~求模型~\eqref{eq:sensitization}的平衡点$g^*$及其稳定条件。(b)~$k_s=0.5$时，分别在无触摸和触摸$\mu=1$时，求使模型失控的疼痛强度。
\end{exercise}

\begin{exercise}[\lvlA]
\label{ex:14:4}
\emph{疼痛的先兆。}时刻$0$的信号预示时刻$T=2$\,s的疼痛，在式~\eqref{eq:ctd}中$r(t)=-P\,\delta(t-T)$，$\tau_\gamma=2$\,s。(a)~求信号时刻$\delta$爆发的面积。(b)~在时刻$t_e=1$\,s的动作取消了疼痛。此刻$\delta$的爆发是多少？它会教给网络什么？
\end{exercise}

\begin{exercise}[\lvlB]
\label{ex:14:5}
\emph{自锁的警报。}给\texttt{models/\allowbreak pain\_gate.py}中的闸门加上动态~\eqref{eq:sensitization}和饱和$z\le4$。触摸$\mu=1$始终存在，损伤$\nu=2$持续时间$T$，然后$\nu=0$。用模拟求出$k_s=4$、$4.6$、$5$、$6$、$8$时，使警报保持开启的最小$T$（以$\tau_s$为单位）。
\end{exercise}

\begin{exercise}[\lvlB]
\label{ex:14:6}
\emph{设计闸门。}在$\beta=1.5$、$w_{z\mu}=0.3$、$g=1$时，选取式~\eqref{eq:gate}中的$q_0$、$w_{q\nu}$、$w_{q\mu}$，使无触摸时疼痛阈值为$0.2$，触摸$\mu=1$时为$1.0$，而在$\nu_\times=2.5$时触摸不再减轻疼痛。增益$g$在多大以内，无痛的触摸不会通过闸门？
\end{exercise}

\begin{exercise}[\lvlC]
\label{ex:14:7}
\emph{音量旋钮。}工作每单位时间带来价值$V$，带着损伤$\nu$工作的代价为$c\nu$，因此当$\nu>V/c$时中断是值得的；当$z>\theta=0.5$时疼痛引起中断。(a)~$V/c=0.25$、$0.5$、$2$时，无触摸的闸门~\eqref{eq:gate}取什么增益$g^*$才能给出这个阈值？(b)~机器可以根据本书中的哪些信号来估计$V$和$c$？
\end{exercise}

\chapter{机器如何学会运动}
\chaptermark{机器如何学会运动}
\label{sec:motorlearning}

\section{三位老师}

在第~\ref{sec:muscles}~章中，我们看到了机器如何驱动肌肉，并停在一个假说上：快速而精确的动作需要身体的内部模型，由它预测指令的结果~\cite{WolpertGhahramani2000}。这个模型从哪里来？它必须学出来，而且学法与我们迄今为止的学习方式不同。

本书已经出现过两位老师。第一位是\emph{没有老师}：赫布规则（第~\ref{sec:dofa}~章）加强经常一起出现的东西，并压缩输入。第二位是\emph{奖励}：\nt{DOPA}向所有神经元广播同一个数，“比预期好还是差”。运动需要第三位老师——\emph{误差}。手够杯子时向左偏了两厘米：这不是“更差”，而是一个向量，它告诉每个输出错在哪个方向、错了多少。工程师会称之为有监督学习。

第二位老师与第三位老师的区别，就是所有输出共用一根导线与每个输出各有一根导线的区别。在命题~\ref{prop:broadcastcost}中，按一个公共数学习时，每多一个其噪声影响结果的神经元，学习就更慢。而如果每个输出$i$都得到自己的误差$e_i$，就可以按最简单的规则改变权重：
\begin{empheq}[box=\keybox]{equation}
\frac{\dd w_{ij}}{\dd t} \;=\; -\eta\,e_i(t)\,x_j(t) .
\label{eq:lms}
\end{empheq}
这是最小均方规则，在自适应滤波器技术中早已为人熟知~\cite{WidrowHoff1960}：权重的变化与其输入和其输出误差的乘积成正比，平均而言沿误差平方的梯度下降。不需要第~\ref{sec:dofa}~章那种靠噪声的搜索：修正的方向直接沿误差导线送来。速度也不随输出数目下降，因为别的输出的误差不会到达这个突触。

\begin{keyblock}
\begin{proposition}[三位老师]
\label{prop:teachers}
赫布规则无需老师，教会的是常见的东西；\nt{DOPA}信号用全网一个数教会的是有利的东西，而且每多一个神经元就更慢；通往每个输出的误差导线按规则~\eqref{eq:lms}教会的是精确的东西，速度与输出数目无关。第三种方式的代价是每个输出一根单独的导线，以及某个知道正确答案的来源。
\end{proposition}
\end{keyblock}

\section{误差导线}

谁能知道动作的正确答案？传感器本身。如果手偏到了目标左边，眼睛和第~\ref{sec:sensors}~章的牵拉传感器会报告这一点。需要一套线路把这个误差送到相应的输出。

大脑中有一个区域，看来正是为此而构造的~\cite{Marr1969,Albus1971}。它的线路异常规整，几乎像晶体，很容易从中认出规则~\eqref{eq:lms}（图~\ref{fig:errorwire}）。

\emph{输入扩展层。}关于指令和身体状态的信号到达由大量小型\nt{GLUT}神经元组成的一层。它们约有七百亿个——比大脑其余部分加起来还多~\cite{Azevedo2009}。每个神经元混合若干输入，信号被展开成一个非常长的向量。工程师熟悉这一手法：把输入的维数升高，在新的空间里，几乎任何所需的关系都能用简单的加权和得到~\cite{Cover1965}。

\emph{输出神经元。}加权和由约三千万个大型输出神经元计算~\cite{Andersen1992cereb}，每个有来自扩展层的约十万个输入。而且它们是\nt{GABA}神经元：学习的结果不是以兴奋、而是以抑制的形式往下传，就像第~\ref{sec:gaba}~章的禁止。

\emph{误差导线。}每个输出神经元还有一个特殊的输入：恰好一根\nt{GLUT}导线，它很少发放，约每秒一次，但每次都在输出神经元中引起强烈的爆发~\cite{Ito2001}。这就是$e_i$：爆发的概率取决于动作误差的方向~\cite{Herzfeld2018}。误差爆发与输入活动同时出现，会削弱这个输入——正是式~\eqref{eq:lms}中的$-\eta\,e_i x_j$项。

\begin{figure}[t]
\centering
\begin{tikzpicture}[>=Stealth,
  gran/.style={circle,draw,thick,fill=blue!6,minimum size=0.32cm,inner sep=0pt},
  outn/.style={circle,draw,very thick,fill=red!10,minimum size=1.1cm,inner sep=0pt,font=\scriptsize},
  inh/.style={-{Bar[width=7pt]},very thick,red!70!black}]
% inputs
\foreach \y/\lab in {1.6/{指令},0.4/{身体状态}}{
  \draw[->,thick,blue!60!black] (-0.6,\y) node[left,font=\scriptsize,black] {\lab} -- (0.35,\y);}
% expansion layer
\foreach \i in {0,...,11}{\node[gran] (g\i) at ({0.8+0.28*mod(\i,2)},{-0.3+0.23*\i}) {};}
\foreach \i in {0,...,11}{\pgfmathsetmacro{\yy}{mod(\i,3)>0 ? 1.6 : 0.4}\draw[gray!60!black] (0.35,\yy) -- (g\i);}
\node[font=\scriptsize,align=center] at (0.95,-1.05) {扩展层\\$\sim7\cdot10^{10}$ \nt{GLUT}};
% parallel wires to the output neuron
\foreach \i in {0,...,11}{\draw[->,gray!70!black] (g\i.east) -- (4.1,{1.0+0.07*(\i-5.5)});}
\node[outn] (P) at (4.6,1.0) {\nt{GABA}};
\node[font=\scriptsize,align=center,above=2pt] at (P.north) {输出神经元\\$\sim10^5$个输入$x_j$，权重$w_{ij}$};
\draw[inh] (P) -- (6.8,1.0) node[right,font=\scriptsize,align=left] {对动作的\\修正};
% error wire
\draw[->,very thick,red!70!black,densely dashed] (4.6,-1.4) node[below,font=\scriptsize,black,align=center] {一根误差导线$e_i$\\\nt{GLUT}，$\sim1$次/秒} -- (P.south);
\end{tikzpicture}
\caption{大脑看来所实现的有监督学习线路。指令和身体状态被庞大的\nt{GLUT}神经元层展开成长向量$x_j$；输出\nt{GABA}神经元以权重$w_{ij}$对其求和。唯一的误差导线送来$e_i$；误差与活跃输入同时出现，就削弱该输入的权重，如规则~\eqref{eq:lms}所示。}
\label{fig:errorwire}
\end{figure}

有一个困难我们在第~\ref{sec:dofa}~章已经见过：误差比导致它的输入来得晚。失误在指令发出几十到几百毫秒后才能看到。因此突触必须记住自己曾经活跃过——需要资格迹，就像规则~\eqref{eq:threefactor}中那样。从实验来看，这条资格迹的长度与每项任务中的反馈延迟相匹配：在误差延迟一百毫秒到达的地方，被削弱得最多的是在误差之前一百毫秒活跃的输入~\cite{Suvrathan2016}。

\begin{keyblock}
\begin{proposition}[误差导线]
\label{prop:errorwire}
在图~\ref{fig:errorwire}所示的线路中，每个输出神经元恰好得到一根误差导线，误差与输入同时出现会削弱该输入。这是应用于被扩展到极高维数的输入上的最小均方规则~\eqref{eq:lms}。线路的结构和同时出现时的削弱已经测量；整条线路作为有监督学习来工作，是主流假说。
\end{proposition}
\end{keyblock}

\section{内部模型即自适应滤波器}

这样的线路学的是什么？最自然的答案是预测。设输入是指令$u(t)$经过一组具有不同时间常数的滤波器，如在扩展层中那样：$x_j(t)=(g_j*u)(t)$。输出是它们的加权和，即对传感器将要报告的内容的预测：
\begin{empheq}[box=\keybox]{equation}
\hat y(t) \;=\; \sum_j w_j\,(g_j*u)(t), \qquad e(t) \;=\; \hat y(t) - y(t) ,
\label{eq:adaptivefilter}
\end{empheq}
而权重按式~\eqref{eq:lms}学习。这是一个\emph{自适应滤波器}：它自行调整自己的传递函数，以复现一个未知系统~\cite{Fujita1982}。这里的未知系统就是自己的身体，连同它的质量、弹性和延迟。学到的滤波器就是第~\ref{sec:muscles}~章的内部模型：它不等传感器回报，就说出传感器将报告什么。

这样的模型有两重用处。第一，它的预测可以代替有延迟的反馈，于是稳定条件~\eqref{eq:delaystable}不再束手束脚：可以依据模型快速纠正。第二，预测与实际之差是关于\emph{意外}的信号。机器自己做的一切，模型都已预测并减去；剩下的是世界做的事。按照一种假说，这就是我们无法给自己挠痒的原因~\cite{Blakemore1998}。

\section{当世界改变时：快与慢}

用一个容易做的实验来检验这套线路。一个人伸手够目标，而世界突然改变：例如手臂受到一个侧向的力。最初几次动作都偏了，随后偏差减小。如果撤去这个力，手臂起初会偏向\emph{另一侧}——模型学到了这个力，仍在补偿它。这直接证明学到的是模型，而不只是“做顺手了”。

实验还揭示了更细微的现象：同时有两个过程在学习——一个快过程，学得快也忘得快；一个慢过程，学得慢却记得久~\cite{Smith2006}。修正在每次动作之后按事件更新：
\begin{empheq}[box=\keybox]{equation}
e = p - (x_f + x_s), \qquad x_f \leftarrow A_f\,x_f + B_f\,e, \qquad x_s \leftarrow A_s\,x_s + B_s\,e ,
\label{eq:tworate}
\end{empheq}
其中$p$是扰动，$x_f$和$x_s$是两个过程学到的修正，$A$表示修正保留到下一次动作的程度，$B$表示它从误差中学习的强度。快过程的$B$大，而$A$明显小于1；慢过程正好相反。

\begin{figure}[t]
\centering
\begin{tikzpicture}[>=Stealth,x=0.034cm,y=1.9cm]
\draw[->,gray!70!black] (0,0) -- (355,0) node[right,font=\small,black] {动作};
\draw[->,gray!70!black] (0,-1.15) -- (0,1.25);
\foreach \v in {-1,1}{\draw[gray!60!black] (-3,\v) -- (3,\v); \node[left,font=\scriptsize] at (0,\v) {$\v$};}
\foreach \n in {100,200,300}{\draw[gray!60!black] (\n,-0.04) -- (\n,0.04); \node[below,font=\scriptsize] at (\n,0) {\n};}
\draw[thin,gray!70!black] plot file {figures/tworate_p.dat};
\draw[thick,orange!85!black] plot file {figures/tworate_fast.dat};
\draw[thick,blue!60!black] plot file {figures/tworate_slow.dat};
\draw[very thick,black] plot file {figures/tworate_net.dat};
\node[font=\scriptsize,gray!50!black] at (120,1.12) {扰动$p$};
\node[font=\scriptsize,black,anchor=west] at (238,0.52) {总和：回升};
\node[font=\scriptsize,blue!60!black,anchor=west] at (150,0.39) {慢过程$x_s$};
\node[font=\scriptsize,orange!85!black,anchor=west] at (60,0.08) {快过程$x_f$};
\draw[densely dashed,gray!60!black] (226,-1.15) -- (226,1.25);
\node[font=\scriptsize,align=center,anchor=south west] at (228,-1.1) {不再有\\误差};
\end{tikzpicture}
\caption{按模型~\eqref{eq:tworate}的两个学习过程，$A_f=0.92$，$B_f=0.1$，$A_s=0.996$，$B_s=0.02$。先是二百次扰动为$p=1$的动作，然后是反向扰动$p=-1$的六次动作，直到总和回到零。虚线之后误差不再送达，总和自行回到原来的修正：快过程忘掉了反向扰动，慢过程还记得正向扰动。}
\label{fig:tworate}
\end{figure}

模型~\eqref{eq:tworate}解释了一个没有它就难以理解的实验（图~\ref{fig:tworate}）。在我们的计算中，二百次带扰动的动作后，修正总和达到$0.84$，其中大部分由慢过程承担，为$0.63$。然后，六次反向扰动的动作使总和回到零：快过程变成负值，$-0.53$，慢过程仍为正值，$0.46$。如果此时不再报告误差，快过程在几十次动作内就忘掉了它的修正，慢过程则要久得多，于是总和在四十次动作内\emph{自行}升到$0.37$：旧技能不经任何训练就回来了。这种自发恢复已在人身上测量；单过程的模型做不到这一点——没有误差时，它只会衰减到零。

为什么要两个过程？第~\ref{sec:acet}~章的最优滤波器给出了一个合理的工程答案。身体因各种原因以不同速度变化：肌肉在几分钟内疲劳，在几个月内增长或衰弱。如果干扰有快有慢，最佳估计就由两个时间常数不同的滤波器组成，各自捕捉自己的那一种~\cite{Kording2007}。快过程是临时修正，“手累了”；慢过程是“手变了”。这目前还是假说，但它解释了为什么一天学会的技能到第二天会部分丢失，以及为什么第二次学得更快。

\section{小结}

\begin{table}[t]
\centering
\footnotesize
\setlength{\tabcolsep}{4pt}
\renewcommand{\arraystretch}{1.15}
\begin{tabular}{>{\raggedright}p{3.0cm}>{\raggedright}p{4.6cm}>{\raggedright\arraybackslash}p{5.2cm}}
\toprule
线路做什么 & 怎么做 & 已知程度 \\
\midrule
有监督学习 & 通往每个输出的误差导线，规则~\eqref{eq:lms} & 线路结构和同时出现时的削弱已测量；有监督学习是主流假说 \\
扩展输入 & 七百亿个\nt{GLUT}神经元 & 神经元数目已测量；扩展的作用是理论 \\
考虑误差延迟 & 与延迟匹配的资格迹 & 已测量 \\
建立内部模型 & 自适应滤波器~\eqref{eq:adaptivefilter} & 有充分依据的假说 \\
以两种速度学习 & 快过程和慢过程~\eqref{eq:tworate} & 后效和自发恢复已测量；两个过程是模型 \\
\bottomrule
\end{tabular}
\caption{机器如何学会运动。}
\label{tab:motorlearning}
\end{table}

现在机器有了三位老师（表~\ref{tab:motorlearning}）。赫布压缩常见的东西；\nt{DOPA}用一个数教会选择有利的东西；误差导线教会精确，而且教得快，因为每个输出都有自己的误差。大脑看来三者都用，并且各用在最便宜的地方：广播信号用于为数不多的决策，单独的导线用于身体的精细调节，在那里正确答案由传感器本身报告。

\section{习题}

\begin{exercise}[\lvlA]
\label{ex:15:1}
\emph{误差导线线路的容量。}线路中有$3\cdot10^7$个输出神经元，每个有$10^5$个输入和自己的误差导线（$1$次/秒）。有$n$个输入的神经元能学会$\approx2n$个向量的任意标注~\cite{Cover1965}。(a)~线路能存储多少个关联？(b)~按式~\eqref{eq:spikeentropy}，取$\Delta t=10\ms$，估计填满一个神经元容量的最短时间。
\end{exercise}

\begin{exercise}[\lvlB]
\label{ex:15:2}
\emph{最小均方规则的速度。}规则~\eqref{eq:lms}的各模式以速率$\eta\lambda_k(C)$衰减。对白噪声输入下的五个滤波器~\eqref{eq:adaptivefilter}，$C_{jk}=2\sqrt{\tau_j\tau_k}/(\tau_j+\tau_k)$。若$\tau_j$按$2$倍递增（$10$--$160\ms$）以及按$4$倍递增，求最快与最慢速率之比。
\end{exercise}

\begin{exercise}[\lvlB]
\label{ex:15:3}
\emph{两个过程的分析。}把取图~\ref{fig:tworate}参数的模型~\eqref{eq:tworate}写成$\mathbf x_{n+1}=M\mathbf x_n+\mathbf b\,p$。(a)~由$M$的特征值求以动作次数计的时间常数。(b)~求恒定$p=1$时的稳态$x_f$、$x_s$及其和。
\end{exercise}

\begin{exercise}[\lvlB]
\label{ex:15:4}
\emph{建模：第二次更快。}按模型~\eqref{eq:tworate}，取\texttt{models/\allowbreak motor\_learning.py}中的参数，求$p=1$时达到$x_f+x_s\ge0.8$所需的动作次数：(a)~未经训练的机器；(b)~像图~\ref{fig:tworate}那样，先做$200$次$p=1$的动作，再用反向扰动使总和归零之后。单过程模型能这样加速吗？
\end{exercise}

\begin{exercise}[\lvlB]
\label{ex:15:5}
\emph{带内部模型的调节器。}对象$\dd x/\dd t=ku$以延迟$d$被观测；调节器$u=-G\hat x$预测$\hat x(t)=x(t-d)+\hat k\int_{t-d}^{t}u\,\dd s$。(a)~$\hat k=k=1$、$d=150\ms$时，求使时间常数为$20\ms$的$G$，并与极限~\eqref{eq:delaystable}比较。(b)~$\hat k=0.4k$时方案是否稳定？
\end{exercise}

\begin{exercise}[\lvlB]
\label{ex:15:6}
\emph{自适应滤波器学习肌肉。}对象是单位面积的抽搐~\eqref{eq:twitch}，$T_c=50\ms$；滤波器~\eqref{eq:adaptivefilter}是$\tau_j=12.5$、$25$、$50$、$100$、$200\ms$的RC电路；输入每$10\ms$取一个新的正态随机数。$\eta=50$和$200$时，规则~\eqref{eq:lms}要多少秒才能把误差降到$0.5\%$？为什么过了$300$\,s权重仍远非最优？
\end{exercise}

\begin{exercise}[\lvlC]
\label{ex:15:7}
\emph{两个过程即最优滤波器。}扰动是两个过程之和$p^{f,s}_{n+1}=A_{f,s}\,p^{f,s}_n+w^{f,s}_n$，噪声方差为$q_f$、$q_s$，误差的观测噪声方差为$R$；卡尔曼滤波器给出式~\eqref{eq:tworate}。在$A_f=0.92$、$A_s=0.996$下，$q_f/R$、$q_s/R$取何值时得到$B_f=0.1$、$B_s=0.02$？若把$q_s$增至四倍，$B_f$、$B_s$如何变化？
\end{exercise}

\chapter{第二个输出：大脑如何控制身体的化学}
\chaptermark{化学输出}
\label{sec:chemoutput}

\section{快输出与慢输出}

在第~\ref{sec:muscles}~章中，我们看到了这台机器的快输出——肌肉。但它还有第二个输出，一个慢输出。大脑把体温、心率、体内的水和糖维持在所需的范围内，让身体准备好奔跑或入睡。为此它不去移动关节，而是释放化学信号。途径有两条。第一条是直接通往内脏器官的神经：通往心脏、血管、肠道和腺体。第二条是血液：一些神经元和腺体不把物质释放到通向邻居的间隙里，而是直接释放进血流。

血液是本书中出现过的最“广播”的总线。第~\ref{sec:serocomputer}--\ref{sec:acet}~章的广播通道把信号发送到大脑的大片区域；血液则把信号发送到身体的每一个细胞。血液绕身体一周大约需要一分钟，因此消息在几十秒内就能到达所有地方，并在物质被分解之前持续几分钟乃至几小时。这种发送根本没有地址。和命题~\ref{prop:receptor}~一样，消息的含义由接收方选择：只有带有该物质受体的细胞才会响应。

\section{油门与刹车：通往器官的两根导线}

最好的例子是心脏。心脏有自己的起搏器，就像第~\ref{sec:serocomputer}~章的 \nt{SERO} 神经元：它自己设定节律，即使切断所有神经，心脏也会以大约每分钟一百次的频率跳动。大脑并不产生节律，只是调节它，而且用的是两根符号相反的导线（图~\ref{fig:gasbrake}）。一根传送 \nt{NORA}：这是\emph{油门}，它加快节律。另一根传送 \nt{ACET}：这是\emph{刹车}，它减慢节律。粗略地说，
\begin{empheq}[box=\keybox]{equation}
f \;\approx\; f_0 \;+\; a_S\,S \;-\; a_P\,P ,
\label{eq:gasbrake}
\end{empheq}
其中 $f_0\approx100$ 次/分是起搏器的固有节律，$S$ 和 $P$ 分别是油门和刹车的活动。静息时刹车是踩着的，心脏通常每分钟跳六七十次；负荷时松开刹车、踩下油门。

\begin{figure}[t]
\centering
\begin{tikzpicture}[>=Stealth,
  blk/.style={draw,very thick,rounded corners=2pt,align=center,font=\small},
  pace/.style={circle,draw,very thick,double,double distance=1.2pt,minimum size=1.8cm,inner sep=0pt,font=\scriptsize,fill=red!8,align=center},
  inh/.style={-{Bar[width=7pt]},thick,red!70!black}]
\node[blk,fill=blue!5,text width=1.6cm] (B) at (0,0) {大脑\\指令};
\node[pace] (H) at (6.2,0) {起搏器};
\draw[->,very thick,orange!85!black] (B.north east) to[bend left=20] node[above,font=\scriptsize,black,align=center] {\nt{NORA}：油门，秒级} (H.north west);
\draw[inh,very thick,blue!60!black] (B.south east) to[bend right=20] node[below,font=\scriptsize,black,align=center] {\nt{ACET}：刹车，亚秒级} (H.south west);
\draw[->,very thick] (H) -- (8.4,0) node[right,font=\small] {频率 $f$};
\node[blk,fill=orange!10,text width=1.6cm] (A) at (0,-2.6) {\nt{ADRE}\\进入血液};
\draw[->,thick,orange!85!black] (B) -- (A);
\foreach \y/\lab in {-2.0/{心脏},-2.6/{血管},-3.2/{肝脏、肌肉}}{
  \draw[->,thick,densely dashed,orange!85!black] (A.east) -- (4.3,\y) node[right,font=\scriptsize,black] {\lab};}
\end{tikzpicture}
\caption{通往心脏起搏器的两根符号相反的导线：油门 \nt{NORA} 慢，刹车 \nt{ACET} 快。下方是走广播总线的同一个油门：\nt{ADRE} 细胞把肾上腺素释放到血液中，所有带有相应受体的器官都会收到信号（虚线箭头）。}
\label{fig:gasbrake}
\end{figure}

这就是第~\ref{sec:glutcomputer}~章的双导线编码和第~\ref{sec:muscles}~章的成对肌肉，只是现在导线控制的不是关节，而是起搏器的节奏。而且两根导线的速度不同。两者都相当于低通滤波器，但刹车传递变化的速度是油门的好几倍\cite{Berger1989}：\nt{ACET} 的路径很短，与受体结合的蛋白质直接打开钾通道，不需要细胞内的第二信使\cite{Logothetis1987gbg}；而 \nt{NORA} 要通过细胞内的一连串反应起作用。工程师会在这里认出双速执行器的手法：快的负责精确的即时微调，慢的负责大而持久的变化。

\nt{ADRE} 的角色也在这里找到了。表~\ref{tab:types}~中写过，它在大脑中的作用尚不清楚。但在大脑之外，它的作用很清楚。油门通路中有一部分细胞不向器官发送导线，而是把肾上腺素直接释放进血液。这是同一个油门，只不过走的是广播总线：几十秒内它同时到达心脏、血管、肝脏和肌肉，并持续几分钟。寻址的油门 \nt{NORA} 调节一个器官；广播的油门 \nt{ADRE} 让整个身体切换到奔跑模式。

\section{带反馈的级联}

许多激素不是由一个腺体释放，而是由一个级联释放。大脑底部的一小群神经元向血液释放第一个信号；它促使中间腺体释放第二个信号；第二个信号再促使终端腺体释放真正作用于身体的激素。而终端激素又回到前面几级并抑制它们。这样就得到一个由三级组成、外加负反馈的放大器——一个电平调节器，就像公式~\eqref{eq:feedback}~中的 \nt{SERO} 系统。

用时间常数为 $\tau$、增益为 $g$ 的 RC 电路描述每一级，用系数 $k$ 描述反馈：
\begin{equation}
\tau\,\frac{\dd x_1}{\dd t} = -x_1 + g\bigl[u - k\,x_3\bigr]_+ ,\qquad
\tau\,\frac{\dd x_2}{\dd t} = -x_2 + g\,x_1 ,\qquad
\tau\,\frac{\dd x_3}{\dd t} = -x_3 + g\,x_2 ,
\label{eq:cascade}
\end{equation}
其中 $u$ 是大脑的指令，$x_3$ 是终端激素的水平。环路增益 $K=g^3k$。平衡时 $x_3=g^3u/(1+K)$；和任何反馈调节器一样，大的环路增益使水平对各级增益的离散不敏感。但每一级都会产生相移：在频率 $\omega=\sqrt3/\tau$ 处，三个相同的级合起来给出半周的相移和八倍的衰减。由此得到控制理论的经典结果：
\begin{empheq}[box=\keybox]{equation}
K \;<\; 8 ,
\label{eq:cascadestable}
\end{empheq}
否则调节器开始振荡，周期约为 $2\pi\tau/\sqrt3\approx3.6\,\tau$。

\begin{keyblock}
\begin{proposition}[精度与稳定性的权衡]
\label{prop:cascade}
带比例反馈的三级级联，其电平误差为 $1/(1+K)$，而只有在 $K<8$ 时才稳定。因此这种调节器保持电平的精度不会优于大约百分之十；试图提高增益会把它变成振荡器。
\end{proposition}
\end{keyblock}

我们在模型~\eqref{eq:cascade}~上验证了这一点（图~\ref{fig:cascade}）。$K=4$ 时，接通后电平稍有振荡，然后稳定下来。$K=7$ 时也会稳定，但很慢。$K=12$ 时振荡既不衰减，也不会无限增长：一级不能输出负信号，于是级联进入平稳的脉动，在平均水平的 $0.83$ 到 $1.24$ 倍之间起伏，周期为 $3.7$--$3.9\,\tau$。

\begin{figure}[t]
\centering
\begin{tikzpicture}[>=Stealth,x=0.3cm,y=1.2cm]
\draw[->,gray!70!black] (0,0) -- (41.5,0) node[right,font=\small,black] {$t/\tau$};
\draw[->,gray!70!black] (0,0) -- (0,2.8) node[left,font=\small,black] {$x_3/x_3^*$};
\foreach \t in {0,10,20,30,40}{\draw[gray!60!black] (\t,-0.05) -- (\t,0.05); \node[below,font=\scriptsize] at (\t,0) {\t};}
\foreach \f in {1,2}{\draw[gray!60!black] (-0.3,\f) -- (0.3,\f); \node[left,font=\scriptsize] at (0,\f) {\f};}
\draw[densely dashed,gray!60!black] (0,1) -- (40,1);
\draw[thick,blue!60!black] plot file {figures/cascade_K4.dat};
\draw[thick,red!70!black] plot file {figures/cascade_K12.dat};
\node[font=\scriptsize,blue!60!black,anchor=west] at (16,0.55) {$K=4$：趋于稳定};
\node[font=\scriptsize,red!70!black,anchor=west] at (16,1.75) {$K=12$：持续脉动};
\end{tikzpicture}
\caption{按~\eqref{eq:cascade}~带反馈的三级级联，终端激素水平以平衡值 $x_3^*$ 为单位。环路增益低于八时电平趋于稳定；高于八时级联自己产生平稳的脉动。}
\label{fig:cascade}
\end{figure}

真实的级联激素确实以脉冲形式释放：例如，应激激素的水平以大约一小时的周期起伏。有一种假说认为，这些脉动正是由各级之间带延迟的反馈本身产生的，并不需要单独的节律发生器\cite{Walker2010}。这与命题~\ref{prop:ei}~中 \nt{GLUT}--\nt{GABA} 环路以及条件~\eqref{eq:delaystable}~中肌肉牵张环路的振荡原因相同：负反馈存在延迟。

\section{激素的脉冲编码}

脉冲不仅可以是副作用，也可以是编码。控制生殖的神经元大约每小时一次，以短促的剂量把信号释放到血液中。如果给接收腺体施加同样的信号，但不是分剂量而是连续施加，它并不会像受到脉冲时那样全力工作，而是沉默下来；如果再改回每小时一次的剂量，功能就会恢复\cite{Belchetz1978}。可见，接收方读取的不是水平，而是剂量的\emph{频率}。

对工程师来说，这里没有什么谜团。在物质长时间作用下会疲劳、不再响应的受体，是一个高通滤波器，就像第~\ref{sec:code}~章中会耗竭的突触~\eqref{eq:depression}。它抑制恒定信号，只让变化通过。对这样的接收方，只能用脉冲传递信息，于是信息从水平转移到剂量的频率和形状中。此外，不同的剂量频率对同一腺体中不同的细胞有不同的作用，所以一条化学线路可以传递不止一个量——就像第~\ref{sec:dofaalt}~章中一根 \nt{DOPA} 导线同时传递快分量和慢分量。

\section{昼夜节律：慢速的模式发生器}

身体最慢的节律是昼夜节律。它由细胞内带延迟的负反馈产生：基因生成的蛋白质在延迟几个小时后抑制自身的生成\cite{TakahashiJS2017}。又是带延迟的负反馈——这是本书中的第四次——又是一个节律，这一次周期约为一天。主发生器是大脑底部一小群神经元，数量为几万个，它们的节律使身体其余细胞同步。

人体这个发生器的固有周期并不正好是一天，平均为 $24.18$ 小时\cite{Czeisler1999}。每天由来自眼睛传感器的光对它进行校准（第~\ref{sec:sensors}~章）。对电子工程师来说，这就是\emph{锁相环}：固有频率为 $\omega_0$ 的振荡器跟踪频率为 $\omega$ 的外部信号，相位差 $\varphi$ 满足方程
\begin{empheq}[box=\keybox]{equation}
\frac{\dd\varphi}{\dd t} \;=\; (\omega-\omega_0) \;-\; K_\varphi\,\sin\varphi .
\label{eq:pll}
\end{empheq}
若 $|\omega-\omega_0|<K_\varphi$，锁定就能保持。周期为 $24.18$ 小时的振荡器每天需要移动约十一分钟；早晨的光通常足以完成这样的移动。在极夜或没有光的洞穴里无法校准，节律就会漂移到自己的固有周期。

昼夜节律发生器不是计算机的时钟发生器。它不为计算计步，也不同步神经元的脉冲。它切换的是\emph{模式}：睡眠与清醒、激素水平、体温、进食准备。它是与第~\ref{sec:serocomputer}--\ref{sec:acet}~章的寄存器同一类的慢速广播寄存器，只是它的值按照一张对准太阳的时间表变化。

\begin{keyblock}
\begin{proposition}[化学输出]
\label{prop:chemoutput}
这台机器的慢输出与快输出由同样的部件构成。各器官接收两根符号相反、速度不同的导线——油门 \nt{NORA} 和刹车 \nt{ACET}——而整个机体同时接收经血液传送的广播信号，其中包括肾上腺素 \nt{ADRE}。激素水平由带负反馈的级联维持，其精度受稳定性限制；一部分激素用剂量频率编码；昼夜节律是一个按光进行相位校准的慢速振荡器。通路的结构、剂量实验和周期实验均已测量；级联脉动由反馈本身产生的解释是假说。
\end{proposition}
\end{keyblock}

\section{小结}

\begin{table}[t]
\centering
\footnotesize
\setlength{\tabcolsep}{4pt}
\renewcommand{\arraystretch}{1.15}
\begin{tabular}{>{\raggedright}p{3.1cm}>{\raggedright}p{4.8cm}>{\raggedright\arraybackslash}p{4.9cm}}
\toprule
化学输出做什么 & 怎样做 & 已知什么 \\
\midrule
调节器官 & 油门 \nt{NORA} 和刹车 \nt{ACET}~\eqref{eq:gasbrake} & 已测量；刹车比油门快 \\
让全身切换到奔跑模式 & 肾上腺素 \nt{ADRE} 进入血液 & 已测量 \\
维持激素水平 & 带反馈的级联，$K<8$~\eqref{eq:cascadestable} & 级联和反馈已测量；脉动源于环路本身是假说 \\
用剂量频率传递信息 & 疲劳的受体相当于高通滤波器 & 对剂量方式的依赖已测量 \\
按昼夜切换模式 & 按光进行相位校准的振荡器~\eqref{eq:pll} & 周期和光校准已测量 \\
\bottomrule
\end{tabular}
\caption{这台机器如何控制身体的化学。}
\label{tab:chemoutput}
\end{table}

这台机器有两个输出（表~\ref{tab:chemoutput}）。快输出是肌肉：毫秒级，寻址，送到每个关节。慢输出是身体的化学：秒、分钟乃至昼夜，经两根导线通往器官，经血液同时送达全身。两个输出的部件是相同的——符号相反的两根导线、起搏器、反馈及其延迟——也正是构成机器内部的那些部件。

\section{习题}

\begin{exercise}[\lvlA]
\label{ex:16:1}
\emph{血液作为滤波器。}某激素从血液中清除的半衰期为 $t_{1/2}=10$~min，每隔 $T=60$~min 以短促剂量释放一次。浓度最大值是最小值的多少倍？剂量的基波被衰减了多少倍？$t_{1/2}$ 为多少时最大值仅为最小值的两倍？
\end{exercise}

\begin{exercise}[\lvlA]
\label{ex:16:2}
\emph{昼夜锁相。}振荡器~\eqref{eq:pll}~的固有周期为 $24.18$~h，光每天移动相位不超过一小时：$K_\varphi=2\pi/24$~rad/d。求稳态相位差 $\varphi^*$（以分钟计）和弛豫时间。外部信号跳变 $\pm6$~h 后，经过多少天失配会小于一小时？
\end{exercise}

\begin{exercise}[\lvlB]
\label{ex:16:3}
\emph{精度与稳定性。}把线性化的级联~\eqref{eq:cascade}~延长为 $n$ 个相同的级。求临界增益 $K_c(n)$，以及 $n=3$ 和 $n=6$ 时可达到的最小误差 $1/(1+K_c)$。当 $n\to\infty$ 时，这个误差趋于多少？
\end{exercise}

\begin{exercise}[\lvlB]
\label{ex:16:4}
\emph{油门与刹车。}在~\eqref{eq:gasbrake}~中，油门和刹车的贡献是 $\tau_S=2$~s 和 $\tau_P=0.4$~s 的 RC 电路的输出，指令分别不超过 $80$ 和 $60$ 次/分。节律 $f_0=100$；静息时刹车 $35$，油门 $0$，$f=65$。最快多久能达到 $f=120$？如果不松开刹车、只用油门呢？
\end{exercise}

\begin{exercise}[\lvlB]
\label{ex:16:5}
\emph{建模：带延迟的级联。}在 \texttt{models/\allowbreak chem\_models.py} 中函数 \texttt{cascade} 的反馈里（复制该函数，不要运行文件）加入延迟 $d$：第一级看到的是 $x_3(t-d)$。求 $d/\tau=0$、$0.5$、$1$、$2$ 时的 $K_c$ 和临界处的周期，并在 $0.9K_c$ 和 $1.1K_c$ 下用模拟加以验证。
\end{exercise}

\begin{exercise}[\lvlB]
\label{ex:16:6}
\emph{读取剂量的接收方。}受体会疲劳：$y=[c-a]_+$，$\tau_a\,\dd a/\dd t=c-a$，$\tau_a=30$~min。激素以每次 $6$~min 的剂量、周期 $T$ 到达，平均剂量 $\bar c$ 不变。求 $T=15$、$60$、$120$~min 以及 $T\to\infty$ 时的平均响应。浓度恒定为 $\bar c$ 时响应是多少？
\end{exercise}

\begin{exercise}[\lvlC]
\label{ex:16:7}
\emph{振荡器还是反馈？}设计一个实验，区分由级联~\eqref{eq:cascade}~的延迟反馈产生的脉动与由独立起搏器产生的脉动。如果 $K$ 只能改变一点五倍，而周期的测量精度为 $5\%$，能否依据周期的偏移区分这两种假说？
\end{exercise}

\chapter{调试这台机器}
\chaptermark{调试这台机器}
\label{sec:debugging}

\section{如何调试一台没有电路图的机器}

拿到一块能工作却没有电路图的电路板，工程师会用四种手法来调试它。他放上\emph{探针}，看导线上传的是什么。他\emph{注入测试信号}，看什么发生了变化。他\emph{断开}电路的一部分，看什么停止了工作。他还\emph{读取控制寄存器}，以了解电路板处于什么模式。整本书都是在尝试读出大脑的电路图；本章来看看，这四种手法中哪些工程师已经会用，以及前面各章对可以期待什么有何启示。

当今最引人注目的这类调试是脑机接口：它们读取神经元的脉冲并把它们变成外部机器的指令，或者反过来，从外部向大脑送入信号。本章大部分篇幅都讲它们，首先是 Neuralink 公司在做的事。

\section{探针：电极听到了什么}

大脑的探针是插入组织中的细电极。它能听到约一百微米半径内神经元的脉冲——通常是一个到几个。电极上的一个脉冲是几十到几百微伏、持续约一毫秒的尖峰，根据其波形可以把相邻的神经元区分开。最容易听到的是大的 \nt{GLUT} 神经元：它们在皮层中占多数（表~\ref{tab:types}），而且电流更强。

主要困难一眼可见。今天好的探针有几百或几千个电极，而大脑中有 $8.6\cdot10^{10}$ 个神经元。我们窃听的大约是这台机器的一亿分之一。为什么从如此微不足道的样本中竟能理解些什么？答案在第~\ref{sec:code}~章。相邻神经元的噪声是相关的，对群体求平均会碰到一个极限（命题~\ref{prop:pooling}）：一百个神经元携带的信息几乎和一千个一样多。此外，运动皮层要解决的问题是低维的：手臂的关节不多（第~\ref{sec:muscles}~章），控制它的几百个神经元的活动可以放进一个仅由一二十个公共变量构成的空间\cite{Gallego2017}。一个覆盖一百个神经元的探针几乎能听到所有这些变量。如果大脑在每个神经元里各存一条独立的消息，这样的窃听就毫无用处。

\section{数据流：直接在探针上压缩}

探针输出的数据流非常大。要看清脉冲的波形，每个电极需要以约十位的精度、每秒约 $20\,000$ 次进行数字化。一千个电极给出
\begin{empheq}[box=\keybox]{equation}
\begin{aligned}
R_{\text{原始}} \;&=\; N\,f_s\,b \;\approx\; 10^3\cdot2\cdot10^4\cdot10 \;\approx\; 2\cdot10^8\ \text{bit/s},\\
R_{\text{脉冲}} \;&\approx\; N\,r\,\log_2\frac{e}{r\,\Delta t} \;\approx\; 8\cdot10^4\ \text{bit/s}.
\end{aligned}
\label{eq:probedata}
\end{empheq}
第二个估计是 $r=10$ 次/秒、精度 $\Delta t=1\ms$ 时脉冲流的容量~\eqref{eq:spikeentropy}：其中真正包含的内容要小两千五百倍。对植入式设备来说这是决定性的差别：在头颅内能承受的、不会让组织发热的功耗下，原始数据流无法通过无线电穿过皮肤传出去。因此植入式设备需要在电极旁边的芯片上直接提取脉冲，只把事件——通道号和脉冲时刻——传出去。Neuralink 系统的第一份描述还没有做到这一点：芯片放大并数字化信号，整个原始数据流经电缆送出，脉冲由外部程序提取；芯片上的压缩和无线连接在那里被列为计划\cite{Musk2019}。

这是一个熟悉的手法。眼睛在把信号送上长电缆之前先将其压缩一百倍（第~\ref{sec:sensors}~章）；事件相机传输的不是帧，而是变化。为了挤进导线，探针不得不做大脑自己在做的事：用脉冲说话，而且只报告新消息。

\section{Neuralink 在做什么}

Neuralink 的设备就是为这些限制而打造的探针\cite{Musk2019}。它不用刚性的针阵列，而是用许多比头发还细的柔性细丝，每根细丝上都排列着电极：在系统的第一份描述中，最多 $3072$ 个电极分布在 $96$ 根细丝上，每根 $32$ 个；而在植入人体的设备中，据公司介绍，约一千个通道分布在 $64$ 根细丝上。柔性细丝对组织的损伤较小，也更能承受大脑在颅内不断轻微移动。细丝由机器人避开血管插入。如上所述，在第一份描述中原始数据流~\eqref{eq:probedata}~经电缆送出。据公司介绍，用于人体的设备结构不同：外壳完全被皮肤覆盖，脉冲在内部提取，数据经无线电送出，电力以无线方式输入。

第一台设备的任务是读取。细丝被放在计划手臂运动的那部分皮层中，解码器把脉冲变成屏幕上光标的运动。双手瘫痪的人通过想象运动来操作计算机。这个想法本身并不新：基于一百个电极刚性阵列的接口早已能让人控制光标\cite{Hochberg2006}、通过想象手写动作来写字\cite{Willett2021}，甚至能以每分钟约六十个词的速度把说话的尝试转换成屏幕上的文字\cite{Willett2023}。Neuralink 的新意在于工程：通道数、无线化、封闭的外壳和机器人插入，也就是尝试做一种可以在实验室外佩戴多年的探针。据我们所知，人体试验的结果主要由公司自己发布，而不是发表在同行评审的论文中；公司报告过，部分细丝在插入后发生了移位，工作的通道数因此减少。对调试来说这是常见的麻烦：相对电路板移动了的探针，听到的已经是别的导线。

公司的第二个方向是写入：一种为失明者向视觉皮层送入信号的视觉设备。下面在写入一节中再谈。

\section{读取：解码器作为接收方}

接口的解码器就是命题~\ref{prop:reader}~意义上的接收方：它自己决定读取消息的哪一部分。几乎所有接口读取的都是\emph{群体频率}：在几十毫秒的窗口内统计每个通道的脉冲数，再以一组权重相加，权重的选择使结果成为想要的运动。这就是第~\ref{sec:rate}~节中的群体频率编码，选择它并非偶然：每个窗口只有几个脉冲时，单个神经元的频率无法确定，而对一百个神经元求和就管用了。

能取出多少信息？用“意念之手”写字的速度约为每分钟九十个字符\cite{Willett2021}，即每秒一个半字符：即使每个字符五比特，也不到每秒八比特。据 Neuralink 介绍，控制光标可达每秒约十比特。而上界~\eqref{eq:spikeentropy}~对\emph{单个}神经元是每秒几十比特，对一百个神经元是几千比特。接口从被窃听的神经元中取出的，只是它们本可承载的一小部分。瓶颈不在导线，而在于群体携带多少个独立变量（命题~\ref{prop:pooling}），以及解码器对编码理解得有多好——正如第~\ref{sec:code}~章所见，这种编码尚未被完全破译。

还有第二个特点，我们在第~\ref{sec:motorlearning}~章已经熟悉。解码器和大脑相互学习。大脑看到光标去了哪里，得到误差并调整自己的指令——这正是经由误差导线的运动学习。而工程师不时重新拟合解码器的权重，因为细丝下的神经元在变化，网络中的连接在重新学习。于是有了两个相互适应的学习者；为了让这一对不失控，最好把它们的学习速度拉开：让一个学得快，另一个学得慢。

\begin{keyblock}
\begin{proposition}[为什么覆盖一千个神经元的探针能工作]
\label{prop:probe}
只窃听极小一部分神经元的接口之所以能控制运动，是因为群体活动是低维的，而且噪声是相关的：几百个神经元就携带了几乎全部公共变量。出于同样的原因，接口每秒只能取出几比特——其数量取决于任务中独立变量的个数，而不是脉冲流的容量。接口的工作性能已有测量；一旦理解了编码，解码器能改进多少仍是开放问题。
\end{proposition}
\end{keyblock}

\section{写入：比读取更难}

向机器送入信号比读取更难。通电的电极激发的不是一个神经元，而是它周围所有的神经元和导线，不分类型和符号。用它无法写入那种重要的是\emph{哪个}神经元、\emph{何时}发放的编码（第~\ref{sec:code}~章）：只能粗略地设定\emph{在哪里}以及\emph{有多强}。

对视觉来说，这在一定程度上够用了。人会把通过视觉皮层电极的电流看成视野中某个位置上的一个光点；当同时点亮多个光点、最多十六个时，失明的人能认出一些字母并分辨物体的边界\cite{Fernandez2021}。这是位置编码，它是可以写入的。但回想第~\ref{sec:sensors}~章：眼睛送进大脑的不是像素，而是压缩的描述——边缘、变化、变亮和变暗，分别走不同的导线。向皮层写入“光点”，就是向一台在输入端等待完全不同编码的机器写入像素。因此，人工视觉目前比按电极数目所能期待的要粗糙。也有一些不需要电极的测试信号：视错觉通过眼睛送入不寻常的输入，使机器偏离正常模式，而每一种错误都指向各自的机制（第~\ref{sec:illusions}~节）。

\section{探针看不到什么}

电极听到的是寻址总线——\nt{GLUT} 神经元的脉冲，以及较难听到的小 \nt{GABA} 神经元的脉冲。但在第~\ref{sec:serocomputer}--\ref{sec:acet}~章中我们看到，整台机器的工作模式由广播寄存器设定：\nt{SERO}、\nt{DOPA}、\nt{NORA}、\nt{ACET} 的浓度。电极听不到它们：源神经元少得可怜，而且它们的信号不是探针旁边的一个脉冲，而是一个体积中的浓度。能看到数据总线却看不到控制寄存器的调试者，总会感到意外：同一个输入会得到不同的回答，因为某处的 $g$、$a$ 或 $\tau_\gamma$ 变了。浓度可以用化学传感器直接在组织中测量\cite{Bunin1998}，但接口中还没有使用这类传感器。完全处于探针视野之外的还有机器的第二个、慢速的输出——血液中的激素（第~\ref{sec:chemoutput}~章）：它们是在血样中测量的，而不是用电极。

探针也看不到权重——而几乎所有记忆和学到的东西恰恰存储在权重中。只能根据回答如何变化来推测它们。探针也看不到人所感受到的疼痛：在第~\ref{sec:pain}~章中我们看到，警报信号的强度由机器自己设定——由脊髓中的闸门和上方的调节器设定，因此无法从传感器读出有多疼。

\section{小结：调试与开放问题}

\begin{table}[t]
\centering
\footnotesize
\setlength{\tabcolsep}{4pt}
\renewcommand{\arraystretch}{1.15}
\begin{tabular}{>{\raggedright}p{2.2cm}>{\raggedright}p{3.6cm}>{\raggedright}p{3.4cm}>{\raggedright\arraybackslash}p{4.0cm}}
\toprule
调试手法 & 接口做什么 & 本书的解释 & 已知什么 \\
\midrule
探针 & 听到几百到几千个神经元 & 低维性和相关噪声（第~\ref{sec:code}~章） & 已测量 \\
探针上的压缩 & 传输事件而不是原始信号 & 脉冲流的容量~\eqref{eq:spikeentropy} & 工程上已解决 \\
读取 & 群体频率 $\to$ 运动、文字、言语 & 解码器是接收方，群体频率编码 & 可用；每秒几比特 \\
共同学习 & 大脑和解码器相互适应 & 误差导线（第~\ref{sec:motorlearning}~章） & 已观察到；适应规则是研究课题 \\
写入 & 电流在视野中点亮光点 & 位置编码可以写入，时间编码不行（第~\ref{sec:sensors}~章） & 光点已测量；有用的视觉尚未实现 \\
读取寄存器 & 几乎没有 & 广播通道（第~\ref{sec:serocomputer}--\ref{sec:acet}~章） & 实验中已有化学传感器 \\
\bottomrule
\end{tabular}
\caption{调试这台机器：脑机接口能做什么，本书各章对它们有何解释。}
\label{tab:debugging}
\end{table}

表~\ref{tab:debugging}~汇总了本章。脑机接口已经能把探针放到寻址总线上，并从中每秒读出几比特；这件事之所以行得通，要归功于第~\ref{sec:code}~章的低维性和相关噪声。它们向机器写入的能力还很粗糙，因为机器的输入编码是压缩的，并且分别寻址到各条导线。而控制寄存器和权重——让机器可学习、可调节的东西——探针目前几乎看不到。

第~\ref{sec:synthesis}~章的每个开放问题同时也是调试者的任务。读取哪种编码，要等探针更密、解码器学会利用脉冲时间而不仅是频率时才能决定。多层网络如何学习，可以把探针同时放进几层，观察它们的回答在学习中如何变化来检验。广播通道传递的是什么，要等电探针加上化学探针后才能看清。本书尝试依据机器的行为和任何工程师都知道的定律来读出它的电路图。正在建造的调试工具将让我们能用探针检验这张电路图。

\section{习题}

\begin{exercise}[\lvlA]
\label{ex:17:1}
\emph{无线链路的预算。}穿过皮肤传输每比特需 $1$~nJ，发射机最多可用 $10$~mW。$N=1000$ 个电极时，原始数据流~\eqref{eq:probedata}~和脉冲流各需要多大功率？原始数据流要求每比特多少能量？
\end{exercise}

\begin{exercise}[\lvlA]
\label{ex:17:2}
\emph{一百个神经元有多少比特。}解码器在 $50\ms$ 窗口内对 $N$ 个神经元的脉冲求和，$r=20$~次/秒，噪声的两两相关 $c=0.1$，信号使群体频率改变 $30\%$（均方根）。估计 $N=10$、$100$、$1000$ 和 $N\to\infty$ 时每个窗口的速率 $\tfrac12\log_2(1+\text{SNR})$。$c=0$、$N=1000$ 时呢？
\end{exercise}

\begin{exercise}[\lvlB]
\label{ex:17:3}
\emph{键盘式接口的速率。}接口每秒一点五次从 $N=32$ 个字符中选一个：以概率 $P$ 选中想要的字符，错误等概率。$P=0.99$、$0.94$、$0.8$ 时它每秒传输多少比特？$P=0.94$ 时要达到 $10$~bit/s，每秒需要多少个字符？
\end{exercise}

\begin{exercise}[\lvlB]
\label{ex:17:4}
\emph{建模：群体解码器。}$N$ 个神经元，$\lambda_i=[b+m\,\mathbf v\cdot\mathbf p_i]_+$，$b=20$，$m=15$~次/秒，窗口 $50\ms$，$\mathbf v$ 每个分量的标准差为 $0.5$；所有神经元看到共同的噪声 $\mathbf v+\boldsymbol\xi$，$\sigma_\xi=0.2$。模拟无偏的群体向量解码器。$N$ 为多少时两种噪声相等？$N\to\infty$ 时每个窗口每个分量给出多少比特？
\end{exercise}

\begin{exercise}[\lvlB]
\label{ex:17:5}
\emph{两个学习者。}大脑输出指令 $m=b\,v^*$，解码器输出 $v=d\,m$，$\langle v^{*2}\rangle=1$。每次尝试后，两者都以速率 $\eta$ 沿 $\tfrac12(v-v^*)^2$ 做一步梯度下降。从 $b=1.2$、$d=1$ 出发，模拟 $\eta=0.8$、$1.1$、$1.5$、$2$。各自单独时在 $\eta<2$ 下稳定；这一对在 $\eta$ 为多少时稳定？
\end{exercise}

\begin{exercise}[\lvlB]
\label{ex:17:6}
\emph{设计：探针流水线。}$1024$ 个通道，每秒 $20\,000$ 个采样，神经元 $10$~次/秒，无线电 $1$~nJ/bit。对采样的白噪声取多高的阈值，才能使每个通道的虚假脉冲不超过 $0.1$~Hz？每 $20\ms$ 传输一次脉冲计数，每次 $2$ 比特：功率是多少？比传输原始 $10$ 位采样小多少倍？
\end{exercise}

\begin{exercise}[\lvlC]
\label{ex:17:7}
\emph{接口速率的极限。}对 $D=10$ 个带宽为 $W=5$~Hz 的变量，在 $\text{SNR}\approx0.9$（类似信号的噪声）和 $90$（$1000$ 个神经元的独立噪声）时估计 $R\le DW\log_2(1+\text{SNR})$。实测约为 $10$~bit/s。什么实验能区分这一差距的两种解释：类似信号的噪声和对编码的无知？
\end{exercise}

\chapter{按电学功能划分的神经元类型}
\chaptermark{作为器件的神经元}
\label{sec:eetypes}

在第~\ref{sec:neurontypes}~章中，我们按照神经元输出端释放的神经递质对它们进行了分类。电子工程师会用另一种方式分类：看神经元对自己的信号做了什么——它作为哪种器件工作。本书的主要器件我们从第~\ref{sec:glutcomputer}~章起就认识了：泄漏积分神经元~\eqref{eq:glutneuron}，即带比较器的RC电路。但大脑中还有别的器件。表~\ref{tab:eetypes}~把它们分成四组，而且几乎每一种都已在书中出现过。

\begin{center}
\footnotesize
\setlength{\tabcolsep}{3pt}
\renewcommand{\arraystretch}{1.15}
\begin{longtable}{>{\raggedright}p{3.1cm}>{\raggedright}p{3.3cm}>{\raggedright}p{4.7cm}>{\raggedright\arraybackslash}p{2.4cm}}
\caption{作为器件的神经元：名称、电子学中的对应物、器件的作用以及它在书中出现的位置。}\label{tab:eetypes}\\
\toprule
神经元 & 器件 & 作用 & 书中位置 \\
\midrule
\endfirsthead
\toprule
神经元 & 器件 & 作用 & 书中位置 \\
\midrule
\endhead
\bottomrule
\endfoot
\multicolumn{4}{l}{\emph{滤波器和积分器}} \\
泄漏积分神经元 & 带比较器的RC积分器 & 在窗口 $\tau$ 内对输入求和，到达阈值时发放 & 第~\ref{sec:glutcomputer}~章，\eqref{eq:glutneuron} \\
符合检测器 & 带时间窗的与门 & 只有输入几乎同时到达时才发放；$\tau$ 很小 & 第~\ref{sec:glutcomputer}、\ref{sec:code}~章，\eqref{eq:window} \\
相位型神经元 & 高通滤波器，边沿检测器 & 对输入的变化作出响应，随后沉默 & 第~\ref{sec:sensors}~章 \\
适应型神经元 & 带自动增益控制的高通滤波器 & 输入恒定时频率下降 & 第~\ref{sec:sensors}~章，\eqref{eq:nakarushton} \\
谐振器 & 振荡回路，带通滤波器 & 对自身频率的输入响应最强 & 第~\ref{sec:resonator}~节 \\
无泄漏积分器 & 运算放大器积分器 & 把速度变成位置并保持它 & 第~\ref{sec:perfectint}~节 \\
\midrule
\multicolumn{4}{l}{\emph{转换器}} \\
紧张型神经元 & 电压—频率转换器 & 频率线性跟随输入，几乎不疲劳 & 第~\ref{sec:code}~章 \\
分级电位神经元（无脉冲） & 模拟放大器，缓冲器 & 在短距离上传递平滑的电压 & 第~\ref{sec:sensors}~章 \\
整流神经元 & 半波整流器 & 频率不会为负，$[u]_+$ & 第~\ref{sec:glutcomputer}、\ref{sec:gaba}~章 \\
“开—关”对 & 一对整流器，双导线编码 & 一个响应“更多”，另一个响应“更少” & 第~\ref{sec:glutcomputer}~章，\eqref{eq:dualrail} \\
带分流抑制的神经元 & 模拟除法器，增益控制 & 对输入做除法，而不是做减法 & 第~\ref{sec:gaba}~章，\eqref{eq:shunt} \\
\midrule
\multicolumn{4}{l}{\emph{振荡器与时间}} \\
起搏器 & 弛张振荡器，多谐振荡器 & 自己以恒定频率输出脉冲 & 第~\ref{sec:serocomputer}、\ref{sec:dofa}~章 \\
带输入的起搏器 & 压控振荡器 & 有自己的节律，输入使其加快或减慢 & 第~\ref{sec:chemoutput}~章 \\
簇放电神经元 & 门控簇脉冲发生器 & 输出密集脉冲的簇——“划” & 第~\ref{sec:code}~章，\eqref{eq:burst} \\
锁相神经元 & 鉴相器，锁相环 & 在输入波的特定相位发放 & 第~\ref{sec:sensors}、\ref{sec:chemoutput}~章 \\
带延迟的轴突 & 延迟线 & 把信号延迟一段给定的时间 & 第~\ref{sec:glutcomputer}~章，图~\ref{fig:itd} \\
\midrule
\multicolumn{4}{l}{\emph{存储与切换}} \\
双稳态神经元 & 施密特触发器，锁存器 & 短暂的输入使它进入持久状态 & 第~\ref{sec:bistable}~节 \\
带树突分支的神经元 & 多路复用器，小型多层网络 & 只有存在允许输入时，分支才让输入通过 & 第~\ref{sec:synthesis}~章 \\
\end{longtable}
\end{center}

有一条说明比整张表更重要：这些不是不同的细胞，而是不同的\emph{模式}。同一个神经元在输入缓慢时可以是积分器，在抑制缩短了它的 $\tau$ 时又可以是符合检测器（第~\ref{sec:code}~章）；清醒时是起搏器，睡眠时是沉默的接收器。最终得到哪种器件，由膜上的离子通道以及调节这些通道的广播通道决定。

\section{书中尚未出现的四种器件}

\subsection{谐振器}
\label{sec:resonator}

有些神经元具有一种慢电流，它会对抗电压的任何变化：电压升高，电流把它往下拉；电压下降，电流把它往上拉，但有延迟。对电子工程师来说，这种电流的行为像一个电感，它与膜电容一起构成振荡回路。神经元对特定频率的输入响应最强，这个频率通常在几赫兹到一二十赫兹之间，对更慢和更快的输入响应较弱\cite{HutcheonYarom2000}。这是单个细胞中的带通滤波器：它从嘈杂的输入中挑出自身频率的节律，例如第~\ref{sec:gaba}~章的局部节律。

\subsection{无泄漏积分器}
\label{sec:perfectint}

泄漏积分器只能记住输入 $\tau$ 这么长的时间——几十毫秒。但眼睛要保持不动，就需要把自己的位置记住好几秒：“转动”指令以短暂的速度形式到来，而眼睛必须保持在新的位置上。网络用与第~\ref{sec:glutcomputer}~章的记忆相同的反馈解决这个问题。如果一群时间常数为 $\tau$ 的神经元以权重 $w$ 自我兴奋，那么 $\tau\,\dd r/\dd t=-r+w\,r+I$，这群神经元的时间常数为
\begin{empheq}[box=\keybox]{equation}
\tau_{\text{eff}} \;=\; \frac{\tau}{1-w} .
\label{eq:perfectint}
\end{empheq}
当 $\tau=0.1\,\text{s}$、$w=0.995$ 时得到 $20\,\text{s}$：用每个都只能记住十分之一秒的神经元，造出了几乎理想的积分器\cite{Seung1996}。代价是精细调节：$w$ 稍大于一，积分器就会饱和；稍小于一，它就很快遗忘。因此这种积分器需要通过学习不断调节，就像第~\ref{sec:motorlearning}~章所描述的那样。

\subsection{双稳态神经元}
\label{sec:bistable}

在第~\ref{sec:glutcomputer}~章和第~\ref{sec:gaba}~章中，触发器是由一群神经元组装的。但单个细胞中也有触发器。有些神经元具有一种电流，一旦开启就会自己维持兴奋：短暂的输入使神经元进入持久的工作状态——\emph{平台}——而抑制使它回到静息。这就是施密特触发器：神经元有两个稳定状态，两者之间有迟滞。例如，控制肌肉的\nt{ACET}神经元就是这样，而这一特性由广播通道开启：当血清素到达神经元时，平台才会出现\cite{Hounsgaard1988}。同一个细胞有\nt{SERO}时是锁存器，没有时是普通的积分器。

\subsection{积分器与谐振器}

理论把可兴奋细胞分为两类\cite{Izhikevich2007}。\emph{积分器}像一个从零开始的压控振荡器：略高于阈值时神经元可以工作得任意慢，频率随输入平滑增加。这样的神经元对到来的一切求和，能用频率很好地传递输入的大小。\emph{谐振器}像一个有最低频率的振荡器：它们不会慢速工作，在阈值输入下立即以有限的频率输出脉冲，并偏好自身频率的输入。前者是第~\ref{sec:code}~章频率编码的天然器件，后者是时间编码的天然器件。

\begin{keyblock}
\begin{proposition}[一个神经元，多种器件]
\label{prop:eetypes}
按电学功能划分，神经元可以作为有泄漏和无泄漏的积分器、符合检测器、高通和带通滤波器、电压—频率转换器、整流器、除法器、振荡器、延迟线和触发器工作。一个给定的细胞成为哪种器件，由它的离子通道以及调节这些通道的广播通道决定。这些模式本身已有测量；大脑在多大程度上利用这种重构进行计算，主要还是假说。
\end{proposition}
\end{keyblock}

对工程师来说，这类似于可编程逻辑阵列：硬件是同一套，器件由配置决定。只不过这里的配置不是在烧录时改变，而是在运行中改变——由第~\ref{sec:serocomputer}--\ref{sec:acet}~章讨论过的慢速广播通道改变。

\section{习题}

\begin{exercise}[\lvlA]
\label{ex:18:1}
\emph{无泄漏积分器的容差。}$\tau=0.1$~s的神经元按~\eqref{eq:perfectint}~在 $T=1$~s内保持眼睛位置，双向误差不超过 $5\%$。(a)~求 $w$ 的允许范围。(b)~$w$ 是 $K$ 个突触权重之和，每个突触有独立的 $10\%$ 误差。$K$ 为多少时和的离散落在容差之内？
\end{exercise}

\begin{exercise}[\lvlB]
\label{ex:18:2}
\emph{作为回路的谐振器。}用变量 $W$ 描述第~\ref{sec:resonator}~节的慢电流：$C\,\dd V/\dd t=-g_LV-g_wW+I$，$\tau_w\,\dd W/\dd t=V-W$。证明这是一个带并联支路 $R_w=1/g_w$、$L_w=\tau_w/g_w$ 的 $RC$ 电路；并在 $C/g_L=10\ms$、$\tau_w=100\ms$、$\alpha=g_w/g_L=4$ 时求谐振频率以及谐振处 $|Z|$ 相对 $|Z(0)|$ 的增益。
\end{exercise}

\begin{exercise}[\lvlB]
\label{ex:18:3}
\emph{积分器如何开始工作。}(a)~神经元 $\tau\,\dd V/\dd t=-V+\mu$，$\tau=10\ms$，阈值 $\theta$，复位到零。$\mu/\theta-1$ 为多少时频率为 $1$~Hz？(b)~模型 $\tau\,\dd v/\dd t=v^2+\mu$（脉冲即 $v$ 趋于 $+\infty$，复位到 $-\infty$）在 $\mu=0.01$ 时给出什么频率？
\end{exercise}

\begin{exercise}[\lvlB]
\label{ex:18:4}
\emph{建模：积分器与谐振器。}习题~\ref{ex:18:2}~的模型取 $g_L=1$、$\tau_m=10\ms$、$\tau_w=100\ms$，阈值 $1$，复位 $V\to0$（$W$ 不变），输入 $\mu=1.5\sin2\pi ft$，$f=1$--$40$~Hz。$\alpha=0$ 和 $\alpha=4$ 时，神经元在哪个频带内输出脉冲？与条件 $1.5\,|Z(f)|>1$ 进行比较。
\end{exercise}

\begin{exercise}[\lvlB]
\label{ex:18:5}
\emph{单细胞锁存器。}带平台电流的神经元：$\tau\,\dd r/\dd t=-r+F(wr+I)$，$F(x)=\min(1,\max(0,x))$，$\tau=10\ms$，$w=1.5$，背景 $I=-0.2$。(a)~幅度 $I=0.3$ 的脉冲至少要持续多久，才能使神经元从 $r=0$ 进入平台？(b)~长脉冲 $I=-0.2-B$ 的幅度 $B$ 至少多大，才能使它回到静息？
\end{exercise}

\begin{exercise}[\lvlB]
\label{ex:18:6}
\emph{积分器的自我调节。}注视时输入 $I=0$，滑动传感器给出漂移速度 $e=\dd r/\dd t$，且 $\dd w/\dd t=-\eta\,e\,r$。在 $w=0.95$、$\tau=0.1$~s、$\langle r^2\rangle=1$ 时，求 $\eta$，使得在 $60$~s的注视内 $|1-w|$ 进入习题~\ref{ex:18:1}~的容差。如果在眼睛转动时也进行学习，会出什么问题？
\end{exercise}

\begin{exercise}[\lvlC]
\label{ex:18:7}
\emph{为什么要重构器件？}广播通道把习题~\ref{ex:18:2}~模型的 $\alpha$ 在 $0$ 和 $4$ 之间切换。对于白噪声电流中的弱正弦信号，在 $11$~Hz和 $1$~Hz处，谐振器的信噪比比积分器高多少倍？设计一个恰好需要切换模式才占优的任务。
\end{exercise}

\chapter{按树突数目划分的神经元}
\chaptermark{树突数目}
\label{sec:dendrites}

\section{输入数目作为电路参数}

在第~\ref{sec:neurontypes}~章中，我们按照输出端释放的物质对神经元分类；在第~\ref{sec:eetypes}~章中，按照它们作为哪种器件工作分类。还有第三个特征，工程师会立刻注意到：\emph{神经元有多少个输入}。输入由树突收集——树突是供其他神经元的轴突附着的分支（第~\ref{sec:dofa}~章，第~\ref{sec:weightstore}~节）。有些神经元根本没有树突，有些有一两个，还有一些有一整棵树，收集十万个输入。对电子工程师来说，这就是输入端的扇入（fan-in），而它几乎总是与任务相称。

树突的数目和大小为什么如此重要，从一个简单的估计就能看出。膜是一个电容器，比电容 $c_m\approx1$~\textmu F/cm$^2$；神经元连同树突的面积 $A$ 越大，要把电压升高 $\Delta V$ 所需送入的电荷就越多。不计泄漏时，输入电流 $I$ 完成这件事所需的时间为
\begin{empheq}[box=\keybox]{equation}
t \;\approx\; \frac{c_m\,A\,\Delta V}{I} .
\label{eq:dendriteload}
\end{empheq}
带几个短树突的小神经元，面积约为 $2\cdot10^{-5}$~cm$^2$，电容约 $20$~pF。一个电流为几纳安的大突触不到十分之一毫秒就能把它升高 $20\mV$。拥有大树突树的神经元面积约大十五倍，电容约 $300$~pF；一个电流约 $0.1$~nA的普通突触需要 $60\ms$ 才能给它充电——远长于泄漏时间常数，也就是说永远充不上。这样的神经元需要几十个同时到达的输入。这里的数字只是数量级，但结论并不依赖于它们：\emph{输入少——快而精确；输入多——慢，靠求和}。

\begin{figure}[t]
\centering
\begin{tikzpicture}[>=Stealth,
  soma/.style={circle,draw,very thick,fill=blue!6,minimum size=0.55cm,inner sep=0pt},
  ax/.style={->,very thick,red!70!black},
  den/.style={very thick,blue!60!black}]
% 0 dendrites: sensor on a line
\begin{scope}[xshift=0cm]
\draw[den,decorate,decoration={snake,amplitude=1pt,segment length=4pt}] (-0.9,0) -- (0,0);
\node[font=\scriptsize] at (-0.9,0.3) {传感器};
\draw[ax] (0,0) -- (1.0,0);
\draw[thick] (0.1,0) -- (0.1,0.45);
\node[soma,minimum size=0.4cm] at (0.1,0.65) {};
\node[font=\scriptsize,align=center] at (0.05,-0.75) {$0$\\带传感器的线路};
\end{scope}
% 1 dendrite
\begin{scope}[xshift=3.3cm]
\draw[den] (-0.9,0) -- (-0.28,0);
\node[soma] at (0,0) {};
\draw[ax] (0.28,0) -- (0.9,0);
\node[font=\scriptsize,align=center] at (0,-0.75) {$1$\\缓冲器，跟随器};
\end{scope}
% 2 dendrites: one per ear
\begin{scope}[xshift=6.2cm]
\draw[den] (-0.28,0) -- (-0.9,0); \draw[den] (0.28,0) -- (0.9,0);
\node[soma] at (0,0) {};
\draw[ax] (0,-0.28) -- (0,-0.6);
\node[font=\scriptsize] at (-0.9,0.25) {左}; \node[font=\scriptsize] at (0.9,0.25) {右};
\node[font=\scriptsize,align=center] at (0,-1.0) {$2$\\时间上的“与”};
\end{scope}
% ~4 short dendrites: mixer
\begin{scope}[xshift=9.0cm]
\foreach \a in {45,135,225,315}{\draw[den] (\a:0.28) -- (\a:0.7); \fill[blue!60!black] (\a:0.72) circle (0.06);}
\node[soma] at (0,0) {};
\draw[ax] (0.28,0) -- (1.0,0);
\node[font=\scriptsize,align=center] at (0,-1.0) {$\sim4$\\混合器};
\end{scope}
% many: summing tree
\begin{scope}[xshift=12.0cm]
\foreach \a in {60,90,120}{\draw[den] (0,0.28) -- ++(\a:0.55) coordinate (b); \foreach \d in {-20,20}{\draw[den] (b) -- ++(\a+\d:0.35);}}
\foreach \a in {-120,-60}{\draw[den] (0,-0.28) -- ++(\a:0.45);}
\node[soma] at (0,0) {};
\draw[ax] (0.28,0) -- (1.0,0);
\node[font=\scriptsize,align=center] at (0,-1.0) {$10^4$--$10^5$\\求和器};
\end{scope}
\end{tikzpicture}
\caption{按输入数目划分的五类。蓝色为树突（输入），红色为轴突（输出）。输入越少，神经元传递单个信号越快越精确；输入越多，它越擅长求和与分类。这是示意图，不是细胞形态图。}
\label{fig:dendriteclasses}
\end{figure}

\section{零个树突：线路末端的传感器}

触觉、痛觉和肌肉牵张的感觉神经元（第~\ref{sec:sensors}、\ref{sec:pain}、\ref{sec:muscles}~章），细胞体上没有一个树突。轴突以一根细枝离开胞体，随即一分为二：一支通往皮肤或肌肉，其末端本身就是传感器；另一支进入脊髓。脉冲从传感器直接跑进脊髓，绕过细胞体。对工程师来说，这是一条传感器位于远端的通信线路，而细胞体挂在一旁，就像电源：它为线路供能，但不参与传输。好处是速度和可靠性：路径上没有任何一处需要对信号求和、再重新决定是否传递。

光和声的传感器是更极端的情况：它们不是收集输入的神经元，而是换能器本身，其输入不是突触，而是受体蛋白（图~\ref{fig:transducer}）。

\section{一个树突：缓冲器}

有一个树突和一个轴突的神经元接收一条输入线路并把它传下去。嗅觉神经元就是这样：唯一的树突伸到鼻腔表面，只携带一种类型的受体\cite{Mombaerts2004}；视网膜中位于光传感器和眼睛输出神经元之间的中继细胞也是这样；连接耳朵传感器和大脑的神经元同样如此。工程师会称它们为缓冲器或跟随器：它们不计算，而是递送一个信号，有时会改变其形式——例如像第~\ref{sec:sensors}~章那样，把传感器的平滑电压变成脉冲。

\section{两个树突：时间上的“与”}

哺乳动物中确定声音方向的符合检测器（图~\ref{fig:itd}）常常恰好有两个树突，朝向相反：一个接收来自左耳的输入，另一个接收来自右耳的输入\cite{Grothe2010}。为什么要把输入分开？回想公式~\eqref{eq:shunt}：当膜的某一处打开了许多兴奋性通道时，那里的电压会接近平衡电位，之后每个输入增加的量越来越少。因此，来自同一只耳朵、汇集在同一个树突上的许多输入会使它饱和，单靠它们无法把神经元推到阈值。而两个树突各来一个输入，则几乎线性相加。两个树突把神经元变成更严格的“与”：只有信号从两侧同时到来时才发放。模型表明，这种分开确实改善了符合检测\cite{AgmonSnir1998}。

\section{几个短树突：混合器与中继器}

\emph{混合器。}在第~\ref{sec:motorlearning}~章的运动学习回路中，约七百亿个小\nt{GLUT}神经元把输入扩展成一个很长的向量。它们每个只有约四个短树突，每个树突末端是一个“爪”，抱住一个输入的末梢。因此，每个混合器只看到众多输入中的约四个，作为其四元组上的一个小“与”元件工作。从 $M$ 条输入线路中可以选出 $\binom{M}{4}$ 个不同的四元组：$M=100$ 时接近四百万。为什么需要这么多？一个有 $n$ 个输入的阈值元件，能正确分为两类的随机模式不超过
\begin{empheq}[box=\keybox]{equation}
P_{\max} \;\approx\; 2n
\label{eq:cover}
\end{empheq}
个\cite{Cover1965}。同一回路的输出神经元从混合器接收约 $10^5$ 个输入，按~\eqref{eq:cover}，它的上界约为 $2\cdot10^5$ 个不同的对应关系。这是理想元件的极限：如果所有权重都为正，如兴奋性突触那样，并且需要留出抗噪余量，那么对同一个神经元的估计约为 $4\cdot10^4$ 个对应关系\cite{Brunel2004}。输入少的混合器，正是为了让大求和器拥有许多不同的、几乎独立的输入\cite{Marr1969,Albus1971}。

\emph{中继器。}在从耳朵通往方向检测器的路径上，有一些只带少数短树突的神经元，它们只接收几个巨大的突触，有些实际上只接收一个。按~\eqref{eq:dendriteload}，这正是所需要的：电容小、电流大，神经元在每个输入脉冲之后不到一毫秒就发放，几乎不损失时间精度。其中一种中继器接收\nt{GLUT}，输出\nt{GLYC}：这是一个精度达几十微秒的反相器，为方向检测器提供快速抑制\cite{BorstSoria2012}。

\section{成千上万个输入：求和器与分类器}

尺度的另一端是皮层中有上万个输入的\nt{GLUT}神经元，以及运动学习回路中有十万个输入的输出\nt{GABA}神经元。按~\eqref{eq:dendriteload}，这样的神经元很慢，不能对单个输入作出响应：它做求和。这就是第~\ref{sec:glutcomputer}~章的积分神经元，也是构建分类器所用的求和器。大树突树还带来新东西：它的分支像各自带阈值的独立子电路一样工作，因此一个神经元的行为就像一个小型多层网络\cite{Beniaguev2021,Gidon2020}。

\section{同时也是输出的树突}

有些神经元根本没有轴突，树突既是输入又是输出：它们接收信号，自己也释放神经递质。研究最多的例子是视网膜中参与确定运动方向的\nt{GABA}神经元（第~\ref{sec:gaba}~章，图~\ref{fig:direction}）。这种神经元的每个分支各自对从细胞体向其末梢方向的运动作出响应，并从这个末梢输出抑制\cite{Euler2002}。一个神经元就是几个独立的方向检测器，每个分支一个。对工程师来说，这不是一个元件，而是一个封装内有多个通道的芯片。

\section{小结}

\begin{table}[t]
\centering
\footnotesize
\setlength{\tabcolsep}{3pt}
\renewcommand{\arraystretch}{1.15}
\begin{tabular}{>{\raggedright}p{1.9cm}>{\raggedright}p{3.4cm}>{\raggedright}p{2.6cm}>{\raggedright}p{1.8cm}>{\raggedright\arraybackslash}p{3.8cm}}
\toprule
树突数 & 例子 & 器件 & 输入数 & 已知什么 \\
\midrule
$0$ & 触觉、痛觉、牵张传感器 & 末端带传感器的线路 & --- & 已确定 \\
$1$ & 嗅觉神经元、视网膜中继细胞、耳的神经元 & 缓冲器，跟随器 & $1$ 到少数 & 已确定 \\
$2$ & 声音方向检测器 & 时间上的“与” & 两束 & 结构已确定；分开输入的好处已在模型中证明 \\
$\sim4$ & 运动学习回路的混合器 & 小“与”元件 & $\sim4$ & 已确定；扩展输入的作用是有充分依据的假说 \\
少数，短 & 耳通路上的中继器 & 中继器，反相器 & $1$ 到几个巨大突触 & 已确定 \\
成千上万 & 皮层\nt{GLUT}神经元，运动学习的输出神经元 & 求和器，分类器 & $10^4$--$10^5$ & 已确定；分支作为子电路已在个别例子中测量 \\
树突即输出 & 视网膜方向\nt{GABA}神经元 & 多通道芯片 & 许多 & 已测量 \\
\bottomrule
\end{tabular}
\caption{按树突数目划分的神经元。}
\label{tab:dendrites}
\end{table}

\begin{keyblock}
\begin{proposition}[输入数目跟随任务]
\label{prop:dendrites}
凡是需要单个信号的速度和精度之处——在传感器、中继器和符合检测器上——神经元的树突少、输入少、突触大：小电容~\eqref{eq:dendriteload}~充电快。凡是需要求和与分类之处，神经元都有带成千上万个输入的大树突树。各类的结构已有测量；计算上的解释有充分依据，但对许多类别来说仍是假说。
\end{proposition}
\end{keyblock}

表~\ref{tab:dendrites}~补充了前两种分类：按神经递质（表~\ref{tab:types}）和按器件（表~\ref{tab:eetypes}）。三种分类都是从不同侧面观察同一个元件：它输出什么，它计算什么，以及它听多少。

\section{习题}

\begin{exercise}[\lvlA]
\label{ex:19:1}
\emph{大神经元需要多少输入。}一个 $C=300$~pF、$\tau_m=20\ms$ 的神经元接收突触，每个突触是 $0.1$~nA的电流源，阈值比静息高 $20\mV$。(a)~要能到达阈值，至少需要多少个突触同时工作？在 $10\ms$ 内呢？$5\ms$ 内呢？(b)~$C=20$~pF的神经元在一个 $4$~nA突触作用下多久到达阈值？
\end{exercise}

\begin{exercise}[\lvlA]
\label{ex:19:2}
\emph{为什么是四个。}混合器是对 $M=100$ 条线路中 $k$ 条的“与”，每个模式中一半线路活跃，混合器共 $7\cdot10^{10}$ 个。(a)~$k=2$、$4$、$40$ 时各有多少个混合器对一个模式响应？(b)~与 $M=100$ 条直接线路相比，混合器使有 $10^5$ 个输入的输出神经元的对应关系数~\eqref{eq:cover}~增加了多少倍？
\end{exercise}

\begin{exercise}[\lvlB]
\label{ex:19:3}
\emph{$2n$ 从何而来。}过原点的超平面把 $n$ 维空间中处于一般位置的 $P$ 个点分成两类，共有 $C(P,n)=2\sum_{k=0}^{n-1}\binom{P-1}{k}$ 种方式。(a)~证明当 $P=2n$ 时，$2^P$ 种划分中恰好一半可分。(b)~$n=100$、$P=180$ 和 $220$ 时，可分的比例是多少？
\end{exercise}

\begin{exercise}[\lvlB]
\label{ex:19:4}
\emph{两个树突作为严格的“与”。}树突是泄漏电导为 $g_L$ 的一段，每个输入在其上打开电导 $g_0=0.5\,g_L$，反转电位为 $E_E$，胞体看到 $\kappa(V_1+V_2)$，阈值 $\theta=1.1\,\kappa E_E$。为什么同一个树突上无论多少输入都无法触发神经元？每个树突上需要多少输入？
\end{exercise}

\begin{exercise}[\lvlB]
\label{ex:19:5}
\emph{设计中继器。}脉冲抖动 $\sigma_t=\sigma_V/\dot V$，其中 $\dot V$ 是阈值附近的电压斜率。要求在噪声 $\sigma_V=1\mV$ 下 $\sigma_t\le20$~\textmu s；忽略泄漏。$C=300$~pF的神经元为此需要多少个 $0.1$~nA的同步突触？$C=20$~pF、只有一个 $4$~nA突触的神经元抖动是多少？
\end{exercise}

\begin{exercise}[\lvlB]
\label{ex:19:6}
\emph{建模：长树突的充电。}模拟长度 $L=\ell/\lambda$、两端封闭的无源电缆 $\tau_m\partial_tV=\lambda^2\partial_x^2V-V+i/g$（$40$ 段，欧拉法）。在远端注入短电流脉冲：$L=0.2$、$1$、$2$ 时，近端何时升到、升到多高（以同样电荷均摊到整个树突时电压的比例计）？
\end{exercise}

\begin{exercise}[\lvlC]
\label{ex:19:7}
\emph{研究：最优输入数目。}电容 $C=C_0+nc_s$，同时有 $m$ 个输入工作，每个电流 $I_s$，神经元记忆 $P\approx2n$ 个对应关系~\eqref{eq:cover}。在 $m$ 恒定和比例 $m=\varphi n$ 恒定时，响应时间如何依赖于 $P$？提出一个代价准则，并分别求中继器和分类器的最优 $n$。
\end{exercise}

\chapter{神经元在睡眠中做什么}
\chaptermark{睡眠}
\label{sec:sleep}

\section{睡眠不是关机，而是另一种模式}

工程师看到一台关掉的计算机，会说：它什么也没做。对沉睡的大脑却不能这么说。神经元在睡眠中并不沉默：在深睡眠中，皮层照样发放脉冲；在快速眼动睡眠中，发放几乎和白天一样多。改变的不是神经元\emph{是否工作}，而是它们\emph{如何}工作。睡眠是这台机器的一组模式，切换这些模式的正是白天也在工作的那些广播通道。

对每种模式，都可以写出“寄存器状态”（表~\ref{tab:sleepmodes}）。白天 \nt{ACET}、\nt{NORA} 和 \nt{SERO} 都在工作，传感器接通，肌肉听从指挥。在慢波深睡眠中，这三个通道都安静下来，来自感官的输入减弱\cite{Hasselmo1999}：皮层中乙酰胆碱的释放下降，\nt{NORA} 和 \nt{SERO} 神经元发放脉冲比白天少\cite{Marrosu1995,AstonJonesBloom1981,McGintyHarper1976}。快速眼动睡眠中的图景很奇特：\nt{ACET} 重新升到白天的水平\cite{Marrosu1995}，而 \nt{NORA} 和 \nt{SERO} 几乎完全沉默\cite{AstonJonesBloom1981,McGintyHarper1976,JacobsAzmitia1992}。快速眼动睡眠中身体的肌肉是关闭的：控制它们的 \nt{ACET} 神经元经 \nt{GABA} 和 \nt{GLYC} 受到强烈抑制\cite{BrooksPeever2012}——这正是第~\ref{sec:gaba}~章中的禁止，只不过一次施加在整个输出上。

\begin{table}[t]
\centering
\footnotesize
\setlength{\tabcolsep}{4pt}
\renewcommand{\arraystretch}{1.15}
\begin{tabular}{>{\raggedright}p{3.1cm}>{\raggedright}p{2.9cm}>{\raggedright}p{3.2cm}>{\raggedright\arraybackslash}p{3.4cm}}
\toprule
& 清醒 & 慢波睡眠 & 快速眼动睡眠 \\
\midrule
\nt{ACET}（写入，第~\ref{sec:acet}~章） & 高 & 低 & 高 \\
\nt{NORA}（增益，第~\ref{sec:nora}~章） & 中等，有爆发 & 低 & 几乎沉默 \\
\nt{SERO}（第~\ref{sec:serocomputer}~章） & 中等 & 低 & 几乎沉默 \\
来自传感器的输入 & 接通 & 减弱 & 减弱 \\
到肌肉的输出 & 接通 & 减弱 & 被抑制关闭 \\
皮层活动 & 平稳 & 慢摆动：工作—静默 & 平稳，如同白天 \\
\bottomrule
\end{tabular}
\caption{这台机器在三种模式下的寄存器状态。所有各行均已测量；\nt{ACET}、\nt{NORA} 和 \nt{SERO} 的水平来自按睡眠阶段所作的直接记录\cite{Marrosu1995,AstonJonesBloom1981,McGintyHarper1976}。}
\label{tab:sleepmodes}
\end{table}

\section{何时入睡：带迟滞的继电器}

是什么决定机器何时切换到睡眠？一个由两部分组成的简单方案效果很好\cite{Borbely1982}。第一部分是\emph{睡眠压力} $S$，它在我们清醒时积累，在睡眠时释放。这是一个白天充电、夜里放电的电容器：
\begin{empheq}[box=\keybox]{equation}
\frac{\dd S}{\dd t} \;=\; \frac{1-S}{\tau_r}\ \ \text{（清醒）},\qquad
\frac{\dd S}{\dd t} \;=\; -\frac{S}{\tau_d}\ \ \text{（睡眠）} .
\label{eq:sleeppressure}
\end{empheq}
第二部分是两个阈值：当 $S$ 达到上阈值 $H(t)$ 时机器入睡，当 $S$ 降到下阈值 $L(t)$ 时醒来。两个阈值都随第~\ref{sec:chemoutput}~章的昼夜节律平缓摆动。这就得到一个带迟滞的继电器——第~\ref{sec:bistable}~节的施密特触发器——它与电容器一起构成一个弛张振荡器，由昼夜节律~\eqref{eq:pll}~进行相位校准。

\begin{figure}[t]
\centering
\begin{tikzpicture}[>=Stealth,x=0.16cm,y=4.2cm]
\foreach \a/\b in {0/4.3,21.2/29.3,45.4/53.4,69.5/72}{\fill[blue!8] (\a,0) rectangle (\b,1);}
\draw[->,gray!70!black] (0,0) -- (75,0) node[right,font=\small,black] {$t$，h};
\draw[->,gray!70!black] (0,0) -- (0,1.08) node[left,font=\small,black] {$S$};
\foreach \t in {0,24,48,72}{\draw[gray!60!black] (\t,-0.02) -- (\t,0.02); \node[below,font=\scriptsize] at (\t,0) {\t};}
\draw[thick,densely dashed,red!70!black] plot file {figures/sleep_H.dat};
\draw[thick,densely dashed,green!45!black] plot file {figures/sleep_L.dat};
\draw[very thick,blue!60!black] plot file {figures/sleep_S.dat};
\node[font=\scriptsize,red!70!black,anchor=west] at (73,0.72) {$H$：入睡};
\node[font=\scriptsize,green!45!black,anchor=west] at (73,0.24) {$L$：醒来};
\node[font=\scriptsize,blue!60!black] at (25.25,0.93) {睡眠};
\node[font=\scriptsize,blue!60!black] at (49.4,0.93) {睡眠};
\end{tikzpicture}
\caption{$\tau_r=18.2$~h、$\tau_d=4.2$~h 时的睡眠压力 $S$~\eqref{eq:sleeppressure}。白天 $S$ 充电到上阈值 $H$，夜里（阴影部分）放电到下阈值 $L$；两个阈值都随昼夜节律摆动。机器自己进入大约 $16$ 小时清醒、$8$ 小时睡眠的模式。}
\label{fig:twoprocess}
\end{figure}

在我们的模型中，取常用值 $\tau_r=18.2$~h 和 $\tau_d=4.2$~h，机器会自己进入 $16.1$ 小时清醒、$8.0$ 小时睡眠的状态（图~\ref{fig:twoprocess}）。模型还解释了一件看似奇怪的事：四十小时不睡之后，压力达到 $0.90$，但模型中的恢复性睡眠只持续 $8.8$ 小时，而不是比平时多一整天。放电是指数式的，高电荷很快就放掉；“欠债”靠睡眠的深度而不是长度来偿还。

物理上扮演 $S$ 角色的是什么？一个有力的候选者是腺苷，即附录~\ref{sec:othermediators}~中的“负荷计数器”：它的浓度在清醒期间上升，在睡眠中下降\cite{PorkkaHeiskanen1997,Dunwiddie2001}。睡眠压力恰恰就是腺苷而不是别的东西，这是假说。

\section{慢波睡眠：整个网络成为弛张振荡器}

在深睡眠中，皮层神经元以整个群体为单位轮流工作：每秒不到一次，它们一起开启几百毫秒（\emph{工作状态}），又一起沉默（\emph{静默状态}）：这种摆动的频率低于 $1$~Hz，在麻醉下最常见的是 $0.3$--$0.4$~Hz\cite{Steriade1993}。这种慢摆动可以用我们已有的部件组装出来。取一群自身反馈很强的 \nt{GLUT} 神经元——即第~\ref{sec:glutcomputer}~章中有两个稳定状态的记忆——再加上慢疲劳，就像第~\ref{sec:muscles}~章的节律发生器那样：
\begin{empheq}[box=\keybox]{equation}
\tau\,\frac{\dd r}{\dd t} = -r + F\bigl(w\,r + I - a\bigr),
\qquad
\tau_a\,\frac{\dd a}{\dd t} = -a + b\,r .
\label{eq:updown}
\end{empheq}
群体开启、工作、疲劳；疲劳 $a$ 吃掉上状态，群体落入静默；在静默中疲劳消退，下状态消失，群体再次开启。这就得到与睡眠压力相同的弛张振荡器，只是快了大约八万倍。

\begin{figure}[t]
\centering
\begin{tikzpicture}[>=Stealth,x=2.9cm,y=1.0cm]
\begin{scope}[yshift=1.9cm]
\draw[->,gray!70!black] (0,0) -- (4.1,0);
\draw[->,gray!70!black] (0,0) -- (0,1.25) node[left,font=\small,black] {$r$};
\draw[thick,blue!60!black] plot file {figures/updown_sleep.dat};
\node[font=\scriptsize,anchor=west] at (4.15,0.5) {睡眠：$b=5$};
\end{scope}
\draw[->,gray!70!black] (0,0) -- (4.1,0) node[right,font=\small,black] {$t$，s};
\draw[->,gray!70!black] (0,0) -- (0,1.25) node[left,font=\small,black] {$r$};
\draw[thick,red!70!black] plot file {figures/updown_wake.dat};
\node[font=\scriptsize,anchor=west] at (4.15,0.5) {白天：$b=1$};
\foreach \t in {0,1,2,3,4}{\draw[gray!60!black] (\t,-0.05) -- (\t,0.05); \node[below,font=\scriptsize] at (\t,0) {\t};}
\end{tikzpicture}
\caption{同一个群体~\eqref{eq:updown}~在两种模式下：$w=5$，$I=2$，$\theta=2.5$，$\tau=10\ms$，$\tau_a=0.3$~s。疲劳强时（$b=5$，如同慢波睡眠），群体在工作和静默之间摆动，周期为 $1.1$~s（模型值；皮层中周期超过一秒，麻醉下约为 $3$~s）。减弱疲劳（$b=1$）——乙酰胆碱就是这样起作用的——同一个群体便平稳地工作。}
\label{fig:updown}
\end{figure}

那么，是什么把网络从摆动转为白天的平稳工作？乙酰胆碱抑制神经元中产生疲劳的慢电流\cite{McCormick1992}。在我们的模型中，疲劳强时（$b=5$），群体以 $1.1$~s 的周期摆动——比实测的摆动快，但数量级相同——若按整数个周期平均，约有 $35\%$ 的时间在工作；疲劳弱时（$b=1$），同一个群体平稳地工作（图~\ref{fig:updown}）。一个寄存器——\nt{ACET} 的水平——就把振荡器切换成了放大器。在睡眠中，慢摆动之上还叠加着每秒 $7$--$15$ 次的较快节律爆发；它们由皮层和一个中间中继核团之间的 \nt{GLUT}--\nt{GABA} 环路产生\cite{SteriadeMcCormickSejnowski1993}，与第~\ref{sec:gaba}~章的节律是近亲。

\section{重放：快进播放白天}

在第~\ref{sec:acet}~章中我们看到，在慢波睡眠中，白天一起工作的神经元会再次一起发放，并且顺序相同。再补充一个工程细节：重放是\emph{快进}的。白天占几秒的序列，在睡眠中只用十分之几秒就播放完\cite{Nadasdy1999}：在海马中大约快二十倍\cite{LeeWilson2002}，在皮层中快六七倍\cite{Euston2007}。而且重放并不是随时发生的：一个区域中短暂的重放爆发会与皮层中工作—静默的摆动以及快速的节律爆发对齐。如果人为增强这种协调，记忆就会改善\cite{Maingret2016}——这已经是因果关系，而不是巧合。

机器为什么要快进播放白天？工程师熟悉一种称为\emph{灾难性遗忘}的困难：一个网络如果一件接一件地连续学习新东西，就会破坏旧的东西。解决办法是\emph{交错}：把新东西和旧东西混在一起，小批量地多次学习。互补学习系统假说\cite{McClelland1995}认为，快速记忆把白天整体记录下来，然后在睡眠中一点一点地、交错地、多次地向慢速的皮层重放。这正是第~\ref{sec:acet}~章中“快速缓冲区—慢速存储器”的组合，也是第~\ref{sec:motorlearning}~章中快慢两个学习者的组合。重放及其协调已有测量；其目的是对慢速记忆进行交错训练，这是有充分依据的假说。

\section{权重的重新归一化}

白天，第~\ref{sec:dofa}~章的赫布规则主要是增强连接。如果只增强，权重就会增长，能耗也随之增长，而灵敏度下降：所有输入都很强的神经元将无法区分它们。按照\emph{突触稳态}假说\cite{TononiCirelli2014}，睡眠的用途之一就是把权重恢复到工作水平：所有连接按同一比例减弱，因此它们的\emph{比值}——也就是学到的东西——得以保留。对工程师来说，这就是归一化，就像规则~\eqref{eq:oja}~那样，只不过不是在运行中进行，而是在单独的时间段进行。

直接测量与此并不矛盾：在小鼠中，睡眠后皮层突触的接触面积比清醒后小约 $18\%$，减小量与大小成正比，而最大、最稳定的接触几乎不变\cite{deVivo2017}。权重的缩放已有测量；这是睡眠的主要任务，则是假说。

\section{快速眼动睡眠：与输入输出断开的机器}

在快速眼动睡眠中，皮层的工作几乎和白天一样，但来自感官的输入减弱，到肌肉的输出被抑制关闭。对工程师来说，这是一个熟悉的模式：\emph{台架测试}。电路满负荷运行，但它的输入接到内部信号发生器上，输出与执行机构断开，以免弄坏任何东西。按照一种假说，梦正是对白天建立起来的世界模型进行的这种内部试运行；网络此时究竟在做什么、是否在学习，尚不清楚。这里测量到的只有表~\ref{tab:sleepmodes}~中记录的内容：哪个通道开着，哪个沉默，以及输出被封锁。

\section{维护与局部睡眠}

长期以来人们认为，睡眠时细胞间隙扩大，大脑能更快地冲洗掉代谢产物\cite{Xie2013}。最近用另一种方法测量冲洗的实验得到了相反的结果：睡眠时冲洗更慢\cite{Miao2024}。这个问题目前悬而未决，我们也让它保持开放。

另一个事实更可靠：睡眠是局部网络的属性，而不是整台机器同时的属性。如果长时间不让动物睡觉，那么在动物清醒并活动时，它皮层的个别区域会短暂地陷入静默，就像在慢波睡眠中一样，而在这些时刻动物更容易出错\cite{Vyazovskiy2011}。每个网络自己疲劳，自己切换；当广播通道把所有网络同时切换时，就形成了全局模式。

\begin{keyblock}
\begin{proposition}[睡眠是一组模式]
\label{prop:sleep}
睡眠中神经元并不关闭：广播通道把机器切换到别的模式。何时睡眠，由经昼夜节律校准的带迟滞继电器~\eqref{eq:sleeppressure}~决定。在慢波睡眠中，网络成为弛张振荡器~\eqref{eq:updown}，并快进重放白天；在快速眼动睡眠中，它在与输入和输出断开的状态下工作。模式、重放和权重缩放都已测量；它们的目的——交错学习和重新归一化——是有充分依据的假说。
\end{proposition}
\end{keyblock}

\section{小结}

\begin{table}[t]
\centering
\footnotesize
\setlength{\tabcolsep}{4pt}
\renewcommand{\arraystretch}{1.15}
\begin{tabular}{>{\raggedright}p{3.2cm}>{\raggedright}p{4.4cm}>{\raggedright\arraybackslash}p{5.2cm}}
\toprule
发生了什么 & 怎样实现 & 已知什么 \\
\midrule
模式切换 & 寄存器 \nt{ACET}、\nt{NORA}、\nt{SERO}（表~\ref{tab:sleepmodes}） & 已测量 \\
何时睡眠 & 电容器 $S$ 和带迟滞的继电器~\eqref{eq:sleeppressure} & 模型很好地描述了实验；腺苷作为 $S$ 是假说 \\
慢摆动 & 带反馈和疲劳的群体~\eqref{eq:updown} & 摆动已测量；模型是最简单的电路 \\
快进播放白天 & 快 $6$--$20$ 倍、与摆动协调的重放 & 已测量；增强协调可改善记忆 \\
为什么要重放 & 对慢速记忆的交错训练 & 有充分依据的假说 \\
权重重新归一化 & 睡眠中连接的缩放 & 接触减小已测量；作为睡眠的主要作用是假说 \\
快速眼动睡眠 & 在输入和输出断开时工作 & 模式已测量；目的未知 \\
冲洗 & 细胞间隙扩大 & 结果相互矛盾 \\
局部睡眠 & 个别区域陷入静默 & 已测量 \\
\bottomrule
\end{tabular}
\caption{神经元在睡眠中做什么。}
\label{tab:sleep}
\end{table}

表~\ref{tab:sleep}~汇总了本章。睡眠原来不是机器工作中的暂停，而是它在运行中的维护：由同样的广播通道切换模式，由与步行节律相同的部件构成振荡器，为慢速记忆快进播放白天，以及重新归一化权重。白天机器写入，夜里它重写并整理写下的东西。

\section{习题}

\begin{exercise}[\lvlA]
\label{ex:20:1}
\emph{没有昼夜节律的继电器。}设阈值恒定：$H=0.67$，$L=0.17$（图~\ref{fig:twoprocess}~模型阈值的平均值）。由~\eqref{eq:sleeppressure}~推出 $\tau_r=18.2$~h、$\tau_d=4.2$~h 时清醒和睡眠的时长以及振荡器的周期。为什么阈值摆动的模型仍然达到 $16.1+8.0=24$~h？这一效应属于第~\ref{sec:chemoutput}~章的哪种器件？
\end{exercise}

\begin{exercise}[\lvlA]
\label{ex:20:2}
\emph{睡眠债。}在压力 $S_0$ 下开始的睡眠持续 $t_s=\tau_d\ln(S_0/L)$，$\tau_d=4.2$~h，$L$ 是醒来时的下阈值。(a)~$40$~h 不睡后 $S_0=0.90$，睡眠持续了 $8.8$~h。$L$ 是多少？在通常的 $L\approx0.092$ 下睡眠会持续多久？(b)~一夜之间突触接触面积减少 $18\%$。白天权重应增长多少？
\end{exercise}

\begin{exercise}[\lvlB]
\label{ex:20:3}
\emph{设计阈值。}选取恒定阈值 $H$ 和 $L$，使 $\tau_r=18.2$~h、$\tau_d=4.2$~h 的继电器~\eqref{eq:sleeppressure}~恰好给出 $16$~h 清醒和 $8$~h 睡眠。如果某天夜里在 $4$ 小时后被叫醒，这样的机器比阈值摆动的机器差在哪里？
\end{exercise}

\begin{exercise}[\lvlB]
\label{ex:20:4}
\emph{慢摆动的周期。}在模型~\eqref{eq:updown}~中，$F(x)=1/\bigl(1+e^{-(x-\theta)/k}\bigr)$，$w=5$，$I=2$，$\theta=2.5$，$k=0.25$，$b=5$，$\tau\ll\tau_a=0.3$~s。(a)~求曲线 $r=F(wr+I-a)$ 的拐点 $a_{\text{上}}$ 和 $a_{\text{下}}$，工作状态和静默状态分别在这里消失。(b)~取工作时 $r\approx1$、静默时 $r\approx0$，求周期和工作时间比例。
\end{exercise}

\begin{exercise}[\lvlB]
\label{ex:20:5}
\emph{摆动何时关闭。}在习题~\ref{ex:20:4}~的模型中，不动点位于曲线 $a=wr+I-F^{-1}(r)$ 与直线 $a=br$ 的交点。只要交点位于中间分支、在两个拐点之间，振荡就会持续。$b$ 为何值时如此？\nt{ACET} 如何通过改变 $b$ 把振荡器切换成放大器？
\end{exercise}

\begin{exercise}[\lvlB]
\label{ex:20:6}
\emph{建模：欠债还是相位。}在 \texttt{sleep\_models.py} 中，取 $48$~h 后的第一次醒来，让机器保持清醒 $D=24$、$32$、$40$、$48$、$64$~h（\texttt{stay\_awake}，\texttt{days=8}）。每个 $D$ 下恢复性睡眠持续多久？是什么决定了它的时长？
\end{exercise}

\begin{exercise}[\lvlB]
\label{ex:20:7}
\emph{建模：遗忘与交错。}线性神经元 $y=\mathbf w\cdot\mathbf x$，$n=40$，按规则 $\mathbf w\leftarrow\mathbf w-0.5\,(y-y^*)\,\mathbf x$ 学习。任务 A 和 B 各有 $10$ 个模式，$x_j\sim\mathcal N(0,1/n)$，$y^*=\pm1$。先学 $200$ 轮 A、再学 $200$ 轮 B 之后，A 上的均方误差是多少？A 和 B 交错学习 $200$ 轮之后呢？
\end{exercise}

\begin{exercise}[\lvlC]
\label{ex:20:8}
\emph{研究：交错还是重新归一化？}习题~\ref{ex:20:7}~的神经元加上阈值；两种夜间程序：(i)~所有权重乘以 $0.82$；(ii)~旧模式与新模式交错重复。每种程序对权重比值和神经元的回答有什么影响？如何根据睡眠后突触的大小区分这两种假说？
\end{exercise}

\chapter{用活神经元搭成的机器}
\chaptermark{用活神经元搭成的机器}
\label{sec:livingmachine}
\begin{chaptergoals}
\item 为什么没有广播通道的培养物会产生群放电，是什么决定其间隔；
\item 活的网络如何在反馈下改变响应，哪些仍未证实；
\item 类器官如何充当储备池，其线性读出在外部的硅片上训练；
\item 迄今培养物是怎样被教的，为什么老师仍站在培养皿之外。
\end{chaptergoals}

\section{电路图公开的试验台}

在第~\ref{sec:debugging}~章里，我们调试的是一台没有电路图的机器：探针在八百亿个神经元中只能听到一千个，其余全靠猜。处在这种境地的工程师会做一件简单的事：在面包板上搭一小块电路，每个元件都看得见。对神经元也可以这样做。

把皮层神经元分散开，培养在带电极阵列的芯片上，电极从几十个到几万个不等。几周后，细胞长出突起，彼此连接成由几千到几十万个神经元组成的网络，其中大部分是\nt{GLUT}，较小的一部分是\nt{GABA}，这一比例取决于标本：在海马培养物中约占神经元的 $6\%$~\cite{Benson1994}。每个电极既能像第~\ref{sec:debugging}~章的探针那样监听，也能像测试信号那样注入电流。试验台的第二种形式是\emph{类器官}：由人类干细胞培养出的三维神经组织小团，大小约一毫米~\cite{Lancaster2013}。在这样的试验台上，活的网络可以工作好几个月；有的平台允许通过网络远程在类器官上做实验，连续一百多天~\cite{Jordan2024}。这一领域被称为\emph{生物计算}或\emph{类器官智能}~\cite{Smirnova2023}。

试验台与大脑有两个重要区别。它极小：神经元数量少十万倍。而且它没有广播通道：皮层培养物中没有\nt{DOPA}、\nt{SERO}、\nt{NORA}和\nt{ACET}神经元。这是一台有数据总线和流量控制、却没有控制寄存器的机器。我们来看看它能做什么。

\section{网络自己在做什么}

放着不管的培养物既不沉默，也不均匀地产生噪声。每隔一段时间，几乎整个网络会同时爆发一次\emph{群放电}，然后又安静下来；群放电之间的间隔从一秒到几分钟不等，因培养物而异~\cite{Wagenaar2006}。对工程师来说这是熟悉的画面：带强正反馈、又没有负载的放大器会变成振荡器。

机制我们已经知道。\nt{GLUT}神经元之间强烈的相互连接会引发雪崩，就像第~\ref{sec:gaba}~章中命题~\ref{prop:ei} 的条件被破坏时那样。雪崩迅速耗尽突触中的神经递质储备~\eqref{eq:depression}，网络随之沉寂。储备以时间常数 $\tau_{\text{rec}}$ 恢复，当恢复到足以引发新雪崩的比例 $x^*$ 时，下一次群放电到来。群放电之间的间隔为
\begin{empheq}[box=\keybox]{equation}
T_{\text{burst}} \;\approx\; \tau_{\text{rec}}\,\ln\frac{1}{1-x^*} .
\label{eq:burstinterval}
\end{empheq}
取第~\ref{sec:code}~章的 $\tau_{\text{rec}}=0.8\,\text{s}$ 和 $x^*=0.9$，约为两秒；$x^*=0.99$ 时约为四秒。这是用本书 $\tau_{\text{rec}}$ 得到的模型估计，落在实测间隔的短端。在微培养物上的直接测量表明，群放电后递质储备要几十秒才能恢复~\cite{CohenSegal2011}，也就是说那里的 $\tau_{\text{rec}}$ 更大；用这样的 $\tau_{\text{rec}}$，公式~\eqref{eq:burstinterval} 给出几十秒乃至几分钟，与慢速培养物一致。这是一个弛张振荡器，与第~\ref{sec:muscles}~章的节律发生器同属一类：那里给它“充电”的是肌群的疲劳，这里是递质储备。

群放电是网络无事可做时所做的事。在活的大脑中，类似的图景出现在深度睡眠中（第~\ref{sec:sleep}~章），也出现在真正的输入尚未接入的发育期大脑中。这种相似性发人深思，但两者机制相同只是一个假说。

\section{没有多巴胺老师，怎样教网络}

培养物中没有\nt{DOPA}，“更好还是更差”的信号只能由我们自己通过电极给出。实验表明，试验台上的网络确实会改变响应，而其中最有意思的是用\emph{什么}来教它。

\emph{会闭嘴的老师。}通过一个电极每一到三秒刺激一次网络，直到另一个电极在刺激后 $50\pm10\ms$ 出现所需的响应。一旦出现，就关掉刺激五分钟。经过几个这样的循环后，所需响应在刺激后立刻出现~\cite{Shahaf2001}。这里的奖励是安静：刺激会扰动网络，而停止刺激则让网络停留在给出正确响应的那个状态。这就是第~\ref{sec:dofa}~章中通过随机偏离做梯度下降（命题~\ref{prop:gradient}），其中扮演误差角色的是“被扰动”这件事本身：一直扰动，直到对为止。

\emph{会打网球的网络。}在一项使用高密度电极阵列试验台的实验中，约一百万个神经元被接入一个单球拍的电脑网球游戏~\cite{Kagan2022}。球的位置以电流形式送到“感觉”电极上：球在哪里用位置编码，距离多远用频率编码。两个“运动”区的活动让球拍上下移动。学习规则很不寻常：如果球拍击中了球，网络得到一个短暂的\emph{可预测}刺激；如果没击中，则得到几秒钟\emph{不可预测}的随机电流。二十分钟的游戏中，连续击球的回合变长了。只有在完整反馈下，一次会话内才有显著提升；把刺激换成安静时，培养物到会话结束时也比没有反馈时打得好，但效果较小；而在几天里跨会话的稳定学习并未出现。对这项实验的详细分析见第~\ref{sec:dishbrain}~章。

作者用一个假说来解释这一结果：网络的行为是为了让自己的输入变得可预测~\cite{Friston2010}。这与第~\ref{sec:code}~章中节省脉冲、第~\ref{sec:recognition}~章中预测画面是同一个思想：输入在意料之中时，机器花费更少。但这项实验本身，尤其是作者关于培养物“有感知能力”的说法，招致了尖锐的批评：效应不大，而且可以用更简单的方式解释~\cite{Balci2023}。诚实的评价是：试验台上的培养物在反馈作用下改变行为，这是已测量的；至于它学的恰恰是预测自己的输入，则是假说。

\emph{会驾驶身体的网络。}以类似方式，人们把培养物接到一个简单的虚拟“身体”上，通过挑选训练刺激的时空模式，教它朝目标移动~\cite{Bakkum2008}。这些实验中都没有多巴胺：充当老师的是工程师选定的电流模式。当老师换成程序或多巴胺本身时会有什么变化，见第~\ref{sec:dishdopamine}--\ref{sec:cartpole}~节。

\begin{figure}[t]
\centering
\begin{tikzpicture}[>=Stealth,
  blk/.style={draw,very thick,rounded corners=2pt,align=center,font=\footnotesize,minimum height=0.8cm,fill=blue!5}]
% (a) closed loop
\node[font=\small] at (1.5,4.3) {(a) 闭环};
\node[blk,text width=2.6cm] (G) at (1.5,3.3) {游戏：球和球拍};
\draw[very thick,fill=orange!10,rounded corners=3pt] (0.5,0.2) rectangle (2.5,1.6);
\foreach \x in {0.7,1.0,...,2.35}{\foreach \y in {0.4,0.7,1.0,1.3}{\fill[gray!60] (\x,\y) circle (0.04);}}
\draw[->,thick,blue!60!black] (G.west) -| (-0.5,0.9) -- (0.5,0.9);
\node[font=\scriptsize,blue!60!black,anchor=east] at (-0.6,2.1) {球 $\to$ 电流};
\draw[->,thick,red!70!black] (2.5,0.9) -- (3.5,0.9) |- (G.east);
\node[font=\scriptsize,red!70!black,anchor=west] at (3.6,2.1) {脉冲 $\to$ 球拍};
\node[font=\scriptsize] at (1.5,-0.2) {电极阵列上的神经元};
\node[font=\scriptsize,align=center] at (1.5,-0.85) {击中：可预测刺激\\未击中：随机电流};
% (b) reservoir
\begin{scope}[xshift=7.9cm]
\node[font=\small] at (1.6,4.3) {(b) 储备池};
\node[blk,text width=1.1cm] (X) at (-0.4,1.0) {输入\\$u(t)$};
\draw[very thick,fill=orange!10] (1.6,1.0) circle (0.8);
\foreach \a/\r in {20/0.35,80/0.5,150/0.3,210/0.55,280/0.4,330/0.2,110/0.15}{\fill[red!60!black] ({1.6+\r*cos(\a)},{1.0+\r*sin(\a)}) circle (0.05);}
\node[font=\scriptsize] at (1.6,-0.2) {类器官：无监督};
\node[blk,text width=1.8cm,fill=green!8] (W) at (4.0,1.0) {读出\\$W_{\text{out}}$};
\draw[->,thick] (X) -- (0.8,1.0);
\draw[->,thick] (2.4,1.0) -- node[above,font=\scriptsize] {$\mathbf r(t)$} (W);
\draw[->,thick] (W) -- (4.0,2.3) node[above,font=\scriptsize] {$y(t)$};
\node[font=\scriptsize,green!40!black] at (4.0,0.3) {有监督训练};
\end{scope}
\end{tikzpicture}
\caption{让活网络做计算的两种方式。(a) 闭环：球的位置以电流输入，运动区的脉冲移动球拍，击球后的可预测刺激或随机刺激充当老师。(b) 储备池~\eqref{eq:reservoir}：不给类器官任何教学信号，它把输入分解成许多带记忆的非线性响应；按误差学习的只有外部的一个线性读出。类器官自身的连接同时也在变化，但是无监督的。}
\label{fig:livingmachine}
\end{figure}

\section{活的储备池}

使用活网络还有另一种方式：根本不去教它。在工程上这称为\emph{储备池计算}~\cite{Maass2002,JaegerHaas2004}。取一个现成的、带记忆的非线性系统，什么系统都行，只要它对输入的响应足够丰富，并且依赖于最近的过去。输入 $u(t)$ 激励它，得到响应向量 $\mathbf r(t)$，只训练线性读出
\begin{empheq}[box=\keybox]{equation}
y(t) \;=\; W_{\text{out}}\,\mathbf r(t) ,
\label{eq:reservoir}
\end{empheq}
其权重用第~\ref{sec:motorlearning}~章的最小二乘规则~\eqref{eq:lms} 求得。储备池做的正是第~\ref{sec:motorlearning}~章里小\nt{GLUT}神经元做的事：把输入展开成一个长向量，使所需的类别在其中可以用平面分开（第~\ref{sec:dendrites}~章）。

电极阵列上的类器官正好可以充当这样的储备池。把语音以电流模式输入给它，只训练一个线性读出后，辨别说话人的准确率约为 $78\%$~\cite{Cai2023}。$78\%$ 这个数字是作者自己报告、媒体跟着转述的。这个结果有用，但要弄清其中的“智能”在哪里：类器官提供非线性和衰减的记忆，按误差学习是在外部、在硅片上完成的（图~\ref{fig:livingmachine}）。同时，类器官本身也并非一成不变：据作者报告，训练期间它自身的功能连接发生了变化，他们称之为类器官内部的无监督学习~\cite{Cai2023}。准确率中有多少来自这种变化、多少来自读出，并未区分。

\section{试验台告诉我们关于这台机器的什么}

\emph{能量。}按估计~\eqref{eq:energy}，一个脉冲约耗 $10^{-10}$~J。一百万个神经元的培养物，每个神经元每秒发放一个脉冲，用于信号传递的功耗为
\begin{empheq}[box=\keybox]{equation}
P \;\approx\; 10^{6}\times1\,\text{s}^{-1}\times10^{-10}\,\text{J} \;=\; 10^{-4}\,\text{W},
\label{eq:dishpower}
\end{empheq}
即十分之一毫瓦。刺激、记录和处理这些脉冲的电子设备功耗要高出几个数量级。与第~\ref{sec:debugging}~章一样，瓶颈不在计算，而在与计算的连接。

\emph{缺少的部件。}试验台展示了带\nt{GABA}流量控制、却没有控制寄存器的\nt{GLUT}总线能做多少事：自组织、产生群放电、在反馈作用下改变响应，以及充当好的储备池。没有控制寄存器，它做不到的是自己决定什么算成功。在试验台上这个信号由工程师给出，而在大脑中，正如我们在第~\ref{sec:dofa}--\ref{sec:acet}~章所见，由广播通道给出。

\emph{伦理。}类器官由人类细胞培养而成，它们越复杂，这种组织能否有所感受的问题就越严肃。工程视角给出一个部分答案：第~\ref{sec:pain}~章中的疼痛不是某一个传感器的信号，而是一整套带闸门、音量调节器和广播通道的警报系统，而试验台上没有这些。但这只是推理，不是证明；该领域本身也主张从一开始就让伦理评估与实验同步进行~\cite{Smirnova2023}。

\begin{keyblock}
\begin{proposition}[试验台上的活网络]
\label{prop:livingmachine}
电极阵列上由\nt{GLUT}和\nt{GABA}神经元组成的网络会自发产生群放电~\eqref{eq:burstinterval}，在反馈作用下改变响应，并可充当储备池~\eqref{eq:reservoir}，供外部训练的读出使用。这些都是已测量的。至于这样的网络按某种普遍规则学习，例如学习预测自己的输入，则是假说；而关于其“有感知能力”的高调说法，大多数专家并不接受。
\end{proposition}
\end{keyblock}

\section{闪光释放的多巴胺}
\label{sec:dishdopamine}

据我们所知，唯一一项直接让多巴胺给培养物当老师的实验，是在一个研究者可远程接入的类器官平台上完成的~\cite{Jordan2024}。在浸泡类器官的液体中加入装在光敏“笼”里的多巴胺：被笼锁住的分子不作用于受体。笼锁多巴胺的浓度为 $1$~mM。需要奖励网络时，用紫外光照射：笼分解，多巴胺释放到周围空间。

回路是这样的。同一个刺激反复施加；如果它引发的脉冲比以前多，就打开闪光、释放多巴胺。在 $240$ 次重复中，诱发脉冲数总体上在增加。作者自己写道，仅凭这一项实验不能断定增长是多巴胺引起的~\cite{Jordan2024}。

对工程师来说，这是在培养皿中搭建三因子规则~\eqref{eq:threefactor} 的第一次尝试：发送者和接收者的活动留下资格迹，由一个公共信号决定是否把变化固定下来。但这里的信号只能为正：有奖励或没有奖励，装置给不出负误差 $\delta<0$。而且多巴胺扩散和清除都很慢，就像~\eqref{eq:lowpass} 中的浓度那样，因此频繁的闪光可能只是抬高了它的总体水平。要把学习与这种效应区分开，需要两个对照：在与网络响应无关的随机时刻给同样多的闪光；以及不加笼锁多巴胺、只给同样的闪光。还需要休息后的检验：多巴胺消散后，变化是否依然存在。

\section{多巴胺教硅片}
\label{sec:biohybrid}

在反方向的实验中，活细胞教电子器件~\cite{Keene2020}。把分泌多巴胺的细胞培养在一个有机神经形态元件上方的微通道中，这个元件是一只晶体管，其沟道电导充当突触权重。细胞释放的多巴胺在元件电极处被氧化，从而改变电导。该器件还模拟了突触间隙中递质的清除机制，因此权重不仅能长期改变，还能复原。

从电路角度看，这是一个化学输入的泄漏积分器：权重累积多巴胺流量，而“清除”决定它遗忘的快慢，与~\eqref{eq:lowpass} 是同一个RC电路。这里的关键是连接方向。第~\ref{sec:debugging}~章的探针看不到广播寄存器，因为它们是化学的。这个元件读取的恰恰是化学寄存器，并把它转换成电学权重。对脑机接口来说，这是通向一种传感器的路径：它不仅能听到数据总线，还能听到控制寄存器。

\section{自带多巴胺神经元的类器官}
\label{sec:midbrainorg}

人们已经能用人类干细胞培养出类器官，它们类似于大脑中\nt{DOPA}神经元所在的那个区域。其中有能工作、能分泌多巴胺的多巴胺神经元~\cite{Jo2016}。这类类器官主要用于研究疾病。我们没有找到用它们的多巴胺给另一个网络当老师的实验。

对于用活神经元搭成的机器，这正是缺失的部件。第~\ref{sec:livingmachine}~章的试验台是没有控制寄存器的数据总线加流量控制；这里已经有了广播信号源。缺的是评论器：一个计算预测误差~\eqref{eq:ctd} 并控制这些神经元的网络。把皮层类器官和这样的类器官连起来，让多巴胺按误差释放，是一个开放问题，目前只是假说，还不是实验。

\section{不靠多巴胺平衡杆子的类器官}
\label{sec:cartpole}

近期最完整的一项实验完全不用多巴胺~\cite{Robbins2026}。把小鼠皮层类器官接入“小车倒立摆”任务的闭环：类器官的活动推动小车，杆子的位置通过刺激反馈回来。两次尝试之间，类器官接受短暂的高频教学刺激。具体给哪些刺激，由一个强化学习程序选择。

结果是测量出来的：
\begin{itemize}
\item 对大多数类器官，程序选择的教学信号比随机选择的信号或没有信号效果更好；
\item 休息 $45$ 分钟后，改进没有保留；
\item 阻断两类谷氨酸受体AMPA和NMDA，改进消失。
\end{itemize}

最后一条说明学到的东西存在哪里：在\nt{GLUT}突触中，并且经由第~\ref{sec:weightstore}~节中充当输入—输出符合检测器的那个受体。第二条说明缺了什么：有快速变化，没有巩固，就像第~\ref{sec:motorlearning}~章中只有快学习者、没有慢学习者的情形。老师的角色则变了：挑选电流模式的工程师被程序取代，但老师仍然站在培养皿之外。

在更早的电极阵列培养物学习实验中，我们也没有找到任何一项使用多巴胺的：教学信号一律是电刺激。

\section{谁在教培养物}

\begin{table}[t]
\centering
\footnotesize
\setlength{\tabcolsep}{4pt}
\renewcommand{\arraystretch}{1.15}
\begin{tabular}{>{\raggedright}p{2.7cm}>{\raggedright}p{2.5cm}>{\raggedright}p{3.2cm}>{\raggedright\arraybackslash}p{4.4cm}}
\toprule
实验 & 谁来教 & 教学信号 & 已知情况 \\
\midrule
出现所需响应时停止刺激 & 工程师 & 刺激结束 & 响应被固定——已测量 \\
DishBrain（第~\ref{sec:dishbrain}~章） & 工程师 & 可预测电流或随机电流 & 一次会话内回合数的显著增长只出现在完整反馈下，安静条件下效应较小——已测量；跨会话不稳定；原因仍有争论 \\
小车倒立摆 & 强化学习程序 & 程序选择的高频刺激 & 优于随机刺激，但休息后不保留——已测量 \\
闪光释放多巴胺 & 规则“脉冲更多即奖励” & 光释放的多巴胺 & 脉冲增加——观察结果；多巴胺的作用未证实 \\
生物混合突触 & 活细胞 & 细胞分泌的多巴胺 & 元件权重的长期改变与复原——已测量 \\
含\nt{DOPA}神经元的类器官 & --- & --- & 有多巴胺来源；未用作老师 \\
\bottomrule
\end{tabular}
\caption{谁在用什么教芯片上的活神经元。}
\label{tab:dishteachers}
\end{table}

在所有由培养物本身学习的实验中，老师目前都在外面（表~\ref{tab:dishteachers}）：工程师、程序，或工程师设定的规则。多巴胺只进过培养皿一次，它的作用还有待证明。在培养皿中搭建真正的广播通道，像第~\ref{sec:dofa}~章那样带有评论器、误差和资格迹，是下一步，而这一步还没有人迈出。


\section{小结}

\begin{table}[t]
\centering
\footnotesize
\setlength{\tabcolsep}{4pt}
\renewcommand{\arraystretch}{1.15}
\begin{tabular}{>{\raggedright}p{2.8cm}>{\raggedright}p{4.2cm}>{\raggedright\arraybackslash}p{5.8cm}}
\toprule
实验 & 活网络做什么 & 已知情况 \\
\midrule
无输入的网络 & 群放电，间隔从一秒到几分钟~\eqref{eq:burstinterval} & 已测量；通过突触耗竭的机制是公认模型 \\
会闭嘴的老师 & 关掉刺激时，所需响应被固定下来 & 已测量 \\
打网球 & 连续击球回合变长；一次会话内的显著提升只出现在完整反馈下 & 行为变化已测量；跨会话学习不稳定；效应大小和“预测输入”的解释有争议 \\
虚拟身体 & 在挑选好的刺激模式下朝目标移动 & 已测量 \\
储备池 & 为读出~\eqref{eq:reservoir} 提供非线性和记忆 & 已测量；按误差学习在外部、在硅片上完成；类器官的连接也在变化 \\
能量 & 每一百万个神经元约 $10^{-4}$~W~\eqref{eq:dishpower} & 按~\eqref{eq:energy} 估计 \\
\bottomrule
\end{tabular}
\caption{用活神经元搭成的机器展示了什么。}
\label{tab:livingmachine}
\end{table}

用活神经元搭成的机器（表~\ref{tab:livingmachine}）是一个试验台，从中可以看出：没有控制寄存器时，数据总线和流量控制能做些什么。它们能做很多事，却不能自己决定什么是好。在试验台上，这个角色由工程师和他的电流模式扮演；在大脑中，由本书中间部分所讲的那几条广播导线扮演。

\section{习题}

\begin{exercise}[\lvlA]
\label{ex:21:1}
\emph{群放电间隔。}一次群放电把储备耗尽到零，随后 $\tau_{\text{rec}}\,\dd x/\dd t=1-x$，当 $x=x^*$ 时开始新的群放电；$\tau_{\text{rec}}=0.8$~s。(a)~求 $x^*=0.9$ 和 $0.99$ 时的间隔。(b)~阈值 $x^*$ 在 $0.99$ 附近随机变化 $\pm0.005$，间隔落在什么范围内？
\end{exercise}

\begin{exercise}[\lvlA]
\label{ex:21:2}
\emph{试验台的能量。}(a)~对整个大脑（$8.6\cdot10^{10}$ 个神经元），估计~\eqref{eq:dishpower} 给出多大功率？(b)~试验台以 $20$~kHz的采样率记录 $1024$ 个电极，每个采样 $10$~bit。按每bit $1$~nJ计，传输这一数据流的功耗是培养物本身功耗的多少倍？
\end{exercise}

\begin{exercise}[\lvlB]
\label{ex:21:3}
\emph{设计：抑制群放电。}通过许多电极的稀疏刺激使每个突触以频率 $r_b$ 释放递质；储备满足 $\tau_{\text{rec}}\,\dd x/\dd t=1-x-Ur_b\tau_{\text{rec}}x$，$U=0.5$，$\tau_{\text{rec}}=0.8$~s。$r_b$ 至少多大时，储备长不到 $x^*=0.9$；$0.99$？此时突触减弱多少？
\end{exercise}

\begin{exercise}[\lvlB]
\label{ex:21:4}
\emph{储备池的记忆不超过神经元数。}有 $N$ 个输出 $\mathbf r(t)$ 的储备池接收零均值、单位方差的白噪声输入 $u(t)$；$m_k$ 是读出 $\mathbf w_k\cdot\mathbf r(t)$ 与 $u(t-k)$ 的相关系数平方的最大值。证明记忆容量 $\mathrm{MC}=\sum_{k\ge0}m_k\le N$。
\end{exercise}

\begin{exercise}[\lvlB]
\label{ex:21:5}
\emph{建模：硅片中的储备池。}储备池 $\mathbf r(t+1)=\tanh(W\mathbf r(t)+\mathbf v\,u(t)+\mathbf c)$，$N=30$，$W$ 为随机矩阵，谱半径 $0.9$，$v_i\sim U(-0.5,0.5)$，$c_i\sim U(-0.2,0.2)$，$u\sim U(-1,1)$，读出用岭回归。求带 $\tanh$ 和不带它时习题~\ref{ex:21:4} 的记忆容量。哪个储备池记得更多？
\end{exercise}

\begin{exercise}[\lvlB]
\label{ex:21:6}
\emph{为活储备池设计读出。}类器官提供 $1024$ 个通道和两类共 $P=240$ 条训练记录。(a)~按习题~\ref{ex:19:3} 的公式，最多保留多少个特征 $n$，才能使随机标注可分的概率不超过 $1\%$？(b)~在 $100$ 条测试记录上 $78\%$ 的准确率比随机猜测高出几个标准差？
\end{exercise}

\begin{exercise}[\lvlC]
\label{ex:21:7}
\emph{研究：没有广播通道的老师。}提出一个刺激方案，通过 $8$--$1024$ 个电极在试验台上实现第~\ref{sec:dofa}~章的三因子规则；再设计一个冻结读出的对照，用来区分权重学习与网络状态漂移。什么结果会否定“学习”这一假说？
\end{exercise}

\chapter{DishBrain}
\chaptermark{DishBrain}
\label{sec:dishbrain}

\section{培养皿里的网络打网球}

DishBrain 是作者给一个试验台起的名字：在这个试验台上，培养在带电极芯片上的活神经元在玩单球拍的简化网球~\cite{Kagan2022}。这是用活神经元搭成的机器中讨论最多的一项实验（这类机器的综述见第~\ref{sec:livingmachine}~章），对本书来说也格外合用。这里的小网络整个都摸得着：每个输入由实验者设定，每个输出都被记录下来。前面各章的所有问题，即信号如何编码、谁来教、探针能看到什么，都可以拿来问一个既没有眼睛、没有肌肉、也没有 \nt{DOPA} 神经元的网络。我们像工程师剖析别人的试验台那样来剖析这项实验：电路、输入、输出、老师、结果、争议。

\section{试验台}

神经元取自小鼠胚胎皮层，每块芯片约 $800\,000$ 个细胞；或由人类干细胞培养而来，每块芯片约一百万个。它们在芯片上生长几周，彼此长入对方，并开始自发地发放脉冲和簇放电。在它们下方 $8$~mm$^2$ 的区域上有 $26\,000$ 个铂电极；最多可同时记录其中的 $1024$ 个，而实际用于刺激的只有八个。

培养物周围闭合着一个回路（图~\ref{fig:dishloop}）。计算机模拟游戏：球飞行、在墙上反弹，球拍上下移动。球的位置被转换成八个“感觉”电极上的电流。两个“运动”区的脉冲按 $10\ms$ 的窗口计数，用来移动球拍。每次击中或未击中后，培养物还会再收到一个刺激，即反馈。$10\ms$ 窗口是试验台计算机的时钟节拍，而不是培养物的：网络本身仍像本书中所有网络一样异步工作。

\begin{figure}[t]
\centering
\begin{tikzpicture}[>=Stealth,
  blk/.style={draw,very thick,rounded corners=3pt,align=center,font=\small,minimum height=1.2cm},
  lab/.style={font=\scriptsize,align=center}]
\node[blk,fill=blue!6,text width=3.4cm] (C) at (0,0) {芯片上的培养物\\\scriptsize $\sim10^6$ 个神经元，$26\,000$ 个电极};
\node[blk,fill=orange!10,text width=3.0cm] (G) at (7.2,0) {计算机：\\网球游戏};
\draw[->,very thick,blue!60!black] (G.north west) to[bend right=25] node[lab,above] {球：\emph{位置}——8 个电极中的哪一个，\\\emph{频率}——4--40 Hz} (C.north east);
\draw[->,very thick,red!70!black] (C.south east) to[bend right=25] node[lab,below] {球拍：“上”区和“下”区\\在 $10\ms$ 内的脉冲计数} (G.south west);
\draw[->,thick,densely dashed,green!45!black] (G.east) -- ++(0.9,0) |- ++(0,-2.6) -| node[lab,pos=0.25,below] {接住：可预测刺激；未接住：不可预测刺激} (C.south);
\end{tikzpicture}
\caption{DishBrain 回路。蓝色箭头是输入：球的位置用位置编码和频率编码表示（第~\ref{sec:code}~章）。红色箭头是输出：两个区的计数之差移动球拍。绿色虚线是每个回合之后的反馈，也是试验台上唯一的“老师”。}
\label{fig:dishloop}
\end{figure}

\section{输入与输出}

\emph{输入}同时用了第~\ref{sec:code}~章的两种编码。刺激八个电极中的哪一个，表示球在什么高度：这是位置编码，和第~\ref{sec:sensors}~章的传感器一样。刺激脉冲有多密，表示球有多远：球在远端墙边时每秒 $4$ 个，离球拍所在的墙越近线性增多，直到每秒 $40$ 个。这是频率编码；球的刺激脉冲幅度为 $75$~mV。培养物事先并不“认识”这两种编码中的任何一种：编码由实验者规定，网络必须学会读懂它。

\emph{输出}的读取方式与第~\ref{sec:debugging}~章的接口相同。一个运动区的脉冲把球拍往上推，另一个往下推；最简单的写法是，每个窗口内
\begin{empheq}[box=\keybox]{equation}
\Delta p \;=\; c\,\bigl(n_\uparrow - n_\downarrow\bigr),
\label{eq:dishreadout}
\end{empheq}
其中 $n_\uparrow$、$n_\downarrow$ 是两个区在 $10\ms$ 内的脉冲数，$c$ 是试验台的系数。这是第~\ref{sec:glutcomputer}~章的双线编码：运动方向不是由信号的符号携带，而是由两根导线中哪一根更响来携带。

来看看这种输出的噪声。如果一个区每秒总共发出 $R$ 个脉冲，那么在 $T=10\ms$ 的窗口里平均有 $RT$ 个，按~\eqref{eq:rateerror}，相对离散度为 $1/\sqrt{RT}$。$R=1000$ 次/秒时每个窗口约为百分之 $30$；$R=100$ 时超过百分之百。球拍不可避免地会抖动，网络只能用一种方式学习：把计数之差的\emph{平均值}往正确方向推。这正是限制第~\ref{sec:debugging}~章中任何接口的同一套规律，只不过现在出现在回路的两侧。

\section{没有多巴胺的老师}

试验台最有意思的部分是老师。皮层神经元培养物里没有 \nt{DOPA} 神经元，试验台也不发送任何“比预期好”的信号。反馈是另一种形式。网络接住了球，所有八个感觉电极同时收到一个短暂的\emph{可预测}刺激：$75$~mV 的脉冲，每秒 $100$ 个，持续 $0.1\,\text{s}$。没接住，就是四秒钟作者称为\emph{不可预测}的刺激：$150$~mV 的脉冲，每秒 $5$ 个，在随机时刻施加到随机选择的电极上；随后是四秒钟不刺激的停顿，球重新发出~\cite{Kagan2022}。

作者的出发点是一个假说：网络的行为是为了让自己的输入变得可预测~\cite{Friston2010}。对工程师来说含义很简单：培养物摆脱四秒随机噪声的办法只有一个，就是接住球。同样的思想我们在第~\ref{sec:code}~章的节省脉冲和第~\ref{sec:recognition}~章的画面预测中都见过。

但这里需要的真是可预测性吗？还可能有更粗糙的机制。假设未接住后的不可预测刺激只是把网络中近期的变化\emph{打乱}，而接住后的平静刺激不去碰它们。用一个向量 $\boldsymbol\psi$ 描述网络的设置，设在设置 $\boldsymbol\psi$ 下网络未接住的概率为 $P_{\text{miss}}(\boldsymbol\psi)$。如果扰动是对称的，即从 $\boldsymbol\psi$ 到 $\boldsymbol\psi'$ 与反方向的可能性相同，那么在稳态下，离开每个设置的流量与进入的流量相平衡，网络停留在设置 $\boldsymbol\psi$ 的时间比例为
\begin{empheq}[box=\keybox]{equation}
\rho(\boldsymbol\psi) \;\propto\; \frac{1}{P_{\text{miss}}(\boldsymbol\psi)} .
\label{eq:dishdensity}
\end{empheq}
未接住的次数少十倍的设置，保持的时间就长十倍。没有人告诉网络该往哪个方向改权重：好的设置只是更少被破坏。

我们在一个模型上检验了这一点，把培养物简化为两个数：球到达高度 $y$ 时，球拍到达 $a\,y+b$，偏差小于 $0.1$ 就算接住。未接住后两个数都受到一次随机扰动，接住后保持不变。在 $2000$ 个回合中，接住的比例从前两百个回合的 $0.21$ 上升到最后五百个回合的 $0.38$（四十个“培养物”，图~\ref{fig:dishmodel}）。如果在\emph{每个}回合后都给扰动，反馈中就没有信息，比例停留在 $0.12$ 左右；如果完全没有扰动，比例保持在 $0.19$。试验台上的实验与此一致：一次会话内的显著提升只出现在完整反馈下。以安静代替刺激时，培养物到会话结束时仍比没有反馈时打得好，但效果较小~\cite{Kagan2022}；需要指出，在这种条件下，未接住后也跟着较长的停顿，接住后则较短，因此两种结局之间的差别并未完全消失。

\begin{figure}[t]
\centering
\begin{tikzpicture}[x=1cm,y=5cm]
\draw[->,gray!70!black] (0,0) -- (10.6,0);
\draw[->,gray!70!black] (0,0) -- (0,0.5) node[above,font=\small,black] {接住比例};
\foreach \v in {0.1,0.2,0.3,0.4}{\draw[gray!60!black] (-0.06,\v) -- (0.06,\v); \node[left,font=\scriptsize] at (0,\v) {\v};}
\foreach \x/\f/\l/\lab in {1.8/0.21/0.38/{未接住后\\给噪声},5.2/0.13/0.11/{每个回合后\\都给噪声},8.6/0.19/0.19/{无噪声}}{
  \fill[blue!35] (\x-0.8,0) rectangle (\x-0.05,\f);
  \fill[red!60!black] (\x+0.05,0) rectangle (\x+0.8,\l);
  \node[below,font=\scriptsize,align=center] at (\x,-0.01) {\lab};
  \node[above,font=\scriptsize] at (\x-0.425,\f) {\f};
  \node[above,font=\scriptsize] at (\x+0.425,\l) {\l};}
\fill[blue!35] (7.4,0.47) rectangle (7.7,0.495); \node[right,font=\scriptsize] at (7.75,0.482) {前 200 个};
\fill[red!60!black] (7.4,0.41) rectangle (7.7,0.435); \node[right,font=\scriptsize] at (7.75,0.422) {后 500 个};
\end{tikzpicture}
\caption{“未接住即打乱”模型（\texttt{dishbrain\_model.py}）：$2000$ 个回合开始和结束时的接住比例，四十个模型培养物的平均值。只有当扰动跟在未接住之后、而不跟在接住之后时，学习才会发生，这与实验一致：一次会话内的显著提升只出现在完整反馈下。}
\label{fig:dishmodel}
\end{figure}

\begin{keyblock}
\begin{proposition}[打乱式老师]
\label{prop:dishscramble}
如果失败会晃动网络、成功则让它保持原样，网络就会自行漂向更少出错的设置：停留在某个设置的时间与失败概率成反比~\eqref{eq:dishdensity}。这种学习既不需要奖励信号，也不需要其符号，也不需要可预测性本身，只需要成功后和失败后的反馈有所不同。培养物的行为变化已测量；培养皿里起作用的是哪种机制，是输入预测还是打乱，尚未确定。
\end{proposition}
\end{keyblock}

与第~\ref{sec:dofa}~章比较一下。那里噪声也用于搜索，但多巴胺带有符号：误差 $\delta$ 指明权重该往哪个方向移动，步子沿梯度走~\eqref{eq:perturbation}。这里没有符号，只有“保留”或“晃动”。这种搜索更粗糙、更慢，但对网络的要求更少：不需要 \nt{DOPA} 神经元，不需要评论家，也不需要延续数秒的资格迹。

\section{测到了什么，争的是什么}

每局游戏持续 $20$ 分钟；比较的是前 $5$ 分钟和后 $15$ 分钟。在有完整反馈的培养物中，连续击球的回合变长，发球即丢的情况变少；在没有细胞、没有游戏以及由计算机控制“球拍”的对照中，都没有类似现象。人类细胞培养物平均能比小鼠培养物坚持更久。这一改进在统计上是可靠的，涉及九个小鼠培养物和十一个人类培养物，每组一百多局游戏，但幅度不大：网络并没有成为好球手，只是比随机稍好一点。它在一次会话内可靠；在几天里跨会话的稳定学习，作者并没有得到~\cite{Kagan2022}。同一小组的后续工作发现，游戏时培养物的活动会趋近\emph{临界}状态，即衰减与雪崩之间的边界，也就是第~\ref{sec:gaba}~章中那种处在边缘的状态，而且越接近临界，游戏打得越好~\cite{Habibollahi2023}。

引起主要争议的不是数字，而是措辞。文章标题宣称培养皿中的神经元“表现出感知能力”（sentience），一组研究者对此提出反驳：闭环中的行为变化还不是智能，也不是感受，结论应当表述得更谨慎~\cite{Balci2023}。从工程角度看，有两个问题悬而未决，都关乎机制。第一，网络是在学习，即它的权重长期改变了，还是反馈只是移动了它的当前状态？第二，需要的真是可预测性，还是只要对成功和失败的响应有任何差别就够了，如命题~\ref{prop:dishscramble} 所述？在每个电极都可访问的试验台上，这两个问题都可以回答，这或许正是它最大的价值。

\section{试验台规格：如何复现}
\label{sec:dishspec}

下面汇总了重新搭建这样一个试验台所需的全部内容：设备、回路、输入编码、输出读取、反馈和实验方案。每个参数都注明了来源。“论文”表示数值取自文献~\cite{Kagan2022} 的正文。“自选”表示论文中没有给出，复现时需要自行设定；我们给出默认值，并说明如何检验。

\subsection*{设备与培养物}

\begin{center}
\footnotesize
\setlength{\tabcolsep}{4pt}
\renewcommand{\arraystretch}{1.15}
\begin{tabular}{>{\raggedright}p{4.0cm}>{\raggedright}p{6.9cm}>{\raggedright\arraybackslash}p{1.6cm}}
\toprule
参数 & 数值 & 来源 \\
\midrule
芯片 & MaxOne 阵列：$8$~mm$^2$ 区域上 $26\,000$ 个铂电极 & 论文 \\
记录 & 最多同时 $1024$ 个通道，每秒 $20\,000$ 个采样 & 论文 \\
刺激 & 技术上最多 $32$ 个电极，实验中用 $8$ 个独立电极 & 论文 \\
细胞 & 小鼠胚胎皮层；由干细胞获得的人类神经元（两种培养方法）；HEK293T 细胞（非神经元）作对照 & 论文 \\
接种 & 每块芯片约 $10^6$ 个细胞 & 论文 \\
实验时的培养龄 & 小鼠：从 $\sim10$ 天起；人类：视培养方法为 $14$--$24$ 天或 $\sim80$ 天 & 论文 \\
\bottomrule
\end{tabular}
\end{center}

\subsection*{回路与时间}

回路以 $10\ms$（$200$ 个采样）为窗口运行：每个窗口内统计运动区电极上的脉冲，更新游戏，并发出下一次刺激。从培养物中的脉冲到响应刺激的延迟约 $5\ms$。游戏期间，每个电极的信号经过截止频率 $100$~Hz 的二阶贝塞尔高通滤波器；信号的绝对值再经 $1$~Hz 的一阶贝塞尔低通滤波器平滑，脉冲阈值与这个平滑后的绝对值成正比。论文中六倍背景标准差的阈值用于单独的活动测量，而不是游戏。为了不把刺激误算为培养物的响应，当同时出现超过 $15$ 个大于 $75$~mV 的大脉冲时，所有信号都被屏蔽。以上均来自论文。

\subsection*{球的编码}

\begin{center}
\footnotesize
\setlength{\tabcolsep}{4pt}
\renewcommand{\arraystretch}{1.15}
\begin{tabular}{>{\raggedright}p{3.4cm}>{\raggedright}p{7.5cm}>{\raggedright\arraybackslash}p{1.6cm}}
\toprule
参数 & 数值 & 来源 \\
\midrule
感觉区 & $8$ 个电极，排列使相邻电极表示球的相邻位置 & 论文 \\
位置编码 & 电极编号 $k=1,\dots,8$ 表示球相对于球拍中心的位置 & 论文 \\
哪个坐标 & 球的竖直位置；复现时取相对于球拍的位置，$y-p\in[-\tfrac12,\tfrac12]$ 时 $k=\lceil 8(y-p+\tfrac12)\rceil$ & 论文；公式为自选 \\
频率编码 & $4$ 到 $40$~Hz 的刺激频率表示到球拍的距离 & 论文 \\
标度方向 & 球在对面墙边时 $4$~Hz，线性增加到球在球拍一侧墙边时的 $40$~Hz：$f=4+36\,(1-x/L)$，$x$ 为到球拍一侧墙的距离，$L$ 为场地宽度 & 论文 \\
脉冲波形 & 短的矩形双相脉冲：先正相，后负相 & 论文 \\
幅度 & $75$~mV；选定这一幅度是为了不强迫受抑制的细胞发放 & 论文 \\
相位时长 & 未给出；采用刺激系统的标准设置，实验期间不改变 & 自选 \\
\bottomrule
\end{tabular}
\end{center}

\subsection*{读取球拍}

论文中有两个运动区：第一个区的脉冲让球拍上移，第二个区的让它下移；两个区的贡献经过加权，以补偿细胞密度和活动水平的差异~\cite{Kagan2022}。确切公式没有给出。复现时我们采用规则~\eqref{eq:dishreadout}，每个 $10\ms$ 窗口 $\Delta p=c\,(n_\uparrow-n_\downarrow)$，用权重按静息平均活动把两个区拉平（自选）：$n_\uparrow$ 和 $n_\downarrow$ 各自除以游戏前测得的本区每窗口平均计数。系数 $c$ 取为使一个平均计数的差每窗口把球拍移动场地高度的 $1\%$（自选）；这样，一个区持续占优时，球拍 $1$~s 就能走完整个场地。

论文正文没有用坐标给出各区在芯片上的位置，只有一张示意图。复现时只需三组互不重叠的电极：中间八个感觉电极，两侧各一组记录电极，一组为“上”，一组为“下”（自选）。

\subsection*{游戏}

论文没有给出场地尺寸、球拍长度和球速；只说球以恒定速度飞行并在墙上反弹。复现时（自选）：场地 $1\times1$，球拍在右侧墙上，长 $0.2$（$|y-p|<0.1$ 即击中，与模型~\texttt{models/dishbrain\_model.py} 一致），球 $1$~s 飞越场地。一个回合中一次都没接住就未接住，记为“ace”；连续击球超过三次的回合记为“长回合”（论文）。

\subsection*{反馈}

\begin{center}
\footnotesize
\setlength{\tabcolsep}{4pt}
\renewcommand{\arraystretch}{1.15}
\begin{tabular}{>{\raggedright}p{2.6cm}>{\raggedright}p{8.3cm}>{\raggedright\arraybackslash}p{1.6cm}}
\toprule
事件 & 施加的刺激 & 来源 \\
\midrule
击中 & 可预测刺激：全部 $8$ 个感觉电极同时，$75$~mV，$100$~Hz，$100\ms$；其间暂停通常的球位置刺激 & 论文 \\
未击中 & 不可预测刺激：$150$~mV，$5$~Hz，$4$~s，在随机时刻施加到从 $8$ 个电极中随机选出的电极上；随后 $4$~s 无刺激，球沿随机方向重新发出 & 论文 \\
随机规律 & 未给出；复现时把脉冲时刻取为平均频率 $5$~Hz 的泊松流 & 自选 \\
\bottomrule
\end{tabular}
\end{center}

\subsection*{实验方案与对照}

一次会话持续 $20$~min；比较前 $5$~min 和后 $15$~min（论文）。三项结果指标：平均回合长度、ace 比例和长回合比例。实验条件（论文）：
\begin{itemize}
\item \emph{刺激}：完整回路，如上所述；
\item \emph{安静}：以同样长的无任何刺激时段代替击中或未击中刺激，然后球沿随机方向发出；
\item \emph{无反馈}：未接住后球直接反弹继续飞行，球位置刺激继续进行；
\item \emph{静息}：游戏进行、球拍随活动移动，但完全没有刺激；
\item \emph{对照}：只有培养基的芯片；HEK293T 细胞；“计算机里的培养物”，即用仿真代替细胞。
\end{itemize}
主系列的样本量：小鼠培养物 $9$ 个（$101$ 次会话），人类培养物 $11$ 个（$138$），只有培养基的芯片 $6$ 块（$80$），静息培养物 $20$ 个（$42$），仿真 $3$ 个（$38$），共 $399$ 次会话。反馈系列：$3$ 天、每天 $3$ 次会话，共 $486$ 次会话，条件为“刺激”“安静”和“无反馈”（$n=27$、$17$ 和 $15$）。平均回合长度从前 $5$ 分钟到后 $15$ 分钟的增长只出现在活的培养物中。在反馈系列中，随时间的显著增长只出现在“刺激”条件下；会话结束时，“安静”条件仍优于“无反馈”，但效应较小；而在三天里跨会话的稳定学习并未出现~\cite{Kagan2022}。

\subsection*{试验台循环的步骤}

每隔 $10\ms$，试验台计算机执行以下操作。
\begin{enumerate}
\item 读取上一窗口内两个运动区的脉冲数 $n_\uparrow$、$n_\downarrow$，并用各区静息平均计数归一化。
\item 把球拍移动 $\Delta p=c\,(n_\uparrow-n_\downarrow)$，并限制在场地边界内。
\item 把球移动一步；在上、下、左三面墙上反弹。
\item 如果球到达右墙：若 $|y-p|<0.1$ 则为击中，回合计数加一，施加击中刺激（$100\ms$）；否则为未击中，记录该回合，施加未击中刺激（$4$~s），再停顿 $4$~s，球重新发出。
\item 否则，按球的竖直位置选择电极 $k$，按到球拍一侧墙的距离选择频率 $f$，在下一个窗口之前向电极 $k$ 施加频率为 $f$、幅度 $75$~mV 的双相脉冲。
\end{enumerate}
这些足以搭建一个与已发表试验台兼容的装置，并检验本章的核心问题：教会培养物的是可预测性，还是仅仅是对击中与未击中响应之间的差别。为此，应在论文的三个条件之外加上论文中没有的第四个：未接住后施加与不可预测刺激时长和能量相同的\emph{可预测}刺激。如果培养物这样也能学会，关键就在于差别，而不在于可预测性。


\section{小结}

\begin{table}[t]
\centering
\footnotesize
\setlength{\tabcolsep}{4pt}
\renewcommand{\arraystretch}{1.15}
\begin{tabular}{>{\raggedright}p{2.1cm}>{\raggedright}p{4.2cm}>{\raggedright}p{1.9cm}>{\raggedright\arraybackslash}p{4.6cm}}
\toprule
试验台部件 & 结构 & 本书章节 & 已知情况 \\
\midrule
输入 & 位置（8 个电极中的哪一个）和频率（4--40 Hz） & \ref{sec:code}, \ref{sec:sensors} & 由实验者规定 \\
输出 & 两个区在 $10\ms$ 内的计数之差~\eqref{eq:dishreadout} & \ref{sec:glutcomputer}, \ref{sec:debugging} & 由实验者规定；噪声按~\eqref{eq:rateerror} \\
老师 & 接住后给可预测刺激，未接住后给不可预测刺激；没有多巴胺 & \ref{sec:dofa} & 一次会话内的显著提升只出现在完整反馈下；安静条件下效应较小；跨会话不稳定 \\
机制 & 输入预测或未接住即打乱~\eqref{eq:dishdensity} & \ref{sec:dofa}, \ref{sec:recognition} & 未确定；两者都能解释结果 \\
网络状态 & 游戏时趋近临界状态 & \ref{sec:gaba} & 与游戏水平的关系已测量 \\
“感知能力” & 作者的论断 & --- & 未被证明；有争议 \\
\bottomrule
\end{tabular}
\caption{工程师眼中的 DishBrain。}
\label{tab:dishbrain}
\end{table}

DishBrain 是最简形式的活神经元机器：输入用编码给定，输出靠计数读取，老师就是对成功与失败的响应之差（表~\ref{tab:dishbrain}）。一个没有眼睛、肌肉和多巴胺的网络稍微学会了打球，单凭这一点，它就成了检验本书思想的试验台。无论最后的机制是预测、打乱还是别的什么，答案都会按第~\ref{sec:debugging}~章建议的方式找到：用探针、测试信号和关闭部件，在一台终于能整个看清的机器上。

\section{习题}

\begin{exercise}[\lvlA]
\label{ex:22:1}
\emph{计数要偏移多少。}两个区合计每秒发出 $R_\uparrow+R_\downarrow=2000$ 个脉冲，均为泊松流。(a)~$R=1000$ 和 $R=100$ 时，一个区在 $10\ms$ 内计数的相对离散度是多少？(b)~球飞行 $1$~s 期间位移~\eqref{eq:dishreadout} 累加。$R_\uparrow-R_\downarrow$ 至少多大，平均位移才是离散度的两倍？
\end{exercise}

\begin{exercise}[\lvlA]
\label{ex:22:2}
\emph{要多少回合才能看出学习。}各回合独立，接住比例服从二项分布。两个时段中各需多少回合，才能使接住比例的增长超过差值的两个标准差：(a)~从 $0.20$ 到 $0.25$；(b)~如模型中那样从 $0.21$ 到 $0.38$？
\end{exercise}

\begin{exercise}[\lvlB]
\label{ex:22:3}
\emph{推导~\eqref{eq:dishdensity}。}设置 $\boldsymbol\psi$ 取值于有限集合；未接住时（概率 $P(\boldsymbol\psi)>0$），网络以对称概率 $K(\boldsymbol\psi,\boldsymbol\psi')$ 转到 $\boldsymbol\psi'$，接住时保持不变。(a)~证明 $\rho\propto1/P$ 是平稳分布。(b)~两个设置分别为 $P=0.9$ 和 $P=0.1$ 时，未接住的比例是多少？
\end{exercise}

\begin{exercise}[\lvlB]
\label{ex:22:4}
\emph{模型的吸收区。}球拍到达 $p=\min\bigl(1,\max(0,ay+b)\bigr)$，球到达高度 $y\sim U(0,1)$，$|p-y|<h=0.1$ 即接住。(a)~证明当 $|b|<h$ 且 $|a-1+b|<h$ 时网络接住所有球，并求该区域的面积。(b)~数值求 $a\sim U(-1,1)$、$b\sim U(0,1)$ 时的平均接住比例。
\end{exercise}

\begin{exercise}[\lvlB]
\label{ex:22:5}
\emph{建模：给设置围上栅栏。}在 \texttt{dishbrain\_model.py} 的副本中，让 $200$ 个培养物各打 $20\,000$ 个回合。最后 $500$ 个回合的接住比例是多少？如果每次扰动后把设置裁剪到矩形 $a\in[-1,2]$、$b\in[-1,1]$ 内呢？用习题~\ref{ex:22:3} 解释差别。
\end{exercise}

\begin{exercise}[\lvlB]
\label{ex:22:6}
\emph{设计输入编码。}高度误差小于 $h=0.1$ 即接住；网络在 $T=0.25$~s 内读取输入，频率不超过 $r$ 时约可区分 $\sqrt{rT}$ 个等级，刺激不超过每秒 $40$ 个脉冲。(a)~八个电极上的位置编码够用吗？(b)~能否用单个电极上的频率来传递高度？
\end{exercise}

\begin{exercise}[\lvlC]
\label{ex:22:7}
\emph{研究：预测还是打乱？}把模型推广到 $d$ 个可调数（$p=\sum_{i<d}a_iy^i$）。在打乱规则下和在带误差符号的规则~\eqref{eq:perturbation} 下，达到接住比例 $0.5$ 所需的回合数如何随 $d$ 增长？试验台上什么实验能区分这两个假说？
\end{exercise}

\appendix
\chapter{不决定神经元类型的物质}
\label{sec:othermediators}

在第~\ref{sec:neurontypes}~章中，我们按主要神经递质对神经元分类，得到了十种类型（表~\ref{tab:types}）。现在该加上一些复杂情况了。神经元释放的不只是主要递质，而且在大脑中，除了主要递质，还有别的物质在起作用。它们改变网络传递和存储信号的方式。

除了寻址总线和广播总线（第~\ref{sec:buses}~节），还有第三种：\emph{反向通道}。有几种物质不是由发送者、而是由\emph{接收者}释放的，它们穿过突触间隙向回走，到达发送者的输出端，改变它下一次释放多少递质。在通信网络中，这类似于接收确认，发送者用它来调整自己的发送。神经元在改变连接权重时用的正是这个通道：发送者的输出端通过它得知接收者的活动，也就是第~\ref{sec:dofa}~章学习规则中的第二个因子。

已知神经递质的完整清单很长：一百多种物质，其中大多数是短的蛋白质链，即神经肽。但它们都可以归入少数几类，同一类中的物质行为相似。表~\ref{tab:othermediators} 收集了从不充当主要递质的物质，并回答工程师会问的同样三个问题：信号走哪条总线，作用快不快，通常是什么符号。它们不决定神经元的类型：它们要么与主要递质一起释放，要么不是由神经元释放的，要么是反向传递的。

\begin{table}[t]
\centering
\footnotesize
\setlength{\tabcolsep}{4pt}
\renewcommand{\arraystretch}{1.12}
\begin{tabular}{>{\raggedright}p{2.25cm}>{\raggedright}p{2.8cm}>{\raggedright}p{1.8cm}>{\raggedright}p{1.5cm}>{\centering}p{0.8cm}>{\raggedright\arraybackslash}p{4.3cm}}
\toprule
类别 & 物质 & 总线 & 速度 & 符号 & 在计算中的作用 \\
\midrule
氨基酸 & D-丝氨酸 & 局部 & 慢 & 与 & 谷氨酸受体的第二把钥匙：“与”条件 \\
\midrule
单胺 & 痕量胺 & 广播 & 慢 & $\pm$ & 单胺通道的微调 \\
\midrule
\multirow{2}{2.25cm}{嘌呤}
 & ATP & 寻址 & 快和慢 & $+$ & 共释放递质；神经元与胶质细胞之间的信号 \\
 & 腺苷 & 容积 & 慢 & $-$ & 随活动积累：负荷计数器~\cite{Dunwiddie2001} \\
\midrule
\multirow{8}{2.25cm}{神经肽（一百多种）}
 & 脑啡肽、内啡肽、强啡肽 & 容积 & 慢 & $-$ & 抑制传递 \\
 & P 物质 & 容积 & 慢 & $+$ & 缓慢增强 \\
 & 神经肽 Y & 容积 & 慢 & $-$ & 缓慢抑制 \\
 & 生长抑素 & 容积 & 慢 & $-$ & 缓慢抑制，与 GABA 共释放 \\
 & 血管活性肠肽（VIP） & 容积 & 慢 & $+$ & 去抑制：抑制那些抑制性神经元 \\
 & 胆囊收缩素 & 容积 & 慢 & $\pm$ & 在部分抑制性神经元中与 GABA 共释放 \\
 & 食欲素 & 广播 & 慢 & $+$ & 觉醒状态的稳定器 \\
 & 催产素、加压素 & 广播 & 慢 & $\pm$ & 行为的全局设置 \\
\midrule
\multirow{2}{2.25cm}{气体}
 & 一氧化氮 NO & 反向、容积 & 慢 & $\pm$ & 能穿过细胞膜；权重改变时的反向信号~\cite{Garthwaite2008} \\
 & CO、H$_2$S & 容积 & 慢 & $\pm$ & 作用研究很少 \\
\midrule
脂质 & 内源性大麻素（花生四烯酸乙醇胺、2-AG） & 反向 & 慢 & $-$ & 接收者减少发送者的释放：按权重的反馈~\cite{WilsonNicoll2001} \\
\bottomrule
\end{tabular}
\caption{不决定神经元类型的物质。\emph{总线}、\emph{速度}和\emph{符号}的含义同表~\ref{tab:types}；此外，容积总线指递质在释放点周围的细胞间隙中扩散，反向总线指从接收者回到发送者，局部总线指在释放点附近起作用。“与”是许可信号：它本身不引起电流，但没有它，另一个输入就不起作用，就像“与”逻辑门的第二个输入。神经肽几乎总是与主要递质一起释放，并且主要在脉冲频率较高时释放~\cite{vandenPol2012}。}
\label{tab:othermediators}
\end{table}

\chapter{符号说明}
\label{sec:notation}

神经元类型用主要神经递质的四字母代码表示：\nt{GLUT}——谷氨酸，\nt{GABA}——$\gamma$-氨基丁酸，\nt{GLYC}——甘氨酸，\nt{ACET}——乙酰胆碱，\nt{DOPA}——多巴胺，\nt{SERO}——血清素（5-羟色胺），\nt{HIST}——组胺，\nt{NORA}——去甲肾上腺素，\nt{ADRE}——肾上腺素，\nt{ASPA}——天冬氨酸（表~\ref{tab:types}）。

字母比物理量少，同一个字母在不同章节里有时表示不同的量。下表中每种含义各占一行；右栏给出引入它的公式。

\begin{center}
\footnotesize
\setlength{\tabcolsep}{4pt}
\renewcommand{\arraystretch}{1.12}
\begin{longtable}{>{\raggedright}p{2.0cm}>{\raggedright}p{9.6cm}>{\raggedright\arraybackslash}p{2.4cm}}
\toprule
符号 & 含义 & 引入处 \\
\midrule
\endfirsthead
\toprule
符号 & 含义 & 引入处 \\
\midrule
\endhead
\bottomrule
\endfoot
$a$ & 线性传递特性的斜率，$f(u)=au$ & \eqref{eq:feedback} \\
$a$ & \nt{ACET} 水平：$0$ 为读，$1$ 为写 & \eqref{eq:acetnet} \\
$a_i$, $a$ & 神经元群的疲劳：在节律发生器中；在慢波睡眠的跷跷板中 & \eqref{eq:halfcentre}, \eqref{eq:updown} \\
$a_S$, $a_P$ & 心率对油门和刹车的敏感度 & \eqref{eq:gasbrake} \\
$A(r)$, $A_0$ & 频率为 $r$ 时突触对一个脉冲的响应；休息充分的突触的响应 & \eqref{eq:depression} \\
$A_1$ & 接收者对一份递质的响应 & \eqref{eq:quantal} \\
$A$ & 神经元连同树突在内的膜面积 & \eqref{eq:dendriteload} \\
$A_f$, $A_s$ & 快过程和慢过程中保留到下一次运动的校正比例 & \eqref{eq:tworate} \\
$b_i$ & 起搏器的自身输入 & \eqref{eq:seroneuron} \\
$b$ & 节律发生器或睡眠跷跷板中，工作使神经元群疲劳的速度 & \eqref{eq:halfcentre}, \eqref{eq:updown} \\
$b$ & 每个数字化采样的比特数 & \eqref{eq:probedata} \\
$B_f$, $B_s$ & 快过程和慢过程中校正从误差中学习的强度 & \eqref{eq:tworate} \\
$c$ & 空间中的递质浓度 & \eqref{eq:seroneuron} \\
$c$ & 相邻神经元噪声的相关系数 & \eqref{eq:correlated} \\
$c$ & 计数之差移动球拍的系数 & \eqref{eq:dishreadout} \\
$c_f$ & $C$ 神经元的响应：所有位置上 $s_{f,p}$ 的最大值 & \eqref{eq:scstage} \\
$c_m$ & 膜的比电容 & \eqref{eq:dendriteload} \\
$C$ & 努力的代价：用时 $\tau_a$ 完成的动作代价为 $C/\tau_a$ & \eqref{eq:vigor} \\
$d'$ & 两个输入的可区分度：差值除以噪声 & \eqref{eq:gainsnr} \\
$d$, $d_j$ & 信号路径上的延迟 & \eqref{eq:glutneuron}, \eqref{eq:vstar} \\
$d$ & 肌肉反馈回路的延迟 & \eqref{eq:delaystable} \\
$D$ & 扩散系数 & \eqref{eq:diffusionlength} \\
$D$ & \nt{GLUT}--\nt{GABA} 网络频率表达式中的分母 & \eqref{eq:isn} \\
$D$ & 延迟奖励的等待时间 & \eqref{eq:patience} \\
$e$, $e_{ij}$ & 突触中的资格迹 & \eqref{eq:threefactor} \\
$e$, $e_i$ & 经误差导线传来的输出误差或运动误差 & \eqref{eq:lms}, \eqref{eq:adaptivefilter}, \eqref{eq:tworate} \\
$E_{\text{spike}}$ & 一个脉冲的能量，包括突触的功耗 & \eqref{eq:energy} \\
$E$ & \nt{GLUT} 群的频率 & \eqref{eq:ei} \\
$E_E$ & 兴奋性突触的平衡电位 & \eqref{eq:shunt} \\
$f$, $F$ & 传递特性：频率作为输入的函数 & \eqref{eq:ratenet}, \eqref{eq:gain} \\
$f$, $f_0$ & 心跳频率；起搏器的固有节律 & \eqref{eq:gasbrake} \\
$f_s$ & 单个电极的采样频率 & \eqref{eq:probedata} \\
$F(t)$, $F_1$ & 肌肉力量；单次抽搐的力量 & \eqref{eq:forcesum} \\
$F_1$, $F_2$ & 向相反方向拉动关节的两块肌肉的力 & \eqref{eq:antagonists} \\
$g$ & 受体对浓度的敏感度 & \eqref{eq:seroneuron} \\
$g_L$, $g_E$, $g_I$ & 泄漏、兴奋和抑制电导 & \eqref{eq:shunt} \\
$g$ & 由 \nt{NORA} 设定的增益 & \eqref{eq:gain}, \eqref{eq:machine} \\
$g$ & 记忆网络中来自 \nt{GABA} 的总抑制 & \eqref{eq:acetnet} \\
$g$ & 由上级设定的疼痛闸门增益 & \eqref{eq:gate} \\
$g$ & 激素级联中一级的增益 & \eqref{eq:cascade} \\
$g_j$ & 自适应滤波器输入端的第 $j$ 个滤波器，有自己的时间常数 & \eqref{eq:adaptivefilter} \\
$G$ & 反馈回路增益 & \eqref{eq:feedback} \\
$G$ & 带延迟回路中的校正速度 & \eqref{eq:delaystable} \\
$h$, $h_I$ & \nt{GLUT} 群和 \nt{GABA} 群的外部输入 & \eqref{eq:ei}, \eqref{eq:isn} \\
$h(t)$ & 肌肉的单次抽搐 & \eqref{eq:twitch} \\
$H(t)$ & 睡眠压力的上阈值：达到它，机器入睡 & \eqref{eq:sleeppressure} \\
$I$, $I_i$ & 神经元或神经元群的外部输入 & \eqref{eq:ratenet}, \eqref{eq:wta}, \eqref{eq:updown} \\
$I$ & \nt{GABA} 群的频率 & \eqref{eq:ei} \\
$I$ & 刺激的亮度或强度 & \eqref{eq:photonnoise}, \eqref{eq:nakarushton} \\
$I$ & 给膜充电的输入电流 & \eqref{eq:dendriteload} \\
$k$ & 簇放电中的脉冲数 & \eqref{eq:burst} \\
$k$ & 增量需超过噪声的倍数 & \eqref{eq:photonnoise} \\
$k$ & 激素级联的反馈系数 & \eqref{eq:cascade} \\
$k_s$ & 疼痛推高敏化的强度 & \eqref{eq:sensitization} \\
$K$ & 结果特征的个数 & \eqref{eq:vectortd} \\
$K$ & 关节刚度 & \eqref{eq:antagonists} \\
$K$ & 激素级联的回路增益，$K=g^3k$ & \eqref{eq:cascade}, \eqref{eq:cascadestable} \\
$K_\varphi$ & 昼夜振荡器向外部信号校准的强度 & \eqref{eq:pll} \\
$L$ & 延迟线的长度 & \eqref{eq:itd} \\
$L$ & 痛觉传感器导线的长度 & \eqref{eq:paintime} \\
$L(t)$ & 睡眠压力的下阈值：达到它，机器醒来 & \eqref{eq:sleeppressure} \\
$M$ & 预兆与奖励之间因果关系的度量 & \eqref{eq:retro} \\
$M_k$ & 特征 $k$ 的期望 & \eqref{eq:vectortd} \\
$M_{ik}$ & 来自 \nt{GABA} 神经元 $k$ 的抑制强度 & \eqref{eq:machine} \\
$M$ & 关节的转矩 & \eqref{eq:antagonists} \\
$M$ & 混合器的输入线数 & \eqref{eq:cover} \\
$n$, $\bar n$ & 一个窗口内的脉冲数；其平均值 & \eqref{eq:rateerror} \\
$n$ & 阈值元件的输入数 & \eqref{eq:cover} \\
$n_\uparrow$, $n_\downarrow$ & “上”区和“下”区在一个窗口内的脉冲计数 & \eqref{eq:dishreadout} \\
$N$ & 神经元数 & \eqref{eq:correlated}, \eqref{eq:energy} \\
$N$ & 探针的电极数 & \eqref{eq:probedata} \\
$N_s$ & 两个神经元之间的释放位点数 & \eqref{eq:quantal} \\
$p$ & 每个脉冲释放递质的概率 & \eqref{eq:burst} \\
$p$ & 记忆网络中活跃神经元的比例 & \eqref{eq:acetnet} \\
$p$ & 需要学会补偿的扰动 & \eqref{eq:tworate} \\
$P$ & 用于信号的功率 & \eqref{eq:energy}, \eqref{eq:dishpower} \\
$P$ & 所存估计的不确定度 & \eqref{eq:kalman} \\
$P$ & “刹车”的活动：通往心脏的 \nt{ACET} 导线 & \eqref{eq:gasbrake} \\
$P_{\max}$ & 阈值元件能分成两类的随机模式的数目 & \eqref{eq:cover} \\
$P_{\text{miss}}$ & 未接住的概率 & \eqref{eq:dishdensity} \\
$q$ & 世界变化的速度 & \eqref{eq:kalman} \\
$q_k$ & \nt{GABA} 神经元 $k$ 的频率 & \eqref{eq:machine} \\
$q_0$ & 抑制性中间神经元的背景活动 & \eqref{eq:gate} \\
$q$ & 疼痛闸门中抑制性中间神经元的频率 & \eqref{eq:gate} \\
$q$ & 下一个运动单位比前一个强多少倍 & \eqref{eq:sizeprinciple} \\
$r$, $r_i$ & 神经元的脉冲频率 & \eqref{eq:ratenet} \\
$r$, $r_t$ & 奖励（奖励流） & \eqref{eq:rpe}, \eqref{eq:ctd} \\
$\mathbf r(t)$ & 储备池的响应向量 & \eqref{eq:reservoir} \\
$R$, $\bar R$ & 获得的奖励；其期望或单位时间平均值 & \eqref{eq:perturbation}, \eqref{eq:vigor} \\
$R$ & 滤波器输入端噪声的方差 & \eqref{eq:kalman} \\
$R$, $r_0$ & 延迟奖励和即时奖励 & \eqref{eq:patience} \\
$R_{\max}$ & 极限传输速率，bit/s & \eqref{eq:spikeentropy} \\
$r_{\max}$ & 最大脉冲频率 & \eqref{eq:gain}, \eqref{eq:nakarushton} \\
$R_{\text{原始}}$, $R_{\text{脉冲}}$ & 探针的数据流：原始数据和仅脉冲，bit/s & \eqref{eq:probedata} \\
$s_j$ & 神经元 $j$ 输出的符号 & \eqref{eq:dalesign} \\
$s_i$ & 接收者受体的符号 & \eqref{eq:seroneuron} \\
$s$, $\hat s$ & 世界的状态；其估计 & \eqref{eq:rpe}, \eqref{eq:kalman} \\
$s_{f,p}$ & $S$ 神经元对位置 $p$ 处特征 $f$ 的响应 & \eqref{eq:scstage} \\
$S$ & “油门”的活动：通往心脏的 \nt{NORA} 导线 & \eqref{eq:gasbrake} \\
$S$ & 睡眠压力 & \eqref{eq:sleeppressure} \\
$t_1$ & 第一个脉冲的时刻 & \eqref{eq:latency} \\
$t_k$ & 第 $k$ 个脉冲的时刻 & \eqref{eq:forcesum} \\
$t_{f,i}$ & 特征 $f$ 的模板 & \eqref{eq:scstage} \\
$T$ & 脉冲计数窗口；光子积累时间 & \eqref{eq:rateerror}, \eqref{eq:photonnoise} \\
$T_c$ & 肌肉抽搐的上升时间 & \eqref{eq:twitch} \\
$T_{\text{burst}}$ & 培养物群放电之间的间隔 & \eqref{eq:burstinterval} \\
$u$ & 神经元的总输入 & \eqref{eq:gain} \\
$u$, $u(t)$ & 大脑的指令：自适应滤波器的输入；给激素级联的指令 & \eqref{eq:adaptivefilter}, \eqref{eq:cascade} \\
$u(t)$ & 储备池的输入 & \eqref{eq:reservoir} \\
$U$ & 恒定输入的水平 & \eqref{eq:latency} \\
$U$ & 每个脉冲释放的递质储备比例 & \eqref{eq:depression} \\
$v$ & 导线中信号的速度；光斑移动的速度 & \eqref{eq:itd}, \eqref{eq:vstar} \\
$V$ & 膜电压 & \eqref{eq:glutneuron}, \eqref{eq:shunt} \\
$V$ & 对未来奖励的期望 & \eqref{eq:rpe}, \eqref{eq:ctd} \\
$w$, $w_{ij}$, $W_{ij}$ & 连接权重 & \eqref{eq:glutneuron}, \eqref{eq:ratenet}, \eqref{eq:acetnet}, \eqref{eq:lms}, \eqref{eq:adaptivefilter} \\
$\mathbf w_k$ & $C$ 神经元 $k$ 的输入权重 & \eqref{eq:traceinvariance} \\
$w_{q\mu}$, $w_{q\nu}$, $w_{z\mu}$ & 疼痛闸门中的连接权重 & \eqref{eq:gate} \\
$W_{\text{out}}$ & 储备池线性读出的权重 & \eqref{eq:reservoir} \\
$x$, $y$ & 逻辑输入和输出 & \eqref{eq:dualrail}, \eqref{eq:veto} \\
$x$, $y$ & 发送者和接收者的活动 & \eqref{eq:hebb} \\
$x$ & 检测器在延迟线上的位置 & \eqref{eq:itd} \\
$x_i$ & 来自感觉器官的输入 & \eqref{eq:acetnet} \\
$x_p$ & 位置 $p$ 处的输入 & \eqref{eq:scstage} \\
$\mathbf x$ & $S$ 神经元的输出，即 $C$ 神经元的输入 & \eqref{eq:traceinvariance} \\
$x_j$ & 自适应滤波器的第 $j$ 个输入：经过滤波器 $g_j$ 的指令 & \eqref{eq:lms}, \eqref{eq:adaptivefilter} \\
$x_f$, $x_s$ & 快过程和慢过程的校正量 & \eqref{eq:tworate} \\
$x_1$, $x_2$, $x_3$ & 激素级联三级上的水平；$x_3$ 为最终激素 & \eqref{eq:cascade} \\
$x^*$ & 引发新雪崩所需的递质储备恢复比例 & \eqref{eq:burstinterval} \\
$y$ & 滤波器的含噪输入 & \eqref{eq:kalman} \\
$y$, $\hat y$ & 传感器信号；自适应滤波器对它的预测 & \eqref{eq:adaptivefilter} \\
$y_k$, $\bar y_k$ & $C$ 神经元 $k$ 的活动；其迹 & \eqref{eq:traceinvariance} \\
$y(t)$ & 储备池读出的输出 & \eqref{eq:reservoir} \\
$z$ & 疼痛闸门输出的频率 & \eqref{eq:gate} \\
\midrule
$\alpha_i^\pm$ & 正误差和负误差的权重 & \eqref{eq:asymmetric} \\
$\beta$ & 相互抑制的强度 & \eqref{eq:wta} \\
$\beta$ & 疼痛闸门输出受抑制的强度 & \eqref{eq:gate} \\
$\gamma$ & 折扣因子 & \eqref{eq:rpe} \\
$\delta(\cdot)$ & $\delta$ 函数：瞬时跳变 & \eqref{eq:glutneuron} \\
$\delta$, $\delta_t$ & 奖励预测误差 & \eqref{eq:rpe}, \eqref{eq:ctd} \\
$\delta t$ & 声音到达两耳的时间差 & \eqref{eq:itd} \\
$\Delta$, $\Delta_{\max}$ & 脉冲间隔；符合窗口 & \eqref{eq:window} \\
$\Delta t$ & 接收者分辨时间的精度 & \eqref{eq:spikeentropy} \\
$\Delta F$ & 下一个运动单位带来的力的增量 & \eqref{eq:sizeprinciple} \\
$\Delta I$ & 可分辨的亮度增量 & \eqref{eq:photonnoise} \\
$\Delta p$ & 一个窗口内球拍的位移 & \eqref{eq:dishreadout} \\
$\Delta u$ & 两个神经元群的输入之差 & \eqref{eq:gainsnr} \\
$\Delta V$ & 膜电压需要升高多少 & \eqref{eq:dendriteload} \\
$\Delta x$ & 相邻传感器之间的距离 & \eqref{eq:vstar} \\
$\varepsilon$ & 输入的相对差值 & \eqref{eq:wtagain} \\
$\eta$ & 学习率 & \eqref{eq:plasticity}, \eqref{eq:lms}, \eqref{eq:traceinvariance} \\
$\eta$ & 放大器之后网络中的噪声 & \eqref{eq:softmax} \\
$\mu$ & 触觉进入疼痛闸门的输入 & \eqref{eq:gate} \\
$\nu$ & 痛觉输入 & \eqref{eq:gate} \\
$\theta$ & 触发阈值 & \eqref{eq:window}, \eqref{eq:scstage} \\
$\kappa$ & 每个脉冲释放的递质量 & \eqref{eq:seroneuron} \\
$\kappa_i$ & 正误差与负误差权重之比 & \eqref{eq:expectile} \\
$\kappa_m$ & 肌肉力量与其刚度之间的关系 & \eqref{eq:antagonists} \\
$\ell$ & 递质扩散所及的距离 & \eqref{eq:diffusionlength} \\
$\ell$ & 肌肉拉动关节的力臂 & \eqref{eq:antagonists} \\
$\xi$ & 神经元输出的随机抖动 & \eqref{eq:perturbation} \\
$\rho(\boldsymbol\psi)$ & 停留在设置 $\boldsymbol\psi$ 的时间比例 & \eqref{eq:dishdensity} \\
$\sigma$ & 离散度，噪声尺度 & \eqref{eq:rateerror}, \eqref{eq:softmax} \\
$\sigma$ & 响应达到最大值一半时的亮度 & \eqref{eq:nakarushton} \\
$\sigma_{\text{pre}}$, $\sigma_{\text{post}}$ & 放大器之前和之后的噪声 & \eqref{eq:gainsnr} \\
$\tau$ & 膜时间常数 & \eqref{eq:glutneuron} \\
$\tau$ & 激素级联中一级的时间常数 & \eqref{eq:cascade} \\
$\tau_{\text{eff}}$ & 自我兴奋神经元群的时间常数 & \eqref{eq:perfectint} \\
$\tau_c$ & 递质在空间中的寿命 & \eqref{eq:seroneuron} \\
$\tau_E$, $\tau_I$ & \nt{GLUT} 群和 \nt{GABA} 群的时间常数 & \eqref{eq:ei} \\
$\tau_e$ & 资格迹的寿命 & \eqref{eq:threefactor} \\
$\tau_{\text{tr}}$ & $C$ 神经元活动迹的寿命 & \eqref{eq:traceinvariance} \\
$\tau_s$ & 损伤后敏感度恢复的时间 & \eqref{eq:sensitization} \\
$\tau_r$, $\tau_d$ & 清醒时睡眠压力积累的时间和睡眠中其释放的时间 & \eqref{eq:sleeppressure} \\
$\tau_{\text{rec}}$ & 递质储备的恢复时间 & \eqref{eq:depression} \\
$\tau_\gamma$ & 规划视野 & \eqref{eq:ctd} \\
$\tau_v$ & 期望更新的速度 & \eqref{eq:ramp} \\
$\tau_a$ & 节律发生器或睡眠跷跷板中的疲劳时间 & \eqref{eq:halfcentre}, \eqref{eq:updown} \\
$\tau_a^*$ & 完成动作的最佳时间 & \eqref{eq:vigor} \\
$\Phi$ & 学习规则的右端 & \eqref{eq:plasticity} \\
$\Phi$ & 光子通量 & \eqref{eq:photonnoise} \\
$\varphi_k$ & 所得结果的特征 $k$ & \eqref{eq:vectortd} \\
$\varphi$ & 昼夜振荡器与外部信号的相位差 & \eqref{eq:pll} \\
$\boldsymbol\psi$ & DishBrain 模型中网络的设置 & \eqref{eq:dishdensity} \\
$\omega$, $\omega_0$ & 外部信号（光）的频率；昼夜振荡器的固有频率 & \eqref{eq:pll} \\
\end{longtable}
\end{center}

\chapter{答案}
\label{sec:answers}

各章末习题的简要答案。完整解答见单独的教师用书。

\answer{ex:01:1}{(a)~$8\cdot10^{10}$ 个\nt{GLUT}，相当于 $10$ 个地球人口；广播神经元 $\approx1.9\cdot10^6$ 个，相当于维也纳大小的城市。(b)~末梢共 $\approx8\cdot10^{14}$ 个，其中广播型 $\approx3\cdot10^{11}$ 个，占比 $\approx4\cdot10^{-4}$，是这些神经元所占比例（$2.2\cdot10^{-5}$）的 $\approx16$ 倍。}
\answer{ex:01:2}{(a)~$\tau_c\ln(c_0/c_*)\approx4.6\ms$。(b)~当 $T\ge\tau_c\ln101$ 时线路才空闲：每秒最多 $\approx22$ 次释放，而不是 $\approx220$ 次。}
\answer{ex:01:3}{$V(s_k)=\gamma^{K-1-k}r$；$\delta_{\text{灯}}=\gamma^Kr=re^{-T/\tau_\gamma}\approx0.60\,r$，与 $\Delta$ 无关，$\tau_\gamma\approx1.95\,\text{s}$。}
\answer{ex:01:4}{$n^*=93$、$47$、$25$、$12$：$\alpha$ 较小时 $n^*\approx5/\alpha$。}
\answer{ex:01:5}{\nt{GLUT}：$1\to10\to10^4\to10^7$，$\approx10^7$ 个神经元，$4$ 级，$4\ms$。\nt{DOPA}：$1\to10\to3.2\cdot10^5$，$\approx3\cdot10^5$ 个神经元，$3$ 级，少 $30$ 倍；与\nt{DOPA}神经元的数目同一量级。}
\answer{ex:01:6}{$\lambda=f_EJ_E-f_IJ_I$，由此 $J_I=37.5$；电流之比 $f_IJ_I/f_EJ_E=0.94$；抑制只要减弱 $6.7\,\%$，网络就进入雪崩（$\lambda\ge1$）。}
\answer{ex:01:7}{每种状态下每个延迟需 $n\ge57$ 次呈现。同样的衰减也可以来自别的机制，例如内部时间估计的误差随 $T$ 增大。}

\answer{ex:02:1}{阵列长 $3.2$~mm，间距 $25\um$：$\approx129$ 个检测器。发放的是宽 $v\,\tau\ln2=1.7$~mm 的一条带，$\approx69$ 个检测器，超过阵列的一半；方向取这条带的中心。}
\answer{ex:02:2}{窗口为 $-\tau\ln\frac{w_2}{\theta-w_1}\le\delta\le\tau\ln\frac{w_1}{\theta-w_2}$，即 $[-0.235,\,0.178]\ms$；中心 $c=\frac\tau2\ln\frac{w_1(\theta-w_1)}{w_2(\theta-w_2)}=-28\,\mu$s。方向向较强的耳朵一侧偏移 $28\,\mu$s，接近三个分辨步长。}
\answer{ex:02:3}{需要满足 $s_iw_{ij}s_j\ge0$ 的 $s_i$；当且仅当所有环都是偶环时它们存在。相互抑制可以化为这种网络（没有持续振荡），\nt{GLUT}--\nt{GABA} 环则不能（可能振荡）。}
\answer{ex:02:4}{六个神经元 $g_k=[x+y+z\ge k]$ 和 $\bar g_k=[\bar x+\bar y+\bar z\ge4-k]$，$k=1,2,3$；$C=g_2$，$\bar C=\bar g_2$，$S=[g_1+\bar g_2+g_3\ge2]$，$\bar S=[\bar g_1+g_2+\bar g_3\ge2]$：共八个神经元。用与、或元件需 $12$ 个。}
\answer{ex:02:5}{$T_{\min}(0.6)=0.718\,\tau$；$I_c=0.225$，在阈值附近 $T_{\min}\propto(I-I_c)^{-1/2}$。}
\answer{ex:02:6}{(a)~$500$ 环，$5\cdot10^4$ 个神经元；散布 $4.5\ms$（$0.45\,\%$），相对误差 $\propto T^{-1/2}$。(b)~所有环共有的随机延迟因子，散布为 $5\,\%$。}
\answer{ex:02:7}{每 $\tau_F\ln2=0.69\,\text{s}$ 刷新一次：每个神经元每秒 $4.3$ 个脉冲，而触发器是 $20$ 个，节省 $4.6$ 倍（每个脉冲 $10^{-10}$~J 时，$4\cdot10^{-7}$~J 对 $2\cdot10^{-6}$~J）。触发器丢失比特；电容器只要压制开始于刷新后不超过 $0.49\,\text{s}$，就能保住比特（概率 $0.71$）。}

\answer{ex:03:1}{$g_E=g_L$ 时 $35\to17.5\mV$；$g_E=4g_L$ 时 $56\to40\mV$，而减法会给出 $38.5\mV$：分流把 $g_E$ 除以 $3$。窗口收窄一半（$\tau_{\text{eff}}$ 减小到 $1/2$）。}
\answer{ex:03:2}{$\operatorname{tr}=\frac{w_{EE}-1}{\tau_E}-\frac1{\tau_I}$，$\det=\frac{1-w_{EE}+w_{EI}w_{IE}}{\tau_E\tau_I}$；在边界上 $\omega=\sqrt{\det}$，$\tau_I=20\ms$，周期 $73\ms$。}
\answer{ex:03:3}{$d_c=\tau_E\arccos(-1/G)/\sqrt{G^2-1}$，周期 $2\pi\tau_E/\sqrt{G^2-1}\in(2d_c,4d_c)$。$G=2$：$d_c=12.1\ms$，周期 $36\ms$；$G=5$：$d_c=3.6\ms$，周期 $12.8\ms$，正好落在快速节律的频带内。}
\answer{ex:03:4}{(a)~上升时在 $I_2=\beta I_1=2$ 处，下降时在 $I_2=I_1/\beta=0.5$ 处，滞回环宽 $1.5$。(b)~$\varepsilon$ 每缩小一个数量级，增加 $\tau\ln10/(\beta-1)=2.3\tau$。}
\answer{ex:03:5}{三个中介，$d_k=5$、$10$、$15\ms$：否决必须覆盖 $15\ms$，每个覆盖 $5\ms$。}
\answer{ex:03:6}{$n+M+1=3+4+1=8$ 个神经元。两个神经元的方案：\nt{GABA} $=[x+y+z\ge2]$，输出 $=[x+y+z-2\,\nt{GABA}\ge1]$。}
\answer{ex:03:7}{$\binom84=70<100\le\binom94=126$：$9$ 个中介（斯珀纳定理）。没有直接连接时需 $100$ 个：$\beta(J-I)$ 的秩为 $n$。}
\answer{ex:03:8}{$h_c=\frac1{2w_{EE}}+\frac{w_{EI}w_{IE}-w_{EE}}{4w_{EE}^2}=\frac7{18}\approx0.39$；模拟显示符号在 $h=0.38$ 和 $0.40$ 之间改变。}

\answer{ex:04:1}{(a)~$\approx1.1\cdot10^{11}$ bit/J。(b)~$\log_2\sqrt{10}\approx1.7$ 比特对 $81$ 比特：少 $\approx50$ 倍。}
\answer{ex:04:2}{$r^*=(P_0/E_{\text{spk}})\ln\bigl(1/(r^*\Delta t)\bigr)$，$P_0/E_{\text{spk}}=1.16\,\text{s}^{-1}$：$r^*\approx6$ 个脉冲每秒，$7.4\cdot10^{10}$ bit/J。}
\answer{ex:04:3}{$\mathcal I=\frac{N}{1+(N-1)c}=9.9$（极限为 $10$），而 $\mathcal I=\frac{N}{1-c}=1111$：相差 $\approx110$ 倍；只有与信号相似的相关才有害。}
\answer{ex:04:4}{$\theta\approx68.7$ 和 $19.9$（以 $w$ 为单位），$m=19$ 和 $m=10$：快接收者所需的符合输入几乎少一半，尽管它的平均输入只有五分之一。}
\answer{ex:04:5}{$U\tau_{\text{rec}}\ge0.725\,\text{s}$ 且 $\tau_{\text{rec}}(1-2U)\le0.05\,\text{s}$，这要求 $U\ge0.483$；最小的 $\tau_{\text{rec}}=0.725\,\text{s}$，此时 $U=1$（$U=0.5$ 时为 $1.45\,\text{s}$）。}
\answer{ex:04:6}{(a)~$\theta\,e/(e-1)\approx1.58\,\theta$，与 $\tau$ 和 $\theta$ 无关。(b)~$14$ 个脉冲。}
\answer{ex:04:7}{(a)~每个词 $1.62$ 比特：$0.77$ 来自脉冲数，$0.84$ 来自它们的排列。(b)~$1.13$ 和 $1.59$ 比特：用 $30$ 个词的估计损失近三分之一。}

\answer{ex:05:1}{$\ell\approx17\um$；$\approx100$~天：向全脑广播靠的是导线分支，扩散只是把每个释放位点模糊到$\sim20\um$。}
\answer{ex:05:2}{$r=ab/(1+G)$，$S=1/(1+G)$精确成立；$+1.7\,\%$，而不是$+20\,\%$。}
\answer{ex:05:3}{$G=2$时$T^*=1.21\,\text{s}$，$G=9$时$0.19\,\text{s}$：实际上$G\sim2$--$4$，抑制$3$--$5$倍。}
\answer{ex:05:4}{总频率为$\lambda$时$\mathrm{CV}=1/\sqrt{2\lambda\tau_c}\approx9\cdot10^{-4}$；单个神经元需要$\tau_c=12.5\,\text{s}$。}
\answer{ex:05:5}{两种情况下都是$c\approx0.40$；$\mathrm{CV}=0.050$对$0.158$，相差$\sqrt{10}$倍。各神经元的频率仍与$a_i$成正比：受调节的只是总和。}
\answer{ex:05:6}{$N_1=\overline{A\vee B}$，$N_2=\overline{A\vee N_1}$，$N_3=\overline{B\vee N_1}$，$N_4=\overline{N_2\vee N_3}$，XOR $=N_5=\overline{N_4}$；每次切换用时$\tau_c\ln2$，经过$4$个门的路径使一次XOR需要$2.8\,\text{s}$。}
\answer{ex:05:7}{(a)~$\tau_\gamma=6/\ln3\approx5.5\,\text{s}$。(b)~$\bar D<4\,\text{s}$，而固定延迟为$D<2.2\,\text{s}$：延迟不可预测反而使动物更有耐心。}
\answer{ex:05:8}{$k_2=1/\tau_d=10\,\text{s}^{-1}$，$k_1=1/\tau_d-1/\tau_\gamma=9.5\,\text{s}^{-1}$；$\tau_\gamma=1/(k_2-mk_1)$：$m$变化$+2.6\,\%$时视野加倍（变化$+5.3\,\%$时视野变为无穷大）。}

\answer{ex:06:1}{$e=(1-e^{-T_a/\tau_e})e^{-d/\tau_e}$。(a)~$0.110$对$4.5\cdot10^{-5}$，相差$\approx2.5\cdot10^3$倍。(b)~$\tau_e^*=T_a/\ln(1+T_a/d)=0.59\,\text{s}$。}
\answer{ex:06:2}{速率为每秒$\nu_x\nu_y(A_+\tau_+-A_-\tau_-)=0.8A_+$，一小时$\approx2900A_+$；$\tau_-=\tau_+/0.6=33\ms$。}
\answer{ex:06:3}{当$\mu^2+s_1^2>s_2^2$时学到$\mathbf e_1$；这里$1.01>0.25$，神经元学到$\mathbf e_1$，尽管沿$\mathbf e_2$的离散大$25$倍：规则找到的是相关矩阵而非协方差矩阵的主向量。}
\answer{ex:06:4}{灯与果汁之间$V=Re^{-(t_c+L-t)/\tau_\gamma}$，因此$\dot V=V/\tau_\gamma$；爆发面积为$Re^{-L/\tau_\gamma}$：$0.61R$和$0.78R$。}
\answer{ex:06:5}{信噪比为$1/(N+1)$，尝试次数$\propto N+1$：$5\cdot10^5$次尝试$\approx1$~个月，$5\cdot10^{11}$次尝试$\approx8\cdot10^4$~年。}
\answer{ex:06:6}{(a)~$k_2=1/\tau_d$，$k_1=1/\tau_d-1/\tau_\gamma$。(b)~虚假误差为$-(V/\tau_\gamma)\,\varepsilon/(1+\varepsilon)$，$\varepsilon=\tau_d/\tau_\gamma$，由此$\tau_d\le0.105\,\text{s}$。}
\answer{ex:06:7}{$0.46$、$0.90$、$0.99$、$0.95$，$d=3\,\text{s}$时为$0.70$。各点落在同一条关于资格迹比例$F=(1-e^{-T_{\text{cue}}/\tau_e})e^{-d/\tau_e}$的曲线上：$F\gtrsim0.1$时能学会，$F<10^{-2}$时为随机选择。}
\answer{ex:06:8}{$\eta^*=1/\bigl(2\sigma^2(N+2)\bigr)$，需$N+2$次尝试；$N$个中只有$m$个抖动时需$N(m+2)/m\ge N+2$次：稀疏化没有帮助。}

\answer{ex:07:1}{$V=\kappa p/\bigl(\kappa p+(1-\kappa)(1-p)\bigr)$：$0.059$、$0.20$、$0.50$；中位数为$0$，而点~\eqref{eq:expectile}随$\kappa$平滑变化。}
\answer{ex:07:2}{$p=1/3$，$x=0.75$，$\sigma_r=0.35$，$V(0.2)=0.083$。}
\answer{ex:07:3}{$V=\sqrt\kappa/(\sqrt\kappa+\sqrt{1-\kappa})$，从$0.25$到$0.75$；离散$\propto\sqrt\eta$；两滴情形下$V=0.25+0.5\kappa$，两种分布在两端的神经元上相差$0.05$，需要$\eta\lesssim0.01$。}
\answer{ex:07:4}{（a）~$J^*=2\sqrt{C\bar R}$。（b）~多付$\bigl(k^{1/2}+k^{-1/2}\bigr)/2-1$：$k=2$时为$6\,\%$，$k=4$时为$25\,\%$——“两倍以内”的精度就够了。}
\answer{ex:07:5}{尖峰面积为$R(e^{-1}-e^{-2})\approx0.23R$——相当于$2.3\,\text{s}$的上升；如果多巴胺报告$V$，就没有尖峰，电平只是台阶式地增大为$e$倍。}
\answer{ex:07:6}{事件使权重移动$\propto A\tau_f/(\tau_e+\tau_f)$，电平带来误差$s\tau_f$：$9\,\text{s}\le\tau_f\le17\,\text{s}$。}
\answer{ex:07:7}{初始时$\Delta P=0.80$，$M=0.90$。（a）~$\Delta P=0.74$（$-8\,\%$），$M=0.43$（$-52\,\%$）；（b）~$\Delta P=0.40$（$-50\,\%$），$M=0.80$（$-11\,\%$）：双重分离。}
\answer{ex:07:8}{$N_{\text{eff}}\to2+\rho^2N\approx10^6$次尝试：比单根导线少$10^4$倍，但$1\,\%$的泄漏仍要花掉一百万次尝试。}

\answer{ex:08:1}{（a）~$\approx1.7\cdot10^6$。（b）~$g=\frac{0.9}{\sqrt{0.19}}\,\sigma_{\text{后}}/\sigma_{\text{前}}\approx16.5$：答案只取决于两种噪声之比。}
\answer{ex:08:2}{（a）~特征值为$-(1\pm g\beta)/\tau$；差在$\tau/(1-g\beta)$内衰减：$g\beta=0.5$时为$40\ms$（输入之差放大三倍），$0.9$时为$200\ms$（放大$19$倍）。（b）~$0.905$和$1.105$；在两个边界之间只有强输入能获胜。}
\answer{ex:08:3}{$P(k)=e^{gu_k/\sigma}/\sum_le^{gu_l/\sigma}$，即~\eqref{eq:softmax}；$g=10\ln9\approx22$。}
\answer{ex:08:4}{$\bar\Delta=1/3$，$g^*\approx3\ln(1/h)$：$21$、$14$、$9$，模型中为$16$--$23$、$11$和$8$。}
\answer{ex:08:5}{$g^*$从$h=0.001$时的$16$--$23$降到$h=0.2$时的$6$--$8$；当$h\le0.05$时，估计$3\ln(1/h)$的误差在一倍半以内；当$h\gtrsim\eta/10$时，规则$Q\leftarrow Q+\eta(R-Q)$来不及吸收探索所得，$g^*$停在$8$左右。}
\answer{ex:08:6}{$g_{\text{lo}}=0$时$T_p=\tau\ln9=44\ms$（模型为$44\ms$），$g_{\text{lo}}=0.25$时为$70\ms$：增益接近零时的一次$2$--$3\tau$的簇放电。}
\answer{ex:08:7}{（a）~最佳恒定值$0.925$（$g=8$），先知$0.929$（$16/8$）——几乎没有什么可赢的。（b）~例如，奖励稀少的平稳时期（$p\in[0,0.3]$）与$h=0.1$的多变时期：$0.850$对$0.872$；$\eta=0.2$时调节拿到约四分之一，$\eta=0.05$时几乎全部拿到。}

\answer{ex:09:1}{(a)~$P_\infty=0.2$，记忆$20\,\text{s}$。(b)~$P_\infty=1$，记忆$1\,\text{s}$。(c)~记忆与噪声幅度成正比：六倍。}
\answer{ex:09:2}{$P(t)=P_\infty\coth(t\sqrt{q/R})$。(a)~$\frac12\sqrt{R/q}=5\,\text{s}$——比记忆短一半。(b)~$t=\frac12\ln21\,\sqrt{R/q}\approx15\,\text{s}$：升高的\nt{ACET}应当维持这么久。}
\answer{ex:09:3}{方差为$qT/2+R/(2T)$；$T^*=\sqrt{R/q}$，方差为$\sqrt{qR}=P_\infty$，与滤波器~\eqref{eq:kalman}相同；增大$(\kappa+1/\kappa)/2$倍：$1.25$和$5.05$。}
\answer{ex:09:4}{$a>a_w=1-\theta/(mw)$：$w=0.041$时为$0.15$，$w=0.045$时为$0.22$。}
\answer{ex:09:5}{在$M\approx40$--$60$时开始出错，$M=80$时多出$1.9$个神经元，$M=100$时模式比例为$0.86$；来自其他模式的干扰方差$\approx(M-1)p(1-p)/N$，所以$M_{\max}\approx80$，$M_{\max}\propto N/p$。}
\answer{ex:09:6}{例如$\eta(a)=\eta\min\bigl(1,\max(0,(a-0.3)/0.6)\bigr)$。用$\eta a$时三个外来神经元全部进入模式；用$\eta(a)$时一个也没有，$a\ge0.9$时的写入不变（比例$1.00$）。}
\answer{ex:09:7}{(a)~$40$次回放；用习题~\ref{ex:09:6}的$\eta(a)$则永远不够。(b)~$200$次回放，$10$个夜晚。}

\answer{ex:10:1}{总线$\approx10^{12}$\,bit/s（每个脉冲$11.4$\,bit），控制系统$\approx40$\,bit/s；需$2.5\cdot10^{10}\,\text{s}$，约八百年。}
\answer{ex:10:2}{$a>1/3$（$a=0$时出现约$67$\,Hz的增幅振荡）。不能：稳定性由乘积$g(1-a)$决定。}
\answer{ex:10:3}{均值$\sigma^2R'x_j$，方差$(N+1)\,x_j^2\sigma^4R'^2$；信噪比$1/(N+1)$，需要$n\approx z^2(N+1)\propto N$次尝试。}
\answer{ex:10:4}{$L_v\in[32.9,\,47.9]\ms$：延迟取$35.4\ms$，窗口$\pm7.5\ms$；错误绑定$2\Delta_{\max}r_ar_v\approx0.38$次/秒。}
\answer{ex:10:5}{$0.098$，而最佳$\alpha=0.2$时为$0.180$：小$45\,\%$。}
\answer{ex:10:6}{$n=\ln\bigl(\ln1.5/(0.6g)\bigr)/\ln(1-\alpha)$：$24$次和$4$次；$g=2$时为$11$次和$2$次。}
\answer{ex:10:7}{$S'=S\bigl(1-4\eta\sigma^2+4\eta^2\sigma^4(G+2)\bigr)$，最佳$\eta\sigma^2=1/(2(G+2))$；约$10^6$次尝试，约一个月。}

\answer{ex:11:1}{(a)~$T=k^2/\bigl(\Phi(\Delta I/I)^2\bigr)=9\,\text{s}$。(b)~$45$个传感器，光斑$\approx7\times7$，分辨率下降约$7$倍。}
\answer{ex:11:2}{(a)~每个脉冲$2$--$3.3$\,bit，为上限（$9.1$--$9.8$\,bit）的$22$--$34\,\%$。(b)~$\log_2\binom{400}{20}\approx111$\,bit；平均有$1$种共同类型。}
\answer{ex:11:3}{(a)~$\sigma/9\le I\le9\sigma$，$1.9$个数量级。(b)~$6$个和$7$个。}
\answer{ex:11:4}{(a)~$f<343/0.44\approx780$\,Hz。(b)~$625$个脉冲，$125$个神经元。}
\answer{ex:11:5}{每个事件$21$\,bit，每秒$4.8\cdot10^5$个事件，每个边缘像素$7$个事件：每秒$136$像素。普通相机每秒$1.25$帧。}
\answer{ex:11:6}{按对数，两个方向都$\approx1.0\,\text{s}$（理论值$0.96\tau$）；按线性，向上$0.2\,\text{s}$，向下$3.0\,\text{s}$（理论值$0.18\tau$和$3.0\tau$）。}
\answer{ex:11:7}{(a)~$\ln I$服从中位数为$\ln\sigma$、尺度为$1$的逻辑斯谛分布：自然对数下的离散度为$\pi/\sqrt3$，即$0.79$个数量级。(b)~$r=r_{\max}\Phi_I(I)^2$。}

\answer{ex:12:1}{$0.60$（若无相关则为$0.18$），极限$\sqrt{c/(rT)}\approx0.58$：一百个神经元只能区分“有/无”。}
\answer{ex:12:2}{(a)~$10^8$个脉冲，$\approx10\,\text{mJ}$：占全脑$3$\,J的$0.3\,\%$。(b)~$3$\,J，约为其$300$倍。}
\answer{ex:12:3}{(a)~$\theta_A\in[2,3)$，$\theta_B\in[1,2)$；另一个图形给出的输入分别不超过$2$和$1$。(b)~$10\cdot96^2\approx9.2\cdot10^4$。}
\answer{ex:12:4}{(a)~$T_d=0.85\,\tau_a=0.34\,\text{s}$。(b)~$T_d\propto\tau_a$，与$I$无关：$\tau_a\approx3.5\,\text{s}$（模拟：$3.1\,\text{s}$）。}
\answer{ex:12:5}{$p\approx0.17$。正确率下降缓慢：$N=8$、$12$、$24$、$48$、$96$、$192$时分别为$0.90$、$0.88$、$0.86$、$0.84$、$0.82$、$0.79$。}
\answer{ex:12:6}{$\tau_{\text{tr}}=20$、$50$和$200\ms$时十次运行全部无误（$30\ms$时十次中有一次失败）；$5$--$10\ms$时约为$0.5$，$0.5$--$2\,\text{s}$时质量降到$0.86$。上限：资格迹须在停顿期间消退，$e^{-0.3\,\text{s}/\tau_{\text{tr}}}\ll1$。}
\answer{ex:12:7}{一个延迟不够：$\pm5\ms$的窗口覆盖不了$5$到$20\ms$的$\Delta x/v$。需两条抑制支路，$5<d_1\le10\ms$，$15\le d_2\le d_1+10\ms$。}
\answer{ex:12:8}{\texttt{two\_stage}：$\Delta=2$，翻转一个点得到平局，翻转两个点才改变答案。数1的分类器：翻转一个点即可；$p=0.1$时为$0.57$。}

\answer{ex:13:1}{(a)~$q=100^{1/99}\approx1.048$，步长$\approx4.8\%$；范围$\approx2.2\cdot10^3$（$67$\,dB）。(b)~步长$22f_1$，即力为$10f_1$时为$220\%$。}
\answer{ex:13:2}{$P_{\text{fail}}\approx1.8\cdot10^{-4}$和$2.8\cdot10^{-16}$：$S=2$时每天约$300$次失败，$S=4$时每天$5\cdot10^{-10}$次。}
\answer{ex:13:3}{(a)~$H(s)=eF_1T_c/(1+sT_c)^2$，二阶滤波器，截止频率$1/(2\pi T_c)\approx3.2$\,Hz。(b)~$r\approx22$次/秒（按一次谐波$\approx20$）。}
\answer{ex:13:4}{(a)~$Gd=1/e$时根$s=-1/d$为二重根，而任何其他$G$都不能使最右侧的根更靠左。(b)~$d=30\ms$：$G\approx12$\,s$^{-1}$，$\tau=30\ms$；$d=150\ms$：$G\approx2.5$\,s$^{-1}$，$\tau=150\ms$——时间常数等于延迟。}
\answer{ex:13:5}{(a)~$E\approx0.427$，$T\approx1.70\,\tau_a=0.68$\,s（$\tau=20\ms$的模型给出$0.84$\,s）。(b)~方程中不含$I$；当$b>\beta-1$时存在节律。}
\answer{ex:13:6}{$b=0.8$时没有节律；$b=1.05$、$1.2$、$1.5$、$2$、$4$时$T\approx2.73$、$1.74$、$1.19$、$0.84$、$0.45$\,s；任何$I$下$T=0.840$\,s：输入只改变幅度，不改变步频。}
\answer{ex:13:7}{代换$s=u-a$给出最大衰减速率$a+1/d$，此时$G=e^{-1-ad}/d$；$a\ge33.3$\,s$^{-1}$，$G=e^{-2}/d\approx4.5$\,s$^{-1}$；共同收缩越强（$a$越大），$G$应取得越小。}
\answer{ex:13:8}{开放题。例如，在疲劳驱动中加入阈值$b[r_i-r_0]_+$，在$b=4$、$r_0=0.2$时得到$I_{\min}=\beta b r_0/(b-\beta+1)\approx0.53$，周期从$I=0.6$时的$1.8$\,s降到$I=3$时的$0.52$\,s。}

\answer{ex:14:1}{(a)~分别拉长$0.17$\,s和$1.5$\,s。(b)~距离小于$11$\,cm时。}
\answer{ex:14:2}{只要$g\le\beta w_{q\mu}/w_{z\mu}=5$，无痛的触摸在任何$\mu$下都不通过；$g>6$时触摸$\mu=1$能通过：强烈的敏化使普通的触摸变成疼痛。}
\answer{ex:14:3}{(a)~$g^*=(1-k_sC)/(1-k_sA)$，当$k_sA<1$时稳定。(b)~无触摸时为$\nu\ge2$，有触摸时$\nu\ge1.7$即失控。}
\answer{ex:14:4}{$V(t)=-Pe^{-(T-t)/\tau_\gamma}$。(a)~$-Pe^{-T/\tau_\gamma}\approx-0.37P$。(b)~$+Pe^{-(T-t_e)/\tau_\gamma}\approx+0.61P$；它随$t_e$增大到$+P$，教会网络回避疼痛。}
\answer{ex:14:5}{$k_s=4$时不会自锁（$k_s\ge4.58$时才有第二个稳定状态）；$k_s=4.6$、$5$、$6$、$8$时$T_{\min}\approx4.35$、$1.40$、$0.65$、$0.32\,\tau_s$。}
\answer{ex:14:6}{$w_{q\nu}=4/9\approx0.444$，$q_0=2/9\approx0.222$，$w_{q\mu}=49/45\approx1.089$；无痛的触摸在$g=49/9\approx5.4$以内是安全的。}
\answer{ex:14:7}{(a)~当$\nu_c\ge q_0/w_{q\nu}$时阈值$\nu_c(g)=\theta/g$，否则为$(\theta+\beta q_0)/(g+\beta w_{q\nu})$；$g^*\approx2.45$、$1.0$、$0.25$。(b)~开放问题。}

\answer{ex:15:1}{(a)~每个神经元$2\cdot10^5$个，整个线路$6\cdot10^{12}$个。(b)~每根导线$\approx8$\,bit/s，至少$2.5\cdot10^4$\,s，约$7$小时。}
\answer{ex:15:2}{倍数$2$：相邻相关$0.943$，$\lambda$从$4.18$到$7.3\cdot10^{-4}$，比值$\approx5.7\cdot10^3$；倍数$4$：相关$0.8$，比值$\approx94$。}
\answer{ex:15:3}{(a)~$\lambda=0.988$和$0.808$：$82$次和$4.7$次动作。(b)~$x_s^*=0.690$，$x_f^*=0.172$，总和$0.862$。}
\answer{ex:15:4}{(a)~$116$次动作。(b)~$19$次。单过程模型在总和为零时没有任何节省。}
\answer{ex:15:5}{(a)~$G=50$\,s$^{-1}$，而无模型时的极限为$10.5$\,s$^{-1}$。(b)~不稳定：$G=50$\,s$^{-1}$时临界延迟$d_c\approx103\ms$；$d=150\ms$时需要$\hat k>0.445k$（对任何$G$和$d$，需要$\hat k>k/2$）。}
\answer{ex:15:6}{误差在$6$--$30$\,s内降到$0.5\%$以下，底限$\approx(6$--$9)\cdot10^{-4}$，最佳权重下为$7.5\cdot10^{-4}$；$\eta=50$时，沿$C$的病态方向，权重与最佳值仍相差可达$0.2$（习题~\ref{ex:15:2}）。}
\answer{ex:15:7}{$\mathbf B=A\mathbf K$；$q_f/R\approx0.035$，$q_s/R\approx8.6\cdot10^{-4}$；$q_s$增至四倍时$B_f\approx0.088$，$B_s\approx0.047$：慢过程学得快了一倍多，快过程则变慢。}

\answer{ex:16:1}{血液是 $\tau=t_{1/2}/\ln2=14.4$~min的RC电路；最大值与最小值之比 $2^{T/t_{1/2}}=64$；基波衰减到 $0.55$；$t_{1/2}=T=60$~min时比值为 $2$。}
\answer{ex:16:2}{$\varphi^*=\arcsin\bigl((\omega-\omega_0)/K_\varphi\bigr)\approx41$~min（$\omega-\omega_0\approx0.047$~rad/d），弛豫时间 $1/(K_\varphi\cos\varphi^*)\approx3.9$~天；跳变 $+6$ 和 $-6$~h后分别需 $7.7$ 和 $7.9$~天。}
\answer{ex:16:3}{$K_c(n)=\sec^n(\pi/n)$：$8$ 和 $2.37$，误差 $11\%$ 和 $30\%$；$n\to\infty$ 时 $K_c\to1$，误差 $\to50\%$。}
\answer{ex:16:4}{极限指令（油门 $80$，刹车 $0$）经 $0.76$~s达到 $f=120$；“只用油门”需 $2.33$~s，慢三倍。}
\answer{ex:16:5}{$3\arctan(\omega\tau)+\omega d=\pi$，$K_c=(1+\omega^2\tau^2)^{3/2}$：$K_c=8$、$3.54$、$2.50$、$1.76$，周期 $3.63$、$5.46$、$6.86$、$9.27\,\tau$；$0.9K_c$ 时振荡衰减，$1.1K_c$ 时振荡稳定持续。}
\answer{ex:16:6}{$0.60\bar c$、$0.87\bar c$、$0.90\bar c$，极限 $0.91\bar c$；浓度恒定时响应为零：这样的接收方只读取剂量。}
\answer{ex:16:7}{开放问题。在反馈机制下，只有 $K>8$ 时才有节律，周期正比于各级的 $\tau$，几乎与 $K$ 无关（$K=8.5$ 时 $3.64\,\tau$，$K=40$ 时 $4.23\,\tau$）：$K$ 改变一点五倍只使周期偏移 $2$--$4\%$，小于测量误差。能作出判决的是断开环路（在第一级输入端维持恒定的终端激素会消灭这种节律），或者观察在周期不同相位加入少量激素后的响应。}

\answer{ex:17:1}{$0.2$~W 对 $81$~\textmu W；原始数据流要求每比特不超过 $50$~pJ。}
\answer{ex:17:2}{$c=0.1$ 时为 $5.6$、$8.7$、$9.2$~bit/s，极限 $9.3$~bit/s；$c=0$、$N=1000$ 时为 $65$~bit/s。}
\answer{ex:17:3}{$B=\log_2N+P\log_2P+(1-P)\log_2\frac{1-P}{N-1}=4.87$、$4.37$、$3.29$ 比特/字符，即 $7.3$、$6.6$、$4.9$~bit/s；要达到 $10$~bit/s，每秒需要 $2.3$ 个字符。}
\answer{ex:17:4}{每个分量的误差方差为 $2b/(Nm^2T)+\sigma_\xi^2$；$N\approx90$ 时两项相等；极限为每个窗口每个分量 $\tfrac12\log_2(1+0.25/0.04)\approx1.4$ 比特。}
\answer{ex:17:5}{当 $\eta_bd^2+\eta_db^2<2$，即大约 $\eta<1$ 时这一对收敛：$0.8$ 时收敛，$1.1$ 时出现周期为 $2$ 的稳定振荡，$1.5$ 时为非周期振荡，$2$ 时发散；两个学习者的速度需要拉开。}
\answer{ex:17:6}{阈值 $\approx4.4\sigma$；$102$~kbit/s，$0.10$~mW，而原始数据流需 $205$~mW，小 $2000$ 倍；只有 $5.7\cdot10^{-5}$ 的计数会饱和（多于 $3$ 个脉冲）。}
\answer{ex:17:7}{开放问题。$\text{SNR}\approx0.9$ 时约 $46$~bit/s，独立噪声时约 $330$~bit/s。如果妨碍的是类似信号的噪声，提取的信息会随通道数增加而饱和；如果是对编码的无知，更好的解码器（例如利用脉冲时间的解码器）会使信息增加。}

\answer{ex:18:1}{(a)~$0.9949\le w\le1.0049$，即 $|1-w|\lesssim0.005$。(b)~$K\ge400$ 个突触。}
\answer{ex:18:2}{$(\omega_R\tau_w)^2=\sqrt{\alpha(\alpha+2+2\varrho)}/\varrho-1$，$\varrho=\tau_m/\tau_w$；$f_R\approx11.1$~Hz；$|Z(\omega_R)|/|Z(0)|=4.6$。}
\answer{ex:18:3}{(a)~$f=1/\bigl(\tau\ln\frac{\mu}{\mu-\theta}\bigr)$；要得到 $1$~Hz，需要 $\mu/\theta-1\approx3.7\cdot10^{-44}$。(b)~$T=\pi\tau/\sqrt\mu$，$f\approx3.2$~Hz：$f\propto\sqrt\mu$。}
\answer{ex:18:4}{积分器在 $f\lesssim17.7$~Hz时工作（低通滤波器），谐振器只在 $5.5$--$22$~Hz频带内工作（带通滤波器）。}
\answer{ex:18:5}{(a)~$T_{\min}=\frac{\tau}{w-1}\ln\frac{0.5}{0.3}\approx10.2\ms$。(b)~$B>w-1-0.2=0.3$。}
\answer{ex:18:6}{调节的时间常数为 $\tau/(\eta\langle r^2\rangle)$，$\eta=\tau\ln10/60\,\text{s}\approx3.8\cdot10^{-3}$；$I\neq0$ 时权重会偏移，需要从 $e$ 中减去指令的副本 $I/\tau$。}
\answer{ex:18:7}{开放问题。谐振器在 $11$~Hz时仅胜出 $1.35$ 倍，在 $1$~Hz时则落后 $17$ 倍；重构的主要收益在于借助阈值选择频带（习题~\ref{ex:18:4}）。}

\answer{ex:19:1}{(a)~一个突触给出 $6.7\mV$（$R=66.7$~M$\Omega$），因此要能到达阈值至少需要 $4$ 个（三个只能渐近地达到）；$10\ms$ 内需要 $8$ 个；$5\ms$ 内需要 $14$ 个。(b)~$0.1\ms\ll\tau_m$，泄漏带来的修正 $\sim0.25\%$。}
\answer{ex:19:2}{(a)~$1.75\cdot10^{10}$（四分之一的混合器对什么都响应）、$4.4\cdot10^9$ 和 $0.06$（$k=40$ 时几乎从不响应）。(b)~$10^3$ 倍：$2\cdot10^5$ 对 $200$。}
\answer{ex:19:3}{(a)~由二项式系数的对称性，$C(2n,n)=2^{2n-1}$。(b)~$0.93$ 和 $0.09$；过渡区在 $P$ 上的宽度 $\sim\sqrt{2n}$，相对宽度 $\sim1/\sqrt n$。}
\answer{ex:19:4}{只靠一个树突，无论输入多少，都有 $V_{\text{胞体}}<\kappa E_E<\theta$；$g_0=0.5g_L$ 时每个树突上需要 $n\ge3$ 个输入。}
\answer{ex:19:5}{$I\ge C\sigma_V/\sigma_t=15$~nA，即 $150$ 个同步突触——不现实；小神经元只需 $1$~nA，而用一个 $4$~nA的突触，抖动为 $5$~\textmu s。}
\answer{ex:19:6}{近端的最大值：$0.03\,\tau_m$ 和 $0.97$（$L=0.2$）；$0.32\,\tau_m$ 和 $0.66$（$L=1$）；$0.79\,\tau_m$ 和 $0.33$（$L=2$）。按总电容所作的估计~\eqref{eq:dendriteload}~只在 $L\ll1$ 时成立。}
\answer{ex:19:7}{开放问题。$m$ 固定时，响应时间随记忆量 $P$ 线性增长；比例 $m=\varphi n$ 固定时，它不再依赖于 $n$，大神经元的缓慢是对许多输入求和的结果，而不是尺寸本身造成的。}

\answer{ex:20:1}{$t_{\text{醒}}=\tau_r\ln\frac{1-L}{1-H}=16.8$~h，$t_{\text{睡}}=\tau_d\ln\frac HL=5.8$~h，周期 $22.5$~h；阈值随昼夜的摆动把振荡器拉到 $24$~h——这就是锁相~\eqref{eq:pll}。}
\answer{ex:20:2}{(a)~$L=0.111$；若 $L=0.092$，睡眠会持续 $9.6$~h。(b)~增长 $22\%$（$1/0.82=1.22$），而不是 $18\%$。}
\answer{ex:20:3}{$H=\frac{p-1}{p-1/q}=0.623$，$L=H/q=0.093$，其中 $p=e^{16/\tau_r}$，$q=e^{8/\tau_d}$。在 $4$~h 后被叫醒，相位会永久偏移（入睡提前 $7.2$~h）；阈值摆动时，相位在两三天内恢复。}
\answer{ex:20:4}{(a)~$r_\pm=\frac12\bigl(1\pm\sqrt{1-4k/w}\bigr)=0.947,\ 0.053$；$a_{\text{上}}=3.51$，$a_{\text{下}}=0.49$。(b)~工作 $\approx0.33$~s，静默 $\approx0.59$~s，周期 $\approx0.93$~s，工作比例 $0.36$。}
\answer{ex:20:5}{$a_{\text{上}}/r_+<b<a_{\text{下}}/r_-$：$3.71<b<9.20$。\nt{ACET} 使 $b$ 降到 $3.71$ 以下：交点移到上分支，工作状态（$r\approx1$）变得稳定。}
\answer{ex:20:6}{$D=24$、$32$、$40$、$48$、$64$ 时分别为 $4.9$、$6.8$、$8.8$、$5.5$、$8.9$~h：时长由入睡时的昼夜相位决定，而不是由欠债决定。}
\answer{ex:20:7}{学完 B 后，A 上的误差平均为 $0.51$（$20$ 组数据中从 $0.19$ 到 $1.08$；零输出的误差为 $1$）。交错学习时两个误差都小于 $10^{-4}$。}
\answer{ex:20:8}{开放问题。均匀缩放保留权重的比值，但只有当阈值按同一比例缩放时，才能保留阈值神经元的回答；交错重放则改变比值。大接触保持不变，这与纯粹的均匀缩放相矛盾。}

\answer{ex:21:1}{(a)~$T=\tau_{\text{rec}}\ln\frac{1}{1-x^*}$：$1.84$~s和 $3.68$~s。(b)~从 $3.36$ 到 $4.24$~s：灵敏度 $\dd T/\dd x^*=\tau_{\text{rec}}/(1-x^*)=80$~s，$x^*\to1$ 时群放电不规则。}
\answer{ex:21:2}{(a)~$8.6$~W，与~\eqref{eq:energy} 一致。(b)~$205$~Mbit/s，$0.2$~W，是培养物（$10^{-4}$~W）的 $2\cdot10^3$ 倍。}
\answer{ex:21:3}{$x_{\text{ss}}=1/(1+Ur_b\tau_{\text{rec}})<x^*$ 的条件是 $r_b>(1/x^*-1)/(U\tau_{\text{rec}})=0.28$~Hz（$x^*=0.9$），$x^*=0.99$ 时为 $0.025$~Hz；突触强度降到 $x_{\text{ss}}$，即减弱 $10\%$ 和 $1\%$。}
\answer{ex:21:4}{$u(t-k)$ 彼此正交归一，$m_k=\|P_Vu(t-k)\|^2$，其中 $P_V$ 是到输出张成的 $N$ 维子空间上的投影；由贝塞尔不等式，$\sum_k m_k\le N$。}
\answer{ex:21:5}{有 $\tanh$ 时 $\mathrm{MC}\approx9$--$11$，线性储备池为 $15$--$18$（三个随机种子），都 $\le30$：非线性以记忆为代价。}
\answer{ex:21:6}{(a)~$n\le102$。(b)~高出 $5.6$ 个标准差。}
\answer{ex:21:7}{开放问题。关键对照：冻结读出，检查在无刺激的停顿之后改进是否依然存在；若不存在，则改变的是状态，而不是权重。}

\answer{ex:22:1}{(a)~$32\%$ 和 $100\%$。(b)~$R_\uparrow-R_\downarrow\ge2\sqrt{2000\,\text{Hz}/1\,\text{s}}\approx89$~Hz，仅为总频率的 $4.5\%$。}
\answer{ex:22:2}{(a)~$n\ge4(p_1q_1+p_2q_2)/\Delta^2\approx560$。(b)~$\approx56$。}
\answer{ex:22:3}{(a)~细致平衡：$\rho PK$ 是对称的。(b)~未接住比例等于调和平均 $1/\langle1/P\rangle=0.18$，而无反馈时为 $0.5$。}
\answer{ex:22:4}{(a)~面积为 $4h^2=0.04$ 的平行四边形（计入对 $p$ 的限制后，数值为 $0.045$）。(b)~$\approx0.18$。}
\answer{ex:22:5}{$0.49$：一部分培养物走进几乎总是未接住的区域，并在那里游荡。加上限制后为 $1.00$：在有界集合上，$\propto1/P$ 的密度集中到吸收区。}
\answer{ex:22:6}{(a)~需要 $\ge5$ 个等级；$8$ 个电极足够。(b)~不能：需要 $r_{\max}\ge25/T=100$~Hz。}
\answer{ex:22:7}{开放问题。随机搜索的时间随 $d$ 大致按体积比 $(L/h)^d$ 增长，带误差符号的搜索则随 $d$ 线性增长。能区分两个假说的是对调实验：未接住后给可预测但强的刺激，接住后给不可预测但弱的刺激；每组需要几十个培养物。}


\chapter{实验补充}
\label{sec:supplement}

这里为每一章汇总了其主要论断及其依据：在什么实验中测得、具体测得了什么、发表在哪里。右栏的状态表示论断有多可靠：\emph{已测量}——有直接实验；\emph{理论}——在正文章节或经典文献中推导出的数学结论；\emph{模型}——用本书模型计算的结果（脚本在 \texttt{models/} 文件夹中）；\emph{估计}——没有直接测量的数量级数值；\emph{假说}——尚未被实验证实的合理解释。没有实验的地方，就直说没有。

\section{第\ref{sec:neurontypes}章　神经元类型}

本章的数字在有直接计数的地方依据人脑细胞的直接计数；“一个神经元，一种递质”的规则依据细胞图谱以及发现了例外的实验；\nt{DOPA}的作用依据脉冲记录；其余广播通道的作用则依据假说。

{\footnotesize
\setlength{\tabcolsep}{3pt}\renewcommand{\arraystretch}{1.15}
\begin{longtable}{>{\raggedright}p{0.5cm}>{\raggedright}p{4.0cm}>{\raggedright}p{5.6cm}>{\raggedright}p{2.3cm}>{\raggedright\arraybackslash}p{1.7cm}}
\toprule
编号 & 命题 & 实验与测量内容 & 来源 & 状态 \\
\midrule\endfirsthead
\toprule
编号 & 命题 & 实验与测量内容 & 来源 & 状态 \\
\midrule\endhead
1 & 人脑约有 $8.6\cdot10^{10}$ 个神经元，\nt{GLUT}神经元几乎全是小脑的小神经元（表~\ref{tab:types}） & 各向同性分离法：把组织溶解成均匀的细胞核悬液，计数带神经元标记物NeuN的细胞核。成年男性有 $(86.1\pm8.1)\times10^9$ 个神经元；其中只有 $19\,\%$ 位于大脑皮层，大部分在小脑 & \cite{Azevedo2009} & 已测量 \\
2 & 几乎每个神经元只有一种主要递质（命题~\ref{prop:dale}），但有例外 & 小鼠脑转录组图谱：约 $4\cdot10^6$ 个细胞，$5322$ 个簇；按合成与包装递质的基因为各簇指定递质。例外有直接测量：脑片中的\nt{DOPA}神经元经同一种囊泡转运体同时释放GABA，抑制纹状体神经元；部分\nt{DOPA}轴突的谷氨酸和多巴胺从同一根导线的不同末梢释放（电子显微镜和光遗传学） & \cite{Strata1999,Yao2023,Tritsch2012,Zhang2015} & 已测量 \\
3 & 整个权重矩阵列同号~\eqref{eq:dalesign} & 本章由命题~\ref{prop:dale} 以及“谷氨酸受体几乎都是兴奋性的、GABA受体几乎都是抑制性的”推出。“同一神经元的所有接收者得到同号权重”作为一般规则没有直接检验；反例是\nt{DOPA}神经元共同释放GABA（第~2行）和带有不同符号受体的接收者（第~11行） & 本章推导 & 理论 \\
4 & 皮层神经元中四分之三到五分之四是\nt{GLUT}，五分之一到四分之一是\nt{GABA} & 在猴皮层 $10$ 个区中，对贯穿皮层全层、宽 $50\um$ 的柱条做GABA免疫染色：$8$ 个区中GABA神经元约占全部神经元的 $25\,\%$，初级视皮层中不到 $20\,\%$ & \cite{Hendry1987} & 已测量 \\
5 & 一个\nt{GLUT}神经元约有 $10^4$ 个输出 & 尸检体视学：年轻男性新皮层中有 $1.64\cdot10^{14}$ 个突触（电子显微镜，体视框法）和约 $2.3\cdot10^{10}$ 个神经元。两者之比为每个神经元 $\sim7\cdot10^3$ 个突触；平均而言输入数等于输出数 & \cite{Tang2001synapses,Pakkenberg1997} & 已测量 \\
6 & \nt{DOPA}神经元的输出分支成 $10^5$--$10^6$ 个末梢；这是根据靶区覆盖范围的估计，而非计数 & 用病毒标记大鼠单个\nt{DOPA}神经元的膜，完整重建轴突：一个神经元覆盖纹状体体积的 $0.45$ 到 $5.7\,\%$（平均 $2.7\,\%$）。测量的是覆盖范围，而非末梢数：末梢数由覆盖范围按量级推出，没有直接计数。\nt{SERO}和\nt{NORA}的末梢数是粗略估计 & \cite{Matsuda2009} & 已测量 \\
7 & 广播神经元很少：\nt{DOPA} $\sim5\cdot10^5$（两群中的主要一群，下限），\nt{SERO} $\sim3\cdot10^5$（两次局部计数之和），\nt{HIST} $\sim6\cdot10^4$（表~\ref{tab:types}，图~\ref{fig:typepie}） & 人体体视学计数：黑质中 $550\,000$ 个含色素神经元（不含腹侧被盖区）；中缝背核 $235\,000$ 个神经元（该核全部神经元），脑桥被盖另有 $125\,000$ 个血清素神经元；下丘脑约 $64\,000$ 个组胺神经元。\nt{NORA}、\nt{GLYC}、\nt{ACET}和\nt{ADRE}没有直接计数：表中是粗略的量级估计 & \cite{Pakkenberg1991,Baker1990,Baker1991,Airaksinen1991} & 已测量 \\
8 & 缝隙中的谷氨酸猛增到约一毫摩尔，在 $\sim1\ms$ 内下降 & 根据突触传递过程中快速解离的竞争性阻断剂从NMDA受体上被置换的情况（海马神经元培养物）：峰值 $1.1$\,mM，衰减时间常数 $1.2\ms$ & \cite{Clements1992} & 已测量 \\
9 & 工作中的皮层里兴奋性电流与抑制性电流几乎平衡，神经元停在阈值附近 & 理论：强连接网络会自行进入平衡状态。实验：雪貂前额叶皮层在体记录中，在“上”状态期间突触电流的反转电位在数百毫秒内保持在 $-37$\,mV左右，而电导变化达 $\sim21$\,nS：兴奋与抑制同升同降 & \cite{vanVreeswijk1996,Haider2006} & 已测量 \\
10 & \nt{DOPA}神经元报告奖励预测误差~\eqref{eq:rpe}（图~\ref{fig:rpe}） & 猴\nt{DOPA}神经元的脉冲：对意外的果汁发出簇放电；学习后对提示信号发出簇放电，对被预测的果汁不响应；许诺的果汁没给时放电频率下陷。在定量实验中，频率与当前奖励和过去奖励加权平均之差成正比，但只对正误差如此。方程~\eqref{eq:rpe} 本身是时间差分算法 & \cite{Schultz1997,Bayer2005,SuttonBarto2018} & 已测量 \\
11 & 符号和速度由受体决定：离子型不到一毫秒，代谢型慢得多（命题~\ref{prop:receptor}） & 向海马神经元的膜片施加短促的谷氨酸脉冲：经离子型受体的电流在 $0.2$--$0.6\ms$ 内上升（从 $20$ 到 $80\,\%$），在 $2$--$3\ms$ 内衰减。锥体神经元的慢抑制电位可被代谢型GABA受体的选择性阻断剂消除。多巴胺受体有两个作用相反的家族；在纹状体中，带D1受体的神经元需要多巴胺来增强突触，带D2受体的神经元需要多巴胺来减弱突触 & \cite{Colquhoun1992,DutarNicoll1988,Beaulieu2011,Shen2008} & 已测量 \\
12 & 广播通道不是一根导线，而是由少数线路组成的线束（\ref{sec:buses}~节） & 用病毒遗传学方法标记小鼠中缝背核的血清素神经元：投射到杏仁核和投射到额叶皮层的神经元位于核内不同位置，分支到互补的区域，对厌恶刺激的反应相反。综述：蓝斑（\nt{NORA}）的神经元亚群编码不同的过程 & \cite{Ren2018,Poe2020} & 已测量 \\
13 & \nt{NORA}设定增益 $g$ & 自适应增益假说；儿茶酚胺提高传递特性陡度的网络模型。“$g$ 随去甲肾上腺素增大”在整个网络中没有直接测量；已测得瞳孔直径跟随猴蓝斑神经元的脉冲变化 & \cite{AstonJones2005,ServanSchreiber1990,Joshi2016} & 假说 \\
14 & \nt{ACET}切换“写入/读取”并设定学习速率 & 假说。依据：在大鼠嗅皮层脑片中，乙酰胆碱激动剂（$100\,\mu$M）使内部“回忆”突触减弱超过 $60\,\%$，而使输入突触减弱不到 $15\,\%$ & \cite{Hasselmo2006,HasselmoBower1992} & 假说 \\
15 & \nt{SERO}设定折扣因子 $\gamma$；血清素神经元放电平稳，每秒几个脉冲，在睡眠的某一阶段几乎沉默 & “血清素 $=\gamma$”假说。最有力的论据：光遗传学激活小鼠中缝背核的血清素神经元，会减少在等待延迟奖励时过早放弃的次数。脉冲频率以及快速眼动睡眠中的沉默是测量过的（记录综述） & \cite{Doya2002,Miyazaki2014,JacobsAzmitia1992} & 假说 \\
\bottomrule
\end{longtable}}

\section{第\ref{sec:glutcomputer}章　\nt{GLUT}神经元计算什么}

本章的逻辑命题是关于非负权重电路的定理，在本章中推导得出；实验支持的是神经元模型，以及按这些定理搭建的电路（符合检测器、延迟线、自我维持的神经元群、链）确实存在于大脑之中。

{\footnotesize
\setlength{\tabcolsep}{3pt}\renewcommand{\arraystretch}{1.15}
\begin{longtable}{>{\raggedright}p{0.5cm}>{\raggedright}p{4.0cm}>{\raggedright}p{5.6cm}>{\raggedright}p{2.3cm}>{\raggedright\arraybackslash}p{1.7cm}}
\toprule
编号 & 命题 & 实验与测量内容 & 来源 & 状态 \\
\midrule\endfirsthead
\toprule
编号 & 命题 & 实验与测量内容 & 来源 & 状态 \\
\midrule\endhead
1 & 神经元是带阈值和复位的RC电路~\eqref{eq:glutneuron} & 向大鼠皮层锥体神经元胞体（脑片）注入类似活体突触流的噪声电流。平均频率随电流均值和散布变化的关系，可以用泄漏积分发放神经元加上慢的频率适应来描述。本章给出的 $\tau$ 和延迟范围没有引用实验 & \cite{Rauch2003} & 已测量 \\
2 & 模型~\eqref{eq:glutneuron} 是一种简化：输入分支本身会产生局部脉冲 & 在体直接记录小鼠视皮层锥体神经元的树突：视觉刺激引发依赖NMDA受体的局部树突脉冲；抑制它们会使神经元的朝向调谐变宽 & \cite{Smith2013} & 已测量 \\
3 & 符合窗口 $\Delta_{\max}=\tau\ln\frac{w}{\theta-w}$~\eqref{eq:window}；$\theta=1.5w$ 时等于 $0.7\tau$ & 模型~\eqref{eq:glutneuron} 的推论。对有快速抑制的神经元，有窄总和窗口的直接测量（第\ref{sec:gaba}章，其表中第~3行） & 本章推导 & 理论 \\
4 & 仅由\nt{GLUT}神经元组成的网络只能计算输入的单调函数（命题~\ref{prop:monotone}） & 本章的证明：从静默出发、右端不减的迭代。这是关于模型的命题；没有在无抑制的活体网络上检验它的实验 & 本章推导 & 理论 \\
5 & 感觉器官提供一对导线：一些细胞对刺激增强起反应，另一些对刺激减弱起反应 & 视网膜：视锥信号在双极细胞处就已分为接通和断开两个通道。在猴身上药理阻断接通型双极细胞：两个通道一直分开通到皮层；阻断后，检测光变强的能力变差，而检测光变弱的能力不受影响 & \cite{Schiller1992} & 已测量 \\
6 & 采用双线编码，\nt{GLUT}神经元网络可以异步地、无需共同节拍地计算任意逻辑函数，代价是神经元数目翻倍（命题~\ref{prop:dualrail}，\eqref{eq:dualrail}，图~\ref{fig:dualxor}） & 双线编码和凭信号本身判断就绪，是延迟不敏感异步电路的标准技术。六个神经元的异或电路在本章搭建。没有实验表明大脑正是用双线编码计算逻辑函数 & \cite{Sparso2001} & 理论 \\
7 & 人耳间声音到达时间差最大 $\approx0.6\ms$，而大脑能分辨约十微秒 & 二选一实验中的听者：时间差分辨阈（正确率75\,\%）对 $150$--$1700$\,Hz频带噪声为 $9\,\mu$s，对 $1$\,kHz纯音为 $11\,\mu$s，对咔嗒声为 $28\,\mu$s。上限 $0.6\ms$ 是本章的计算，$0.22\,\text{m}/343\,\text{m/s}$ & \cite{KlumppEady1956} & 已测量 \\
8 & 鸟类通过延迟线和符合检测器把时间差变成位置~\eqref{eq:itd}（图~\ref{fig:itd}） & 谷仓猫头鹰：来自两耳的轴突从相对两侧进入符合检测器核团；脉冲沿轴突的到达时间随深度有序变化，穿过核团的延迟范围（$700\um$ 上 $160\,\mu$s）与两耳间延迟范围一致 & \cite{Carr1990} & 已测量 \\
9 & 哺乳动物中没有找到延迟阵列；同样的任务由带有精确定时的\nt{GLYC}神经元抑制的符合检测器完成 & 沙鼠内侧上橄榄核在体记录：响应不形成方向图；用微量移液管施加甘氨酸及其阻断剂表明，精确定时的\emph{甘氨酸}抑制决定了延迟的工作范围。这里的抑制来自\nt{GLYC}，而不是\nt{GABA} & \cite{Brand2002,Grothe2010} & 已测量 \\
10 & 强反馈神经元群是触发器；自我维持的活动在刺激后持续数秒（图~\ref{fig:bistable}） & 猴前额叶皮层，延迟眼动任务，眼睛需移向记住的位置（延迟 $1$--$6$\,s）：$170$ 个任务相关神经元中有 $87$ 个在延迟期改变频率，其中 $79\,\%$ 依赖于记住的方向。这种活动由有两个稳定状态的反馈维持，是理论 & \cite{Funahashi1989,Wang2001} & 已测量 \\
11 & 驱动持续超过 $\approx0.7\tau$ 才能写入比特（图~\ref{fig:recurrentgroup}b） & 用欧拉法计算方程 $\tau\,\dd r/\dd t=-r+f(wr+I)$，$w=1$，$I=0.6$；\texttt{models/} 中没有对应脚本。没有实验 & 本章计算 & 模型 \\
12 & 记忆可以存放在突触易化中，几乎不需要脉冲 & 理论：末梢中的残余钙使记录保持约一秒。依据：雪貂内侧前额叶皮层（多个神经元同时记录）中存在一个由易化突触连接、后效持久的锥体神经元子网。关于记忆保持期间“安静”间歇的观察综述。皮层何时用哪种方式，是悬而未决的问题 & \cite{Mongillo2008,WangMarkram2006,Stokes2015} & 假说 \\
13 & 记忆由疲劳擦除：递质储备耗竭 & 皮层锥体神经元成对记录：高频脉冲下突触响应下降，下降速度由释放概率决定。这正是熄灭自我维持神经元群的机制，尚无直接证明 & \cite{Tsodyks1997} & 已测量 \\
14 & 神经元群链是由事件启动的定时器；鸣禽的神经元严格依次发放 & 在睡眠和鸣唱时记录斑胸草雀高级发声中枢中已鉴定的神经元：每个投射到运动核团的神经元在歌曲的一个精确时刻发出一个短簇放电，各神经元依次发放 & \cite{Hahnloser2002} & 已测量 \\
\bottomrule
\end{longtable}}

\section{第\ref{sec:gaba}章　\nt{GLUT}与\nt{GABA}协同工作}

本章的逻辑命题和动力学命题都在本章中由线性分析和欧姆定律推出；实验表明，它们所依据的电路——带延迟的否决、紧随兴奋之后的前馈抑制、由抑制实现的稳定、\nt{GLUT}--\nt{GABA}环路的节律——确实在大脑中工作。

{\footnotesize
\setlength{\tabcolsep}{3pt}\renewcommand{\arraystretch}{1.15}
\begin{longtable}{>{\raggedright}p{0.5cm}>{\raggedright}p{4.0cm}>{\raggedright}p{5.6cm}>{\raggedright}p{2.3cm}>{\raggedright\arraybackslash}p{1.7cm}}
\toprule
编号 & 命题 & 实验与测量内容 & 来源 & 状态 \\
\midrule\endfirsthead
\toprule
编号 & 命题 & 实验与测量内容 & 来源 & 状态 \\
\midrule\endhead
1 & 皮层中五分之一到四分之一的神经元是\nt{GABA}神经元 & 猴皮层 $10$ 个区的GABA免疫染色：$8$ 个区中约占神经元的 $25\,\%$，初级视皮层中不到 $20\,\%$。皮层抑制性神经元多样性综述 & \cite{Hendry1987,Markram2004} & 已测量 \\
2 & 抑制的作用在很大程度上取决于落点：树突、胞体、脉冲产生处、别的导线的末梢 & 综述：皮层\nt{GABA}神经元的各个类别按其在接收者上的靶点区分。同时记录锥体神经元的胞体和树突：前馈抑制在胞体上比在树突上强得多。脊髓中作用于别的导线末梢的抑制减少一条输入线路的递质释放（测量综述） & \cite{Markram2004,Pouille2001,Rudomin1999} & 已测量 \\
3 & 经中介改变符号的代价是一个延迟，约一毫秒；紧随兴奋而来的抑制收窄符合窗口 & 大鼠海马脑片中的锥体神经元：经一个中介的抑制以很短的延迟紧随直接兴奋到达，两个兴奋性输入能总和到阈值的窗口小于 $2\ms$。在这种抑制较弱的树突上，窗口较宽 & \cite{Pouille2001} & 已测量 \\
4 & 否决 $a\wedge\bar b$~\eqref{eq:veto} 加上或和常数1是函数完备的（命题~\ref{prop:gabacomplete}） & 本章的证明。经典结果：带抑制性输入的阈值元件网络可以实现任意逻辑表达式。这是关于模型的命题，没有实验 & \cite{McCullochPitts1943} & 理论 \\
5 & 带延迟的否决给出运动方向检测器，最优速度 $v^*=\Delta x/d$~\eqref{eq:vstar} & 兔视网膜：方向选择性由抑制产生，抑制在“零”方向运动时抵消兴奋。分别记录方向选择细胞的兴奋性和抑制性输入，表明星爆无长突细胞（\nt{GABA}）对它们的抑制是不对称的，只来自朝向“零”方向一侧的突起。$v^*$ 的公式是本章的推导 & \cite{Barlow1965,Fried2002} & 已测量 \\
6 & 分流抑制对电压做除法而非减法~\eqref{eq:shunt}（图~\ref{fig:shunt}） & 三个电导构成的分压器的欧姆定律，氯的平衡电位在静息附近；针对不发放脉冲的神经元 & 本章推导 & 理论 \\
7 & 分流对脉冲频率的作用是减法，除法来自抑制加上平衡的背景噪声 & 模型：发放中的神经元流过分流通道的电流几乎恒定，对频率的效应是减法。实验：用动态钳向大鼠体感皮层锥体神经元（脑片）注入兴奋性和抑制性背景电导；两者同时平衡地增加，会除以响应斜率，而频率没有明显平移。综述：除以总活动的运算见于大脑许多部位 & \cite{HoltKoch1997,Chance2002,Carandini2012} & 已测量 \\
8 & $\beta>1$ 时相互抑制选出唯一赢家（命题~\ref{prop:wta}，\eqref{eq:wta}） & 特征值 $-(1\pm\beta)/\tau$ 是本章的推导。$\beta=0.5$ 时的频率 $0.73$ 和 $0.53$ 是不动点 $r_{1,2}=(I_{1,2}-\beta I_{2,1})/(1-\beta^2)$；\texttt{models/} 中没有对应脚本。本章未引用直接实验 & 本章推导 & 理论 \\
9 & 经\nt{GABA}用信号复位存储；两群相互抑制的神经元构成锁存器 & 由图~\ref{fig:bistable} 和命题~\ref{prop:wta} 推出。没有实验 & 本章推导 & 理论 \\
10 & 若抑制足够强且足够快，\nt{GLUT}--\nt{GABA}环路的平衡点稳定（命题~\ref{prop:ei}） & 双群体模型~\eqref{eq:ei}；条件 (i) 和 (ii) 是其矩阵的行列式和迹，在本章中推导 & \cite{WilsonCowan1972} & 理论 \\
11 & $w_{EE}=1.5$、$w_{EI}w_{IE}=2$ 时，快速抑制（$\tau_I=2\ms$）维持 $E^*=0.67h$，慢速抑制（$30\ms$）激起振荡（图~\ref{fig:ei}） & \eqref{eq:ei} 的数值解，脚本 \texttt{figures/make\_ei.py} & 计算 & 模型 \\
12 & 皮层工作在抑制稳定状态；其特征是推一下\nt{GABA}神经元，它们的频率反而下降~\eqref{eq:isn} & 理论预言了这种反常响应。实验：猫初级视皮层的细胞内记录——当响应被周围刺激抑制时，抑制性电导不增反降，与兴奋性电导一同下降。在小鼠视皮层、体感皮层和运动皮层中光遗传学兴奋PV神经元，会降低抑制性神经元的频率，包括无刺激时和浅麻醉下。图~\ref{fig:ei} 数值下的系数 $-1/3$ 是本章的推导 & \cite{Tsodyks1997isn,Ozeki2009,Sanzeni2020} & 已测量 \\
13 & “\nt{GLUT}兴奋\nt{GABA}，\nt{GABA}抑制\nt{GLUT}”环路产生周期为 $12$--$50\ms$ 的快速节律 & 在小鼠体感皮层以 $8$--$200$\,Hz光遗传学刺激快速\nt{GABA}神经元，选择性地增强 $\gamma$ 节律（$20$--$80$\,Hz，周期 $12$--$50\ms$）；刺激\nt{GLUT}神经元只增强较慢的节律 & \cite{Cardin2009,Buzsaki2012} & 已测量 \\
\bottomrule
\end{longtable}}

\section{第\ref{sec:code}章　摩尔斯电码}

本章的极限估计来自信息论和计算；测量过的是信道中噪声产生在哪里、大脑工作得多快、一个脉冲要花多少代价，而皮层实际使用哪种编码，仍是悬而未决的问题。

{\footnotesize
\setlength{\tabcolsep}{3pt}\renewcommand{\arraystretch}{1.15}
\begin{longtable}{>{\raggedright}p{0.5cm}>{\raggedright}p{4.0cm}>{\raggedright}p{5.6cm}>{\raggedright}p{2.3cm}>{\raggedright\arraybackslash}p{1.7cm}}
\toprule
编号 & 命题 & 实验与测量内容 & 来源 & 状态 \\
\midrule\endfirsthead
\toprule
编号 & 命题 & 实验与测量内容 & 来源 & 状态 \\
\midrule\endhead
1 & 皮层\nt{GLUT}神经元的频率从每秒零点几个到几十个脉冲，神经元大部分时间工作在下限附近 & 在清醒大鼠听皮层中用贴附细胞的移液管记录（神经元的选择与其活动无关）：任一时刻频率高于每秒 $20$ 个脉冲的神经元不到 $5\,\%$；部分神经元的背景频率低于每秒 $0.01$；频率分布为对数正态 & \cite{Hromadka2008} & 已测量 \\
2 & 脉冲信道容量 $R_{\max}\approx r\log_2\frac{e}{r\Delta t}$~\eqref{eq:spikeentropy}，几乎全部在于脉冲的时间（命题~\ref{prop:capacity}） & 由长度为 $\Delta t$ 的格子组成的二进制串的熵。本章的数字：$r=10$ 每秒、$\Delta t=1\ms$ 时每个脉冲 $8$ 比特，$\Delta t=10\ms$ 时约 $5$ 比特，脉冲计数器每秒不到 $2$ 比特 & \cite{MacKay1952} & 理论 \\
3 & 感觉神经元每个脉冲传送一到几比特，最好的情况下可达上限的一半 & 对已知刺激的感觉神经元，根据其脉冲重建刺激；测量汇总见该书 & \cite{Rieke1997} & 已测量 \\
4 & 脉冲发生器是精确的；精确时间由输入的快速跳变决定 & 大鼠皮层脑片中的神经元：对重复的、快速波动的电流，脉冲以优于 $1\ms$ 的精度重复；对恒定电流，脉冲时刻漂移开来 & \cite{Mainen1995} & 已测量 \\
5 & 突触不可靠：一半以上的脉冲因释放的随机性而到不了接收者 & 对海马锥体神经元的输入做最小刺激（脑片）：一半以上的脉冲不引起响应。排除了刺激阈值波动、轴突分支处的传导故障和实验温度低等因素；剩下的是概率性释放 & \cite{AllenStevens1994} & 已测量 \\
6 & 活体皮层中的脉冲流看起来是随机的：间隔散布如泊松流，脉冲数方差接近均值；部分“噪声”是关于动物运动的消息 & 清醒猴视区V1和MT的记录：间隔的变异系数约为 $1$，脉冲数方差与均值之比如同随机过程；把独立小输入相加的模型给出的散布要小得多。在小鼠视皮层中，相当大一部分散布可以由动物自身运动的录像预测 & \cite{Softky1993,Stringer2019} & 已测量 \\
7 & 频率编码很慢：误差 $1/\sqrt{rT}$~\eqref{eq:rateerror}，而大脑在 $\sim150\ms$ 内认出图片 & 公式是泊松统计，本章推导。实验：受试者判断呈现 $20\ms$ 的陌生照片中是否有动物；“是”和“否”的脑电位约在 $150\ms$ 后分开。把这段时间分成十来级、每级 $10$--$20\ms$，是本章的估计，不是测量 & \cite{Thorpe1996} & 已测量 \\
8 & 对神经元群求平均受相关性限制：误差最多减小到 $1/\sqrt c$ 分之一~\eqref{eq:correlated}（命题~\ref{prop:pooling}） & 同时记录的猴MT区神经元对：脉冲数相关系数平均为 $0.12$，理论分析表明即便这样的相关也限制了神经元群的收益。理论：只有与信号相似的那部分相关设定极限。对立观点的模型：$50$--$100$ 个神经元的群体频率在一个脉冲间隔（$10$--$50\ms$）内就能读出，皮层不利用单个脉冲的时间 & \cite{Zohary1994,MorenoBote2014,ShadlenNewsome1998} & 已测量 \\
9 & 数可以写在首个脉冲的时间~\eqref{eq:latency}、神经元群的脉冲顺序或局部节律的相位中 & 蝾螈视网膜：短暂呈现图像的空间结构记录在输出神经元首个脉冲的相对延迟中，几乎与对比度无关。大鼠海马位置细胞：穿过“自己的”位置时，脉冲在 $\theta$ 节律（$7$--$12$\,Hz）周期中越来越早，移动 $100^\circ$ 到 $355^\circ$。皮层在多大程度上利用脉冲时间，是悬而未决的问题（第~8行） & \cite{Gollisch2008,OKeefe1993} & 已测量 \\
10 & 读出哪种编码由接收者的时间常数决定~\eqref{eq:reader}（命题~\ref{prop:reader}）；活跃皮层中的神经元更接近符合检测器 & 公式~\eqref{eq:reader} 是本章推导。在体记录综述：活跃皮层中膜电导比静息时高数倍，时间常数相应缩短。皮层正是以这种方式切换编码，尚无直接证明 & \cite{Destexhe2003} & 假说 \\
11 & 压抑型突触传送频率的变化，本质上是 $\Delta r/r$~\eqref{eq:depression}（图~\ref{fig:depression}） & 皮层锥体神经元成对记录：压抑速度由释放概率决定；压抑慢时响应反映频率，压抑快时反映符合。基于大鼠皮层实测压抑的模型：快输入和慢输入上相同的频率相对变化给出相同的响应，即突触传送 $\Delta r/r$。数字 $1.7\to2.2$、$6.7$ 和 $90\ms$ 是本章的计算 & \cite{Tsodyks1997,Abbott1997depression} & 已测量 \\
12 & 簇放电是“划”：它丢失的概率 $(1-p)^k$~\eqref{eq:burst} 很小，易化使之更小 & 综述：在不可靠的突触上，单个脉冲常常丢失，而簇放电借助易化得以可靠传送；簇放电引起突触的长时程变化。大脑把簇放电当作独立符号来读，是假说 & \cite{Lisman1997,AllenStevens1994} & 假说 \\
13 & 快速\nt{GABA}神经元设定节律和“标点”；另一类抑制抑制性神经元并打开闸门（命题~\ref{prop:punctuation}） & 皮层\nt{GABA}神经元各类别特性的综述。光遗传学：节律性兴奋快速\nt{GABA}神经元会诱发 $\gamma$ 节律，对刺激的响应取决于周期的相位。在小鼠视皮层中，PV神经元相互抑制，SST神经元抑制其他所有类别，VIP神经元主要抑制SST。在清醒小鼠中兴奋VIP神经元会解除对\nt{GLUT}神经元的抑制；开启VIP神经元的是强化信号。作为一般规则的角色分工是假说 & \cite{Tremblay2016,Cardin2009,Pfeffer2013,Pi2013} & 假说 \\
14 & 能量把平均频率限制在每秒约一个脉冲的量级~\eqref{eq:energy} & 根据解剖和生理数据得出的啮齿类灰质能量预算：脉冲和突触电流分别占信号传递能量的 $47\,\%$ 和 $34\,\%$；同时活跃的神经元不超过 $\sim15\,\%$。人类皮层：同时强烈活跃的神经元不超过 $\sim1\,\%$。每个脉冲 $\sim10^9$ 个ATP、信号传递功率 $\sim10$\,W，是本章基于这些计算的估计 & \cite{Attwell2001,Lennie2003} & 理论 \\
15 & 编码是稀疏的，并被调整到每焦耳比特数最多；皮层以精度换能量（命题~\ref{prop:economy}） & 节约编码假说。依据：限制食物的小鼠视皮层中，AMPA受体电导和突触ATP消耗下降（$-29\,\%$），脉冲频率保持不变，而朝向调谐变宽 $32\,\%$ & \cite{Barlow1961,Padamsey2022} & 假说 \\
\bottomrule
\end{longtable}}

\section{第~\ref{sec:serocomputer}~章：\nt{SERO}神经元}

\nt{SERO}神经元本身的性质（节律、自抑制、由受体决定的符号、子系统）是直接测量的；本章的计算——调节器、滤波器、逻辑——由本章的方程推出；$\tau_c$、血清素的扩散系数和释放位点数是估计；回路在脑中充当电平调节器是假说（在中缝核内自抑制是点对点的，只造成短暂停顿），而\nt{SERO}充当视野则是有行为实验支持的假说。

{\footnotesize
\setlength{\tabcolsep}{3pt}\renewcommand{\arraystretch}{1.15}
\begin{longtable}{>{\raggedright}p{0.5cm}>{\raggedright}p{4.0cm}>{\raggedright}p{5.6cm}>{\raggedright}p{2.3cm}>{\raggedright\arraybackslash}p{1.7cm}}
\toprule
序 & 命题 & 实验与测量内容 & 来源 & 状态 \\
\midrule\endfirsthead
\toprule
序 & 命题 & 实验与测量内容 & 来源 & 状态 \\
\midrule\endhead
1 & 血清素作用的符号由接收者的受体选择（命题~\ref{prop:receptor}） & 血清素受体按药理和机制分为若干家族：5-HT$_{1A}$经 G 蛋白打开钾通道并抑制细胞，5-HT$_{2A}$和离子型的 5-HT$_3$则起兴奋作用。同一种递质在不同细胞中产生相反的响应。 & \cite{BarnesSharp1999} & 已测量 \\
2 & 清醒时\nt{SERO}神经元平稳地每秒发放几个脉冲 & 记录自由活动猫的中缝背核神经元：安静清醒时为缓慢而有节律的放电，$2.8\pm0.2$次/秒；慢波睡眠中频率下降$34$--$68\,\%$，快速眼动睡眠中下降$98\,\%$。麻醉大鼠中为$1.7\pm0.5$次/秒，非常规则。 & \cite{TrulsonJacobs1979, GagnonParent2014, JacobsAzmitia1992} & 已测量 \\
3 & \nt{SERO}神经元是起搏器：不需要每个脉冲都有一次推动，但需要平稳的背景输入（脑片中是经$\alpha_1$受体的去甲肾上腺素，体内来自\nt{NORA}） & 在脑片中细胞缓慢而平稳地放电，每个脉冲之后有持续$200$--$800\ms$的后超极化。脑片中自发工作的细胞少于体内：平稳节律需要经$\alpha_1$受体的去甲肾上腺素紧张性输入，阻断这些受体则节律熄灭。 & \cite{VandermaelenAghajanian1983} & 已测量 \\
4 & 受体几乎都是代谢型的；响应缓慢：在大约一秒内升起又衰减 & 除 5-HT$_3$外，所有受体家族都与 G 蛋白偶联。在脑片中，诱发的血清素释放经 5-HT$_{1A}$产生抑制性电流，在一秒内升起又衰减。 & \cite{BarnesSharp1999, CourtneyFord2016} & 已测量 \\
5 & 在目标区域血清素释放到体积中，而不只是释放到间隙中；在中缝核内部，对邻居的抑制是点对点的 & 脑片中的快速循环伏安法：每个脉冲的释放量在有突触的区域和几乎没有突触的中缝核中相同；不受摄取阻断和受体阻断的影响；每个末梢的分子数远少于转运体和受体的数目，而细胞外浓度与主要受体的亲和力相符。在中缝核内，经 5-HT$_{1A}$的抑制是点对点的：转运体不让血清素在间隙外积累。 & \cite{Bunin1998, CourtneyFord2016} & 已测量 \\
6 & 按~\eqref{eq:seroneuron}，浓度因泵回而下降；$\tau_c$量级为一秒是估计 & 活体脑中的伏安法：细胞外血清素主要受再摄取和降解限制，而不是受合成限制。所引文献中没有$\tau_c$的直接测量；其量级是本章的估计。 & \cite{Hashemi2012serotonin} & 已测量；$\tau_c$为估计 \\
7 & 一个起搏器实现“非”和“或非”；\nt{SERO}逻辑的完备性（命题~\ref{prop:serocomplete}，图~\ref{fig:serogates}） & 由~\eqref{eq:seroneuron}推出：受抑制的起搏器是“或非”，而“或非”在功能上是完备的。没有用\nt{SERO}神经元搭建逻辑的实验；在脑中，体积彼此分离的条件并不成立。 & 本章推导 & 理论 \\
8 & 血清素经\nt{SERO}神经元自身的受体抑制它 & 在麻醉大鼠中，静脉注射和离子电泳给予的 5-HT$_{1A}$激动剂抑制中缝背核神经元放电，其剂量—反应关系与血清素本身相同；5-HT$_{1B}$激动剂几乎无作用。在脑片中，邻近细胞的内源性血清素经 5-HT$_{1A}$产生电流和放电中的短暂停顿，不改变平均频率。 & \cite{Sprouse1987, CourtneyFord2016} & 已测量 \\
9 & 在模型中，经公共体积的回路是电平调节器：离散和时间常数缩小为$1/(1+G)$~\eqref{eq:feedback}；回路在脑中这样工作是假说 & 负反馈放大器的经典公式；本章重新推导。\nt{SERO}系统的环路增益$G$没有测量过。已测量的是：在脑片中自抑制是点对点的，造成短暂停顿而不改变平均频率。 & \cite{Black1934feedback, CourtneyFord2016} & 理论；在脑中为假说 \\
10 & 浓度是频率的滑动平均，即低通滤波器~\eqref{eq:lowpass}，图~\ref{fig:lowpass} & 线性方程~\eqref{eq:seroneuron}的解；图是按该公式在$\tau_c=1\,\text{s}$时的计算。\texttt{models/}中没有单独的模型。 & 本章推导 & 理论 \\
11 & 间隙的曲折使小分子的扩散减慢约$2.6$倍 & 点源法：离子电泳注入四甲基铵，在脑片和活体脑中用离子选择电极记录其浓度。曲折度$\lambda\approx1.6$，即有效$D$比自由扩散小$\lambda^2\approx2.6$倍；细胞间隙体积分数约为$0.2$。这些工作没有测量血清素本身在组织中的扩散系数；本章中它是估计。 & \cite{Sykova2008} & 已测量 \\
12 & 独立“导线”的长度$\ell\sim\sqrt{D\tau}\approx20\um$~\eqref{eq:diffusionlength}；导线数受这类区域的数目限制（命题~\ref{prop:volumewires}） & 扩散定律并代入$D$和$\tau$的估计（第~6、11~行）；命题~\ref{prop:volumewires}是推论。没有直接计数容积传递独立通道的实验。 & 本章推导 & 理论 \\
13 & 一个\nt{SERO}神经元的输出散布在脑的巨大区域上；释放位点估计约$\sim10^5$个 & 大鼠中缝背核单个神经元的完整重建：所有上行轴突都高度分支，一个细胞的轴突总长可达$18.7$~cm，一个细胞的分支可到达纹状体、前额叶皮层和杏仁核。每个细胞的释放位点数没有直接计数。 & \cite{GagnonParent2014} & 已测量；$10^5$为估计 \\
14 & \nt{SERO}神经元分为几个子系统，输出和响应各不相同 & 小鼠的病毒—遗传标记：投射到杏仁核和投射到额叶皮层的神经元在分支上几乎不重叠，接收不同的输入，对厌恶刺激的响应相反；激活或抑制每一组对行为的改变各不相同。 & \cite{Ren2018} & 已测量 \\
15 & 指数折扣下的耐心条件$D<\tau_\gamma\ln(R/r_0)$~\eqref{eq:patience}；在实测的双曲折扣下为$D<\tau_\gamma(R/r_0-1)$ & 由折扣$e^{-D/\tau_\gamma}$或$1/(1+D/\tau_\gamma)$推出；本章推导。猴子在秒级延迟上的选择：折扣更接近双曲线而非指数；\nt{DOPA}神经元对延迟奖励预示信号的响应随延迟下降，同样符合双曲线。 & 本章推导；\cite{KobayashiSchultz2008} & 理论 \\
16 & \nt{SERO}设定视野$\tau_\gamma$（第~\ref{sec:horizon}~节） & 在小鼠中用光遗传学激活中缝背核\nt{SERO}神经元，会延长对延迟奖励的等待，并提高在奖励变稀少时寻找奖励的坚持程度。在人类中，降低血清素（无色氨酸饮食）使延迟奖励的折扣更陡。浓度如何变成$\tau_\gamma$尚不清楚。 & \cite{Doya2002, Miyazaki2014, Lottem2018, Schweighofer2008} & 假说 \\
\bottomrule
\end{longtable}}

\section{第~\ref{sec:dofa}~章：借助\nt{DOPA}学习}

突触的机制（量子、符合检测器、钙、可塑性窗口）以及多巴胺作为预测误差的行为是直接测量的；规则导向梯度以及广播的代价是定理；选择学习的结果是按本书模型的计算；计算误差的电路是假说。

{\footnotesize
\setlength{\tabcolsep}{3pt}\renewcommand{\arraystretch}{1.15}
\begin{longtable}{>{\raggedright}p{0.5cm}>{\raggedright}p{4.0cm}>{\raggedright}p{5.6cm}>{\raggedright}p{2.3cm}>{\raggedright\arraybackslash}p{1.7cm}}
\toprule
序 & 命题 & 实验与测量内容 & 来源 & 状态 \\
\midrule\endfirsthead
\toprule
序 & 命题 & 实验与测量内容 & 来源 & 状态 \\
\midrule\endhead
1 & 在脑中没有发现人工网络那种形式的误差反向传播 & 没有直接实验：这是一个关于“不存在”的命题。综述分析了该算法需要什么（与前向连接对应的反向连接、单独的阶段），以及已提出的哪些生物学近似。 & \cite{Lillicrap2020, WhittingtonBogacz2019} & 假说 \\
2 & 权重分解为释放位点数、释放概率和对一份递质的响应：$w=N_s\,p\,A_1$~\eqref{eq:quantal} & 释放减少时的蛙神经肌肉突触：接收者的响应呈台阶式波动，每级是自发的单份微小响应的整数倍；每个脉冲释放的份数服从泊松分布。 & \cite{delCastillo1954} & 已测量 \\
3 & 接收者的受体是符合检测器；钙启动权重变化 & 小鼠神经元的膜片钳：在静息电位下，Mg$^{2+}$离子堵塞由谷氨酸打开的通道，去极化时离开。海马脑片：接收者中的钙螯合剂阻断长时程增强，而用光在细胞内释放钙本身就能增强传递。 & \cite{Nowak1984, Malenka1988calcium} & 已测量 \\
4 & 决定可由任何一侧执行：接收者一侧$A_1$增大，发送者一侧$p$改变 & 在单个树突棘上方用光释放谷氨酸，引起树突棘及其受体电流迅速而持久地增长；增强伴随着 AMPA 受体的嵌入。接收者去极化会暂时抑制发送者的释放；这一效应由逆向穿过间隙的内源性大麻素传递（在\nt{GABA}输入上证明）。短时耗竭取决于发送者的释放概率。 & \cite{Matsuzaki2004, Malinow2002, WilsonNicoll2001, Tsodyks1997} & 已测量 \\
5 & 发送者和接收者同时活跃时突触增强~\eqref{eq:hebb} & 麻醉兔：沿海马输入通路施加短串脉冲后，响应增强可持续$30$~min至$10$~h。在海马脑片中，接收者的脉冲只增强同一时刻活跃的那些突触。 & \cite{Bliss1973LTP, Kelso1986hebb} & 已测量 \\
6 & 规则~\eqref{eq:oja}找到输入相关的主特征向量（命题~\ref{prop:oja}） & 定理；证明见本章。没有神经元学会其输入主成分的实验；突触中限制权重增长的机制尚未确定。 & \cite{Oja1982} & 理论 \\
7 & 按时间的赫布规则：窗口约$\pm20\ms$（图~\ref{fig:stdp}）；重复的序列中会长出链 & 培养的海马神经元对：接收者的脉冲在发送者脉冲之后$20\ms$内产生持久增强，之前$20\ms$内产生减弱；二者都依赖 NMDA 受体。图中$A_-=0.6A_+$的比例只是示意。网络中如此长出链，没有直接测量。 & \cite{BiPoo1998} & 已测量 \\
8 & 资格迹~\eqref{eq:threefactor}：共同活动后$1$--$2\,\text{s}$到达的多巴胺增强连接，更晚则不然（命题~\ref{prop:threefactor}） & 纹状体脑片，分别光遗传刺激谷氨酸和多巴胺输入：只有当多巴胺在谷氨酸输入后$0.3$--$2\,\text{s}$到达时，树突棘才长大；窗口由 cAMP 的快速上升和分解决定。在皮层中，配对刺激留下分别用于增强和减弱的迹，单胺受体在秒级窗口内把它们转化为长期变化。 & \cite{Yagishita2014, He2015eligibility} & 已测量 \\
9 & 没有多巴胺也有长达数秒的窗口：第三个信号是接收者的强烈兴奋 & 活体小鼠，CA1 区：接收者中的一次平台事件增强在其前后数秒内到达的输入，一次经过就形成位置野。 & \cite{Bittner2017} & 已测量 \\
10 & 三因子规则的平均步长沿奖励梯度~\eqref{eq:perturbation}（命题~\ref{prop:gradient}） & 基于随机偏离的学习的梯度定理；本章在小噪声下推导。 & \cite{Williams1992} & 理论 \\
11 & 噪声即搜索：网络从自身的随机偏离中学习 & 幼年斑胸草雀：关闭基底神经节回路的输出核，会急剧且可逆地消除鸣唱的变异性。成年斑胸草雀：在音高偏向平均值一侧的音节演唱上施加干扰，鸟会迅速把音高移向另一侧——从自身的离散中学习。 & \cite{Olveczky2005, TumerBrainard2007} & 已测量 \\
12 & 当$\tau_e=1\,\text{s}$、$\delta=R-\bar R$时能学会二选一；$\tau_e=50\ms$时学不会；不减去期望时权重无限增长，学习速率大时网络可能固定在更差的选择上 & \texttt{models/dofa\_learning.py}，$\eta=0.05$，$500$次尝试，$6$个随机种子。$\tau_e=1\,\text{s}$：选择$A$的比例$0.65$--$0.79\to0.96$--$1.00$，权重$0.85$--$1.33$。$\tau_e=50\ms$：$0.38$--$0.55$，权重不变。$\delta=R$：获胜者权重$860$--$1190$。$\delta=R$，$\eta=0.5$，$3$个种子：一个固定在$B$上（$0.04\to0$）。脚本输出全部四种情况；重算结果与本章一致。 & \texttt{models/\allowbreak dofa\_\allowbreak learning.py} & 模型 \\
13 & 依靠一个公共信号学习，所需时间与噪声神经元数目成正比地增长（命题~\ref{prop:broadcastcost}） & 线性网络的学习曲线：扰动神经元比精确梯度慢的倍数等于输出神经元数，扰动权重还要再慢输入数那么多倍。 & \cite{Werfel2005} & 理论 \\
14 & \nt{DOPA}的输出主要通向动作选择区域 & 大鼠单个黑质纹状体\nt{DOPA}神经元轴突的完整重建：一个细胞的分支覆盖纹状体体积的$0.45$--$5.7\,\%$（平均$2.7\,\%$），纹状体外几乎没有分支。 & \cite{Matsuda2009} & 已测量 \\
15 & 皮层学习预测自己的输入，误差在原地计算 & 综述：感觉皮层中存在对预期输入与实际输入失配的响应。这是否为皮层的普遍学习规则，尚未证明。 & \cite{Keller2018} & 假说 \\
16 & 多巴胺的行为像误差~\eqref{eq:ctd}：尖峰、向预示信号的转移、低谷（图~\ref{fig:ctdcases}） & 猴子的\nt{DOPA}神经元：对意外的果汁出现尖峰；学习后对预示信号出现尖峰，对承诺的果汁无响应；果汁未到时出现低谷。连续形式~\eqref{eq:ctd}是理论。 & \cite{Schultz1997, Doya2000} & 已测量 \\
17 & 计算$\dd V/\dd t$并减去期望的电路 & 小鼠腹侧被盖区的记录和光遗传学：\nt{DOPA}神经元从奖励中减去期望；邻近的\nt{GABA}神经元在奖励可以预期时抑制它们，兴奋或抑制这些神经元会改变误差。导数如何计算，尚未证明。 & \cite{Eshel2015arithmetic} & 假说 \\
18 & \nt{DOPA}神经元是起搏器；低谷比尖峰浅 & 在脑片中\nt{DOPA}神经元在缓慢的起搏电位上自发地、非常规则地放电。在猴子中，频率定量地跟随正误差，而负误差只得到微弱的传递。 & \cite{GraceOnn1989, Bayer2005} & 已测量 \\
19 & 行动者与评论器（图~\ref{fig:actorcritic}） & 算法来自强化学习理论。在人类（fMRI）中，腹侧纹状体无论是否需要选择都表现得像评论器，背侧纹状体只在需要选择动作时如此。 & \cite{SuttonBarto2018, ODoherty2004actor} & 假说 \\
20 & 在带第二族受体的神经元中，规则里$\delta$的符号是反的 & 纹状体脑片：在带 D1 受体的神经元中，配对刺激产生需要 D1 的增强；在带 D2 的神经元中产生需要 D2 的减弱。 & \cite{Shen2008} & 已测量 \\
\bottomrule
\end{longtable}}

\section{第~\ref{sec:dofaalt}~章：\nt{DOPA}的其他理论}

每一种扩展图景都有自己对多巴胺的直接测量作为依据（翻转点的离散、对厌恶和新异事件的响应、缓慢上升、对味道替换的响应），但解释它们的公式是理论，其中两种理论之间的实验结果目前相互矛盾。

{\footnotesize
\setlength{\tabcolsep}{3pt}\renewcommand{\arraystretch}{1.15}
\begin{longtable}{>{\raggedright}p{0.5cm}>{\raggedright}p{4.0cm}>{\raggedright}p{5.6cm}>{\raggedright}p{2.3cm}>{\raggedright\arraybackslash}p{1.7cm}}
\toprule
序 & 命题 & 实验与测量内容 & 来源 & 状态 \\
\midrule\endfirsthead
\toprule
序 & 命题 & 实验与测量内容 & 来源 & 状态 \\
\midrule\endhead
1 & 不对称规则~\eqref{eq:asymmetric}收敛到点~\eqref{eq:expectile}；$\kappa=1/2$时为平均值，对概率为$1/2$的一滴奖励有$V=\kappa$（图~\ref{fig:expectiles}） & 点~\eqref{eq:expectile}是期望分位数（expectile），即正负偏差权重不同的最小二乘问题的解；本章例子的推导见本章。 & \cite{NeweyPowell1987}；本章推导 & 理论 \\
2 & 不同\nt{DOPA}神经元的翻转点不同，并按不对称性排序（命题~\ref{prop:expectile}） & 小鼠，腹侧被盖区经光遗传标记的\nt{DOPA}神经元，不同大小的奖励：爆发变为低谷的点因神经元而异；响应不对称性越大的神经元，翻转点越高；由一组神经元的响应可重建奖励分布的形状。这是离散的主要功能，则是假说。 & \cite{Dabney2020} & 已测量 \\
3 & 部分\nt{DOPA}神经元对厌恶事件以爆发响应 & 猴子：有的\nt{DOPA}神经元被奖励预示信号兴奋、被眼部吹气的预示信号抑制，另一些则对二者都兴奋；这两组位于中脑的不同部分。 & \cite{Matsumoto2009, BrombergMartin2010} & 已测量 \\
4 & 误差的模或新异性设定学习速率 & 所引文献中没有\nt{DOPA}神经元的重要性信号改变学习速率的直接实验；综述提出了“注意，重要”信号的作用。 & \cite{BrombergMartin2010} & 假说 \\
5 & 危险维度有单独的导线 & 小鼠：投射到纹状体尾部的\nt{DOPA}神经元对新异的强刺激响应，而不对奖励响应；毁损它们会减弱对新物体的回避。 & \cite{Menegas2018} & 已测量 \\
6 & 最优动作时间$\tau_a^*=\sqrt{C/\bar R}$~\eqref{eq:vigor} & 和式$C/\tau_a+\bar R\tau_a$的最小值；本章推导。在原始模型中，动作速度由等于平均奖励的时间代价决定。 & \cite{Niv2007}；本章推导 & 理论 \\
7 & 多巴胺慢电平携带$\bar R$并设定动作速度（命题~\ref{prop:multiplex}） & 大鼠，伏隔核中的微透析和伏安法：多巴胺水平逐分钟跟随奖励频率和动作速度。在人类中，L-DOPA增强反应时间对平均奖励频率的依赖。慢电平恰好等于$\bar R$，尚未证明。 & \cite{Hamid2016, Beierholm2013vigor} & 假说 \\
8 & 慢电平由两个来源组成：输出端的释放可在脉冲频率不变时改变，但频率本身也有缓慢上升 & 大鼠，同一任务：伏隔核中的释放跟随变化的奖励期望，而腹侧被盖区已鉴定\nt{DOPA}神经元的脉冲频率则不然。虚拟走廊中的小鼠（第~10~行）：已鉴定\nt{DOPA}神经元的脉冲中也有接近奖励时的平缓上升。 & \cite{Mohebi2019, Kim2020} & 已测量 \\
9 & 动物接近远处奖励时多巴胺平缓上升 & 迷宫中的大鼠，纹状体伏安法：随着接近目标，多巴胺浓度在几秒内上升，斜率取决于距离和奖励大小。 & \cite{Howe2013ramp, Hamid2016} & 已测量 \\
10 & 上升是估计修正时的误差$\delta$~\eqref{eq:ramp}；走廊中的跳变产生爆发（图~\ref{fig:ramp}） & 公式和每秒$\delta=0.75\,V$这个数是本章推导。实验：虚拟走廊中的小鼠，瞬间移到离目标更近处；胞体脉冲、胞体和输出端的钙信号以及纹状体中的浓度都表现为误差，而不是期望$V$。两种解释哪一种在所有输出上都正确，仍在讨论中。 & \cite{Kim2020} & 已测量 \\
11 & 向后看：多巴胺报告因果联系的度量~\eqref{eq:retro} & 小鼠，在该理论与误差~\eqref{eq:ctd}预测不同的实验中测量伏隔核的多巴胺释放：在若干实验中释放遵循前者。 & \cite{Jeong2022} & 假说 \\
12 & 学习中多巴胺的响应逐渐从奖励移到预示信号 & 小鼠，学习过程中腹侧纹状体内\nt{DOPA}神经元输出端的信号：爆发逐渐从奖励时刻移到预示信号时刻，正如按~\eqref{eq:ctd}学习所要求的。 & \cite{Amo2022} & 已测量 \\
13 & 在价值不变时，\nt{DOPA}神经元对奖励味道的替换有响应 & 大鼠，腹侧被盖区神经元记录：把果汁换成价值相等的另一种会引起响应，尽管误差~\eqref{eq:ctd}为零。 & \cite{Takahashi2017} & 已测量 \\
14 & 人为诱发的爆发能教会两个中性刺激之间的联系 & 大鼠，无奖励的两刺激预配对实验：在从第一个刺激过渡到第二个时光遗传激活\nt{DOPA}神经元会建立联系，抑制它们则妨碍学习。 & \cite{Sharpe2017} & 已测量 \\
15 & 按特征的误差向量~\eqref{eq:vectortd} & 特征预测误差理论；向量形式没有直接检验。 & \cite{Gardner2018} & 假说 \\
16 & 同一任务中不同\nt{DOPA}神经元携带不同的量 & 虚拟走廊中的小鼠，\nt{DOPA}神经元双光子成像：有的与奖励相关，有的与速度和转向相关，有的与预示信号和选择准确性相关。 & \cite{Engelhard2019} & 已测量 \\
17 & 线束代替导线（命题~\ref{prop:bundle}） & 第~2--16~行的汇总：每种分解都已分别测得；它们都是同一信号的组成部分，这一统一图景没有实验支持。 & --- & 假说 \\
\bottomrule
\end{longtable}}

\section{第~\ref{sec:nora}~章：\nt{NORA}}

\nt{NORA}系统的构造、它的两种模式、与瞳孔的联系（不仅经由\nt{NORA}）、信号与背景之比的提高以及行为中的倒U形都是直接测量的；乘法增益$g$是模型；增益切换选择模式并起逆温度的作用，是本章的推论；调节增益的收益是按本书模型的计算；“\nt{NORA}是探索的旋钮”和“\nt{NORA}是中断”是假说。

{\footnotesize
\setlength{\tabcolsep}{3pt}\renewcommand{\arraystretch}{1.15}
\begin{longtable}{>{\raggedright}p{0.5cm}>{\raggedright}p{4.0cm}>{\raggedright}p{5.6cm}>{\raggedright}p{2.3cm}>{\raggedright\arraybackslash}p{1.7cm}}
\toprule
序 & 命题 & 实验与测量内容 & 来源 & 状态 \\
\midrule\endfirsthead
\toprule
序 & 命题 & 实验与测量内容 & 来源 & 状态 \\
\midrule\endhead
1 & 人类的\nt{NORA}神经元有几万个，几乎都在一个核团中 & 人脑干连续切片，酪氨酸羟化酶染色：蓝斑中有$53\,900$个细胞，相邻亚核中另有$6\,260$个。 & \cite{Baker1989LC} & 已测量 \\
2 & 输出几乎散布到全脑，但核团的不同部分在一定程度上服务于不同区域 & 小鼠的病毒—遗传标记：蓝斑的\nt{NORA}神经元接收来自脑内广泛区域的输入，并向脑和脊髓的广泛区域发出输出。大鼠：投射到杏仁核的一组增强恐惧的习得，而投射到前额叶皮层的另一组增强恐惧的消退。 & \cite{Schwarz2015LC, Uematsu2017LC, Poe2020} & 已测量 \\
3 & 紧张性模式：平稳的脉冲，频率从每秒零点几个到几个，随觉醒程度变化 & 大鼠，在睡眠—觉醒周期中记录蓝斑神经元：清醒时频率最高，慢波睡眠中较低，快速眼动睡眠中几乎为零；频率的变化先于状态的转换。 & \cite{AstonJonesBloom1981} & 已测量 \\
4 & 时相性模式：对任务中的重要事件产生短促的簇放电，大约在决策时刻 & 猴子，寻找稀有目标的任务：所有蓝斑神经元只对目标以簇放电响应，在目标出现后$\sim90\ms$、手部反应前$\sim200\ms$；表现差时簇放电较弱。在二选一任务中，簇放电与反应的时间关系比与刺激的更紧密，正确和错误反应时都有。 & \cite{AstonJones1994vigilance, Clayton2004LC} & 已测量 \\
5 & 瞳孔直径是\nt{NORA}的不精确测试点：它既跟随\nt{NORA}，也跟随\nt{ACET}和其他区域 & 猴子：蓝斑的脉冲，直至单个细胞的单个脉冲，都能预测瞳孔扩大；上丘和扣带皮层中也有类似但更晚的联系。小鼠：瞳孔的波动既跟随皮层中去甲肾上腺素能输出的活动，也跟随胆碱能输出的活动。 & \cite{Joshi2016, Reimer2016} & 已测量 \\
6 & 去甲肾上腺素对背景活动的抑制强于对诱发活动的抑制；乘法增益~\eqref{eq:gain}是再现这一效应的模型 & 清醒猴子，听觉皮层，离子电泳去甲肾上腺素：神经元背景活动受到的抑制强于对同类叫声的响应——信噪比升高。乘法S形函数~\eqref{eq:gain}取自一个网络模型；斜率被真的乘上系数，没有测量过。 & \cite{Foote1975NE, ServanSchreiber1990} & 已测量；$g$为模型 \\
7 & 增益只对放大器之后的噪声提高$d'$~\eqref{eq:gainsnr}（命题~\ref{prop:gainnoise}） & 本章推导；在网络模型中增益改善信号检测。没有把放大器前后的噪声分开并检验这一区别的实验。 & 本章推导；\cite{ServanSchreiber1990} & 理论 \\
8 & 边界$g\beta=1$把选择网络从“斟酌”切换到“已决断”~\eqref{eq:wtagain}（命题~\ref{prop:gainwta}，图~\ref{fig:pitchfork}） & 两个相互抑制的组的线性分析；本章推导。脑中没有对这一阈值的直接测量。 & 本章推导 & 理论 \\
9 & 增益是选择的逆温度~\eqref{eq:softmax} & 放大器之后的噪声服从逻辑斯谛分布时，选择概率是$T=\sigma/g$的softmax；本章推导。 & 本章推导 & 理论 \\
10 & \nt{NORA}决定探索多少：高紧张性活动对应小的有效$g$（探索），决策时刻的时相性簇放电对应大的（决断） & 人类：在随机选择之前的基线瞳孔比选择已知最佳项之前的大；瞳孔较大的人探索倾向更高。大鼠：转入随机选择模式受\nt{NORA}向前扣带皮层的输入控制。簇放电提高有效$g$，没有直接测量。 & \cite{JepmaNieuwenhuis2011, Tervo2014stochastic, AstonJones2005} & 假说 \\
11 & 奖励与$g$呈倒U形关系；在快速变化的世界中最佳$g$更小（命题~\ref{prop:explore}，图~\ref{fig:invertedU}） & \texttt{models/nora\_explore.py}，$60$次运行，每次$4000$次尝试。每$1000$次尝试变化一次：最佳$g=16$，比例$0.976$；每$20$次：$g=8$，比例$0.886$；$g=0.5$时为$0.77$。重算结果与本章一致。 & \texttt{models/\allowbreak nora\_\allowbreak explore.py} & 模型 \\
12 & 行为中的倒U形：紧张性活动适中、簇放电清晰时表现最好 & 猴子，与第~4~行相同的任务：表现好的时期紧张性频率适中，对目标的簇放电清晰；紧张性频率高时没有簇放电，错误反应更多。紧张性和诱发活动的变化跟随表现质量的波动。 & \cite{Usher1999LC, AstonJones2005} & 已测量 \\
13 & \nt{NORA}报告非预期的不确定性，\nt{ACET}报告预期的不确定性 & 含两种不确定性的贝叶斯推断理论；所引文献中没有按递质区分这两种信号的直接实验。 & \cite{YuDayan2005} & 假说 \\
14 & 在突然变化之后人学得更快，瞳孔扩大 & 人类，带突变的预测任务：瞳孔的短暂变化跟随对变化概率的估计，并与基线瞳孔一起预测新数据对估计的改变程度（学习速率）；与光照和任务无关的瞳孔变化会改变这一速率。瞳孔在这里显示的正是\nt{NORA}，是根据第~5~行间接数据的推论。 & \cite{Nassar2012} & 已测量 \\
15 & 时相性簇放电起中断作用：网络复位状态 & 没有\nt{NORA}簇放电使网络状态复位的直接实验；只测得了对重要事件的簇放电（第~4~行）。 & \cite{BouretSara2005} & 假说 \\
16 & 按意外调节$g$或学习速率，几乎不比最佳恒定设置更好 & \texttt{models/nora\_explore.py}，每个时期$1000$次尝试的世界（每$1000$次和每$20$次变化一次）。最佳恒定$g=8$：$0.925$；调节$g$：最高$0.927$；调节学习速率：最高$0.926$；恒定速率$0.4$：$0.928$。重算结果与本章一致。 & \texttt{models/\allowbreak nora\_\allowbreak explore.py} & 模型 \\
\bottomrule
\end{longtable}}

\section{第~\ref{sec:acet}~章：加入\nt{ACET}}

乙酰胆碱的三种作用已在脑片和活体脑中测得，写入与读出的冲突在本书模型中得到展示，而\nt{ACET}充当写使能线并报告预期的不确定性，是有部分实验支持的假说。

{\footnotesize
\setlength{\tabcolsep}{3pt}\renewcommand{\arraystretch}{1.15}
\begin{longtable}{>{\raggedright}p{0.5cm}>{\raggedright}p{4.0cm}>{\raggedright}p{5.6cm}>{\raggedright}p{2.3cm}>{\raggedright\arraybackslash}p{1.7cm}}
\toprule
序 & 命题 & 实验与测量内容 & 来源 & 状态 \\
\midrule\endfirsthead
\toprule
序 & 命题 & 实验与测量内容 & 来源 & 状态 \\
\midrule\endhead
1 & 脑中的\nt{ACET}神经元不多，集中在几个核团中，而输出遍布整个皮层：这是广播通道 & 大鼠，乙酰胆碱合成酶抗体染色和逆行标记：皮层、海马、杏仁核、嗅球和丘脑的乙酰胆碱主要来自基底前脑和脑干的六个紧凑细胞群 & \cite{Mesulam1983} & 已测量 \\
2 & 皮层中乙酰胆碱水平在清醒时高，在慢波睡眠中低 & 自由活动猫的皮层和海马微透析并同步记录脑电：安静清醒时乙酰胆碱释放比慢波睡眠高$\sim100\%$，活跃清醒时高$\sim175\%$；快速眼动睡眠中皮层的释放与清醒时相当 & \cite{Marrosu1995,Hasselmo1999} & 已测量 \\
3 & 信号的含义由受体决定：乙酰胆碱既有快速的通道型受体，也有在细胞内启动反应的慢速受体（命题~\ref{prop:receptor}） & 测量综述：乙酰胆碱对兴奋性、传递和可塑性的作用取决于释放部位、受体亚型（烟碱型为离子通道，毒蕈碱型经 G 蛋白）和接收细胞 & \cite{Picciotto2012} & 已测量 \\
4 & 第一种作用：削弱内部连接，几乎不影响外部输入（\eqref{eq:acetnet}中的因子$1-a$） & 大鼠嗅皮层脑片：$100$\,\textmu M 乙酰胆碱激动剂使同一细胞中内部连接（Ib 层）的电位下降$60\%$以上，外部输入（Ia 层）的下降不到$15\%$；抑制发生在突触前。海马脑片（CA1）：内部通路和输入通路分别为$89\%$和$40\%$ & \cite{HasselmoBower1992,HasselmoSchnell1994} & 已测量 \\
5 & 第二种作用：增强输入（因子$\frac{1+a}{2}$） & 新皮层脑片：烟碱型受体增强来自丘脑的输入突触，对皮层内突触无作用；毒蕈碱型受体削弱两类突触。乙酰胆碱提高神经元兴奋性（综述） & \cite{Gil1997,Hasselmo2006} & 已测量 \\
6 & 第三种作用：使权重更容易改变（速率$\eta a$） & 大鼠：电刺激基底核（皮层乙酰胆碱的来源）并与纯音配对，使成年动物的听觉皮层图谱按所呈现的声音发生显著重组 & \cite{KilgardMerzenich1998,Hasselmo2006} & 已测量 \\
7 & \eqref{eq:acetnet}第二个方程中针对稀疏模式的赫布规则给出可根据部分补全模式的联想记忆 & 采用稀疏模式和协方差规则的网络容量计算：活跃神经元比例$p$小时，网络存储的模式明显多于$p=1/2$时 & \cite{Tsodyks1988} & 理论 \\
8 & 低$a$时的写入损坏记忆：网络写下的是回忆出的而不是看到的；$a=0.1$时的损坏不比$a=0$时轻 & \texttt{models/\allowbreak acet\_memory.py}：$200$个神经元，五个各含$20$个神经元的模式，新模式与其中一个共有$10$个神经元，$10$次运行。$a=0$时新模式未被记住，旧模式带出$10$个外来神经元中的$5.9$个；$a=0.1$时新模式可回忆（$0.97$），但两者融合：新模式带出$7.3$个外来神经元，旧模式带出$7.9$个；从$a=0.2$起不到一个 & 本章推导 & 模型 \\
9 & 命题~\ref{prop:writeread}：不存在使写入和读出同样良好的$a$值（图~\ref{fig:acetsim}） & 同一脚本，学习速率$\eta a$：$a\ge0.9$时新模式一次呈现即被完整记住，$a=0.75$时为$0.8$，$a\le0.5$时记不住。凭一半模式读出：$a\le0.2$时为$1.0$，$a=0.3$时为$0.96$，$a=0.4$时为$0.04$ & 本章推导 & 模型 \\
10 & 在脑中，一根广播导线——\nt{ACET}——切换写入和读出 & 这一想法来自根据脑片测量建立的模型。最直接的实验在人类身上：毒蕈碱受体阻断剂东莨菪碱损害新词对的学习，对与旧词对重叠最多的影响最大，但不损害已学词对的回忆。没有对网络两种模式的直接测量 & \cite{HasselmoAndersonBower1992,Atri2004} & 假说 \\
11 & 滤波器~\eqref{eq:kalman}是最优的；输入权重和学习速率是同一个量$P/R$，稳态下$P=\sqrt{qR}$，记忆为$\sqrt{R/q}$（命题~\ref{prop:kalman}） & 不需要实验：卡尔曼—布西滤波器的最优性是估计理论的经典结果；稳态在本章推导 & \cite{KalmanBucy1961} & 理论 \\
12 & \nt{ACET}向网络报告一个类似$P$的量——预期的不确定性 & 在一个注意的概率模型中提出。没有乙酰胆碱水平跟随估计不确定性的直接测量 & \cite{YuDayan2005} & 假说 \\
13 & 部分\nt{ACET}神经元在二十毫秒左右内对奖励和惩罚作出响应，强化越出乎意料，响应越强 & 小鼠，经光遗传鉴定的基底前脑胆碱能神经元：对奖励和惩罚的响应潜伏期为$18\pm3$\,ms，幅度随强化的意外程度增大，两个核团的神经元响应相似 & \cite{Hangya2015} & 已测量 \\
14 & \nt{NORA}发现世界的意外变化，\nt{ACET}使网络保持在写入模式 & 这种分工是在模型中提出的。间接证据：人类与\nt{NORA}相关的瞳孔直径跟踪变化和学习速率。没有对\nt{NORA}--\nt{ACET}这一对的直接检验 & \cite{YuDayan2005,Nassar2012} & 假说 \\
15 & 在慢波睡眠中，处于读出模式的网络回放白天记录的内容，这些回放把它转写到长期记忆中 & 大鼠海马神经元集群记录：白天一起工作的细胞在随后的慢波睡眠中更常一起发放，且顺序相同——已测量。关于转写：学习后抑制海马涟漪会损害空间记忆，延长自发涟漪则改善它；原因是\nt{ACET}水平低，这一点未经证明 & \cite{WilsonMcNaughton1994,Skaggs1996,Girardeau2009,FernandezRuiz2019,Hasselmo1999} & 假说 \\
\bottomrule
\end{longtable}}

\section{第~\ref{sec:synthesis}~章　整台机器}

作为汇总的这一章不引入新的实验：表中每一行都依托于详细讨论该命题的那一章，以及同样的文献；\nt{GLUT}总线和经由\nt{GABA}的数据流控制已经确立，而各广播通道的角色及其协同工作主要还是假说。

{\footnotesize
\setlength{\tabcolsep}{3pt}\renewcommand{\arraystretch}{1.15}
\begin{longtable}{>{\raggedright}p{0.5cm}>{\raggedright}p{4.0cm}>{\raggedright}p{5.6cm}>{\raggedright}p{2.3cm}>{\raggedright\arraybackslash}p{1.7cm}}
\toprule
序号 & 命题 & 实验及测量内容 & 文献 & 状态 \\
\midrule\endfirsthead
\toprule
序号 & 命题 & 实验及测量内容 & 文献 & 状态 \\
\midrule\endhead
1 & $8.6\cdot10^{10}$个神经元只对应四类控制信号~\eqref{eq:machine} & 在脑匀浆中计数细胞核（各向同性分离法）：成年男性有$(86.1\pm8.1)\cdot10^9$个神经元 & \cite{Azevedo2009} & 已测量 \\
2 & 命题~\ref{prop:architecture}的前两层：数据沿\nt{GLUT}寻址总线传送，连接的符号由发送端决定，数据流由\nt{GABA}神经元引导和归一化（第~\ref{sec:neurontypes}、\ref{sec:gaba}~章） & 每个神经元一种主要递质（测量综述）；前馈抑制缩窄锥体神经元整合输入的时间窗；除以总活动的归一化在许多脑区都已测到 & \cite{Strata1999,Pouille2001,Carandini2012} & 已测量 \\
3 & 元件不可靠：突触丢失大部分脉冲（表~\ref{tab:computer}，第~\ref{sec:code}~章） & 海马切片，最小刺激：一半以上的脉冲在CA1不引起反应；排除了阈值失败和轴突传导失败，原因是递质释放的随机性 & \cite{AllenStevens1994} & 已测量 \\
4 & 大脑以$\sim20$瓦工作，平均每个神经元每秒约一个脉冲（第~\ref{sec:code}~章）；工程师尚未造出这样的机器 & 由实测能耗得出估计~\eqref{eq:energy}：每个脉冲连同突触工作约耗$10^9$个ATP分子（啮齿类皮层）；对皮层的细致估计给出更低的频率。技术现状依据综述 & \cite{Attwell2001,Lennie2003,MehonicKenyon2022} & 理论 \\
5 & 大多数突触的学习是局部资格迹与一个公共数$\delta$的乘积（式~\eqref{eq:machine}第二个方程，第~\ref{sec:dofa}~章） & 猴的\nt{DOPA}神经元对奖励预测误差作出反应；在纹状体切片中，只有当多巴胺在谷氨酸输入之后$0.3$--$2$\,s到达时，树突棘才会增大——资格迹已被测到 & \cite{Schultz1997,Yagishita2014} & 已测量 \\
6 & 只有动作学习回路才有通往每个输出的单独误差导线（第~\ref{sec:motorlearning}~章） & 小脑：攀缘纤维信号与某一输入同时出现会削弱该输入；在猴身上，浦肯野细胞的复合脉冲携带动作误差的方向并引起校正 & \cite{Ito2001,Herzfeld2018} & 已测量 \\
7 & 每条广播通道都是由几条线路组成的线束（命题~\ref{prop:bundle}） & 对\nt{DOPA}：同一任务中不同神经元分别携带奖励、运动和刺激特征；快速误差与慢速多巴胺水平各不相同。对\nt{SERO}、\nt{NORA}、\nt{ACET}，这种分解研究得较少 & \cite{Engelhard2019,Hamid2016,Mohebi2019} & 已测量 \\
8 & \nt{SERO}：水平调节器，按照假说还设定时间范围$\tau_\gamma$（第~\ref{sec:serocomputer}~章） & 自抑制：自身5-HT$_{1A}$受体的激动剂抑制中缝核血清素神经元——已测量。时间范围由理论提出；最有力的实验是在小鼠中光遗传激活这些神经元，延长了对延迟奖励的等待 & \cite{Sprouse1987,Doya2002,Miyazaki2014} & 假说 \\
9 & \nt{NORA}：增益$g$在探索与决断之间选择（第~\ref{sec:nora}~章） & 信噪比提高：在清醒猴的听觉皮层中，去甲肾上腺素对背景活动的抑制强于对同类叫声反应的抑制——已测量；乘性增益是模型；在探索与决断之间的选择作用是理论 & \cite{Foote1975NE,ServanSchreiber1990,AstonJones2005} & 假说 \\
10 & \nt{ACET}：切换写入和读取（第~\ref{sec:acet}~章） & 三种作用（削弱内部连接、增强输入、促进可塑性）已测量；模式切换得到人类东莨菪碱实验的间接支持 & \cite{HasselmoBower1992,Gil1997,Atri2004} & 假说 \\
11 & 公式~\eqref{eq:machine}描述整台机器 & 这是第~\ref{sec:glutcomputer}--\ref{sec:acet}~章各模型的汇总，每一部分都依托于各自的章节；作为整体未经实验检验 & 本章推导 & 假说 \\
12 & 旋钮在没有统一控制的情况下协同工作：误差、意外、写入、复原（图~\ref{fig:timeline}） & 没有实验：示意图是定性的。单个环节有依据：\nt{NORA}与\nt{ACET}分工的理论，以及人类瞳孔与学习速度的关联 & \cite{YuDayan2005,Nassar2012} & 假说 \\
13 & 按一个公共信号学习时，影响结果的神经元越多，学得越慢（命题~\ref{prop:broadcastcost}） & 线性网络学习曲线的计算：按输出扰动学习比梯度下降慢，慢的倍数等于输出的个数 & \cite{Werfel2005} & 理论 \\
14 & 皮层如何调整多层回路尚不清楚；候选机制有簇放电、神经元分支中的计算、预测式自监督学习 & 簇放电作为误差载体——模型；只有$5$--$8$层的网络才能复现皮层神经元细致模型的响应——模型；人类皮层神经元分支中产生异或的钙脉冲——在切片中已测量；自监督学习——证据综述 & \cite{Payeur2021,Beniaguev2021,Gidon2020,Keller2018} & 假说 \\
\bottomrule
\end{longtable}}

\section{第~\ref{sec:sensors}~章　机器的输入}

传感器的结构已被直接测量——通过单细胞电流记录、耳蜗振动测量和视网膜细胞计数；灵敏度极限和散粒噪声公式来自物理学，而传感器编码被调节到信息最大，则是有有力支撑的假说。

{\footnotesize
\setlength{\tabcolsep}{3pt}\renewcommand{\arraystretch}{1.15}
\begin{longtable}{>{\raggedright}p{0.5cm}>{\raggedright}p{4.0cm}>{\raggedright}p{5.6cm}>{\raggedright}p{2.3cm}>{\raggedright\arraybackslash}p{1.7cm}}
\toprule
序号 & 命题 & 实验及测量内容 & 文献 & 状态 \\
\midrule\endfirsthead
\toprule
序号 & 命题 & 实验及测量内容 & 文献 & 状态 \\
\midrule\endhead
1 & 传感器是换能器：受体产生电流，眼睛和耳朵感觉细胞的输出是送上\nt{GLUT}总线的谷氨酸；最前端的细胞不发放脉冲（图~\ref{fig:transducer}） & 龟的视杆和视锥细胞释放兴奋性氨基酸：它被一小片切下的膜上的谷氨酸通道捕获。大鼠耳蜗内毛细胞：神经纤维末梢的电流经谷氨酸AMPA受体流过，释放频率随细胞的平缓去极化上升到$150$\,Hz & \cite{CopenhagenJahr1989,GlowatzkiFuchs2002} & 已测量 \\
2 & 暗处的光传感器以约一皮安的电流响应单个光子 & 用吸引电极记录蟾蜍视杆外段：对暗弱闪光的响应是量子化的，单个事件$\approx1$\,pA（暗电流的$5\%$），离散度$0.2$\,pA，效率$0.5$。人报告单个光子到达的正确率高于随机水平 & \cite{Baylor1979,Tinsley2016} & 已测量 \\
3 & 声音传感器能察觉小于一纳米的位移 & 穆斯堡尔法，豚鼠耳蜗基底膜$16$--$19$\,kHz处：在听神经响应阈值处，膜的速度约为$0.04$\,mm/s，即位移$\approx0.35$\,nm（我们的换算）。测量的是基底膜，而不是传感器的纤毛束 & \cite{Sellick1982,Hudspeth1989} & 已测量 \\
4 & 暗处的精度受光子散粒噪声~\eqref{eq:photonnoise}限制；光弱时积累时间更长 & 公式来自泊松统计。在视杆中，暗光下的噪声与独立量子事件之和相符；在明亮背景上，对闪光的响应更小、更短。“零点几秒”这个数字在此没有单独检验 & \cite{Baylor1979} & 理论 \\
5 & 输入范围：光为十个数量级，声为十二个数量级，而脉冲频率不到三个数量级 & 声音：基底膜在特征频率上的响应在$40$--$80$\,dB区间内斜率低至$0.2$\,dB/dB，$100$\,dB以上仍可见压缩。数量级本身是标准估计，本章未给出处 & \cite{Ruggero1997} & 已测量 \\
6 & 传感器响应为式~\eqref{eq:nakarushton}，而$\sigma$跟随平均亮度，使曲线沿对数轴移动（图~\ref{fig:adaptation}） & 公式最早拟合于鱼视网膜的S电位。龟的视锥细胞：响应服从$V=I V_m/(I+\sigma)$，覆盖$\sim3.5$个数量级的亮度；背景持续$2$--$3$分钟后，曲线水平移动而宽度不变。与除以总活动的联系见综述 & \cite{NakaRushton1966,NormannPerlman1979,Carandini2012} & 已测量 \\
7 & 传感器传送变化，其编码被调节到每个脉冲携带最多信息（命题~\ref{prop:sensors}） & 这一原则由理论提出。最有力的证据：苍蝇第一层神经元的“对比度$\to$响应”特性与自然场景对比度的累积分布一致——这样可获得最大信息 & \cite{Barlow1961,Laughlin1981} & 假说 \\
8 & 接通传感器和断开传感器构成双线编码；事件相机在硅片上复现了它 & 实验综述：视网膜的接通通道和断开通道在第一个突触处就已分开，并可分别关闭。相机：$128\times128$像素，无帧，动态范围$120$\,dB，延迟$15$\,\textmu s & \cite{Schiller1992,Lichtsteiner2008,Gallego2022} & 已测量 \\
9 & 眼睛有$\sim10^8$个光传感器和$\sim10^6$个输出神经元：压缩$\sim100$倍 & 在完整的人视网膜标本上计数：平均$92\cdot10^6$个视杆和$4.6\cdot10^6$个视锥；不同眼睛的输出（神经节）细胞为$0.7\cdot10^6$到$1.5\cdot10^6$个 & \cite{Curcio1990,CurcioAllen1990} & 已测量 \\
10 & 小鼠眼睛的输出通道超过三十种 & 小鼠：对$11\,000$多个输出细胞做双光子钙成像并对响应聚类——明显多于$30$种功能类型。人的数目尚未测量 & \cite{Baden2016} & 已测量 \\
11 & 豚鼠眼睛传送$\sim875\,000$\,bit/s；换算到人眼约为$\sim10^7$\,bit/s & 豚鼠，多电极阵列，自然场景中的运动：每个细胞$6$--$13$\,bit/s，$\sim10^5$个输出细胞共$\sim875\,000$\,bit/s。人（$\sim10^6$个细胞）的$\sim10^7$\,bit/s是作者的换算 & \cite{Koch2006} & 已测量 \\
12 & 耳朵是带通滤波器组，频率图近似对数（图~\ref{fig:filterbank}） & 基底膜最大振动的位置取决于频率；“频率—位置”函数近似指数型，用同一形式描述人和活体动物的耳蜗 & \cite{Greenwood1990,Robles2001} & 已测量 \\
13 & 谐振器是主动的：一部分传感器把微弱振动的振幅放大数千倍（$66$--$76$\,dB），响度增大时增益下降 & 激光测振，南美栗鼠耳蜗底部：弱纯音在特征频率上的增益相对镫骨为$66$--$76$\,dB；死亡使灵敏度降低$60$--$81$\,dB（振幅$10^3$--$10^4$倍）并消除压缩。动力来源是外毛细胞 & \cite{Ruggero1997,Robles2001} & 已测量 \\
14 & 耳朵的放大器工作在自激边缘 & 计算：在霍普夫分岔点附近，范围压缩、尖锐调谐和组合音不可避免——这正是听觉所有已知的非线性。耳蜗确实如此调节，尚无直接证明 & \cite{Eguiluz2000} & 假说 \\
15 & 低频时脉冲与声音同相（豚鼠的相位锁定从$\sim600$\,Hz开始减弱，直到$\sim3.5$\,kHz仍可察觉），并据此确定方向；高频时只剩位置编码 & 豚鼠听神经：相位锁定在约$600$\,Hz开始下降，$3.5$\,kHz以上消失；猫和猴的界限约高一个倍频程。双耳到达时间差见综述 & \cite{PalmerRussell1986,Grothe2010} & 已测量 \\
16 & 其他传感器（表~\ref{tab:sensors}）：嗅觉有$\sim400$种受体，每个神经元一种，组合编码；味觉部分经由ATP和血清素；触觉有快适应和慢适应传感器；冷和热各自独立 & 小鼠：一种受体识别多种气味，一种气味激活多种受体，不同气味对应不同集合。没有ATP受体，味觉神经不响应；味蕾释放血清素。人手皮肤单根纤维记录：按适应速度分为四类。冷通道不同于热通道 & \cite{Mombaerts2004,Malnic1999,Finger2005,Huang2005,JohanssonVallbo1979,McKemy2002} & 已测量 \\
\bottomrule
\end{longtable}}

\section{第~\ref{sec:recognition}~章　识别}

识别速度、最初几级的结构、响应的线性读出以及从时间连续性中学习，已在人和猴身上测量；把整条链归结为两种操作以及从视频中学习，已在本书的模型上检验，而对错觉的解释则从有充分依据到存在争议不等。

{\footnotesize
\setlength{\tabcolsep}{3pt}\renewcommand{\arraystretch}{1.15}
\begin{longtable}{>{\raggedright}p{0.5cm}>{\raggedright}p{4.0cm}>{\raggedright}p{5.6cm}>{\raggedright}p{2.3cm}>{\raggedright\arraybackslash}p{1.7cm}}
\toprule
序号 & 命题 & 实验及测量内容 & 文献 & 状态 \\
\midrule\endfirsthead
\toprule
序号 & 命题 & 实验及测量内容 & 文献 & 状态 \\
\midrule\endhead
1 & 平移时，逐像素与样板比较的正确率不比硬币好；一对级~\eqref{eq:scstage}能在任何位置认出图形 & \texttt{models/\allowbreak recognition\_models.py}，$24$点“视网膜”，$4000$次试验：翻转点比例为$0$、$0.1$、$0.2$时，样板法为$0.49$、$0.51$、$0.52$；$S$与$C$级对为$1.00$、$0.86$、$0.70$ & 本章推导 & 模型 \\
2 & 图片上的动物在$\sim150$\,ms内被认出，沿约十级走一遍，每级$10$--$20$\,ms & 人，对呈现$20$\,ms的照片判断“有无动物”：两种回答的诱发电位在$\sim150$\,ms后分开。猴：首次响应的时间逐级增加（V1最早，然后是V2和V4）。级数和“每级零个或一个脉冲”是由这些数字推出的 & \cite{Thorpe1996,Schmolesky1998} & 已测量 \\
3 & $S$级是对部件的“与”，$C$级是对位置取最大值（命题~\ref{prop:conveyor}） & 猫，初级视皮层：简单细胞响应特定位置、特定朝向的边缘，复杂细胞在更大区域内响应同一朝向。复杂细胞的细胞内记录：许多细胞对两根条纹的响应接近两个响应中较大的那个，平均而言最接近max运算。通过相互抑制取最大值见命题~\ref{prop:wta}，除以总活动见综述 & \cite{HubelWiesel1962,Lampl2004,Carandini2012} & 已测量 \\
4 & 沿链条向上，组合越来越大、越来越复杂，响应越来越不依赖位置 & 猴，把刺激简化为“关键特征”：V4和下颞叶皮层后部的细胞以较小的感受野响应形状、颜色和纹理的复杂组合；下颞叶皮层前部的感受野更大。模型中$S$与$C$的交替再现了这一图景 & \cite{KobatakeTanaka1994,Riesenhuber1999} & 已测量 \\
5 & 卷积网络重复了这种交替，并且对高层响应的预测最好；物体可以从神经元群体的频率中线性读出 & 为识别而训练、未对大脑做拟合的网络：顶层预测猴下颞叶皮层神经元的响应，中间层预测V4的响应。基于$\sim100$个随机选取细胞的活动、时间窗从$12.5$\,ms起的线性分类器，能在不同位置和大小下认出物体及其类别 & \cite{Fukushima1980,Yamins2014,Hung2005} & 已测量 \\
6 & 若资格迹不短于$\sim20$\,ms，带资格迹的赫布规则~\eqref{eq:traceinvariance}能从无标注视频中学会不变性 & 慢特征和资格迹规则的思想是理论。\texttt{models/\allowbreak recognition\_models.py}，两个$C$神经元，$16$个输入，运行$10$次：物体间停顿$0.3$\,s。$\tau_{\text{tr}}=5$\,ms时与图形的吻合平均为$0.52$，$10$\,ms时为$0.56$；$20$、$50$和$200$\,ms时所有运行均为$1.00$（$30$\,ms时十次中有一次出错） & \cite{Foldiak1991,WiskottSejnowski2002} & 模型 \\
7 & 大脑中的不变性是从时间的连续性中学到的 & 猴：在眼睛跳向视野某一位置时悄悄替换物体；下颞叶皮层神经元的位置不变性随替换而改变，一小时后即已显著。这是视觉学习的主要规则，则是假说 & \cite{LiDiCarlo2008} & 已测量 \\
8 & 运动：方向检测器，然后在大区域上使用同样的“与/或”对 & 兔视网膜中的方向检测器；猴MT区记录到的全部$110$个神经元都有方向选择性，对运动的响应强于V1 & \cite{Barlow1965,Albright1984} & 已测量 \\
9 & 高层向下发送预测，向上传送误差 & 该模型解释了初级视皮层感受野的非经典效应；个别观察与之相符（综述），但作为流水线的总体结构未经证明 & \cite{RaoBallard1999,Keller2018} & 假说 \\
10 & 困难图片需要反馈和绕几圈 & 猴：对“困难”图片，下颞叶皮层的判决晚$\sim30$\,ms形成；这些较晚的响应用前馈网络预测得差，用带反馈的或更深的网络预测得好 & \cite{Kar2019} & 已测量 \\
11 & 难以察觉的像素篡改能欺骗网络；人的视觉稳健得多，但并非无懈可击 & 在网络上：专门挑选的微小扰动改变答案；“熊猫$\to$长臂猿”的例子出自第二篇文献。对人，在短暂呈现时，能在网络间良好迁移的篡改会改变选择 & \cite{Szegedy2014,Goodfellow2015,Elsayed2018} & 已测量 \\
12 & 预期错觉：不可靠的输入让位于预期——暗淡的物体移动“更慢”，连衣裙被看成不同颜色（第~\ref{sec:illusions}~节） & 对比度较低的光栅看起来移动得较慢。对$1401$人的调查：连衣裙被稳定地看成三种样子，光照提示可切换颜色；在$\sim13\,000$名观察者中，起决定作用的是对光照的假定。预期慢速运动的贝叶斯模型解释了一系列运动错觉 & \cite{Thompson1982,LaferSousa2015,Wallisch2017,Weiss2002,Purves2011} & 假说 \\
13 & 增益调节：瀑布、倾斜和对比度；赫尔曼栅格不是邻域抑制 & 兔视网膜的方向检测器在长时间运动中降低频率，运动停止后降到背景以下。倾斜错觉可由除以邻居活动的模型再现。不改变视网膜输出神经元响应的栅格改动消除了暗斑——邻域抑制的解释被否定 & \cite{BarlowHill1963,Schwartz2009,SchillerCarvey2005} & 已测量 \\
14 & 延迟补偿：运动物体看起来在闪光的前方 & 错觉已测量。心理物理学结果既与运动外推矛盾，也与延迟差矛盾；有人提出事后解释——依据闪光后$\sim80$\,ms内的事件。争论尚无定论 & \cite{Nijhawan1994,EaglemanSejnowski2000} & 假说 \\
15 & 静止的“蛇”在旋转：延迟差~\eqref{eq:latency}被方向检测器读出 & 错觉在无色时也存在；相邻元素对单独就能产生它；猴皮层中具有方向选择性的神经元对这些静止的元素对给出方向性响应。在人身上，旋转由眼睛的微跳视和眨眼触发 & \cite{Conway2005,OteroMillan2012} & 已测量 \\
16 & 双稳图形：相互抑制、疲劳和噪声；切换主要由噪声引起 & 竞争神经元群体的模型：其中的切换由\emph{噪声}引起，没有噪声就没有切换；它解释了切换频率随刺激强度增加。疲劳在这里作用较小 & \cite{MorenoBote2007} & 假说 \\
17 & 一部分错觉是机器所学任务的结果 & 训练来预测视频下一帧的网络“看到”静止圆环在旋转，而在对照图片上看不到；用于图像去噪和锐化的网络像人一样受颜色和亮度错觉影响 & \cite{Watanabe2018,GomezVilla2020} & 已测量 \\
\bottomrule
\end{longtable}}

\section{第~\ref{sec:muscles}~章　通往肌肉的输出}

输出的结构——运动单位、突触的安全系数、按大小启动、成对的肌肉、牵拉环路和节律发生器——都已直接测量；有延迟环路的稳定条件是一条定理；身体的内部模型是有充分依据的假说；图中的数字是按本书模型计算的。

{\footnotesize
\setlength{\tabcolsep}{3pt}\renewcommand{\arraystretch}{1.15}
\begin{longtable}{>{\raggedright}p{0.5cm}>{\raggedright}p{4.0cm}>{\raggedright}p{5.6cm}>{\raggedright}p{2.3cm}>{\raggedright\arraybackslash}p{1.7cm}}
\toprule
序号 & 命题 & 实验及测量内容 & 文献 & 状态 \\
\midrule\endfirsthead
\toprule
序号 & 命题 & 实验及测量内容 & 文献 & 状态 \\
\midrule\endhead
1 & 运动单位：一个\nt{ACET}神经元控制一组纤维，从几百根到两千来根；在精细调节的肌肉中单位较小 & 人的肌肉，计数运动轴突和肌纤维：每个神经元平均约有$340$根纤维（手的骨间肌）、约$410$根（肱桡肌）、约$1930$根（腓肠肌内侧头）。本章没有给出最小单位的确切数字。 & \cite{Feinstein1955mu} & 已测量 \\
2 & 肌肉上的突触释放的递质比使纤维兴奋所需的多好几倍 & 神经肌肉突触测量的综述：成年哺乳动物的安全系数（释放的递质与阈值所需之比）通常为$3$--$5$；在人身上，安全系数更多依靠接收方膜的褶皱，而不是释放量。 & \cite{WoodSlater2001} & 已测量 \\
3 & 抽搐~\eqref{eq:twitch}在量级为$50\ms$的时间$T_c$内上升 & 人手的单个运动单位，自主用力，按该单位的脉冲对力作平均：上升时间$30$--$100\ms$，$80\,\%$的单位小于$70\ms$。函数$(t/T_c)e^{1-t/T_c}$是运动单位池模型中的近似。 & \cite{MilnerBrown1973rec, Fuglevand1993} & 已测量 \\
4 & 抽搐叠加~\eqref{eq:forcesum}：$5$、$15$和$40$次/秒时，力分别为$0.2$--$1.1F_1$、$1.8$--$2.2F_1$和约$5.4F_1$（图~\ref{fig:tetanus}） & 按\texttt{models/\allowbreak muscle\_\allowbreak models.py}计算，$T_c=50\ms$：最后$0.3\,\text{s}$内为$0.21$--$1.10$、$1.76$--$2.19$和$5.28$--$5.49\,F_1$。线性叠加是简化：模型中没有真实肌肉的饱和。 & \texttt{models/\allowbreak muscle\_\allowbreak models.py} & 模型 \\
5 & 单位按大小启动；最弱的比最强的弱一百倍 & 猫的腹根：反射输入逐渐增强时，在腹根中脉冲较小的神经元，即小神经元，最先放电。人手骨间肌：各单位的抽搐力从$0.1$到$10$\,g，并随该单位启动时的力近似线性增长。 & \cite{Henneman1957, MilnerBrown1973rec} & 已测量 \\
6 & 力的步长与力成正比，$\Delta F/F\approx q-1$~\eqref{eq:sizeprinciple}：浮点 & 本章由力的等比数列推出。与实验相符：用力大时，同样的力增量所需启动的新单位要少得多。 & 本章推导；\cite{MilnerBrown1973rec} & 理论 \\
7 & 力的离散程度与力本身成正比 & 拇指肌肉的自主等长用力：力的标准差随平均力线性增长；同样的用力若由电刺激引起，离散程度不变。运动单位池模型：按力的顺序启动单位导致线性增长。 & \cite{Jones2002sdn} & 已测量 \\
8 & 动作的平滑性是信号依赖噪声的结果 & 模型：在这种噪声下使终点离散最小的轨迹，再现了实测的眼动和手臂轨迹。尚无把这一原因与其他原因区分开来的直接实验。 & \cite{HarrisWolpert1998} & 假说 \\
9 & 一对肌肉的力之差决定力矩，力之和决定刚度~\eqref{eq:antagonists} & 通过共同收缩进行阻抗控制的理论。人的手臂：用小推力测得的每个关节的刚度大致与该关节的力矩成正比；不同的共同收缩比例改变手的刚度椭圆的形状和方向。 & \cite{Hogan1984, GomiOsu1998stiff} & 已测量 \\
10 & 牵拉传感器兴奋自己的\nt{ACET}神经元，并经抑制性中间神经元关闭拮抗肌（图~\ref{fig:antagonists}） & 猫的脊髓：由牵拉传感器纤维兴奋的单个中间神经元，在拮抗肌的\nt{ACET}神经元中产生单突触抑制电位，突触延迟$0.28$--$0.42\ms$。该实验中没有确定中间神经元的递质（甘氨酸）。 & \cite{Jankowska1972Ia} & 已测量 \\
11 & 环路延迟：脊髓$25$--$30\ms$，经皮层长一倍半到三倍，经眼睛$100$--$200\ms$ & 人的手臂，对关节施加推力：肌肉的短潜伏期响应在$20$--$45\ms$，长潜伏期响应在$45$--$105\ms$，自主响应在$120$--$180\ms$。运动中移动手指的可见位置：平均$160\ms$后作出校正。 & \cite{Pruszynski2008rapid, SaundersKnill2003vis} & 已测量 \\
12 & 当$Gd<\pi/2$时，$\dd x/\dd t=-Gx(t-d)$稳定~\eqref{eq:delaystable}；在边界上频率为$1/(4d)$ & 关于时滞方程根的定理。用\texttt{models/\allowbreak muscle\_\allowbreak models.py}检验，$d=30\ms$：到$3\,\text{s}$时，$G=40$、$50$、$55$每秒对应的幅度为$3\cdot10^{-6}$、$0.13$、$38$（边界为$52$）。 & \cite{Hayes1950} & 理论 \\
13 & 生理性震颤$8$--$12$\,Hz；牵拉环路是其来源之一 & 测量综述：健康人都有约$10$\,Hz范围的细微震颤；其来源是混合的——中枢振荡、单位放电特性、机械共振和反射环路共振，不同条件下主导因素不同。 & \cite{McAuleyMarsden2000} & 假说 \\
14 & 大脑依据身体的内部模型预测指令的结果 & 人用手指捏住重物移动手臂：在惯性、黏性和弹性负载下，握力与负载同时变化，而不是滞后一个环路延迟，也就是说负载被预测到了。综述汇总了这类实验；尚未在神经元中直接观察到模型本身。 & \cite{FlanaganWing1997grip, WolpertGhahramani2000} & 假说 \\
15 & 步行、呼吸、咀嚼的节律由局部网络在恒定输入下产生 & 实验综述：离体神经系统（脊髓、神经节）在没有节律性输入、也没有肌肉反馈的情况下输出节律性的运动模式。 & \cite{MarderBucher2001} & 已测量 \\
16 & 带疲劳的相互抑制~\eqref{eq:halfcentre}产生周期为$0.84\,\text{s}\approx2\tau_a$的节律（图~\ref{fig:halfcentre}） & 定理：在恒定输入下没有稳定状态、带有相互抑制和适应的网络必然振荡。按\texttt{models/\allowbreak muscle\_\allowbreak models.py}计算（$I=1$，$\beta=2$，$b=2$，$\tau=20\ms$，$\tau_a=0.4\,\text{s}$）：周期$0.84\,\text{s}$。真实的发生器正是如此构成，尚未得到证明。 & \cite{Matsuoka1985osc} & 模型 \\
\bottomrule
\end{longtable}}

\section{第~\ref{sec:pain}~章　疼痛}

痛觉传感器、两根导线、脊髓闸门、下行的音量旋钮、敏化和对注意力的夺取都已直接测量；闸门和敏化的公式是本书的简化模型；机器选择警报音量的规则是假说。

{\footnotesize
\setlength{\tabcolsep}{3pt}\renewcommand{\arraystretch}{1.15}
\begin{longtable}{>{\raggedright}p{0.5cm}>{\raggedright}p{4.0cm}>{\raggedright}p{5.6cm}>{\raggedright}p{2.3cm}>{\raggedright\arraybackslash}p{1.7cm}}
\toprule
序号 & 命题 & 实验及测量内容 & 文献 & 状态 \\
\midrule\endfirsthead
\toprule
序号 & 命题 & 实验及测量内容 & 文献 & 状态 \\
\midrule\endhead
1 & 高阈值：不同痛觉传感器的热阈值从$41^\circ$到$46$--$48^\circ$ & 单纤维记录。猴的皮肤：慢（C）传感器的中位阈值为$41^\circ$，第二类快传感器为$46^\circ$，第一类高于$53^\circ$。人的皮肤：同时响应压力的C传感器为$40.7^\circ$，不响应压力的为$48^\circ$。 & \cite{Treede1995heat, Weidner1999cfib} & 已测量 \\
2 & 痛觉传感器的适应远弱于触觉传感器，轻度损伤后变得更敏感；强烈加热后的减弱只在一类中测到 & 皮肤轻度烫伤（$50^\circ$，$100\,\text{s}$）：$5$--$10$\,min后，人的痛阈和猴C传感器的阈值降低$2$--$6^\circ$；其他皮肤传感器没有这种偏移。第二类快传感器在强烈加热后响应受到抑制。与触觉传感器适应的比较是一般性估计，此处没有单独的实验支持。 & \cite{LaMotte1982hyper, Treede1995heat} & 已测量 \\
3 & 两根导线：快的$5$--$30$\,m/s，慢的$0.5$--$2$\,m/s；到达时间$t=L/v$~\eqref{eq:paintime} & 传导速度：猴的快痛觉传感器平均$15$和$25$\,m/s（两类）；人的C传感器$0.8$--$1.0$\,m/s。$L=1$\,m时，分别为几十毫秒和约一秒。 & \cite{Treede1995heat, Weidner1999cfib} & 已测量 \\
4 & 疼痛来两次：先快而准，后钝而久 & 人前臂受热的反应时间：随温度升高从$1100$降到$400\ms$；强烈加热有毛皮肤产生双重疼痛，手掌则不会。第一痛只能由快于$6$\,m/s的纤维传送——按潜伏期与之对应的是第二类快传感器，而手掌中没有这类传感器。角色分工（“哪里”和“有多糟”）是一种解读。 & \cite{CampbellLaMotte1983first, Treede1995heat} & 已测量 \\
5 & 闸门：脊髓输出神经元接收疼痛和触摸，抑制性中间神经元被触摸兴奋（图~\ref{fig:gate}） & 闸门方案作为理论提出。小鼠，遗传标记并去除神经元群：表达生长抑素的兴奋性神经元是机械痛所必需的；表达强啡肽的抑制性神经元阻止来自触觉传感器的输入到达它们——没有这些神经元，轻触就会引起疼痛。 & \cite{MelzackWall1965, Duan2014} & 已测量 \\
6 & 触摸减轻疼痛 & 人，手臂上的激光痛觉脉冲：同时触摸降低疼痛评分和分辨脉冲能量的能力；在同一皮节内，触摸点与疼痛点距离越远，效应越弱。 & \cite{Mancini2014touch} & 已测量 \\
7 & 闸门理论的假设：剧痛自己关闭抑制性中间神经元，闸门大开 & 本章标注为原始理论的假设。小鼠研究描述了闸门如何阻止触觉进入疼痛通道；其中没有显示疼痛关闭中间神经元。 & \cite{MelzackWall1965} & 假说 \\
8 & 闸门~\eqref{eq:gate}：无触摸阈值$0.18$，有触摸$0.86$，$g=2$时$0.11$；$\nu=1$时触摸使输出从$1$降到$0.25$，$\nu=2$时无效；触摸本身不通过（图~\ref{fig:gatecurves}） & 用正文中的权重按\texttt{models/\allowbreak pain\_\allowbreak gate.py}计算，再现了所有这些数字。 & \texttt{models/\allowbreak pain\_\allowbreak gate.py} & 模型 \\
9 & 来自脑干的下行通路双向改变闸门增益 & 大鼠延髓，加热尾部时记录：一些神经元在缩尾之前放电（“开”细胞），另一些沉默（“关”细胞）；微弱刺激该区域可抑制反射。综述：同样的回路既易化也抑制疼痛传递，阿片类经由它们起作用。 & \cite{Fields1983rvm, Fields2004} & 已测量 \\
10 & 对缓解的预期经阿片通道减轻疼痛 & 急性疼痛患者：对那些因预期缓解而疼痛减轻的人，阻断阿片受体使疼痛回到其他人的水平；对其他人，阻断不改变疼痛。 & \cite{Levine1978} & 已测量 \\
11 & 大脑根据中断的收益选择警报音量 & 没有测量过音量选择规则的实验。 & --- & 假说 \\
12 & 敏化：损伤后闸门输出增大，阈值降低，部分来自脊髓；在损伤性刺激停止后仍然持续 & 大鼠：皮肤损伤后屈曲反射阈值下降，响应增大；电生理分析显示，部分增长产生于脊髓本身。损伤性刺激引起持久的感觉改变，比刺激本身持续得更久。本章未给出确切的消退时间。 & \cite{Woolf1983} & 已测量 \\
13 & 敏化~\eqref{eq:sensitization}是正反馈；环路强时，警报可能在愈合后自锁 & 来自模型~\eqref{eq:sensitization}（本章末的习题）。没有实验表明愈合后的持续疼痛正是这种形式的自锁环路。 & 本章推导 & 假说 \\
14 & 疼痛是负奖励，从疼痛中学习遵循时间差分误差 & 人，断层扫描，疼痛先兆序列：腹侧纹状体和前岛叶皮层的活动与时间差分学习信号吻合。猴：一部分\nt{DOPA}神经元被不愉快的吹气及其先兆兴奋，另一部分被抑制。 & \cite{Seymour2004, Matsumoto2009} & 已测量 \\
15 & 疼痛中断当前工作 & 人，电痛刺激加声音探测：疼痛开始后$100\ms$时对探测的反应明显变慢，$1.5\,\text{s}$时已与对照无差别。综述把这类实验归纳为注意力夺取模型。 & \cite{Crombez1996disrupt, EcclestonCrombez1999} & 已测量 \\
16 & 缩回反射在脊髓中闭合 & 大鼠：背角深层神经元重复了单块肌肉缩回反射的空间模式；无法从颈段逆向激活它们，即它们是脊髓中间神经元。 & \cite{Schouenborg1995wr} & 已测量 \\
\bottomrule
\end{longtable}}

\section{第~\ref{sec:motorlearning}~章　运动学习}

最小均方规则和单独误差导线的优势是定理；误差导线线路的结构、同时出现时的削弱、资格迹与延迟的匹配、后效以及技能的自发恢复都已测量；整条线路作为有监督学习工作并建立内部模型，是主流假说；双过程模型的数字是按本书模型计算的。

{\footnotesize
\setlength{\tabcolsep}{3pt}\renewcommand{\arraystretch}{1.15}
\begin{longtable}{>{\raggedright}p{0.5cm}>{\raggedright}p{4.0cm}>{\raggedright}p{5.6cm}>{\raggedright}p{2.3cm}>{\raggedright\arraybackslash}p{1.7cm}}
\toprule
序号 & 命题 & 实验及测量内容 & 文献 & 状态 \\
\midrule\endfirsthead
\toprule
序号 & 命题 & 实验及测量内容 & 文献 & 状态 \\
\midrule\endhead
1 & 规则~\eqref{eq:lms}平均而言沿误差平方的梯度下降 & 自适应滤波器理论的结果：与输入和输出误差成正比的修正，就是对均方误差梯度的随机下降。 & \cite{WidrowHoff1960} & 理论 \\
2 & 每个输出有误差导线时，速度与输出数目无关；只有一个公共数时，速度随噪声神经元数目下降（命题~\ref{prop:teachers}） & 线性网络的学习曲线：按神经元的随机扰动学习，比精确梯度慢的倍数等于被扰动神经元的个数。 & \cite{Werfel2005} & 理论 \\
3 & 误差导线线路作为有监督学习工作（命题~\ref{prop:errorwire}） & 由线路结构推出的理论。第7--10行的实验与之相符；整条线路正是这样工作，尚未证明。 & \cite{Marr1969, Albus1971} & 假说 \\
4 & 扩展层：约$7\cdot10^{10}$个小型\nt{GLUT}神经元，比大脑其余部分加起来还多 & 在溶解组织的悬液中计数细胞核：人小脑有$69\cdot10^9$个神经元，全脑为$86\cdot10^9$个。体视学：老年男性小脑中有$101\cdot10^9$个小神经元。 & \cite{Azevedo2009, Andersen1992cereb} & 已测量 \\
5 & 升高输入维数使所需的关系变为线性 & 关于划分数目的定理：$n$个输入的带阈值加权和能实现约$2n$个一般位置向量的任意标注。小脑正是利用了这一点，属于第3行假说的一部分。 & \cite{Cover1965} & 理论 \\
6 & 约$3\cdot10^7$个输出\nt{GABA}神经元，每个约有$10^5$个输入 & 人小脑体视学：$30.5\cdot10^6$个输出神经元。电子显微镜，大鼠：每个输出神经元约有$1.75\cdot10^5$个来自扩展层的输入。输出神经元的抑制性递质见综述。 & \cite{Andersen1992cereb, NapperHarvey1988, Ito2001} & 已测量 \\
7 & 每个输出神经元恰好一根误差导线，约每秒一次，每次引起强烈爆发 & 记录综述：每个输出神经元接受一个攀缘\nt{GLUT}输入，它以约$1$\,Hz的频率引起复合脉冲。 & \cite{Ito2001} & 已测量 \\
8 & 爆发的概率取决于动作误差的方向 & 猴，小脑动眼区，扫视适应：视觉误差的方向决定复合脉冲的概率，误差大小决定其时间；爆发改变下一次尝试中的简单脉冲，并使动作沿该细胞偏好的误差方向偏移。 & \cite{Herzfeld2018} & 已测量 \\
9 & 误差爆发与活跃输入同时出现会削弱该输入（$-\eta e_i x_j$） & 实验综述：同时刺激扩展层输入和误差导线，引起该输入的长时程削弱；单独刺激一个输入则不会。 & \cite{Ito2001} & 已测量 \\
10 & 资格迹的长度与误差延迟匹配：延迟$100\ms$时，被削弱最多的是$100\ms$前活跃的输入 & 切片和活体小鼠：在负责动眼学习的区域，削弱被精确地调节到输入与误差之间约$120\ms$的间隔——正是这项任务中误差信号的传送时间；在另一区域，不同细胞调节到不同的间隔。 & \cite{Suvrathan2016} & 已测量 \\
11 & 线路是学习身体内部模型的自适应滤波器~\eqref{eq:adaptivefilter} & 由线路结构推出的模型。尚无把学到的滤波器作为身体传递函数来测量的直接实验。 & \cite{Fujita1982} & 假说 \\
12 & 预测减去自身动作的后果，所以不能给自己挠痒 & 人：自己的触摸比别人同样的触摸感觉更不痒；在扫描仪中，躯体感觉皮层对自己的触摸响应较弱，而当动作触及皮肤时小脑的活动较小。用减去预测来解释是假说。 & \cite{Blakemore1998} & 假说 \\
13 & 在外力作用下偏差减小，撤去外力后手臂偏向另一侧 & 人在力场中握着机器人手柄伸手：轨迹恢复为直线；撤去力场后，后效几乎是最初畸变的镜像；它还会迁移到未经训练的空间区域。 & \cite{ShadmehrMussaIvaldi1994} & 已测量 \\
14 & 两个过程在学习~\eqref{eq:tworate}；反向扰动之后，技能会自行恢复 & 人，力场，然后是短暂的反向力场和误差被钳制的试次：修正自行回到原来的值；双过程模型能再现这一点，单过程模型不能。 & \cite{Smith2006} & 已测量 \\
15 & 图~\ref{fig:tworate}的数字：$200$次动作后为$0.84$（慢过程$0.63$）；$6$次反向；快过程$-0.53$，慢过程$0.46$；$40$次动作内回升到$0.37$ & 按\texttt{models/\allowbreak motor\_\allowbreak learning.py}计算，$A_f=0.92$，$B_f=0.1$，$A_s=0.996$，$B_s=0.02$：$0.840$（$0.634$和$0.205$），$6$次动作，$-0.531$和$0.457$，第$40$次动作时峰值$0.371$。 & \texttt{models/\allowbreak motor\_\allowbreak learning.py} & 模型 \\
16 & 需要两个过程，是因为身体以不同速度变化 & 理论：在具有多种寿命的干扰下，最优估计给出多个时间常数不同的滤波器，并解释了自发恢复和加速的再学习。 & \cite{Kording2007} & 假说 \\
\bottomrule
\end{longtable}}

\section{第~\ref{sec:chemoutput}~章　化学输出}

通往心脏的两根导线及其不同的速度、激素级联中的负反馈、剂量频率编码、昼夜节律发生器及其光校准，都有直接测量；界限 $K<8$ 是本章推导的控制理论定理；由级联反馈本身产生的脉动是一个有强有力实验支持的假说。

{\footnotesize
\setlength{\tabcolsep}{3pt}\renewcommand{\arraystretch}{1.15}
\begin{longtable}{>{\raggedright}p{0.5cm}>{\raggedright}p{4.0cm}>{\raggedright}p{5.6cm}>{\raggedright}p{2.3cm}>{\raggedright\arraybackslash}p{1.7cm}}
\toprule
编号 & 论断 & 实验及测量内容 & 来源 & 状态 \\
\midrule\endfirsthead
\toprule
编号 & 论断 & 实验及测量内容 & 来源 & 状态 \\
\midrule\endhead
1 & 心脏起搏器自身每分钟跳动约 $100$ 次；静息时刹车踩着，频率通常约为 $60$--$70$ & 人体，完全阻断通往心脏的两根导线：固有心率高于静息心率，并随年龄下降，成人约为每分钟 $100$ 次。可见静息时刹车占主导。静息心率 $60$--$70$ 是典型值，未单独引用。 & \cite{JoseCollison1970ihr} & 已测量 \\
2 & 油门 \nt{NORA} 和刹车 \nt{ACET} 是低通滤波器，刹车更快~\eqref{eq:gasbrake} & 狗，对心脏迷走神经或交感神经进行调频刺激，测量到心房频率的传递函数：两条通路都是低通滤波器，交感通路的截止频率低得多，并附加约 $1.7\,\text{s}$ 的纯延迟。特性依赖于平均张力，因此线性公式~\eqref{eq:gasbrake}~只是近似。 & \cite{Berger1989} & 已测量 \\
3 & 刹车之所以快，是因为 \nt{ACET} 不经第二信使直接打开钾通道，而 \nt{NORA} 要经过一连串反应 & 心房细胞膜片：乙酰胆碱开启的钾通道由与受体偶联的 G 蛋白 $\beta\gamma$ 亚基打开，不需要细胞内的第二信使。\nt{NORA} 经 cAMP 的通路此处未引用。 & \cite{Logothetis1987gbg} & 已测量 \\
4 & 血液绕身体一周约需一分钟；血流中的 \nt{ADRE} 到达所有带受体的器官 & 数量级的一般估计；本章即如此表述，未引用来源。 & --- & 估计 \\
5 & 激素级联被负反馈包围：终端激素抑制前面几级 & 实验综述：肾上腺皮质激素在快、中、慢三个时间范围内抑制垂体释放 ACTH 和下丘脑释放释放激素。 & \cite{KellerWood1984fb} & 已测量 \\
6 & 三个相同的级在 $K<8$ 时稳定~\eqref{eq:cascadestable}，临界处周期为 $2\pi\tau/\sqrt3\approx3.6\tau$ & 本章推导：在频率 $\sqrt3/\tau$ 处，三级给出 $180^\circ$ 相移和 $8$ 倍衰减。 & 本章推导 & 理论 \\
7 & 电平误差为 $1/(1+K)$，因此这种调节器的精度不会优于约 $10\,\%$（命题~\ref{prop:cascade}） & 第~6 行对比例反馈的推论。真实的轴在多个级上、在多个时间范围内都有反馈（第~5 行），因此这一界限属于模型~\eqref{eq:cascade}，并非测量所得。 & 本章推导 & 理论 \\
8 & $K=4$ 和 $7$ 时电平趋于稳定，$K=12$ 时出现周期为 $3.7$--$3.9\,\tau$ 的平稳脉动（图~\ref{fig:cascade}） & 用 \texttt{models/\allowbreak chem\_\allowbreak models.py} 计算：$K=12$ 时相继的周期为 $3.9$、$3.7$、$3.8$、$3.7\,\tau$；$K=4$ 时计算结束时的幅度为 $0.003$。 & \texttt{models/\allowbreak chem\_\allowbreak models.py} & 模型 \\
9 & 应激激素大约每小时释放一次脉冲；脉冲由反馈本身产生，无需单独的发生器 & 大鼠，持续输注释放激素：ACTH 和皮质酮以接近每小时一次的生理频率振荡——脉冲并不需要来自下丘脑的节律性输入。带延迟反馈的垂体—肾上腺模型能给出这样的振荡。 & \cite{Walker2012glc, Walker2010} & 假说 \\
10 & 接收方读取的是剂量频率，而不是水平 & 自身不能释放促性腺激素释放激素的猴子：持续输注不能恢复垂体激素的释放，每小时一次的剂量能恢复；改回持续输注后响应再次消失。受体疲劳相当于高通滤波器是一种解释。 & \cite{Belchetz1978} & 已测量 \\
11 & 不同的剂量频率对同一腺体的不同细胞作用不同 & 同一批猴子：剂量频率提高到每小时 $2$--$5$ 次时，两种垂体激素都下降；降低到每 $3$~h 一次时，一种（LH）下降而另一种（FSH）上升，因此两者的比值由频率决定。 & \cite{Wildt1981gnrh} & 已测量 \\
12 & 昼夜节律由细胞内延迟数小时的负反馈产生 & 综述：时钟基因的蛋白质在延迟后抑制自身基因的转录；由转录和翻译构成的环路给出约一天的周期。 & \cite{TakahashiJS2017} & 已测量 \\
13 & 主发生器是几万个神经元组成的一群细胞，它使身体其余部分同步 & 综述：视交叉上核是主昼夜节律发生器；培养中的单个神经元就是独立的振荡器，网络使它们同步；小鼠每侧约有 $10^4$ 个神经元。 & \cite{Welsh2010scn} & 已测量 \\
14 & 人体的固有周期为 $24.18$~h & 在昏暗光线下、按脱离昼夜的作息生活的人：褪黑素、体温和皮质醇节律的周期在年轻人和老年人中平均均为 $24.18$~h，离散很小。 & \cite{Czeisler1999} & 已测量 \\
15 & 光像锁相环一样校准发生器~\eqref{eq:pll}；每天需要移动约 $11$~min & 人体，在一天中不同时间给予一次 $6.7$~h 的强光：得到幅度为 $5$~h 的第一类相位响应曲线——体温最低点之前给光造成延迟，之后给光造成提前。方程~\eqref{eq:pll}~是标准模型；$11$~min $=0.18$~h 是算术。 & \cite{Khalsa2003prc} & 已测量 \\
16 & 没有光就没有校准，节律漂移到固有周期 & 完全失明、没有光感、在正常社会中生活的人：约有一半人的褪黑素和皮质醇节律按自己的周期运行，与昼夜不一致。 & \cite{Sack1992blind} & 已测量 \\
\bottomrule
\end{longtable}}

\section{第~\ref{sec:debugging}~章　调试这台机器}

电极听到了什么、活动的低维性、基于刚性阵列的接口的工作、大脑与解码器的共同学习以及刺激引起的视觉光点，都已在同行评审的实验中测量；数据流的估计是本章的算术；关于植入人体的 Neuralink 设备，就我们所能核实的，数据主要由公司自己发布。

{\footnotesize
\setlength{\tabcolsep}{3pt}\renewcommand{\arraystretch}{1.15}
\begin{longtable}{>{\raggedright}p{0.5cm}>{\raggedright}p{4.0cm}>{\raggedright}p{5.6cm}>{\raggedright}p{2.3cm}>{\raggedright\arraybackslash}p{1.7cm}}
\toprule
编号 & 论断 & 实验及测量内容 & 来源 & 状态 \\
\midrule\endfirsthead
\toprule
编号 & 论断 & 实验及测量内容 & 来源 & 状态 \\
\midrule\endhead
1 & 电极能听到约一百微米半径内神经元的脉冲，通常是一个到几个；可按波形区分它们 & 大鼠 CA1 区，同时进行细胞内和细胞外记录：细胞外脉冲的波形重现了细胞内脉冲的宽度和幅度。估计：四极电极本可区分 $60$--$100$ 个神经元，实际只分出约六个。 & \cite{Henze2000extra} & 已测量 \\
2 & 大脑中有 $8.6\cdot10^{10}$ 个神经元；探针听到的约为一亿分之一 & 对溶解组织悬液中的细胞核计数：人脑中有 $86\cdot10^9$ 个神经元。比例是本章的算术。 & \cite{Azevedo2009} & 已测量 \\
3 & 运动皮层几百个神经元的活动可以放进一二十个公共变量 & 猴子运动皮层记录的综述：群体活动由少数为许多神经元所共有的隐变量描述。 & \cite{Gallego2017} & 已测量 \\
4 & 相邻神经元的噪声是相关的，一百个神经元携带的信息几乎和一千个一样多（命题~\ref{prop:pooling}） & 猴子视觉皮层：相邻神经元噪声的相关系数约为 $0.1$，平均化的收益在一百个神经元时饱和。限制信息的相关性理论。 & \cite{Zohary1994, MorenoBote2014} & 已测量 \\
5 & 原始数据流 $2\cdot10^8$~bit/s，而脉冲中只有 $8\cdot10^4$~bit/s~\eqref{eq:probedata} & 本章依据脉冲流容量~\eqref{eq:spikeentropy}~所作的算术。数字化参数是真实的：Neuralink 芯片为 $19.3$~kHz、每通道 $10$ 位。 & 本章推导；\cite{Rieke1997, Musk2019} & 理论 \\
6 & Neuralink 的第一个系统：芯片放大并数字化，原始数据流经电缆送出，脉冲由外部程序提取；据公司介绍，用于人体的设备在芯片上提取脉冲，采用封闭外壳、无线电以及无线供电 & 公司的论文：完整的原始数据流经 USB-C 电缆送出，脉冲由芯片外的程序实时提取；板上压缩、无线供电和无线传输被列为临床设备的计划。关于植入人体的设备，只有公司的介绍。必须在芯片上提取脉冲是本章由~\eqref{eq:probedata}~得出的结论。 & \cite{Musk2019}；Neuralink 报告，无同行评审论文 & 已测量（公司论文）；人体部分为公司报告 \\
7 & 最多 $3072$ 个电极，分布在 $96$ 根细丝上，每根 $32$ 个；细丝由机器人避开血管插入 & 公司的论文：$3072$ 个电极分布在 $96$ 根细丝上，每根 $32$ 个；机器人每分钟插入 $6$ 根细丝（$192$ 个电极）；长期植入阵列的大鼠中，多达 $70\,\%$ 的电极能听到脉冲。 & \cite{Musk2019} & 已测量（公司论文） \\
8 & 人体中约一千个通道分布在 $64$ 根细丝上；部分细丝移位；光标约 $10$~bit/s & 仅有公司的介绍。在 PubMed 中没有检索到报告人体结果的同行评审论文；这与本章正文中的保留说法一致。 & Neuralink 报告；无同行评审论文 & 已测量（公司报告） \\
9 & 基于一百个电极阵列的接口可以控制光标 & 瘫痪患者，初级运动皮层中 $96$ 个电极：受伤三年后，移动手臂的意图仍会改变脉冲；解码器提供“神经光标”（电子邮件、电视），并能控制假手和机械臂。 & \cite{Hochberg2006} & 已测量 \\
10 & 用“意念之手”写字：每分钟约 $90$ 个字符，低于 $8$~bit/s & 手部瘫痪患者尝试手写，使用循环神经网络解码器：实时每分钟 $90$ 个字符，准确率 $94.1\,\%$，自动校正后超过 $99\,\%$。以比特计的估计是本章的算术。 & \cite{Willett2021} & 已测量 \\
11 & 说话的尝试可以每分钟约 $60$ 个词的速度转换成文字 & 失去清晰言语能力的患者，皮层中植入阵列：每分钟 $62$ 个词；词汇量 $50$ 个词时词错误率 $9.1\,\%$，词汇量 $125\,000$ 个词时为 $23.8\,\%$。 & \cite{Willett2023} & 已测量 \\
12 & 接口只取出容量的一小部分：单个神经元每秒几十比特，一百个神经元每秒几千比特（命题~\ref{prop:probe}） & 上界~\eqref{eq:spikeentropy}~是理论；与第~8、10 行的比较是本章的推论。瓶颈在于低维性和对编码的无知而不在导线，这一点未经直接实验检验。 & \cite{Rieke1997} & 理论 \\
13 & 大脑和解码器相互学习 & 猴子控制光标；解码器在闭环中适应，神经元同时重新学习：准确率提高，技能得以保持，受自己手臂运动的干扰也更小。把学习速度拉开的建议是工程经验法则，本章没有单独的实验。 & \cite{Orsborn2014coad} & 已测量 \\
14 & 通过电极的电流不分类型地激发周围的神经元和导线 & 小鼠和猫的皮层，双光子钙成像：即使是弱电流，也会通过直接激发直径几十微米体积内的轴突，激活分散在远至数毫米处的稀疏神经元。也就是说，激发没有选择性，但也不是连成一片的。 & \cite{Histed2009stim} & 已测量 \\
15 & 刺激视觉皮层会点亮光点；同时点亮最多 $16$ 个光点可以认出一些字母和物体边界 & 盲人，视觉皮层中植入 $96$ 个电极 $6$ 个月：单个电极的阈值为 $66.8\pm36.5$~\textmu A；同时刺激电极组（最多 $16$ 个，这是刺激器的上限）会降低阈值并产生可分辨的图案，能认出一些字母和物体边界。 & \cite{Fernandez2021} & 已测量 \\
16 & Neuralink 的第二个方向是向视觉皮层送入信号的设备 & 只有公司关于研发的介绍；没有找到关于其工作的同行评审数据。 & Neuralink 报告 & 假说 \\
17 & 广播型神经递质的浓度可以用化学传感器在组织中测量 & 大鼠脑片，碳纤维快速循环伏安法：测量了血清素的释放和再摄取；该物质会扩散到突触之外。 & \cite{Bunin1998} & 已测量 \\
\bottomrule
\end{longtable}}

\section{第~\ref{sec:eetypes}~章　作为器件的神经元}

单个细胞的工作模式已通过电极注入电流和记录电压直接测量；积分器和分类的数学是经典理论；而大脑利用模式重构进行计算，仍是假说。

{\footnotesize
\setlength{\tabcolsep}{3pt}\renewcommand{\arraystretch}{1.15}
\begin{longtable}{>{\raggedright}p{0.5cm}>{\raggedright}p{4.0cm}>{\raggedright}p{5.6cm}>{\raggedright}p{2.3cm}>{\raggedright\arraybackslash}p{1.7cm}}
\toprule
编号 & 论断 & 实验及测量内容 & 来源 & 状态 \\
\midrule\endfirsthead
\toprule
编号 & 论断 & 实验及测量内容 & 来源 & 状态 \\
\midrule\endhead
1 & 谐振器：对抗电压变化的慢电流使神经元成为带通滤波器，峰值在几赫兹到几十赫兹（第~\ref{sec:resonator}~节） & 在脑片中向\nt{GLUT}神经元注入频率平滑上升的正弦电流并测量阻抗 $|Z(f)|$：$33^\circ$C时峰值在 $2$--$5$~Hz（$38^\circ$C时约 $7$~Hz）；谐振由慢钾电流和慢阳离子电流产生，并被持续钠电流增强。综述给出了同一手法在许多类细胞中的结果 & \cite{HuVervaekeStorm2002,HutcheonYarom2000} & 已测量 \\
2 & 慢电流的行为像电感，与膜电容一起构成振荡回路 & 测量了轴突对电流的阈下响应，并用线性化的电导方程描述；等效电路包含一个“唯象”电感，其大小依赖于电压 & \cite{Mauro1970} & 理论 \\
3 & 带反馈的群体的时间常数为 $\tau_{\text{eff}}=\tau/(1-w)$~\eqref{eq:perfectint}；$\tau=0.1$~s、$w=0.995$ 时为 $20$~s & 本章从群体的线性方程推出；作为保持眼睛位置的机制，由一篇理论工作提出 & 本章推导；\cite{Seung1996} & 理论 \\
4 & 眼位积分器靠网络而不是单个细胞的特性来保持输入 & 在鱼的眼位保持核团中进行细胞内记录：快速转动后，神经元的频率阶跃变化并保持；向单个细胞注入短暂电流只引起短暂偏移，而电压阶跃在阈下也保持——它们由突触输入维持 & \cite{Aksay2001} & 已测量 \\
5 & 无泄漏积分器需要通过学习不断调节；失调时它要么遗忘，要么饱和 & 用与眼位成 $\pm$ 比例的视觉反馈训练鱼：保持变得不稳定或出现泄漏，神经元的时间常数下降一个数量级以上，降到 $<1$~s；正常的视觉反馈逐渐恢复了稳定 & \cite{Major2004} & 已测量 \\
6 & 双稳态神经元：短暂输入开启持久平台，抑制将其关闭；这一特性由血清素开启（第~\ref{sec:bistable}~节） & 猫控制肌肉的\nt{ACET}神经元的细胞内记录：短暂的突触兴奋使神经元进入持久工作，短暂的抑制使其复位；在脊髓猫中，给予血清素前体后这种状态才出现 & \cite{Hounsgaard1988} & 已测量 \\
7 & 两类：积分器（频率从零开始增长）和谐振器（在阈值处立即出现有限频率） & 两类由分岔理论导出。实验：在皮层脑片中，\nt{GLUT}神经元能以任意小的频率开始工作（第~1型），快速\nt{GABA}神经元则从有限频率跳变开始（第~2型） & \cite{Izhikevich2007,Tateno2004} & 已测量 \\
8 & 同一个神经元可以是积分器或符合检测器：抑制缩短求和窗口 & 在海马脑片中，兴奋性输入加和到阈值的窗口短于 $2$~ms；紧随兴奋而来的前馈抑制使它缩短；在抑制较弱的树突中，窗口较宽 & \cite{Pouille2001} & 已测量 \\
9 & 由轴突构成的延迟线（表~\ref{tab:eetypes}） & 在谷仓猫头鹰中测量了通往符合检测器的轴突延迟：不同的轴突长度给出不同的延迟，检测器被调谐到声音到达的时间差 & \cite{Carr1990} & 已测量 \\
10 & 清醒时是起搏器，睡眠时是沉默的细胞 & \nt{SERO}神经元白天几乎规则地工作，在慢波睡眠中减慢，在快速眼动睡眠中几乎沉默（详见第~\ref{sec:sleep}~章） & \cite{JacobsAzmitia1992,McGintyHarper1976} & 已测量 \\
11 & 器件由离子通道以及调节它们的广播通道决定（命题~\ref{prop:eetypes}） & 已在小型网络中证明：调节物质改变神经元的固有特性和突触强度，使同一电路产生不同的节律 & \cite{Marder2012} & 已测量 \\
12 & 大脑利用模式重构进行计算（命题~\ref{prop:eetypes}） & 没有直接实验表明解决某个任务需要重构细胞的模式；只有调节物改变电路输出的实验 & \cite{Marder2012} & 假说 \\
\bottomrule
\end{longtable}}

\section{第~\ref{sec:dendrites}~章　树突数目}

各类的结构和输入数目已通过解剖和电记录测量；估计~\eqref{eq:dendriteload}~是简单的电容器物理；而输入数目的计算解释依赖于理论和模型。

{\footnotesize
\setlength{\tabcolsep}{3pt}\renewcommand{\arraystretch}{1.15}
\begin{longtable}{>{\raggedright}p{0.5cm}>{\raggedright}p{4.0cm}>{\raggedright}p{5.6cm}>{\raggedright}p{2.3cm}>{\raggedright\arraybackslash}p{1.7cm}}
\toprule
编号 & 论断 & 实验及测量内容 & 来源 & 状态 \\
\midrule\endfirsthead
\toprule
编号 & 论断 & 实验及测量内容 & 来源 & 状态 \\
\midrule\endhead
1 & 膜的比电容 $c_m\approx1$~\textmu F/cm$^2$ & 用吸管从神经元胞体上取下膜片，施加电压阶跃，根据电容电流的衰减求出总电容，再除以面积：三类神经元均为 $0.9$~\textmu F/cm$^2$ & \cite{Gentet2000} & 已测量 \\
2 & 充电时间 $t\approx c_mA\Delta V/I$~\eqref{eq:dendriteload}：大神经元需要几十个同时输入，小神经元一个大突触就够 & 按电容器充电定律所作的估计；数字（$20$ 和 $300$~pF，$0.1$ 和几nA）是数量级 & 本章推导 & 理论 \\
3 & 一个巨大突触给出几纳安的电流，立即把小神经元推到阈值 & 在脑片中同时记录末梢和接收者：单个突触的电流为 $2$--$13$~nA，每个输入脉冲都引起阈上响应，$36^\circ$C时延迟约 $0.4$~ms；接收者输出\nt{GLYC} & \cite{Borst1995,BorstSoria2012} & 已测量 \\
4 & 零个树突：触觉和痛觉感觉神经元的脉冲从传感器绕过细胞体进入脊髓 & 结构（在胞体处一分为二的轴突）已由解剖确定。传导不需要胞体的兴奋性，这一点在这类神经元的经过验证的模型中得到证明：没有可兴奋的胞体，脉冲同样可靠地通过分叉 & \cite{AmirDevor2003} & 理论 \\
5 & 一个树突：嗅觉神经元携带一种受体，传递一条输入线路 & 受体基因实验的综述：每个嗅觉神经元只表达大家族中的一个受体基因 & \cite{Mombaerts2004} & 已测量 \\
6 & 两个树突：声音方向检测器中，来自左右耳的输入落在不同的树突上 & 综述汇集的哺乳动物解剖和记录：来自两耳的输入分布在两个相反方向的树突上 & \cite{Grothe2010} & 已测量 \\
7 & 把输入分到两个树突上使“与”更严格 & 在模型中证明：一个树突上的饱和~\eqref{eq:shunt}~使来自一侧的输入无法到达阈值，符合检测得到改善。没有改变输入排布的直接实验 & \cite{AgmonSnir1998} & 理论 \\
8 & 运动学习回路中约 $7\cdot10^{10}$ 个小\nt{GLUT}神经元，每个 $\sim4$ 个输入 & 用各向同性分离法计数细胞：人脑这一部分约有 $69\cdot10^9$ 个神经元。在活体大鼠中记录一个这样的细胞：触碰时来自少数输入的离散电流成簇到来，它们相加时细胞发放 & \cite{Azevedo2009,Chadderton2004} & 已测量 \\
9 & 有 $n$ 个输入的阈值元件最多分开 $P_{\max}\approx2n$ 个随机模式~\eqref{eq:cover}；对运动学习回路的输出神经元，上界 $\sim2\cdot10^5$，现实估计 $\sim4\cdot10^4$ 个对应关系 & 上界是组合几何的经典结果。现实估计是只有正权重并留有抗噪余量的容量理论，应用于这个神经元实测的输入数 & \cite{Cover1965,Brunel2004} & 理论 \\
10 & 输入少的混合器用于扩展输入；“四个”接近最优 & 理论：混合器层给出的表示维数，大约在解剖学上观察到的输入数目处达到最大；扩展的最初假说见经典工作 & \cite{Marr1969,Albus1971,LitwinKumar2017} & 假说 \\
11 & 大树突树：$10^4$--$10^5$ 个输入 & 按电子显微照片计数突触：大鼠运动学习回路的一个输出\nt{GABA}神经元上有 $\approx1.7\cdot10^5$ 个突触；海马\nt{GLUT}神经元上约有 $30\,000$ 个兴奋性输入和 $1\,700$ 个抑制性输入 & \cite{NapperHarvey1988,Megias2001} & 已测量 \\
12 & 大树突树的分支是各自带阈值的独立子电路；神经元是一个小型多层网络 & 在细树突的两个位置点刺激：同一分支上的输入按S形曲线相加，不同分支上的输入线性相加。在人类皮层神经元的树突中发现了分级的钙脉冲，使单个细胞能解决线性不可分的任务。模型：要重现一个神经元需要一个深层网络 & \cite{Polsky2004,Gidon2020,Beniaguev2021} & 已测量 \\
13 & 树突即输出：视网膜\nt{GABA}神经元的每个分支都是独立的方向检测器 & 在单个分支中进行双光子钙成像：分支中的信号对从细胞体向末梢的运动有选择性，而胞体电压没有 & \cite{Euler2002} & 已测量 \\
14 & 输入数目跟随任务（命题~\ref{prop:dendrites}） & 各类的结构已测量（第~3--13行）；输入数目是为了速度或分类而选定的，这一点只有针对个别类别的理论和模型给出 & \cite{LitwinKumar2017,AgmonSnir1998} & 假说 \\
\bottomrule
\end{longtable}}

\section{第~\ref{sec:sleep}~章　睡眠}

睡眠的模式、重放和突触的缩小已通过神经元记录、微透析和电子显微镜测量；睡眠的时间安排和慢摆动由本书的模型描述，而睡眠的目的主要还是假说。

{\footnotesize
\setlength{\tabcolsep}{3pt}\renewcommand{\arraystretch}{1.15}
\begin{longtable}{>{\raggedright}p{0.5cm}>{\raggedright}p{4.0cm}>{\raggedright}p{5.6cm}>{\raggedright}p{2.3cm}>{\raggedright\arraybackslash}p{1.7cm}}
\toprule
编号 & 论断 & 实验及测量内容 & 来源 & 状态 \\
\midrule\endfirsthead
\toprule
编号 & 论断 & 实验及测量内容 & 来源 & 状态 \\
\midrule\endhead
1 & 睡眠中皮层神经元并不沉默；改变的是工作模式 & 大鼠皮层神经元的长时间记录：慢波睡眠中频率不为零，工作期与静默期交替；长时间清醒后，各状态下的频率都更高，睡眠后则更低 & \cite{Vyazovskiy2009} & 已测量 \\
2 & \nt{ACET} 在白天和快速眼动睡眠中高，在慢波睡眠中低（表~\ref{tab:sleepmodes}） & 自由活动的猫，皮层微透析：清醒时乙酰胆碱的释放比慢波睡眠高 $100$--$175\%$，快速眼动睡眠中约与安静清醒时持平 & \cite{Marrosu1995} & 已测量 \\
3 & \nt{NORA} 和 \nt{SERO} 在慢波睡眠中安静下来，在快速眼动睡眠中几乎沉默 & 按睡眠阶段记录大鼠的单个 \nt{NORA} 神经元和猫的 \nt{SERO} 神经元：频率在清醒时最高，慢波睡眠中较低，快速眼动睡眠中几乎为零 & \cite{AstonJonesBloom1981,McGintyHarper1976} & 已测量 \\
4 & 快速眼动睡眠中肌肉经 \nt{GABA} 和 \nt{GLYC} 被抑制关闭 & 在大鼠中测量咀嚼肌 \nt{ACET} 神经元的活动，并把阻断剂直接施加到它们上：只有两类 \nt{GABA} 受体和 \nt{GLYC} 受体都被阻断时，麻痹才解除 & \cite{BrooksPeever2012} & 已测量 \\
5 & 睡眠压力 $S$ 是带两个阈值的电容器~\eqref{eq:sleeppressure}，$\tau_r=18.2$~h，$\tau_d=4.2$~h & 双过程模型；时间常数由脑电图慢波功率随清醒增长、随睡眠下降的方式得出，阈值的形状由自发醒来的时间得出 & \cite{Borbely1982,Daan1984} & 理论 \\
6 & 机器自己进入 $16.1$~h 清醒和 $8.0$~h 睡眠（图~\ref{fig:twoprocess}） & 按~\eqref{eq:sleeppressure}~加上摆动阈值计算，脚本 \texttt{models/sleep\_models.py}：清醒 $16.1$~h，睡眠 $7.9$--$8.0$~h & --- & 模型 \\
7 & $40$~h 不睡后 $S=0.90$，但睡眠只持续 $8.8$~h：欠债靠深度而不是长度偿还 & 模型：\texttt{models/sleep\_models.py} 给出 $S=0.90$ 和 $8.8$~h。实验：八个人 $40.5$~h 不睡后，第一个恢复夜中脑电图慢波功率显著高于平时 & \cite{Borbely1981} & 模型 \\
8 & 物理上 $S$ 就是腺苷 & 猫的微透析：在一个控制清醒的区域中，细胞外腺苷在清醒期间上升，睡眠剥夺时继续上升，睡眠中下降。$S$ 仅仅是腺苷这一点尚未证明 & \cite{PorkkaHeiskanen1997,Dunwiddie2001} & 假说 \\
9 & 慢波睡眠中皮层摆动：工作几百 ms——静默，每秒不到一次（$<1$~Hz，麻醉下常为 $0.3$--$0.4$~Hz） & 麻醉猫的细胞内记录：$0.8$--$1.5$~s 的去极化与长时间的超极化交替，节律 $<1$~Hz，最常见为 $0.3$--$0.4$~Hz；自然睡眠中也能看到同样的工作与静默交替（第~1 行） & \cite{Steriade1993,Vyazovskiy2009} & 已测量 \\
10 & 摆动来自带反馈和疲劳的群体~\eqref{eq:updown}；$b=5$ 时周期 $1.1$~s（模型值），约 $35\%$ 的时间工作；$b=1$ 时平稳工作（图~\ref{fig:updown}） & 计算，脚本 \texttt{models/sleep\_models.py}：周期 $1.11$~s，工作时间比例 $0.35$（按整数个周期平均）；$b=1$ 时工作比例为 $1.0$。模型的周期短于实测值（第~9 行） & --- & 模型 \\
11 & \nt{ACET} 关闭摆动，使皮层转入平稳工作 & 刺激猫的 \nt{ACET} 神经元核团，阻断了 $79\%$ 皮层神经元的慢节律，使其转入平稳工作，主要是长时间超极化消失所致。乙酰胆碱抑制产生疲劳的慢钾电流 & \cite{SteriadeAmzica1993,McCormick1992} & 已测量 \\
12 & $7$--$15$~Hz 的节律爆发由皮层和中继核团之间的 \nt{GLUT}--\nt{GABA} 环路产生 & 在雪貂视觉中继核团的脑片中，爆发由中继细胞对来自相邻 \nt{GABA} 核团的抑制所作的反跳簇放电产生，而该核团又由中继细胞自己兴奋 & \cite{vonKrosigk1993,SteriadeMcCormickSejnowski1993} & 已测量 \\
13 & 白天的重放是快进的：海马中约快 $20$ 倍，皮层中快 $6$--$7$ 倍 & 慢波睡眠中，海马神经元的重复序列在时间上被压缩。跑道奔跑后，位置神经元的序列以同样的顺序在 $\sim100$~ms 的短爆发中重复，压缩约 $20$ 倍；皮层中压缩 $6$--$7$ 倍 & \cite{Nadasdy1999,LeeWilson2002,Euston2007} & 已测量 \\
14 & 增强重放与皮层的协调可改善记忆（因果关系） & 大鼠在训练后接受与重放锁定的刺激，以加强重放与皮层慢波和节律爆发的联系：第二天记忆良好，对照组处于随机水平 & \cite{Maingret2016} & 已测量 \\
15 & 重放的目的是对慢速记忆进行交错训练 & 互补学习系统理论和人工网络上的计算：顺序学习会破坏旧知识，交错学习则不会。没有直接的脑实验表明重放恰恰为此所需 & \cite{McClelland1995} & 假说 \\
16 & 睡眠后突触按大小成比例缩小，大的几乎不变 & 对小鼠皮层 $6920$ 个突触的三维电子显微镜观察：睡眠后接触面积比清醒后小 $\sim18\%$，缩小与大小成比例，大而稳定的突触不变 & \cite{deVivo2017} & 已测量 \\
17 & 睡眠的主要任务是权重的重新归一化 & 突触稳态假说；第~1 和第~16 行与之一致，但不能证明这是主要功能 & \cite{TononiCirelli2014} & 假说 \\
18 & 快速眼动睡眠是对世界模型的台架测试 & 模式（表~\ref{tab:sleepmodes}）已测量；网络在运行世界模型是理论推测，没有实验 & \cite{Hobson2009} & 假说 \\
19 & 睡眠中大脑的冲洗：问题悬而未决 & 一项实验：睡眠中细胞间隙宽 $60\%$，清除更快。另一项采用不同测量方法的实验：睡眠中和麻醉下清除明显更慢。结果相互矛盾 & \cite{Xie2013,Miao2024} & 假说 \\
20 & 局部睡眠：清醒动物的皮层区域会短暂陷入静默 & 大鼠长时间清醒后，个别皮层区域的神经元短暂沉默，就像在慢波睡眠中一样，局部脑电图中出现慢波；此时大鼠在任务中更容易出错 & \cite{Vyazovskiy2011} & 已测量 \\
\bottomrule
\end{longtable}}

\section{第~\ref{sec:livingmachine}~章　用活神经元搭成的机器}

电极阵列上培养物的行为是在带记录和刺激的直接实验中测得的；群放电机制依据突触耗竭模型和微培养物实验，而关于培养皿中多巴胺的论断依据的是个别实验，其中一部分连作者本人也只视为演示。

{\footnotesize
\setlength{\tabcolsep}{3pt}\renewcommand{\arraystretch}{1.15}
\begin{longtable}{>{\raggedright}p{0.5cm}>{\raggedright}p{4.0cm}>{\raggedright}p{5.6cm}>{\raggedright}p{2.3cm}>{\raggedright\arraybackslash}p{1.7cm}}
\toprule
序号 & 论断 & 实验及测得的结果 & 来源 & 状态 \\
\midrule\endfirsthead
\toprule
序号 & 论断 & 实验及测得的结果 & 来源 & 状态 \\
\midrule\endhead
1 & 芯片上的培养物：主要是\nt{GLUT}，较小比例的\nt{GABA}，比例取决于标本；海马培养物中约 $6\%$ & 电极阵列上由几千个神经元组成的网络是标准实验（第~3~行）。海马培养物中\nt{GABA}神经元的比例用GABA染色计数：约占神经元的 $6\%$。对皮层培养物，我们不给出经过核实的数字 & \cite{Benson1994} & 已测量 \\
2 & 类器官：由人类干细胞培养、大小约一毫米的三维神经组织 & 用干细胞培养出含有多个脑区的三维类器官，其中包括带有前体细胞层和成熟神经元的皮层 & \cite{Lancaster2013} & 已测量 \\
3 & 无输入时培养物以群放电形式爆发，间隔从一秒到几分钟，因培养物而异 & 对密度为 $3\,000$ 到 $50\,000$ 个神经元的 $58$ 个皮层培养物记录了五周（$963$ 段半小时记录）：两周后几乎所有培养物都以全网络群放电为主；群放电间隔为 $1$ 到 $300$~s，并强烈依赖于标本 & \cite{Wagenaar2006} & 已测量 \\
4 & 群放电因递质耗竭而结束，间隔 $T_{\text{burst}}\approx\tau_{\text{rec}}\ln\frac{1}{1-x^*}$~\eqref{eq:burstinterval}；$2$--$4$~s是 $\tau_{\text{rec}}=0.8$~s时的模型估计，实测的恢复更慢 & 公式在本章推导。模型：带耗竭性突触的网络会自发产生全网络群放电。在 $4$--$30$ 个神经元的微培养物上的实验：群放电时长取决于距上次群放电的时间，递质囊泡耗竭会取消群放电；作者的结论是群放电受递质储备限制。那里储备恢复需要几十秒，代入同一公式得到几十秒乃至几分钟的间隔 & 本章推导；\cite{Tsodyks2000,CohenSegal2011,Tsodyks1997} & 理论 \\
5 & “会闭嘴的老师”：出现所需响应时关掉刺激，响应被固定下来 & 以 $0.3$--$1$~Hz刺激培养物，直到刺激后 $50\pm10$~ms出现选定的响应，然后关掉刺激 $5$~min；几个循环后所需响应立刻出现；绘制了学习曲线 & \cite{Shahaf2001} & 已测量 \\
6 & 培养物打网球；一次会话内连续击球回合的显著增长只出现在完整反馈下 & 详见第~\ref{sec:dishbrain}~章及其补充：以安静代替刺激时，培养物到会话结束时也比无反馈时好，但效应较小；跨会话学习不稳定 & \cite{Kagan2022} & 已测量 \\
7 & 培养物学会预测自己的输入是假说；关于其“有感知能力”的高调说法遭到反驳 & 自由能原理是一个理论设想；没有直接实验能把它与更简单的解释区分开。三十位研究者在回应文章中指出，闭环中的行为变化既不表明智能，也不表明感知 & \cite{Friston2010,Balci2023} & 假说 \\
8 & 在挑选好的刺激模式下，培养物驾驶虚拟身体走向目标 & 程序根据当前结果自行挑选训练刺激模式；几十分钟内网络达到指定状态，训练 $2$~h后可塑性保持 $80$~min。同样的刺激若与行为无关地施加，则不起作用 & \cite{Bakkum2008} & 已测量 \\
9 & 储备池：只训练线性读出 $y=W_{\text{out}}\mathbf r$~\eqref{eq:reservoir} & 储备池计算理论：任意带衰减记忆的非线性系统加上一个训练好的线性输出 & \cite{Maass2002,JaegerHaas2004} & 理论 \\
10 & 作为储备池的类器官辨别说话人的准确率约 $78\%$（作者数据）；按误差学习的是读出，类器官的连接在无监督地变化 & 把几位说话人的语音以电极阵列上的电流模式输入类器官；训练好的读出辨别说话人的准确率约 $78\%$，这是作者报告的。作者还报告训练期间类器官的功能连接发生了变化，称之为类器官内部的无监督学习 & \cite{Cai2023} & 已测量 \\
11 & 一百万个神经元用于脉冲的功耗约 $10^{-4}$~W~\eqref{eq:dishpower} & 估计：由皮层能量预算~\eqref{eq:energy} 得出的每个脉冲 $10^{-10}$~J，乘以脉冲数；培养物功率没有直接测量 & 本章推导；\cite{Attwell2001} & 理论 \\
12 & 闪光释放多巴胺：$1$~mM笼锁多巴胺，$240$ 次重复，诱发脉冲数增加 & 在远程类器官平台上，当刺激引发更多脉冲时，用紫外光（$365$~nm，$0.8$~s）释放光敏笼中的多巴胺（$1$~mM）；$240$ 步（约 $5$~h）中脉冲数总体增加。作者写道，仅凭一次实验无法确立与多巴胺的关联；没有对照 & \cite{Jordan2024} & 假说 \\
13 & 活细胞分泌的多巴胺能长期、可逆地改变有机晶体管的权重 & 分泌多巴胺的细胞培养在有机神经形态元件上方；多巴胺在电极处氧化改变沟道电导；展示了权重的长期改变及其复原 & \cite{Keene2020} & 已测量 \\
14 & 存在含有工作\nt{DOPA}神经元的类器官；其多巴胺未被用作老师 & 在人类干细胞类器官中发现了电活动正常的成熟\nt{DOPA}神经元以及多巴胺的产生。文献中我们没有找到用这种多巴胺去教另一个网络的实验 & \cite{Jo2016} & 已测量 \\
15 & 类器官平衡杆子：程序给出的教学刺激优于随机刺激，不能巩固，经由\nt{GLUT}突触 & 小鼠皮层类器官与“小车倒立摆”任务构成闭环：对大多数类器官，强化学习选择的信号比随机信号或无信号效果更好；休息 $45$~min后改进不保留；阻断AMPA和NMDA受体则改进消失 & \cite{Robbins2026} & 已测量 \\
16 & 没有控制寄存器，试验台上的网络无法自己决定什么算成功（命题~\ref{prop:livingmachine}） & 在第~5--15~行的所有实验中，教学信号都由实验者或程序设定；没有培养物自己形成成功标准的实验——这是从“不存在”得出的推论，不是实验 & --- & 假说 \\
\bottomrule
\end{longtable}}

\section{第~\ref{sec:dishbrain}~章　DishBrain}

试验台的结构和游戏结果取自一项实验工作及其后续研究；噪声公式和向好设置漂移的公式在本章推导，并用本书的模型做了检验，而培养皿中的学习机制尚未确定。

{\footnotesize
\setlength{\tabcolsep}{3pt}\renewcommand{\arraystretch}{1.15}
\begin{longtable}{>{\raggedright}p{0.5cm}>{\raggedright}p{4.0cm}>{\raggedright}p{5.6cm}>{\raggedright}p{2.3cm}>{\raggedright\arraybackslash}p{1.7cm}}
\toprule
序号 & 论断 & 实验及测得的结果 & 来源 & 状态 \\
\midrule\endfirsthead
\toprule
序号 & 论断 & 实验及测得的结果 & 来源 & 状态 \\
\midrule\endhead
1 & 试验台：每块芯片约 $800\,000$ 个小鼠皮层细胞或约 $10^6$ 个人类细胞；$8$~mm$^2$ 上 $26\,000$ 个电极，最多记录 $1024$ 个，经 $8$ 个电极刺激 & 该工作的方法：MaxOne芯片，$8$~mm$^2$ 上 $26\,000$ 个铂电极，记录 $1024$ 个通道，每秒 $20\,000$ 个采样；技术上最多可刺激 $32$ 个电极，能独立刺激的是 $8$ 个；小鼠培养物约从第 $10$ 天起使用 & \cite{Kagan2022} & 已测量 \\
2 & 节拍 $10$~ms的回路：运动区的脉冲移动球拍（图~\ref{fig:dishloop}） & 脉冲按 $10$~ms窗口计数；从脉冲到响应刺激的延迟约 $5$~ms，由硬件缓冲决定 & \cite{Kagan2022} & 已测量 \\
3 & 输入：位置编码（$8$ 个电极中的哪一个）和 $4$--$40$~Hz的频率编码 & 主实验中，刺激频率从球在远端墙边时的 $4$~Hz线性增加到球拍处的 $40$~Hz；球的刺激脉冲幅度为 $75$~mV；脉冲为双相矩形脉冲 & \cite{Kagan2022} & 已测量 \\
4 & 输出 $\Delta p=c\,(n_\uparrow-n_\downarrow)$~\eqref{eq:dishreadout}，窗口计数噪声 $1/\sqrt{RT}$ & 该工作描述了两区差值规则（各区按活动加权），但没有给出确切公式。噪声估计是泊松计数的推论，在本章推导 & \cite{Kagan2022}；本章推导 & 理论 \\
5 & 老师：击中后给可预测刺激 $75$~mV、$100$~Hz、$0.1$~s；未击中后给不可预测刺激 $150$~mV、$5$~Hz、$4$~s，随机位置和时刻，随后 $4$~s停顿 & 方法：击中后 $75$~mV、$100$~Hz、$100$~ms，同时施加到全部 $8$ 个电极；未击中后 $150$~mV、$5$~Hz，在 $4$~s内于随机时刻施加到随机电极，随后 $4$~s无刺激。培养物中没有多巴胺 & \cite{Kagan2022} & 已测量 \\
6 & 网络的行为是为了让自己的输入变得可预测 & 自由能原理是一个理论设想，作者据此设计了方案；实验与它相符，但不能把它与更简单的机制区分开（第~7--8~行） & \cite{Friston2010,Kagan2022} & 假说 \\
7 & 若失败晃动网络而成功不晃动，停留在某设置的时间 $\rho\propto1/P_{\text{miss}}$~\eqref{eq:dishdensity}（命题~\ref{prop:dishscramble}） & 由对称扰动马尔可夫链中的流量平衡得出，在本章推导。没有在培养物上直接检验 & 本章推导 & 理论 \\
8 & “未击中即打乱”模型会学习：击中比例 $0.21\to0.38$；每个回合后都扰动时 $\approx0.12$，无扰动时 $0.19$（图~\ref{fig:dishmodel}） & 计算，脚本 \texttt{models/dishbrain\_model.py}，$40$ 个模型培养物，各 $2000$ 个回合：$0.21\to0.38$；$0.13\to0.11$；$0.19\to0.19$ & --- & 模型 \\
9 & 只有处在完整回路中的活培养物，连续击球回合才会变长；对照组不学习 & $399$ 次 $20$~min的会话：小鼠培养物 $9$ 个（$101$ 次会话），人类 $11$ 个（$138$），只有培养基的芯片 $6$ 块（$80$），静息培养物 $20$ 个（$42$），仿真 $3$ 个（$38$）。只有活培养物在后 $15$~min的平均回合长度大于前 $5$~min；非神经元细胞（HEK293T）和只有培养基的芯片没有效应 & \cite{Kagan2022} & 已测量 \\
10 & 一次会话内的显著增长只出现在完整反馈下；以安静代替刺激时，会话结束时的表现好于无反馈，但效应较小 & $3$ 天内 $486$ 次会话：“刺激”（$n=27$）、“安静”（$17$）、“无反馈”（$15$）。随时间的增长只在“刺激”条件下显著；会话结束时，“安静”条件优于“无反馈”，但效应较小 & \cite{Kagan2022} & 已测量 \\
11 & 效应不大；网络是否能长期学习是一个开放问题 & 一次会话内的改进是可靠的，但据作者本人所说，几天里跨会话的学习并未稳定出现 & \cite{Kagan2022} & 已测量 \\
12 & 游戏时网络趋近临界状态，越接近，游戏打得越好 & 对同一批培养物的活动雪崩分析：结构化输入使网络趋近临界状态，游戏水平与接近程度相关；但单有临界性、而没有关于行动后果的信息，不足以产生学习 & \cite{Habibollahi2023} & 已测量 \\
13 & 培养物的“感知能力”未被证明 & 三十位研究者的回应文章：闭环中的行为变化既不证明智能，也不证明感受；没有能表明这一点的实验 & \cite{Balci2023} & 假说 \\
14 & 建议的第四个对照：未击中后施加时长和能量相同的可预测刺激（第~\ref{sec:dishspec}~节） & 这样的实验尚未做过；这是本书的建议 & --- & 假说 \\
\bottomrule
\end{longtable}}


\begin{thebibliography}{99}
\bibitem{Azevedo2009} F. A. C. Azevedo et al., \emph{Equal numbers of neuronal and nonneuronal cells make the human brain an isometrically scaled-up primate brain}, J. Comp. Neurol. \textbf{513}, 532 (2009).
\bibitem{Schultz1997} W. Schultz, P. Dayan, and P. R. Montague, \emph{A neural substrate of prediction and reward}, Science \textbf{275}, 1593 (1997).
\bibitem{SuttonBarto2018} R. S. Sutton and A. G. Barto, \emph{Reinforcement Learning: An Introduction}, 2nd ed. (MIT Press, Cambridge, MA, 2018).
\bibitem{Doya2002} K. Doya, \emph{Metalearning and neuromodulation}, Neural Networks \textbf{15}, 495 (2002).
\bibitem{Strata1999} P. Strata and R. Harvey, \emph{Dale's principle}, Brain Res. Bull. \textbf{50}, 349 (1999).
\bibitem{vanVreeswijk1996} C. van Vreeswijk and H. Sompolinsky, \emph{Chaos in neuronal networks with balanced excitatory and inhibitory activity}, Science \textbf{274}, 1724 (1996).
\bibitem{AstonJones2005} G. Aston-Jones and J. D. Cohen, \emph{An integrative theory of locus coeruleus-norepinephrine function: adaptive gain and optimal performance}, Annu. Rev. Neurosci. \textbf{28}, 403 (2005).
\bibitem{Beaulieu2011} J.-M. Beaulieu and R. R. Gainetdinov, \emph{The physiology, signaling, and pharmacology of dopamine receptors}, Pharmacol. Rev. \textbf{63}, 182 (2011).
\bibitem{Sparso2001} J. Spars{\o} and S. Furber (eds.), \emph{Principles of Asynchronous Circuit Design: A Systems Perspective} (Kluwer, Boston, 2001).
\bibitem{Carr1990} C. E. Carr and M. Konishi, \emph{A circuit for detection of interaural time differences in the brain stem of the barn owl}, J. Neurosci. \textbf{10}, 3227 (1990).
\bibitem{Wang2001} X.-J. Wang, \emph{Synaptic reverberation underlying mnemonic persistent activity}, Trends Neurosci. \textbf{24}, 455 (2001).
\bibitem{Hahnloser2002} R. H. R. Hahnloser, A. A. Kozhevnikov, and M. S. Fee, \emph{An ultra-sparse code underlies the generation of neural sequences in a songbird}, Nature \textbf{419}, 65 (2002).
\bibitem{Barlow1965} H. B. Barlow and W. R. Levick, \emph{The mechanism of directionally selective units in rabbit's retina}, J. Physiol. \textbf{178}, 477 (1965).
\bibitem{Carandini2012} M. Carandini and D. J. Heeger, \emph{Normalization as a canonical neural computation}, Nat. Rev. Neurosci. \textbf{13}, 51 (2012).
\bibitem{Pouille2001} F. Pouille and M. Scanziani, \emph{Enforcement of temporal fidelity in pyramidal cells by somatic feed-forward inhibition}, Science \textbf{293}, 1159 (2001).
\bibitem{Tsodyks1997isn} M. V. Tsodyks, W. E. Skaggs, T. J. Sejnowski, and B. L. McNaughton, \emph{Paradoxical effects of external modulation of inhibitory interneurons}, J. Neurosci. \textbf{17}, 4382 (1997).
\bibitem{Buzsaki2012} G. Buzs\'aki and X.-J. Wang, \emph{Mechanisms of gamma oscillations}, Annu. Rev. Neurosci. \textbf{35}, 203 (2012).
\bibitem{Rieke1997} F. Rieke, D. Warland, R. de Ruyter van Steveninck, and W. Bialek, \emph{Spikes: Exploring the Neural Code} (MIT Press, Cambridge, MA, 1997).
\bibitem{Mainen1995} Z. F. Mainen and T. J. Sejnowski, \emph{Reliability of spike timing in neocortical neurons}, Science \textbf{268}, 1503 (1995).
\bibitem{Softky1993} W. R. Softky and C. Koch, \emph{The highly irregular firing of cortical cells is inconsistent with temporal integration of random EPSPs}, J. Neurosci. \textbf{13}, 334 (1993).
\bibitem{Thorpe1996} S. Thorpe, D. Fize, and C. Marlot, \emph{Speed of processing in the human visual system}, Nature \textbf{381}, 520 (1996).
\bibitem{Zohary1994} E. Zohary, M. N. Shadlen, and W. T. Newsome, \emph{Correlated neuronal discharge rate and its implications for psychophysical performance}, Nature \textbf{370}, 140 (1994).
\bibitem{MorenoBote2014} R. Moreno-Bote et al., \emph{Information-limiting correlations}, Nat. Neurosci. \textbf{17}, 1410 (2014).
\bibitem{Gollisch2008} T. Gollisch and M. Meister, \emph{Rapid neural coding in the retina with relative spike latencies}, Science \textbf{319}, 1108 (2008).
\bibitem{OKeefe1993} J. O'Keefe and M. L. Recce, \emph{Phase relationship between hippocampal place units and the EEG theta rhythm}, Hippocampus \textbf{3}, 317 (1993).
\bibitem{Destexhe2003} A. Destexhe, M. Rudolph, and D. Par\'e, \emph{The high-conductance state of neocortical neurons in vivo}, Nat. Rev. Neurosci. \textbf{4}, 739 (2003).
\bibitem{Tsodyks1997} M. V. Tsodyks and H. Markram, \emph{The neural code between neocortical pyramidal neurons depends on neurotransmitter release probability}, Proc. Natl. Acad. Sci. USA \textbf{94}, 719 (1997).
\bibitem{Lisman1997} J. E. Lisman, \emph{Bursts as a unit of neural information: making unreliable synapses reliable}, Trends Neurosci. \textbf{20}, 38 (1997).
\bibitem{Attwell2001} D. Attwell and S. B. Laughlin, \emph{An energy budget for signaling in the grey matter of the brain}, J. Cereb. Blood Flow Metab. \textbf{21}, 1133 (2001).
\bibitem{Lennie2003} P. Lennie, \emph{The cost of cortical computation}, Curr. Biol. \textbf{13}, 493 (2003).
\bibitem{Yagishita2014} S. Yagishita et al., \emph{A critical time window for dopamine actions on the structural plasticity of dendritic spines}, Science \textbf{345}, 1616 (2014).
\bibitem{Werfel2005} J. Werfel, X. Xie, and H. S. Seung, \emph{Learning curves for stochastic gradient descent in linear feedforward networks}, Neural Comput. \textbf{17}, 2699 (2005).
\bibitem{Doya2000} K. Doya, \emph{Reinforcement learning in continuous time and space}, Neural Comput. \textbf{12}, 219 (2000).
\bibitem{Oja1982} E. Oja, \emph{Simplified neuron model as a principal component analyzer}, J. Math. Biol. \textbf{15}, 267 (1982).
\bibitem{BiPoo1998} G.-Q. Bi and M.-M. Poo, \emph{Synaptic modifications in cultured hippocampal neurons: dependence on spike timing, synaptic strength, and postsynaptic cell type}, J. Neurosci. \textbf{18}, 10464 (1998).
\bibitem{Williams1992} R. J. Williams, \emph{Simple statistical gradient-following algorithms for connectionist reinforcement learning}, Mach. Learn. \textbf{8}, 229 (1992).
\bibitem{Bayer2005} H. M. Bayer and P. W. Glimcher, \emph{Midbrain dopamine neurons encode a quantitative reward prediction error signal}, Neuron \textbf{47}, 129 (2005).
\bibitem{Shen2008} W. Shen, M. Flajolet, P. Greengard, and D. J. Surmeier, \emph{Dichotomous dopaminergic control of striatal synaptic plasticity}, Science \textbf{321}, 848 (2008).
\bibitem{Dabney2020} W. Dabney, Z. Kurth-Nelson, N. Uchida, C. K. Starkweather, D. Hassabis, R. Munos, and M. Botvinick, \emph{A distributional code for value in dopamine-based reinforcement learning}, Nature \textbf{577}, 671 (2020).
\bibitem{Matsumoto2009} M. Matsumoto and O. Hikosaka, \emph{Two types of dopamine neuron distinctly convey positive and negative motivational signals}, Nature \textbf{459}, 837 (2009).
\bibitem{BrombergMartin2010} E. S. Bromberg-Martin, M. Matsumoto, and O. Hikosaka, \emph{Dopamine in motivational control: rewarding, aversive, and alerting}, Neuron \textbf{68}, 815 (2010).
\bibitem{Menegas2018} W. Menegas, K. Akiti, R. Amo, N. Uchida, and M. Watabe-Uchida, \emph{Dopamine neurons projecting to the posterior striatum reinforce avoidance of threatening stimuli}, Nat. Neurosci. \textbf{21}, 1421 (2018).
\bibitem{Niv2007} Y. Niv, N. D. Daw, D. Joel, and P. Dayan, \emph{Tonic dopamine: opportunity costs and the control of response vigor}, Psychopharmacology \textbf{191}, 507 (2007).
\bibitem{Mohebi2019} A. Mohebi et al., \emph{Dissociable dopamine dynamics for learning and motivation}, Nature \textbf{570}, 65 (2019).
\bibitem{Hamid2016} A. A. Hamid et al., \emph{Mesolimbic dopamine signals the value of work}, Nat. Neurosci. \textbf{19}, 117 (2016).
\bibitem{Kim2020} H. R. Kim, A. N. Malik et al., \emph{A unified framework for dopamine signals across timescales}, Cell \textbf{183}, 1600 (2020).
\bibitem{Jeong2022} H. Jeong et al., \emph{Mesolimbic dopamine release conveys causal associations}, Science \textbf{378}, eabq6740 (2022).
\bibitem{Amo2022} R. Amo et al., \emph{A gradual temporal shift of dopamine responses mirrors the progression of temporal difference error in machine learning}, Nat. Neurosci. \textbf{25}, 1082 (2022).
\bibitem{Takahashi2017} Y. K. Takahashi et al., \emph{Dopamine neurons respond to errors in the prediction of sensory features of expected rewards}, Neuron \textbf{95}, 1395 (2017).
\bibitem{Sharpe2017} M. J. Sharpe et al., \emph{Dopamine transients are sufficient and necessary for acquisition of model-based associations}, Nat. Neurosci. \textbf{20}, 735 (2017).
\bibitem{Gardner2018} M. P. H. Gardner, G. Schoenbaum, and S. J. Gershman, \emph{Rethinking dopamine as generalized prediction error}, Proc. R. Soc. B \textbf{285}, 20181645 (2018).
\bibitem{Engelhard2019} B. Engelhard et al., \emph{Specialized coding of sensory, motor and cognitive variables in VTA dopamine neurons}, Nature \textbf{570}, 509 (2019).
\bibitem{Clements1992} J. D. Clements, R. A. J. Lester, G. Tong, C. E. Jahr, and G. L. Westbrook, \emph{The time course of glutamate in the synaptic cleft}, Science \textbf{258}, 1498 (1992).
\bibitem{Hasselmo2006} M. E. Hasselmo, \emph{The role of acetylcholine in learning and memory}, Curr. Opin. Neurobiol. \textbf{16}, 710 (2006).
\bibitem{JacobsAzmitia1992} B. L. Jacobs and E. C. Azmitia, \emph{Structure and function of the brain serotonin system}, Physiol. Rev. \textbf{72}, 165 (1992).
\bibitem{Schiller1992} P. H. Schiller, \emph{The ON and OFF channels of the visual system}, Trends Neurosci. \textbf{15}, 86 (1992).
\bibitem{Grothe2010} B. Grothe, M. Pecka, and D. McAlpine, \emph{Mechanisms of sound localization in mammals}, Physiol. Rev. \textbf{90}, 983 (2010).
\bibitem{Markram2004} H. Markram et al., \emph{Interneurons of the neocortical inhibitory system}, Nat. Rev. Neurosci. \textbf{5}, 793 (2004).
\bibitem{Joshi2016} S. Joshi, Y. Li, R. M. Kalwani, and J. I. Gold, \emph{Relationships between pupil diameter and neuronal activity in the locus coeruleus, colliculi, and cingulate cortex}, Neuron \textbf{89}, 221 (2016).
\bibitem{ServanSchreiber1990} D. Servan-Schreiber, H. Printz, and J. D. Cohen, \emph{A network model of catecholamine effects: gain, signal-to-noise ratio, and behavior}, Science \textbf{249}, 892 (1990).
\bibitem{YuDayan2005} A. J. Yu and P. Dayan, \emph{Uncertainty, neuromodulation, and attention}, Neuron \textbf{46}, 681 (2005).
\bibitem{Nassar2012} M. R. Nassar et al., \emph{Rational regulation of learning dynamics by pupil-linked arousal systems}, Nat. Neurosci. \textbf{15}, 1040 (2012).
\bibitem{BouretSara2005} S. Bouret and S. J. Sara, \emph{Network reset: a simplified overarching theory of locus coeruleus noradrenaline function}, Trends Neurosci. \textbf{28}, 574 (2005).
\bibitem{Hebb1949} D. O. Hebb, \emph{The Organization of Behavior: A Neuropsychological Theory} (Wiley, New York, 1949).
\bibitem{MacKay1952} D. M. MacKay and W. S. McCulloch, \emph{The limiting information capacity of a neuronal link}, Bull. Math. Biophys. \textbf{14}, 127 (1952).
\bibitem{WilsonCowan1972} H. R. Wilson and J. D. Cowan, \emph{Excitatory and inhibitory interactions in localized populations of model neurons}, Biophys. J. \textbf{12}, 1 (1972).
\bibitem{AllenStevens1994} C. Allen and C. F. Stevens, \emph{An evaluation of causes for unreliability of synaptic transmission}, Proc. Natl. Acad. Sci. USA \textbf{91}, 10380 (1994).
\bibitem{Barlow1961} H. B. Barlow, \emph{Possible principles underlying the transformations of sensory messages}, in \emph{Sensory Communication}, ed. W. A. Rosenblith (MIT Press, Cambridge, MA, 1961), pp.~217--234.
\bibitem{Cardin2009} J. A. Cardin et al., \emph{Driving fast-spiking cells induces gamma rhythm and controls sensory responses}, Nature \textbf{459}, 663 (2009).
\bibitem{Lillicrap2020} T. P. Lillicrap, A. Santoro, L. Marris, C. J. Akerman, and G. Hinton, \emph{Backpropagation and the brain}, Nat. Rev. Neurosci. \textbf{21}, 335 (2020).
\bibitem{Fremaux2016} N. Fr\'emaux and W. Gerstner, \emph{Neuromodulated spike-timing-dependent plasticity, and theory of three-factor learning rules}, Front. Neural Circuits \textbf{9}, 85 (2016).
\bibitem{Haider2006} B. Haider, A. Duque, A. R. Hasenstaub, and D. A. McCormick, \emph{Neocortical network activity in vivo is generated through a dynamic balance of excitation and inhibition}, J. Neurosci. \textbf{26}, 4535 (2006).
\bibitem{HoltKoch1997} G. R. Holt and C. Koch, \emph{Shunting inhibition does not have a divisive effect on firing rates}, Neural Comput. \textbf{9}, 1001 (1997).
\bibitem{Chance2002} F. S. Chance, L. F. Abbott, and A. D. Reyes, \emph{Gain modulation from background synaptic input}, Neuron \textbf{35}, 773 (2002).
\bibitem{Ozeki2009} H. Ozeki, I. M. Finn, E. S. Schaffer, K. D. Miller, and D. Ferster, \emph{Inhibitory stabilization of the cortical network underlies visual surround suppression}, Neuron \textbf{62}, 578 (2009).
\bibitem{HasselmoBower1992} M. E. Hasselmo and J. M. Bower, \emph{Cholinergic suppression specific to intrinsic not afferent fiber synapses in rat piriform (olfactory) cortex}, J. Neurophysiol. \textbf{67}, 1222 (1992).
\bibitem{HasselmoAndersonBower1992} M. E. Hasselmo, B. P. Anderson, and J. M. Bower, \emph{Cholinergic modulation of cortical associative memory function}, J. Neurophysiol. \textbf{67}, 1230 (1992).
\bibitem{Hasselmo1999} M. E. Hasselmo, \emph{Neuromodulation: acetylcholine and memory consolidation}, Trends Cogn. Sci. \textbf{3}, 351 (1999).
\bibitem{Tsodyks1988} M. V. Tsodyks and M. V. Feigel'man, \emph{The enhanced storage capacity in neural networks with low activity level}, Europhys. Lett. \textbf{6}, 101 (1988).
\bibitem{KalmanBucy1961} R. E. Kalman and R. S. Bucy, \emph{New results in linear filtering and prediction theory}, J. Basic Eng. \textbf{83}, 95 (1961).
\bibitem{WilsonMcNaughton1994} M. A. Wilson and B. L. McNaughton, \emph{Reactivation of hippocampal ensemble memories during sleep}, Science \textbf{265}, 676 (1994).
\bibitem{BarnesSharp1999} N. M. Barnes and T. Sharp, \emph{A review of central 5-HT receptors and their function}, Neuropharmacology \textbf{38}, 1083 (1999).
\bibitem{Bunin1998} M. A. Bunin and R. M. Wightman, \emph{Quantitative evaluation of 5-hydroxytryptamine (serotonin) neuronal release and uptake: an investigation of extrasynaptic transmission}, J. Neurosci. \textbf{18}, 4854 (1998).
\bibitem{Sprouse1987} J. S. Sprouse and G. K. Aghajanian, \emph{Electrophysiological responses of serotoninergic dorsal raphe neurons to 5-HT$_{1A}$ and 5-HT$_{1B}$ agonists}, Synapse \textbf{1}, 3 (1987).
\bibitem{Sykova2008} E. Syková and C. Nicholson, \emph{Diffusion in brain extracellular space}, Physiol. Rev. \textbf{88}, 1277 (2008).
\bibitem{WilsonNicoll2001} R. I. Wilson and R. A. Nicoll, \emph{Endogenous cannabinoids mediate retrograde signalling at hippocampal synapses}, Nature \textbf{410}, 588 (2001).
\bibitem{Garthwaite2008} J. Garthwaite, \emph{Concepts of neural nitric oxide-mediated transmission}, Eur. J. Neurosci. \textbf{27}, 2783 (2008).
\bibitem{vandenPol2012} A. N. van den Pol, \emph{Neuropeptide transmission in brain circuits}, Neuron \textbf{76}, 98 (2012).
\bibitem{Dunwiddie2001} T. V. Dunwiddie and S. A. Masino, \emph{The role and regulation of adenosine in the central nervous system}, Annu. Rev. Neurosci. \textbf{24}, 31 (2001).
\bibitem{Skaggs1996} W. E. Skaggs and B. L. McNaughton, \emph{Replay of neuronal firing sequences in rat hippocampus during sleep following spatial experience}, Science \textbf{271}, 1870 (1996).
\bibitem{Baylor1979} D. A. Baylor, T. D. Lamb, and K.-W. Yau, \emph{Responses of retinal rods to single photons}, J. Physiol. \textbf{288}, 613 (1979).
\bibitem{Hudspeth1989} A. J. Hudspeth, \emph{How the ear's works work}, Nature \textbf{341}, 397 (1989).
\bibitem{NakaRushton1966} K. I. Naka and W. A. H. Rushton, \emph{S-potentials from colour units in the retina of fish (Cyprinidae)}, J. Physiol. \textbf{185}, 536 (1966).
\bibitem{Lichtsteiner2008} P. Lichtsteiner, C. Posch, and T. Delbruck, \emph{A $128\times128$ 120 dB 15 $\mu$s latency asynchronous temporal contrast vision sensor}, IEEE J. Solid-State Circuits \textbf{43}, 566 (2008).
\bibitem{Koch2006} K. Koch et al., \emph{How much the eye tells the brain}, Curr. Biol. \textbf{16}, 1428 (2006).
\bibitem{Robles2001} L. Robles and M. A. Ruggero, \emph{Mechanics of the mammalian cochlea}, Physiol. Rev. \textbf{81}, 1305 (2001).
\bibitem{Mombaerts2004} P. Mombaerts, \emph{Genes and ligands for odorant, vomeronasal and taste receptors}, Nat. Rev. Neurosci. \textbf{5}, 263 (2004).
\bibitem{WoodSlater2001} S. J. Wood and C. R. Slater, \emph{Safety factor at the neuromuscular junction}, Prog. Neurobiol. \textbf{64}, 393 (2001).
\bibitem{Henneman1957} E. Henneman, \emph{Relation between size of neurons and their susceptibility to discharge}, Science \textbf{126}, 1345 (1957).
\bibitem{HarrisWolpert1998} C. M. Harris and D. M. Wolpert, \emph{Signal-dependent noise determines motor planning}, Nature \textbf{394}, 780 (1998).
\bibitem{Hogan1984} N. Hogan, \emph{Adaptive control of mechanical impedance by coactivation of antagonist muscles}, IEEE Trans. Autom. Control \textbf{29}, 681 (1984).
\bibitem{McAuleyMarsden2000} J. H. McAuley and C. D. Marsden, \emph{Physiological and pathological tremors and rhythmic central motor control}, Brain \textbf{123}, 1545 (2000).
\bibitem{WolpertGhahramani2000} D. M. Wolpert and Z. Ghahramani, \emph{Computational principles of movement neuroscience}, Nat. Neurosci. \textbf{3}, 1212 (2000).
\bibitem{MarderBucher2001} E. Marder and D. Bucher, \emph{Central pattern generators and the control of rhythmic movements}, Curr. Biol. \textbf{11}, R986 (2001).
\bibitem{Miyazaki2014} K. W. Miyazaki et al., \emph{Optogenetic activation of dorsal raphe serotonin neurons enhances patience for future rewards}, Curr. Biol. \textbf{24}, 2033 (2014).
\bibitem{Fuglevand1993} A. J. Fuglevand, D. A. Winter, and A. E. Patla, \emph{Models of recruitment and rate coding organization in motor-unit pools}, J. Neurophysiol. \textbf{70}, 2470 (1993).
\bibitem{Curcio1990} C. A. Curcio, K. R. Sloan, R. E. Kalina, and A. E. Hendrickson, \emph{Human photoreceptor topography}, J. Comp. Neurol. \textbf{292}, 497 (1990).
\bibitem{Hayes1950} N. D. Hayes, \emph{Roots of the transcendental equation associated with a certain difference-differential equation}, J. London Math. Soc. \textbf{25}, 226 (1950).
\bibitem{MelzackWall1965} R. Melzack and P. D. Wall, \emph{Pain mechanisms: a new theory}, Science \textbf{150}, 971 (1965).
\bibitem{Fields2004} H. Fields, \emph{State-dependent opioid control of pain}, Nat. Rev. Neurosci. \textbf{5}, 565 (2004).
\bibitem{Levine1978} J. D. Levine, N. C. Gordon, and H. L. Fields, \emph{The mechanism of placebo analgesia}, Lancet \textbf{312}, 654 (1978).
\bibitem{Woolf1983} C. J. Woolf, \emph{Evidence for a central component of post-injury pain hypersensitivity}, Nature \textbf{306}, 686 (1983).
\bibitem{Seymour2004} B. Seymour et al., \emph{Temporal difference models describe higher-order learning in humans}, Nature \textbf{429}, 664 (2004).
\bibitem{EcclestonCrombez1999} C. Eccleston and G. Crombez, \emph{Pain demands attention: a cognitive-affective model of the interruptive function of pain}, Psychol. Bull. \textbf{125}, 356 (1999).
\bibitem{WidrowHoff1960} B. Widrow and M. E. Hoff, \emph{Adaptive switching circuits}, IRE WESCON Convention Record, part 4, 96 (1960).
\bibitem{Marr1969} D. Marr, \emph{A theory of cerebellar cortex}, J. Physiol. \textbf{202}, 437 (1969).
\bibitem{Albus1971} J. S. Albus, \emph{A theory of cerebellar function}, Math. Biosci. \textbf{10}, 25 (1971).
\bibitem{Cover1965} T. M. Cover, \emph{Geometrical and statistical properties of systems of linear inequalities with applications in pattern recognition}, IEEE Trans. Electron. Comput. \textbf{EC-14}, 326 (1965).
\bibitem{Ito2001} M. Ito, \emph{Cerebellar long-term depression: characterization, signal transduction, and functional roles}, Physiol. Rev. \textbf{81}, 1143 (2001).
\bibitem{Suvrathan2016} A. Suvrathan, H. L. Payne, and J. L. Raymond, \emph{Timing rules for synaptic plasticity matched to behavioral function}, Neuron \textbf{92}, 959 (2016).
\bibitem{Fujita1982} M. Fujita, \emph{Adaptive filter model of the cerebellum}, Biol. Cybern. \textbf{45}, 195 (1982).
\bibitem{Blakemore1998} S.-J. Blakemore, D. M. Wolpert, and C. D. Frith, \emph{Central cancellation of self-produced tickle sensation}, Nat. Neurosci. \textbf{1}, 635 (1998).
\bibitem{Smith2006} M. A. Smith, A. Ghazizadeh, and R. Shadmehr, \emph{Interacting adaptive processes with different timescales underlie short-term motor learning}, PLoS Biol. \textbf{4}, e179 (2006).
\bibitem{Kording2007} K. P. K\"ording, J. B. Tenenbaum, and R. Shadmehr, \emph{The dynamics of memory as a consequence of optimal adaptation to a changing body}, Nat. Neurosci. \textbf{10}, 779 (2007).
\bibitem{Yao2023} Z. Yao et al., \emph{A high-resolution transcriptomic and spatial atlas of cell types in the whole mouse brain}, Nature \textbf{624}, 317 (2023).
\bibitem{Sanzeni2020} A. Sanzeni, B. Akitake, H. C. Goldbach, C. E. Leedy, N. Brunel, and M. H. Histed, \emph{Inhibition stabilization is a widespread property of cortical networks}, eLife \textbf{9}, e54875 (2020).
\bibitem{Padamsey2022} Z. Padamsey, D. Katsanevaki, N. Dupuy, and N. L. Rochefort, \emph{Neocortex saves energy by reducing coding precision during food scarcity}, Neuron \textbf{110}, 280 (2022).
\bibitem{Stringer2019} C. Stringer, M. Pachitariu, N. Steinmetz, C. B. Reddy, M. Carandini, and K. D. Harris, \emph{Spontaneous behaviors drive multidimensional, brainwide activity}, Science \textbf{364}, eaav7893 (2019).
\bibitem{Bittner2017} K. C. Bittner, A. D. Milstein, C. Grienberger, S. Romani, and J. C. Magee, \emph{Behavioral time scale synaptic plasticity underlies CA1 place fields}, Science \textbf{357}, 1033 (2017).
\bibitem{WhittingtonBogacz2019} J. C. R. Whittington and R. Bogacz, \emph{Theories of error back-propagation in the brain}, Trends Cogn. Sci. \textbf{23}, 235 (2019).
\bibitem{Reimer2016} J. Reimer et al., \emph{Pupil fluctuations track rapid changes in adrenergic and cholinergic activity in cortex}, Nat. Commun. \textbf{7}, 13289 (2016).
\bibitem{Beniaguev2021} D. Beniaguev, I. Segev, and M. London, \emph{Single cortical neurons as deep artificial neural networks}, Neuron \textbf{109}, 2727 (2021).
\bibitem{Gidon2020} A. Gidon et al., \emph{Dendritic action potentials and computation in human layer 2/3 cortical neurons}, Science \textbf{367}, 83 (2020).
\bibitem{Gallego2022} G. Gallego et al., \emph{Event-based vision: a survey}, IEEE Trans. Pattern Anal. Mach. Intell. \textbf{44}, 154 (2022).
\bibitem{Berger1989} R. D. Berger, J. P. Saul, and R. J. Cohen, \emph{Transfer function analysis of autonomic regulation. I. Canine atrial rate response}, Am. J. Physiol. \textbf{256}, H142 (1989).
\bibitem{Walker2010} J. J. Walker, J. R. Terry, and S. L. Lightman, \emph{Origin of ultradian pulsatility in the hypothalamic--pituitary--adrenal axis}, Proc. R. Soc. B \textbf{277}, 1627 (2010).
\bibitem{Belchetz1978} P. E. Belchetz, T. M. Plant, Y. Nakai, E. J. Keogh, and E. Knobil, \emph{Hypophysial responses to continuous and intermittent delivery of hypothalamic gonadotropin-releasing hormone}, Science \textbf{202}, 631 (1978).
\bibitem{Czeisler1999} C. A. Czeisler et al., \emph{Stability, precision, and near-24-hour period of the human circadian pacemaker}, Science \textbf{284}, 2177 (1999).
\bibitem{TakahashiJS2017} J. S. Takahashi, \emph{Transcriptional architecture of the mammalian circadian clock}, Nat. Rev. Genet. \textbf{18}, 164 (2017).
\bibitem{Musk2019} E. Musk and Neuralink, \emph{An integrated brain-machine interface platform with thousands of channels}, J. Med. Internet Res. \textbf{21}, e16194 (2019).
\bibitem{Hochberg2006} L. R. Hochberg et al., \emph{Neuronal ensemble control of prosthetic devices by a human with tetraplegia}, Nature \textbf{442}, 164 (2006).
\bibitem{Willett2021} F. R. Willett et al., \emph{High-performance brain-to-text communication via handwriting}, Nature \textbf{593}, 249 (2021).
\bibitem{Willett2023} F. R. Willett et al., \emph{A high-performance speech neuroprosthesis}, Nature \textbf{620}, 1031 (2023).
\bibitem{Gallego2017} J. A. Gallego, M. G. Perich, L. E. Miller, and S. A. Solla, \emph{Neural manifolds for the control of movement}, Neuron \textbf{94}, 978 (2017).
\bibitem{Fernandez2021} E. Fern\'andez et al., \emph{Visual percepts evoked with an intracortical 96-channel microelectrode array inserted in human occipital cortex}, J. Clin. Invest. \textbf{131}, e151331 (2021).
\bibitem{Tritsch2012} N. X. Tritsch, J. B. Ding, and B. L. Sabatini, \emph{Dopaminergic neurons inhibit striatal output through non-canonical release of GABA}, Nature \textbf{490}, 262 (2012).
\bibitem{Zhang2015} S. Zhang et al., \emph{Dopaminergic and glutamatergic microdomains in a subset of rodent mesoaccumbens axons}, Nat. Neurosci. \textbf{18}, 386 (2015).
\bibitem{Mongillo2008} G. Mongillo, O. Barak, and M. Tsodyks, \emph{Synaptic theory of working memory}, Science \textbf{319}, 1543 (2008).
\bibitem{Stokes2015} M. G. Stokes, \emph{`Activity-silent' working memory in prefrontal cortex: a dynamic coding framework}, Trends Cogn. Sci. \textbf{19}, 394 (2015).
\bibitem{Smith2013} S. L. Smith, I. T. Smith, T. Branco, and M. H\"ausser, \emph{Dendritic spikes enhance stimulus selectivity in cortical neurons in vivo}, Nature \textbf{503}, 115 (2013).
\bibitem{Baden2016} T. Baden et al., \emph{The functional diversity of retinal ganglion cells in the mouse}, Nature \textbf{529}, 345 (2016).
\bibitem{Lottem2018} E. Lottem et al., \emph{Activation of serotonin neurons promotes active persistence in a probabilistic foraging task}, Nat. Commun. \textbf{9}, 1000 (2018).
\bibitem{Payeur2021} A. Payeur, J. Guerguiev, F. Zenke, B. A. Richards, and R. Naud, \emph{Burst-dependent synaptic plasticity can coordinate learning in hierarchical circuits}, Nat. Neurosci. \textbf{24}, 1010 (2021).
\bibitem{MehonicKenyon2022} A. Mehonic and A. J. Kenyon, \emph{Brain-inspired computing needs a master plan}, Nature \textbf{604}, 255 (2022).
\bibitem{Poe2020} G. R. Poe et al., \emph{Locus coeruleus: a new look at the blue spot}, Nat. Rev. Neurosci. \textbf{21}, 644 (2020).
\bibitem{Hangya2015} B. Hangya, S. P. Ranade, M. Lorenc, and A. Kepecs, \emph{Central cholinergic neurons are rapidly recruited by reinforcement feedback}, Cell \textbf{162}, 1155 (2015).
\bibitem{FernandezRuiz2019} A. Fern\'andez-Ruiz et al., \emph{Long-duration hippocampal sharp wave ripples improve memory}, Science \textbf{364}, 1082 (2019).
\bibitem{Herzfeld2018} D. J. Herzfeld, Y. Kojima, R. Soetedjo, and R. Shadmehr, \emph{Encoding of error and learning to correct that error by the Purkinje cells of the cerebellum}, Nat. Neurosci. \textbf{21}, 736 (2018).
\bibitem{Duan2014} B. Duan et al., \emph{Identification of spinal circuits transmitting and gating mechanical pain}, Cell \textbf{159}, 1417 (2014).
\bibitem{Tremblay2016} R. Tremblay, S. Lee, and B. Rudy, \emph{GABAergic interneurons in the neocortex: from cellular properties to circuits}, Neuron \textbf{91}, 260 (2016).
\bibitem{Ren2018} J. Ren et al., \emph{Anatomically defined and functionally distinct dorsal raphe serotonin sub-systems}, Cell \textbf{175}, 472 (2018).
\bibitem{Pfeffer2013} C. K. Pfeffer, M. Xue, M. He, Z. J. Huang, and M. Scanziani, \emph{Inhibition of inhibition in visual cortex: the logic of connections between molecularly distinct interneurons}, Nat. Neurosci. \textbf{16}, 1068 (2013).
\bibitem{Pi2013} H.-J. Pi et al., \emph{Cortical interneurons that specialize in disinhibitory control}, Nature \textbf{503}, 521 (2013).
\bibitem{Keller2018} G. B. Keller and T. D. Mrsic-Flogel, \emph{Predictive processing: a canonical cortical computation}, Neuron \textbf{100}, 424 (2018).
\bibitem{HubelWiesel1962} D. H. Hubel and T. N. Wiesel, \emph{Receptive fields, binocular interaction and functional architecture in the cat's visual cortex}, J. Physiol. \textbf{160}, 106 (1962).
\bibitem{Riesenhuber1999} M. Riesenhuber and T. Poggio, \emph{Hierarchical models of object recognition in cortex}, Nat. Neurosci. \textbf{2}, 1019 (1999).
\bibitem{Fukushima1980} K. Fukushima, \emph{Neocognitron: a self-organizing neural network model for a mechanism of pattern recognition unaffected by shift in position}, Biol. Cybern. \textbf{36}, 193 (1980).
\bibitem{Yamins2014} D. L. K. Yamins et al., \emph{Performance-optimized hierarchical models predict neural responses in higher visual cortex}, Proc. Natl. Acad. Sci. USA \textbf{111}, 8619 (2014).
\bibitem{Hung2005} C. P. Hung, G. Kreiman, T. Poggio, and J. J. DiCarlo, \emph{Fast readout of object identity from macaque inferior temporal cortex}, Science \textbf{310}, 863 (2005).
\bibitem{WiskottSejnowski2002} L. Wiskott and T. J. Sejnowski, \emph{Slow feature analysis: unsupervised learning of invariances}, Neural Comput. \textbf{14}, 715 (2002).
\bibitem{Foldiak1991} P. F\"oldi\'ak, \emph{Learning invariance from transformation sequences}, Neural Comput. \textbf{3}, 194 (1991).
\bibitem{LiDiCarlo2008} N. Li and J. J. DiCarlo, \emph{Unsupervised natural experience rapidly alters invariant object representation in visual cortex}, Science \textbf{321}, 1502 (2008).
\bibitem{RaoBallard1999} R. P. N. Rao and D. H. Ballard, \emph{Predictive coding in the visual cortex: a functional interpretation of some extra-classical receptive-field effects}, Nat. Neurosci. \textbf{2}, 79 (1999).
\bibitem{Kar2019} K. Kar, J. Kubilius, K. Schmidt, E. B. Issa, and J. J. DiCarlo, \emph{Evidence that recurrent circuits are critical to the ventral stream's execution of core object recognition behavior}, Nat. Neurosci. \textbf{22}, 974 (2019).
\bibitem{Szegedy2014} C. Szegedy et al., \emph{Intriguing properties of neural networks}, in \emph{Proc. ICLR} (2014), arXiv:1312.6199.
\bibitem{Weiss2002} Y. Weiss, E. P. Simoncelli, and E. H. Adelson, \emph{Motion illusions as optimal percepts}, Nat. Neurosci. \textbf{5}, 598 (2002).
\bibitem{Purves2011} D. Purves, W. T. Wojtach, and R. B. Lotto, \emph{Understanding vision in wholly empirical terms}, Proc. Natl. Acad. Sci. USA \textbf{108}, Suppl.~3, 15588 (2011).
\bibitem{LaferSousa2015} R. Lafer-Sousa, K. L. Hermann, and B. R. Conway, \emph{Striking individual differences in color perception uncovered by `the dress' photograph}, Curr. Biol. \textbf{25}, R545 (2015).
\bibitem{Wallisch2017} P. Wallisch, \emph{Illumination assumptions account for individual differences in the perceptual interpretation of a profoundly ambiguous stimulus in the color domain: ``The dress''}, J. Vis. \textbf{17}(4), 5 (2017).
\bibitem{Schwartz2009} O. Schwartz, T. J. Sejnowski, and P. Dayan, \emph{Perceptual organization in the tilt illusion}, J. Vis. \textbf{9}(4), 19 (2009).
\bibitem{SchillerCarvey2005} P. H. Schiller and C. E. Carvey, \emph{The Hermann grid illusion revisited}, Perception \textbf{34}, 1375 (2005).
\bibitem{Nijhawan1994} R. Nijhawan, \emph{Motion extrapolation in catching}, Nature \textbf{370}, 256 (1994).
\bibitem{Conway2005} B. R. Conway, A. Kitaoka, A. Yazdanbakhsh, C. C. Pack, and M. S. Livingstone, \emph{Neural basis for a powerful static motion illusion}, J. Neurosci. \textbf{25}, 5651 (2005).
\bibitem{OteroMillan2012} J. Otero-Millan, S. L. Macknik, and S. Martinez-Conde, \emph{Microsaccades and blinks trigger illusory rotation in the ``rotating snakes'' illusion}, J. Neurosci. \textbf{32}, 6043 (2012).
\bibitem{MorenoBote2007} R. Moreno-Bote, J. Rinzel, and N. Rubin, \emph{Noise-induced alternations in an attractor network model of perceptual bistability}, J. Neurophysiol. \textbf{98}, 1125 (2007).
\bibitem{Watanabe2018} E. Watanabe, A. Kitaoka, K. Sakamoto, M. Yasugi, and K. Tanaka, \emph{Illusory motion reproduced by deep neural networks trained for prediction}, Front. Psychol. \textbf{9}, 345 (2018).
\bibitem{GomezVilla2020} A. G\'omez-Villa, A. Mart\'in, J. Vazquez-Corral, M. Bertalm\'io, and J. Malo, \emph{Color illusions also deceive CNNs for low-level vision tasks: analysis and implications}, Vision Res. \textbf{176}, 156 (2020).
\bibitem{delCastillo1954} J. del Castillo and B. Katz, \emph{Quantal components of the end-plate potential}, J. Physiol. \textbf{124}, 560 (1954).
\bibitem{Nowak1984} L. Nowak, P. Bregestovski, P. Ascher, A. Herbet, and A. Prochiantz, \emph{Magnesium gates glutamate-activated channels in mouse central neurones}, Nature \textbf{307}, 462 (1984).
\bibitem{Malinow2002} R. Malinow and R. C. Malenka, \emph{AMPA receptor trafficking and synaptic plasticity}, Annu. Rev. Neurosci. \textbf{25}, 103 (2002).
\bibitem{Matsuzaki2004} M. Matsuzaki, N. Honkura, G. C. R. Ellis-Davies, and H. Kasai, \emph{Structural basis of long-term potentiation in single dendritic spines}, Nature \textbf{429}, 761 (2004).
\bibitem{Rudomin1999} P. Rudomin and R. F. Schmidt, \emph{Presynaptic inhibition in the vertebrate spinal cord revisited}, Exp. Brain Res. \textbf{129}, 1 (1999).
\bibitem{HutcheonYarom2000} B. Hutcheon and Y. Yarom, \emph{Resonance, oscillation and the intrinsic frequency preferences of neurons}, Trends Neurosci. \textbf{23}, 216 (2000).
\bibitem{Seung1996} H. S. Seung, \emph{How the brain keeps the eyes still}, Proc. Natl. Acad. Sci. USA \textbf{93}, 13339 (1996).
\bibitem{Hounsgaard1988} J. Hounsgaard, H. Hultborn, B. Jespersen, and O. Kiehn, \emph{Bistability of alpha-motoneurones in the decerebrate cat and in the acute spinal cat after intravenous 5-hydroxytryptophan}, J. Physiol. \textbf{405}, 345 (1988).
\bibitem{Izhikevich2007} E. M. Izhikevich, \emph{Dynamical Systems in Neuroscience: The Geometry of Excitability and Bursting} (MIT Press, Cambridge, MA, 2007).
\bibitem{AgmonSnir1998} H. Agmon-Snir, C. E. Carr, and J. Rinzel, \emph{The role of dendrites in auditory coincidence detection}, Nature \textbf{393}, 268 (1998).
\bibitem{BorstSoria2012} J. G. G. Borst and J. Soria van Hoeve, \emph{The calyx of Held synapse: from model synapse to auditory relay}, Annu. Rev. Physiol. \textbf{74}, 199 (2012).
\bibitem{Euler2002} T. Euler, P. B. Detwiler, and W. Denk, \emph{Directionally selective calcium signals in dendrites of starburst amacrine cells}, Nature \textbf{418}, 845 (2002).
\bibitem{AstonJonesBloom1981} G. Aston-Jones and F. E. Bloom, \emph{Activity of norepinephrine-containing locus coeruleus neurons in behaving rats anticipates fluctuations in the sleep-waking cycle}, J. Neurosci. \textbf{1}, 876 (1981).
\bibitem{BrooksPeever2012} P. L. Brooks and J. H. Peever, \emph{Identification of the transmitter and receptor mechanisms responsible for REM sleep paralysis}, J. Neurosci. \textbf{32}, 9785 (2012).
\bibitem{Borbely1982} A. A. Borb\'ely, \emph{A two process model of sleep regulation}, Hum. Neurobiol. \textbf{1}, 195 (1982).
\bibitem{PorkkaHeiskanen1997} T. Porkka-Heiskanen et al., \emph{Adenosine: a mediator of the sleep-inducing effects of prolonged wakefulness}, Science \textbf{276}, 1265 (1997).
\bibitem{Steriade1993} M. Steriade, A. N\'u\~nez, and F. Amzica, \emph{A novel slow ($<1$ Hz) oscillation of neocortical neurons in vivo: depolarizing and hyperpolarizing components}, J. Neurosci. \textbf{13}, 3252 (1993).
\bibitem{McCormick1992} D. A. McCormick, \emph{Neurotransmitter actions in the thalamus and cerebral cortex and their role in neuromodulation of thalamocortical activity}, Prog. Neurobiol. \textbf{39}, 337 (1992).
\bibitem{SteriadeMcCormickSejnowski1993} M. Steriade, D. A. McCormick, and T. J. Sejnowski, \emph{Thalamocortical oscillations in the sleeping and aroused brain}, Science \textbf{262}, 679 (1993).
\bibitem{Nadasdy1999} Z. N\'adasdy et al., \emph{Replay and time compression of recurring spike sequences in the hippocampus}, J. Neurosci. \textbf{19}, 9497 (1999).
\bibitem{Maingret2016} N. Maingret et al., \emph{Hippocampo-cortical coupling mediates memory consolidation during sleep}, Nat. Neurosci. \textbf{19}, 959 (2016).
\bibitem{McClelland1995} J. L. McClelland, B. L. McNaughton, and R. C. O'Reilly, \emph{Why there are complementary learning systems in the hippocampus and neocortex}, Psychol. Rev. \textbf{102}, 419 (1995).
\bibitem{TononiCirelli2014} G. Tononi and C. Cirelli, \emph{Sleep and the price of plasticity: from synaptic and cellular homeostasis to memory consolidation and integration}, Neuron \textbf{81}, 12 (2014).
\bibitem{deVivo2017} L. de Vivo et al., \emph{Ultrastructural evidence for synaptic scaling across the wake/sleep cycle}, Science \textbf{355}, 507 (2017).
\bibitem{Xie2013} L. Xie et al., \emph{Sleep drives metabolite clearance from the adult brain}, Science \textbf{342}, 373 (2013).
\bibitem{Miao2024} A. Miao et al., \emph{Brain clearance is reduced during sleep and anesthesia}, Nat. Neurosci. \textbf{27}, 1046 (2024).
\bibitem{Vyazovskiy2011} V. V. Vyazovskiy et al., \emph{Local sleep in awake rats}, Nature \textbf{472}, 443 (2011).
\bibitem{Lancaster2013} M. A. Lancaster et al., \emph{Cerebral organoids model human brain development and microcephaly}, Nature \textbf{501}, 373 (2013).
\bibitem{Jordan2024} F. D. Jordan et al., \emph{Open and remotely accessible Neuroplatform for research in wetware computing}, Front. Artif. Intell. \textbf{7}, 1376042 (2024).
\bibitem{Smirnova2023} L. Smirnova et al., \emph{Organoid intelligence (OI): the new frontier in biocomputing and intelligence-in-a-dish}, Front. Sci. \textbf{1}, 1017235 (2023).
\bibitem{Wagenaar2006} D. A. Wagenaar, J. Pine, and S. M. Potter, \emph{An extremely rich repertoire of bursting patterns during the development of cortical cultures}, BMC Neurosci. \textbf{7}, 11 (2006).
\bibitem{Shahaf2001} G. Shahaf and S. Marom, \emph{Learning in networks of cortical neurons}, J. Neurosci. \textbf{21}, 8782 (2001).
\bibitem{Kagan2022} B. J. Kagan et al., \emph{In vitro neurons learn and exhibit sentience when embodied in a simulated game-world}, Neuron \textbf{110}, 3952 (2022).
\bibitem{Friston2010} K. Friston, \emph{The free-energy principle: a unified brain theory?}, Nat. Rev. Neurosci. \textbf{11}, 127 (2010).
\bibitem{Balci2023} F. Balci et al., \emph{A response to claims of emergent intelligence and sentience in a dish}, Neuron \textbf{111}, 604 (2023).
\bibitem{Bakkum2008} D. J. Bakkum, Z. C. Chao, and S. M. Potter, \emph{Spatio-temporal electrical stimuli shape behavior of an embodied cortical network in a goal-directed learning task}, J. Neural Eng. \textbf{5}, 310 (2008).
\bibitem{Maass2002} W. Maass, T. Natschl\"ager, and H. Markram, \emph{Real-time computing without stable states: a new framework for neural computation based on perturbations}, Neural Comput. \textbf{14}, 2531 (2002).
\bibitem{JaegerHaas2004} H. Jaeger and H. Haas, \emph{Harnessing nonlinearity: predicting chaotic systems and saving energy in wireless communication}, Science \textbf{304}, 78 (2004).
\bibitem{Cai2023} H. Cai et al., \emph{Brain organoid reservoir computing for artificial intelligence}, Nat. Electron. \textbf{6}, 1032 (2023).
\bibitem{Habibollahi2023} F. Habibollahi, B. J. Kagan, A. N. Burkitt, and C. French, \emph{Critical dynamics arise during structured information presentation within embodied in vitro neuronal networks}, Nat. Commun. \textbf{14}, 5287 (2023).
\bibitem{Robbins2026} A. Robbins et al., \emph{Goal-directed learning in cortical organoids}, Cell Rep. \textbf{45}, 116984 (2026).
\bibitem{Keene2020} S. T. Keene et al., \emph{A biohybrid synapse with neurotransmitter-mediated plasticity}, Nat. Mater. \textbf{19}, 969 (2020).
\bibitem{Jo2016} J. Jo et al., \emph{Midbrain-like organoids from human pluripotent stem cells contain functional dopaminergic and neuromelanin-producing neurons}, Cell Stem Cell \textbf{19}, 248 (2016).
\bibitem{Hendry1987} S. H. Hendry, H. D. Schwark, E. G. Jones, and J. Yan, \emph{Numbers and proportions of GABA-immunoreactive neurons in different areas of monkey cerebral cortex}, J. Neurosci. \textbf{7}, 1503 (1987).
\bibitem{Tang2001synapses} Y. Tang, J. R. Nyengaard, D. M. G. De Groot, and H. J. G. Gundersen, \emph{Total regional and global number of synapses in the human brain neocortex}, Synapse \textbf{41}, 258 (2001).
\bibitem{Pakkenberg1997} B. Pakkenberg and H. J. G. Gundersen, \emph{Neocortical neuron number in humans: effect of sex and age}, J. Comp. Neurol. \textbf{384}, 312 (1997).
\bibitem{Matsuda2009} W. Matsuda, T. Furuta, K. C. Nakamura, H. Hioki, F. Fujiyama, R. Arai, and T. Kaneko, \emph{Single nigrostriatal dopaminergic neurons form widely spread and highly dense axonal arborizations in the neostriatum}, J. Neurosci. \textbf{29}, 444 (2009).
\bibitem{Pakkenberg1991} B. Pakkenberg, A. M{\o}ller, H. J. G. Gundersen, A. Mouritzen Dam, and H. Pakkenberg, \emph{The absolute number of nerve cells in substantia nigra in normal subjects and in patients with Parkinson's disease estimated with an unbiased stereological method}, J. Neurol. Neurosurg. Psychiatry \textbf{54}, 30 (1991).
\bibitem{Baker1990} K. G. Baker, G. M. Halliday, and I. T\"ork, \emph{Cytoarchitecture of the human dorsal raphe nucleus}, J. Comp. Neurol. \textbf{301}, 147 (1990).
\bibitem{Baker1991} K. G. Baker, G. M. Halliday, P. Halasz, J.-P. Hornung, L. B. Geffen, R. G. H. Cotton, and I. T\"ork, \emph{Cytoarchitecture of serotonin-synthesizing neurons in the pontine tegmentum of the human brain}, Synapse \textbf{7}, 301 (1991).
\bibitem{Airaksinen1991} M. S. Airaksinen, A. Paetau, L. Palj\"arvi, K. Reinikainen, P. Riekkinen, R. Suomalainen, and P. Panula, \emph{Histamine neurons in human hypothalamus: anatomy in normal and Alzheimer diseased brains}, Neuroscience \textbf{44}, 465 (1991).
\bibitem{Colquhoun1992} D. Colquhoun, P. Jonas, and B. Sakmann, \emph{Action of brief pulses of glutamate on AMPA/kainate receptors in patches from different neurones of rat hippocampal slices}, J. Physiol. \textbf{458}, 261 (1992).
\bibitem{DutarNicoll1988} P. Dutar and R. A. Nicoll, \emph{A physiological role for GABA$_B$ receptors in the central nervous system}, Nature \textbf{332}, 156 (1988).

\bibitem{Rauch2003} A. Rauch, G. La Camera, H.-R. L\"uscher, W. Senn, and S. Fusi, \emph{Neocortical pyramidal cells respond as integrate-and-fire neurons to in vivo-like input currents}, J. Neurophysiol. \textbf{90}, 1598 (2003).
\bibitem{KlumppEady1956} R. G. Klumpp and H. R. Eady, \emph{Some measurements of interaural time difference thresholds}, J. Acoust. Soc. Am. \textbf{28}, 859 (1956).
\bibitem{Brand2002} A. Brand, O. Behrend, T. Marquardt, D. McAlpine, and B. Grothe, \emph{Precise inhibition is essential for microsecond interaural time difference coding}, Nature \textbf{417}, 543 (2002).
\bibitem{Funahashi1989} S. Funahashi, C. J. Bruce, and P. S. Goldman-Rakic, \emph{Mnemonic coding of visual space in the monkey's dorsolateral prefrontal cortex}, J. Neurophysiol. \textbf{61}, 331 (1989).
\bibitem{WangMarkram2006} Y. Wang, H. Markram, P. H. Goodman, T. K. Berger, J. Ma, and P. S. Goldman-Rakic, \emph{Heterogeneity in the pyramidal network of the medial prefrontal cortex}, Nat. Neurosci. \textbf{9}, 534 (2006).

\bibitem{McCullochPitts1943} W. S. McCulloch and W. Pitts, \emph{A logical calculus of the ideas immanent in nervous activity}, Bull. Math. Biophys. \textbf{5}, 115 (1943).
\bibitem{Fried2002} S. I. Fried, T. A. M\"unch, and F. S. Werblin, \emph{Mechanisms and circuitry underlying directional selectivity in the retina}, Nature \textbf{420}, 411 (2002).

\bibitem{Hromadka2008} T. Hrom\'adka, M. R. DeWeese, and A. M. Zador, \emph{Sparse representation of sounds in the unanesthetized auditory cortex}, PLoS Biol. \textbf{6}, e16 (2008).
\bibitem{Abbott1997depression} L. F. Abbott, J. A. Varela, K. Sen, and S. B. Nelson, \emph{Synaptic depression and cortical gain control}, Science \textbf{275}, 221 (1997).
\bibitem{ShadlenNewsome1998} M. N. Shadlen and W. T. Newsome, \emph{The variable discharge of cortical neurons: implications for connectivity, computation, and information coding}, J. Neurosci. \textbf{18}, 3870 (1998).

\bibitem{TrulsonJacobs1979} M. E. Trulson and B. L. Jacobs, \emph{Raphe unit activity in freely moving cats: correlation with level of behavioral arousal}, Brain Res. \textbf{163}, 135 (1979).
\bibitem{GagnonParent2014} D. Gagnon and M. Parent, \emph{Distribution of VGLUT3 in highly collateralized axons from the rat dorsal raphe nucleus as revealed by single-neuron reconstructions}, PLoS One \textbf{9}, e87709 (2014).
\bibitem{VandermaelenAghajanian1983} C. P. Vandermaelen and G. K. Aghajanian, \emph{Electrophysiological and pharmacological characterization of serotonergic dorsal raphe neurons recorded extracellularly and intracellularly in rat brain slices}, Brain Res. \textbf{289}, 109 (1983).
\bibitem{CourtneyFord2016} N. A. Courtney and C. P. Ford, \emph{Mechanisms of 5-HT$_{1A}$ receptor-mediated transmission in dorsal raphe serotonin neurons}, J. Physiol. \textbf{594}, 953 (2016).
\bibitem{Hashemi2012serotonin} P. Hashemi et al., \emph{Brain dopamine and serotonin differ in regulation and its consequences}, Proc. Natl. Acad. Sci. USA \textbf{109}, 11510 (2012).
\bibitem{Black1934feedback} H. S. Black, \emph{Stabilized feedback amplifiers}, Bell Syst. Tech. J. \textbf{13}, 1 (1934).
\bibitem{KobayashiSchultz2008} S. Kobayashi and W. Schultz, \emph{Influence of reward delays on responses of dopamine neurons}, J. Neurosci. \textbf{28}, 7837 (2008).
\bibitem{Schweighofer2008} N. Schweighofer et al., \emph{Low-serotonin levels increase delayed reward discounting in humans}, J. Neurosci. \textbf{28}, 4528 (2008).

\bibitem{Malenka1988calcium} R. C. Malenka, J. A. Kauer, R. S. Zucker, and R. A. Nicoll, \emph{Postsynaptic calcium is sufficient for potentiation of hippocampal synaptic transmission}, Science \textbf{242}, 81 (1988).
\bibitem{Bliss1973LTP} T. V. P. Bliss and T. L{\o}mo, \emph{Long-lasting potentiation of synaptic transmission in the dentate area of the anaesthetized rabbit following stimulation of the perforant path}, J. Physiol. \textbf{232}, 331 (1973).
\bibitem{Kelso1986hebb} S. R. Kelso, A. H. Ganong, and T. H. Brown, \emph{Hebbian synapses in hippocampus}, Proc. Natl. Acad. Sci. USA \textbf{83}, 5326 (1986).
\bibitem{He2015eligibility} K. He et al., \emph{Distinct eligibility traces for LTP and LTD in cortical synapses}, Neuron \textbf{88}, 528 (2015).
\bibitem{Olveczky2005} B. P. \"Olveczky, A. S. Andalman, and M. S. Fee, \emph{Vocal experimentation in the juvenile songbird requires a basal ganglia circuit}, PLoS Biol. \textbf{3}, e153 (2005).
\bibitem{TumerBrainard2007} E. C. Tumer and M. S. Brainard, \emph{Performance variability enables adaptive plasticity of `crystallized' adult birdsong}, Nature \textbf{450}, 1240 (2007).
\bibitem{Eshel2015arithmetic} N. Eshel et al., \emph{Arithmetic and local circuitry underlying dopamine prediction errors}, Nature \textbf{525}, 243 (2015).
\bibitem{GraceOnn1989} A. A. Grace and S.-P. Onn, \emph{Morphology and electrophysiological properties of immunocytochemically identified rat dopamine neurons recorded in vitro}, J. Neurosci. \textbf{9}, 3463 (1989).
\bibitem{ODoherty2004actor} J. O'Doherty et al., \emph{Dissociable roles of ventral and dorsal striatum in instrumental conditioning}, Science \textbf{304}, 452 (2004).

\bibitem{NeweyPowell1987} W. K. Newey and J. L. Powell, \emph{Asymmetric least squares estimation and testing}, Econometrica \textbf{55}, 819 (1987).
\bibitem{Beierholm2013vigor} U. Beierholm et al., \emph{Dopamine modulates reward-related vigor}, Neuropsychopharmacology \textbf{38}, 1495 (2013).
\bibitem{Howe2013ramp} M. W. Howe, P. L. Tierney, S. G. Sandberg, P. E. M. Phillips, and A. M. Graybiel, \emph{Prolonged dopamine signalling in striatum signals proximity and value of distant rewards}, Nature \textbf{500}, 575 (2013).

\bibitem{Baker1989LC} K. G. Baker, I. T\"ork, J.-P. Hornung, and P. Halasz, \emph{The human locus coeruleus complex: an immunohistochemical and three dimensional reconstruction study}, Exp. Brain Res. \textbf{77}, 257 (1989).
\bibitem{Schwarz2015LC} L. A. Schwarz et al., \emph{Viral-genetic tracing of the input--output organization of a central noradrenaline circuit}, Nature \textbf{524}, 88 (2015).
\bibitem{Uematsu2017LC} A. Uematsu et al., \emph{Modular organization of the brainstem noradrenaline system coordinates opposing learning states}, Nat. Neurosci. \textbf{20}, 1602 (2017).
\bibitem{AstonJones1994vigilance} G. Aston-Jones, J. Rajkowski, P. Kubiak, and T. Alexinsky, \emph{Locus coeruleus neurons in monkey are selectively activated by attended cues in a vigilance task}, J. Neurosci. \textbf{14}, 4467 (1994).
\bibitem{Clayton2004LC} E. C. Clayton, J. Rajkowski, J. D. Cohen, and G. Aston-Jones, \emph{Phasic activation of monkey locus ceruleus neurons by simple decisions in a forced-choice task}, J. Neurosci. \textbf{24}, 9914 (2004).
\bibitem{Foote1975NE} S. L. Foote, R. Freedman, and A. P. Oliver, \emph{Effects of putative neurotransmitters on neuronal activity in monkey auditory cortex}, Brain Res. \textbf{86}, 229 (1975).
\bibitem{JepmaNieuwenhuis2011} M. Jepma and S. Nieuwenhuis, \emph{Pupil diameter predicts changes in the exploration--exploitation trade-off: evidence for the adaptive gain theory}, J. Cogn. Neurosci. \textbf{23}, 1587 (2011).
\bibitem{Tervo2014stochastic} D. G. R. Tervo et al., \emph{Behavioral variability through stochastic choice and its gating by anterior cingulate cortex}, Cell \textbf{159}, 21 (2014).
\bibitem{Usher1999LC} M. Usher, J. D. Cohen, D. Servan-Schreiber, J. Rajkowski, and G. Aston-Jones, \emph{The role of locus coeruleus in the regulation of cognitive performance}, Science \textbf{283}, 549 (1999).

\bibitem{Mesulam1983} M.-M. Mesulam, E. J. Mufson, B. H. Wainer, and A. I. Levey, \emph{Central cholinergic pathways in the rat: an overview based on an alternative nomenclature (Ch1--Ch6)}, Neuroscience \textbf{10}, 1185 (1983).
\bibitem{Marrosu1995} F. Marrosu et al., \emph{Microdialysis measurement of cortical and hippocampal acetylcholine release during sleep-wake cycle in freely moving cats}, Brain Res. \textbf{671}, 329 (1995).
\bibitem{Picciotto2012} M. R. Picciotto, M. J. Higley, and Y. S. Mineur, \emph{Acetylcholine as a neuromodulator: cholinergic signaling shapes nervous system function and behavior}, Neuron \textbf{76}, 116 (2012).
\bibitem{HasselmoSchnell1994} M. E. Hasselmo and E. Schnell, \emph{Laminar selectivity of the cholinergic suppression of synaptic transmission in rat hippocampal region CA1: computational modeling and brain slice physiology}, J. Neurosci. \textbf{14}, 3898 (1994).
\bibitem{Gil1997} Z. Gil, B. W. Connors, and Y. Amitai, \emph{Differential regulation of neocortical synapses by neuromodulators and activity}, Neuron \textbf{19}, 679 (1997).
\bibitem{KilgardMerzenich1998} M. P. Kilgard and M. M. Merzenich, \emph{Cortical map reorganization enabled by nucleus basalis activity}, Science \textbf{279}, 1714 (1998).
\bibitem{Atri2004} A. Atri et al., \emph{Blockade of central cholinergic receptors impairs new learning and increases proactive interference in a word paired-associate memory task}, Behav. Neurosci. \textbf{118}, 223 (2004).
\bibitem{Girardeau2009} G. Girardeau, K. Benchenane, S. I. Wiener, G. Buzs\'aki, and M. B. Zugaro, \emph{Selective suppression of hippocampal ripples impairs spatial memory}, Nat. Neurosci. \textbf{12}, 1222 (2009).

\input{supplement/bib_10}
\bibitem{CopenhagenJahr1989} D. R. Copenhagen and C. E. Jahr, \emph{Release of endogenous excitatory amino acids from turtle photoreceptors}, Nature \textbf{341}, 536 (1989).
\bibitem{GlowatzkiFuchs2002} E. Glowatzki and P. A. Fuchs, \emph{Transmitter release at the hair cell ribbon synapse}, Nat. Neurosci. \textbf{5}, 147 (2002).
\bibitem{Tinsley2016} J. N. Tinsley et al., \emph{Direct detection of a single photon by humans}, Nat. Commun. \textbf{7}, 12172 (2016).
\bibitem{Sellick1982} P. M. Sellick, R. Patuzzi, and B. M. Johnstone, \emph{Measurement of basilar membrane motion in the guinea pig using the M\"ossbauer technique}, J. Acoust. Soc. Am. \textbf{72}, 131 (1982).
\bibitem{Ruggero1997} M. A. Ruggero, N. C. Rich, A. Recio, S. S. Narayan, and L. Robles, \emph{Basilar-membrane responses to tones at the base of the chinchilla cochlea}, J. Acoust. Soc. Am. \textbf{101}, 2151 (1997).
\bibitem{NormannPerlman1979} R. A. Normann and I. Perlman, \emph{The effects of background illumination on the photoresponses of red and green cones}, J. Physiol. \textbf{286}, 491 (1979).
\bibitem{Laughlin1981} S. Laughlin, \emph{A simple coding procedure enhances a neuron's information capacity}, Z. Naturforsch. C \textbf{36}, 910 (1981).
\bibitem{CurcioAllen1990} C. A. Curcio and K. A. Allen, \emph{Topography of ganglion cells in human retina}, J. Comp. Neurol. \textbf{300}, 5 (1990).
\bibitem{Greenwood1990} D. D. Greenwood, \emph{A cochlear frequency-position function for several species --- 29 years later}, J. Acoust. Soc. Am. \textbf{87}, 2592 (1990).
\bibitem{Eguiluz2000} V. M. Egu\'iluz, M. Ospeck, Y. Choe, A. J. Hudspeth, and M. O. Magnasco, \emph{Essential nonlinearities in hearing}, Phys. Rev. Lett. \textbf{84}, 5232 (2000).
\bibitem{PalmerRussell1986} A. R. Palmer and I. J. Russell, \emph{Phase-locking in the cochlear nerve of the guinea-pig and its relation to the receptor potential of inner hair-cells}, Hear. Res. \textbf{24}, 1 (1986).
\bibitem{Malnic1999} B. Malnic, J. Hirono, T. Sato, and L. B. Buck, \emph{Combinatorial receptor codes for odors}, Cell \textbf{96}, 713 (1999).
\bibitem{Finger2005} T. E. Finger et al., \emph{ATP signaling is crucial for communication from taste buds to gustatory nerves}, Science \textbf{310}, 1495 (2005).
\bibitem{Huang2005} Y.-J. Huang et al., \emph{Mouse taste buds use serotonin as a neurotransmitter}, J. Neurosci. \textbf{25}, 843 (2005).
\bibitem{JohanssonVallbo1979} R. S. Johansson and A. B. Vallbo, \emph{Tactile sensibility in the human hand: relative and absolute densities of four types of mechanoreceptive units in glabrous skin}, J. Physiol. \textbf{286}, 283 (1979).
\bibitem{McKemy2002} D. D. McKemy, W. M. Neuhausser, and D. Julius, \emph{Identification of a cold receptor reveals a general role for TRP channels in thermosensation}, Nature \textbf{416}, 52 (2002).

\bibitem{Schmolesky1998} M. T. Schmolesky et al., \emph{Signal timing across the macaque visual system}, J. Neurophysiol. \textbf{79}, 3272 (1998).
\bibitem{Lampl2004} I. Lampl, D. Ferster, T. Poggio, and M. Riesenhuber, \emph{Intracellular measurements of spatial integration and the MAX operation in complex cells of the cat primary visual cortex}, J. Neurophysiol. \textbf{92}, 2704 (2004).
\bibitem{KobatakeTanaka1994} E. Kobatake and K. Tanaka, \emph{Neuronal selectivities to complex object features in the ventral visual pathway of the macaque cerebral cortex}, J. Neurophysiol. \textbf{71}, 856 (1994).
\bibitem{Albright1984} T. D. Albright, \emph{Direction and orientation selectivity of neurons in visual area MT of the macaque}, J. Neurophysiol. \textbf{52}, 1106 (1984).
\bibitem{Goodfellow2015} I. J. Goodfellow, J. Shlens, and C. Szegedy, \emph{Explaining and harnessing adversarial examples}, in \emph{Proc. ICLR} (2015), arXiv:1412.6572.
\bibitem{Elsayed2018} G. F. Elsayed et al., \emph{Adversarial examples that fool both computer vision and time-limited humans}, in \emph{Advances in Neural Information Processing Systems 31} (2018), arXiv:1802.08195.
\bibitem{Thompson1982} P. Thompson, \emph{Perceived rate of movement depends on contrast}, Vision Res. \textbf{22}, 377 (1982).
\bibitem{BarlowHill1963} H. B. Barlow and R. M. Hill, \emph{Evidence for a physiological explanation of the waterfall phenomenon and figural after-effects}, Nature \textbf{200}, 1345 (1963).
\bibitem{EaglemanSejnowski2000} D. M. Eagleman and T. J. Sejnowski, \emph{Motion integration and postdiction in visual awareness}, Science \textbf{287}, 2036 (2000).

\bibitem{Feinstein1955mu} B. Feinstein, B. Lindeg{\aa}rd, E. Nyman, and G. Wohlfart, \emph{Morphologic studies of motor units in normal human muscles}, Acta Anat. \textbf{23}, 127 (1955).
\bibitem{MilnerBrown1973rec} H. S. Milner-Brown, R. B. Stein, and R. Yemm, \emph{The orderly recruitment of human motor units during voluntary isometric contractions}, J. Physiol. \textbf{230}, 359 (1973).
\bibitem{Jones2002sdn} K. E. Jones, A. F. de C. Hamilton, and D. M. Wolpert, \emph{Sources of signal-dependent noise during isometric force production}, J. Neurophysiol. \textbf{88}, 1533 (2002).
\bibitem{GomiOsu1998stiff} H. Gomi and R. Osu, \emph{Task-dependent viscoelasticity of human multijoint arm and its spatial characteristics for interaction with environments}, J. Neurosci. \textbf{18}, 8965 (1998).
\bibitem{Jankowska1972Ia} E. Jankowska and W. J. Roberts, \emph{Synaptic actions of single interneurones mediating reciprocal Ia inhibition of motoneurones}, J. Physiol. \textbf{222}, 623 (1972).
\bibitem{Pruszynski2008rapid} J. A. Pruszynski, I. Kurtzer, and S. H. Scott, \emph{Rapid motor responses are appropriately tuned to the metrics of a visuospatial task}, J. Neurophysiol. \textbf{100}, 224 (2008).
\bibitem{SaundersKnill2003vis} J. A. Saunders and D. C. Knill, \emph{Humans use continuous visual feedback from the hand to control fast reaching movements}, Exp. Brain Res. \textbf{152}, 341 (2003).
\bibitem{FlanaganWing1997grip} J. R. Flanagan and A. M. Wing, \emph{The role of internal models in motion planning and control: evidence from grip force adjustments during movements of hand-held loads}, J. Neurosci. \textbf{17}, 1519 (1997).
\bibitem{Matsuoka1985osc} K. Matsuoka, \emph{Sustained oscillations generated by mutually inhibiting neurons with adaptation}, Biol. Cybern. \textbf{52}, 367 (1985).

\bibitem{Treede1995heat} R.-D. Treede, R. A. Meyer, S. N. Raja, and J. N. Campbell, \emph{Evidence for two different heat transduction mechanisms in nociceptive primary afferents innervating monkey skin}, J. Physiol. \textbf{483}, 747 (1995).
\bibitem{Weidner1999cfib} C. Weidner et al., \emph{Functional attributes discriminating mechano-insensitive and mechano-responsive C nociceptors in human skin}, J. Neurosci. \textbf{19}, 10184 (1999).
\bibitem{CampbellLaMotte1983first} J. N. Campbell and R. H. LaMotte, \emph{Latency to detection of first pain}, Brain Res. \textbf{266}, 203 (1983).
\bibitem{LaMotte1982hyper} R. H. LaMotte, J. G. Thalhammer, H. E. Torebj{\"o}rk, and C. J. Robinson, \emph{Peripheral neural mechanisms of cutaneous hyperalgesia following mild injury by heat}, J. Neurosci. \textbf{2}, 765 (1982).
\bibitem{Mancini2014touch} F. Mancini, T. Nash, G. D. Iannetti, and P. Haggard, \emph{Pain relief by touch: a quantitative approach}, Pain \textbf{155}, 635 (2014).
\bibitem{Fields1983rvm} H. L. Fields, J. Bry, I. Hentall, and G. Zorman, \emph{The activity of neurons in the rostral medulla of the rat during withdrawal from noxious heat}, J. Neurosci. \textbf{3}, 2545 (1983).
\bibitem{Crombez1996disrupt} G. Crombez, C. Eccleston, F. Baeyens, and P. Eelen, \emph{The disruptive nature of pain: an experimental investigation}, Behav. Res. Ther. \textbf{34}, 911 (1996).
\bibitem{Schouenborg1995wr} J. Schouenborg, H.-R. Weng, J. Kalliom{\"a}ki, and H. Holmberg, \emph{A survey of spinal dorsal horn neurones encoding the spatial organization of withdrawal reflexes in the rat}, Exp. Brain Res. \textbf{106}, 19 (1995).

\bibitem{Andersen1992cereb} B. B. Andersen, L. Korbo, and B. Pakkenberg, \emph{A quantitative study of the human cerebellum with unbiased stereological techniques}, J. Comp. Neurol. \textbf{326}, 549 (1992).
\bibitem{ShadmehrMussaIvaldi1994} R. Shadmehr and F. A. Mussa-Ivaldi, \emph{Adaptive representation of dynamics during learning of a motor task}, J. Neurosci. \textbf{14}, 3208 (1994).

\bibitem{JoseCollison1970ihr} A. D. Jose and D. Collison, \emph{The normal range and determinants of the intrinsic heart rate in man}, Cardiovasc. Res. \textbf{4}, 160 (1970).
\bibitem{Logothetis1987gbg} D. E. Logothetis, Y. Kurachi, J. Galper, E. J. Neer, and D. E. Clapham, \emph{The $\beta\gamma$ subunits of GTP-binding proteins activate the muscarinic K$^+$ channel in heart}, Nature \textbf{325}, 321 (1987).
\bibitem{KellerWood1984fb} M. E. Keller-Wood and M. F. Dallman, \emph{Corticosteroid inhibition of ACTH secretion}, Endocr. Rev. \textbf{5}, 1 (1984).
\bibitem{Walker2012glc} J. J. Walker et al., \emph{The origin of glucocorticoid hormone oscillations}, PLoS Biol. \textbf{10}, e1001341 (2012).
\bibitem{Wildt1981gnrh} L. Wildt et al., \emph{Frequency and amplitude of gonadotropin-releasing hormone stimulation and gonadotropin secretion in the rhesus monkey}, Endocrinology \textbf{109}, 376 (1981).
\bibitem{Welsh2010scn} D. K. Welsh, J. S. Takahashi, and S. A. Kay, \emph{Suprachiasmatic nucleus: cell autonomy and network properties}, Annu. Rev. Physiol. \textbf{72}, 551 (2010).
\bibitem{Khalsa2003prc} S. B. S. Khalsa, M. E. Jewett, C. Cajochen, and C. A. Czeisler, \emph{A phase response curve to single bright light pulses in human subjects}, J. Physiol. \textbf{549}, 945 (2003).
\bibitem{Sack1992blind} R. L. Sack, A. J. Lewy, M. L. Blood, L. D. Keith, and H. Nakagawa, \emph{Circadian rhythm abnormalities in totally blind people: incidence and clinical significance}, J. Clin. Endocrinol. Metab. \textbf{75}, 127 (1992).

\bibitem{Henze2000extra} D. A. Henze et al., \emph{Intracellular features predicted by extracellular recordings in the hippocampus in vivo}, J. Neurophysiol. \textbf{84}, 390 (2000).
\bibitem{Histed2009stim} M. H. Histed, V. Bonin, and R. C. Reid, \emph{Direct activation of sparse, distributed populations of cortical neurons by electrical microstimulation}, Neuron \textbf{63}, 508 (2009).
\bibitem{Orsborn2014coad} A. L. Orsborn et al., \emph{Closed-loop decoder adaptation shapes neural plasticity for skillful neuroprosthetic control}, Neuron \textbf{82}, 1380 (2014).

\bibitem{HuVervaekeStorm2002} H. Hu, K. Vervaeke, and J. F. Storm, \emph{Two forms of electrical resonance at theta frequencies, generated by M-current, h-current and persistent Na$^+$ current in rat hippocampal pyramidal cells}, J. Physiol. \textbf{545}, 783 (2002).
\bibitem{Mauro1970} A. Mauro, F. Conti, F. Dodge, and R. Schor, \emph{Subthreshold behavior and phenomenological impedance of the squid giant axon}, J. Gen. Physiol. \textbf{55}, 497 (1970).
\bibitem{Aksay2001} E. Aksay et al., \emph{In vivo intracellular recording and perturbation of persistent activity in a neural integrator}, Nat. Neurosci. \textbf{4}, 184 (2001).
\bibitem{Major2004} G. Major et al., \emph{Plasticity and tuning of the time course of analog persistent firing in a neural integrator}, Proc. Natl. Acad. Sci. USA \textbf{101}, 7745 (2004).
\bibitem{Tateno2004} T. Tateno, A. Harsch, and H. P. C. Robinson, \emph{Threshold firing frequency-current relationships of neurons in rat somatosensory cortex: type 1 and type 2 dynamics}, J. Neurophysiol. \textbf{92}, 2283 (2004).
\bibitem{Marder2012} E. Marder, \emph{Neuromodulation of neuronal circuits: back to the future}, Neuron \textbf{76}, 1 (2012).
\bibitem{McGintyHarper1976} D. J. McGinty and R. M. Harper, \emph{Dorsal raphe neurons: depression of firing during sleep in cats}, Brain Res. \textbf{101}, 569 (1976).

\bibitem{Gentet2000} L. J. Gentet, G. J. Stuart, and J. D. Clements, \emph{Direct measurement of specific membrane capacitance in neurons}, Biophys. J. \textbf{79}, 314 (2000).
\bibitem{Borst1995} J. G. G. Borst, F. Helmchen, and B. Sakmann, \emph{Pre- and postsynaptic whole-cell recordings in the medial nucleus of the trapezoid body of the rat}, J. Physiol. \textbf{489}, 825 (1995).
\bibitem{AmirDevor2003} R. Amir and M. Devor, \emph{Electrical excitability of the soma of sensory neurons is required for spike invasion of the soma, but not for through-conduction}, Biophys. J. \textbf{84}, 2181 (2003).
\bibitem{Chadderton2004} P. Chadderton, T. W. Margrie, and M. H\"ausser, \emph{Integration of quanta in cerebellar granule cells during sensory processing}, Nature \textbf{428}, 856 (2004).
\bibitem{LitwinKumar2017} A. Litwin-Kumar et al., \emph{Optimal degrees of synaptic connectivity}, Neuron \textbf{93}, 1153 (2017).
\bibitem{NapperHarvey1988} R. M. A. Napper and R. J. Harvey, \emph{Number of parallel fiber synapses on an individual Purkinje cell in the cerebellum of the rat}, J. Comp. Neurol. \textbf{274}, 168 (1988).
\bibitem{Megias2001} M. Meg\'{\i}as, Z. Emri, T. F. Freund, and A. I. Guly\'as, \emph{Total number and distribution of inhibitory and excitatory synapses on hippocampal CA1 pyramidal cells}, Neuroscience \textbf{102}, 527 (2001).
\bibitem{Polsky2004} A. Polsky, B. W. Mel, and J. Schiller, \emph{Computational subunits in thin dendrites of pyramidal cells}, Nat. Neurosci. \textbf{7}, 621 (2004).
\bibitem{Brunel2004} N. Brunel et al., \emph{Optimal information storage and the distribution of synaptic weights: perceptron versus Purkinje cell}, Neuron \textbf{43}, 745 (2004).

\bibitem{Vyazovskiy2009} V. V. Vyazovskiy et al., \emph{Cortical firing and sleep homeostasis}, Neuron \textbf{63}, 865 (2009).
\bibitem{Daan1984} S. Daan, D. G. Beersma, and A. A. Borb\'ely, \emph{Timing of human sleep: recovery process gated by a circadian pacemaker}, Am. J. Physiol. \textbf{246}, R161 (1984).
\bibitem{Borbely1981} A. A. Borb\'ely et al., \emph{Sleep deprivation: effect on sleep stages and EEG power density in man}, Electroencephalogr. Clin. Neurophysiol. \textbf{51}, 483 (1981).
\bibitem{SteriadeAmzica1993} M. Steriade, F. Amzica, and A. Nu\~nez, \emph{Cholinergic and noradrenergic modulation of the slow ($\approx0.3$ Hz) oscillation in neocortical cells}, J. Neurophysiol. \textbf{70}, 1385 (1993).
\bibitem{vonKrosigk1993} M. von Krosigk, T. Bal, and D. A. McCormick, \emph{Cellular mechanisms of a synchronized oscillation in the thalamus}, Science \textbf{261}, 361 (1993).
\bibitem{LeeWilson2002} A. K. Lee and M. A. Wilson, \emph{Memory of sequential experience in the hippocampus during slow wave sleep}, Neuron \textbf{36}, 1183 (2002).
\bibitem{Euston2007} D. R. Euston, M. Tatsuno, and B. L. McNaughton, \emph{Fast-forward playback of recent memory sequences in prefrontal cortex during sleep}, Science \textbf{318}, 1147 (2007).
\bibitem{Hobson2009} J. A. Hobson, \emph{REM sleep and dreaming: towards a theory of protoconsciousness}, Nat. Rev. Neurosci. \textbf{10}, 803 (2009).

\bibitem{CohenSegal2011} D. Cohen and M. Segal, \emph{Network bursts in hippocampal microcultures are terminated by exhaustion of vesicle pools}, J. Neurophysiol. \textbf{106}, 2314 (2011).
\bibitem{Tsodyks2000} M. Tsodyks, A. Uziel, and H. Markram, \emph{Synchrony generation in recurrent networks with frequency-dependent synapses}, J. Neurosci. \textbf{20}, RC50 (2000).
\bibitem{Benson1994} D. L. Benson, F. H. Watkins, O. Steward, and G. Banker, \emph{Characterization of GABAergic neurons in hippocampal cell cultures}, J. Neurocytol. \textbf{23}, 279 (1994).

\input{supplement/bib_22}
\end{thebibliography}


\end{document}
